\documentclass[11pt,a4paper]{article}
\usepackage[T1]{fontenc}
\usepackage[utf8]{inputenc}
\usepackage{lmodern,amsmath,amssymb,amsthm,mathtools,mathrsfs}
\usepackage[margin=25mm,headheight=15pt]{geometry}
\usepackage{microtype,booktabs,longtable,array,enumitem,xcolor,fancyhdr}
\usepackage{listings,fancyvrb,needspace}
\usepackage[hyphens]{url}
\usepackage{hyperref}
\definecolor{inkblue}{RGB}{27,55,83}
\definecolor{ink}{HTML}{163E4A}
\definecolor{muted}{HTML}{52616B}
\hypersetup{colorlinks=true,linkcolor=inkblue,citecolor=inkblue,urlcolor=inkblue,
 bookmarksdepth=3,pdftitle={Euler's Constant: Arithmetic, Monodromy, and Gamma-Related Structures},
 pdfauthor={Compiled for KEIO},pdfsubject={Complete eight-manuscript collection; banana volume excluded}}
\newtheorem{theorem}{Theorem}[section]
\newtheorem{lemma}[theorem]{Lemma}
\newtheorem{proposition}[theorem]{Proposition}
\newtheorem{corollary}[theorem]{Corollary}
\theoremstyle{definition}
\newtheorem{definition}[theorem]{Definition}
\newtheorem{example}[theorem]{Example}
\newtheorem{observation}[theorem]{Observation}
\theoremstyle{remark}
\newtheorem{remark}[theorem]{Remark}
\newcommand{\Q}{\mathbb Q}
\newcommand{\Z}{\mathbb Z}
\newcommand{\N}{\mathbb N}
\newcommand{\C}{\mathbb C}
\newcommand{\R}{\mathbb R}
\newcommand{\Qbar}{\overline{\mathbb Q}}
\newcommand{\dd}{\,\mathrm d}
\newcommand{\bG}{\mathbf G}
\newcommand{\bE}{\mathbf E}
\newcommand{\bD}{\mathbf D}
\DeclareMathOperator{\Ein}{Ein}
\DeclareMathOperator{\Var}{Var}
\DeclareMathOperator{\Log}{Log}
\DeclareMathOperator{\Res}{Res}
\DeclareMathOperator{\den}{den}
\DeclareMathOperator{\lc}{lc}
\DeclareMathOperator{\lcm}{lcm}
\DeclareMathOperator{\Frac}{Frac}
\DeclareMathOperator{\ord}{ord}
\DeclareMathOperator{\diag}{diag}
\DeclareMathOperator{\adj}{adj}
\numberwithin{equation}{section}
\renewcommand{\thesection}{\Roman{part}.\arabic{section}}
\renewcommand{\theHsection}{\arabic{part}.main.\arabic{section}}
\renewcommand{\theHsubsection}{\theHsection.\arabic{subsection}}
\renewcommand{\theHsubsubsection}{\theHsubsection.\arabic{subsubsection}}
\renewcommand{\theHequation}{\theHsection.\arabic{equation}}
\providecommand{\theHtheorem}{}
\renewcommand{\theHtheorem}{\theHsection.\arabic{theorem}}
\newcommand{\PartAppendix}{%
 \setcounter{section}{0}\setcounter{subsection}{0}%
 \renewcommand{\thesection}{\Roman{part}.\Alph{section}}%
 \renewcommand{\theHsection}{\arabic{part}.appendix.\arabic{section}}}
\setlength{\emergencystretch}{3em}
\setlength{\parskip}{2pt}
\setlist{itemsep=2pt,topsep=4pt}
\setcounter{tocdepth}{1}
\makeatletter
\renewcommand{\l@section}{\@dottedtocline{1}{0em}{4.2em}}
\makeatother
\pagestyle{fancy}\fancyhf{}
\fancyhead[L]{\small Euler's constant: complete research collection}
\fancyhead[R]{\small KEIO}
\fancyfoot[C]{\thepage}
\renewcommand{\headrulewidth}{0.3pt}
\lstset{language=Python,basicstyle=\ttfamily\scriptsize,
 breaklines=true,columns=fullflexible,keepspaces=true,showstringspaces=false,
 frame=single,rulecolor=\color{black!20},aboveskip=8pt,belowskip=8pt}
\title{\bfseries Euler's Constant\\[5pt]
 \Large Arithmetic, Monodromy,\\and Gamma-Related Structures\\[10pt]
 \large Complete integrated research collection}
\author{Compiled for KEIO\\\small Eight manuscripts, with source provenance and original proof status}
\date{19 October 2026}
\begin{document}
\maketitle
\begin{abstract}
This single-source volume collects all six Euler/Gamma research manuscripts
produced in the present conversation and the two newly supplied manuscripts.
It preserves their mathematical statements, proofs, failed attempts,
limitations, exact checks, appendices, and bibliographies. The separate
61-page Bessel--periods--tableaux--banana volume is excluded in full.
The material covers asymptotic coefficient arithmetic, special-function
inequalities, stochastic stopping boundaries, conservative matrix fields,
moment determinants, differential and difference algebra, monodromy,
primitive lattices, rational positivity certificates, logarithmic Meijer
identities, complete-monotonicity degrees, and Dickman-related asymptotics.
The editorial integration provides one pagination, one contents list,
part-qualified cross-references, and a source-to-part concordance.
It is a content-preserving compilation, not a new claim of arithmetic
irrationality or transcendence. Source-level conjecture counts are not added
together: several results overlap, some are restricted cases, and some are
historical or classical consequences. No proof of the irrationality or
transcendence of Euler's constant is claimed.
\end{abstract}
\bigskip
\noindent\textbf{Edition.} Original source versions dated 18--19 October 2026.
Internal verification and scope notes are retained; external peer review and
comprehensive priority certification are not asserted.
\clearpage
\phantomsection\addcontentsline{toc}{section}{Editorial scope and source concordance}
\section*{Editorial scope and source concordance}
The phrase ``this conversation'' is interpreted as the entire conversation,
not only the immediately preceding 29-page manuscript. The six generated
Euler/Gamma manuscripts are distinct continuations, not interchangeable
versions. In particular, \emph{Arithmetic Transfers II} does not contain or
replace \emph{Arithmetic Transfers Around Euler's Constant}. Both are
included, together with the scalar-convergence, monodromy,
rational-certificate, and fractional-degree manuscripts.

The two attached TeX manuscripts are included in full as Parts IV and VII.
The attached \emph{Twelve further attempts} version is byte-identical to
the earlier attachment used in the research; it is included once.
The source PDFs total 236 pages. Re-typesetting and the shared front matter
change the final page count without defining the scope of the contents.

The 61-page \emph{Bessel Periods, Tableaux, and Banana} integrated volume is
excluded as a whole. No chapter or appendix is extracted from that volume.
The Bessel stopping results and inverse Bessel-polynomial recurrence already
contained in the eight Euler/Gamma manuscripts remain included: they belong
to the requested sources and are not chapters of the excluded banana volume.

\begin{center}\small
\begin{tabular}{@{}p{.06\textwidth}p{.63\textwidth}p{.10\textwidth}p{.12\textwidth}@{}}
\toprule
Part & Source manuscript / topic & Source pages & This volume\\
\midrule
I & Arithmetic transfers around Euler's constant & 36 & \pageref{part1:start}--\pageref{part1:end}\\[5pt]
II & Arithmetic transfers II: units, parameter jets, and source assertions & 35 & \pageref{part2:start}--\pageref{part2:end}\\[5pt]
III & Scalar convergence, moment problems, and differential algebra & 37 & \pageref{part3:start}--\pageref{part3:end}\\[5pt]
IV & Twelve further irrationality attempts and their failure points & 14 & \pageref{part4:start}--\pageref{part4:end}\\[5pt]
V & Monodromy, primitive lattices, and positive rational forms & 30 & \pageref{part5:start}--\pageref{part5:end}\\[5pt]
VI & Rational certificates, ramification, and arithmetic obstructions & 27 & \pageref{part6:start}--\pageref{part6:end}\\[5pt]
VII & A logarithmic Meijer identity and exact arithmetic structure & 28 & \pageref{part7:start}--\pageref{part7:end}\\[5pt]
VIII & Fractional degrees and Dickman-related asymptotics & 29 & \pageref{part8:start}--\pageref{part8:end}\\[5pt]
\bottomrule
\end{tabular}
\end{center}

\subsection*{How to read the integration}
Statements and proofs keep their source wording. References to ``this paper'',
``this manuscript'', or ``earlier manuscripts'' inside a part retain their
source-local meaning. Internal labels and citations are qualified by part;
for example, a reference prefixed VI belongs to Part VI. Bibliographies
remain with their respective parts so that the original citation scope is
preserved. Abstracts and preparation credits are retained; duplicate title
pages and repeated local contents lists are replaced by this shared front
matter. Reproducibility code is embedded directly in this one TeX file.

Some arguments recur in later parts in a stronger form or under different
hypotheses. They are retained because dropping them would also drop their
local proof context. They are not counted as independent new resolutions.
The limits and corrections in later parts must be read together with the
earlier statements; combining manuscripts does not upgrade a conditional,
restricted, or unsuccessful argument into an unconditional arithmetic proof.

\subsection*{Version-sensitive correction note}
Part VII's discussion of corrections to the supplied attempts was written
against earlier formulations. The Part IV attachment included here already
states the nonzero target-constant restriction, the change in a recurrence
under a Gamma twist, and the qualification about modification of the final
segment of a continuation path. Thus the historical correction discussion
does not assert that those particular qualifications are absent from the
current Part IV version. Its original reasoning is retained for context.

The broader ``discarded'' verdict about arithmetic holonomy in Part IV is
preserved as an assessment of that attempted route, not endorsed as a
universal impossibility theorem. Proposition~\ref{p4:prop:cancel} has an
identifiability hypothesis. The affine cancellation-locus analysis and
the logical boundary explained in Section~\ref{p6:sec:holonomy-scope},
together with Part VII's correction discussion, delimit its valid use.
In particular, an obstruction to an unconditional cancellation formula
does not by itself exclude a contradiction constructed under an
algebraicity hypothesis.

\subsection*{Proof status and scope of this edition}
This task is editorial integration. It does not add a new global theorem,
rerun a comprehensive literature search, or certify every earlier claim
anew. The proofs and source audits have the status stated in their own parts.
In particular, source-level ledgers containing ten entries are not a claim
of ten independent presently open major conjectures newly solved, either
within that part or across this volume. The arithmetic nature of
$\gamma$ and of the Euler--Gompertz constant remains unresolved by these
manuscripts.
\clearpage
\tableofcontents
\clearpage
\clearpage
\setcounter{section}{0}\setcounter{subsection}{0}
\renewcommand{\thesection}{\Roman{part}.\arabic{section}}
\renewcommand{\theHsection}{\arabic{part}.main.\arabic{section}}
\renewcommand{\refname}{References for Part I}
\fancyhead[L]{\small Euler\textquotesingle s constant: Part I}
\part{Arithmetic transfers around Euler's constant}
\label{part1:start}
\noindent\textbf{Source:} {\small\path{KEIO_Euler_Arithmetic_Transfers_20261018.tex}}.

Original PDF: 36 pages. Mathematical content and source-level qualifications are retained.


% BEGIN PRESERVED SOURCE BODY 1
\begin{abstract}
We investigate transfers from the Euler--Gompertz and $E$-function setting to
asymptotic arithmetic, orthogonal polynomials, special-function inequalities,
and optimal stopping. For the coefficients $d_k$ in the
Brent--Glasser--Guttmann expansion, we obtain the exact reduced denominator
$2^{6k-s_2(k)}$, where $s_2(k)$ is the binary digit sum, strengthening the
recovered integrality result. We prove both all-order coefficient sign
observations in that paper, using explicit analytic tail bounds and exact
finite certificates. A radial-oscillator transfer also gives exact
Gevrey order two and the Borel radius. We prove divergence of the factorially cleared
Euler--Gompertz approximants discussed by Murabayashi--Yoshida and a uniform
obstruction for all nonnegative integer-polynomial filters under that
normalization. A polynomial--Gamma monotonicity question receives sharp
answers with the coefficient assumptions made explicit. For rational dimension
and reward, Bessel-bridge stopping boundaries are algebraic exactly on the
lines $q-d\in2\mathbb Z_{\ge0}$; an exact parameter sensitivity recovers
$\Ein(1)=\gamma+\delta/e$. Every finite dimension jet of the quadratic
optimal value at dimension two is algebraically independent. Further results concern gamma-cancelling linear
forms and the limitations of moment positivity and Gevrey zero removal.
The manuscript distinguishes recovered results, extensions, and unproved
targets. It does not prove irrationality or transcendence of $\gamma$ and does
not establish ten new resolutions of published open problems.
\end{abstract}
\noindent\textbf{Keywords:} Euler--Mascheroni constant; Euler--Gompertz constant;
dyadic denominator; Kummer function; optimal stopping; positive polynomial;
arithmetic Gevrey series.


\section{Scope, source ledger, and the arithmetic question}
We use
\[
 \gamma=\lim_{n\to\infty}(H_n-\log n),\qquad
 \delta=\int_0^\infty\frac{e^{-t}}{1+t}\,dt,\qquad
 \eta=\Ein(1)=\gamma+\delta/e.
\]
The intended research target was ten genuinely new solutions to previously
published open questions. That target is \emph{not met by this manuscript}.
The number of theorem statements below must not be mistaken for that target
count. Parameter instances, a stronger form of an already recovered result,
and a new problem formulated during the investigation are not independent
published conjectures.

The published statements examined are listed below. ``Proved here'' refers to
the argument supplied in this manuscript, not to completed external peer review
or established historical priority. The older project records
\cite{p1:priorBGG,p1:priorTransfer} already contain BGG integrality and its mod-$32$
refinement; both are identified as recovered, not first discoveries of this
research round. An unnumbered numerical observation is kept separate from a
numbered conjecture.

\small
\begin{longtable}{@{}p{.065\textwidth}p{.48\textwidth}p{.385\textwidth}@{}}
\toprule ID&Published target and location&Outcome in this manuscript\\\midrule
\endfirsthead
\toprule ID&Published target and location&Outcome in this manuscript\\\midrule
\endhead
A1&BGG \cite{p1:BrentGlasserGuttmann2019}, Conjecture 10: $64^kd_k\in\Z$.&Recovered proof; strengthened to an exact reduced-denominator formula.\\
A2&BGG, Remark 11: $64^kd_k\equiv\binom{2k}{k}\pmod{32}$ for $k\ge3$.&Recovered all-order proof; not a new discovery in this round.\\
A3&BGG, Remark 11: $d_k<0$ for every $k\ge3$.&Proved for all indices by an analytic tail bound and an exact finite certificate.\\
A4&BGG, Remark 8: $c_m<0$ except $m=0,4$.&Proved for all indices by a second dominance argument and an exact finite certificate.\\
A5&Murabayashi--Yoshida \cite{p1:MurabayashiYoshida2024}, before Theorem 7.1: factorially cleared error tends to infinity.&Proved with explicit bounds. An unnumbered prediction, not Conjecture 7.5.\\
A6&Same paper, Conjecture 7.5: monotone fractional-part subsequences.&Not proved; the needed direction of monotonicity is not supplied by the monotone-subsequence theorem.\\
A7&Bulboac\u a--Zayed \cite{p1:bz2026}, Remark 3(2): polynomial-numerator shifts.&Fixed-polynomial answer and sharp nonnegative-coefficient answer; unrestricted degree-only version is false. One source question.\\
A8&Same remark: sharp constant-numerator shift.&Reproved; a prior public claim exists \cite{p1:pudim2026}, so excluded from new resolutions.\\
A9&Lagarias \cite{p1:Lagarias2013}, Conjecture 1.0.1: $\gamma\notin\Q$.&Not proved.\\
A10&Lagarias, Conjecture 1.0.2: $\gamma$ is not a Kontsevich--Zagier period.&Not proved, nor is its transcendence consequence.\\
A11&Fischler--Rivoal \cite{p1:FischlerRivoal2024}, Conjecture 1: $\mathcal E\cap\mathcal D=\Qbar$.&Not proved.\\
A12&Same paper, Conjecture 2: pure Gevrey zero removal.&Not proved.\\
A13&Same paper, Conjecture 3: mixed Gevrey zero removal.&Not proved; an explicit restriction is equivalent to the arithmetic target.\\
A14&Diamond--Ford \cite{p1:DiamondFord2008}: lower edge of the generalized Euler spectrum.&Not proved; the stated equality is equivalent to RH.\\
A15&Ford--Luca--Moree \cite{p1:FordLucaMoree2014}, Conjecture 1: cyclotomic Euler--Kronecker limit points.&Not proved.\\
A16&Dixit--Murty \cite{p1:DixitMurty2023}, introduction: infinitely many negative cyclotomic Euler--Kronecker constants.&Not proved unconditionally here.\\
\bottomrule
\end{longtable}
\normalsize

Four source-level targets receive proofs in this round: the two BGG sign
observations (A3--A4), the unnumbered divergence observation (A5), and the
polynomial shift question (A7), with its interpretations distinguished.
This is neither ten new conjecture resolutions nor a claim of verified
historical priority. The bounded literature search found no earlier proof
of the two sign assertions; that search is not an exhaustive priority audit.

The polynomial question A7 lacks a coefficient-sign specification in its
original remark. Adding nonnegativity silently would change the problem; the
three interpretations are therefore treated separately. A5 is a modest
analytic observation whose proof clarifies the denominator obstruction. Neither
should be advertised as an arithmetic breakthrough for $\gamma$.

Two useful boundaries organize the investigation. A result about the
\emph{function} $\Gamma$ is not automatically a result about the arithmetic
of its logarithmic derivative at one. Likewise, an algebraic-independence
theorem for finite $E$-function values does not automatically pass to the
singular limits in which $\gamma$ appears. The following sections make these
distinctions mathematically explicit while retaining the transfers that do work.


\section{Exact arithmetic of the BGG product coefficients}
\label{p1:sec:bgg-arithmetic}

Let $M(a,b,x)={}_1F_1(a;b;x)$ and let $U(a,b,x)$ be Tricomi's
confluent hypergeometric function. Brent--Glasser--Guttmann
\cite{p1:BrentGlasserGuttmann2019} define Taylor coefficients $a_n,b_n$ by
\[
 e^{z/(1-z)}=\sum_{n\ge0}a_nz^n,\qquad
 e^{1/(1-z)}E_1\!\left(\frac1{1-z}\right)=\sum_{n\ge0}b_nz^n.
\]
For $n\ge1$ their exact formulas are
\[
 a_n=e^{-1}M(n+1,2,1),\quad b_n=-\Gamma(n)U(n,0,1),\quad
 \rho_n=a_nb_n.
\]
Their product has the asymptotic expansion
\[
 -2n^{3/2}\rho_n\sim\sum_{k\ge0}d_kn^{-k},\qquad d_0=1.
\]
The integrality of $r_k=64^kd_k$ is their Conjecture~10. Remark~11
records both a mod-$32$ pattern and the negative signs from $k=3$ onward.
We first recover the integrality mechanism and sharpen it to an exact
denominator formula. The sign requires the independent analytic argument
of Section~\ref{p1:sec:bgg-sign}.

\subsection{An exact differential identity and its normalization}
For $x>0$ and $n>0$, set
\[
 \begin{gathered}
 u=e^{-x/2}xM(n+1,2,x),\qquad v=e^{-x/2}U(n,0,x),\\
 H_n=\Gamma(n+1)uv,\qquad F(n,x)=2\sqrt{n/x}\,H_n(x).
 \end{gathered}
\]
Both $u,v$ satisfy $y''=(n/x+1/4)y$. At zero, $u\sim x$ and
$v\to1/\Gamma(n+1)$, so their constant Wronskian is
$-1/\Gamma(n+1)$. The logarithmic term allowed in $v'$ contributes
$O(x\log x)$ to this limit and vanishes. Direct multiplication gives
\[
 (H_n')^2-2H_nH_n''+4\left(\frac nx+\frac14\right)H_n^2=1.
\]
After inserting the definition of $F$, this becomes
\begin{equation}\label{p1:eq:bgg-exact-F}
 F^2+\frac1{4n}\left[
 \left(x+\frac3{4x}\right)F^2-FF'
 +x\bigl((F')^2-2FF''\bigr)\right]=1.
\end{equation}
At $x=1$, $H_n(1)=-n\rho_n$ and $F(n,1)=-2n^{3/2}\rho_n$.

The heat-kernel representation proved in
Section~\ref{p1:sec:bgg-sign} supplies
$F(n,x)\sim\sum_{k\ge0}d_k(x)n^{-k}$, $d_0(x)=1$, locally uniformly
for $x>0$, together with any fixed number of $x$-derivatives.
Indeed the Bessel expansion in the explicit kernel can be truncated
uniformly on compact positive $x$-intervals and differentiated a fixed
number of times; Watson's lemma then applies to the differentiated
Laplace integrals. This justifies coefficient comparison in
\eqref{p1:eq:bgg-exact-F}; no differentiation of a merely pointwise
asymptotic assertion is being assumed. In particular $d_k(1)=d_k$.
Define the formal series
\[
 R(t,x)=\sum_{k\ge0}r_k(x)t^k,\qquad r_k(x)=64^kd_k(x).
\]
Substitution of $n^{-1}=64t$ into \eqref{p1:eq:bgg-exact-F} gives
\begin{equation}\label{p1:eq:bgg-r-scaled}
 R^2=1+32txRR_{xx}+16tRR_x-16txR_x^2
 -(16x+12x^{-1})tR^2.
\end{equation}
Its constant term is $R(0,x)=1$, and its coefficient of degree $k\ge1$
determines $2r_k$ from preceding coefficients. Thus it has a unique
formal solution over $\Q[x,x^{-1}]$.


\subsection{Exact dyadic denominators and exclusion of integral dyadic roots}
\label{p1:subsec:bgg-exact-dyadic}

Write
\[
 R(x,t)=\sum_{n\geq0}r_n(x)t^n,
 \qquad r_n(x)=64^n d_n(x),\qquad r_0(x)=1.
\]
Throughout this subsection, a prime denotes differentiation with respect to
$x$.  We use the scaled identity \eqref{p1:eq:bgg-r-scaled} as an identity of
formal power series; no convergence of the asymptotic expansion is required.
For a nonnegative integer $n$, let $s_2(n)$ denote the sum of its binary
digits.  The notation $v_2$ means the additive valuation normalized by
$v_2(2)=1$, also on $\overline{\mathbb Q}_2$.

\begin{theorem}[Exact dyadic valuation]
\label{p1:thm:bgg-exact-dyadic}
For every $n\geq0$, define
\[
 A_n(x)=\frac{n!}{2^n}\,r_n(x).
\]
Then
\begin{equation}
 \label{p1:eq:bgg-hurwitz-parity}
 A_n(x)\in\mathbb Z[x,x^{-1}],\qquad
 A_n(x)\equiv x^{-n}\pmod {2\mathbb Z[x,x^{-1}]}.
\end{equation}
More precisely,
\begin{equation}
 \label{p1:eq:bgg-polynomial-parity}
 P_n(x):=x^n A_n(x)\in\mathbb Z[x^2],
 \qquad P_n(x)\equiv1\pmod {2\mathbb Z[x^2]}.
\end{equation}
In particular, for every $n\geq1$,
\begin{equation}
 \label{p1:eq:bgg-exact-denominator}
 v_2\bigl(r_n(1)\bigr)=s_2(n),\qquad
 \operatorname{den}\bigl(d_n(1)\bigr)=2^{6n-s_2(n)},
\end{equation}
where $\operatorname{den}$ denotes the positive denominator in lowest terms.
\end{theorem}

\begin{proof}
Put $V(x)=16x+12/x$.  Extracting the coefficient of $t^n$ from
\eqref{p1:eq:bgg-r-scaled} gives, for $n\geq1$,
\begin{align}
 2r_n={}&
 \sum_{i+j=n-1}
 \bigl(32x r_i r_j''+16r_i r_j'
             -16x r_i'r_j'-V r_i r_j\bigr)
 \notag\\
 &\hspace{8mm}-\sum_{i=1}^{n-1}r_i r_{n-i}.
 \label{p1:eq:bgg-r-coefficient}
\end{align}
This triangular recurrence first shows that
\begin{equation}
 \label{p1:eq:bgg-r-even}
 r_n\in2\mathbb Z[x,x^{-1}]\qquad(n\geq1).
\end{equation}
Indeed, all coefficients in the first sum of
\eqref{p1:eq:bgg-r-coefficient} are divisible by $4$; the second sum has the
same divisibility if the preceding positive-index $r_i$ are even.  Dividing
by $2$ completes the induction, including the base case with an empty
second sum.  The same recurrence shows inductively that the exponents
occurring in $r_n$ belong to
\[
 \{-n,-n+2,\ldots,n-2,n\}.
\]

After substituting $r_i=2^i A_i/i!$ and multiplying
\eqref{p1:eq:bgg-r-coefficient} by $n!/2^{n+1}$, we obtain
\begin{align}
 A_n={}&n\sum_{i+j=n-1}\binom{n-1}{i}
 \left[
 8x A_i A_j''+4A_i A_j'-4x A_i'A_j'
 -\left(4x+\frac3x\right)A_i A_j
 \right]
 \notag\\
 &\hspace{8mm}-\frac12\sum_{i=1}^{n-1}\binom ni A_i A_{n-i}.
 \label{p1:eq:bgg-hurwitz-recurrence}
\end{align}
The last half-sum has integer coefficients whenever the preceding $A_i$
do.  To see this, pair its terms with indices $i$ and $n-i$.  Each pair
contributes $\binom ni A_i A_{n-i}$.  If $n=2m$, the remaining middle term
has integer coefficient
\[
 \frac12\binom{2m}{m}=\binom{2m-1}{m-1}.
\]
Since $A_0=1$, recurrence \eqref{p1:eq:bgg-hurwitz-recurrence} therefore proves
$A_n\in\mathbb Z[x,x^{-1}]$ for every $n$.

It remains to determine the parity, and it is essential here to reduce
the half-sum only after the preceding pairing.  The first positive-index
term is
\[
 A_1(x)=-4x-\frac3x\equiv x^{-1}\pmod2.
\]
Suppose now that $n\geq2$ and that
$A_i\equiv x^{-i}\pmod2$ for $i<n$.  All derivative terms and the $4x$
term in the first sum of \eqref{p1:eq:bgg-hurwitz-recurrence} vanish modulo
$2$.  Its remaining contribution is
\[
 -\frac{3n}{x}\sum_{i+j=n-1}\binom{n-1}{i}x^{-i-j}
 =-3n\,2^{n-1}x^{-n}\equiv0\pmod2.
\]
In the paired half-sum, each product $A_iA_{n-i}$ may be replaced modulo
$2$ by $x^{-n}$.  The sum of the resulting integer coefficients is
\[
 \frac12\sum_{i=1}^{n-1}\binom ni
 =\frac{2^n-2}{2}=2^{n-1}-1,
\]
which is odd.  Thus $A_n\equiv x^{-n}\pmod2$, proving
\eqref{p1:eq:bgg-hurwitz-parity}.  The exponent range already established
implies \eqref{p1:eq:bgg-polynomial-parity}.

In particular $A_n(1)$ is an odd integer.  Since
$r_n(1)=2^n A_n(1)/n!$, Legendre's formula gives
\[
 v_2\bigl(r_n(1)\bigr)=n-v_2(n!)=s_2(n).
\]
Finally, \eqref{p1:eq:bgg-r-even} ensures that $r_n(1)$ is an integer, so
$d_n(1)=r_n(1)/2^{6n}$ has no odd denominator factor.  Cancelling its exact
power of $2$ proves \eqref{p1:eq:bgg-exact-denominator}.
\end{proof}

\begin{corollary}[Dyadic nonvanishing]
\label{p1:cor:bgg-dyadic-roots}
Let $n\geq1$ and let $0\ne\alpha\in\overline{\mathbb Q}_2$ satisfy
$v_2(\alpha)\geq0$.  Then
\begin{equation}
 \label{p1:eq:bgg-dyadic-argument-valuation}
 v_2\bigl(d_n(\alpha)\bigr)
 =s_2(n)-6n-nv_2(\alpha).
\end{equation}
Consequently $d_n$ has no nonzero algebraic root that is integral at a
place above $2$.  Equivalently, every algebraic root of $d_n$, under every
embedding into $\overline{\mathbb Q}_2$, has strictly negative $2$-adic
valuation.
\end{corollary}

\begin{proof}
By \eqref{p1:eq:bgg-polynomial-parity}, write $P_n=1+2Q_n$ with
$Q_n\in\mathbb Z[x^2]$.  Integrality of $\alpha$ gives
$v_2(Q_n(\alpha))\geq0$, and therefore $v_2(P_n(\alpha))=0$.
The identity
\[
 d_n(\alpha)=\frac{2^{-5n}}{n!}\,\alpha^{-n}P_n(\alpha)
\]
now yields \eqref{p1:eq:bgg-dyadic-argument-valuation}.  Its right-hand side
is finite, proving the asserted nonvanishing and the statement about
algebraic roots.
\end{proof}

The theorem establishes the exact denominator and a family of nonvanishing
statements.  It does not determine the real sign of $d_n(1)$; in particular,
the negativity assertion for $n\geq3$ is proved by the separate argument in Section~\ref{p1:sec:bgg-sign}.

\subsection{Recovered congruence and its Laurent refinement}
The following all-order congruence is recovered from the earlier project
record \cite{p1:priorTransfer}. Its proof is included to make the present
arithmetic account self-contained. Congruences are coefficientwise in
$\Z[x,x^{-1}][[t]]$.
Since positive-index $r_k(x)$ are even, write $R=1+2A$.

\begin{theorem}[BGG congruence, with its two initial exceptions]
For the normalized BGG sequence $r_k=64^k d_k$, one has
\[
 r_k\equiv\binom{2k}{k}\pmod {32}\qquad(k\geq3).
\]
More precisely,
\[
 R(t,1)\equiv(1-4t)^{-1/2}+16t+16t^2\pmod {32}.
\]
\end{theorem}
\begin{proof}
Substitute $R=1+2A$ into \eqref{p1:eq:bgg-r-scaled}, subtract $1$, and divide
the resulting integer identity by $4$. Reduction modulo $2$ gives
\[
 A+A^2\equiv t/x\pmod2.
\]
The unique solution with zero constant coefficient is
\[
 A\equiv\sum_{j\geq0}(t/x)^{2^j}\pmod2.
\]
Indeed, the square shifts this sum by one index, and uniqueness follows
coefficientwise from the equation. Differentiation with respect to $x$
is legitimate for Laurent series and yields
\[
 A_x\equiv t/x^2\pmod2,
\]
since the derivative of every summand with $j\geq1$ vanishes in
characteristic two.

Because $R_x=2A_x$ and $R_{xx}=2A_{xx}$, reduction of
\eqref{p1:eq:bgg-r-scaled} modulo $64$ now gives
\[
 \bigl(1+(16x+12x^{-1})t\bigr)R^2
 \equiv1+32tA_x\equiv1+32t^2x^{-2}\pmod{64}.
\]
Specializing to $x=1$, and inverting the denominator of constant
coefficient $1$, we obtain
\begin{equation}\label{p1:eq:bgg-square-mod64}
 R(t,1)^2\equiv\frac{1+32t^2}{1+28t}\pmod{64}.
\end{equation}

Set
\[
 C(t)=(1-4t)^{-1/2}=\sum_{k\geq0}\binom{2k}{k}t^k,
 \qquad S(t)=C(t)+16t+16t^2.
\]
For $k\geq1$, the central binomial coefficient is even, so $C\equiv1$
modulo $2$. Consequently
\[
 S^2\equiv\frac1{1-4t}+32t+32t^2\pmod{64}.
\]
Multiplication by $1+28t$ gives $1+32t^2$ modulo $64$, since
\[
 \frac{1+28t}{1-4t}=1+\frac{32t}{1-4t}\equiv1+32t\pmod{64}.
\]
Together with \eqref{p1:eq:bgg-square-mod64}, this proves
$R(t,1)^2-S(t)^2\in64\mathbb Z[[t]]$.
Both $R(t,1)$ and $S(t)$ are congruent to $1$ modulo $2$. Therefore
\[
 U(t):=\frac{R(t,1)+S(t)}2\in\mathbb Z[[t]],\qquad U(0)=1.
\]
It is an integral unit. Factoring the difference of squares and cancelling
this unit gives
\[
 (R(t,1)-S(t))U(t)\in32\mathbb Z[[t]],
 \qquad R(t,1)-S(t)\in32\mathbb Z[[t]].
\]
Extracting coefficients proves the theorem.
\end{proof}

The two corrections are necessary: $r_1=-14\equiv2+16$ and
$r_2=86\equiv6+16$ modulo $32$. This theorem proves the all-order
congruence observed numerically in BGG Remark 11, using the normalization and integrality proved above. The sign
observation is established separately in Section~\ref{p1:sec:bgg-sign}.

\begin{proposition}[Laurent refinement]
Put $u=t/x$ and $h=x+x^{-1}$. Then
\[
 R(t,x)\equiv(1-4u)^{-1/2}
 -8th(1-4u)^{-3/2}
 +16u^2(1-4u)^{-1/2}\pmod{32}.
\]
\end{proposition}
\begin{proof}
Denote the displayed right-hand side by $S_x$. Its coefficients are
integral Laurent polynomials: those of $(1-4u)^{-3/2}$ are
$(2k+1)\binom{2k}{k}$. Squaring modulo $64$ yields
\[
 S_x^2\equiv
 \frac{1+32u^2}{1-4u}-\frac{16th}{(1-4u)^2}\pmod{64}.
\]
Multiply by
$1+(16x+12x^{-1})t=(1-4u)+16th$.
The cross terms cancel modulo $64$, leaving $1+32u^2$.
Thus $S_x^2\equiv R^2\pmod{64}$. Both series are congruent to $1$
modulo $2$, so $(S_x+R)/2$ is an integral unit, and the same
factorization argument proves $S_x\equiv R\pmod{32}$.
\end{proof}


\section{The full sign assertion for the Gompertz asymptotic coefficients}
\label{p1:sec:bgg-sign}

Brent, Glasser and Guttmann \cite{p1:BrentGlasserGuttmann2019}, Remark~11,
reported $d_k<0$ for $3\leq k\leq1000$ by computation. The following
argument establishes the infinite assertion. It combines their
special-function product with the heat kernel of a radial oscillator.
The finite part uses exact integer arithmetic, with a self-contained
verification program below.

\begin{theorem}\label{p1:thm:bgg-sign}
For the coefficients defined by
\[
 -2n^{3/2}\rho_n\sim\sum_{k\geq0}d_kn^{-k},\qquad
 \rho_n=-e^{-1}\Gamma(n)M(n+1,2,1)U(n,0,1),
\]
one has $d_k<0$ for every integer $k\geq3$.
\end{theorem}

\subsection{A heat-kernel coefficient identity}

Let
\[
 \mathscr H=-\frac{d^2}{dr^2}+\frac{r^2}{16}+\frac3{4r^2}
\]
be the positive self-adjoint radial-oscillator operator on $(0,\infty)$
with its regular realization at zero. Its heat kernel is
\begin{equation}\label{p1:eq:radial-kernel}
 K(s;r,z)=\frac{\sqrt{rz}}{4\sinh(s/2)}
 \exp\left(-\frac{r^2+z^2}{8}\coth(s/2)\right)
 I_1\left(\frac{rz}{4\sinh(s/2)}\right).
\end{equation}
This is the Hermite-type Laguerre kernel with parameter $\alpha=1$,
after the changes of variables $x=r/2$, $y=z/2$, and $t=s/4$;
see \cite{p1:NowakStempak2016}, Section~2.2 and equation~(1).
The factor $1/2$ contributed by dilation of the Lebesgue integration
variable is included in \eqref{p1:eq:radial-kernel}.

To connect this kernel to the coefficients above, put $x=r^2/4$ and
\[
 u(x)=e^{-x/2}xM(n+1,2,x),\qquad
 v(x)=e^{-x/2}U(n,0,x),\qquad
 H_n(x)=\Gamma(n+1)u(x)v(x).
\]
Both $u,v$ solve $y''=(n/x+1/4)y$, and
$W_x(u,v)=-1/\Gamma(n+1)$. The functions
$u(r^2/4)/\sqrt r$ and $v(r^2/4)/\sqrt r$ are respectively the
regular and decaying solutions of $(\mathscr H+n)\psi=0$.
Their $r$-Wronskian is $-1/(2\Gamma(n+1))$.
The Green-function formula therefore gives the diagonal resolvent
\begin{equation}\label{p1:eq:radial-resolvent}
 G_n(r,r):=(\mathscr H+n)^{-1}(r,r)
 =\frac{2H_n(x)}r
 =\frac{H_n(x)}{\sqrt x}
 =\int_0^\infty e^{-ns}K(s;r,r)\,ds.
\end{equation}
At $x=1$, $H_n(1)=-n\rho_n$. Thus
\begin{equation}\label{p1:eq:resolvent-bgg-match}
 2\sqrt n\,G_n(2,2)=-2n^{3/2}\rho_n.
\end{equation}

For $s$ near zero, set
\[
 S(s)=\frac{\sinh(s/2)}{s/2},\qquad
 B_l(s)=S(s)^{l-1/2}\exp\bigl(-\tanh(s/4)\bigr),
\]
where the square root is the analytic branch equal to one at zero.
The standard large-positive-argument expansion of $I_1$ \cite[10.40.1]{p1:DLMF} is
\[
 I_1(z)\sim\frac{e^z}{\sqrt{2\pi z}}
       \sum_{l\geq0}b_lz^{-l},\qquad
 b_l=\frac{(-1/2)_l(3/2)_l}{2^l l!}.
\]
Here $b_0=1$ and $b_l<0$ for every $l\geq1$.
Substitution in \eqref{p1:eq:radial-kernel} at $r=2$ gives
\begin{equation}\label{p1:eq:kernel-formal}
 K(s;2,2)\sim\frac1{\sqrt{4\pi s}}\,
 S(s)^{-1/2}e^{-\tanh(s/4)}
       \sum_{l\geq0}b_l\sinh(s/2)^l.
\end{equation}
Indeed $\coth(s/2)-1/\sinh(s/2)=\tanh(s/4)$.
Only terms with $l\leq k$ affect the coefficient of $s^k$.
Applying Watson's lemma to \eqref{p1:eq:radial-resolvent} and using
\eqref{p1:eq:resolvent-bgg-match} yields the exact coefficient formula
\begin{equation}\label{p1:eq:sign-coefficients}
 d_k=(1/2)_k\left(
 [s^k]B_0(s)-\sum_{l=1}^k a_l[s^{k-l}]B_l(s)\right),
 \qquad
 a_l=-\frac{(-1/2)_l(3/2)_l}{4^l l!}>0.
\end{equation}
The use of the divergent Bessel expansion here is purely asymptotic:
for each fixed $k$ it is truncated after finitely many terms, with a
remainder of the requisite order as $s\downarrow0$. The heat-kernel
tail away from zero is exponentially small in $n$, so it does not
change any coefficient. Thus no convergence of the infinite formal
sum is asserted or needed.

\subsection{Uniform dominance of the negative term}

Fix an integer $k\geq144$ and take $r=3/\sqrt k\leq1/4$ as a Cauchy
radius in the complex $s$-plane. For $|s|\leq r$,
\begin{equation}\label{p1:eq:S-bound}
 |S(s)-1|\leq\frac{r^2}{20}.
\end{equation}
For example, the first nonconstant term of its absolutely convergent
series is $r^2/24$, and for $r\leq1$ the ratio of each subsequent
term to its predecessor is at most $1/80$.
In particular the branch defining $S^{-1/2}$ is analytic on the disk.
Writing $u=r^2/20\leq1/2$ and using
$-\tfrac12\log(1-u)\leq u$ gives
\[
 |S(s)^{l-1/2}|\leq\exp((k+1)r^2/20)
       \qquad(0\leq l\leq k).
\]
Also $|\tanh(s/4)|\leq r/3$ on this disk. To see this explicitly,
put $z=s/4$, so $|z|\leq1/16$. Then
\[
 |\sinh z|\leq |z|\cosh|z|,\quad
 \cosh|z|\leq e^{|z|}\leq\frac{16}{15},\quad
 |\cosh z|\geq2-\cosh|z|\geq\frac{14}{15},
\]
whence $|\tanh z|\leq(8/7)|z|<r/3$.
Consequently
\[
 |B_l(s)|
 \leq\exp\left(\frac{(k+1)r^2}{20}+\frac r3\right)
 \leq\exp\left(\frac9{20}+\frac1{320}+\frac1{12}\right)<2.
\]
The last strict inequality follows, for instance, because the
exponent is less than $2/3$, while $e<11/4$ implies $e^2<8$.
Cauchy's estimate now gives
\begin{equation}\label{p1:eq:B-Cauchy}
 |[s^m]B_l(s)|\leq2r^{-m}\qquad(m\geq0,\ 0\leq l\leq k).
\end{equation}

The coefficients $a_l$ satisfy $a_1=3/16$ and
\[
 \frac{a_{j+1}}{a_j}
 =\frac{(j-1/2)(j+3/2)}{4(j+1)}
 \geq\frac{j-1/2}{4}\qquad(j\geq1).
\]
Concavity of the logarithm gives the midpoint bound
\[
 \sum_{j=1}^{k-1}\log(j-1/2)
 \geq\int_0^{k-1}\log t\,dt
 =(k-1)\log(k-1)-(k-1).
\]
Since the terms $\log(j-1/2)$ are increasing, the mean of a terminal
segment is at least the mean of the whole segment. Hence, for
$1\leq l<k$,
\begin{equation}\label{p1:eq:a-ratio}
 \frac{a_l}{a_k}
 \leq\left(\frac{4e}{k-1}\right)^{k-l}.
\end{equation}
Put
\[
 b=\frac{4e}{(k-1)r}
   =\frac{4e\sqrt k}{3(k-1)}.
\]
The last expression decreases for $k>1$; using $e<11/4$ shows that
\[
 b<\frac{11\sqrt{144}}{3(144-1)}=\frac4{13}.
\]
By \eqref{p1:eq:B-Cauchy} and \eqref{p1:eq:a-ratio}, all the terms with
$1\leq l<k$, divided by $a_k$, have total absolute value at most
\begin{equation}\label{p1:eq:tail-one}
 2\sum_{m=1}^{k-1}b^m
 <\frac{2b}{1-b}\leq\frac89.
\end{equation}
The term with $l=0$ is bounded separately. Since $a_1=3/16$,
\begin{align*}
 \frac{|[s^k]B_0(s)|}{a_k}
 &\leq\frac{32}{3r}b^{k-1}
 =\frac{32\sqrt k}{9}b^{k-1}\\
 &<\frac{32\sqrt k}{9}\left(\frac13\right)^{k-1}
 \leq\frac{128}{3}\,3^{-143}<\frac19.
\end{align*}
The sequence $\sqrt k\,3^{-(k-1)}$ decreases for $k\geq1$.
Finally $B_k(0)=1$, so the $l=k$ term in
\eqref{p1:eq:sign-coefficients} is exactly $-a_k$.
The remaining terms have combined absolute value strictly less than
$a_k$ by \eqref{p1:eq:tail-one} and the last estimate. Therefore $d_k<0$
for every $k\geq144$.

\subsection{Finite exact certificate}

The remaining values $3\leq k\leq143$ can be verified using the scalar
recurrence of \cite{p1:BrentGlasserGuttmann2019}, Lemma~14 and Corollary~15.
The following program uses only integer arithmetic. Here
$D_k=64^kd_k$, and
$T_{j,r}=4^r(j+1/2)_r/r!$ is an integer. The four terms defining
\texttt{beta} are exactly $64^{k-j}\alpha_{j,k}$ in that recurrence.
The assertions verify divisibility at every step and all required
strict signs. The first seven terms provide a normalization check.

\begin{lstlisting}
K = 143
T = []
for j in range(K):
    row = [1]
    for r in range(1, K-j+3):
        v, rem = divmod(row[-1]*(4*j+4*r-2), r)
        assert rem == 0
        row.append(v)
    T.append(row)
p2 = [2**m for m in range(K+1)]
p3 = [3**m for m in range(K+1)]
D = [1]
for k in range(1, K+1):
    total = 0
    for j in range(k):
        m = k-j
        row = T[j]
        beta = ((-1+3*p2[m-1]-2*p3[m])*row[m-1]) << (4*m+2)
        beta += ((7-17*p2[m]+17*p3[m])*row[m]) << (4*m)
        beta += ((-13+38*p2[m]-33*p3[m])*row[m+1]) << (4*m-2)
        beta += (6*(1-4*p2[m]+3*p3[m])*row[m+2]) << (4*m-4)
        total += beta*D[j]
    val, rem = divmod(total, 8*k)
    assert rem == 0
    D.append(val)
assert D[:7] == [1, -14, 86, -3660, -1042202,
                 -247948260, -108448540420]
assert all(D[k] < 0 for k in range(3, K+1))
print("Exact finite certificate verified: 3 <= k <= 143")
\end{lstlisting}

This finite certificate, together with the analytic argument for all
$k\geq144$, proves Theorem~\ref{p1:thm:bgg-sign}. The same scalar
calculation was also run through $k=1000$ as a consistency check;
that larger range is not needed by the proof.

\subsection{Uniform eventual sign and late-coefficient growth}

The same argument also gives information not limited to evaluation at
$x=1$. Define $d_k(x)$ by
\[
 2\sqrt{n/x}\,H_n(x)\sim\sum_{k\geq0}d_k(x)n^{-k}
 \qquad(x>0).
\]
The kernel derivation gives the analogue of
\eqref{p1:eq:sign-coefficients} with
\[
 B_{l,x}(s)=S(s)^{l-1/2}e^{-x\tanh(s/4)},\qquad
 a_l\ \hbox{replaced by}\ a_lx^{-l}.
\]

\begin{corollary}\label{p1:cor:uniform-sign}
For every $x>0$ and every integer $k\geq144\max(1,x^2)$, one has
$d_k(x)<0$. For each fixed $x>0$, as $k\to\infty$,
\begin{align}
 d_k(x)
 &=\frac{(-1/2)_k(1/2)_k(3/2)_k}{4^k k!x^k}
       \bigl(1+O_x(k^{-1})\bigr)\label{p1:eq:late-coefficient}\\
 &=-\frac{(k!)^2}{\pi^{3/2}(4x)^k k^{3/2}}
       \bigl(1+O_x(k^{-1})\bigr).\nonumber
\end{align}
\end{corollary}
\begin{proof}
Keep $r=3/\sqrt k$. Under the stated restriction,
$x r/3=x/\sqrt k\leq1/12$, so the same bound
$|B_{l,x}|<2$ holds on the Cauchy circle. The geometric ratio is now
\[
 b_x=\frac{4ex\sqrt k}{3(k-1)}
 <\frac{11}{3}\frac{x}{\sqrt k}\frac{k}{k-1}
 \leq\frac4{13}.
\]
The terms with $1\leq l<k$ have total absolute value, relative to the
dominant $-a_kx^{-k}$, less than $8/9$. The $l=0$ term is bounded by
\[
 \frac{32x\sqrt k}{9}\,b_x^{k-1}
 \leq\frac{8k}{27}\,3^{-(k-1)}<\frac19.
\]
This proves the sign assertion.

For fixed $x$, $b_x=O_x(k^{-1/2})$. Separate the $l=k-1$ term and use
$[s]B_{k-1,x}(s)=-x/4$. Its relative absolute value is
\[
 \frac{x^2}{4}\frac{a_{k-1}}{a_k}
 =\frac{x^2k}{(k-3/2)(k+1/2)}=O_x(k^{-1}).
\]
The remaining terms with $1\leq l\leq k-2$ have total relative
absolute value at most $2b_x^2/(1-b_x)=O_x(k^{-1})$; the $l=0$ term
is smaller still. This proves the first formula in
\eqref{p1:eq:late-coefficient}. The second follows from the gamma-ratio
form of Stirling's formula and
\[
 \Gamma(-1/2)\Gamma(1/2)\Gamma(3/2)=-\pi^{3/2}.
\]
\end{proof}

\begin{corollary}\label{p1:cor:second-factorial-borel}
For fixed $x>0$, the formal series $\sum_{k\geq0}d_k(x)z^k$ has
exact Gevrey order $2$: there are constants $C,A>0$ with
$|d_k(x)|\leq CA^k(k!)^2$, and no analogous bound with $(k!)^s$
for any $s<2$. Its second-factorial Borel transform
\[
 \mathcal B_x(z)=\sum_{k\geq0}\frac{d_k(x)}{(k!)^2}z^k
\]
has radius of convergence exactly $4x$. The value
$\mathcal B_x(4x)$ exists as an absolutely convergent series, but
\[
 \mathcal B_x'(z)\longrightarrow-\infty
       \quad\hbox{as }z\uparrow4x.
\]
In particular $z=4x$ is a singular point of this transform.
\end{corollary}
\begin{proof}
All growth and radius assertions follow directly from
\eqref{p1:eq:late-coefficient}. At $z=4x$ the series terms are
asymptotic to $-\pi^{-3/2}k^{-3/2}$ and hence are absolutely summable.
The differentiated terms are asymptotic to
$-(4x\pi^{3/2})^{-1}k^{-1/2}$. They are negative from some point
onward, so monotone convergence of their absolute values gives the
one-sided derivative limit. An analytic continuation across $4x$
would have a finite derivative there, which is impossible.
\end{proof}

\section{The individual Kummer coefficients: resolving BGG Remark 8}

Brent--Glasser--Guttmann (BGG) \cite{p1:BrentGlasserGuttmann2019}, Remark 8 of
\emph{A conjectured integer sequence arising from the exponential integral}
(arXiv:1812.00316v4, p.~9), observed that the coefficients $c_m$ in their
Theorem 5 appear negative except for $c_0$ and $c_4$, and reported numerical
verification through $m=1000$. This section proves that sign statement.
The coefficients here are exactly their $c_m$, not the Hadamard-product
coefficients $d_m$.

Define the analytic function
\[
 \mu(z)=z^{-1}-(e^z-1)^{-1}-\frac12,
 \qquad A(z)=e^{\mu(z)}=\sum_{\ell\geq0}a_\ell z^\ell.
\]
The apparent singularity at zero is removable. Both functions are analytic
in $|z|<2\pi$. Reindexing BGG Theorem 5, equation (12), gives the exact
finite formula
\begin{equation}\label{p1:eq:bgg-c-reindexed}
 (-1)^m c_m=\sum_{\ell=0}^m a_\ell T_{m,\ell},\qquad
 T_{m,\ell}=
 \frac{(2\ell-m+3/2)_{2m-2\ell}}{4^{m-\ell}(m-\ell)!}.
\end{equation}
The subscript on the parenthesized expression denotes the rising
factorial. In particular $a_0=1$.

\begin{lemma}[A uniform bound for the analytic multiplier]
For every $\ell\geq0$, $|a_\ell|<2^{1-\ell}$.
\end{lemma}
\begin{proof}
Using the Bernoulli expansion of $\mu$ and the usual alternating signs
of the even Bernoulli numbers, on $|z|=2$ one has
\[
 |\mu(z)|\leq\sum_{j\geq1}\frac{|B_{2j}|2^{2j-1}}{(2j)!}
 =\frac12-\frac12\cot1<\frac12.
\]
The identity follows from the power series of $\cot(z/2)$, and the last
inequality uses $0<1<\pi/2$. Hence $|A(z)|<e^{1/2}<2$ on this circle.
Cauchy's coefficient estimate gives the result.
\end{proof}

\begin{lemma}[Uniform suppression away from the leading Bessel term]
If $m\geq10$ and $1\leq\ell\leq m$, then
\[
 \left|\frac{T_{m,\ell}}{T_{m,0}}\right|
 \leq\left(\frac{216}{m}\right)^\ell.
\]
\end{lemma}
\begin{proof}
Cancellation of rising factorials, or the gamma-function expression for
them, gives
\begin{equation}\label{p1:eq:bgg-c-ratio}
 \left|\frac{T_{m,\ell}}{T_{m,0}}\right|
 =\frac{4^\ell\,m!/(m-\ell)!}
 {|(3/2-m)_{2\ell}|}.
\end{equation}
None of the denominator factors is zero.
If $\ell\leq3m/8$, all its $2\ell$ factors are negative, and their
absolute values are at least
\[
 m-2\ell-\frac12\geq\frac m4-\frac12\geq\frac m5.
\]
Since $m!/(m-\ell)!\leq m^\ell$, the ratio is at most
$(100/m)^\ell$.

If $\ell>3m/8$, the denominator consists of $2\ell$ consecutive
half-integers. Each positive absolute half-integer can occur at most
twice, so
\[
 |(3/2-m)_{2\ell}|\geq\bigl((1/2)_\ell\bigr)^2.
\]
Concavity of $\log$ and the midpoint rule on each unit interval imply
\[
 \log((1/2)_\ell)=\sum_{j=0}^{\ell-1}\log(j+1/2)
 \geq\int_0^\ell\log u\,du=\ell\log\ell-\ell.
\]
Consequently the ratio in \eqref{p1:eq:bgg-c-ratio} is at most
\[
 \left(\frac{4me^2}{\ell^2}\right)^\ell
 <\left(\frac{256e^2}{9m}\right)^\ell
 <\left(\frac{216}{m}\right)^\ell,
\]
where $e<11/4$ suffices for the final numerical bound.
\end{proof}

\begin{theorem}[All-order sign of the BGG coefficients $c_m$]
One has
\[
 c_0=1>0,\qquad c_4=\frac{710401}{127401984}>0,
 \qquad c_m<0\quad(m\geq1,\ m\ne4).
\]
\end{theorem}
\begin{proof}
For $m\geq325$, the preceding lemmas give
\[
 \left|\sum_{\ell=1}^m a_\ell
 \frac{T_{m,\ell}}{T_{m,0}}\right|
 <2\sum_{\ell=1}^{\infty}\left(\frac{108}{m}\right)^\ell
 =\frac{216}{m-108}<1.
\]
Thus the sum in \eqref{p1:eq:bgg-c-reindexed} has the same sign as
$T_{m,0}$. Its numerator $(3/2-m)_{2m}$ has precisely $m-1$ negative
factors for $m\geq1$. Therefore $(-1)^mT_{m,0}<0$, proving $c_m<0$
for every $m\geq325$.

The remaining finite range $1\leq m\leq324$ is checked in exact rational
arithmetic by the certificate algorithm below. It gives only the exception
$m=4$; no floating-point comparisons are used. The six initial outputs are
\[
 1,-\frac5{48},-\frac{479}{4608},-\frac{15313}{3317760},
 \frac{710401}{127401984},-\frac{3532731539}{214035333120},
\]
which agree with the published normalization.
\end{proof}

\subsection{Reproducible finite certificate}
The following Python program uses only the standard library. It implements
the Bernoulli recurrence, the coefficient recurrence for $e^\mu$, and the
exact finite sum \eqref{p1:eq:bgg-c-reindexed}. To avoid repeatedly forming
large rising factorials, it uses
\[
 \frac{T_{m+1,0}}{T_{m,0}}=
 \frac{4-(2m+1)^2}{16(m+1)},\qquad
 \frac{T_{m,\ell+1}}{T_{m,\ell}}=
 \frac{16(m-\ell)}{(3-2m+4\ell)(5-2m+4\ell)}.
\]
All divisors appearing here are nonzero. Execution verifies all required
signs and produces a canonical text certificate with rows
$m$, numerator, denominator. Its SHA-256 digest is
\begin{center}\ttfamily\small
7be8eb1341cb2153157d2859c6719dc48\\
e1e6288a0daf5f76d7f62ac26dd36fc
\end{center}
where the digest is the concatenation of the two displayed lines.
\begin{lstlisting}
from fractions import Fraction as Q
from math import factorial, comb
from hashlib import sha256

N = 324
B = [Q(1)]
for n in range(1, N + 2):
    B.append(-sum((comb(n+1,j)*B[j]
                   for j in range(n)), Q(0))/(n+1))
mu = [Q(0)] * (N+1)
for j in range(1, N+1, 2):
    mu[j] = -B[j+1] / factorial(j+1)
a = [Q(1)]
for n in range(1, N+1):
    a.append(sum((j*mu[j]*a[n-j]
                  for j in range(1,n+1,2)), Q(0))/n)
c = []
t0 = Q(1)
for m in range(N+1):
    if m:
        t0 *= Q(4-(2*m-1)**2, 16*m)
    term, value = t0, Q(0)
    for ell in range(m+1):
        value += a[ell] * term
        if ell < m:
            term *= Q(16*(m-ell),
                (3-2*m+4*ell)*(5-2*m+4*ell))
    c.append((-1)**m * value)
assert c[:6] == [Q(1), Q(-5,48), Q(-479,4608),
    Q(-15313,3317760), Q(710401,127401984),
    Q(-3532731539,214035333120)]
assert all(c[m] < 0 for m in range(1,N+1) if m != 4)
assert c[4] > 0
cert = '\n'.join(f'{m} {v.numerator} {v.denominator}'
                 for m,v in enumerate(c)) + '\n'
print(sha256(cert.encode()).hexdigest())
\end{lstlisting}


\section{Euler--Gompertz approximants and a positive-filter obstruction}
\label{p1:sec:gompertz-filter}

Put
\[
 \delta=\int_0^\infty\frac{e^{-t}}{1+t}\,dt.
\]
In this section $p_m=P_m(1)$ and $q_m=Q_m(1)$ retain the
normalization of Murabayashi and Yoshida
\cite{p1:MurabayashiYoshida2024}. To avoid confusion between an index and
its even subsequence, write
\[
 \ell_j=q_{2j},\qquad h_j=p_{2j}\qquad(j\geq0).
\]
The source's Conjecture~7.5 concerns certain infinite monotone
subsequences of fractional-part expressions. We do not prove that
conjecture. We do prove the separate, unnumbered numerical prediction
stated immediately before its Theorem~7.1: the factorially cleared even
linear forms diverge in absolute value. We then establish a uniform
obstruction for a much larger class of positive polynomial filters.

\subsection{Exact normalization of the remainder}

\begin{lemma}\label{p1:lem:gompertz-remainder}
For every integer $j\geq0$,
\begin{equation}\label{p1:eq:gompertz-remainder}
 R_j:=\ell_j\delta-h_j
 =\int_0^\infty e^{-t}\frac{t^j}{(1+t)^{j+1}}\,dt>0.
\end{equation}
Moreover $j!\ell_j,j!h_j\in\mathbb Z$.
\end{lemma}
\begin{proof}
Eliminating the odd indices in the continuant recurrence, whose partial
quotients at $x=1$ are $c_{2j}=1/j$ and $c_{2j+1}=1$, gives
\[
 (j+1)u_{j+1}=(2j+2)u_j-ju_{j-1}\qquad(j\geq1)
\]
for both even subsequences, with
\[
 (\ell_0,\ell_1)=(1,2),\qquad(h_0,h_1)=(0,1).
\]
Multiplication by $j!$ turns this into
\[
 U_{j+1}=(2j+2)U_j-j^2U_{j-1},
\]
which proves the integrality assertion.

Denote the integral on the right of
\eqref{p1:eq:gompertz-remainder} temporarily by $I_j$. Its first two values
are $I_0=\delta$ and $I_1=2\delta-1$. For $j\geq1$ the expression
$(j+1)I_{j+1}-(2j+2)I_j+jI_{j-1}$ equals
\[
 \int_0^\infty
  \frac{e^{-t}t^{j-1}}{(1+t)^{j+2}}(j-2t-t^2)\,dt
 =\int_0^\infty\frac{d}{dt}
   \left(\frac{e^{-t}t^j}{(1+t)^{j+1}}\right)\,dt=0.
\]
The boundary terms vanish. Equality with $\ell_j\delta-h_j$ follows
from the recurrence and initial values. Positivity follows from the
integrand.
\end{proof}

\begin{theorem}[The published divergence observation]
\label{p1:thm:my-divergence}
The even approximants in \cite{p1:MurabayashiYoshida2024} satisfy
\[
 \left|j!q_{2j}\delta-j!p_{2j}\right|\longrightarrow+\infty.
\]
More quantitatively, for $j\geq1$,
\begin{equation}\label{p1:eq:remainder-lower}
 R_j\geq\frac{e^{-2\sqrt j-1}}{\sqrt j+2},
\end{equation}
and for $j\geq4$,
\begin{equation}\label{p1:eq:remainder-upper}
 R_j\leq\sqrt j\,e^{3-2\sqrt j}.
\end{equation}
Consequently
\[
 \log(j!R_j)=j\log j-j-2\sqrt j+O(\log j).
\]
\end{theorem}
\begin{proof}
Restrict the integral in \eqref{p1:eq:gompertz-remainder} to
$[\sqrt j,\sqrt j+1]$. On this interval,
\[
 e^{-t}\geq e^{-\sqrt j-1},\qquad
 \left(\frac{t}{1+t}\right)^j
 =e^{-j\log(1+1/t)}\geq e^{-\sqrt j},\qquad
 \frac1{1+t}\geq\frac1{\sqrt j+2}.
\]
This proves \eqref{p1:eq:remainder-lower}; multiplying by $j!$ already
proves the claimed divergence.

For the upper bound, $\log(1+1/t)\geq1/(t+1)$ for $t>0$.
Set $\varepsilon=j^{-1/2}$ and $a=1-\varepsilon$. By the
arithmetic--geometric mean inequality,
\[
 t+\frac j{t+1}
 =\varepsilon t+a(t+1)+\frac j{t+1}-a
 \geq\varepsilon t+2\sqrt{aj}-a.
\]
It follows that
\[
 R_j\leq e^{a-2\sqrt{aj}}
       \int_0^\infty\frac{e^{-\varepsilon t}}{1+t}\,dt
 \leq\sqrt j\,e^{a-2\sqrt{aj}}
 \leq\sqrt j\,e^{3-2\sqrt j}.
\]
In the last step we used $a\leq1$ and
$\sqrt{aj}\geq\sqrt j-1$. The two estimates give
$\log R_j=-2\sqrt j+O(\log j)$; Stirling's formula finishes the proof.
\end{proof}

This establishes a numerical prediction explicitly left unproved in
the source. Its significance is the failure of the stated factorial
normalization, not a resolution of the arithmetic nature of $\delta$.
It should not be reported as a proof of the source's numbered
Conjecture~7.5.

\subsection{All nonnegative integer-polynomial filters}

The substitution $y=t/(1+t)$ transforms the remainder into a moment:
\begin{equation}\label{p1:eq:gompertz-moments}
 R_j=\int_0^1 y^j w(y)\,dy,
 \qquad
 w(y)=\frac{\exp(-y/(1-y))}{1-y}>0.
\end{equation}
A polynomial filter changes the moment index through a finite linear
combination. The following theorem allows its degree and every
coefficient to vary with the approximation index.

\begin{theorem}[Uniform positive-filter obstruction]
\label{p1:thm:positive-filter}
Let $N\geq0$ and let
\[
 P(y)=\sum_{j=0}^d c_jy^j\in\mathbb Z[y]\setminus\{0\},
 \qquad d\leq N,
\]
be nonnegative on $[0,1]$. Define
\[
 A_N=N!\sum_{j=0}^dc_j\ell_j,
 \qquad
 B_N=N!\sum_{j=0}^dc_jh_j.
\]
Then $A_N,B_N\in\mathbb Z$ and
\begin{equation}\label{p1:eq:filter-bound}
 A_N\delta-B_N
 =N!\int_0^1P(y)w(y)\,dy
 \geq\frac{N!}{2e(2N+1)8^N}.
\end{equation}
In particular, these positive integer linear forms tend to $+\infty$
uniformly over all the permitted polynomial filters as $N\to\infty$.
\end{theorem}
\begin{proof}
The integrality follows from Lemma~\ref{p1:lem:gompertz-remainder}, and
the integral identity follows from \eqref{p1:eq:gompertz-moments}.

Write $a_d=c_d$ for the leading coefficient; then $|a_d|\geq1$.
Let $\mathcal L_d$ denote the Legendre polynomial normalized by
$\mathcal L_d(1)=1$. On $I=[0,1/2]$, the shifted polynomial
$\mathcal L_d(4y-1)$ satisfies
\[
 |\mathcal L_d(4y-1)|\leq1,
 \qquad
 \int_I\mathcal L_d(4y-1)^2\,dy=\frac1{2(2d+1)}.
\]
It is orthogonal on $I$ to all polynomials of degree less than $d$.
Since its leading coefficient is
$2^d\binom{2d}{d}$, subtracting its appropriate multiple from $P$
gives
\[
 \int_I P(y)\mathcal L_d(4y-1)\,dy
 =\frac{a_d}{2^{d+1}(2d+1)\binom{2d}{d}}.
\]
Nonnegativity of $P$ and the bound on $\mathcal L_d$ imply
\[
 \int_I P(y)\,dy
 \geq\left|\int_I P(y)\mathcal L_d(4y-1)\,dy\right|
 \geq\frac1{2^{d+1}(2d+1)\binom{2d}{d}}
 \geq\frac1{2(2d+1)8^d}.
\]
The weight satisfies $w(y)\geq e^{-1}$ on $I$. The contribution from
$[1/2,1)$ is nonnegative, so
\[
 \int_0^1P(y)w(y)\,dy
 \geq\frac1{2e(2d+1)8^d}
 \geq\frac1{2e(2N+1)8^N}.
\]
Multiplication by $N!$ proves \eqref{p1:eq:filter-bound}. Its lower bound
tends to infinity by Stirling's formula.
\end{proof}

\begin{corollary}[Unbounded orders of finite differences]
\label{p1:cor:finite-difference-barrier}
For all integers $n,m\geq0$,
\[
 \sum_{j=0}^m(-1)^j\binom mj R_{n+j}
 =\int_0^\infty e^{-t}
       \frac{t^n}{(1+t)^{n+m+1}}\,dt>0.
\]
For every path of pairs $(n,m)$ with $n+m\to\infty$,
\[
 (n+m)!\sum_{j=0}^m(-1)^j\binom mj R_{n+j}
 \longrightarrow+\infty.
\]
\end{corollary}
\begin{proof}
Apply Theorem~\ref{p1:thm:positive-filter} to
$P(y)=y^n(1-y)^m$, of degree $N=n+m$. Alternatively, integrating only
over $t\in[1,2]$ gives the direct lower bound
\[
 (n+m)!\int_0^\infty e^{-t}
       \frac{t^n}{(1+t)^{n+m+1}}\,dt
 \geq\frac{e^{-2}(n+m)!}{3^{n+m+1}}\longrightarrow+\infty.
\]
\end{proof}

\paragraph{Precise scope.}
Theorem~\ref{p1:thm:positive-filter} concerns the specified factorial
normalization. It does not exclude exceptional savings from common
divisors of $(A_N,B_N)$, smaller denominator-clearing factors,
sign-changing polynomial filters, or other approximation families.
Indeed, if $G_N$ is a positive common divisor and
$(A_N\delta-B_N)/G_N\to0$, then \eqref{p1:eq:filter-bound} forces the
necessary condition
\[
 \frac{G_N}{N!/[2e(2N+1)8^N]}\longrightarrow+\infty.
\]
This last condition is necessary, not sufficient. The theorem is a
newly formulated structural extension in this work, not an additional
previously published conjecture counted as solved.


\section{Polynomial--Gamma thresholds: a sharp answer with explicit scope}
\label{p1:sec:gamma-thresholds}

Bulboac\u{a} and Zayed~\cite{p1:bz2026}, Remark~3(2), ask for shifts ensuring decrease of a polynomial divided by $\Gamma(s+\rho)\Gamma(s)$ on $[1,\infty)$. The remark does not specify signs of the polynomial coefficients. We separate a fixed positive polynomial, a uniform nonnegative-coefficient class, and an unrestricted uniform interpretation. These are three readings of one published question. The following results do not establish irrationality or transcendence of $\gamma$.

Write $\psi=\Gamma'/\Gamma$ and let $\chi:\mathbb R\to(0,\infty)$ be the inverse of $\psi$. This inverse exists because
\[
 \psi'(x)=\sum_{k=0}^{\infty}\frac1{(x+k)^2}>0,
 \qquad \psi(0+)=-\infty,\qquad\psi(\infty)=\infty.
\]
For a real polynomial $P$ positive on $[1,\infty)$ and $\rho>-1$, put
\[
 Q_{P,\rho}(s)=\frac{P(s)}{\Gamma(s+\rho)\Gamma(s)}.
\]

\begin{theorem}[Exact threshold for a fixed positive polynomial]
\label{p1:thm:fixed-poly-threshold}
The quantity
\[
 B_P=\max_{s\ge1}\left\{
 \chi\!\left(\frac{P'(s)}{P(s)}-\psi(s)\right)-s\right\}
\]
is finite, the maximum is attained, and $B_P>-1$. The function $Q_{P,\rho}$ is strictly decreasing on $[1,\infty)$ if and only if $\rho\ge B_P$.
\end{theorem}
\begin{proof}
The expression being maximized is continuous. If $d=\deg P$, then
\[
 \frac{P'(s)}{P(s)}=\frac ds+O(s^{-2}),\qquad
 \psi(s)=\log s+O(s^{-1}).
\]
Its argument inside $\chi$ tends to $-\infty$, so $\chi$ tends to zero. The expression therefore equals $-s+o(1)$ at infinity. A maximum exists, and its value is at least
$\chi(P'(1)/P(1)+\gamma)-1>-1$.

The logarithmic derivative is
\[
 L_{P,\rho}(s):=\frac{Q_{P,\rho}'(s)}{Q_{P,\rho}(s)}
 =\frac{P'(s)}{P(s)}-\psi(s+\rho)-\psi(s).
\]
By strict increase of $\psi$, $L_{P,\rho}(s)\le0$ for every $s\ge1$ is equivalent to $\rho\ge B_P$. If $\rho<B_P$, the derivative is positive at some point, or by continuity immediately to the right of $1$, and decrease is impossible.

Suppose $\rho\ge B_P$. The function $L_{P,\rho}$ is nonpositive and real analytic on $(1,\infty)$. Its behavior $L_{P,\rho}(s)=-2\log s+o(\log s)$ prevents it from vanishing identically. The analytic identity theorem implies that it cannot vanish on an interval. Its integral over every nondegenerate interval is strictly negative. Hence $Q_{P,\rho}$ is strictly decreasing, also at the critical shift.
\end{proof}

This is an attained variational characterization, rather than an elementary closed formula for an arbitrary polynomial. For nonnegative coefficients a uniform closed formula is available.

\begin{theorem}[Sharp shift uniform over nonnegative coefficients]
\label{p1:thm:uniform-poly-threshold}
Let
\[
 \mathcal P_m=\left\{s^m+\sum_{j=1}^{m-1}a_js^j+1:a_j\ge0\right\},
 \qquad
 c_m=\begin{cases}\tfrac12,&m=1,\\m-1,&m\ge2.\end{cases}
\]
The smallest shift $\rho>-1$ for which $Q_{P,\rho}$ is strictly decreasing on $[1,\infty)$ for every $P\in\mathcal P_m$ is
\[
 R_m=\chi(\gamma+c_m)-1.
\]
\end{theorem}
\begin{proof}
For $m=1$, $P'/P=1/(s+1)\le1/2$, and the increasing function $\psi(s+\rho)+\psi(s)$ has left endpoint value $\psi(1+\rho)-\gamma$. This proves both sufficiency and necessity.

For $m\ge2$, put $A_m(s)=(m-1)(s^m+1)-ms^{m-1}$. Then
\[
 A_m(1)=m-2\ge0,\qquad
 A_m'(s)=m(m-1)s^{m-2}(s-1)\ge0\quad(s\ge1).
\]
Every internal monomial contributes
\[
 a_j\big((m-1)s^j-js^{j-1}\big)
 =a_js^{j-1}\big((m-1)s-j\big)\ge0.
\]
Therefore $P'/P\le m-1$. If $\rho\ge R_m$, then for $s>1$,
\[
 \psi(s+\rho)+\psi(s)>\psi(1+\rho)-\gamma
 \ge m-1\ge\frac{P'(s)}{P(s)},
\]
which proves strict decrease.

For $m=2$, every $P\in\mathcal P_2$ satisfies $P'(1)/P(1)=1$. Thus $\rho<R_2$ gives a positive derivative at $1$. For $m\ge3$, choose
$P_A(s)=s^m+As^{m-1}+1$. Then
\[
 \frac{P_A'(1)}{P_A(1)}=
 \frac{m+(m-1)A}{A+2}\longrightarrow m-1.
\]
If $\rho<R_m$, a sufficiently large finite $A$ makes this ratio exceed $\psi(1+\rho)-\gamma$, again giving a positive derivative. This proves sharpness.
\end{proof}

\begin{proposition}[No uniform threshold without coefficient restrictions]
\label{p1:prop:negative-coeff-obstruction}
For each $m\ge2$ and every fixed $\rho>-1$, some monic degree-$m$ real polynomial with constant coefficient $1$, positive on $[1,\infty)$, makes $Q_{P,\rho}$ fail to decrease. Thus no finite threshold depending only on $m$ exists for this unrestricted class.
\end{proposition}
\begin{proof}
For $m=2$, take $P_t(s)=(s-1)^2+t^2s$ with $t>0$. At $s=1+t$,
\[
 P_t(1+t)=t^2(2+t),\quad P_t'(1+t)=t(2+t),\quad
 \frac{P_t'(1+t)}{P_t(1+t)}=\frac1t.
\]
The quantity $\psi(1+t+\rho)+\psi(1+t)$ stays bounded as $t\downarrow0$, so $Q_{P_t,\rho}'(1+t)>0$ for small $t$. For $m\ge3$, multiply $P_t$ by $s^{m-2}+1$. The product retains the required normalization and positivity, while its logarithmic derivative only increases.
\end{proof}

The proposition refutes the uniform unrestricted reading. It does not refute a reading in which the shift may depend on the coefficients; Theorem~\ref{p1:thm:fixed-poly-threshold} covers that case.

\subsection{An interior transition for sparse numerators}
Let $S_m=B_{P_m}$ for $P_m(s)=s^m+1$, and define
\[
 r_m=\chi(\gamma+m/2)-1.
\]
The value $r_m$ makes the logarithmic derivative vanish at $s=1$. It is not always the optimal shift.

\begin{theorem}[The transition occurs at degree four]
\label{p1:thm:sparse-transition}
For integers $m=1,2,3$, $S_m=r_m$. For every integer $m\ge4$, $S_m>r_m$, and the maximum defining $S_m$ is attained at an interior point $s>1$.
\end{theorem}
\begin{proof}
Put
\[
 J_m(s)=\frac{ms^{m-1}}{1+s^m},\qquad
 G_{m,\rho}(s)=\psi(s+\rho)+\psi(s)-J_m(s).
\]
For $m=1,2$, $J_m$ is nonincreasing on $[1,\infty)$, hence $G_{m,\rho}$ strictly increases and its endpoint determines its sign. For $m=3$,
\[
 J_3'(s)=\frac{3s(2-s^3)}{(1+s^3)^2}<\frac1s<\psi'(s),\qquad s\ge1.
\]
The first inequality follows from positivity of
$f(s)=s^6+3s^5+2s^3-6s^2+1$: $f(1)=1$ and $f'(s)\ge15s>0$ for $s\ge1$. The second follows by comparing the trigamma series with its integral. Thus $G_{3,\rho}'>0$ as well. This gives $S_m=r_m$ for these three degrees.

For $m\ge4$ and $\rho=r_m$, one has $G_{m,r_m}(1)=0$ and
\[
 G_{m,r_m}'(1)=\psi'(1+r_m)+\zeta(2)-\frac{m(m-2)}4.
\]
Because $\psi(4)=11/6-\gamma<\gamma+m/2$, $1+r_m>4$. Since $\psi'$ is decreasing,
\[
 \psi'(1+r_m)<\psi'(4)=\zeta(2)-\frac{49}{36}.
\]
Also $\zeta(2)<5/3$; for instance the first five terms plus the integral tail give $\zeta(2)<5989/3600<5/3$. Hence
\[
 G_{m,r_m}'(1)
 <2\zeta(2)-\frac{49}{36}-2<-\frac1{36}.
\]
Thus $G_{m,r_m}<0$ just to the right of $1$, where the quotient increases. Consequently $S_m>r_m$. Since the endpoint value in the maximizing formula equals $r_m$, its attained maximum must be interior.
\end{proof}

\begin{theorem}[Euler's constant as a sharp-threshold scale]
\label{p1:thm:sparse-scale}
As $m\to\infty$ through integers,
\[
 S_m=\frac{e^{m+\gamma-1}}m
 \left(1+O\!\left(\frac{(\log m)^2}{m}\right)\right),
\]
and
\[
 \frac{R_m}{S_m}
 =m\left(1+O\!\left(\frac{(\log m)^2}{m}\right)\right).
\]
In particular, $\displaystyle\lim_{m\to\infty}m e^{-m}S_m=e^{\gamma-1}$.
\end{theorem}
\begin{proof}
For $m\ge2$, $J_m$ attains its maximum on $[1,\infty)$ at
\[
 s_m=(m-1)^{1/m},\qquad M_m=J_m(s_m)=(m-1)^{1-1/m}.
\]
The threshold formula and $\psi(s)\ge-\gamma$ give the two-sided bound
\[
 \chi(M_m-\psi(s_m))-s_m\ \le S_m\ \le\chi(M_m+\gamma)-1.
\]
Elementary expansion gives
\[
 s_m=1+O\!\left(\frac{\log m}{m}\right),\qquad
 \psi(s_m)=-\gamma+O\!\left(\frac{\log m}{m}\right),
\]
and
\[
 M_m=m-1-\log m+O\!\left(\frac{(\log m)^2}{m}\right).
\]
The standard digamma expansion implies
$\chi(y)=e^y+\tfrac12+O(e^{-y})$ as $y\to\infty$.
Applying it to both bounds yields
\[
 S_m=e^{M_m+\gamma}\left(1+O\!\left(\frac{\log m}{m}\right)\right)
 =\frac{e^{m+\gamma-1}}m
 \left(1+O\!\left(\frac{(\log m)^2}{m}\right)\right).
\]
For $m\ge2$, $R_m=\chi(m-1+\gamma)-1
=e^{m+\gamma-1}-\tfrac12+O(e^{-m})$, which proves the ratio formula.
\end{proof}

The last two theorems are extensions developed here, not additional independently posed open problems. They connect $\gamma$ to a change from endpoint to interior optimization and to its exponential growth scale. Recovering $e^\gamma$ as a limit gives no arithmetic classification of that number.

\paragraph{Prior-art qualification.}
The constant-numerator threshold $\chi(\gamma)-1$ is immediate from the same endpoint argument. A public repository already records this result as APP-0009~\cite{p1:pudim2026}, dated May~23, 2026, so it is not counted as a new solution here. That public claim is distinguished from peer-reviewed validation. No matching polynomial-family result was found in the searches made for this note; that limited search does not establish priority. The polynomial analysis above should be counted at most once against the published question.

\section{Arithmetic of Bessel-bridge stopping boundaries}
\label{p1:sec:bessel-arithmetic}

The analytic stopping problem in this section is already solved: for reflected
Brownian bridges see Ekstr\"om--Wanntorp \cite{p1:EkstromWanntorp2009}, and for
arbitrary positive-dimensional Bessel bridges see Hobson--Liu
\cite{p1:HobsonLiu2026}. We prove an arithmetic classification of their boundary
constants. No explicit published arithmetic conjecture was identified, so this
classification is not counted as the resolution of a previously stated open
problem. Its novelty relative to the full literature remains unverified.

\subsection{A Kummer arithmetic lemma}

For positive rational numbers $a,b$, put
\[
 F_{a,b}(z)=M(a,b,z)
 =\sum_{n=0}^{\infty}\frac{(a)_n}{(b)_n}\frac{z^n}{n!}.
\]
Here $(a)_0=1$ and $(a)_n=a(a+1)\cdots(a+n-1)$.

\begin{lemma}[Strictly increasing elasticity]
\label{p1:lem:kummer-elasticity}
For $a,b>0$, the function $zF_{a,b}'(z)/F_{a,b}(z)$ increases strictly
from zero to infinity on $(0,\infty)$. Consequently, for every $c>0$,
\begin{equation}\label{p1:eq:kummer-boundary}
 zF_{a,b}'(z)=cF_{a,b}(z)
\end{equation}
has exactly one positive solution.
\end{lemma}
\begin{proof}
Write $A_n=(a)_n/((b)_n n!)>0$ and define
\[
 \mathbb P_z(N=n)=\frac{A_nz^n}{F_{a,b}(z)}.
\]
Termwise differentiation of the entire series yields
\[
 \frac{zF_{a,b}'(z)}{F_{a,b}(z)}=\mathbb E_zN,
 \qquad
 \frac{\mathrm d}{\mathrm dz}\mathbb E_zN
 =\frac{\operatorname{Var}_z(N)}{z}>0.
\]
The variance is positive since both $n=0$ and $n=1$ have positive
probability. The limit at zero is zero. For each fixed integer $K$, the
polynomial $\sum_{n\le K}A_nz^n$ has smaller degree than
$A_{K+1}z^{K+1}$, so $\mathbb P_z(N\le K)\to0$ as $z\to\infty$.
Thus $\liminf_{z\to\infty}\mathbb E_zN\ge K+1$ for every $K$.
\end{proof}

\begin{lemma}[Arithmetic growth of the coefficients]
\label{p1:lem:kummer-E}
If $a,b$ are positive rational numbers, then $F_{a,b}$ and $F_{a,b}'$
are $E$-functions. They form a system with no nonzero finite singularity:
\begin{equation}\label{p1:eq:kummer-system}
 \frac{\mathrm d}{\mathrm dz}
 \begin{pmatrix}F_{a,b}\\F_{a,b}'\end{pmatrix}
 =\begin{pmatrix}0&1\\a/z&1-b/z\end{pmatrix}
 \begin{pmatrix}F_{a,b}\\F_{a,b}'\end{pmatrix}.
\end{equation}
\end{lemma}
\begin{proof}
Write $a=r/s$, $b=u/v$ in lowest positive terms. The factorial-normalized
coefficients are
\[
 C_n=\frac{(a)_n}{(b)_n}
 =\left(\frac vs\right)^n
   \frac{\prod_{j=0}^{n-1}(r+sj)}{\prod_{j=0}^{n-1}(u+vj)},
 \qquad |C_n|\le\max(1,a/b)^n.
\]
We verify the denominator condition. Fix an integer $L\ge u+v$.
For a prime $p\nmid sv$, each progression has precisely one root modulo
$p^k$. The counts of roots in $n$ consecutive indices differ by at most
one. Since denominator factors for $n\le N$ are at most $LN$, the
negative $p$-adic valuation of any $C_n$, $n\le N$, is at most
$\lfloor\log_p(LN)\rfloor$. All these denominator prime powers therefore
divide $\operatorname{lcm}(1,\ldots,LN)$.
For each of the finitely many primes dividing $sv$, progression
valuations and the prefactor contribute at most $O_{a,b}(N)$: a
progression with step prime to $p$ contributes at most
$\sum_{k\ge1}(N/p^k+\mathbf1_{p^k\le L'N})=O_{a,b}(N)$ for a fixed
$L'$, whereas a progression whose step is divisible by $p$ has initial
term prime to $p$ and contributes zero.
The elementary Chebyshev estimate
$\log\operatorname{lcm}(1,\ldots,m)=O(m)$ now gives a common denominator
$\exp(O_{a,b}(N))$. The derivative has coefficients $C_{n+1}$ and obeys
the same estimates. Finally, direct substitution of the series gives
$zF''+(b-z)F'-aF=0$, hence \eqref{p1:eq:kummer-system}.
\end{proof}

\begin{lemma}[The only first-order cases]
\label{p1:lem:kummer-logder}
For $a,b>0$, the functions $F_{a,b}$ and $F_{a,b}'$ are linearly
dependent over $\mathbb C(z)$ if and only if $a-b\in\mathbb Z_{\ge0}$.
In that case, with $m=a-b$,
\begin{equation}\label{p1:eq:kummer-poly}
 F_{a,b}(z)=e^zP_m(z),\qquad
 P_m(z)=\sum_{j=0}^{m}
 \frac{m!}{(m-j)!\,(b)_j\,j!}z^j.
\end{equation}
\end{lemma}
\begin{proof}
Suppose $R=F_{a,b}'/F_{a,b}$ is rational. Because $F_{a,b}$ is entire
and $F_{a,b}(0)=1$, the finite poles of $R$ are simple, with positive
integer residues equal to zero multiplicities. If $P$ is the polynomial
formed from these finitely many zeros, then $R-P'/P$ is a polynomial.
The Riccati equation
\[
 z(R'+R^2)+(b-z)R-a=0
\]
excludes any positive-degree polynomial part, by comparing degrees at
infinity. Its constant part $c$ satisfies $c^2-c=0$. Hence
$F_{a,b}=Ce^{cz}P$ with $c=0$ or $1$.
If $c=0$, leading coefficients in Kummer's equation give
$\deg P+a=0$, impossible. If $c=1$, the polynomial satisfies
\[
 zP''+(b+z)P'+(b-a)P=0,
\]
so $\deg P=a-b\in\mathbb Z_{\ge0}$.
Conversely, the polynomial in \eqref{p1:eq:kummer-poly} satisfies this last
equation, and $e^zP_m(z)$ is the unique analytic solution of Kummer's
equation with value $1$ at zero. This proves the converse and the formula.
\end{proof}

\begin{theorem}[Arithmetic classification of Kummer elasticity roots]
\label{p1:thm:kummer-root-arithmetic}
Let $a,b$ be positive rational numbers and let $c$ be a positive algebraic
number. The unique positive solution $z_c$ of
\eqref{p1:eq:kummer-boundary} is algebraic if $a-b\in\mathbb Z_{\ge0}$,
and is transcendental otherwise.
\end{theorem}
\begin{proof}
In the first case, \eqref{p1:eq:kummer-poly} reduces the boundary equation
to the nonzero polynomial equation
\[
 z(P_m+P_m')-cP_m=0.
\]
In the second case, Lemmas~\ref{p1:lem:kummer-E} and
\ref{p1:lem:kummer-logder} allow the use of Beukers's
Corollary~1.4 \cite{p1:Beukers2006}: linearly independent $E$-functions in a
rational first-order system take linearly independent values over
$\overline{\mathbb Q}$ at every nonzero algebraic ordinary point of the
system. If $z_c$ were algebraic, the relation
$z_cF_{a,b}'(z_c)-cF_{a,b}(z_c)=0$ would contradict this corollary.
\end{proof}

\subsection{Application to Bessel bridges}

Let $d>0$. Define the squared Bessel bridge $Y$ pinned at zero at time
$1$ through
\[
 \mathrm dY_s=\left(d-\frac{2Y_s}{1-s}\right)\mathrm ds
             +2\sqrt{Y_s}\,\mathrm dW_s,
 \qquad R_s=\sqrt{Y_s}.
\]
This formulation also covers $0<d<1$, where a naive radial SDE at zero
is unsuitable. For $q>0$, write
\[
 V_{d,q}(t,r)=\sup_{t\le\tau\le1}\mathbb E_{t,r}[R_\tau^q],
 \qquad F(z)=M(q/2,d/2,z).
\]
Let $z_{d,q}$ be the unique positive root of
\begin{equation}\label{p1:eq:bessel-root}
 2z_{d,q}F'(z_{d,q})=qF(z_{d,q}),
 \qquad B_{d,q}=\sqrt{2z_{d,q}}.
\end{equation}
The analytic solution of Hobson--Liu \cite[Theorem~3.1]{p1:HobsonLiu2026}
is
\begin{equation}\label{p1:eq:bessel-value}
 V_{d,q}(t,r)=
 \begin{cases}
 (1-t)^{q/2}B_{d,q}^q
 \dfrac{F(r^2/[2(1-t)])}{F(z_{d,q})},
       &r<B_{d,q}\sqrt{1-t},\\[6pt]
 r^q,&r\ge B_{d,q}\sqrt{1-t}.
 \end{cases}
\end{equation}
Their parameters are $\alpha=d$, $n=q$, and their squared boundary
constant is $Z_{\alpha,n}=B_{d,q}^2=2z_{d,q}$.
Indeed their series $\psi(y)$ has coefficient ratio
$(q+2k)/[2(k+1)(d+2k)]$, so
$\psi(y)=M(q/2,d/2,y/2)$. This checks the factor of two in
\eqref{p1:eq:bessel-root}.

\begin{theorem}[Boundary dichotomy for rational dimension and reward]
\label{p1:thm:bessel-arithmetic}
For $d,q\in\mathbb Q_{>0}$,
\[
 B_{d,q}\in\overline{\mathbb Q}
 \quad\Longleftrightarrow\quad
 q-d\in2\mathbb Z_{\ge0}.
\]
In every other case $B_{d,q}$ is transcendental.
\end{theorem}
\begin{proof}
Apply Theorem~\ref{p1:thm:kummer-root-arithmetic} with
$a=q/2$, $b=d/2$, $c=q/2$. Algebraicity of $B_{d,q}$ is equivalent to
algebraicity of $z_{d,q}$.
\end{proof}

\begin{proposition}[A strict boundary estimate and verification]
\label{p1:prop:bessel-boundary-strict}
For all real $d,q>0$,
\[
 B_{d,q}^2>\frac{q+d-2}{2}.
\]
For integer $d\ge1$, formula~\eqref{p1:eq:bessel-value} can be verified
directly by the following argument.
\end{proposition}
\begin{proof}
Put $h(y)=F(y^2/2)$ and $\lambda(y)=yh'(y)/h(y)$.
Lemma~\ref{p1:lem:kummer-elasticity} gives $\lambda'(y)>0$ for $y>0$.
The radial differential equation is
\[
 h''+\left(\frac{d-1}{y}-y\right)h'-qh=0.
\]
At $B=B_{d,q}$, where $\lambda(B)=q$, it follows that
\[
 0<\lambda'(B)=\frac qB(2B^2+2-d-q).
\]
More precisely, for the probability law in
Lemma~\ref{p1:lem:kummer-elasticity} at $z=B^2/2$,
$\lambda'(B)=4\operatorname{Var}(N)/B$, and hence
\[
 B^2=\frac{q+d-2}{2}+\frac{2}{q}\operatorname{Var}(N).
\]
This proves the strict estimate and identifies its exact positive gap.
In particular, it strengthens the
nonstrict bound in \cite[Proposition~3.1]{p1:HobsonLiu2026}, which is an
already-proved proposition rather than an open conjecture.

To verify the value formula, set $v(y)=B^qh(y)/h(B)$ for $y<B$ and
$v(y)=y^q$ for $y\ge B$. Because $\lambda(y)<q$ for $0<y<B$, the
ratio $h(y)/y^q$ decreases there. Thus $v(y)\ge y^q$.
The candidate $(1-t)^{q/2}v(r/\sqrt{1-t})$ satisfies the continuation
PDE, and the boundary equation gives continuous and smooth fit.
In the stopping region its generator is
\[
 q r^{q-2}\left(\frac{q+d-2}{2}-\frac{r^2}{1-t}\right)\le0.
\]
Applying It\^o's formula to the squared radius avoids an issue at
$r=0$, since the continuation function is analytic in $Y=r^2$.
Smooth fit removes the boundary local-time term. The candidate along
the bridge is therefore a nonnegative local supermartingale, bounding
all expected stopping rewards from above. Up to its first boundary
hit it is a local martingale. For integer dimension, the supremum of
the bridge radius has all positive moments, and the candidate is
bounded by a constant plus $R_s^q$; localization and dominated
convergence give equality at the first boundary hit (or at time $1$).
This proves optimality. The result for arbitrary positive dimension
is supplied by the squared-Bessel verification in
\cite[Theorem~3.1]{p1:HobsonLiu2026}.
\end{proof}

\begin{corollary}[Arithmetic of the value in the algebraic-boundary cases]
\label{p1:cor:bessel-value-arithmetic}
Suppose $d,q\in\mathbb Q_{>0}$ and $q=d+2m$, $m\in\mathbb Z_{\ge0}$.
Then $V_{d,q}(0,0)$ is transcendental. More generally, if $t,r$ are
algebraic, $0\le t<1$, $r\ge0$, and
$r<B_{d,q}\sqrt{1-t}$, then $V_{d,q}(t,r)$ is transcendental.
On and outside the stopping boundary its value is algebraic.
\end{corollary}
\begin{proof}
Let $z=z_{d,q}$ and $w=r^2/[2(1-t)]$. Formula~\eqref{p1:eq:kummer-poly}
rewrites the continuation value as
\[
 (1-t)^{q/2}\frac{B_{d,q}^qP_m(w)}{P_m(z)}e^{w-z}.
\]
The prefactor is nonzero algebraic; all coefficients of $P_m$ are
positive. The exponent is nonzero algebraic because $0\le w<z$.
Lindemann's theorem proves transcendence. Outside the continuation
region the value is the algebraic number $r^q$.
\end{proof}

\begin{corollary}[Reflected bridges and excursions]
For reflected one-dimensional Brownian bridges, $B_{1,q}$ is
algebraic exactly when $q$ is a positive odd integer. In particular,
\[
 B_{1,1}=1,\qquad B_{1,3}=3^{1/4},\qquad
 B_{1,5}^{6}+5B_{1,5}^{4}-15B_{1,5}^{2}-15=0.
\]
The corresponding first two initial values are
\[
 V_{1,1}(0,0)=e^{-1/2},\qquad
 V_{1,3}(0,0)=
 \frac{3^{3/4}}{1+\sqrt3}e^{-\sqrt3/2}.
\]
A Brownian excursion is a three-dimensional Bessel bridge. Its linear
reward boundary $B_{3,1}$ is transcendental. More generally, for
rational $q>0$, its boundary is algebraic exactly for
$q\in\{3,5,7,\ldots\}$. At $q=3$,
\[
 B_{3,3}=\sqrt3,\qquad V_{3,3}(0,0)=3^{3/2}e^{-3/2}.
\]
\end{corollary}

\paragraph{Scope of the transfer.}
These results transfer an established $E$-function theorem to a family
of free-boundary constants. They do not imply irrationality or
transcendence of Euler's constant. The individual examples are
instances of one classification theorem, not separate resolutions of
published conjectures.


\subsection{Euler--Gompertz constants as a dimension sensitivity}
\label{p1:subsec:dimension-sensitivity}

The previous arithmetic classification does not use $\gamma$ directly.
There is, however, an exact sensitivity identity involving the classical
combination $\eta=\operatorname{Ein}(1)=\gamma+\delta/e$.
Here $\delta=\int_0^\infty e^{-u}\log(1+u)\,\mathrm du$ and
\[
 \operatorname{Ein}(z)=\int_0^z\frac{1-e^{-u}}{u}\,\mathrm du.
\]
Let $J(d,q)=V_{d,q}(0,0)$, with the dimension $d$ treated as a positive
real parameter in the squared-Bessel model.

\begin{theorem}[Arithmetic independence of value and dimension sensitivity]
\label{p1:thm:value-dimension-sensitivity}
The function $J$ is real analytic for $d,q>0$. At $(d,q)=(2,2)$,
\begin{align*}
 J(2,2)&=\frac2e,\\
 \left.\frac{\partial J}{\partial d}\right|_{(2,2)}
 &=\frac{\operatorname{Ein}(1)}e
   =\frac\gamma e+\frac\delta{e^2},\\
 \left.\frac{\partial J}{\partial q}\right|_{(2,2)}
 &=\frac{\log2-\operatorname{Ein}(1)}e.
\end{align*}
Moreover, $J(2,2)$ and
$\left.\partial J/\partial d\right|_{(2,2)}$ are algebraically independent
over $\overline{\mathbb Q}$.
\end{theorem}
\begin{proof}
The Kummer function is analytic in its positive parameters. Strict
elasticity and the implicit function theorem show that $z=z_{d,q}$
is locally real analytic, hence so is
\[
 J(d,q)=\frac{(2z)^{q/2}}{M(q/2,d/2,z)}.
\]
At $d=q=2$, one has $M(1,1,z)=e^z$ and $z=1$.
For $H_n=\sum_{k=1}^n1/k$ and $H_0=0$, termwise parameter
differentiation gives
\[
 \left.\partial_aM(a,1,z)\right|_{a=1}
 =\sum_{n=0}^\infty H_n\frac{z^n}{n!}
 =e^z\operatorname{Ein}(z).
\]
The last identity follows by writing
$H_n=\int_0^1(1-t^n)/(1-t)\,\mathrm dt$ and summing the exponential
series. Differentiating $M(a,a,z)=e^z$ along the diagonal gives
\[
 \left.\partial_bM(1,b,z)\right|_{b=1}
 =-e^z\operatorname{Ein}(z).
\]
Now
\[
 \log J(d,q)=\frac q2\log(2z)-\log M(q/2,d/2,z).
\]
Its coefficient of every derivative of $z$ vanishes, since smooth fit
is precisely $M_z/M=q/(2z)$. At $(2,2)$ the remaining terms yield
\[
 \partial_d\log J=\frac12\operatorname{Ein}(1),\qquad
 \partial_q\log J=\frac12(\log2-\operatorname{Ein}(1)),
\]
which prove the identities.

Use the established algebraic independence of $e$ and
$\eta=\operatorname{Ein}(1)$; see
\cite[Section~5(ii)]{p1:Rivoal2012}, where the independent pair
$e^{-1},E(-1)$ equals $e^{-1},-\eta$ in the notation used there. Set $U=J(2,2)$ and
$D=\partial_dJ(2,2)$. The displayed identities give
\[
 U=2/e,\quad D=\eta/e,\qquad e=2/U,\quad\eta=2D/U.
\]
Thus $\overline{\mathbb Q}(U,D)=
\overline{\mathbb Q}(e,\eta)$, of transcendence degree two.
\end{proof}

The independence theorem concerns an optimal value together with its
response to a change in the Bessel dimension. It leaves the separate
arithmetic nature of $\gamma$ and $\delta$ unresolved. In particular,
no transcendence assertion about the reward-exponent derivative
$(\log2-\eta)/e$ is made here.


\subsection{Algebraic independence of every finite dimension jet}

Define $\mathcal J(d)=J(d,2)$ and, for integers $n\ge1$, put
\[
 H_n(z)=\sum_{k=1}^{\infty}\frac{(-1)^{k+1}z^k}{k^n k!},
 \qquad H_0(z)=1-e^{-z}.
\]
Thus $H_1=\operatorname{Ein}$ and $H_n'=H_{n-1}/z$ for $n\ge1$,
with the quotient interpreted analytically at zero.

\begin{theorem}[All finite dimension jets are algebraically independent]
\label{p1:thm:all-dimension-jets}
For every integer $N\ge0$, the $N+1$ real numbers
\[
 \mathcal J(2),\mathcal J'(2),\ldots,\mathcal J^{(N)}(2)
\]
are algebraically independent over $\overline{\mathbb Q}$.
More precisely,
\begin{equation}\label{p1:eq:jet-field-equality}
 \overline{\mathbb Q}\bigl(\mathcal J(2),\ldots,
                  \mathcal J^{(N)}(2)\bigr)
 =\overline{\mathbb Q}\bigl(e,H_1(1),\ldots,H_N(1)\bigr).
\end{equation}
\end{theorem}
\begin{proof}
Let $\varepsilon=(d-2)/2$. An exact parameter expansion is
\begin{equation}\label{p1:eq:kummer-parameter-series}
 M(1,1+\varepsilon,z)
 =e^zS(\varepsilon,z),\qquad
 S(\varepsilon,z)=1+\sum_{n\ge1}(-1)^nH_n(z)\varepsilon^n.
\end{equation}
To verify it first take $\varepsilon>0$ and use
\[
 e^{-z}M(1,1+\varepsilon,z)
 =\varepsilon\int_0^1e^{-zt}t^{\varepsilon-1}\,\mathrm dt
 =1+\varepsilon\sum_{k\ge1}\frac{(-z)^k}{k!(k+\varepsilon)}.
\]
Expanding $\varepsilon/(k+\varepsilon)$ proves
\eqref{p1:eq:kummer-parameter-series}; analytic continuation supplies a
neighborhood of $\varepsilon=0$.

Write $z(\varepsilon)=z_{2+2\varepsilon,2}$, so $z(0)=1$ and
\[
 \mathcal J(2+2\varepsilon)
 =\frac{2z(\varepsilon)e^{-z(\varepsilon)}}
              {S(\varepsilon,z(\varepsilon))}.
\]
The stationary boundary equation is
\begin{equation}\label{p1:eq:jet-stationarity}
 \frac1{z(\varepsilon)}-1
 -\frac{S_z(\varepsilon,z(\varepsilon))}
       {S(\varepsilon,z(\varepsilon))}=0.
\end{equation}
Its derivative with respect to $z$ at $(\varepsilon,z)=(0,1)$ is
$-1$, so it determines a unique convergent local power series for
$z(\varepsilon)$.

Put $\eta_n=H_n(1)$ and
$K_n=\overline{\mathbb Q}(e,\eta_1,\ldots,\eta_n)$, with
$K_0=\overline{\mathbb Q}(e)$. Every positive-order $z$-derivative of
$H_n$ at $z=1$ belongs to $K_{n-1}$: differentiate
$H_n'=H_{n-1}/z$ repeatedly, using that all derivatives of $H_0$
at $1$ belong to $K_0$. In particular, induction on the coefficient of
$\varepsilon^n$ in \eqref{p1:eq:jet-stationarity} shows that
\begin{equation}\label{p1:eq:root-jet-triangular}
 [\varepsilon^n]z(\varepsilon)\in K_{n-1}\qquad(n\ge1).
\end{equation}
Indeed the coefficient involving the new root coefficient is $-1$
times that coefficient, whereas every other term involves derivatives
of $H_k$ with $k\le n$, at least one derivative when $k=n$, and lower
root coefficients.

Next write
\[
 L(\varepsilon)=
 \log\frac{\mathcal J(2+2\varepsilon)}{2/e}
 =\log z(\varepsilon)-z(\varepsilon)+1
   -\log S(\varepsilon,z(\varepsilon)).
\]
The function $\log z-z+1$ has zero first derivative at $z=1$.
For each $n\ge1$, the coefficient of $\varepsilon^n$ in $L$ is
therefore
\begin{equation}\label{p1:eq:log-jet-triangular}
 [\varepsilon^n]L(\varepsilon)
 =(-1)^{n+1}\eta_n+R_n,\qquad R_n\in K_{n-1}.
\end{equation}
For completeness, the unique term containing $H_n(1)$ comes from the
linear term $-S+1$ in $-\log S$; nonlinear terms use only $H_k$ with
$k<n$. Substituting $z(\varepsilon)$ adds only the lower-field terms
specified in \eqref{p1:eq:root-jet-triangular}. The zero derivative of
$\log z-z+1$ prevents a new unrestricted root coefficient from entering.
Exponentiating \eqref{p1:eq:log-jet-triangular} gives
\[
 [\varepsilon^n]\mathcal J(2+2\varepsilon)
 =\frac2e\bigl((-1)^{n+1}\eta_n+\widetilde R_n\bigr),
 \qquad\widetilde R_n\in K_{n-1}.
\]
On the other hand this coefficient is
$2^n\mathcal J^{(n)}(2)/n!$. Starting with
$\mathcal J(2)=2/e$, the displayed triangular relations recursively
recover every $\eta_n$ from the dimension jet and vice versa. This
proves \eqref{p1:eq:jet-field-equality}.

Finally, the established algebraic independence of
$e,H_1(1),\ldots,H_N(1)$ follows from
\cite[Theorem~2]{p1:Powers2026}; its constants are
$1-e^{-1}$ and $n!H_n(1)$. It is also the fixed-$\alpha=0$, $z=1$
case of the integral family discussed in \cite[Section~7]{p1:Rivoal2012},
since
\[
 H_n(1)=\frac1{n!}\int_0^1(-\log t)^n e^{-t}\,\mathrm dt.
\]
The common field in \eqref{p1:eq:jet-field-equality} consequently has
transcendence degree $N+1$.
\end{proof}

This is one theorem about an infinite family of stochastic-optimization
sensitivities. Taking ten or more derivatives does not turn it into ten
independently posed historical questions. The first nontrivial member
is the exact Euler--Gompertz sensitivity proved above.


\begin{corollary}[No arithmetic algebraic differential equation]
\label{p1:cor:dimension-differential-transcendence}
The function $\mathcal J(d)=V_{d,2}(0,0)$ is differentially
transcendental over $\overline{\mathbb Q}(d)$: for no $N\ge0$ and
nonzero polynomial
$P\in\overline{\mathbb Q}(d)[X_0,\ldots,X_N]$ does
\[
 P(d,\mathcal J(d),\mathcal J'(d),\ldots,
                 \mathcal J^{(N)}(d))=0
\]
hold on an interval.
\end{corollary}
\begin{proof}
Clear the rational-function denominators. Divide the resulting
polynomial by the highest common power of $d-2$ dividing all its
coefficient polynomials. The relation still holds and its
specialization at $d=2$ is a nonzero polynomial relation over
$\overline{\mathbb Q}$ among a finite dimension jet, contradicting
Theorem~\ref{p1:thm:all-dimension-jets}. Real analyticity extends any
interval relation throughout the connected domain $d>0$ before this
evaluation.
\end{proof}

The coefficient field in this corollary is expressly
$\overline{\mathbb Q}(d)$. No claim over $\mathbb C(d)$ is inferred from
this specialization argument.

\section{Prime valuations, cyclotomic units, and rank certificates}
\label{p1:nt-section}

This section derives linear-independence certificates for generalized
Euler--Briggs constants from their character decomposition and Baker's
theorem.  The inputs are classical.  The results are presented as proved
consequences, without a claim of priority or a claim to have resolved a
published open conjecture.  In particular, the argument does not determine
the arithmetic nature of Euler's constant itself.

\subsection{The arithmetic decomposition}
\label{p1:nt-decomposition-section}

Fix a finite set $\Omega$ of rational primes and write
\[
 M_\Omega=\prod_{p\in\Omega}p,
 \qquad d_\Omega=\prod_{p\in\Omega}\left(1-\frac1p\right),
 \qquad C_\Omega=\sum_{p\in\Omega}\frac{\log p}{p-1}.
\]
Empty products and sums have their usual values $1$ and $0$.
For $q>1$, $(q,M_\Omega)=1$, and $1\le a\le q$ with $(a,q)=1$, define
\begin{equation}
\label{p1:nt-generalized-constant}
 G_\Omega(a,q)=
 \lim_{x\to\infty}\left(
 \sum_{\substack{n\le x\\ n\equiv a\pmod q\\(n,M_\Omega)=1}}\frac1n
 -\frac{d_\Omega}{q}\log x\right),
 \qquad
 T_\Omega(a,q)=\frac{q}{d_\Omega}G_\Omega(a,q).
\end{equation}
The limiting constant exists: the indicator of the summation set is
periodic, has mean $d_\Omega/q$, and has bounded discrepancy from this
mean.  Partial summation then gives convergence with error $O_{q,M_\Omega}(x^{-1})$.

\begin{lemma}[Character decomposition]
\label{p1:nt-character-decomposition}
With the preceding notation,
\begin{equation}
\label{p1:nt-normalized-decomposition}
 T_\Omega(a,q)=\gamma+C_\Omega+
 \sum_{p\mid q}\frac{\log p}{p-1}+U_\Omega(a,q),
\end{equation}
where
\begin{equation}
\label{p1:nt-u-definition}
 U_\Omega(a,q)=\frac{q}{d_\Omega\varphi(q)}
 \sum_{\substack{\chi\bmod q\\\chi\ne\chi_0}}
 \overline{\chi}(a)L(1,\chi)
 \prod_{p\in\Omega}\left(1-\frac{\chi(p)}p\right).
\end{equation}
In particular, $U_\Omega(a,q)$ is an algebraic linear combination of
nonprincipal Dirichlet $L$-values at $1$.
\end{lemma}

\begin{proof}
This is the normalization of formula~(2) in \cite{p1:GMS2016}.
For completeness, character orthogonality gives, for $\Re s>1$,
\[
 \sum_{\substack{n\ge1\\n\equiv a\pmod q\\(n,M_\Omega)=1}}n^{-s}
 =\frac1{\varphi(q)}\sum_{\chi\bmod q}
 \overline\chi(a)L(s,\chi)
 \prod_{p\in\Omega}(1-\chi(p)p^{-s}).
\]
The principal-character contribution equals
\[
 \frac{\zeta(s)}{\varphi(q)}
 \prod_{p\mid q}(1-p^{-s})\prod_{p\in\Omega}(1-p^{-s}).
\]
Its residue at $s=1$ is $d_\Omega/q$, and its constant term is
\[
 \frac{d_\Omega}{q}
 \left(\gamma+C_\Omega+
 \sum_{p\mid q}\frac{\log p}{p-1}\right).
\]
The other character contributions are holomorphic at $1$ and give the
terms in \eqref{p1:nt-u-definition} after normalization.  The bounded
discrepancy noted above identifies the constant term of the Dirichlet
series with the limit in \eqref{p1:nt-generalized-constant}, for example by
partial summation.  This proves the formula.
\end{proof}

\begin{lemma}[Prime logarithms modulo unit logarithms]
\label{p1:nt-prime-unit-separation}
Let $S$ be a finite set of distinct rational primes, and let
$u_1,\ldots,u_m$ be algebraic units with chosen complex logarithms.
If the coefficients $b_p,c_j,c_0$ are algebraic and at least one $b_p$
is nonzero, then
\begin{equation}
\label{p1:nt-prime-unit-form}
 \sum_{p\in S}b_p\log p+\sum_{j=1}^m c_j\log u_j+c_0\pi
\end{equation}
is transcendental.  Consequently, if this expression vanishes, every
$b_p$ vanishes.
\end{lemma}

\begin{proof}
Use $\log(-1)=i\pi$ to incorporate $c_0\pi$ into the unit-logarithm
sum.  Choose a basis $\ell_1,\ldots,\ell_r$ over $\mathbb Q$ for the
span of these finitely many unit logarithms, choosing the basis from the
logarithms themselves.  We first show that
\[
 \{\log p:p\in S\}\cup\{\ell_1,\ldots,\ell_r\}
\]
is linearly independent over $\mathbb Q$.
Clear denominators in a proposed rational relation and exponentiate.
The result has the form
\[
 \prod_{p\in S}p^{m_p}\,v=1,
\]
where $m_p\in\mathbb Z$ and $v$ is an algebraic unit.  Place all
these numbers in a number field $K$.  At a prime ideal $\mathfrak p$
of $K$ above a fixed $p\in S$, valuation of the displayed identity gives
$m_p v_{\mathfrak p}(p)=0$.  The other rational primes and every
algebraic unit have zero valuation there.  Hence every $m_p=0$,
and the remaining relation vanishes by the choice of the $\ell_j$.

Baker's theorem lifts the resulting rational linear independence of
logarithms of algebraic numbers to linear independence over
$\overline{\mathbb Q}$ \cite{p1:Baker1966}.  Thus
\eqref{p1:nt-prime-unit-form} is nonzero whenever some $b_p\ne0$.
The zero-or-transcendental part of the same theorem then makes it
transcendental.  One may equivalently use the formulation of Baker's
theorem recorded in \cite[Theorem~14]{p1:GMS2016}.
\end{proof}

\begin{lemma}[The unit contribution of nonprincipal characters]
\label{p1:nt-l-values-units}
For every nonprincipal Dirichlet character $\chi$, the number
$L(1,\chi)$ belongs to the $\overline{\mathbb Q}$-linear span of
$\pi$ and logarithms of algebraic units.
\end{lemma}

\begin{proof}
First suppose that $\chi$ is primitive of conductor $f$.  For odd
$\chi$, the finite formula for $L(1,\chi)$ expresses it as an
algebraic multiple of $\pi$.  For even $\chi$, the finite Fourier
formula expresses it as an algebraic multiple of
\[
 \sum_{a\in(\mathbb Z/f\mathbb Z)^*}
 \overline\chi(a)\log|1-\zeta_f^a|,
 \qquad \zeta_f=e^{2\pi i/f}.
\]
If $f$ has at least two distinct prime factors, then
$\Phi_f(1)=1$, so the algebraic integers $1-\zeta_f^a$ are units.
Their absolute values are also algebraic units: the square of such an
absolute value is the product of a unit with its complex conjugate,
and both the positive square root and its inverse are algebraic
integers.

If $f$ is a prime power, use
\[
 \log|1-\zeta_f^a|=
 \log|1-\zeta_f|+
 \log\left|\frac{1-\zeta_f^a}{1-\zeta_f}\right|.
\]
The quotients are cyclotomic units.  Since $\chi$ is nonprincipal,
$\sum_a\overline\chi(a)=0$, and the common nonunit logarithm
disappears from the character sum.  This proves the primitive case.

Finally, if $\chi$ modulo $q$ is induced by a primitive character
$\chi^*$ of conductor $f$, then $\chi^*$ is nonprincipal and
\[
 L(1,\chi)=L(1,\chi^*)
 \prod_{\substack{p\mid q\\p\nmid f}}
 \left(1-\frac{\chi^*(p)}p\right).
\]
The finite factor is algebraic, which proves the general case.
These are the same standard cyclotomic-unit inputs used in
\cite[Section~3]{p1:GMS2016}.
\end{proof}

\subsection{Coprime moduli and transcendental differences}
\label{p1:nt-coprime-section}

\begin{theorem}[A coprime-modulus certificate]
\label{p1:nt-coprime-certificate}
Let $q_1,\ldots,q_n>1$ be pairwise coprime, with
$(q_i,M_\Omega)=1$, and let $a_i$ be any reduced residue modulo
$q_i$.  Put $T_i=T_\Omega(a_i,q_i)$.  If
$c_1,\ldots,c_n\in\overline{\mathbb Q}$ are not all zero and
$\sum_i c_i=0$, then
\[
 \sum_{i=1}^n c_iT_i
\]
is transcendental.
\end{theorem}

\begin{proof}
By Lemma~\ref{p1:nt-character-decomposition},
\[
 T_i=\gamma+C_\Omega+
 \sum_{p\mid q_i}\frac{\log p}{p-1}+U_\Omega(a_i,q_i).
\]
The condition $\sum_i c_i=0$ cancels $\gamma+C_\Omega$.
Lemma~\ref{p1:nt-l-values-units} puts all remaining $U$-terms into
a finite algebraic linear combination of unit logarithms and $\pi$.
Because the moduli are pairwise coprime, a prime $p\mid q_i$ occurs
in no other modulus.  Its coefficient is therefore $c_i/(p-1)$.
Some $c_i$ is nonzero, and that $q_i>1$ has a prime divisor.
Lemma~\ref{p1:nt-prime-unit-separation} proves the assertion.
\end{proof}

\begin{corollary}[Dimension and the relation space]
\label{p1:nt-coprime-rank}
Under the hypotheses of Theorem~\ref{p1:nt-coprime-certificate},
\begin{equation}
\label{p1:nt-rank-bound}
 \dim_{\overline{\mathbb Q}}
 \operatorname{span}_{\overline{\mathbb Q}}
 \{G_\Omega(a_i,q_i):1\le i\le n\}\ge n-1.
\end{equation}
The space of algebraic linear relations among $T_1,\ldots,T_n$
has dimension at most one.  Moreover,
\[
 1,\ T_2-T_1,\ldots,T_n-T_1
\]
are linearly independent over $\overline{\mathbb Q}$.
An infinite family of pairwise coprime admissible moduli consequently
gives an infinite-dimensional span over $\overline{\mathbb Q}$.
\end{corollary}

\begin{proof}
Let
\[
 R=\left\{c\in\overline{\mathbb Q}^{\,n}:
 \sum_i c_iT_i=0\right\}.
\]
The map $R\longrightarrow\overline{\mathbb Q}$ given by
$c\longmapsto\sum_i c_i$ is injective by
Theorem~\ref{p1:nt-coprime-certificate}.  Hence $\dim R\le1$, and
rank--nullity proves the bound for the $T_i$.  Multiplication by the
nonzero rational factors $d_\Omega/q_i$ gives
\eqref{p1:nt-rank-bound}.

For the stronger affine statement, suppose
$b+\sum_{i=2}^n b_i(T_i-T_1)=0$ with algebraic coefficients.
If any $b_i\ne0$, the sum of differences is transcendental by
Theorem~\ref{p1:nt-coprime-certificate}, whereas the equation makes
it the algebraic number $-b$.  Thus every $b_i=0$ and then $b=0$.
Applying the rank bound to arbitrarily long finite subfamilies proves
the infinite-dimensional assertion.
\end{proof}

For a concrete normalization check, when $\Omega$ is empty one has
\[
 2G_\varnothing(1,2)=\gamma+\log2,
 \qquad
 3G_\varnothing(1,3)=\gamma+\frac12\log3+
 \frac{\pi}{2\sqrt3}.
\]
Their difference is therefore
\[
 3G_\varnothing(1,3)-2G_\varnothing(1,2)
 =\frac12\log3-\log2+\frac{\pi}{2\sqrt3},
\]
a transcendental number by the theorem.  This is an illustration of
the certificate, not a new claim about a previously unresolved named
constant.

\subsection{A prime-divisibility incidence matrix}
\label{p1:nt-incidence-section}

\begin{theorem}[Incidence rank without coprimality]
\label{p1:nt-incidence-certificate}
Let $q_1,\ldots,q_n>1$ satisfy $(q_i,M_\Omega)=1$, without assuming
pairwise coprimality.  Choose reduced residues $a_i$ and put
$T_i=T_\Omega(a_i,q_i)$.  Let $S$ be the set of primes dividing
$q_1\cdots q_n$, and define the matrix
\[
 A_{p i}=\frac{\mathbf1_{p\mid q_i}}{p-1}
 \quad(p\in S,\ 1\le i\le n),
 \qquad
 H=\left\{c\in\overline{\mathbb Q}^{\,n}:\sum_i c_i=0\right\}.
\]
For every $c\in H$ with $Ac\ne0$, the number
$\sum_i c_iT_i$ is transcendental.  Consequently,
\begin{equation}
\label{p1:nt-incidence-rank-bound}
 \dim_{\overline{\mathbb Q}}
 \left\{\sum_i c_iT_i:c\in H\right\}
 \ge \operatorname{rank}_{\overline{\mathbb Q}}(A|_H).
\end{equation}
Thus a finite prime-divisibility rank calculation supplies an
arithmetic linear-independence certificate.
\end{theorem}

\begin{proof}
For $c\in H$, the decomposition
\eqref{p1:nt-normalized-decomposition} cancels $\gamma+C_\Omega$ and
leaves the coefficient vector $Ac$ for the rational-prime logarithms.
All other terms are unit logarithms and $\pi$ by
Lemma~\ref{p1:nt-l-values-units}.  If $Ac\ne0$,
Lemma~\ref{p1:nt-prime-unit-separation} makes the resulting number
transcendental.

Write $r=\operatorname{rank}(A|_H)$ and choose
$c^{(1)},\ldots,c^{(r)}\in H$ such that their images under $A$
form a basis of $A(H)$.  Any nonzero algebraic linear combination of
these vectors has nonzero image under $A$, and its corresponding
combination of the $T_i$ is therefore nonzero.  Their $r$ values are
linearly independent, proving \eqref{p1:nt-incidence-rank-bound}.
\end{proof}

The restriction to coefficients cancelling $\gamma$ is substantive.
The method identifies prime-logarithm components by valuations, but
it supplies no corresponding independence assertion for a linear
form retaining $\gamma$.  For this reason, even the infinite-dimensional
family conclusion does not establish the transcendence of a specified
individual constant.

\subsection{Position relative to the external questions}
\label{p1:nt-position}

The character decomposition above is established prior work
\cite{p1:GMS2016}; the proofs combine it with Baker's theorem
\cite{p1:Baker1966} and elementary valuations.  In particular,
\cite[Theorems~8--9]{p1:GMS2016} already studies algebraic-multiple
classes and infinite-dimensional spaces of these constants.  No
assertion that the present rank formulations were previously open is
made here.

The analytic questions in the accompanying ledger remain separate.
Diamond and Ford identify the spectral equality
$\inf_\Omega\gamma(\Omega)=e^{-\gamma}$ with the Riemann hypothesis,
and already disprove monotonicity along initial prime sets
\cite{p1:DiamondFord2008}.  The cyclotomic Euler--Kronecker limit-point
conjecture of Ford, Luca, and Moree \cite{p1:FordLucaMoree2014}, and the
unconditional infinitude of negative cyclotomic Euler--Kronecker
constants discussed by Dixit and Murty \cite{p1:DixitMurty2023}, are
not resolved by any result in this section.

\section{Arithmetic Gevrey transfer: an exact obstruction}

Put
\[
 \operatorname{Ein}(z)=\sum_{n=1}^{\infty}\frac{(-1)^{n+1}z^n}{n\,n!},
 \qquad \eta=\operatorname{Ein}(1)=\gamma+\frac{\delta}{e},
 \qquad \delta=\int_0^\infty\frac{e^{-t}}{1+t}\,dt.
\]
The algebraic independence of $e$ and $\eta$ is established in Rivoal's
2012 paper \cite[Section 5]{p1:Rivoal2012}; it is not a result first
proved here. We use $\overline{\mathbb Q}$ for the algebraic numbers.

The relevant open statements are recorded in the source ledger;
we use the mixed-function framework of \cite{p1:FischlerRivoal2024}.
The coefficient-comparison argument below specializes the mechanism
of that paper's Proposition 2. It is included as an exact obstruction,
not as an independently new resolution.

\subsection{Definitions needed for the obstruction}
An $E$-function has an expansion $\sum_{n\geq0}a_nz^n/n!$ with algebraic
coefficients satisfying the standard exponential conjugate and common-
denominator bounds, and satisfies a linear differential equation over
$\overline{\mathbb Q}(z)$. An $\mathfrak E_1$-function is a formal
series $\sum_{n\geq0}n!a_nw^n$ whose associated $\sum a_nw^n$ is a
$G$-function. A mixed function is a formal sum
\[
 A(z)+\sum_{n\geq1}b_nz^{-n},
\]
where $A$ is an $E$-function and the negative-power part is an
$\mathfrak E_1$-function in $1/z$. In particular, all $b_n$ are
algebraic. Equality of mixed functions means coefficientwise equality;
multiplication by $z-1$ is consequently unambiguous.

Define
\[
 \widehat D(z)=\sum_{n\geq0}(-1)^n n!z^{-n-1},\qquad
 W_r(z)=e^z\bigl(\operatorname{Ein}(z)-r\bigr)-\widehat D(z)
 \quad(r\in\overline{\mathbb Q}).
\]
The series $\widehat D$ is of the required arithmetic Gevrey type. Indeed,
its ordinary Borel transform is $1/(1+t)$; alternatively the coefficient
sequence after dividing the coefficient of $z^{-m}$ by $m!$ is
$(-1)^{m-1}/m$, a $G$-function sequence. The positive part of $W_r$ is an
$E$-function because this class is closed under products and algebraic
linear combinations.

\begin{proposition}[Arithmetic nondivisibility]
For every algebraic $r$, there is no mixed function $V$ satisfying
\[
 W_r(z)=(z-1)V(z).
\]
\end{proposition}
\begin{proof}
Suppose $V=A(z)+\sum_{n\geq1}b_nz^{-n}$ is mixed. Comparing the
nonnegative powers of $z$ gives
\[
 e^z\bigl(\operatorname{Ein}(z)-r\bigr)=(z-1)A(z)+b_1.
\]
Since $A$ is entire, evaluation at $z=1$ implies
$e(\eta-r)=b_1\in\overline{\mathbb Q}$. This contradicts the algebraic
independence of $e$ and $\eta$: the polynomial
$X(Y-r)-b_1$ is nonzero in $\overline{\mathbb Q}[X,Y]$.
\end{proof}

\begin{proposition}[An exact threshold for a restricted zero-removal rule]
The following statements are equivalent:
\begin{enumerate}
\item $\gamma$ is transcendental.
\item For every $r\in\overline{\mathbb Q}$, the implication
\[
 (W_r)_0(1)=0\quad\Longrightarrow\quad
 W_r=(z-1)V\text{ for some mixed }V
\]
holds.
\end{enumerate}
Replacing $\overline{\mathbb Q}$ by $\mathbb Q$ in statement (2) makes it
equivalent to the irrationality of $\gamma$.
\end{proposition}
\begin{proof}
On the positive real axis the Borel--Laplace sum is
\[
 D_0(z)=\int_0^\infty\frac{e^{-zt}}{1+t}\,dt
       =e^z E_1(z),\qquad
 E_1(z)=\int_z^\infty\frac{e^{-u}}u\,du.
\]
The classical identity
$\operatorname{Ein}(z)=\gamma+\log z+E_1(z)$, valid for $z>0$, yields
\[
 (W_r)_0(z)=e^z(\gamma+\log z-r),\qquad
 (W_r)_0(1)=e(\gamma-r).
\]
If $\gamma$ is transcendental, the antecedent in (2) is false for every
algebraic $r$, so the implication is true. Conversely, if $\gamma$ is
algebraic, choose $r=\gamma$. Its antecedent then holds, but its conclusion
contradicts the preceding proposition. The rational version is identical.
\end{proof}

This establishes an exact limitation of the proposed analytic-to-arithmetic
transfer. A removable analytic zero does not automatically preserve
arithmetic coefficient conditions. For this explicit family, establishing
that preservation is equivalent to establishing the desired arithmetic
nature of $\gamma$. This is an equivalence and obstruction theorem, not
a proof of the conjecture. It must not be counted as a solved previously
open problem.

\subsection{A field-theoretic constraint on coefficient searches}
Let $K$ denote the algebraic closure of $\mathbb Q(e)$ inside $\mathbb C$.
The known algebraic independence of $e,\eta$ implies $\eta\notin K$.

\begin{proposition}[At most one exceptional coefficient]
There is at most one $a\in K$ for which $\gamma+a\delta\in K$.
Moreover $a=1/e$ is not exceptional. If an exceptional $a_0$ exists, then
$\delta\notin K$ and, for every $a\in K\setminus\{a_0\}$,
$\gamma+a\delta$ is transcendental over $K$.
\end{proposition}
\begin{proof}
If $a\ne b$ and both combinations belong to $K$, subtraction gives
$\delta\in K$, and then $\gamma\in K$. Consequently
$\eta=\gamma+\delta/e\in K$, a contradiction. The identity
$\gamma+\delta/e=\eta\notin K$ excludes $1/e$. If an exceptional $a_0$
exists and $\delta$ were in $K$, the same contradiction would follow.
For $a\ne a_0$, the first assertion excludes membership in $K$; since
$K$ is algebraically closed, this is equivalent to transcendence over $K$.
\end{proof}

This elementary consequence considerably narrows the possible exceptional
coefficient set, but does not identify it. In particular, the method does
not rule out $a_0=0$, which is exactly the obstruction to deciding whether
$\gamma$ is algebraic over $\mathbb Q(e)$. No novelty or open-problem
resolution is asserted for this field-theoretic deduction.


\section{What moment positivity can and cannot decide}
\label{p1:sec:gumbel-moments}

The Gumbel distribution provides a direct probability model for $\gamma$.
The following exact translation calculation explains a limitation of
an arithmetic strategy that fixes the higher cumulants and tests only
moment positivity.

Let $T$ have density $e^{-t}$ on $(0,\infty)$ and put $X=-\log T$.
For real $t<1$,
\[
 \mathbb E e^{tX}=\Gamma(1-t),\qquad
 \log\Gamma(1-t)=\gamma t+\sum_{k\ge2}\frac{\zeta(k)}k t^k
 \quad(|t|<1).
\]
Thus its cumulants are $\kappa_1=\gamma$ and
$\kappa_k=(k-1)!\zeta(k)$ for $k\ge2$.
For an arbitrary real number $g$ define $X_g=X+g-\gamma$ and
$\mu_j(g)=\mathbb E X_g^j$. Its first cumulant is $g$, while all
higher cumulants are unchanged.

\begin{theorem}[Translation invariance of the leading Hankel determinants]
For $N\ge0$, put $M_N(g)=(\mu_{i+j}(g))_{0\le i,j\le N}$.
Every $M_N(g)$ is positive definite, for every $g\in\R$. Moreover,
\[
 \det M_N(g)=\det M_N(\gamma)>0.
\]
Consequently the moment-positivity conditions with these fixed higher
cumulants cannot distinguish a rational, algebraic irrational, or
transcendental value of the first cumulant.
\end{theorem}
\begin{proof}
Set $h=g-\gamma$ and let $T_N(h)$ be the lower triangular matrix
whose $(i,j)$ entry is $\binom ijh^{i-j}$ for $j\le i$, and zero otherwise.
The binomial theorem gives
\[
 (1,X_g,\ldots,X_g^N)^{\mathsf T}
 =T_N(h)(1,X,\ldots,X^N)^{\mathsf T}.
\]
Taking Gram matrices yields
$M_N(g)=T_N(h)M_N(\gamma)T_N(h)^{\mathsf T}$.
All diagonal entries of $T_N(h)$ equal one, so its determinant is one,
which proves determinant invariance. For a nonzero real polynomial $p$
of degree at most $N$, the continuous everywhere-positive density of
$X_g$ implies $\mathbb E p(X_g)^2>0$. This proves positive definiteness.
The construction works for every real $g$, proving the final assertion
within the stated class of tests.
\end{proof}

For example, the first two nontrivial determinants are
\[
 \det M_1=\zeta(2),\qquad
 \det M_2=2\zeta(2)^3+6\zeta(2)\zeta(4)-4\zeta(3)^2>0.
\]
The statement is not that every minor is invariant: the one-entry minor
$\mu_2(g)=g^2+\zeta(2)$ already varies with $g$. Nor does it exclude
arithmetic information from integral representations, denominator
control, or additional identities tying the mean to the higher
cumulants. It excludes only an inference from positivity alone after
free translation of the first cumulant.


\section{What remains to reach the arithmetic target}

The strongest transfers in this manuscript use structures that survive
exactly: a Wronskian identity survives as an integral Laurent recurrence;
a radial heat kernel exposes the dominant sign of asymptotic coefficients;
and a stationary optimization formula transports algebraic independence
of $E$-values to parameter derivatives. Each transfer includes the relevant
normalization and its proof.

None produces a nonzero integer linear form in $1,\gamma$ tending to
zero. The positive-filter theorem identifies a quantitative denominator
obstruction for one broad approximation family. The mixed Gevrey theorem
identifies a restricted zero-removal statement that is already equivalent
to the arithmetic target. The Gumbel calculation shows why moment
positivity alone cannot supply the missing arithmetic input.

For future work, three concrete mathematical gaps remain useful targets:
\begin{enumerate}
\item Find and prove denominator cancellation beyond factorial clearing
for sign-changing or multivariate Euler--Gompertz filters, while retaining
a strictly nonzero, sufficiently small remainder.
\item Determine whether the exact BGG coefficient recurrences admit
additional congruences or root-location theorems that constrain rational
approximants, rather than only the asymptotic series coefficients.
\item Extend the stochastic parameter-jet correspondence to other
$E$-function families with established functional independence, and
separately determine whether any singular parameter limit preserves
enough arithmetic information to reach $\gamma$ itself.
\end{enumerate}
These are research directions, not claims that the remaining conjectures
have been settled. The present document is a research draft with complete
arguments for its stated results and reproducible finite certificates;
it has not received external peer review. No journal-level or prize-level
assessment follows from the number of its theorems.

\subsection{Verification and reproducibility}
All source targets were checked against the primary text, with the
coefficient conventions retained. The finite sign certificates use exact
integer or rational arithmetic and are embedded in this document, so the
single accompanying \texttt{.tex} file is self-contained. For the product
coefficients, the scalar certificate was independently compared with the
nonlinear Laurent recurrence through degree $143$. Both methods gave
identical integers. Symbolic recurrence proofs, rather than finite tests,
establish the denominator and congruence assertions at every degree.

The mathematical arguments were separately checked in parallel research
passes. Such checks are useful error detection, not external validation.
The source ledger deliberately separates numbered conjectures, unnumbered
observations, recovered project results, and new extensions. The requested
ten new solutions of existing open questions have not been established.


\small
\begin{thebibliography}{VIII.99}
\bibitem[I.1]{p1:BrentGlasserGuttmann2019}
R. P. Brent, M. L. Glasser, and A. J. Guttmann,
\emph{A conjectured integer sequence arising from the exponential integral},
Journal of Integer Sequences \textbf{22} (2019), Article 19.4.7.
\url{https://arxiv.org/abs/1812.00316}.

\bibitem[I.2]{p1:NowakStempak2016}
A. Nowak and K. Stempak,
\emph{Sharp estimates for potential operators associated with Laguerre
and Dunkl--Laguerre expansions}, preprint arXiv:1402.2522 (2014),
equation (1) and Section 2.2.
\url{https://arxiv.org/abs/1402.2522}.

\bibitem[I.3]{p1:DLMF}
NIST Digital Library of Mathematical Functions,
Section 10.40, asymptotic expansions of modified Bessel functions;
Section 13, confluent hypergeometric functions.
\url{https://dlmf.nist.gov/10.40}.

\bibitem[I.4]{p1:MurabayashiYoshida2024}
N. Murabayashi and H. Yoshida,
\emph{Rational function approximations of the special function
$e^xE_1(x)$ and applications to irrationality of Euler--Gompertz constant
$\delta$}, arXiv:2210.06768v2, revised 22 September 2024.
\url{https://arxiv.org/abs/2210.06768}.

\bibitem[I.5]{p1:bz2026}
T. Bulboac\u a and H. M. Zayed,
\emph{Monotonic nature of the Gamma function},
Journal of Inequalities and Applications \textbf{2026}, Article 27,
published 22 January 2026.
\url{https://doi.org/10.1186/s13660-025-03425-0}.

\bibitem[I.6]{p1:Beukers2006}
F. Beukers, \emph{A refined version of the Siegel--Shidlovskii theorem},
Annals of Mathematics (2) \textbf{163} (2006), 369--379.
\url{https://doi.org/10.4007/annals.2006.163.369}.

\bibitem[I.7]{p1:EkstromWanntorp2009}
E. Ekstr\"om and H. Wanntorp,
\emph{Optimal stopping of a Brownian bridge},
Journal of Applied Probability \textbf{46} (2009), 170--180.

\bibitem[I.8]{p1:HobsonLiu2026}
D. Hobson and J. Liu,
\emph{Optimal stopping of a Brownian excursion and an
$\alpha$-dimensional Bessel bridge}, Journal of Applied Probability,
published online 15 June 2026, 1--18.
\url{https://doi.org/10.1017/jpr.2026.10102}.

\bibitem[I.9]{p1:Rivoal2012}
T. Rivoal, \emph{On the arithmetic nature of the values of the gamma
function, Euler's constant, and Gompertz's constant},
Michigan Mathematical Journal \textbf{61} (2012), 239--254.
\url{https://rivoal.perso.math.cnrs.fr/articles/gammater.pdf}.

\bibitem[I.10]{p1:Powers2026}
M. R. Powers, \emph{Transcendence results for $\Gamma^{(n)}(1)$ and
related sequences of generalized constants}, arXiv:2511.01849v7
(2026), Theorem 2.
\url{https://arxiv.org/abs/2511.01849v7}.

\bibitem[I.11]{p1:Lagarias2013}
J. C. Lagarias, \emph{Euler's constant: Euler's work and modern
developments}, Bulletin of the American Mathematical Society
\textbf{50} (2013), 527--628.
\url{https://arxiv.org/abs/1303.1856}.

\bibitem[I.12]{p1:FischlerRivoal2024}
S. Fischler and T. Rivoal,
\emph{Relations between values of arithmetic Gevrey series, and
applications to values of the Gamma function},
arXiv:2301.13518 (2023), Conjectures 1--3 and Proposition 2.
\url{https://arxiv.org/abs/2301.13518}.

\bibitem[I.13]{p1:GMS2016}
S. Gun, V. Kumar Murty, and E. Saha,
\emph{Linear and algebraic independence of generalized Euler--Briggs
constants}, Journal of Number Theory \textbf{166} (2016), 117--136.
\url{https://arxiv.org/abs/1604.02896}.

\bibitem[I.14]{p1:Baker1966}
A. Baker, \emph{Linear forms in the logarithms of algebraic numbers},
Mathematika \textbf{13} (1966), 204--216.
The formulation used here is also recorded in
\cite[Theorem 14]{p1:GMS2016}.

\bibitem[I.15]{p1:DiamondFord2008}
H. G. Diamond and K. Ford, \emph{Generalized Euler constants},
Mathematical Proceedings of the Cambridge Philosophical Society
\textbf{145} (2008), 27--41.
\url{https://arxiv.org/abs/math/0703508}.

\bibitem[I.16]{p1:FordLucaMoree2014}
K. Ford, F. Luca, and P. Moree,
\emph{Values of the Euler phi-function not divisible by a given odd
prime, and the distribution of Euler--Kronecker constants for cyclotomic
fields}, Mathematics of Computation \textbf{83} (2014), 1447--1476.
\url{https://arxiv.org/abs/1108.3805}.

\bibitem[I.17]{p1:DixitMurty2023}
A. Dixit and M. R. Murty, \emph{On Ihara's conjectures for
Euler--Kronecker constants}, Acta Arithmetica \textbf{210} (2023), 95--123.
\url{https://doi.org/10.4064/aa220729-19-3}.

\bibitem[I.18]{p1:pudim2026}
Pudim project, public application ledger,
\emph{APPLICATIONS.md}, entry APP-0009, dated 23 May 2026;
prior public claim concerning the constant-numerator shift,
not used as a proof input.
\url{https://github.com/pudim-project/pudim-ai-demo-zetalaw/blob/main/APPLICATIONS.md}.

\bibitem[I.19]{p1:priorBGG}
KEIO project record,
\emph{BGG Conjecture 10: Wronskian--Laurent proof},
private working manuscript, 26 August 2026. Retrieved and checked
in the present project; records the integrality argument.

\bibitem[I.20]{p1:priorTransfer}
KEIO project record,
\emph{Ten Methods Transfer Report}, 23 September 2026,
private working manuscript. Retrieved and checked in the present project;
records the discrete integrality and mod-$32$ proofs.
\end{thebibliography}
% END PRESERVED SOURCE BODY 1
\phantomsection\label{part1:end}
\clearpage
\setcounter{section}{0}\setcounter{subsection}{0}
\renewcommand{\thesection}{\Roman{part}.\arabic{section}}
\renewcommand{\theHsection}{\arabic{part}.main.\arabic{section}}
\renewcommand{\refname}{References for Part II}
\fancyhead[L]{\small Euler\textquotesingle s constant: Part II}
\part{Arithmetic transfers II: units, parameter jets, and source assertions}
\label{part2:start}
\noindent\textbf{Source:} {\small\path{KEIO_Euler_Arithmetic_Frontiers_20261018.tex}}.

Original PDF: 35 pages. Mathematical content and source-level qualifications are retained.


% BEGIN PRESERVED SOURCE BODY 2
\begin{abstract}
We classify which positive rational-parameter Kummer values at positive
algebraic arguments are units in the ring of $E$-values, using the recent
phase morphism of Fischler--Rivoal. At argument one the only cases are
the exponential-polynomial family and $M(2,3,1)=2$. The same morphism
gives a denominator obstruction for
$\Ein(1)=\gamma+\delta/e$. Independent transfers give algebraically
independent dimension derivatives of a Bessel stopping boundary and
differential transcendence over $\mathbb C(d)$.
We also supply proofs or counterexamples for ten distinct assertions
left unproved in precisely identified source versions: a Whittaker
real-zero branch conjecture, Vargas's exterior-zero conjecture,
G\"orgens' local-minimum conjecture and two stated endpoint expectations,
a sharp digamma exponent question, two Gurland-ratio questions,
Bolbachan's digamma-limit conjecture, and an odd-totient sign question.
The last limit conjecture fails exponentially, with growth rate
$1/\log2$. Further results sharpen periodic-Dirichlet descent bounds.
The source count includes minor questions and preprints; it is not a
claim of ten major or historically first breakthroughs. Current priority
is not certified, especially where a revised journal version could not
be inspected. Neither irrationality nor transcendence of $\gamma$ is
proved.
\end{abstract}
\noindent\textbf{Keywords:} Euler--Mascheroni constant; Euler--Gompertz constant;
$E$-function; Kummer function; optimal stopping; differential transcendence;
digamma inequality; exponential-integral sections; periodic Dirichlet series.

\section{What is proved, and what the count means}
\label{p2:sec:round2-ledger}
Write
\[
 \gamma=\lim_{n\to\infty}\left(\sum_{k=1}^n\frac1k-\log n\right),
 \qquad \delta=\int_0^\infty\frac{e^{-t}}{1+t}\,dt,
 \qquad \eta=\Ein(1)=\gamma+\delta/e.
\]
This second research stage starts from the arithmetic of these constants
and transfers to special-function differential equations, stochastic
stopping, complex zeros, convexity, and periodic Dirichlet series.
The strongest arithmetic statements concern the \emph{ring} containing
$\eta$ and the finite jets of optimized values. They do not separate
the arithmetic of $\gamma$ from that of $\delta$.

The following ledger contains ten distinct mathematical targets stated
as conjectures, open questions, or unproved expectations in the cited
source versions. For each target the manuscript gives a full proof or
counterexample. Two equivalent convergence formulations in one source
are counted once. Parameter cases and endpoints of one conjecture count once; corollaries
count only when separately posed in a source.
No result from the preceding research stage is reused to increase this
count.

\small
\begin{longtable}{@{}p{.045\textwidth}p{.43\textwidth}p{.46\textwidth}@{}}
\toprule ID & Exact source target & Answer and proof location \\\midrule
\endfirsthead
\toprule ID & Exact source target & Answer and proof location \\\midrule
\endhead
Q1 & Boussaada--Mazanti--Niculescu \cite[p.~6, after Fig.~1]{p2:BMN2022}:
maximal domain and endpoint limits of a real Whittaker zero branch.
& Proved by a Sturm identity; a sharper Bessel-scaled endpoint is obtained.
Theorem~\ref{p2:thm:whittaker-branch}. \\
Q2 & Vargas \cite[Conjecture~5.1]{p2:Vargas2014}: all finite limits of
normalized Taylor-section zeros lie on the proposed curves.
& False for the stated signed/complex-integrand class. Any prescribed
point outside the unit disk can occur. Theorem~\ref{p2:thm:vargas-counterexample}. \\
Q3 & G\"orgens \cite[Section~4]{p2:Gorgens2014}: the observed local minimum
occurs exactly at $\alpha=1/2$.
& Proved by a two-sided expansion with a positive quadratic coefficient.
Theorem~\ref{p2:thm:gorgens-minimum}. \\
Q4 & G\"orgens \cite[Section~4]{p2:Gorgens2014}: the optimal stopping time
tends to $1$ almost surely as $\alpha\downarrow0$.
& Proved under the natural common-Brownian coupling.
Theorem~\ref{p2:thm:gorgens-limits}. This is an elementary stated expectation. \\
Q5 & G\"orgens \cite[Introduction]{p2:Gorgens2014}: the value tends to zero
as $\alpha\to\infty$.
& Proved, with the stronger rate
$V(\alpha)\sim\sqrt{\log\alpha/\alpha}$ and second-order squared asymptotics.
Theorems~\ref{p2:thm:gorgens-limits} and \ref{p2:thm:gorgens-large-drift}.
One expectation, not two targets. \\
Q6 & Qi--Lim--Nantomah \cite[Open Problem~3]{p2:qln2025}: the largest
exponent range for a reciprocal digamma inequality.
& Exactly $[-1,B)$, where $B\approx5.97283686$ is characterized by a
proved unique tangency point. Theorem~\ref{p2:thm:qln-op3}. \\
Q7 & Wi\'sniewska \cite[Open Problems~1 and~2]{p2:Wis2025}: global convergence
and equality of the Gurland-ratio series.
& Both equivalent assertions hold for all positive arguments, with a
geometric remainder bound. Section~\ref{p2:sec:gurland2}; counted once. \\
Q8 & Wi\'sniewska \cite[Open Problem~3]{p2:Wis2025}: location of the
mean-value parameter and its distance from the midpoint.
& Sharp bound $t>g+\frac23(s-g)$, plus a sharper ratio-dependent bound
and exact scaling limits. Theorem~\ref{p2:thm:gurland-location}. \\
Q9 & Bolbachan \cite[Conjecture~1.4]{p2:Bolbachan2011}: a Bernoulli--Stirling
limit represents the digamma function.
& False as printed: the approximants diverge, with exponential rate
$1/\log2$. Theorem~\ref{p2:thm:bol-counterexample}. \\
Q10 & Bharadwaj \cite[Section~8, pp.~24--25]{p2:Bharadwaj2022}: existence of
an odd $N$ with a signed totient sum equal to zero.
& Explicit nontrivial answer $N=105$, also the least odd $N>1$.
Section~\ref{p2:r2nt-section}. \\
\bottomrule
\end{longtable}
\normalsize

\paragraph{The distinction between a proof and a priority claim.}
The ledger is a count of resolved \emph{source assertions}, not a
certification that ten problems were still globally open immediately
before this work. The proofs were checked internally against the stated
hypotheses; they have not undergone external peer review or formal
proof-assistant verification. Searches and source comparisons did not
locate earlier resolutions of the listed exact assertions, but cannot
establish historical priority. In particular, the revised journal text
of Bharadwaj's paper was not accessible; our question match and numerical
comparison concern its 2022 arXiv version. Bolbachan's source is a 2011
preprint prepared for a high-school mathematical conference. G\"orgens'
two endpoint expectations and the totient example are modest results.
Thus this document does not certify the stronger request for ten new
major breakthroughs, or an Annals/Fields-level achievement.

\paragraph{Additional results excluded from that count.}
The complete positive Kummer unit classification
(Theorem~\ref{p2:thm:positive-kummer-units}) is a family-wide verification
of a larger units conjecture, not a solution of the full conjecture.
The rational-denominator obstruction
(Theorem~\ref{p2:thm:ein-denominators}), boundary-jet algebraic
independence, and differential-transcendence statements are additional
transfers. The adaptive descent bound improves an explicitly suggested
direction rather than closing the full Erd\H{o}s conjecture.
The higher-polygamma range in Section~\ref{p2:sec:qln-op4-independent}
has an earlier public claim and is excluded. The value-jet proposition
is repeated only to make the later proof self-contained.

\paragraph{How the fields connect.}
The phase morphism converts a special value into a finite exponential
polynomial, so algebraic factorization detects forbidden reciprocals.
Kummer differential equations also encode stopping boundaries; parameter
expansion then exposes a tower of arithmetically independent constants.
Sturm identities control real zeros, while reversing Taylor sections
reveals complex zeros missed by a proposed limiting curve. Positivity
and curvature settle the sharp inequality questions. In the periodic
Dirichlet setting, mean zero cancels $\gamma$ and leaves exact
prime-logarithm obstructions. Each transfer has a different conclusion;
none is an implicit proof of the original irrationality problem.



% Source fragment: round2_arithmetic.tex
% Standalone mathematical fragment; requires amsmath,amssymb,amsthm,hyperref.
% Independent second-round work.  Baseline BGG results are deliberately absent.
\section{Units throughout a positive Kummer family}
\label{p2:sec:kummer-units-round2}

Let $\mathcal E$ be the ring of values of Siegel $E$-functions at algebraic
arguments.  Fischler and Rivoal, \emph{The ring of values of $E$-functions is
not a field}, arXiv:2610.16603v1 (12 October 2026; manuscript dated
14 October 2026), Conjecture~1.2, propose
\[
 \mathcal E^\times=\{\alpha e^\beta:
     \alpha,\beta\in\overline{\mathbb Q},\ \alpha\ne0\}.
\]
Their Theorem~1.1 supplies nonunits $M(a,b,1)$ when
$0<a<1$ and $0<b-a<1$ are rational.  Here
\[
 M(a,b,z)={}_1F_1(a;b;z)
       =\sum_{j\ge0}\frac{(a)_j}{(b)_j j!}z^j.
\]
We classify units in this family for \emph{every} positive rational
$a,b$ and every positive algebraic argument.  This is a restricted-family
result concerning Conjecture~1.2, not a proof of that conjecture.
Rivoal--Roques, \emph{$E$-functions of order 2 and units of $E$-values}
(2015), Theorem~5, does not already give this exclusion: it requires the
function and the witness for its inverse value to be components of one
joint two-dimensional differential system.  Here an inverse is excluded
even if its possible $E$-function witness has arbitrarily large order.

\begin{theorem}[Positive Kummer unit classification]
\label{p2:thm:positive-kummer-units}
Let $a,b\in\mathbb Q_{>0}$ and
$x\in\overline{\mathbb Q}\cap\mathbb R_{>0}$.  For $n\ge1$ put
\[
 A_{n,b}(z)=\frac1{(n-1)!}
  \sum_{j=0}^{n-1}
  \frac{(b-n)_j(1-n)_j}{j!}z^{n-1-j}.
\]
Then $M(a,b,x)$ is a unit of $\mathcal E$ if and only if at least one
of the following conditions holds:
\begin{enumerate}
\item $a-b\in\mathbb Z_{\ge0}$;
\item $a=n\in\mathbb Z_{>0}$ and $A_{n,b}(x)=0$.
\end{enumerate}
In the first case the value is a nonzero algebraic multiple of $e^x$.
In the second case it is a positive algebraic number.
In particular,
\[
 M(a,b,1)\in\mathcal E^\times
 \quad\Longleftrightarrow\quad
 a-b\in\mathbb Z_{\ge0}\quad\hbox{or}\quad(a,b)=(2,3).
\]
The exceptional value is $M(2,3,1)=2$.
\end{theorem}

\begin{proof}
We use the injective phase morphism of Fischler--Rivoal,
Proposition~3.1.  It sends an $E$-value $f(1)$ to the exponential
polynomial
\[
 \Phi_\theta(f(1))(X)=\sum_\rho f_\rho(1)e^{\rho X},
\]
where $f(z)=\sum_\rho f_\rho(z)e^{\rho z}$ is its canonical
large-sector, Gevrey-order-one decomposition, in a direction just above
$\theta$.  It fixes algebraic constants and sends $e^\alpha$ to
$e^{\alpha X}$.  A unit has only one nonzero phase, by their
Corollary~3.3 (or by the units of the ring of exponential polynomials).
We apply it to the $E$-function $z\mapsto M(a,b,xz)$.

The standard Kummer connection formula (DLMF~13.2.41), on the upper
side of the positive real axis, is
\begin{equation}
\label{p2:eq:positive-kummer-connection}
 M(a,b,z)=
 \frac{\Gamma(b)e^{i\pi a}}{\Gamma(b-a)}U(a,b,z)
 +\frac{\Gamma(b)e^{-i\pi(b-a)}}{\Gamma(a)}
     e^z U(b-a,b,e^{-i\pi}z).
\end{equation}
Both summands are the canonical terms with respective phases $0$ and
$1$: this follows from the Gevrey asymptotic expansion
$U(c,b,z)\sim z^{-c}\sum_{j\ge0}
(-1)^j(c)_j(1+c-b)_j/(j!z^j)$ in the usual sectors, together with
uniqueness of the large-sector decomposition.  For clarity, a sector
bisected by a sufficiently small positive angle has opening greater
than $\pi$ and lies in both relevant $U$-asymptotic sectors; hence the
choice matches $\Phi_0$.  After replacing $z$ by $xz$, write
\[
 \Phi_0(M(a,b,x))(X)=C_0+C_1e^{xX}.
\]
Its constant coefficient is
\begin{equation}
\label{p2:eq:positive-kummer-C0}
 C_0=\frac{\Gamma(b)e^{i\pi a}}{\Gamma(b-a)}U(a,b,x).
\end{equation}
The integral (DLMF~13.4.4)
\[
 U(a,b,x)=\frac1{\Gamma(a)}
       \int_0^\infty e^{-xt}t^{a-1}(1+t)^{b-a-1}\,dt
\]
is strictly positive for the parameters in the theorem.

First suppose $a-b=m\in\mathbb Z_{\ge0}$.  Kummer's transformation
gives $M(a,b,x)=e^xM(-m,b,-x)$, an algebraic multiple of $e^x$.
It is nonzero because the defining series for $M(a,b,x)$ is strictly
positive.  Such values and their reciprocals belong to $\mathcal E$.

Suppose henceforth that $b-a\notin\mathbb Z_{\le0}$, so that $C_0\ne0$.
If $a$ is not an integer, $C_0$ has nonzero imaginary part, since
$\Gamma(b)/\Gamma(b-a)$ is real and nonzero and
$\sin(\pi a)\ne0$.  But $M(a,b,x)$ is real and
$M(a,b,x)=C_0+C_1e^x$.  Thus $C_1\ne0$.  The two nonzero phases
exclude a unit.

It remains to consider $a=n\in\mathbb Z_{>0}$.  The phase-$1$
formal series terminates because $(1-n)_j=0$ for $j\ge n$;
its sum is therefore exact.  Equivalently, substituting the terminating
$U$-expansion into \eqref{p2:eq:positive-kummer-connection} gives
\begin{equation}
\label{p2:eq:positive-kummer-C1}
 C_1=\Gamma(b)x^{1-b}A_{n,b}(x).
\end{equation}
Consequently $A_{n,b}(x)\ne0$ implies two nonzero phases and excludes a
unit.

To treat $A_{n,b}(x)=0$ without making an unsupported arithmetic claim
about $U$-values, use the contiguous identity (DLMF~13.3.1)
\[
 nM(n+1,b,z)=(b-n)M(n-1,b,z)+(2n-b+z)M(n,b,z).
\]
It gives polynomials $B_{n,b}\in\mathbb Q[z]$ such that
\begin{equation}
\label{p2:eq:positive-kummer-polynomial-reduction}
 M(n,b,z)=A_{n,b}(z)M(1,b,z)+B_{n,b}(z).
\end{equation}
Here the initial data are $A_{0,b}=0$, $A_{1,b}=1$,
$B_{0,b}=1$, $B_{1,b}=0$, and both sequences satisfy
\[
 nY_{n+1,b}=(b-n)Y_{n-1,b}+(2n-b+z)Y_{n,b}.
\]
The displayed explicit formula for $A_{n,b}$ satisfies this recurrence;
alternatively it follows by comparing phase-$1$ coefficients with
$M(1,b,z)$, whose phase-$1$ coefficient is $\Gamma(b)z^{1-b}$.
If $A_{n,b}(x)=0$, equation
\eqref{p2:eq:positive-kummer-polynomial-reduction} shows that the value is
$B_{n,b}(x)\in\overline{\mathbb Q}$.  Positivity again makes it nonzero.
This proves the general classification.

For $x=1$, the polynomial
\[
 P_n(b)=(n-1)!A_{n,b}(1)
    =\sum_{j=0}^{n-1}(-1)^j\binom{n-1}{j}(b-n)_j
\]
has integer coefficients and leading coefficient $(-1)^{n-1}$.
For $n=1$ it is the constant $1$.  For $n\ge2$ its rational roots are
integers.  Thus a zero with rational $b>0$ requires $b=m\in\mathbb Z_{>0}$.
The case $m\le n$ is already covered by the first alternative, so assume
$q=m-n\ge1$, and put $k=n-1$.  If $T_q$ has the gamma density
$e^{-t}t^{q-1}/\Gamma(q)$ on $(0,\infty)$, then
\[
 P_n(m)=\mathbb E(1-T_q)^k.
\]
For even $k$ this is strictly positive.  For odd $k$, the function
$t\mapsto(1-t)^k$ is strictly decreasing.  Coupling
$T_q=T_1+T_{q-1}$ with independent positive summands proves, for $q>1$,
\[
 \mathbb E(1-T_q)^k<\mathbb E(1-T_1)^k
 =\sum_{j=0}^{k}(-1)^j\binom{k}{j}j!
 =(-1)^k\, !k,
\]
where $!k$ is the number of derangements of $k$ elements.
For odd $k\ge3$, $!k>0$; the same formula handles $q=1$.
For $k=1$ the expectation is $1-q$ and vanishes exactly when $q=1$.
Hence the only further zero is $n=2,m=3$.
Finally $M(2,b,z)=(z+2-b)M(1,b,z)+b-1$, giving $M(2,3,1)=2$.
\end{proof}

\begin{remark}[A genuine exceptional curve at variable argument]
For every rational $b>2$ the algebraic point $x=b-2$ satisfies
\[
 M(2,b,b-2)=b-1.
\]
Thus a blanket statement that non-polynomial Kummer values are nonunits
would be false.  The second alternative in
Theorem~\ref{p2:thm:positive-kummer-units} is essential.
\end{remark}

\section{A denominator obstruction involving Euler's constant}
\label{p2:sec:ein-denominator-round2}

Write
\[
 \operatorname{Ein}(z)=\sum_{n\ge1}\frac{(-1)^{n-1}z^n}{n\,n!},
 \qquad \eta=\operatorname{Ein}(1)=\gamma+\frac\delta e,
 \qquad \delta=\int_0^\infty\frac{e^{-t}}{1+t}\,dt.
\]
Thus $\eta\in\mathcal E$ and $\gamma,\delta>0$.
The identity on the positive real axis
\[
 \operatorname{Ein}(z)=\gamma+\log z+E_1(z),
 \qquad
 E_1(z)=\int_z^\infty\frac{e^{-t}}{t}\,dt
\]
and the Gevrey expansion
$e^zE_1(z)\sim\sum_{j\ge0}(-1)^j j!z^{-j-1}$
show that its canonical phases are $0$ and $-1$.
In particular, Fischler--Rivoal's morphism gives
\begin{equation}
\label{p2:eq:ein-phase-round2}
 \Phi_0(\eta)(X)=\gamma+\delta e^{-X}.
\end{equation}
The logarithm in the phase-$0$ term vanishes when its coefficient
function is evaluated at $z=1$.  Both coefficients in
\eqref{p2:eq:ein-phase-round2} are nonzero, so $\eta$ is not a unit of
$\mathcal E$.  More generally:

\begin{theorem}[Rational denominators in $\operatorname{Ein}(1)$]
\label{p2:thm:ein-denominators}
Let $P,Q\in\overline{\mathbb Q}[Y]$ be relatively prime,
$P\ne0$, $Q\ne0$, and assume $Q(\eta)\ne0$.  If
\[
 \frac{P(\eta)}{Q(\eta)}\in\mathcal E,
\]
then either $Q$ is constant or
$Q(Y)=c(Y-\gamma)^m$ for some $c\ne0$ and $m\ge1$.
In particular,
\[
 \frac1{\gamma+\delta/e}\notin\mathcal E.
\]
More generally, $1/(\eta+\alpha)\notin\mathcal E$ for every
nonnegative real algebraic $\alpha$.
If $\gamma$ is transcendental, the exact intersection is
\[
 \overline{\mathbb Q}(\eta)\cap\mathcal E
       =\overline{\mathbb Q}[\eta].
\]
The final assertion is explicitly conditional; no irrationality or
transcendence claim for $\gamma$ is being made.
\end{theorem}

\begin{proof}
Let $\xi=P(\eta)/Q(\eta)\in\mathcal E$.  Apply the ring morphism
to $Q(\eta)\xi=P(\eta)$ and use
\eqref{p2:eq:ein-phase-round2}.  As an identity of meromorphic functions,
\[
 \Phi_0(\xi)(X)=
  \frac{P(\gamma+\delta e^{-X})}
       {Q(\gamma+\delta e^{-X})}.
\]
We use the elementary fact that a rational function of $T=e^{-X}$
which is a finite exponential polynomial in $X$ must be a Laurent
polynomial in $T$.  Indeed the rational expression is $2\pi i$-periodic.
If its exponential-polynomial expression is
$\sum_{\rho}c_\rho e^{\rho X}$ with distinct phases and nonzero
coefficients, periodicity and linear independence of exponentials give
$e^{2\pi i\rho}=1$ for every $\rho$.  Hence $\rho\in\mathbb Z$,
which proves the assertion.

Consequently $P(\gamma+\delta T)/Q(\gamma+\delta T)$ has no finite
pole except possibly at $T=0$.  The numerator and denominator remain
relatively prime under the invertible affine substitution
$Y=\gamma+\delta T$ because $\delta\ne0$.  Therefore every root of
$Q$ is $\gamma$.  This gives the stated shape of $Q$.
For $Q(Y)=Y+\alpha$, its root $-\alpha$ differs from the positive
number $\gamma$, proving the explicit exclusions.
If $\gamma$ is transcendental, a nonconstant polynomial over
$\overline{\mathbb Q}$ cannot have it as a root; hence $Q$ must be
constant.  The reverse inclusion in the intersection statement follows
from $\eta\in\mathcal E$ and the ring property.
\end{proof}

\paragraph{Exact scope and provenance.}
The Kummer classification is a new extension derived here from the
published phase morphism and classical Kummer identities.  The
$\operatorname{Ein}$ result is a new application of the same morphism.
Neither is the full units conjecture, nor Fischler--Rivoal's
$\mathcal E\cap\mathcal D=\overline{\mathbb Q}$ conjecture, nor a
resolution of the irrationality of $\gamma$ or $\delta$.  Multiple
parameter choices or denominators are not counted as independent
published open questions.

\paragraph{Primary references.}
S.~Fischler and T.~Rivoal,
\emph{The ring of values of $E$-functions is not a field},
\url{https://arxiv.org/abs/2610.16603}, Proposition~3.1,
Corollary~3.3, Conjecture~1.2 and Theorem~1.1.
The phase convention is also displayed explicitly in their Section~4.3
(printed page~12).  Compare T.~Rivoal and J.~Roques,
\emph{$E$-functions of order 2 and units of $E$-values},
\url{https://rivoal.perso.math.cnrs.fr/articles/unitse.pdf}, Theorem~5
(printed page~3).
The relevant special-function formulas are
\url{https://dlmf.nist.gov/13.2.E39},
\url{https://dlmf.nist.gov/13.2.E41},
\url{https://dlmf.nist.gov/13.3.E1},
\url{https://dlmf.nist.gov/13.4.E4}, and
\url{https://dlmf.nist.gov/13.7.E3}.


% Source fragment: round2_whittaker_fragment.tex
\section{The real-zero branch conjecture for Whittaker functions}
\label{p2:sec:whittaker-branch}

Boussaada, Mazanti and Niculescu, \emph{Some remarks on the location of
non-asymptotic zeros of Whittaker and Kummer hypergeometric functions},
Bulletin des Sciences Math\'ematiques \textbf{174} (2022), 103093,
arXiv:2106.03378, conjecture immediately after Figure~1 that the real-zero
branch of $M_{k,l}$ passing through $(k,z)=(l+3/2,2l+1)$ has maximal
parameter interval $(l+1/2,\infty)$, and tends to infinity and zero at its
respective endpoints. The following proof establishes this statement for
every $l>-1/2$, and gives a sharper right-endpoint limit.

\begin{theorem}\label{p2:thm:whittaker-branch}
Fix $b>0$. For every $\lambda>0$, the function
$y_\lambda(x)={}_1F_1(-\lambda;b;x)$ has a first positive zero
$\rho_b(\lambda)$. This zero is simple, depends real analytically on
$\lambda$, and satisfies
\[
 \rho_b'(\lambda)=
 -\frac{\displaystyle\int_0^{\rho_b(\lambda)}
 x^{b-1}e^{-x}y_\lambda(x)^2\,dx}
 {\rho_b(\lambda)^b e^{-\rho_b(\lambda)}
   y_\lambda'(\rho_b(\lambda))^2}<0.
\]
Moreover,
\[
 \lim_{\lambda\downarrow0}\rho_b(\lambda)=\infty,
 \qquad
 \lim_{\lambda\to\infty}\lambda\rho_b(\lambda)
   =\frac{j_{b-1,1}^{2}}4,
\]
where $j_{\nu,1}$ is the first positive zero of $J_\nu$.
Consequently, for $l>-1/2$, the branch in the cited conjecture is
\[
 z_l(k)=\rho_{2l+1}(k-l-1/2),\qquad k>l+1/2.
\]
\end{theorem}

\begin{proof}
Set $p(x)=x^be^{-x}$ and $w(x)=x^{b-1}e^{-x}$. Kummer's equation is
\[
 (p y_\lambda')'+\lambda w y_\lambda=0,\qquad
 y_\lambda(0)=1,
\]
and integration from zero gives
\[
 p(x)y_\lambda'(x)=-\lambda\int_0^xw(t)y_\lambda(t)\,dt.
\]
If $y_\lambda$ had no positive zero, it would remain positive. Its
integral on any fixed initial interval would then be a positive constant,
and for all sufficiently large $x$ the last identity would imply
$y_\lambda'(x)\leq-Ce^xx^{-b}$ with $C>0$. Integration would force
$y_\lambda$ to become negative, a contradiction. At the first zero the
same identity gives $y_\lambda'<0$, so that zero is simple.

The Kummer series is jointly analytic in $\lambda$ and $x$ for fixed
$b>0$. The implicit function theorem, together with positivity on each
closed subinterval before the first zero, shows that the first-zero
function is locally real analytic. These local branches agree and cover
$(0,\infty)$.

Write $v=\partial_\lambda y_\lambda$. Differentiating the equation and
forming the Lagrange identity gives
\[
 [p(y_\lambda v'-v y_\lambda')]'=-w y_\lambda^2.
\]
The boundary term at zero vanishes because $b>0$ and the functions and
their derivatives are bounded there. At $\rho=\rho_b(\lambda)$ the
identity therefore reads
$p(\rho)v(\rho)y_\lambda'(\rho)=\int_0^\rho w y_\lambda^2$.
Implicit differentiation of $y_\lambda(\rho)=0$ proves the displayed
strictly negative derivative formula.

As $\lambda\downarrow0$, $y_\lambda\to1$ uniformly on every compact
interval, proving the first limit. For the second limit, the defining
series gives, uniformly on compact subsets of the complex $t$-plane,
\[
 y_\lambda(t/\lambda)\longrightarrow
 {}_0F_1(;b;-t)
 =\Gamma(b)t^{(1-b)/2}J_{b-1}(2\sqrt t).
\]
For completeness, when $\lambda\geq1$,
$|(-\lambda)_m|/\lambda^m\leq m!$, so the summands are dominated on
$|t|\leq T$ by $T^m/(b)_m$, a summable sequence. The limiting function
is positive before $j_{b-1,1}^2/4$ and changes sign at that simple zero.
Uniform convergence therefore locates the first zero of the approximants
between any two sufficiently close points on its opposite sides.
This proves the asserted scaled limit.

Finally,
$M_{k,l}(x)=e^{-x/2}x^{l+1/2}
{}_1F_1(l+1/2-k;2l+1;x)$ for $x>0$. At $\lambda=1$ the Kummer function
is $1-x/b$, so the branch passes through the specified point. For
$\lambda\leq0$ all its series coefficients are nonnegative and its
constant term is one; there are no positive zeros. Since the branch
diverges as $\lambda\downarrow0$, its maximal interval as a finite real
zero branch is exactly the one stated.
\end{proof}

\paragraph{Status and scope.}
This is one conjecture, with two endpoint assertions, and should be
counted as one problem. The proof uses classical Sturm--Liouville
identities. Recent related spectral results include R.~Vanlaere,
arXiv:2511.20025v4 (2026), Theorem~1.6, for zeros in the first Kummer
parameter when $b\geq1$. They do not state the branch conjecture above.
No claim of priority beyond the checked sources is implicit here.


% Source fragment: round2_vargas_fragment.tex
\section{An exterior-zero counterexample to a conjecture on exponential integrals}
\label{p2:sec:vargas-counterexample}

Vargas, \emph{Limit curves for zeros of sections of exponential
integrals}, Constructive Approximation \textbf{40} (2014), 219--239,
arXiv:1302.3695v4, Conjecture~5.1, asserts that every finite limit point
of the zeros of $s_n[F](nz)$ lies on $D_{a,b}\cup D_{\mathrm{imag}}$.
Here $F(z)=\int_{-a}^b\phi(t)e^{zt}\,dt$ belongs to the class defined in
Section~2 of that paper, and $s_n[F]$ denotes its Taylor polynomial of
degree $n$. Signed and complex-valued integrands are expressly allowed.
For $a=b=1$ the proposed set is contained in $\{|z|\leq1\}$.

\begin{theorem}\label{p2:thm:vargas-counterexample}
For every complex $c$ with $|c|>1$, let
\[
 F_c(z)=\int_{-1}^1(1+ct)e^{zt}\,dt.
\]
There are zeros $z_n$ of $s_n[F_c](nz)$, indexed by all sufficiently
large even integers, such that $z_n\to-c$. Thus any prescribed point
strictly outside the unit disk can occur as a zero limit. In particular,
$c=3/2$ is a real signed-integrand counterexample to Vargas's Conjecture~5.1.
\end{theorem}

\begin{proof}
The integrand is integrable and analytic, with nonzero endpoint values
$1-c$ and $1+c$. It satisfies the source's hypotheses with
$a=b=1$, $\mu=\nu=0$,
$f_1(u)=1-c+cu$ and $f_2(u)=1+c-cu$. The extra noncancellation
condition (5) of that paper is also satisfied: its expression equals
$2$ for even indices and $2c$ for odd indices.

Its moments are
\[
 m_k=\int_{-1}^1(1+ct)t^k\,dt
 =\begin{cases}2/(k+1),&k\text{ even},\\
                 2c/(k+2),&k\text{ odd}.
   \end{cases}
\]
For even $n$, reverse the normalized Taylor polynomial:
\[
 Q_n(z):=\frac{n!}{m_n(nz)^n}s_n[F_c](nz)
 =\sum_{j=0}^n\frac{m_{n-j}}{m_n}
    \frac{n!}{(n-j)!n^j}z^{-j}.
\]
For fixed $j$, the coefficient tends to $1$ if $j$ is even, and to
$c$ if $j$ is odd. Also,
\[
 0<\frac{n!}{(n-j)!n^j}\leq1,
 \qquad
 \left|\frac{m_{n-j}}{m_n}\right|
 \leq\max(1,|c|)(j+1),
\]
where the latter inequality follows from
$(n+1)/(n-j+1)\leq j+1$ for $0\leq j\leq n$.
Dominated convergence of the geometric majorant shows that, locally
uniformly on $|z|>1$,
\[
 Q_n(z)\longrightarrow
 \sum_{h\geq0}z^{-2h}+c\sum_{h\geq0}z^{-2h-1}
 =\frac{1+c/z}{1-z^{-2}}.
\]
The last function has a simple zero at $z=-c$, lying strictly outside
the unit disk. Choose a small disk about $-c$ contained in $|z|>1$.
Rouch\'e's theorem gives exactly one zero of $Q_n$ in that disk for
all sufficiently large even $n$, and these zeros converge to $-c$.
The normalizing factor is nonzero there, so they are zeros of
$s_n[F_c](nz)$ as well. This contradicts the stated conjecture.
\end{proof}

\paragraph{Scope.}
This disproves the conjecture for the signed-integrand class actually
specified in its source. It does not settle the source's separate
Conjecture~5.2 on positive angular densities. Nor is it an assertion
about nonnegative integrands. The construction produces a tunable
isolated exterior zero, a phenomenon not detected by the limiting
normalized zero-counting measure.


% Source fragment: round2_probability.tex
% Round 2 probability branch.  Fragment for inclusion in the main article.
% Bibliography key needed: Gorgens2014 (arXiv:1409.5285v1).

\section{A conjecture about the scaling of a Brownian bridge}
\label{p2:sec:gorgens-conjecture}

This section addresses an independently posed conjecture. In
\cite[Section~4]{p2:Gorgens2014}, G\"orgens explicitly conjectures that the
optimal expected reward of an $\alpha$-Brownian bridge has a local minimum
at $\alpha=1/2$, and reports that he could not prove it. The same paper
states two further limiting expectations, in its introduction and
Section~4. We prove the local-minimum assertion and both limits below.
The latter are elementary accompanying results; they should not be
presented as two major conjectures. A search locating no later proof is
not a proof of priority, and the newness of these arguments remains to
be checked by specialists.

Let $W$ be a standard Brownian motion and, for $\alpha>0$, let
\[
 dX_t^{(\alpha)}=dW_t-\frac{\alpha X_t^{(\alpha)}}{1-t}\,dt,
 \qquad X_0^{(\alpha)}=X_1^{(\alpha)}=0.
\]
Set
\[
 V(\alpha)=\sup_{0\leq\tau\leq1}\mathbb E X_\tau^{(\alpha)}.
\]
G\"orgens' established solution has an optimal threshold
$B(\alpha)\sqrt{1-t}$. Write $\varepsilon=2\alpha-1$, and normalize the
positive increasing continuation solution by $h_\varepsilon(0)=1$.
Then
\begin{equation}\label{p2:eq:gorgens-ode}
 h_\varepsilon''(y)-\varepsilon y h_\varepsilon'(y)
     -h_\varepsilon(y)=0,\qquad
 B h_\varepsilon'(B)=h_\varepsilon(B),\qquad
 V(\alpha)=\frac{B}{h_\varepsilon(B)}.
\end{equation}
At $\varepsilon=0$ these are $h_0(y)=e^y$, $B=1$, and $V(1/2)=e^{-1}$.

\begin{theorem}[The local-minimum conjecture]\label{p2:thm:gorgens-minimum}
For real $u\to0$ from either side,
\begin{align}
 B(\tfrac12+u)&=1-\frac u2+\frac{5u^2}{8}+O(|u|^3),
       \label{p2:eq:gorgens-boundary-expansion}\\
 V(\tfrac12+u)&=\frac1e+\frac{u^2}{12e}+O(|u|^3).
       \label{p2:eq:gorgens-value-expansion}
\end{align}
In particular $\alpha=1/2$ is a strict local minimum of $V$.
\end{theorem}
\begin{proof}
The initial slope obtained from the explicit Kummer formula in
\cite[Theorem~1]{p2:Gorgens2014} is
\begin{equation}\label{p2:eq:gorgens-initial-slope}
 h_\varepsilon'(0)=
 \begin{cases}
  \sqrt{2\varepsilon}\,
    \dfrac{\Gamma(x+1/2)}{\Gamma(x)},
       &\varepsilon>0,\quad x=(2\varepsilon)^{-1},\\[6pt]
  \sqrt{-2\varepsilon}\,
    \dfrac{\Gamma(x+1)}{\Gamma(x+1/2)},
       &\varepsilon<0,\quad x=-(2\varepsilon)^{-1}.
 \end{cases}
\end{equation}
Stirling's formula with remainder gives, as $x\to+\infty$,
\[
 \log\frac{\Gamma(x+1/2)}{\Gamma(x)}
   =\frac12\log x-\frac1{8x}+O(x^{-3}),
\]
and hence
\[
 \frac{\Gamma(x+1/2)}{\Gamma(x)\sqrt{x}}
   =1-\frac1{8x}+\frac1{128x^2}+O(x^{-3}).
\]
Substitution in the first case of
\eqref{p2:eq:gorgens-initial-slope}, and inversion of the same expansion
in the second case using $\Gamma(x+1)=x\Gamma(x)$, yield the identical
two-sided expansion
\begin{equation}\label{p2:eq:gorgens-slope-expansion}
 h_\varepsilon'(0)=1-\frac\varepsilon4
                       +\frac{\varepsilon^2}{32}
                       +O(|\varepsilon|^3).
\end{equation}
The uniform remainder here is essential: the continuation solution is
selected by a condition at $y=-\infty$, so merely substituting a formal
power series in its differential equation would not suffice.

On any fixed compact interval containing $[0,1]$, standard continuous
dependence for the first-order system associated with
\eqref{p2:eq:gorgens-ode} shows that $h_\varepsilon$ and its first
derivative converge uniformly to $e^y$. Thus
$r_\varepsilon=h_\varepsilon'/h_\varepsilon$ is well defined there for
small $|\varepsilon|$ and satisfies
\[
 r_\varepsilon'+r_\varepsilon^2
       -\varepsilon y r_\varepsilon-1=0.
\]
Let
\[
 p_1(y)=\frac y2-\frac14,\qquad
 p_2(y)=\frac{y^2}{8}-\frac y8+\frac1{32}.
\]
They obey
\[
 p_1'+2p_1=y,\qquad
 p_2'+2p_2=yp_1-p_1^2,
\]
with $p_1(0)=-1/4$ and $p_2(0)=1/32$. Subtracting the Riccati
equations for $r_\varepsilon$ and
$1+\varepsilon p_1+\varepsilon^2p_2$, and applying the integral
Gronwall inequality, proves uniformly on that compact interval that
\begin{equation}\label{p2:eq:gorgens-riccati-expansion}
 r_\varepsilon(y)=1+\varepsilon p_1(y)
                         +\varepsilon^2p_2(y)
                         +O(|\varepsilon|^3).
\end{equation}
Indeed the approximate solution has residual $O(|\varepsilon|^3)$,
its initial discrepancy has that size by
\eqref{p2:eq:gorgens-slope-expansion}, and all coefficients remain bounded.

At $(\varepsilon,y)=(0,1)$ the function
$y r_\varepsilon(y)-1$ has $y$-derivative $1$. There is consequently
a root near $1$ for small $|\varepsilon|$, and its $y$-derivative
remains bounded away from zero. This is the unique positive stopping
threshold in G\"orgens' solution. Substitution of
$B=1+b_1\varepsilon+b_2\varepsilon^2+O(|\varepsilon|^3)$ in
$Br_\varepsilon(B)=1$ gives
$b_1=-1/4$ and $b_2=5/32$. More formally, substitution of this
quadratic polynomial leaves residual $O(|\varepsilon|^3)$, and the
mean-value theorem gives the asserted error for the actual root.

Since $h_\varepsilon(0)=1$,
\[
 \log V(\alpha)=\log B-\int_0^B r_\varepsilon(y)\,dy.
\]
The necessary integrals are
\[
 \int_0^1p_1(y)\,dy=0,\qquad
 p_1(1)=\frac14,\qquad
 \int_0^1p_2(y)\,dy=\frac1{96}.
\]
With $B-1=-\varepsilon/4+O(\varepsilon^2)$ this yields
\begin{align*}
 \log V(\alpha)
 &= -1-\frac12(B-1)^2
       -\varepsilon p_1(1)(B-1)
       -\varepsilon^2\int_0^1p_2(y)\,dy
       +O(|\varepsilon|^3)\\
 &= -1+\frac{\varepsilon^2}{48}+O(|\varepsilon|^3).
\end{align*}
Exponentiating and setting $\varepsilon=2u$ proves both expansions.
The positive quadratic term gives strict local minimality.
\end{proof}

\begin{remark}
Differentiating Stirling's expansion with its differentiated remainder
also gives the ordinary derivatives $V'(1/2)=0$ and
$V''(1/2)=1/(6e)$. The strict-minimum conclusion above requires only
the explicitly proved two-sided value expansion, so it does not
depend on this additional regularity assertion.
\end{remark}

\begin{theorem}[Two stated limiting expectations]\label{p2:thm:gorgens-limits}
Under the coupling by a common Brownian motion, the optimal stopping
times in \cite[equation~(10)]{p2:Gorgens2014} satisfy
\[
 \tau^*(\alpha)\longrightarrow1\quad\hbox{almost surely as }\alpha\downarrow0.
\]
Furthermore
\[
 \sup_{0\leq t\leq1}|X_t^{(\alpha)}|\longrightarrow0
 \quad\hbox{almost surely and in every }L^p,\quad 0<p<\infty,
 \qquad \alpha\to\infty.
\]
In particular $V(\alpha)\to0$ as $\alpha\to\infty$.
\end{theorem}
\begin{proof}
Integration by parts gives a common pathwise version, for $t<1$,
\begin{equation}\label{p2:eq:gorgens-path-identity}
 X_t^{(\alpha)}
   =W_t-\alpha(1-t)^\alpha\int_0^t
                 \frac{W_s}{(1-s)^{\alpha+1}}\,ds.
\end{equation}
The positive kernel in the integral has total mass
$1-(1-t)^\alpha\leq1$. Therefore, writing
$M=\sup_{0\leq s\leq1}|W_s|$, we have simultaneously for every
$\alpha>0$ and $t\in[0,1]$,
\begin{equation}\label{p2:eq:gorgens-uniform-path-bound}
 |X_t^{(\alpha)}|\leq2M.
\end{equation}
The known limit $B(\alpha)\to\infty$ from
\cite[Theorem~2(i)]{p2:Gorgens2014} also follows directly from the ODE.
As $\alpha\downarrow0$, its parameter $\varepsilon$ tends to $-1$
and its initial slope in \eqref{p2:eq:gorgens-initial-slope} tends to
$\sqrt{\pi/2}$. Thus the normalized solutions converge on compact
$y$-intervals to a solution $h$ of $h''+yh'-h=0$, with $h(0)=1$.
The residual $Q=h-yh'$ satisfies $Q'=-yQ$ and $Q(0)=1$, hence
$Q(y)=e^{-y^2/2}>0$. Uniform compact convergence therefore precludes
a bounded subsequence of roots $B(\alpha)$ of
$h_\varepsilon-yh_\varepsilon'=0$, proving that limit.
At any threshold hit strictly before $1$,
continuity gives
$B(\alpha)\sqrt{1-\tau^*(\alpha)}=X_{\tau^*(\alpha)}^{(\alpha)}$.
Consequently, also when $\tau^*(\alpha)=1$,
\[
 0\leq1-\tau^*(\alpha)\leq\frac{4M^2}{B(\alpha)^2},
\]
which proves the first assertion and, additionally, convergence in
every $L^p$.

For the second assertion fix $\delta\in(0,1)$ and set
\[
 \omega_W(\delta)=\sup\{|W_t-W_s|:s,t\in[0,1],\ |t-s|\leq\delta\}.
\]
Rewrite \eqref{p2:eq:gorgens-path-identity} as
\[
 X_t^{(\alpha)}=(1-t)^\alpha W_t+
 \alpha(1-t)^\alpha\int_0^t
             \frac{W_t-W_s}{(1-s)^{\alpha+1}}\,ds.
\]
The part of the integral with $s\geq(t-\delta)^+$ is bounded in
absolute value by $\omega_W(\delta)$. For $t\geq\delta$, the kernel
mass of the other part is at most
\[
 \left(\frac{1-t}{1-t+\delta}\right)^\alpha
 \leq(1+\delta)^{-\alpha}.
\]
The first term is at most $\omega_W(\delta)$ if $t<\delta$, and
at most $(1-\delta)^\alpha M$ otherwise. Thus
\[
 \sup_t|X_t^{(\alpha)}|
 \leq2\omega_W(\delta)
        +M(1-\delta)^\alpha+2M(1+\delta)^{-\alpha}.
\]
First let $\alpha\to\infty$ and then $\delta\downarrow0$.
Uniform continuity of every Brownian path gives the almost-sure
limit. The domination \eqref{p2:eq:gorgens-uniform-path-bound} and the
finiteness of every moment of $M$ give all the $L^p$ limits. Finally,
$0\leq V(\alpha)\leq\mathbb E\sup_t|X_t^{(\alpha)}|$.
\end{proof}

\begin{remark}[Scope and arithmetic connection]
These results concern a drift parameter $\alpha$, not the dimension
parameter of the Bessel bridge considered below. The common mechanism is
parameter sensitivity of a special-function free boundary. The gamma
function enters \eqref{p2:eq:gorgens-initial-slope}, and the critical
value and curvature are $e^{-1}$ and $(6e)^{-1}$. None of this proves
irrationality of the Euler--Mascheroni constant. The local-minimum
result resolves the precise assertion made in the cited preprint;
it neither establishes a global minimum nor locates the other
numerically observed critical point.
\end{remark}

\subsection{A sharp form of the large-drift limit}

The limiting expectation in Theorem~\ref{p2:thm:gorgens-limits} admits
the following quantitative strengthening. It is part of the same
endpoint analysis, not an additional historical question.

\begin{theorem}[Large-drift asymptotics]\label{p2:thm:gorgens-large-drift}
Put $\varepsilon=2\alpha-1$. As $\alpha\to\infty$,
\begin{align}
 \varepsilon B(\alpha)^2
   &=2\log\varepsilon-\log\log\varepsilon-\log(4\pi)+o(1),
          \label{p2:eq:gorgens-large-B}\\
 \varepsilon V(\alpha)^2
   &=2\log\varepsilon-\log\log\varepsilon-\log(4\pi)-2+o(1).
          \label{p2:eq:gorgens-large-V}
\end{align}
In particular
$B(\alpha)\sim V(\alpha)\sim\sqrt{\log\alpha/\alpha}$.
\end{theorem}
\begin{proof}
For $\varepsilon>0$ put $a=1/\varepsilon$ and define
\[
 I_a(u)=\int_0^\infty s^{a-1}\exp(us-s^2/2)\,ds,
 \qquad F_a(u)=I_a'(u)
        =\int_0^\infty s^a\exp(us-s^2/2)\,ds.
\]
Integration by parts in $s$ shows
$I_a''-uI_a'-aI_a=0$. Hence
$h_\varepsilon(y)=I_a(\sqrt\varepsilon y)/I_a(0)$
solves \eqref{p2:eq:gorgens-ode}. Its initial slope is
$\sqrt{2\varepsilon}\,\Gamma((a+1)/2)/\Gamma(a/2)$,
which agrees with \eqref{p2:eq:gorgens-initial-slope}; uniqueness for
the ordinary initial-value problem identifies it with the normalized
continuation solution. With $u=\sqrt\varepsilon B(\alpha)$ its threshold equation
and value are therefore
\begin{equation}\label{p2:eq:gorgens-large-equation}
 uF_a(u)=I_a(u),\qquad
 V(\alpha)=B(\alpha)\frac{I_a(0)}{I_a(u)}.
\end{equation}
Also
\[
 I_a(0)=2^{a/2-1}\Gamma(a/2)\sim\frac1a\qquad(a\downarrow0).
\]
For $u$ in any fixed bounded positive interval, $F_a(u)$ is bounded
uniformly for $0<a\leq1$, since $s^a\leq1+s$. On the other hand
$I_a(u)\geq I_a(0)\to\infty$. The root in
\eqref{p2:eq:gorgens-large-equation} must consequently satisfy
$u\to\infty$ as $a\downarrow0$.

We require two estimates, uniform for $0<a\leq1$ as $u\to\infty$:
\begin{equation}\label{p2:eq:gorgens-laplace-estimates}
 F_a(u)=\sqrt{2\pi}\,u^a e^{u^2/2}(1+o(1)),
 \qquad
 \int_0^uF_a(v)\,dv=\frac{F_a(u)}u(1+o(1)).
\end{equation}
For the first, completion of the square gives
\[
 \frac{F_a(u)}{u^ae^{u^2/2}}
 =\int_{-u}^\infty(1+t/u)^a e^{-t^2/2}\,dt.
\]
On bounded $t$-intervals the factor converges to $1$ uniformly for
$a\in(0,1]$. For $t\geq0$ it is bounded by $1+t$ when $u\geq1$;
for $-u<t<0$ it is bounded by $1$. Gaussian domination gives the
first estimate uniformly.

For completeness, the endpoint integral used in the second estimate
has uniform asymptotics
\[
 \int_A^u v^a e^{v^2/2}\,dv
       \sim u^{a-1}e^{u^2/2}\qquad(0<a\leq1).
\]
Indeed integration by parts writes this integral as
\[
 u^{a-1}e^{u^2/2}-A^{a-1}e^{A^2/2}
       +(1-a)\int_A^u v^{a-2}e^{v^2/2}\,dv.
\]
For $A>1$ the final integral is between zero and $A^{-2}$ times
the original integral. After division by the upper-end term,
letting $u\to\infty$ and then $A\to\infty$ squeezes the ratio to
$1$, uniformly in $a$. The first estimate of
\eqref{p2:eq:gorgens-laplace-estimates}, applied above $A$, transfers
this conclusion to $F_a$. The integral below $A$ is uniformly
bounded and negligible. This proves both estimates.

Let $J_a(u)=\int_0^uF_a(v)\,dv$, so
$I_a(u)=I_a(0)+J_a(u)$. At the threshold,
\[
 \frac{J_a(u)}{uF_a(u)}=\frac{1+o(1)}{u^2},
 \qquad uF_a(u)=I_a(0)(1+o(1))\sim\frac1a.
\]
Writing $L=\log(1/a)=\log\varepsilon$ and taking logarithms in
the first asymptotic estimate for $F_a$ now gives
\begin{equation}\label{p2:eq:gorgens-u-log-equation}
 \frac{u^2}{2}+(a+1)\log u+\frac12\log(2\pi)=L+o(1).
\end{equation}
Since $u\to\infty$, this first implies $u^2\leq2L+O(1)$;
substitution back implies $u^2=2L+O(\log L)$. Consequently
$\log u^2=\log(2L)+o(1)$ and $a\log u^2=o(1)$.
Equation~\eqref{p2:eq:gorgens-u-log-equation} therefore yields
\[
 u^2=2L-\log L-\log(4\pi)+o(1),
\]
which is \eqref{p2:eq:gorgens-large-B}.

Finally, \eqref{p2:eq:gorgens-large-equation} gives the exact relation
\[
 \frac{V(\alpha)}{B(\alpha)}
 =1-\frac{J_a(u)}{uF_a(u)}
 =1-\frac{1+o(1)}{u^2}.
\]
It follows that $\varepsilon V(\alpha)^2=u^2-2+o(1)$, proving
\eqref{p2:eq:gorgens-large-V} and the stated equivalences.
\end{proof}

\section{Positive-order jets of the Bessel stopping boundary}
\label{p2:sec:boundary-jets-round2}

We give all definitions and inputs so that this section can be read
independently. In a Bessel bridge of dimension $d>0$, the squared radius
solves
\[
 dY_t=\left(d-\frac{2Y_t}{1-t}\right)dt+2\sqrt{Y_t}\,dW_t,
 \qquad Y_0=0.
\]
For the reward $Y_\tau$, the solution of
\cite[Theorem~3.1]{p2:HobsonLiu2026} has a radius threshold
$B(d,2)\sqrt{1-t}$ and initial value $\mathcal J(d)$, given by
\begin{equation}\label{p2:eq:round2-bessel-definitions}
 F_d(z)=M(1,d/2,z),\qquad
 z_dF_d'(z_d)=F_d(z_d),\qquad
 B(d,2)=\sqrt{2z_d},\qquad
 \mathcal J(d)=\frac{2z_d}{F_d(z_d)}.
\end{equation}
Here $M(a,b,z)=\sum_{k\geq0}(a)_kz^k/((b)_k k!)$.
The root $z_d>0$ is unique and depends real analytically on $d>0$.
To see this directly, write $F_d(z)=\sum_{k\geq0}c_kz^k$ with
$c_k>0$ and give the integer $k$ probability $c_kz^k/F_d(z)$.
Then $zF_d'/F_d$ is its mean and its derivative is its strictly
positive variance divided by $z$. The mean increases from zero to
infinity: for any fixed $m$, the mass on $\{0,\ldots,m\}$ tends to
zero as $z\to\infty$, by comparison with a higher coefficient.
The analytic implicit function theorem gives the last assertion.

Define the entire functions
\[
 H_0(z)=1-e^{-z},\qquad
 H_n(z)=\sum_{k\geq1}\frac{(-1)^{k+1}z^k}{k^n k!}\quad(n\geq1),
 \qquad \eta_n=H_n(1).
\]
They satisfy $H_n'=H_{n-1}/z$ for $n\geq1$, with the quotient
understood by continuity at zero. We use the established arithmetic
input that, for every $N\geq1$, the numbers
\begin{equation}\label{p2:eq:round2-known-independent-constants}
 e,H_1(1),\ldots,H_N(1)
 \quad\hbox{are algebraically independent over }\overline{\mathbb Q}.
\end{equation}
One explicit source is \cite[Theorem~2]{p2:Powers2026}; its constants
are $1-e^{-1}$ and $n!H_n(1)$. The same finite-family statement is
the $\alpha=0$, $z=1$ case of the integral family discussed in
\cite[Section~7]{p2:Rivoal2012}, using
$H_n(1)=\frac1{n!}\int_0^1(-\log t)^n e^{-t}\,dt$.

Put $\varepsilon=(d-2)/2$,
$z(\varepsilon)=z_{2+2\varepsilon}=B(2+2\varepsilon,2)^2/2$, and
$K_n=\overline{\mathbb Q}(e,\eta_1,\ldots,\eta_n)$, with
$K_0=\overline{\mathbb Q}(e)$. The parameter identity
\begin{equation}\label{p2:eq:round2-parameter-identity}
 M(1,1+\varepsilon,z)=e^z S(\varepsilon,z),\qquad
 S(\varepsilon,z)=1+\sum_{n\geq1}(-1)^nH_n(z)\varepsilon^n
\end{equation}
holds analytically near $\varepsilon=0$. Indeed for
$\varepsilon>0$, Euler's integral gives
\[
 e^{-z}M(1,1+\varepsilon,z)
 =\varepsilon\int_0^1 e^{-zt}t^{\varepsilon-1}\,dt
 =1+\varepsilon\sum_{k\geq1}\frac{(-z)^k}{k!(k+\varepsilon)}.
\]
Expansion in $\varepsilon$ proves the identity, and analytic
continuation supplies a neighborhood of zero.

The next theorem is an additional transfer result, not an answer to a
separately identified historical open question.

\begin{theorem}[Boundary-jet independence]\label{p2:thm:boundary-jets-round2}
For every integer $N\geq1$,
\[
 \overline{\mathbb Q}\bigl(
      \partial_dB(2,2),\ldots,\partial_d^N B(2,2)\bigr)
 =\overline{\mathbb Q}\bigl(e,H_1(1),\ldots,H_{N-1}(1)\bigr).
\]
In particular the $N$ positive-order dimension derivatives on the
left are algebraically independent over $\overline{\mathbb Q}$.
The zero-order value $B(2,2)=\sqrt2$ is algebraic and is deliberately
excluded.
\end{theorem}
\begin{proof}
Write $z(\varepsilon)=1+\sum_{n\geq1}z_n\varepsilon^n$. As before,
\[
 S(\varepsilon,z)=1+\sum_{n\geq1}(-1)^nH_n(z)\varepsilon^n,
 \qquad \frac1z-1-\frac{S_z}{S}=0.
\]
The first coefficient gives $z_1=H_0(1)=1-e^{-1}$. We show, for
$n\geq2$, the sharpened triangular formula
\begin{equation}\label{p2:eq:boundary-jet-leading}
 z_n=(-1)^{n+1}e^{-1}\eta_{n-1}+R_n,
 \qquad R_n\in K_{n-2}.
\end{equation}
In the coefficient of $\varepsilon^n$ in $S_z/S$ at $z=1$, the
only terms containing $\eta_{n-1}$ are the numerator term
$(-1)^nH_n'(1)=(-1)^n\eta_{n-1}$ and the product of
$-H_1'(1)=-H_0(1)$ with the term
$(-1)^n\eta_{n-1}$ in $[\varepsilon^{n-1}]S^{-1}$.
Their sum is
\[
 (-1)^n\bigl(1-H_0(1)\bigr)\eta_{n-1}
       =(-1)^ne^{-1}\eta_{n-1}.
\]
Every other term lies in $K_{n-2}$, since positive-order
$z$-differentiation lowers the index of the $H$ constants.
Substitution of $z=z(\varepsilon)$ does not change that leading
coefficient: terms resulting from the substitution involve either
lower coefficients $z_j\in K_{j-1}$ or a further positive derivative
of some $H_k$, and hence lie in $K_{n-2}$. The coefficient $z_n$
appears in $1/z-1$ with multiplier $-1$. This proves
\eqref{p2:eq:boundary-jet-leading} inductively.

Now $e=(1-z_1)^{-1}$, and
\eqref{p2:eq:boundary-jet-leading} recursively recovers
$\eta_1,\ldots,\eta_{N-1}$ from $z_1,\ldots,z_N$ and vice versa.
The conversion $B=\sqrt{2z}$ is triangular with nonzero algebraic
diagonal coefficient, since $z(0)=1$. The change of variables
$d=2+2\varepsilon$ only adds nonzero rational factors. This proves
the field equality. The known algebraic independence
\eqref{p2:eq:round2-known-independent-constants} gives the last assertion
(for $N=1$, only the transcendence of $e$ is needed).
\end{proof}

\begin{proposition}[Value jets, restated from the preceding research stage]
\label{p2:prop:round2-value-jets}
For every $N\geq0$,
\[
 \overline{\mathbb Q}\bigl(\mathcal J(2),\mathcal J'(2),
                     \ldots,\mathcal J^{(N)}(2)\bigr)
 =\overline{\mathbb Q}\bigl(e,H_1(1),\ldots,H_N(1)\bigr).
\]
Thus the $N+1$ displayed values of $\mathcal J$ and its derivatives
are algebraically independent. This proposition repeats a prior-stage
result to make the differential-transcendence argument self-contained;
it is not counted as a new result of this stage.
\end{proposition}
\begin{proof}
Equations \eqref{p2:eq:round2-bessel-definitions} and
\eqref{p2:eq:round2-parameter-identity} give $z(0)=1$,
$\mathcal J(2)=2/e$, and
\[
 \mathcal J(2+2\varepsilon)
   =\frac{2z(\varepsilon)e^{-z(\varepsilon)}}
          {S(\varepsilon,z(\varepsilon))}.
\]
Every positive-order derivative of $H_n$ at $1$ belongs to
$K_{n-1}$, by repeated differentiation of $H_n'=H_{n-1}/z$.
The root-coefficient argument above gives $z_n\in K_{n-1}$.
Now set
\[
 L(\varepsilon)=\log\frac{\mathcal J(2+2\varepsilon)}{2/e}
 =\log z(\varepsilon)-z(\varepsilon)+1
       -\log S(\varepsilon,z(\varepsilon)).
\]
For each $n\geq1$, its coefficient has the triangular form
\[
 [\varepsilon^n]L(\varepsilon)=(-1)^{n+1}\eta_n+R_n,
 \qquad R_n\in K_{n-1}.
\]
Indeed the only occurrence of $H_n(1)$ is the linear term
$-(-1)^nH_n(1)\varepsilon^n$ in $-\log S$. All its nonlinear
terms use lower $H$ indices. Substitution of $z(\varepsilon)$ uses
positive derivatives of the $H$ functions and lower root coefficients,
which belong to the indicated lower field. The derivative of
$\log z-z+1$ at $z=1$ is zero, so the newest root coefficient
does not appear in that term. Exponentiation consequently gives
\[
 \frac{2^n\mathcal J^{(n)}(2)}{n!}
   =\frac2e\bigl((-1)^{n+1}\eta_n+\widetilde R_n\bigr),
 \qquad\widetilde R_n\in K_{n-1}.
\]
Starting from $e=2/\mathcal J(2)$, these relations recover the
constants $\eta_n$ successively from the derivatives and vice versa.
This proves the field identity, and
\eqref{p2:eq:round2-known-independent-constants} proves independence.
\end{proof}

\begin{remark}[The first Euler-constant sensitivity]
The first value derivative is
\[
 \mathcal J'(2)=\frac{H_1(1)}e
      =\frac\gamma e+\frac\delta{e^2},
 \qquad \delta=\int_0^\infty\frac{e^{-t}}{1+t}\,dt,
\]
where $H_1(1)=\gamma+\delta/e$ is the classical
Euler--Gompertz identity. Thus these jets do encode an exact
combination involving Euler's constant. The independence theorem
does not separate the arithmetic of $\gamma$ from that of $\delta$.
The first boundary derivative is
$\partial_dB(2,2)=(1-e^{-1})/(2\sqrt2)$; Euler's constant enters
the next derivative through
$z_2=e^{-1}(1-e^{-1}-H_1(1))$.
\end{remark}

\section{The coefficient field in differential transcendence}
\label{p2:sec:constants-elimination-round2}

The earlier conclusion over $\overline{\mathbb Q}(d)$ in fact implies
the usual conclusion over $\mathbb C(d)$. This is a correction and
strengthening of its formulation, using the standard transitivity of
differential algebraicity; it is not a separate arithmetic discovery.
Here is a self-contained version of the needed fact.

\begin{lemma}[Elimination of finitely many constant coefficients]
\label{p2:lem:constant-elimination-round2}
An analytic germ which is differentially algebraic over
$\mathbb C(t)$ is differentially algebraic over $\mathbb Q(t)$.
\end{lemma}
\begin{proof}
An algebraic differential equation uses only finitely many complex
constants, say $c_1,\ldots,c_m$. Put
$K=\mathbb Q(c_1,\ldots,c_m)(t)$, a differential field with the usual
derivative and with constant $c_j$. Choose a nonzero differential
polynomial relation for the germ $f$ having minimal order $r$, and
then minimal degree in its highest derivative. Its separant, the
partial derivative with respect to the indeterminate representing
$f^{(r)}$, does not vanish after substitution: otherwise it would
give a nonzero relation of lower order or of smaller degree at the
same order, in characteristic zero. Differentiating the chosen
relation therefore expresses $f^{(r+1)}$ rationally in
$K(f,f',\ldots,f^{(r)})$. This field is stable under differentiation
and contains every derivative of $f$.

As an ordinary field it has finite transcendence degree over
$\mathbb Q(t)$, at most $m+r+1$. Sufficiently many derivatives of
$f$ must therefore be algebraically dependent over $\mathbb Q(t)$.
That is an algebraic differential equation with rational-function
coefficients over $\mathbb Q$, as asserted.
\end{proof}

\begin{corollary}[Differential transcendence over complex rational functions]
Both $d\mapsto V_{d,2}(0,0)$ and $d\mapsto B(d,2)$ are
differentially transcendental over $\mathbb C(d)$.
\end{corollary}
\begin{proof}
Proposition~\ref{p2:prop:round2-value-jets} and specialization at $d=2$
exclude an algebraic differential equation over
$\overline{\mathbb Q}(d)$ for the value function. In detail, clear
rational-function denominators in an alleged relation and divide its
polynomial coefficients by their largest common power of $d-2$.
Their specialization at $2$ is then a nonzero polynomial relation
over $\overline{\mathbb Q}$ among a finite jet, a contradiction.
For the boundary, Theorem
\ref{p2:thm:boundary-jets-round2} gives the same exclusion for
$d\mapsto\partial_dB(d,2)$: its complete finite jets at $2$ consist
of the algebraically independent positive-order derivatives of $B$.
If $B$ were differentially algebraic, so would be its derivative;
this follows also from the derivative-stable finite-transcendence-
degree field constructed in Lemma~\ref{p2:lem:constant-elimination-round2}.
Finally that lemma excludes equations over $\mathbb C(d)$ for both
functions. Analytic continuation on $d>0$ permits any asserted
interval relation to be evaluated near $d=2$.
\end{proof}

% Suggested bibliography entry:
% \bibitem{Gorgens2014} M. G\"orgens,
% Optimal stopping of an $\alpha$-Brownian bridge,
% arXiv:1409.5285v1 (18 September 2014), Sections 1 and 4.


% Source fragment: round2_inequalities.tex
% Fragment: requires amsmath, amssymb, amsthm, and theorem/proposition/lemma environments.
% Bibliography keys: qln2025, shuang2026, pudim2026.
\section{A sharp exponent range for a reciprocal digamma inequality}
\label{p2:sec:qln-op3}

Qi, Lim, and Nantomah \cite[Open Problem~3]{p2:qln2025} ask for the largest
range of real $\beta$ such that
\[
 F_\beta(x)=x^\beta h(x)-P(x)>0\quad(x>0),\qquad
 h(x)=\psi(x)+x\psi'(x),\quad P(x)=\psi(x)\psi(1/x).
\]
The following theorem gives the range and specifies its non-elementary
endpoint by one equation having exactly one solution.  Its proof does not
rely on numerical optimization or interval subdivision.

\begin{theorem}[Exact range for the digamma problem]
\label{p2:thm:qln-op3}
Let $r$ be the unique positive zero of $\psi$, and, for $1/r<x<1$, put
\[
 L(x)=\log\frac{\psi(x)\psi(1/x)}{\psi(x)+x\psi'(x)},\qquad
 G(x)=x\log x\,L'(x)-L(x).
\]
There is exactly one $\xi\in(1/r,1)$ with $G(\xi)=0$.  Define
\[
 B=\frac{L(\xi)}{\log\xi}.
\]
Then
\[
 \boxed{\quad F_\beta(x)>0\text{ for every }x>0
                 \quad\Longleftrightarrow\quad -1\le\beta<B.\quad}
\]
Moreover $B>0$, and
\[
 (x\partial_x)^2 L(x)<-\frac{22}{81}\qquad(1/r<x<1).
\]
Numerically, for orientation only,
$\xi\approx0.76859976$ and $B\approx5.97283686$.
\end{theorem}

\begin{proof}
We give the elementary estimates used below explicitly.  The standard series
and integral representations imply
\begin{align}
 \psi(x)&=-\gamma-\frac1x+
       \sum_{k=1}^{\infty}\frac{x}{k(k+x)},\label{p2:eq:qln-psi-series}\\
 h(x)&=-\gamma+\sum_{k=1}^{\infty}
       \frac{x(2k+x)}{k(k+x)^2},\label{p2:eq:qln-h-series}\\
 h'(x)&=2\sum_{k=1}^{\infty}\frac{k}{(x+k)^3}
       =\int_0^\infty e^{-xt}
         \left(\frac{t}{2\sinh(t/2)}\right)^2\,dt>0,\label{p2:eq:qln-hprime}\\
 h''(x)&=-6\sum_{k=1}^{\infty}\frac{k}{(x+k)^4}<0.
\end{align}
All differentiations are locally uniform on $x>0$.  We use the elementary
rational bounds
\[
 0<\gamma<\frac{29}{50},\qquad \log2<\frac7{10},\qquad
 \frac85<\zeta(2)<\frac53,\qquad \zeta(3)>\frac65.
\]
For completeness, here are finite certificates for these constants.
Integral comparison gives
\[
 \zeta(2)>\sum_{k=1}^3\frac1{k^2}+\frac14=\frac{29}{18}>\frac85,
 \quad
 \zeta(2)<\sum_{k=1}^5\frac1{k^2}+\frac15=\frac{5989}{3600}<\frac53,
\]
and
\[
 \zeta(3)>\sum_{k=1}^6\frac1{k^3}+\frac1{98}
 =\frac{1411783}{1176000}>\frac65.
\]
The positive series $\log[(1+z)/(1-z)]=2\sum_{j\ge0}z^{2j+1}/(2j+1)$
gives $\log2<2/3+1/36=25/36<7/10$ by taking $z=1/3$
and bounding the denominators of all terms after the first by $3$.
The midpoint integral comparison for $1/t$ yields
$\gamma<H_4-\log(9/2)$.  The same logarithm series, now at $z=7/11$, gives
\[
 \log(9/2)>2\sum_{j=0}^5\frac{(7/11)^{2j+1}}{2j+1}
 =\frac{451}{300}+\frac{213407380807}{2824585539048900}
 >\frac{451}{300}.
\]
Since $H_4=25/12$, this proves $\gamma<29/50$.
Finally $\gamma\ge1-\log2>0$.  No decimal approximation is needed.
Since $\psi(1)=-\gamma<0$ and
$\psi(3/2)=2-\gamma-2\log2>0$, we have $1<r<3/2$.
The first three terms of \eqref{p2:eq:qln-h-series} give
\[
 h(1/3)>\frac7{16}+\frac{13}{98}+\frac{19}{300}-\frac{29}{50}
       =\frac{629}{11760}>0.
\]
Also
\[
 h(1/2)=-\gamma-2\log2+\tfrac32\zeta(2)
     >\frac{21}{50}>\gamma^2,
 \qquad \frac{h(1)}2=\frac{\zeta(2)-\gamma}{2}
     >\frac{51}{100}>\gamma^2.                 \tag{\(*\)}
\]
We shall need the slightly more precise, still rational, estimate
\[
 h(2/3)>\frac7{10}.                            \tag{\(**\)}
\]
Indeed, for $x=2/3$ the summand
$f(t)=x(2t+x)/[t(t+x)^2]=1/t-t/(t+x)^2$ is decreasing on $t>0$.
Thus its tail starting at $k=7$ exceeds
\[
 \int_7^\infty f(t)\,dt
 =\log(1+x/7)+\frac{x}{7+x}>\frac{2x}{7+x}=\frac4{23}.
\]
The sum of its first six terms is $1523252659/1370784800$, whence
\[
 h(2/3)>\frac{889267245}{1261122016}
       =\frac7{10}+\frac{32409169}{6305610080}>\frac7{10}.
\]

We next record two curvature estimates.  For $z<r$ define
\[
 u(z)=\frac{z\psi'(z)}{-\psi(z)},\qquad
 R(z)=\frac{-z\psi''(z)}{\psi'(z)},\qquad
 K_\psi(z)=(z\partial_z)^2\log[-\psi(z)].
\]
The polygamma series give $R(z)<2$, because
\[
 \frac{R(z)}2=
 \frac{\sum_{k\ge0}z(z+k)^{-3}}{\sum_{k\ge0}(z+k)^{-2}}<1.
\]
Direct differentiation therefore yields
\[
 K_\psi(z)=u(z)(R(z)-1)-u(z)^2<u(z)-u(z)^2.    \tag{\(\dagger\)}
\]
If $1/r<z<r$, then $z>2/3$, so $h(z)>0$ and $u(z)>1$.
Consequently $K_\psi(z)<0$ throughout that interval.
In particular, $t\mapsto\log P(e^t)$ is even and strictly concave for
$|t|<\log r$.  It attains its maximum at zero.  Outside $(1/r,r)$
the product $P$ is nonpositive, so in fact
\[
 P(x)\le\gamma^2\quad\text{for every }x>0.     \tag{\(\ddagger\)}
\]

When $1\le z<r<3/2$, strict convexity in the trapezoidal comparison for
$\sum_{k\ge0}(z+k)^{-2}$ gives
\[
 \psi'(z)>\frac1z+\frac1{2z^2}.
\]
Using $-\psi(z)\le\gamma<3/5$, we obtain $u(z)>20/9$.  It follows from
\((\dagger)\) that
\[
 K_\psi(z)<-\frac{220}{81}\qquad(1\le z<r).   \tag{\(\dagger\dagger\)}
\]

Let $2/3\le x\le1$ and set $v=xh'(x)/h(x)$.  The function
$K(t)=[t/(2\sinh(t/2))]^2$ is decreasing on $(0,\infty)$.
After $u=xt$ in \eqref{p2:eq:qln-hprime},
\[
 xh'(x)=\int_0^\infty e^{-u}K(u/x)\,du
\]
is increasing in $x$.  Therefore
\[
 xh'(x)\le h'(1)=2\zeta(2)-2\zeta(3)<\frac{14}{15},
 \qquad 0<v<\frac43,
\]
where \((**)\) was used.  The series for $h'$ and $h''$ also give
\[
 \frac{-xh''(x)}{h'(x)}<\frac{3x}{x+1}\le\frac32.
\]
Hence
\[
 K_h(x):=(x\partial_x)^2\log h(x)
 =v\left(1+\frac{xh''(x)}{h'(x)}\right)-v^2
 >-\frac v2-v^2>-\frac{22}{9}.                \tag{\(\star\)}
\]
For $1/r<x<1$, both arguments in the product defining $P$ are below $r$,
and $1/x\ge1$.  Combining the curvature bounds gives
\[
 (x\partial_x)^2 L(x)
 =K_\psi(x)+K_\psi(1/x)-K_h(x)
 <0-\frac{220}{81}+\frac{22}{9}
 =-\frac{22}{81}.                            \tag{\(\star\star\)}
\]

We now prove the admissible lower endpoint directly.  If $0<x<1/3$, put
$w=-\psi(x)>0$.  From \eqref{p2:eq:qln-psi-series},
$w<1/x+\gamma$, while $\psi'(x)>1/x^2$.  Thus
\[
 \frac{\psi'(x)}w>\frac1{x(1+\gamma x)}>\frac1x-\gamma.
\]
Since $\psi(1/x)>\psi(3)=3/2-\gamma>\gamma$, it follows that
\[
 F_{-1}(x)=w\left(\frac{\psi'(x)}w+\psi(1/x)-\frac1x\right)>0.
\]
For $1/3\le x\le1/2$, one has $h(x)>0$ and $P(x)<0$.
For $1/2\le x\le1$, \((*)\) and \((\ddagger)\) give
$h(x)/x\ge h(1/2)>\gamma^2\ge P(x)$.
For $1\le x\le r<2$, they give
$h(x)/x>h(1)/2>\gamma^2\ge P(x)$.
Finally, for $x\ge r$, we have $h(x)>0$ and $P(x)\le0$.
Consequently $F_{-1}(x)>0$ for every $x>0$.

Let $a\in(0,1/3)$ be the unique zero of $h$.  If $\beta\ge-1$, then
$F_\beta\ge F_{-1}>0$ on $(0,a)$, because $x<1$ and $h(x)<0$;
the same comparison holds on $[1,\infty)$, where $h>0$.
On $[a,1/r]$, the two terms have opposite signs, and $F_\beta>0$.
The only remaining interval is $(1/r,1)$, where
\[
 F_\beta(x)>0\quad\Longleftrightarrow\quad
 \beta<Q(x):=\frac{L(x)}{\log x}.             \tag{\(\clubsuit\)}
\]
On this interval $h>7/10$ and $P\le\gamma^2$, so $L<0$ and $Q>0$.
Moreover $Q(x)\to+\infty$ at both endpoints: at $1/r$ the numerator
$P$ has a simple zero, while at $1$ the finite limit
$L(1)=\log[\gamma^2/(\zeta(2)-\gamma)]$ is strictly negative.
By \((\star\star)\),
\[
 G'(x)=\frac{\log x}{x}(x\partial_x)^2L(x)>0.
\]
The simple zero of $P$ gives $G(x)\to-\infty$ as $x\downarrow1/r$;
as $x\uparrow1$, $G(x)\to-L(1)>0$.
Thus $G$ has precisely one zero $\xi$, and
\[
 Q'(x)=\frac{G(x)}{x(\log x)^2}
\]
shows that $Q$ decreases up to $\xi$ and increases thereafter.
Its minimum is $B=Q(\xi)>0$.  Relation \((\clubsuit)\) yields exactly
$-1\le\beta<B$, including the strict exclusion of $B$ by equality at $\xi$.
Finally, if $\beta<-1$, then as $x\downarrow0$,
\[
 h(x)=-\gamma+O(x),\qquad
 P(x)=\frac{\log x}{x}+O(|\log x|+1),
\]
so $F_\beta(x)\sim-\gamma x^\beta<0$.
This proves necessity of the lower endpoint and completes the proof.
\end{proof}

\paragraph{Interpretation and provenance.}
The endpoint is a unique tangency slope for the strictly concave function
$t\mapsto L(e^t)$; the quantitative curvature bound prevents additional
contacts.  This is an analytic optimization result, not a claim about the
irrationality or transcendence of $\gamma$.
A targeted literature search on 18 October 2026 found the original problem
and the 2026 alternative proofs of its $\beta=0$ case \cite{p2:shuang2026}, but
no previous solution of this exponent range.  This is not an exhaustive
priority claim.

\section{The higher derivative exponent range: an independent proof}
\label{p2:sec:qln-op4-independent}
The following is independently proved here, but a public repository
\cite{p2:pudim2026} already claims the corresponding classification outside
$n=2$.  Its priority is unresolved, and it is excluded from any strict count
of newly resolved open questions.  All orders below are parts of the single
Open Problem~4 in \cite{p2:qln2025}.

\begin{theorem}[Odd orders and all even orders at least four]
For an integer $n\ge1$, let
\[
 C_n(x)=\psi^{(n)}(x)+x\psi^{(n+1)}(x),\qquad
 P_n(x)=\psi^{(n)}(x)\psi^{(n)}(1/x).
\]
If $n$ is odd, $x^\beta C_n(x)-P_n(x)<0$ for every real $\beta$ and every
$x>0$.  If $n\ge4$ is even, this inequality holds for every $x>0$ exactly
when $n\le\beta\le n+1$.  For those admissible exponents,
\[
 0<\frac{x^\beta C_n(x)}{P_n(x)}
 <\frac{n^2+2n-1}{n!}\le\frac{23}{24}\qquad(x>0).
\]
\end{theorem}
\begin{proof}
Write $Z_s(x)=\sum_{k\ge0}(x+k)^{-s}$ and
$A_n(x)=(n+1)xZ_{n+2}(x)-Z_{n+1}(x)$.  Then
\[
 C_n=(-1)^n n!A_n,\qquad
 P_n=(n!)^2Z_{n+1}(x)Z_{n+1}(1/x).
\]
Integration by parts in the Hurwitz-zeta integral gives
\[
 A_n(x)=\frac1{n!}\int_0^\infty e^{-xt}
       \frac{t^n}{1-e^{-t}}
       \left(n-\frac{t}{e^t-1}\right)dt>0.
\]
The boundary terms vanish even when $n=1$.  Thus the odd case is immediate.
For $s>1$, integral comparison gives
\[
 \frac{x^{1-s}}{s-1}<Z_s(x)
 <x^{-s}+\frac{x^{1-s}}{s-1},\qquad Z_s(x)>x^{-s}.
\]
Consequently
\[
 A_n(x)<(n+1)x^{-n-1}+\frac{n-1}{n}x^{-n},
 \quad P_n(x)>\frac{(n!)^2}{nx},
 \quad P_n(x)>\frac{(n!)^2x}{n}.
\]
For even $n\ge4$ these imply, respectively,
\[
 \frac{x^nC_n(x)}{P_n(x)}
 <\frac{n(n+1)+(n-1)x}{n!}
 \le\frac{n^2+2n-1}{n!}\quad(0<x\le1),
\]
and
\[
 \frac{x^{n+1}C_n(x)}{P_n(x)}
 <\frac{n(n+1)/x+(n-1)}{n!}
 \le\frac{n^2+2n-1}{n!}\quad(x\ge1).
\]
The last coefficient equals $23/24$ at $n=4$ and decreases thereafter.
The two displayed estimates prove sufficiency for every $\beta\in[n,n+1]$.
For necessity, the standard integral comparison asymptotics yield
\[
 \frac{x^\beta C_n(x)}{P_n(x)}
 \sim\frac{n}{(n-1)!}x^{\beta-n}\quad(x\downarrow0),
 \qquad
 \frac{x^\beta C_n(x)}{P_n(x)}
 \sim\frac{n-1}{n!}x^{\beta-n-1}\quad(x\to\infty).
\]
The first diverges for $\beta<n$, and the second for $\beta>n+1$.
\end{proof}

% Suggested bibliography:
% \bibitem{qln2025} F. Qi, D. Lim, K. Nantomah,
% Monotonicity and positivity of several functions involving ratios and products
% of polygamma functions, J. Inequal. Appl. 2025, 5 (2025),
% doi:10.1186/s13660-024-03245-8.
% \bibitem{shuang2026} Y. Shuang, D. Lim, F. Qi,
% Alternative proofs for the positivity of a function involving the product
% of two digamma functions, Aust. J. Math. Anal. Appl.23(1),Art.5 (2026),
% https://ajmaa.org/searchroot/files/pdf/v23n1/v23i1p5.pdf.
% \bibitem{pudim2026} Pudim Project, public research archive (2026),
% https://github.com/pudim-project/pudim-ai-demo-zetalaw;
% node wiki/latest/wiki/nodes/T-Q2-op4-reduction-n2.md,
% and APP-0011. Accessed 18 October 2026. Public claims, not peer review.


% Source fragment: round2_gurland_fragment.tex
\section{Gurland's ratio: global convergence and a sharp location theorem}
\label{p2:sec:gurland2}

Wi\'sniewska \cite[Section 2, Open problems 1--3]{p2:Wis2025} asks about
global convergence of a Hurwitz-zeta expansion and the position of its
mean-value parameter. Problems 1 and 2 are equivalent through the exact
finite expansion and are treated as one mathematical question here.
The location question has the following quantitative answer.

Put $s=(x+y)/2$, $d=|x-y|/2$, and $g=\sqrt{xy}=\sqrt{s^2-d^2}$.
Throughout this section $x,y>0$. Define
\[
 L=\log\frac{\Gamma(1+x)\Gamma(1+y)}{\Gamma(1+s)^2},\qquad
 S_m=\sum_{k=1}^{m-1}\frac{d^{2k}}{k}\zeta(2k,1+s).
\]

\begin{theorem}[Unconditional convergence of the published expansion]
For every $x,y>0$,
\[
 L=\sum_{k=1}^{\infty}\frac{d^{2k}}{k}\zeta(2k,1+s).
\]
Writing $q=d/(1+s)<1$, the remainder satisfies, for $m\ge2$,
\[
 0\le L-S_m\le
 \frac{d^{2m}\zeta(2m,1+s)}{m(1-q^2)}
 \le\frac{q^{2m}}{m(1-q^2)}
       \left(1+\frac{1+s}{2m-1}\right).
\]
In particular it tends to zero, without the extra sufficient condition
imposed in \cite[Corollary 1]{p2:Wis2025}.
\end{theorem}
\begin{proof}
The Weierstrass product, with the linear terms cancelling, gives
\[
 L=\sum_{j=1}^{\infty}-\log\left(1-\frac{d^2}{(j+s)^2}\right).
\]
Since $d<s<1+s$, every summand can be expanded into its positive
logarithmic series. Tonelli's theorem interchanges the two sums.
For $0\le u\le q^2$,
$\sum_{k=m}^{\infty}u^k/k\le u^m/[m(1-q^2)]$.
Summing this estimate proves the first bound. The integral test gives
\[
 \zeta(2m,1+s)\le(1+s)^{-2m}
       +\frac{(1+s)^{1-2m}}{2m-1},
\]
which proves the second. The diagonal case $d=0$ is immediate.
\end{proof}

For $d>0$ let $t$ be the unique number satisfying
$L=d^2\psi'(1+t)$. Uniqueness follows because $\psi'$ is strictly
decreasing on $(0,\infty)$.
For $0<r<1$ set
\[
 h(r)=(1+r)\log(1+r)+(1-r)\log(1-r),\qquad
 H(s,d)=\frac{s r^2}{h(r)},\quad r=d/s.
\]

\begin{theorem}[Sharp uniform location and distance from the midpoint]
\label{p2:thm:gurland-location}
For $x\ne y$,
\[
 \frac{2s+g}{3}<H(s,d)<t<s.
\]
Thus the parameter always lies strictly in the upper third of $(g,s)$.
In particular its distance from the midpoint is positive and obeys
\[
 \frac{s-g}{6}
 <H(s,d)-\frac{s+g}{2}
 <t-\frac{s+g}{2}<\frac{s-g}{2}.
\]
The best possible constant $c$ in the uniform inequality
$t>g+c(s-g)$ is $c=2/3$.
More precisely, for each fixed $r\in(0,1)$, with $d=rs$,
\[
 \lim_{s\to\infty}\frac{t}{s}=\frac{r^2}{h(r)},\qquad
 \lim_{s\downarrow0}\frac{t}{s}=1.
\]
Consequently both the shape-dependent lower bound $H(s,d)/s$
and the upper bound $1$ are best possible when $d/s$ is fixed.
\end{theorem}
\begin{proof}
Let $U$ have probability density $1-|u|$ on $[-1,1]$, and put
$X=s+dU$. Twice integrating the second derivative of $\log\Gamma$
gives the exact identity
\[
 \frac{L}{d^2}=\mathbb E\, f(X),\qquad f(v)=\psi'(1+v).
\]
Since $f$ is strictly convex and decreasing,
$\mathbb E f(X)>f(s)$ and hence $t<s$.

We next prove that $F(w)=f(1/w)$ is strictly concave for $w>0$.
The derivative test reduces this assertion to
\[
 2\zeta(3,1+v)-3v\zeta(4,1+v)>0\quad(v>0).
\]
For completeness, write $a(u)=1/(e^u-1)$ and use the Laplace
integral for Hurwitz zeta. Integration by parts, whose boundary
terms vanish, yields
\[
 2\zeta(3,1+v)-3v\zeta(4,1+v)
 =\frac12\int_0^\infty e^{-vu}u^2[-a(u)-u a'(u)]\,du>0.
\]
Indeed $u/(e^u-1)$ is strictly decreasing on $(0,\infty)$:
its derivative has numerator $e^u(1-u)-1<0$.
Equivalently,
$f''(v)/f'(v)-(-2/v)=2/v-3\zeta(4,1+v)/\zeta(3,1+v)>0$,
which gives $F''<0$ by the chain rule.
Strict Jensen inequality now gives
\[
 \mathbb E f(X)=\mathbb E F(1/X)
 <F(\mathbb E(1/X))
 = f\bigl((\mathbb E(1/X))^{-1}\bigr).
\]
Since $f$ decreases, $t>(\mathbb E(1/X))^{-1}$.
Direct integration of the triangular density shows that this
last reciprocal is exactly $H(s,d)$.

To prove the remaining explicit bound, put $q=\sqrt{1-r^2}$ and
$A(r)=3r^2/(2+q)-h(r)$. Direct differentiation gives
\[
 A''(r)=\frac{(1-q)^2(6+5q-2q^2)}{q^3(q+2)^3}>0,
 \qquad A(0)=A'(0)=0.
\]
Thus $h(r)<3r^2/(2+q)$, or $H(s,d)>(2s+g)/3$.
Subtracting the midpoint proves the asserted distance bounds.

For the limiting assertions use
$\psi'(1+v)=v^{-1}+O(v^{-2})$ as $v\to\infty$.
Because $r<1$ is fixed, this is uniform over
$v=s(1+ru)$, $|u|\le1$. Therefore
\[
 \mathbb E f(X)=\frac{h(r)}{s r^2}+O(s^{-2}).
\]
The proved inequalities bound $t/s$ away from zero, so the same
asymptotic for $f(t)$ proves the first limit. For $s\downarrow0$,
Taylor expansion around zero, using $\mathbb E X=s$ and
$\mathbb E(X-s)^2=d^2/6$, gives
\[
 t=s-\frac{\zeta(4)}{4\zeta(3)}d^2+O(s^3),
\]
which proves the second limit. Finally,
$r^2/h(r)=1-r^2/6+O(r^4)$ and
$\sqrt{1-r^2}=1-r^2/2+O(r^4)$. First taking $s\to\infty$
and then $r\downarrow0$ makes $(t-g)/(s-g)$ tend to $2/3$.
No larger uniform constant is possible.
\end{proof}

\paragraph{Interpretation and scope.}
The convergence result answers both of the first two source questions;
it is one closure, not two independent conjectures. The location theorem
answers the third by a sharp universal fraction and sharper bounds
depending on the ratio of the arguments. These are special-function
results, not an irrationality theorem for $\gamma$. Their connection
with $\gamma$ is through logarithmic derivatives of the gamma function;
the constant term cancels in the symmetric ratio.


% Source fragment: round2_gompertz.tex
% Main independently audited source-conjecture closure:
\section{A counterexample to Bolbachan's digamma limit conjecture}
\label{p2:sec:bolbachan-counterexample}

Bolbachan's preprint \cite[Conjecture~1.4]{p2:Bolbachan2011}
proposes the following limit for every $u>0$:
\begin{equation}\label{p2:eq:bol-source}
 \psi(u)=\log u+\lim_{m\to\infty} L_m(u),\qquad
 L_m(u)=\frac1{m!}\sum_{k=1}^m(-1)^k\binom mk
 \frac{A_{k,m+1}}{k!}
 \int_0^\infty x^{k-1}e^{-x}\log(1+x/u)\,dx,
\end{equation}
where
\[
 A_{k,n}=\sum_{t=2}^n\genfrac{\{}{\}}{0pt}{}{n}{t}
 \sum_{w=1}^{t-1}(-k)^{t-w}
 \sum_{j=1}^w(-1)^j B_j\genfrac{[}{]}{0pt}{}{w}{j}.
\]
Here the Stirling numbers of the first kind in square brackets are
unsigned.  We first use the convention $B_1=-1/2$.

\begin{theorem}\label{p2:thm:bol-counterexample}
For every fixed real $u>0$, the sequence in \eqref{p2:eq:bol-source}
satisfies
\[
 L_m(u)\longrightarrow+\infty,
 \qquad
 \lim_{m\to\infty}L_m(u)^{1/m}=\frac1{\log2}.
\]
Consequently Conjecture~1.4 is false as printed.  If instead Bernoulli
numbers with $B_1=+1/2$ are intended, then $L_m(u)\to-\infty$ and
$|L_m(u)|^{1/m}\to1/\log2$; the limit assertion still fails.
\end{theorem}

\subsection{Exact reduction to a coefficient integral}
Put
\[
 c_w=\sum_{j=1}^w(-1)^jB_j\genfrac{[}{]}{0pt}{}{w}{j},
 \qquad C(y)=\sum_{w\ge1}c_w\frac{y^w}{w!}.
\]
The Stirling exponential generating function and the Bernoulli generating
function give, as convergent identities for $|y|<1$,
\[
 1+C(y)=\sum_{j\ge0}\frac{B_j\log(1-y)^j}{j!}
 =\frac{-\log(1-y)}{y}.
\]
In particular
\begin{equation}\label{p2:eq:bol-C}
 c_w=\frac{w!}{w+1},\qquad
 C(y)=-\frac{\log(1-y)}{y}-1.
\end{equation}
The apparent singularity at $y=0$ is removable.

The expression $A_{k,m+1}/k$ extends polynomially to $k=0$ and has degree
at most $m-1$.  Define
\[
 T_m(z)=\frac1{m!}\sum_{k=0}^m(-1)^k\binom mk
 \left.\frac{A_{k,m+1}}k\right|_{\text{polynomial extension}}z^k.
\]
The $m$th finite difference of a polynomial of degree at most $m-1$
vanishes, so
\begin{equation}\label{p2:eq:bol-Tone}T_m(1)=0.\end{equation}
Writing $Y(v)=e^v-1$, set
\begin{equation}\label{p2:eq:bol-J}
 J_m(v,z)=\int_0^{Y(v)}C(w)
       \left(1-ze^{w-Y(v)}\right)^m\,dw.
\end{equation}
Then the following exact identity holds:
\begin{equation}\label{p2:eq:bol-TJ}
 T_m(z)=-(m+1)[v^{m+1}]J_m(v,z).
\end{equation}
For verification, let
$H_t(k)=\sum_{w=1}^{t-1}(-k)^{t-w}c_w$, with $H_0=H_1=0$.
The recurrence $H_{t+1}=-kH_t-kc_t$ gives
\[
 \sum_{t\ge0}H_t(k)\frac{y^t}{t!}
 =-k e^{-ky}\int_0^y e^{kw}C(w)\,dw.
\]
Composition with $y=e^v-1$ performs the Stirling transform, hence
\[
 \frac{A_{k,m+1}}k
 =-(m+1)![v^{m+1}]\int_0^{Y(v)}
                  e^{k(w-Y(v))}C(w)\,dw.
\]
Summing the binomial expansion proves \eqref{p2:eq:bol-TJ}, including $k=0$
by polynomial extension.

The elementary identity
\[
 \frac1{k!}\int_0^\infty x^{k-1}e^{-x}\log(1+x/u)\,dx
 =\frac1k\int_0^\infty\frac{e^{-ut}}t
                   \bigl(1-(1+t)^{-k}\bigr)\,dt
\]
follows from the Laplace integral for the logarithm and Tonelli's theorem.
Therefore, putting $z(t)=(1+t)^{-1}$, the original finite sum becomes
\begin{equation}\label{p2:eq:bol-LT}
 L_m(u)=-\int_0^\infty\frac{e^{-ut}}t
                 \bigl(T_m(z(t))-T_m(1)\bigr)\,dt.
\end{equation}
Keeping the subtraction at $z=1$ is essential for uniform estimates near
$t=0$.

\subsection{A disk estimate with rational constants}
Let $R=7/10$ and $0\le s\le1$.  We claim
\begin{equation}\label{p2:eq:bol-disk}
 \left|1-\exp\{-s(e^v-1)\}\right|
 \le\frac{318799}{320000}<1\qquad (|v|\le R).
\end{equation}
Write $\mathcal B_n(x)=\sum_{j=0}^n\genfrac{\{}{\}}{0pt}{}{n}{j}x^j$
for the Touchard polynomials and $b_n=\mathcal B_n(1)$ for the Bell numbers.
The exponential generating function is
\[
 e^{-s(e^v-1)}=\sum_{n\ge0}\mathcal B_n(-s)\frac{v^n}{n!}.
\]
For $1\le n\le5$, the Bernstein coefficients of $\mathcal B_n(-s)$
on $[0,1]$, in increasing order, are
\[
\begin{array}{c|l|c}
 n&\text{Bernstein coefficients}&M_n\\ \hline
 1&0,-1&1\\
 2&0,-1/2,0&1/2\\
 3&0,-1/3,1/3,1&1\\
 4&0,-1/4,2/3,5/4,1&5/4\\
 5&0,-1/5,11/10,7/5,1/5,-2&2
\end{array}
\]
and their absolute values are bounded by the last column.  For $n\ge6$,
$|\mathcal B_n(-s)|\le b_n$.  Consequently the left side of
\eqref{p2:eq:bol-disk} is at most
\[
 e^{e^R-1}-1+\sum_{n=1}^5(M_n-b_n)\frac{R^n}{n!}
 <\frac{14}{5}-1+\sum_{n=1}^5(M_n-b_n)\frac{(7/10)^n}{n!}
 =\frac{318799}{320000}.
\]
The elementary bounds $e^{7/10}<101/50$ and $e^{51/50}<14/5$ justify
$e^{e^R-1}<14/5$; they can also be checked directly from the exponential
series with its geometric remainder bound.

For $z\in[0,1]$, convexity of the closed unit disk now gives
\begin{equation}\label{p2:eq:bol-kernelbound}
 |1-z e^{-sY(v)}|\le1 \qquad (|v|\le R,\ 0\le s\le1).
\end{equation}
Indeed $1-z e^{-sY}=(1-z)+z(1-e^{-sY})$.

\subsection{The branch cut gives a positive dominant contribution}
Put $\rho=\log2<R$.  Use the principal branch of $\log(1-w)$ in
\eqref{p2:eq:bol-C}.  Formula \eqref{p2:eq:bol-J}, with the straight path
$w=\tau Y(v)$, defines an analytic function on
\[
 \{v:|v|<R\}\setminus[\rho,R).
\]
There are no other cuts in this disk: since $|\operatorname{Im}v|<\pi$,
$Y(v)$ is real and at least $1$ exactly when $v$ is real and at least
$\rho$.  For real $\rho<v<R$, its upper minus lower jump is
\begin{equation}\label{p2:eq:bol-jump}
 J_m(v+i0,z)-J_m(v-i0,z)
 =2\pi i\int_1^{e^v-1}
            \frac{(1-z e^{w-(e^v-1)})^m}{w}\,dw.
\end{equation}
The positive sign follows because $\log(1-w)$ has upper boundary value
$\log(w-1)-i\pi$ for $w>1$.

There exists a finite constant $M$, independent of $m$ and $z$, such that
on the circle $|v|=R$, with either boundary value at its positive real
point,
\begin{equation}\label{p2:eq:bol-outerbound}
 |J_m(v,z)-J_m(v,1)|\le M m(1-z).
\end{equation}
To see this, factor the difference of the two $m$th powers in
\eqref{p2:eq:bol-J}, use \eqref{p2:eq:bol-kernelbound}, and bound
$|e^{w-Y(v)}|$ by $e^{|Y(v)|}\le e^{e^R-1}$.  The remaining integral
$\int_0^1|Y(v)C(\tau Y(v))|\,d\tau$ is uniformly bounded: $C$ is regular
at zero and has only an integrable logarithmic singularity at $1$.
The same assertion holds on any smaller circle $|v|=r<\rho$.

Deform the coefficient contour in \eqref{p2:eq:bol-TJ} to the circle of
radius $R$ and the two sides of the cut.  Small indentations at $\rho$
make no contribution because the logarithmic singularity is integrable.
Equations \eqref{p2:eq:bol-LT}--\eqref{p2:eq:bol-outerbound} give
\begin{align}\label{p2:eq:bol-positivecut}
 L_m(u)={}&(m+1)\int_\rho^R\frac{dv}{v^{m+2}}
 \int_1^{e^v-1}\frac{dw}{w}\int_0^\infty\frac{e^{-ut}}t\\
 &\quad\times\left[
  \left(1-\frac{e^{w-(e^v-1)}}{1+t}\right)^m
  -\left(1-e^{w-(e^v-1)}\right)^m\right]dt
 +O_u(m^2R^{-m}).\notag
\end{align}
All three integrations in the displayed main term have nonnegative
integrands.  The error follows from \eqref{p2:eq:bol-outerbound} and
$\int_0^\infty e^{-ut}(1-z(t))\,dt/t
 =\int_0^\infty e^{-ut}/(1+t)\,dt<\infty$.
These estimates also justify all contour--integral interchanges.

Choose $\varepsilon>0$ with $\rho+2\varepsilon<R$, put
$h=e^\varepsilon-1$, and restrict the main integral to
\[
 \rho+\varepsilon\le v\le\rho+2\varepsilon,
 \qquad 1\le w\le1+h,
 \qquad T\le t\le T+1.
\]
This rectangle lies inside its integration domain.  On it, set
\[
 a=\frac{T}{T+1},\qquad
 b=1-\exp(2-2e^{2\varepsilon}).
\]
The bracket in \eqref{p2:eq:bol-positivecut} is at least $a^m-b^m$.
Choose $T$ sufficiently large that $a>b$.  It follows that for all
sufficiently large $m$,
\begin{equation}\label{p2:eq:bol-lower}
 L_m(u)\ge c_{u,\varepsilon,T}(m+1)
       \left(\frac{a}{\rho+2\varepsilon}\right)^m
       -C_u m^2R^{-m},
 \qquad c_{u,\varepsilon,T}>0.
\end{equation}
By taking $\varepsilon$ small and $T$ large we can make
$a/(\rho+2\varepsilon)$ arbitrarily close to $1/\rho$ from below,
and in particular strictly greater than $1/R$.  Thus the positive term
in \eqref{p2:eq:bol-lower} dominates the error, proving divergence and
\[
 \liminf_{m\to\infty}L_m(u)^{1/m}\ge\frac1\rho.
\]
On the other hand, Cauchy's estimate on any circle $|v|=r<\rho$, together
with the smaller-circle version of \eqref{p2:eq:bol-outerbound}, gives
$|L_m(u)|\le C_{u,r}m(m+1)r^{-m-1}$.  Letting $r\uparrow\rho$ proves the
matching upper bound.

Finally, replacing $B_1=-1/2$ by $B_1=+1/2$ changes $c_w$ by
$-(w-1)!$, hence changes $C(w)$ to
\[
 C_+(w)=\left(1-\frac1w\right)\log(1-w)-1.
\]
Its upper minus lower jump for $w>1$ is
$-2\pi i(1-1/w)$.  The same proof applies with a negative cut term and
the positive weight $1-1/w$ in place of $1/w$.  In the lower bound one
may restrict $w$ to $[1+h/2,1+h]$, where this weight is bounded below by
a positive constant.  This proves the alternative-convention assertion
and completes the theorem.


\section{The fractional-part and gcd obstructions that remain}
Murabayashi--Yoshida's Conjecture~7.5 asks for an infinite nondecreasing
subsequence of $\delta\{q_{2k}\}-\{p_{2k}\}$ and an infinite nonincreasing
subsequence of $\delta\{q_{2k-1}\}-\{p_{2k-1}\}$.  Their Theorem~7.4 shows
that either would imply irrationality of $\delta$.  We do not prove this
conjecture.  A general monotone-subsequence theorem cannot choose the
required direction, and fixed-order Laguerre asymptotics do not control
fractional parts.

\begin{proposition}[Finite-chain obstruction for rational constants]
Let $\alpha=a/b>0$ be rational in lowest terms, with $a,b$ positive
integers. If
$r_n=\alpha q_n-p_n$ strictly decreases inside $(0,1/b)$, every
nondecreasing subsequence of $s_n=\alpha\{q_n\}-\{p_n\}$ has at most
$a+b$ terms.  If instead $r_n$ strictly increases inside $(-1/b,0)$,
the same bound holds for nonincreasing subsequences of $s_n$.
\end{proposition}
\begin{proof}
The number $h_n=b s_n-b r_n=b\lfloor p_n\rfloor-a\lfloor q_n\rfloor$
is integral.  In the first case it lies in $\{-b,\ldots,a-1\}$ and
strictly increases along the specified subsequence.  In the second it
lies in $\{-b+1,\ldots,a\}$ and strictly decreases.  Both sets have
$a+b$ elements.
\end{proof}

For the integer Laguerre approximants $Q_n=n!q_{2n}$ and $P_n=n!p_{2n}$,
we have
\[
 U_{n+1}=2(n+1)U_n-n^2U_{n-1},\qquad
 (Q_0,Q_1)=(1,2),\quad(P_0,P_1)=(0,1).
\]
Write $g_n=\gcd(Q_n,P_n)$ and $\Lambda_n=\operatorname{lcm}(1,\ldots,n)$.
\begin{proposition}[Exact congruence and reset identities]
For nonnegative integers $n,r$ and any positive integer $M$,
\[
 n\equiv r\pmod M\quad\Longrightarrow\quad Q_n\equiv Q_r\pmod M.
\]
For $n\ge M$ one has
$P_n\equiv P_MQ_{n\bmod M}\pmod M$.
Moreover, for $n\ge1$,
\[
 \gcd(Q_n,n)=1,\qquad
 \gcd(Q_n,\Lambda_n^2)\mid g_n\mid((n-1)!)^2.
\]
\end{proposition}
\begin{proof}
Writing $(n)_{\underline j}=n(n-1)\cdots(n-j+1)$, the explicit Laguerre
formula gives
$Q_n=\sum_{j\ge0}(n)_{\underline j}^2/j!$, with all $j>n$ terms zero.
For $n\equiv r\pmod M$, each difference of summands equals
\[
 \bigl((n)_{\underline j}-(r)_{\underline j}\bigr)
 \left(\binom nj+\binom rj\right)
\]
and is divisible by $M$.  Taking $r=0,M=n$ gives $Q_n\equiv1\pmod n$.
Modulo $M$, the recurrence at index $M$ gives $P_{M+1}\equiv2P_M$;
uniqueness for the forward recurrence and periodicity of its
coefficients give $P_{M+r}\equiv P_M Q_r\pmod M$ for $r\ge0$.

For a prime power $p^a$, put $j=p^{\lceil a/2\rceil}$.  If $j\le n$,
then $j^2\equiv0\pmod{p^a}$ and $Q_j\equiv1\pmod p$ is invertible
modulo $p^a$.  At indices $j,j+1$, the vectors $(P_j,P_{j+1})$ and
$(Q_j,Q_{j+1})$ are proportional modulo $p^a$; the recurrence preserves
this proportionality at every later index.  Hence
$p^{\min(a,v_p(Q_n))}\mid g_n$ whenever $j\le n$.  Taking
$a=2\lfloor\log_p n\rfloor$ for each prime $p\le n$ proves the lower
divisibility bound.  Finally, the recurrence gives
$Q_nP_{n-1}-P_nQ_{n-1}=-((n-1)!)^2$, yielding the upper bound.
\end{proof}

The gap between these bounds is substantial.  They neither establish
sufficient gcd cancellation to obtain an irrationality proof nor bound
all cancellation strongly enough to exclude it.  In particular, the
factorial normalization in the positive-polynomial obstruction must not
be replaced without proof by a reduced common denominator.


% Source fragment: round2_number_theory.tex
\section{Sharper descent for periodic Dirichlet series}
\label{p2:r2nt-section}

The argument below strengthens the explicit large-prime bound in
Theorem~2 of the 2022 version of Bharadwaj's paper
\cite{p2:Bharadwaj2022}.  Section~8 of that version expressly suggests
improving this bound.  Our comparison is with that identified version;
priority relative to subsequent work is not asserted.  The full
Erd\H{o}s nonvanishing conjecture remains unresolved here.

For a rational-valued periodic function $h$ of period $K$ and mean
zero, write
\[
 L(1,h)=\sum_{n\ge1}\frac{h(n)}n,
 \qquad
 \mathcal L_K=\operatorname{span}_{\mathbb Q}
       \{\log r:r\text{ prime},\ r\mid K\}.
\]
Following \cite{p2:Bharadwaj2022}, say that $h$ has Property~U if
$L(1,h)\notin\mathcal L_K$.  An Erd\H{o}s function of period $K$
has values $\pm1$ outside $K\mathbb Z$, is zero on
$K\mathbb Z$, and has mean zero.  Property~U implies the
nonvanishing sought in the conjecture, but is a stronger condition.

\subsection{The conductor criterion and its integral form}

All moduli in the descent argument are odd and squarefree.  If
$m\mid t$, let $P_m^t$ denote averaging over the fibers of
\[
 (\mathbb Z/t\mathbb Z)^*\longrightarrow
 (\mathbb Z/m\mathbb Z)^*.
\]
Thus $P_m^t$ has operator norm at most one for the supremum norm.
For $r$ coprime to $m$, let $S_rv(x)=v(rx)$ on units modulo $m$.
For each prime $r\mid m$, choose a generator $\tau_r$ of its
prime-coordinate group in
$(\mathbb Z/m\mathbb Z)^*\cong\prod_{r\mid m}(\mathbb Z/r\mathbb Z)^*$,
and let it act by translation.  Put
\[
 \Delta_m=\prod_{r\mid m}(I-\tau_r).
\]
In the character basis, $\Delta_m$ is zero precisely on the
characters imprimitive modulo $m$.  Indeed, a character is primitive
for a squarefree modulus exactly when every prime-coordinate
character is nonprincipal, and a generator detects each such
nonprincipal character.

For $d\mid K$, define the function $h_d$ on units modulo $K/d$ by
$h_d(x)=h(dx)$.  We use the following established conductor criterion
from the proof of Theorem~4 in \cite{p2:Bharadwaj2022}:
\begin{equation}
\label{p2:r2nt-conductor-criterion}
\begin{aligned}
 L(1,h)\in\mathcal L_K
 \quad\Longleftrightarrow\quad &
 \Delta_{K/D}\sum_{d\mid D}\frac1d
 \prod_{r\mid D/d}\left(I-\frac{S_r}{r}\right)
 P_{K/D}^{K/d}h_d=0
 \\
 &\text{for every proper }D\mid K.
\end{aligned}
\end{equation}
Here and below products indexed by $r$ are over primes.  To explain
the input's scope, the character expansion of
$L(s,h)=\sum_{d\mid K}d^{-s}L(s,h_d)$ groups its nonprincipal
terms by primitive conductor.  The averaging operators select
characters whose conductors divide $K/D$, and the displayed Euler
factors remove the additional prime factors of $K/d$.  The
nonprincipal conductor contributions vanish exactly under the
conditions on the right of \eqref{p2:r2nt-conductor-criterion}; their
separation is the cited theorem.  The remaining principal contribution
is a rational combination of prime logarithms, since the mean of
$h$ is zero.  This is the Property~U formulation of that criterion;
no prime-logarithm vanishing equations are imposed in it.  More
formally, set $u_r(n)=1-r\mathbf1_{r\mid n}$ for each $r\mid K$.
This is a mean-zero $K$-periodic function with
$L(1,u_r)=\log r$, as follows by taking the limit at $s=1$ in
$(1-r^{1-s})\zeta(s)$.  Every $(u_r)_d$ is constant on the relevant
unit group, so all its nonprincipal conductor expressions vanish.
If $L(1,h)=\sum_{r\mid K}c_r\log r$, applying the cited vanishing
criterion to $h-\sum_rc_ru_r$ proves necessity in
\eqref{p2:r2nt-conductor-criterion}.  For sufficiency, grouping by
primitive conductor leaves only the principal contribution, which
belongs to $\mathcal L_K$.

For an integer-valued $f$ of period $N$, define, for each proper
$D\mid N$ and $m=N/D$,
\begin{equation}
\label{p2:r2nt-gd-definition}
 G_D(f)=D\varphi(D)\Delta_m
 \sum_{d\mid D}\frac1d
 \prod_{r\mid D/d}\left(I-\frac{S_r}{r}\right)
 P_m^{N/d}f_d.
\end{equation}

\begin{lemma}[Integrality and the exact product bound]
\label{p2:r2nt-gd-norm}
If $N$ is odd and squarefree and $\|f\|_\infty\le B$, then
$G_D(f)$ is integer-valued and
\begin{equation}
\label{p2:r2nt-gd-bound}
 \|G_D(f)\|_\infty
 \le B\,2^{\omega(N/D)}\varphi(D)
            \prod_{r\mid D}(r+2).
\end{equation}
If $r_0$ is the least prime divisor of $N>1$, the maximum of the
factor multiplying $B$ over proper divisors $D$ is
\begin{equation}
\label{p2:r2nt-cn}
 C(N)=2\varphi(N/r_0)\prod_{r\mid N/r_0}(r+2)
     =2\prod_{r\mid N/r_0}(r-1)(r+2).
\end{equation}
\end{lemma}

\begin{proof}
The denominator introduced by $P_m^{N/d}$ divides
$\varphi(D/d)$.  The factor $d^{-1}$ and the Euler operators have
combined denominator dividing $D$, since $D$ is squarefree.
Thus multiplication by $D\varphi(D)$ makes each term integral.
Translation differences preserve integrality.

The norm of $\Delta_m$ is at most $2^{\omega(m)}$, and
$\|I-S_r/r\|\le1+1/r$.  Consequently,
\begin{align*}
 \|G_D(f)\|_\infty
 &\le B D\varphi(D)2^{\omega(m)}
   \sum_{d\mid D}\frac1d\prod_{r\mid D/d}\left(1+\frac1r\right)\\
 &= B\varphi(D)2^{\omega(m)}\prod_{r\mid D}(r+2).
\end{align*}
The equality is the squarefree factorization
$D\sum_{d\mid D}d^{-1}\prod_{r\mid D/d}(1+1/r)=\prod_{r\mid D}(r+2)$.
In the factor on the right of \eqref{p2:r2nt-gd-bound}, a prime
$r\mid D$ contributes $(r-1)(r+2)$, whereas a prime
$r\mid N/D$ contributes $2$.  The ratio
$(r-1)(r+2)/2$ exceeds one and increases with $r\ge3$.
Since $D$ must be proper, the maximum is obtained by omitting just
the least prime $r_0$ from $D$.  This proves \eqref{p2:r2nt-cn}.
\end{proof}

\begin{theorem}[Adaptive and uniform large-prime descent]
\label{p2:r2nt-descent}
Let $N>1$ be odd and squarefree.  Suppose $f$ is an integer-valued
mean-zero function of period $N$ with Property~U.  Define
\[
 C_f=\min_{\substack{D\mid N,\ D<N\\G_D(f)\ne0}}
                  \|G_D(f)\|_\infty.
\]
The minimum is over a nonempty finite set.  If $p\nmid N$ is prime
and $p>1+C_f$, then every integer-valued mean-zero function $g$ of
period $pN$ satisfying $g(pn)=f(n)$ has Property~U.
In particular, if $\|f\|_\infty\le B$, the uniform condition
\begin{equation}
\label{p2:r2nt-uniform-prime-bound}
 p>1+B C(N)
\end{equation}
suffices.  No size bound on $g$ is required.
\end{theorem}

\begin{proof}
The nonemptiness follows from
\eqref{p2:r2nt-conductor-criterion}.  Suppose, towards a contradiction,
that $g$ fails Property~U.  Fix a proper $D\mid N$, put $m=N/D$,
and define a function on the units modulo $pm$ by
\[
 h_D=D\varphi(D)\sum_{d\mid D}\frac1d
       \prod_{r\mid D/d}\left(I-\frac{S_r}{r}\right)
       P_{pm}^{pN/d}g_d.
\]
The same denominator argument as in
Lemma~\ref{p2:r2nt-gd-norm} shows that $h_D$ is integer-valued;
the averaging fibers have size $\varphi(D/d)$, independent of $p$.
Apply the conductor criterion to $g$ with divisor $D$.  It gives
\[
 \Delta_m(I-\tau_p)h_D=0.
\]
Hence $\Delta_m h_D$ is invariant under a generator of the
$p$-coordinate group and is constant along that coordinate.
It follows that
\[
 F_D=\Delta_m P_m^{pm}h_D
\]
is integer-valued: its average is simply the value at any chosen
$p$-coordinate.

Next apply the conductor criterion to $g$ with divisor $pD$.
Split its divisor sum into $d\mid D$ and $pd$ with $d\mid D$.
In the second part, $g_{pd}=f_d$.  In the first part, the additional
Euler factor is $I-S_p/p$.  Averaging is transitive and commutes
with the translation operators.  After multiplication by
$pD\varphi(D)$, the resulting identity is therefore
\begin{equation}
\label{p2:r2nt-integral-descent-equation}
 (pI-S_p)F_D+G_D(f)=0.
\end{equation}
If $G_D(f)\ne0$, then $F_D\ne0$, so its integer values give
$\|F_D\|_\infty\ge1$.  Since $S_p$ permutes the coordinates,
the reverse triangle inequality gives
\[
 \|G_D(f)\|_\infty
 =\|(pI-S_p)F_D\|_\infty
 \ge(p-1)\|F_D\|_\infty
 \ge p-1.
\]
Choose $D$ attaining $C_f$.  This contradicts $p>1+C_f$.
Finally, Lemma~\ref{p2:r2nt-gd-norm} gives $C_f\le B C(N)$,
proving the uniform statement.
\end{proof}

\begin{corollary}[Congruence and sign refinement]
\label{p2:r2nt-sign-refinement}
In the setting of Theorem~\ref{p2:r2nt-descent}, suppose additionally
that $g(n)\equiv a\pmod M$ whenever $pN\nmid n$, where $M\ge1$
is an integer.  Then $p>1+C_f/M$ suffices.  In particular, for
Erd\H{o}s functions $f$ and $g$, the sufficient uniform condition is
\begin{equation}
\label{p2:r2nt-sign-bound}
 p>1+\frac{C(N)}2
   =1+\prod_{r\mid N/r_0}(r-1)(r+2).
\end{equation}
\end{corollary}
\begin{proof}
Every $g_d$ occurring in $h_D$ takes values congruent to $a$
modulo $M$, because $d\mid D<N$.  On expanding the averaging and
Euler operators in the definition of $h_D$, their coefficients are
integers after multiplication by $D\varphi(D)$, as established above.
Their row sums are independent of the unit-group coordinate: a
translation has row sum one, and a fiber sum has its fixed fiber
size as row sum.  Therefore $h_D$ is constant modulo $M$.
Since $m=N/D>1$, the operator $\Delta_m$ contains a translation
difference, so $\Delta_mh_D$ takes values in $M\mathbb Z$.
Its constancy along the $p$-coordinate implies that $F_D$ also
takes values in $M\mathbb Z$.  Thus, if $G_D(f)\ne0$, equation
\eqref{p2:r2nt-integral-descent-equation} gives
$\|G_D(f)\|_\infty\ge M(p-1)$.
Taking a divisor attaining $C_f$ proves the first assertion.
For sign functions, use $M=2$, $a=1$, and $B=1$.
\end{proof}

\begin{corollary}[A primitive-component shortcut]
\label{p2:r2nt-primitive-shortcut}
Under the hypotheses of Theorem~\ref{p2:r2nt-descent}, suppose that
$\Delta_N f_1\ne0$.  Then $p>1+B2^{\omega(N)}$ is sufficient.
\end{corollary}
\begin{proof}
Take $D=1$ in \eqref{p2:r2nt-gd-definition}.  Then
$G_1(f)=\Delta_N f_1\ne0$ and its supremum norm is at most
$B2^{\omega(N)}$.
\end{proof}

For orientation, the sign-function thresholds in
\eqref{p2:r2nt-sign-bound} are compared below with
$2^{\omega(N)}\prod_{r\mid N}(r^2+2r+2)$, the threshold in
Theorem~2 of the cited 2022 version.  These are sufficient bounds,
not claimed optimal bounds.
\begin{center}
\begin{tabular}{r|r|r}
$N$ & sign-function condition & 2022 condition\\\hline
$3$ & $p>2$ & $p>34$\\
$15$ & $p>29$ & $p>2516$\\
$105$ & $p>1513$ & $p>327080$\\
$195$ & $p>5041$ & $p>991304$\\
$1155$ & $p>196561$ & $p>94853200$
\end{tabular}
\end{center}

\subsection{Prime marginals and the precise remaining obstruction}

\begin{lemma}[Prime marginal obstruction]
\label{p2:r2nt-prime-marginals}
Let $K$ be squarefree and let $h$ be a rational-valued mean-zero
function of period $K$.  For $r\mid K$, define
\[
 S_r(h)=\sum_{1\le a\le K,\ r\mid a}h(a).
\]
In the decomposition of $L(1,h)$ into rational-prime logarithms
and an algebraic combination of unit logarithms and $\pi$, the
coefficient of $\log r$ is
\[
 -\frac{r S_r(h)}{K(r-1)}.
\]
If any $S_r(h)\ne0$, then $L(1,h)$ is transcendental and nonzero.
\end{lemma}
\begin{proof}
Put $M_d=\sum_{(a,K/d)=1}h(da)$ for $d\mid K$.
The principal-character contributions to the constant term at $1$
give
\[
 \frac1K\sum_{d\mid K}M_d
 \left(\gamma-\log d+
             \sum_{r\mid K/d}\frac{\log r}{r-1}\right).
\]
The coefficient of $\gamma$ vanishes because
$\sum_dM_d=0$.  Squarefreeness gives
$S_r(h)=\sum_{r\mid d}M_d$, and therefore the coefficient of
$\log r$ is
\[
 \frac1K\left(\frac{\sum_{r\nmid d}M_d}{r-1}
                    -\sum_{r\mid d}M_d\right)
 =-\frac{rS_r(h)}{K(r-1)}.
\]
The other character terms are algebraic combinations of logarithms
of algebraic units and $\pi$.  Prime valuations and Baker's theorem
separate the nonzero rational-prime component from these terms,
and give the stated transcendence conclusion.
\end{proof}

For example, if every Erd\H{o}s function of period $N$ has
Property~U, the uniform bound proves nonvanishing for every
Erd\H{o}s function $g$ of period $pN$.  Indeed, put $f(n)=g(pn)$.
If $f$ has nonzero mean, the $p$-marginal in
Lemma~\ref{p2:r2nt-prime-marginals} already excludes $L(1,g)=0$.
If its mean is zero, apply Theorem~\ref{p2:r2nt-descent}.
The known $N\equiv3\pmod4$ argument of Murty and Saradha supplies
the required Property~U through a nonzero $\pi$ component;
see the proof of Corollary~1 in \cite{p2:Bharadwaj2022}.
Thus this specialization yields periods $pN\equiv1\pmod4$ when
also $p\equiv3\pmod4$.

\subsection{An exact answer to a stated existence question}
\label{p2:r2nt-existence-subsection}

Section~8 of the 2022 version of \cite{p2:Bharadwaj2022} asks whether
there exists an odd natural number $N$ and signs $a_d\in\{\pm1\}$ for
$d\mid N$, $d>1$, such that
\begin{equation}
\label{p2:r2nt-external-existence-question}
 a_1=0,\qquad \sum_{d\mid N}a_d\varphi(d)=0.
\end{equation}
For $N=1$ the displayed condition is vacuous; we answer its
nontrivial interpretation $N>1$.  The following supplies an explicit
example and the least nontrivial odd period.
The later journal version was identified but its full
text was unavailable for comparison, so current priority is not
asserted.

\begin{proposition}[An odd example and its minimality]
\label{p2:r2nt-smallest-example}
The least odd integer $N>1$ satisfying
\eqref{p2:r2nt-external-existence-question} is $105$.
\end{proposition}
\begin{proof}
For $N=105$, take $a_5=a_{105}=1$ and all other $a_d=-1$
for $d>1$.  Then
\begin{align*}
 &-\varphi(3)+\varphi(5)-\varphi(7)-\varphi(15)
 -\varphi(21)-\varphi(35)+\varphi(105)\\
 &\hspace{2em}=-2+4-6-8-12-24+48=0.
\end{align*}
If an odd integer $N>1$ has at most two distinct prime factors,
then
\[
 \frac{\varphi(N)}N\ge
 \left(1-\frac13\right)\left(1-\frac15\right)
 =\frac8{15}>\frac12.
\]
The identity $\sum_{d\mid N}\varphi(d)=N$ shows that the
$d=N$ term has absolute value greater than the sum of the absolute
values of all other terms with $d>1$:
$\varphi(N)>N-1-\varphi(N)$.  Cancellation is impossible.
Every odd integer below $105=3\cdot5\cdot7$ has at most two
distinct prime factors.  This proves minimality.
\end{proof}

This example also shows why Property~U cannot simply be replaced
by nonvanishing in the descent proof.  At period $105$, define a
function constant on gcd
classes by the following table:
\begin{center}
\begin{tabular}{c|rrrrrrrr}
$d=(n,105)$ & $1$ & $3$ & $5$ & $7$ & $15$ & $21$ & $35$ & $105$\\\hline
$f(n)$ & $1$ & $-1$ & $-1$ & $-1$ & $-1$ & $1$ & $-1$ & $0$
\end{tabular}
\end{center}
It has mean zero, since
$48-24-12-8-6+4-2=0$.  Its Dirichlet series factors as
$L(s,f)=F(s)\zeta(s)$, with
\begin{align*}
 F(s)={}&1-2\,3^{-s}-2\,5^{-s}-2\,7^{-s}
       +2\,15^{-s}+4\,21^{-s}+2\,35^{-s}-3\,105^{-s}.
\end{align*}
This follows by checking the eight gcd classes, or by M\"obius
inversion on the divisors of $105$.  Direct evaluation gives
$F(1)=0$ and
\begin{equation}
\label{p2:r2nt-u-counterexample}
 L(1,f)=F'(1)=\frac{13}{35}\log3+
        \frac5{21}\log5+\frac1{15}\log7>0.
\end{equation}
Thus it fails Property~U despite having a positive transcendental
$L$-value.  This explicit example identifies a genuine limitation
of iterating Property~U descent.

\begin{proposition}[Elementary exclusion of a double zero]
\label{p2:r2nt-gcd-invariant}
Let $N>1$ be odd and not a square.  Every Erd\H{o}s function of
period $N$ that is constant on gcd classes has nonzero $L(1,f)$.
Equivalently, its Dirichlet-polynomial factor $F(s)$ has a simple
zero at $s=1$.  Nonvanishing here has an elementary proof that
does not use Baker's theorem.
\end{proposition}
\begin{proof}
Write $f(n)=a_t$ when $t=N/(n,N)$, so $a_1=0$ and
$a_t\in\{\pm1\}$ for $t>1$.  Partitioning by gcd gives
\[
 L(s,f)=F(s)\zeta(s),\qquad
 F(s)=\sum_{t\mid N}a_t(N/t)^{-s}
                         \prod_{r\mid t}(1-r^{-s}).
\]
The mean-zero condition gives $F(1)=N^{-1}\sum_ta_t\varphi(t)=0$.
Differentiating and using this identity yields
\begin{equation}
\label{p2:r2nt-log-coefficients}
 F'(1)=\frac1N\sum_{r\mid N}B_r\log r,\qquad
 B_r=\sum_{t\mid N}a_t\varphi(t)
       \left(v_r(t)+\frac{\mathbf1_{r\mid t}}{r-1}\right).
\end{equation}
Each $B_r$ is an integer.  Choose a prime $r$ whose exponent
$v=v_r(N)$ is odd.  Terms with $r\nmid t$ make no contribution.
For $t=r^je$, with $j\ge1$ and $(e,r)=1$, the unsigned weight is
\[
 \varphi(r^je)\left(j+\frac1{r-1}\right)
 =r^{j-1}\bigl((r-1)j+1\bigr)\varphi(e).
\]
It is even when $e>1$, and odd when $e=1$.  Hence
$B_r\equiv v\equiv1\pmod2$, independently of the signs.
Thus at least one coefficient in \eqref{p2:r2nt-log-coefficients}
is nonzero.  The logarithms of distinct rational primes are
linearly independent over $\mathbb Q$: exponentiating an integer
relation would contradict unique prime factorization.  Therefore
$F'(1)\ne0$, and $L(1,f)=F'(1)\ne0$.
\end{proof}

The source also asks for an approach to excluding a double zero
of this Dirichlet-polynomial factor without Baker's theorem.  The
last proposition supplies such an approach for every odd nonsquare
period, including every odd squarefree period; odd square periods
are outside its conclusion.  We distinguish the exact existence
answer from this partial response to a broader method request,
and from the quantitative descent improvement.  Neither the
full Erd\H{o}s conjecture nor current novelty is claimed.  Nor does the argument
prove the irrationality or transcendence of $\gamma$: throughout,
the mean-zero condition has cancelled its coefficient before the
arithmetic separation is applied.


% Source fragment: round2_closing.tex
\section{Remaining arithmetic barriers and excluded routes}
\label{p2:sec:round2-remaining}

The following central targets remain unproved by this manuscript:
$\gamma\notin\mathbb Q$; transcendence of $\gamma$; irrationality of
$\delta$; the full description of the units of $\mathcal E$;
the conjectured intersection of $E$-values with arithmetic Gevrey values;
and the full Erd\H{o}s nonvanishing conjecture. A finite family of
$E$-values can be algebraically independent even though its singular
limiting constants cannot currently be separated arithmetically.
The condition on the intersection
$\overline{\mathbb Q}(\eta)\cap\mathcal E$ in
Theorem~\ref{p2:thm:ein-denominators} remains explicitly conditional on
the transcendence of $\gamma$.

Murabayashi--Yoshida's Conjecture~7.5 \cite{p2:MY2024} was examined again.
It prescribes particular directions for monotone subsequences of
fractional-part combinations arising from their even and odd
Euler--Gompertz approximants. The general monotone-subsequence theorem
does not supply those prescribed directions. Each proposed direction
would already have strong arithmetic consequences for $\delta$.
Likewise, divergence after factorial denominator clearing does not
rule out substantial cancellation by common divisors or by
sign-changing filters. No such cancellation theorem was established.

Several seemingly promising historical questions were excluded after
checking later sources. The incomplete-gamma question in a 2016 paper
is answered by the same authors' 2017 sharp Gautschi result
\cite{p2:Yang2017Gautschi}. A rational gamma-bound parameter problem from
2017 has a 2020 solution \cite{p2:Gamma2020}. Higher-order polygamma
claims in a public 2026 archive \cite{p2:pudim2026} prevent assigning
unqualified novelty to our independent proof. None of these results
is silently counted as a new solution.

The manuscript therefore separates three issues: validity of each
displayed argument, exact agreement with the source statement, and
historical novelty or mathematical importance. The first two are
addressed explicitly. The third requires further specialist comparison
and review. In particular, resolving ten source assertions of mixed
difficulty does not establish the ten major new breakthroughs requested
as the original research goal.

\PartAppendix
\section{Exact checks used in the proofs}
All principal theorems are analytic or algebraic proofs. Numerical
values of tangency points are illustrations, not proof inputs.
The following short Python 3 certificate verifies the rational
arithmetic in the contour bound, the sharp-exponent proof, and the
totient example. It uses only the standard library.
\begin{lstlisting}
from fractions import Fraction as F
from math import comb, factorial

# Touchard polynomials and Bernstein coefficients, degrees 1,...,5.
S = [[1]]
for n in range(1, 6):
    row = [0]*(n+1)
    for j in range(1, n+1):
        row[j] = S[n-1][j-1]
        if j < n:
            row[j] += j*S[n-1][j]
    S.append(row)
maxima, bells = [], []
for n in range(1, 6):
    coeff = [(-1)**j*S[n][j] for j in range(n+1)]
    bern = [sum(F(comb(k,j),comb(n,j))*coeff[j]
                for j in range(k+1)) for k in range(n+1)]
    maxima.append(max(map(abs, bern)))
    bells.append(sum(S[n]))
assert maxima == [F(1), F(1,2), F(1), F(5,4), F(2)]
R = F(7,10)
bound = F(14,5)-1 + sum((a-b)*R**n/factorial(n)
           for n,(a,b) in enumerate(zip(maxima,bells),1))
assert bound == F(318799,320000) < 1

# Rigorous exponential-series upper bound.
def exp_upper(x, N):
    partial = sum(x**k/factorial(k) for k in range(N+1))
    tail = x**(N+1)/factorial(N+1)/(1-x/F(N+2))
    return partial + tail
assert exp_upper(F(7,10),20) < F(101,50)
assert exp_upper(F(51,50),20) < F(14,5)

x = F(2,3)
lower = sum(x*(2*k+x)/(k*(k+x)**2) for k in range(1,7))
lower += F(4,23)-F(29,50)
assert lower == F(889267245,1261122016) > F(7,10)
assert -2+4-6-8-12-24+48 == 0
print("All exact rational checks passed.")
\end{lstlisting}



% Source fragment: round2_bibliography.tex
\begin{thebibliography}{VIII.99}
\bibitem[II.1]{p2:FR2026}
S.~Fischler and T.~Rivoal, \emph{The ring of values of $E$-functions
is not a field}, arXiv:2610.16603v1 (2026), Proposition~3.1,
Corollary~3.3, and Conjecture~1.2.
\url{https://arxiv.org/abs/2610.16603}.

\bibitem[II.2]{p2:RR2015}
T.~Rivoal and J.~Roques, \emph{$E$-functions of order 2 and units of
$E$-values}, preprint dated 6 November 2015, Theorem~5.
\url{https://rivoal.perso.math.cnrs.fr/articles/unitse.pdf}.

\bibitem[II.3]{p2:DLMF}
NIST Digital Library of Mathematical Functions, Chapter~13,
equations 13.2.39, 13.2.41, 13.3.1, 13.4.4, and 13.7.3.
\url{https://dlmf.nist.gov/13}.

\bibitem[II.4]{p2:BMN2022}
I.~Boussaada, G.~Mazanti, and S.-I.~Niculescu,
\emph{Some remarks on the location of non-asymptotic zeros of
Whittaker and Kummer hypergeometric functions}, Bulletin des Sciences
Math\'ematiques \textbf{174} (2022), 103093.
\url{https://arxiv.org/abs/2106.03378}.
\url{https://doi.org/10.1016/j.bulsci.2021.103093}.

\bibitem[II.5]{p2:Vargas2014}
A.~R.~Vargas, \emph{Limit curves for zeros of sections of exponential
integrals}, Constructive Approximation \textbf{40} (2014), 219--239.
\url{https://arxiv.org/abs/1302.3695}, version 4, Conjecture~5.1.
\url{https://doi.org/10.1007/s00365-014-9241-7}.

\bibitem[II.6]{p2:Gorgens2014}
M.~G\"orgens, \emph{Optimal stopping of an $\alpha$-Brownian bridge},
arXiv:1409.5285v1 (18 September 2014), Introduction and Section~4.
\url{https://arxiv.org/abs/1409.5285}.

\bibitem[II.7]{p2:qln2025}
F.~Qi, D.~Lim, and K.~Nantomah, \emph{Monotonicity and positivity of
several functions involving ratios and products of polygamma functions},
Journal of Inequalities and Applications (2025), article 5.
\url{https://doi.org/10.1186/s13660-024-03245-8}.

\bibitem[II.8]{p2:shuang2026}
Y.~Shuang, D.~Lim, and F.~Qi, \emph{Alternative proofs for the
positivity of a function involving the product of two digamma functions},
Australian Journal of Mathematical Analysis and Applications
\textbf{23}(1) (2026), article 5.
\url{https://ajmaa.org/searchroot/files/pdf/v23n1/v23i1p5.pdf}.

\bibitem[II.9]{p2:Wis2025}
H.~Wi\'sniewska, \emph{Some inequalities for Gurland's ratio of the
gamma functions}, arXiv:2512.07028v1 (7 December 2025),
Section~2, Open Problems~1--3.
\url{https://arxiv.org/abs/2512.07028}.

\bibitem[II.10]{p2:Bolbachan2011}
V.~Bolbachan, \emph{Rational approximants for the Euler--Gompertz
constant}, arXiv:1104.4721v2 (10 November 2011), Conjecture~1.4.
The preprint identifies itself as prepared for the Moscow Mathematical
Conference for High-School Students.
\url{https://arxiv.org/abs/1104.4721v2}.

\bibitem[II.11]{p2:Bharadwaj2022}
A.~Bharadwaj, \emph{On primitivity and vanishing of Dirichlet series},
arXiv:2204.03674v1 (2022), Theorems~2 and~4 and Section~8.
\url{https://arxiv.org/abs/2204.03674}.
Later journal publication: Michigan Mathematical Journal
\textbf{75}(4) (2025), 725--759,
\url{https://doi.org/10.1307/mmj/20226269}.
The revised full text was not inspected; comparisons in this work
are to the specified arXiv version.

\bibitem[II.12]{p2:HobsonLiu2026}
D.~Hobson and J.~Liu, \emph{Optimal stopping of a Brownian excursion
and an $\alpha$-dimensional Bessel bridge}, Journal of Applied
Probability, published online 15 June 2026, 1--18.
\url{https://doi.org/10.1017/jpr.2026.10102}.

\bibitem[II.13]{p2:Rivoal2012}
T.~Rivoal, \emph{On the arithmetic nature of the values of the gamma
function, Euler's constant, and Gompertz's constant},
Michigan Mathematical Journal \textbf{61} (2012), 239--254, Section~7.
\url{https://rivoal.perso.math.cnrs.fr/articles/gammater.pdf}.

\bibitem[II.14]{p2:Powers2026}
M.~R.~Powers, \emph{Transcendence results for $\Gamma^{(n)}(1)$ and
related sequences of generalized constants}, arXiv:2511.01849v7
(2026), Theorem~2.
\url{https://arxiv.org/abs/2511.01849v7}.

\bibitem[II.15]{p2:MY2024}
N.~Murabayashi and H.~Yoshida, \emph{Rational function approximations
of the special function $e^xE_1(x)$ and applications to irrationality
of Euler--Gompertz constant $\delta$}, arXiv:2210.06768v2,
revised 22 September 2024, Conjecture~7.5.
\url{https://arxiv.org/abs/2210.06768}.

\bibitem[II.16]{p2:Yang2017Gautschi}
Z.-H.~Yang, W.~Zhang, and Y.-M.~Chu,
\emph{Sharp Gautschi inequality for parameter $0<p<1$ with applications},
Mathematical Inequalities \& Applications \textbf{20}(4) (2017),
1107--1120.
\url{https://doi.org/10.7153/mia-2017-20-71}.
\url{https://files.ele-math.com/articles/mia-20-71.pdf}.

\bibitem[II.17]{p2:Gamma2020}
J.-M.~Shen, Z.-H.~Yang, W.-M.~Qian, W.~Zhang, and Y.-M.~Chu,
\emph{Sharp rational bounds for the gamma function},
Mathematical Inequalities \& Applications \textbf{23}(3) (2020),
843--853.
\url{https://doi.org/10.7153/mia-2020-23-68}.
\url{https://files.ele-math.com/articles/mia-23-68.pdf}.

\bibitem[II.18]{p2:pudim2026}
Pudim Project, public research archive (2026), node
\emph{OP4 solved except $n=2$ exact lower endpoint unresolved},
and APP-0011; accessed 18 October 2026. These are public claims,
not evidence of external peer review.
\url{https://github.com/pudim-project/pudim-ai-demo-zetalaw}.
\end{thebibliography}
% END PRESERVED SOURCE BODY 2
\phantomsection\label{part2:end}
\clearpage
\setcounter{section}{0}\setcounter{subsection}{0}
\renewcommand{\thesection}{\Roman{part}.\arabic{section}}
\renewcommand{\theHsection}{\arabic{part}.main.\arabic{section}}
\renewcommand{\refname}{References for Part III}
\fancyhead[L]{\small Euler\textquotesingle s constant: Part III}
\part{Scalar convergence, moment problems, and differential algebra}
\label{part3:start}
\noindent\textbf{Source:} {\small\path{Keio_Euler_Constant_Structures_20261018.tex}}.

Original PDF: 37 pages. Mathematical content and source-level qualifications are retained.


% BEGIN PRESERVED SOURCE BODY 3
\begin{abstract}
We distinguish two asymptotic assertions that are combined in a recent
conjecture on conservative matrix fields for Euler's constant. For every
rank-three trajectory $(-A_1,-A_2;\beta,\beta,\beta+d)$ with
$A_i\ge0$ and $\beta,d\ge1$ integral, we prove that almost every fixed
real seed has dual projective limit $(1,-1,0)^T$, contradicting the proposed
full-state isolation. Nevertheless, for every nonzero fixed complex seed,
the prescribed scalar ratio converges to $\gamma$ and its denominator is
eventually nonzero. The proof combines exact contiguity determinants,
factorial separation, and concentration of positive finite sums.

We also prove the all-order coefficient pairing explicitly asked for by
Xu and You, and evaluate the bordered factorial Hankel determinants
conjectured by Ferguson; these two proofs are applications of classical
Stieltjes and Laguerre theory. Further results determine exact differential
and difference orders, a solvable differential Galois group, denominator
obstructions, and rational intersections in the ring of $E$-function values.
An application of a theorem of Bostan--Di Vizio--Raschel establishes
differential transcendence in the parameter of fixed-cutoff incomplete
gamma functions. We separate source-level resolutions, restricted-family
results, and classical consequences. No irrationality or transcendence
of Euler's constant or the Euler--Gompertz constant is established.
\end{abstract}
\vspace{4pt}
\noindent\textbf{Status.} Research manuscript prepared with AI assistance.
Proofs received internal independent checks; external peer review and a
comprehensive priority search have not been completed.


\section{Results, source questions, and scope}
\label{p3:sec:introduction}

Euler's constant and the Euler--Gompertz constant are
\[
 \gamma=\lim_{n\to\infty}(H_n-\log n),\qquad
 \delta=\int_0^\infty\frac{e^{-t}}{1+t}\,dt,
 \qquad H_n=\sum_{j=1}^n\frac1j.
\]
Their approximation theory touches several distinct structures: polynomial
recurrences, differential equations, positive moments, and rings of special
values. The distinction between these structures matters. A function can
have a completely factored differential operator and still have maximal
minimal differential order. A dual vector can converge in a direction
annihilated by both extraction forms while their ratio has a meaningful
limit. Positive approximation can coexist with insufficient denominator
control for an irrationality argument.

\subsection{The principal recurrence result}

Kalisch et al.~\cite[Definitions 20--21, Conjecture 1]{p3:KMLWSK26}
prove full-state isolation and scalar convergence for a rank-three family
with lower-parameter increments $(1,1,1)$. They conjecture both statements
for all ranks and nonpositive upper increments, with positive lower
increments whose minimum is repeated. Our main result treats every
unequal rank-three lower tuple with a repeated minimum, up to permutation.
For
\[
 t=(-A_1,-A_2;\beta,\beta,\beta+d),\quad
 A_1,A_2\in\mathbb Z_{\ge0},\quad \beta,d\in\mathbb Z_{\ge1},
\]
let $\Phi_n$ be the source transition from the origin to $nt$, and put
$p=(1,1,0)^T$, $q=(1,1,1)^T$. We prove the following two conclusions.
\begin{enumerate}
\item For almost every fixed real $v$, the projective limit of
      $\Phi_n^{-T}v$ is $u=(1,-1,0)^T$, which is not the canonical
      Meijer state. This refutes the full-state part of Conjecture~1
      (Theorem~\ref{p3:thm:full-state-counterexample}).
\item For every fixed $v\in\mathbb C^3\setminus\{0\}$,
      $q^T\Phi_n^{-T}v\ne0$ eventually and
      \[
       \frac{p^T\Phi_n^{-T}v}{q^T\Phi_n^{-T}v}\longrightarrow\gamma.
      \]
      Thus the scalar conclusion survives on the whole stated family
      (Theorem~\ref{p3:thm:unequal-rank-three-scalar}).
\end{enumerate}
The mechanism is $p^Tu=q^Tu=0$. Exact factorial estimates first separate
the dominant dual direction. The ratio is then recovered from the
subdominant coordinates using a positive weight distribution concentrated
at an index of order $\sqrt n$. This gives a proof in the relevant
quotient without relying on the false full-state extension.

\subsection{What is resolved in the cited sources}

The following accounting uses a source assertion as the unit. Parameter
instances, corollaries, and algorithmic consequences are not counted as
separate open problems.

\begin{longtable}{@{}>{\raggedright\arraybackslash}p{0.25\textwidth}>{\raggedright\arraybackslash}p{0.39\textwidth}>{\raggedright\arraybackslash}p{0.28\textwidth}@{}}
\toprule
Source assertion & Conclusion in this manuscript & Scope \\
\midrule
\endhead
Kalisch et al., Conjecture 1~\cite{p3:KMLWSK26}
& Full-state extension refuted; scalar convergence proved on all unequal
rank-three repeated-minimum trajectories.
& One conjecture separated into a false assertion and a true restricted-family assertion.\\[6pt]
Xu--You, Open problem after Theorem 1~\cite{p3:XuYou2014}
& $a_{2j+1}=-a_{2j}$ for every $j\ge1$, with positive rational even
coefficients and a determinant algorithm.
& Exact answer by classical Stieltjes theory; first priority not asserted.\\[6pt]
Ferguson, final example of Section 2~\cite{p3:Ferguson2020}
& $P_N/Q_N=H_N$ and positive product formulas for both determinants.
& Exact answer by classical Laguerre theory; first priority not asserted.\\[6pt]
Kalisch et al., Section 9.2, item 1~\cite{p3:KMLWSK26}
& Explicit $\gamma+\log z$ identity and a shifted-parameter extension.
& The identity component only; the unshifted identity also follows directly
from their equation and evaluation at one.\\[6pt]
Rivoal, Problem 6~\cite{p3:RivoalXUPS2019}; Fischler--Rivoal, Conjecture 1.2~\cite{p3:FischlerRivoal2026Units}
& Polynomial field intersections and exclusion of nonconstant polynomial
units for specified exponential-integral values and combinations.
& Restricted subfields; the full unit-group conjecture is not proved.\\
\bottomrule
\end{longtable}

There are consequently two explicit source questions answered affirmatively
and one general conjecture refuted as stated. The remaining entries are
partial results or consequences, not additional full conjecture resolutions.
The first three conclusions do not certify that the corresponding
identities have no equivalent earlier proofs. In particular, the positive
digamma continued fraction predates the Xu--You question.

\subsection{Structural results and their logical limits}

Sections~\ref{p3:sec:meijer-extension}--\ref{p3:sec:diff-structure} connect
Mellin--Barnes residues with gamma and logistic random variables, and
determine a Picard--Vessiot group
$\mathbb G_a^p\rtimes(\mathbb G_m\times\mathbb G_a)$ of dimension $p+2$.
The minimal differential order of the selected Meijer solution is $p+1$.
Section~\ref{p3:sec:gamma-arithmetic} proves that Rivoal's fixed-parameter
third-order recurrence is irreducible exactly when its parameter is nonzero.
Its integral solution lattice vanishes; cumulative denominators of nonzero
rational solutions grow exponentially along a subsequence, and the original CMF vector normalization
has a factorial denominator barrier.

Sections~\ref{p3:sec:factorial-hankel}--\ref{p3:sec:gompertz-nonholonomic}
classify a missing zeroth moment and show that logarithmic Gompertz moments
do not satisfy polynomial-coefficient linear recurrences. A stronger
functional differential-transcendence result follows from the established
parabolic difference-equation theorem of Bostan--Di Vizio--Raschel.
Section~\ref{p3:ef:section} uses the phase morphism of Fischler--Rivoal to
prove complete rational-intersection statements for the stated families.

These conclusions address different mathematical categories. In
particular, functional differential transcendence is not numerical
transcendence; a common denominator of a vector is not the height of a
reduced scalar ratio; and minimal order of one recurrence does not exclude
a different lower-order construction. These distinctions are used in the
proofs and in the unresolved-target ledger at the end.

\paragraph{Conventions.}
$(x)_m=x(x+1)\cdots(x+m-1)$, with $(x)_0=1$;
$H_n^{(r)}=\sum_{j=1}^n j^{-r}$, with $H_0^{(r)}=0$;
$\theta=z\,d/dz$. The symbols $\Gamma(s,a)$ and $\gamma(s,a)$ denote
the upper and lower incomplete gamma functions; their two-argument notation
distinguishes the latter from Euler's constant. Differential order is over
$\mathbb C(z)$ unless specified; recurrence order is over $\mathbb C(n)$.



\section{A counterexample to full-state isolation in Conjecture~1}
\label{p3:sec:cmf-counterexample}

\paragraph{Scope.}
This section concerns Conjecture~1 of Kalisch et al.,
arXiv:2610.13919v1~\cite{p3:KMLWSK26}, specifically its assertion that Theorem~3 (full-state
isolation) extends to every trajectory in $\mathcal T_r$.
It does not refute the separate scalar-ratio conclusion of their Theorem~1.
The distinction is essential: the two linear forms used to extract $\gamma$
both annihilate the limiting direction exhibited below.

Write $\theta=z\,d/dz$. In rank three at the source point, the differential
operator is
\[
 L_0=\theta^3+z(\theta+1)^2=\theta^2(\theta+z).
\]
For a germ $f$ at $z=1$, put
\[
 J(f)=(f(1),\theta f(1),\theta^2f(1)).
\]
All matrix conventions below use row states, exactly as the source paper:
$J(f_n)=J(f_0)\Phi_n$.
Equivalently, the matrix $\Phi_n$ can be defined intrinsically by
\eqref{p3:eq:cmf-contiguity} on the three-dimensional solution space of $L_0$;
the polynomial $P_n$ is displayed below. This characterizes the transition
without needing a separate list of elementary CMF matrices.

\begin{theorem}[Failure of full-state isolation on unequal rank-three trajectories]
\label{p3:thm:full-state-counterexample}
Let $A_1,A_2\geq0$ and $\beta,d\geq1$ be integers, and consider the rank-three
CMF trajectory
\[
 t=(-A_1,-A_2;\beta,\beta,\beta+d).
\]
Let $\Phi_n=M_{0,t}(0)\cdots M_{0,t}(n-1)$ be its transition matrix.
Every matrix in this product is defined and invertible.
For Lebesgue-almost every fixed $v\in\mathbb R^3$ there are nonzero real
scalars $c_n(v)$ such that
\[
 c_n(v)\Phi_n^{-T}v\longrightarrow u:=(1,-1,0)^T.
\]
In particular, the conclusion of Theorem~3 in arXiv:2610.13919v1 cannot
hold for all generic seeds on this trajectory. Consequently Conjecture~1,
which explicitly includes that conclusion, is false as stated.
The smallest example is $t=(0,0;1,1,2)$.
\end{theorem}

\begin{proof}
Set $K=\beta+d$ and
\[
 P_n(X)=(X+1)_{A_1n}(X+1)_{A_2n}
              (-X)_{\beta n}^{2}(-X)_{Kn},
\]
where $(x)_m$ is the rising factorial. Upper-parameter decrements are
implemented by $\theta-a_i+1$, and lower-parameter increments by
$-\theta+b_j$. All these scalar operators are polynomials in $\theta$
and commute. It follows from the contiguous relations, without any
asymptotic argument, that
\begin{equation}
 J(P_n(\theta)f)=J(f)\Phi_n\qquad(L_0f=0).
 \label{p3:eq:cmf-contiguity}
\end{equation}
Indeed each contiguous matrix is obtained by reducing the operators
$\theta^jD$ modulo the old monic differential equation, where $D$ is
the corresponding first-order contiguous operator. Thus the identity
holds for every solution germ, not only the canonical Meijer $G$ branch.

We first record the exact determinant. At arbitrary parameters, the
rank-three scalar operator has polynomial in $\theta$
\[
 \prod_{j=1}^3(\theta-b_j)+z\prod_{i=1}^2(\theta-a_i+1).
\]
At $z=1$, an upper decrement has determinant
$\prod_j(b_j-a_i+1)$; a lower increment has determinant
$\prod_i(b_j-a_i+1)$. Along a path consisting first of upper decrements
and then lower increments all these factors are positive integers.
Multiplying them gives
\begin{equation}
 \det\Phi_n=\prod_{i=1}^2
       \bigl(((A_i+\beta)n)!\bigr)^2
       ((A_i+\beta+d)n)!.
 \label{p3:eq:cmf-determinant}
\end{equation}
Conservation gives the same transition for any admissible path. In
particular, there is no singular-step issue on the stated trajectory.

Use the principal logarithm near $z=1$, and set $H_0=H_0^{(2)}=0$.
Three solutions of $L_0f=0$ are
\begin{align*}
 f_E(z)&=e\,e^{-z},\\
 f_B(z)&=\sum_{k\geq0}\frac{(-1)^kz^k}{k!}
                      (\log z-H_k),\\
 f_A(z)&=\sum_{k\geq0}\frac{(-1)^kz^k}{k!}
       \left(\frac{(\log z)^2}{2}-H_k\log z
                         +\frac{H_k^2+H_k^{(2)}}2\right).
\end{align*}
Indeed $(\theta+z)f_E=0$, $(\theta+z)f_B=1$, and
$(\theta+z)f_A=\log z$. These identities follow either by termwise
calculation, or by multiplying the integrals of $e^z/z$ and
$e^z\log z/z$ by $e^{-z}$. If $b=f_B(1)$ and $a=f_A(1)$, their initial
jet matrix is
\[
 C_0=\begin{pmatrix}1&-1&0\\b&1-b&-1\\a&-a&1\end{pmatrix},
 \qquad\det C_0=1.
\]
Let $E_n,B_n,A_n$ be the corresponding rows after applying $P_n(\theta)$,
and let $U_n$ be the matrix with these three rows. Then
\[
 U_n=C_0\Phi_n,\qquad \det U_n=\det\Phi_n,
 \qquad\Phi_n^{-T}=C_0^T U_n^{-T}.
\]

Put $s=A_1+A_2+2\beta$. We establish bounds, with constants depending
only on the fixed trajectory,
\begin{equation}
 \|E_n\|\leq C^n n^h(n!)^s,\qquad
 \|B_n\|+\|A_n\|\leq C^n n^h(n!)^{s+d}
 \quad(n\geq1).
 \label{p3:eq:cmf-row-bounds}
\end{equation}
The same letter $C>1$ and integer $h$ may be enlarged below.

For the elementary row there is an exact series
\[
 P_n(\theta)e^{-z}
 =z^{Kn}\frac{((K+A_1)n)!((K+A_2)n)!}{(dn)!^2}
 \sum_{j\geq0}\frac{(-z)^j}{j!}
 \frac{((K+A_1)n+1)_j((K+A_2)n+1)_j}{(dn+1)_j^2}.
\]
The quotient of rising factorials is at most
$R^j$, where $R=(K+A_1)(K+A_2)/d^2$.
Applying $\theta^m$ for $m=0,1,2$ inserts $(Kn+j)^m$ at $z=1$.
The resulting absolute sum is $O(n^2)$ times the prefactor. The elementary
multinomial bounds
\[
 (qn)!\geq(n!)^q,\qquad (qn)!\leq q^{qn}(n!)^q\quad(q\in\mathbb N)
\]
then give the first assertion in \eqref{p3:eq:cmf-row-bounds}.

Here are details for the other two rows. Put $Q_{m,n}(X)=X^mP_n(X)$.
Termwise application of the finite-order differential operator is valid
on a neighborhood of $z=1$. The entries of $B_n$ and $A_n$, for
$m=0,1,2$, are respectively
\begin{align*}
 &\sum_{k\geq0}\frac{(-1)^k}{k!}
       \bigl(Q_{m,n}'(k)-H_k Q_{m,n}(k)\bigr),\\
 &\sum_{k\geq0}\frac{(-1)^k}{k!}
       \left(\frac{Q_{m,n}''(k)}2-H_kQ_{m,n}'(k)
              +\frac{H_k^2+H_k^{(2)}}2Q_{m,n}(k)\right).
\end{align*}
For $0\leq k<\beta n$, $P_n$ has a zero of multiplicity at least three,
so every displayed summand vanishes. For $\beta n\leq k<Kn$ the zero
is simple. Writing $k=\beta n+j$ with $0\leq j<dn$, one obtains
\begin{equation}
 \frac{|P_n'(k)|}{k!}
 =\frac{((\beta+A_1)n+j)!((\beta+A_2)n+j)!}{j!^2}
        (dn-1-j)!.
 \label{p3:eq:cmf-simple-root}
\end{equation}
Also $|P_n''(k)|\leq C\log(n+1)|P_n'(k)|$: remove the simple root
and use the logarithmic derivative of the remaining linear factors.
Their roots are integers, with multiplicity at most four, and their
nonzero distances from $k$ give harmonic sums of length $O(n)$.
The factors $X^m$ contribute only powers of $n$ and logarithms.

To bound \eqref{p3:eq:cmf-simple-root}, put
$m_i=(\beta+A_i)n$. Since $j<dn$,
\[
 (m_i+j)!\leq m_i!\,((K+A_i)n)^j,
 \qquad (dn-1-j)!\leq(dn)!.
\]
The extra factor $(Rn^2)^j/(j!)^2$ is at most $\exp(2\sqrt R\,n)$,
by the exponential series. Thus the sum of these $dn$ terms is bounded
by a polynomial in $n$ times $C^n(n!)^{s+d}$.

For $k=Kn+j$, all roots of $P_n$ are a nonzero integer distance from
$k$. Logarithmic differentiation gives
\[
 |P_n'(k)|\leq C\log(n+j+2)|P_n(k)|,\qquad
 |P_n''(k)|\leq C\log^2(n+j+2)|P_n(k)|.
\]
Moreover
\[
 \frac{|P_n(Kn+j)|}{(Kn+j)!}
 =\frac{((K+A_1)n)!((K+A_2)n)!}{(dn)!^2}
   \frac1{j!}
   \frac{((K+A_1)n+1)_j((K+A_2)n+1)_j}{(dn+1)_j^2}.
\]
The tail is therefore at most the elementary prefactor times a polynomial
in $n$ and $\log(n+2)$, since $\sum R^j(j+1)^h/j!$ is finite for
fixed $h,R$. This proves \eqref{p3:eq:cmf-row-bounds}.

Let $w_{1,n},w_{2,n},w_{3,n}$ denote the rows of the cofactor matrix of
$U_n$. Thus $w_{1,n}=B_n\times A_n$, and the other two rows are cross
products involving $E_n$. The determinant identity and Cauchy--Schwarz
give
\[
 \|w_{1,n}\|\geq\frac{\det U_n}{\|E_n\|}.
\]
On the other hand,
\[
 \|w_{2,n}\|+\|w_{3,n}\|
 \leq\|E_n\|(\|A_n\|+\|B_n\|).
\]
By \eqref{p3:eq:cmf-determinant} and $(qn)!\geq(n!)^q$,
$\det U_n\geq(n!)^{3s+2d}$. Combining the last inequalities yields
\begin{equation}
 \frac{\|w_{2,n}\|+\|w_{3,n}\|}{\|w_{1,n}\|}
 \leq\varepsilon_n:=\frac{C^n n^h}{(n!)^d}.
 \label{p3:eq:cmf-cofactor-separation}
\end{equation}
In particular $\sum_n\sqrt{\varepsilon_n}<\infty$.

For completeness, the following elementary measure argument turns this
operator estimate into a statement for fixed seeds. In any fixed bounded
ball in $\mathbb R^3$, the set
\[
 \left\{v:|w_{1,n}\cdot v|<
                  \sqrt{\varepsilon_n}\,\|w_{1,n}\|\right\}
\]
is a slab of volume at most a constant times $\sqrt{\varepsilon_n}$.
The first Borel--Cantelli lemma, followed by a countable union over balls,
shows that almost every $v$ belongs to only finitely many such slabs.
For each such $v$,
\[
 \left|\frac{w_{i,n}\cdot v}{w_{1,n}\cdot v}\right|
 \leq\|v\|\sqrt{\varepsilon_n}\longrightarrow0
 \qquad(i=2,3).
\]
Since $U_n^{-T}$ is its cofactor matrix divided by $\det U_n$, it follows
that
\[
 \frac{\det U_n}{w_{1,n}\cdot v}\Phi_n^{-T}v
 =C_0^T\begin{pmatrix}1\\
    (w_{2,n}\cdot v)/(w_{1,n}\cdot v)\\
    (w_{3,n}\cdot v)/(w_{1,n}\cdot v)\end{pmatrix}
 \longrightarrow C_0^T e_1=u.
\]
The prefactor is nonzero for all sufficiently large $n$ and may be
chosen arbitrarily nonzero for the finitely many earlier indices.

Finally, the source's Definition~12 excludes a proper linear subspace
of seeds when it requires a nonzero dominant-solution component.
That subspace has Lebesgue measure zero. Hence the preceding almost-everywhere
conclusion includes generic seeds in that definition. The canonical
state $S_3(0)$ satisfies
\[
 q^TS_3(0)=1,\qquad q=(1,1,1)^T,
\]
by the source's Proposition~1, whereas $q^Tu=0$. These two nonzero vectors
are not proportional. No scalar normalization of a sequence whose
projective limit is $u$ can converge to $S_3(0)$. This contradicts the
full-state assertion on the stated trajectories.
\end{proof}

\begin{remark}[Why a convergent scalar ratio does not rescue full-state isolation]
The two extraction vectors in rank three are
$p=(1,1,0)^T$ and $q=(1,1,1)^T$. Both satisfy $p^Tu=q^Tu=0$.
Thus the preceding theorem does not determine
$(p^T\Phi_n^{-T}v)/(q^T\Phi_n^{-T}v)$: it is a ratio of two
subdominant components after their common dominant component cancels.
Numerical convergence of this ratio to $\gamma$ is fully compatible
with failure of the asserted vector limit. Scalar convergence on these trajectories requires an argument in the
quotient by the $u$-direction or an equivalent analysis of subdominant
components; it cannot be justified by extending Theorem~3 unchanged.
The next section supplies such a proof using positive finite sums.
\end{remark}



\section{Scalar convergence on unequal rank-three trajectories}
\label{p3:sec:cmf-scalar}

\begin{theorem}[Scalar convergence survives on the same trajectories]
\label{p3:thm:unequal-rank-three-scalar}
Under the hypotheses of Theorem~\ref{p3:thm:full-state-counterexample}, let
$p=(1,1,0)^T$ and $q=(1,1,1)^T$. For every fixed nonzero
$v\in\mathbb C^3$, the denominator below is nonzero for all sufficiently
large $n$, and
\[
 \lim_{n\to\infty}\frac{p^T\Phi_n^{-T}v}{q^T\Phi_n^{-T}v}=\gamma.
\]
Thus the scalar conclusion of the source's Conjecture~1 is true on
these rank-three trajectories, although its full-state conclusion is false.
\end{theorem}

\begin{proof}
Retain the basis, matrices, and notation from the preceding proof.
Set $m_i=(\beta+A_i)n$, and, for $0\leq j<dn$, define
\[
 k_j=\beta n+j,\qquad
 W_{n,j}=\frac{(m_1+j)!(m_2+j)!(dn-1-j)!}{(j!)^2}>0,
 \qquad Z_n=\sum_{j=0}^{dn-1}W_{n,j}.
\]
Let $b(k)=(1,k,k^2)$ and $b'(k)=(0,1,2k)$. The simple-root
calculation in \eqref{p3:eq:cmf-simple-root}, including its sign, gives
\begin{align}
 B_n&=-\sum_{j=0}^{dn-1}W_{n,j}b(k_j)+T_{B,n},\\
 A_n&=-\sum_{j=0}^{dn-1}W_{n,j}
              \{R_{n,j}b(k_j)+b'(k_j)\}+T_{A,n},
 \label{p3:eq:cmf-weighted-rows}
\end{align}
where
\[
 R_{n,j}=H_{m_1+j}+H_{m_2+j}-2H_j-H_{dn-1-j}.
\]
To verify the latter formula, the logarithmic derivative of $P_n(X)/(X-k_j)$
at $k_j$ gives $P_n''(k_j)/(2P_n'(k_j))$; subtracting $H_{k_j}$
cancels the intermediate harmonic numbers. The derivative of the factor
$X^m$ supplies $b'(k_j)$. The remainders are precisely the terms
$k\geq Kn$ in the preceding proof and satisfy
\[
 \|T_{B,n}\|+\|T_{A,n}\|\leq C^n n^h(n!)^s.
\]
Since $Z_n\geq W_{n,0}=m_1!m_2!(dn-1)!\geq(n!)^{s+d}/(dn)$,
these remainders are negligible relative to $Z_n$ even after multiplication
by any fixed power of $n$.

We need two elementary concentration facts for the probability weights
$W_{n,j}/Z_n$. Put
\[
 c=\sqrt{\frac{(\beta+A_1)(\beta+A_2)}d}>0.
\]
Then
\begin{equation}
 \frac{j}{\sqrt n}\longrightarrow c\quad\hbox{in probability},
 \qquad
 \frac1{Z_n}\sum_j W_{n,j}|R_{n,j}+\gamma|\longrightarrow0.
 \label{p3:eq:cmf-concentration}
\end{equation}
Here are direct estimates proving both assertions. The exact successive
ratio is
\[
 \frac{W_{n,j+1}}{W_{n,j}}
 =\frac{(m_1+j+1)(m_2+j+1)}{(j+1)^2(dn-1-j)}.
\]
For fixed $0<\epsilon<c$, this ratio is uniformly greater than
$1+\delta$ below $(c-\epsilon/2)\sqrt n$, for a suitable $\delta>0$
and all sufficiently large $n$. Hence the total weight below
$(c-\epsilon)\sqrt n$, divided by $Z_n$, is
$O(\sqrt n\,e^{-c_\epsilon\sqrt n})$.
Choose $\eta>0$ so small that
\[
 \frac{(\beta+A_1+\eta)(\beta+A_2+\eta)}{d-\eta}
 <(c+\epsilon/2)^2.
\]
After an arbitrarily small enlargement of the numerator to absorb the
$1/n$ terms, the same ratio is uniformly less than a constant $q<1$
between $(c+\epsilon/2)\sqrt n$ and $\eta n$.
Thus the total weight from $(c+\epsilon)\sqrt n$ to $\eta n$ is
$O(e^{-c_\epsilon\sqrt n})Z_n$.
For $j\geq\eta n$, use
\[
 \frac{W_{n,j}}{W_{n,0}}
 =\frac{(m_1+1)_j(m_2+1)_j}{(j!)^2}
       \frac{(dn-1-j)!}{(dn-1)!}
 \leq\frac{e^{Cn}}{j!}.
\]
Indeed each rising-factorial quotient is at most $(Cn)^j/j!\leq e^{Cn}$,
and $(dn-1)!/(dn-1-j)!\geq j!$. The sum of this last range is at most
$dn\,e^{Cn}/\lfloor\eta n\rfloor!$ times $W_{n,0}$, which is negligible.
This proves concentration, with tails decreasing faster than any power
of $n$.

Uniformly for $|j/\sqrt n-c|\leq\epsilon$, the elementary asymptotic
$H_N=\log N+\gamma+o(1)$ gives
\[
 R_{n,j}+\gamma
 =\log\frac{c^2}{(j/\sqrt n)^2}+o(1).
\]
Throughout the full range $0\leq j<dn$ one has
$|R_{n,j}|\leq C\log(n+2)$. The tail estimates, followed by
$\epsilon\downarrow0$, prove the second assertion in
\eqref{p3:eq:cmf-concentration}.

Next write $E_n=E_{0,n}(1,e_{1,n},e_{2,n})$. The exact series for the
elementary solution in the previous proof, together with dominated
convergence, implies
\begin{gather*}
 E_{0,n}\sim e^{1-R}\frac{((K+A_1)n)!((K+A_2)n)!}{(dn)!^2},
 \qquad R=\frac{(K+A_1)(K+A_2)}{d^2},\\
 \frac{e_{1,n}}n\longrightarrow K,
 \qquad \frac{e_{2,n}}{n^2}\longrightarrow K^2.
\end{gather*}
For the moment estimates, domination follows from
$\sum_{j\geq0}R^j(j+1)^2/j!<\infty$; the leading exponential is nonzero.
In particular $E_{0,n}\ne0$ eventually.

For the fixed seed $v=(v_0,v_1,v_2)^T\ne0$, set
\[
 h_n(k)=\frac{(E_n\times b(k))\cdot v}{E_{0,n}}
 =v_0(e_{1,n}k^2-e_{2,n}k)
       +v_1(e_{2,n}-k^2)+v_2(k-e_{1,n}).
\]
Let $m=3$ if $v_0\ne0$, $m=2$ if $v_0=0,v_1\ne0$, and $m=1$ otherwise.
Uniformly on every fixed bounded range of $j/\sqrt n$,
\[
 \frac{h_n(k_j)}{n^m}\longrightarrow h_*(v)\ne0,
\]
where respectively
\[
 h_*(v)=v_0K\beta(\beta-K),\qquad
 h_*(v)=v_1(K^2-\beta^2),\qquad
 h_*(v)=v_2(\beta-K).
\]
These constants do not vanish because $K>\beta>0$.
On the entire summation range, $|h_n(k_j)|\leq C_v n^m$ and
$|h_n'(k_j)|\leq C_v n^{m-1}$.
Consequently \eqref{p3:eq:cmf-concentration} yields
\begin{align*}
 \frac{\sum_j W_{n,j}h_n(k_j)}{Z_n n^m}&\longrightarrow h_*(v),\\
 \frac{\sum_j W_{n,j}\{R_{n,j}h_n(k_j)+h_n'(k_j)\}}
      {Z_n n^m}&\longrightarrow-\gamma h_*(v).
\end{align*}
The remainders in \eqref{p3:eq:cmf-weighted-rows} do not affect these limits:
$\|E_n\|/|E_{0,n}|=O(n^2)$ and their factorial saving exceeds every
polynomial factor.

Finally $p^TC_0^T=(0,1,0)$ and $q^TC_0^T=(0,0,1)$, so Cramer's rule gives
\[
 \frac{p^T\Phi_n^{-T}v}{q^T\Phi_n^{-T}v}
 =\frac{(A_n\times E_n)\cdot v}{(E_n\times B_n)\cdot v}
 =-\frac{(E_n\times A_n)\cdot v}{(E_n\times B_n)\cdot v}
 \longrightarrow\gamma.
\]
The denominator is eventually nonzero by $h_*(v)\ne0$ and the first
weighted limit. This completes the proof.
\end{proof}


% Fragment; requires amsmath, amssymb, amsthm and theorem environments.
% Bibliography keys: KMLWSK26 (arXiv:2610.13919), DLMF.
\section{An exact variable-argument identity and a probability interpretation}
\label{p3:sec:meijer-extension}

Kalisch et al.~\cite[Proposition~1]{p3:KMLWSK26} evaluate two differential
forms in a Meijer $G$-function at $z=1$. Their Section~9.2 asks for an
extension involving $\gamma+\log z$. The identity component has the exact
extension below. It does not require a limiting argument in a matrix
recurrence, and does not by itself prove the dynamical assertions or the
identification with Rivoal's recurrences also discussed in that section.

Throughout, $\Omega=\mathbb C\setminus(-\infty,0]$, $\Log$ is the principal
logarithm, and $\theta=z\,d/dz$. For an integer $p\ge1$ and a real number
$a>-1$, define
\begin{align}
 \Phi_{p,a}(s)&=\Gamma(a-s)
       \bigl(\Gamma(-s)\Gamma(1+s)\bigr)^p,\label{p3:eq:meijer-kernel-a}\\
 G_{p,a}(z)&=\frac1{2\pi i}\int_{c-i\infty}^{c+i\infty}
       \Phi_{p,a}(s)z^s\,ds,
       \qquad -1<c<\min(0,a).\label{p3:eq:meijer-G-a}
\end{align}
Here $z^s=\exp(s\Log z)$. This is
$G^{p+1,p}_{p,p+1}\bigl(z\mid 0,\ldots,0;
0,\ldots,0,a\bigr)$ in the Mellin--Barnes convention
of~\cite[\S16.17]{p3:DLMF}. In particular $G_{p,0}$ is the $G_p$ of
\cite{p3:KMLWSK26}. The value of the integral is independent of the allowed
choice of $c$.

Put $\alpha=a+1>0$, and define polynomials $P_{p,a,j}(x)$ by
\begin{equation}\label{p3:eq:meijer-P-generating}
 \frac{\Gamma(\alpha-t)}{\Gamma(\alpha)}
 \left(\frac{\pi t}{\sin\pi t}\right)^p e^{xt}
   =\sum_{j\ge0}P_{p,a,j}(x)t^j
 \quad\bigl(|t|<\min(1,\alpha)\bigr).
\end{equation}

\begin{theorem}[Exact two-parameter identity]\label{p3:thm:meijer-exact-extension}
For $p\ge1$, $a>-1$, $0\le k\le p-1$ and $z\in\Omega$,
\begin{equation}\label{p3:eq:meijer-exact-extension}
 \left[z(1+\theta)^k+(-1)^p(\theta-a)\theta^k\right]G_{p,a}(z)
   =\Gamma(a+1)P_{p,a,p-k-1}(\Log z).
\end{equation}
In particular, for $p\ge2$,
\begin{align}
 \left[z(1+\theta)^{p-2}
       +(-1)^p(\theta-a)\theta^{p-2}\right]G_{p,a}(z)
   &=\Gamma(a+1)\bigl(\Log z-\psi(a+1)\bigr),\label{p3:eq:meijer-digamma}\\
 \left[z(1+\theta)^{p-1}
       +(-1)^p(\theta-a)\theta^{p-1}\right]G_{p,a}(z)
   &=\Gamma(a+1).\label{p3:eq:meijer-unit}
\end{align}
Consequently the quotient of the two left sides is exactly
$\Log z-\psi(a+1)$, with a nonzero constant denominator.
At $a=0$ this quotient is $\gamma+\Log z$.
\end{theorem}

\begin{remark}[A direct consequence of the existing equation]
For $a=0$, the two displayed identities can alternatively be read off
from the Meijer differential equation and the values at $z=1$ already
proved in~\cite{p3:KMLWSK26}. Indeed, $\theta z=z(\theta+1)$ gives
\[
 A_{p,k}=\theta^k\bigl(z+(-1)^p\theta\bigr),\qquad
 \theta^p\bigl(z+(-1)^p\theta\bigr)G_p=0.
\]
Thus $A_{p,p-1}G_p$ is constant and equals $1$ at $z=1$, whereas
$\theta(A_{p,p-2}G_p)=1$ and its value at $z=1$ is $\gamma$.
This proves the original variable-argument extension as a direct
consequence of their existing formulas. The contour argument below
also provides the shifted parameter $a$ and every polynomial in
\eqref{p3:eq:meijer-P-generating}.
\end{remark}

\begin{proof}
Euler's reflection and recurrence formulas give
\begin{equation}\label{p3:eq:meijer-shift-a}
 \Phi_{p,a}(s-1)=(-1)^p(a-s)\Phi_{p,a}(s).
\end{equation}
On any fixed vertical strip, away from its finitely many real poles,
Stirling's estimate and the reflection formula imply
\[
 |\Phi_{p,a}(\sigma+iT)|
 =O\left((1+|T|)^C
        e^{-(p+1/2)\pi|T|}\right).
\]
Thus the integral in~\eqref{p3:eq:meijer-G-a}, its derivatives of every fixed
order, and all contour shifts used below converge absolutely and locally
uniformly on $\Omega$. In particular, applying a polynomial in $\theta$
amounts to multiplication by that polynomial in $s$. Horizontal sides of
the rectangles with real parts $c-1,c$ tend to zero: the additional factor
$z^s$ has modulus $|z|^\sigma e^{-T\arg z}$ and $|\arg z|<\pi$.

Let $Q(s)=(s+1)^k$ and $A(s)=Q(s)\Phi_{p,a}(s)z^s$.
By~\eqref{p3:eq:meijer-shift-a}, the left side of
\eqref{p3:eq:meijer-exact-extension}, divided by $z$, is
\[
 \frac1{2\pi i}\int_{\Re s=c}
      \bigl(A(s)-A(s-1)\bigr)\,ds
 =\frac1{2\pi i}
      \left(\int_{\Re s=c}-\int_{\Re s=c-1}\right)A(s)\,ds.
\]
Both vertical lines are directed upward. The difference is the positively
oriented boundary integral of the intervening strip. The only pole in
that strip is $s=-1$: the poles of $\Gamma(a-s)$ lie to the right of $c$,
and the other integer poles are outside it. Therefore the expression is
$\operatorname{Res}_{s=-1}A(s)$.

Writing $s=-1+t$, the reflection formula gives the exact local expression
\[
 A(-1+t)=z^{-1}t^{k-p}\Gamma(a+1-t)
       \left(\frac{\pi t}{\sin\pi t}\right)^p e^{t\Log z}.
\]
Taking the coefficient of $t^{-1}$ proves
\eqref{p3:eq:meijer-exact-extension}. The coefficients of $t^0$ and $t^1$
in~\eqref{p3:eq:meijer-P-generating} are $1$ and $x-\psi(a+1)$.
Finally $\Gamma(a+1)>0$ and $\psi(1)=-\gamma$, proving the last assertions.
\end{proof}

\begin{remark}[Parameters and source scope]
For a nonnegative integer $a=m$, formula~\eqref{p3:eq:meijer-digamma}, divided
by $m!$, gives $\gamma+\Log z-H_m$. For $a=-1/2$, division by $\sqrt\pi$
gives $\gamma+\Log z+2\log2$; this is $\gamma+\log(4z)$ on the positive
axis. These special values of $\psi$ are classical. The theorem supplies
their exact occurrence in this family of differential forms. For fixed
rational $a,z$ the operator coefficients are rational, but no assertion
about rationality of the resulting constant or convergence of a dual
matrix system follows merely from this observation.
\end{remark}

\begin{proposition}[Probability and Bell-polynomial interpretation]
\label{p3:prop:meijer-probability}
Let $T$ have the gamma distribution of shape $a+1$ and rate $1$, and let
$L_1,\ldots,L_p$ be independent standard logistic variables, also independent
of $T$. Set $X_{p,a}=-\log T+L_1+\cdots+L_p$. Then
\begin{equation}\label{p3:eq:meijer-probability}
 P_{p,a,j}(x)=\frac1{j!}\mathbb E[(x+X_{p,a})^j].
\end{equation}
The cumulants of $X_{p,a}$ are
\[
 \kappa_1=-\psi(a+1),\qquad
 \kappa_m=(-1)^m\psi^{(m-1)}(a+1)
       +2p\,\mathbf1_{2\mid m}(m-1)!\zeta(m)\quad(m\ge2).
\]
Equivalently, normalize the complete exponential Bell polynomials by
\[
 \exp\left(\sum_{m\ge1}u_mt^m/m!\right)
 =\sum_{j\ge0}B_j(u_1,\ldots,u_j)t^j/j!.
\]
Then
\[
 P_{p,a,j}(x)=\frac{B_j(x+\kappa_1,\kappa_2,\ldots,\kappa_j)}{j!}.
\]
For $a=0$ and $w=\gamma+x$, the first four polynomials are
\[
 1,\qquad w,\qquad
 \frac{w^2+(2p+1)\zeta(2)}2,\qquad
 \frac{w^3+3(2p+1)\zeta(2)w+2\zeta(3)}6.
\]
The formula of index $j$ occurs in
Theorem~\ref{p3:thm:meijer-exact-extension} whenever $j\le p-1$.
\end{proposition}

\begin{proof}
The gamma density gives
$\mathbb E[e^{-t\log T}]=\Gamma(a+1-t)/\Gamma(a+1)$ for $t<a+1$.
A standard logistic variable has density $e^{-v}/(1+e^{-v})^2$; substituting
$u=e^{-v}$ and using Euler's beta integral yields
\[
 \mathbb E[e^{tL_i}]
   =\int_0^\infty\frac{u^{-t}}{(1+u)^2}\,du
   =\Gamma(1-t)\Gamma(1+t)=\frac{\pi t}{\sin\pi t}
       \qquad(|t|<1).
\]
Independence identifies~\eqref{p3:eq:meijer-P-generating} as the moment
generating function of $x+X_{p,a}$, proving~\eqref{p3:eq:meijer-probability}.
All moments exist because that generating function is finite in a
neighborhood of zero. Its logarithm has the expansion
\[
 xt+\log\Gamma(a+1-t)-\log\Gamma(a+1)
       +p\sum_{\ell\ge1}\frac{\zeta(2\ell)}{\ell}t^{2\ell}.
\]
The cumulant and Bell formulas follow by Taylor expansion. At $a=0$ use
$\psi^{(m-1)}(1)=(-1)^m(m-1)!\zeta(m)$ for $m\ge2$.
\end{proof}

\begin{proposition}[Classical incomplete-gamma endpoint]\label{p3:prop:meijer-p1}
For $a>-1$ and $z\in\Omega$,
\[
 G_{1,a}(z)=\int_0^\infty\frac{e^{-u}u^a}{u+z}\,du
   =\Gamma(a+1)z^a e^z\Gamma(-a,z).
\]
In particular $G_{1,0}(z)=e^z E_1(z)$, where $E_1(z)=\Gamma(0,z)$.
\end{proposition}
\begin{proof}
First take $z>0$. Mellin inversion for $(1+x)^{-1}$, followed by absolute
integration on $\Re s=c$, gives the first equality; the inner integral
is $\Gamma(a-s)$ and $c<a$. Alternatively its Mellin transform uniquely
identifies the same absolutely convergent Barnes integral. To evaluate the
real integral, write $(u+z)^{-1}=\int_0^\infty e^{-v(u+z)}\,dv$ and use
Tonelli's theorem. The result is
\[
 \Gamma(a+1)\int_0^\infty e^{-zv}(1+v)^{-a-1}\,dv
   =\Gamma(a+1)z^ae^z\Gamma(-a,z),
\]
by $w=z(1+v)$. All three expressions are holomorphic on $\Omega$, so
analytic continuation proves the stated domain. This is the classical
incomplete-gamma integral representation; see~\cite[\S8.6]{p3:DLMF}.
\end{proof}

\subsection*{An exact finite telescoping remainder}
The following certificate may be useful for checking computations. It is
an identity of analytic expressions, not a rational approximation theorem:
the residues themselves can contain $\gamma$ and zeta values.
For $a=0$ put
\[
 F_{p,N}(z)=-\sum_{n=0}^N
       \operatorname{Res}_{s=n}\bigl(\Phi_{p,0}(s)z^s\bigr),
 \qquad A_{p,k}=z(1+\theta)^k+(-1)^p\theta^{k+1}.
\]
For every $N\ge0$, $0\le k\le p-1$ and $z\in\Omega$,
\begin{equation}\label{p3:eq:meijer-finite-error}
 P_{p,0,p-k-1}(\Log z)-A_{p,k}F_{p,N}(z)
 =z\operatorname{Res}_{s=N}
       \bigl((s+1)^k\Phi_{p,0}(s)z^s\bigr).
\end{equation}
Indeed, applying $A_{p,k}/z$ turns each summand into the difference of the
residues of $A(s)$ and $A(s-1)$ from the preceding proof. The finite sum
telescopes to the residues at $-1$ and $N$, with exactly the signs in
\eqref{p3:eq:meijer-finite-error}.

More explicitly, define $H_N^{(m)}=\sum_{\nu=1}^N\nu^{-m}$, $H_N=H_N^{(1)}$
(empty sums are zero), and $w_N=\Log z+\gamma-H_N$. The right side of
\eqref{p3:eq:meijer-finite-error} is exactly
\begin{equation}\label{p3:eq:meijer-error-coeff}
 \frac{(-1)^{(N+1)(p+1)}z^{N+1}}{N!}[t^p](N+1+t)^k
 \exp\left(w_Nt+
  \sum_{m=2}^{p}
   \frac{(1+2p\mathbf1_{2\mid m})\zeta(m)+(-1)^mH_N^{(m)}}{m}t^m
       \right).
\end{equation}
To verify this, use
\[
 \Gamma(-N-t)=\frac{(-1)^{N+1}\Gamma(1-t)}
                    {N!\,t\prod_{\nu=1}^N(1+t/\nu)}
\]
and expand the logarithm at $t=0$; powers higher than $p$ cannot affect
$[t^p]$. For fixed $p,k,z$ this gives the rigorous bound
\[
 P_{p,0,p-k-1}(\Log z)-A_{p,k}F_{p,N}(z)
 =O_{p,k,z}\left(
     \frac{|z|^{N+1}(N+1)^k(1+\log(N+1))^p}{N!}\right).
\]
Here $H_N=O(\log(N+1))$, while $H_N^{(m)}\le\zeta(m)$ for $m\ge2$;
the coefficient in~\eqref{p3:eq:meijer-error-coeff} is a fixed finite
polynomial in these quantities. Thus no unproved asymptotic summation or
interchange of infinite residue sums is needed for this remainder bound.


\section{Exact differential order and the full solvable Galois group}
\label{p3:sec:diff-structure}

The identities in the preceding section expose a particularly simple
differential module behind the Meijer functions.  This section concerns
differentiation in $z$, not the difference equations in the trajectory
index.  Reducibility in one of these settings does not imply reducibility
in the other.

Put $\theta=z\,d/dz$, $\epsilon=(-1)^p$, and retain the function
$G_p=G_{p,0}$ defined by the Mellin--Barnes integral above.
The standard Meijer equation factors in the Ore ring $\mathbb C(z)[\theta]$:
\begin{equation}\label{p3:eq:factor-meijer}
 \mathcal L_p=z(1+\theta)^p+(-1)^p\theta^{p+1}
      =\theta^p\bigl(z+(-1)^p\theta\bigr).
\end{equation}
The identity uses $\theta z=z(\theta+1)$; factors must not be commuted.
The factorization itself is an elementary consequence of the standard
equation.  We give the precise order and Galois group, without a claim that
this application of classical differential algebra resolves an external
open conjecture.

\begin{theorem}[Minimality despite complete factorization]
\label{p3:thm:meijer-minimal}
For each integer $p\ge2$, the minimal homogeneous differential order of
$G_p$ over $\mathbb C(z)$ is exactly $p+1$.  Its local monodromy at $z=0$
has a single unipotent Jordan block of size $p+1$.
\end{theorem}
\begin{proof}
Equation~\eqref{p3:eq:factor-meijer} supplies the upper bound.  At the pole
$s=0$ of the Mellin kernel,
\[
 \Gamma(-s)^{p+1}\Gamma(1+s)^p
     =(-1)^{p+1}s^{-p-1}(1+O(s)).
\]
Closing the contour to the right gives a minus sign.  Thus the coefficient
of $(\log z)^p$ in the local Frobenius expansion of $G_p$ at zero has
constant term $(-1)^p/p!\ne0$.  More explicitly, the poles at
$s=0,1,2,\ldots$ give
\[
 G_p(z)=\sum_{j=0}^{p} h_j(z)(\log z)^j,
 \qquad h_j\ \hbox{holomorphic near }0,
 \qquad h_p(0)=(-1)^p/p!.
\]
The convergence follows equally from the regular singular equation and
its Frobenius construction; the residue at zero identifies the nonzero
top coefficient.  If $T$ denotes continuation once around zero, then
$(T-1)^{p+1}G_p=0$ and
$(T-1)^pG_p=(2\pi i)^p p!h_p\ne0$.
Consequently $G_p,(T-1)G_p,\ldots,(T-1)^pG_p$ are linearly independent
over $\mathbb C$.  Indeed, applying a suitable power of $T-1$ to a
shortest nontrivial linear relation gives a contradiction.  Every
rational-coefficient differential equation annihilating $G_p$ also
annihilates these continuations, and so has order at least $p+1$.
Their span is the full solution space of $\mathcal L_p$, proving the
Jordan assertion as well.
\end{proof}

Fix a simply connected domain avoiding zero, put $\ell=\log z$ and
$E=e^{\epsilon z}$, and choose primitives
\begin{equation}\label{p3:eq:J-primitives}
 J_j'(z)=\frac{E(z)\ell(z)^j}{z},\qquad 0\le j<p.
\end{equation}
All differential fields below consist of analytic or meromorphic germs
on such a domain and have constant field $\mathbb C$.

\begin{lemma}[A rational integration obstruction]\label{p3:lem:primitive-obstruction}
Let $\epsilon\in\mathbb C^*$, $E=e^{\epsilon z}$,
$F=\mathbb C(z,\ell)$ and $K=F(E)$.  For every nonzero polynomial
$P\in\mathbb C[X]$, the function $EP(\ell)/z$ has no primitive in $K$.
\end{lemma}
\begin{proof}
First, no nonzero $a\in F$ satisfies $a'/a=c$ for a nonzero constant
$c$.  Expand at $\ell=\infty$:
$a=r(z)\ell^d(1+O(\ell^{-1}))$ with $r\in\mathbb C(z)^*$ and
$d\in\mathbb Z$.  Comparing leading powers gives $r'/r=c$, which is
impossible for a rational function.  This also proves that $E$ is
transcendental over $F$: if it were algebraic, its field norm would be
a nonzero element of $F$ with logarithmic derivative a positive integer
multiple of $\epsilon$.  The transcendence of $\ell$ over $\mathbb C(z)$
follows from its infinite additive monodromy.

Suppose $f\in F(E)$ satisfies $f'=EP(\ell)/z$.  Expand the rational
function $f$ as a formal Laurent series in the transcendental variable
$E$, with coefficients in $F$.  Differentiation takes the coefficient
of $E^j$ to $a_j'+j\epsilon a_j$.  For $j\ne0,1$ the preceding
observation forces $a_j=0$, and for $j=0$ it forces $a_0$ to be constant.
Thus $f=c+Ea$ with $a\in F$, and
\begin{equation}\label{p3:eq:rational-a}
 a'+\epsilon a=P(\ell)/z.
\end{equation}

We claim that $a$ is polynomial in $\ell$.  A finite pole in that
variable, along an irreducible polynomial $h\in\mathbb C(z)[\ell]$,
would acquire one higher order under differentiation unless $h$ divided
$h'$.  Such a divisor cannot occur.  An algebraic root $\alpha(z)$ of
$h$ would satisfy $\alpha'=1/z$ (differentiate
$h(z,\alpha(z))=0$), whereas an algebraic function has a Puiseux
expansion at zero whose derivative has zero coefficient of $z^{-1}$.
The right side of~\eqref{p3:eq:rational-a} has no finite poles in $\ell$,
so all such poles are excluded.

Write $a=\sum_{j=0}^d a_j(z)\ell^j$.  If $d>\deg P$, its top
coefficient satisfies $a_d'+\epsilon a_d=0$, again impossible.
If $d<\deg P$, the top power on the right is absent on the left.
Therefore $d=\deg P$ and the top coefficient must satisfy
$a_d'+\epsilon a_d=c/z$ with $c\ne0$.  A rational solution can have
no pole away from zero, since differentiation would raise its order;
a pole at zero would similarly produce a pole of order at least two.
With no pole at zero, the left side is holomorphic there and cannot
equal $c/z$.  This final contradiction proves the lemma.
\end{proof}

\begin{theorem}[Full differential Galois group]\label{p3:thm:full-meijer-group}
For every $p\ge2$ the Picard--Vessiot field of $\mathcal L_p$ is
\[
 \mathcal P_p=\mathbb C(z)(E,\ell,J_0,\ldots,J_{p-1}).
\]
The $p+2$ displayed generators are algebraically independent over
$\mathbb C(z)$.  Every differential automorphism is uniquely specified
by $a\in\mathbb C^*$, $b\in\mathbb C$, and
$c_0,\ldots,c_{p-1}\in\mathbb C$ through
\begin{align*}
 E&\longmapsto aE,&\ell&\longmapsto\ell+b,\\
 J_j&\longmapsto
 a\sum_{k=0}^j\binom jk b^{j-k}J_k+c_j
 &&(0\le j<p).
\end{align*}
In particular the group is
$\mathbb G_a^p\rtimes(\mathbb G_m\times\mathbb G_a)$,
with the action specified by these formulas, and has dimension $p+2$.
\end{theorem}
\begin{proof}
The lemma proved the algebraic independence of $E$ and $\ell$.
The Kolchin--Ostrowski theorem for antiderivatives
\cite[Theorem~2.3]{p3:Srinivasan2008} says that algebraic dependence among
$J_0,\ldots,J_{p-1}$ over $K$ would give a nonzero constant linear
combination belonging to $K$.  Its derivative would contradict
Lemma~\ref{p3:lem:primitive-obstruction}.  This proves the asserted
algebraic independence.

A full fundamental family of solutions is
\[
 E^{-1},\ E^{-1}J_0,\ldots,E^{-1}J_{p-1}.
\]
Indeed $(\theta+\epsilon z)(E^{-1}J_j)=\ell^j$ and its $p$th
Euler derivative vanishes.  Applying $\theta+\epsilon z$ to a linear
relation proves independence.  Their differential field contains $E$
and every $J_j$; since $p\ge2$, it also contains
$\ell=(\theta+\epsilon z)(E^{-1}J_1)$.
It is therefore exactly the displayed field.

The differential equations force $E\mapsto aE$ and
$\ell\mapsto\ell+b$.  Differentiating the asserted formula for $J_j$
shows that it is the general possible image, up to the constant $c_j$.
Conversely, algebraic independence ensures that every indicated
assignment extends to an automorphism of the rational function field,
and the equations verify that it commutes with differentiation.
The subgroup $a=1,b=0$ is the translation group $\mathbb G_a^p$;
setting all $c_j=0$ gives the complementary
$\mathbb G_m\times\mathbb G_a$.  This proves the group description.
\end{proof}

\begin{remark}[Functions versus values]
The preceding independence is over $\mathbb C(z)$ and concerns
functions.  It makes no claim about independence of their values at
$z=1$.  In fact the Meijer identities explicitly produce relations
among such values involving $\gamma$ and zeta values.  No specialization
principle applicable to arbitrary Liouvillian functions turns this
theorem into a proof of the irrationality of $\gamma$.
\end{remark}


\section{Difference irreducibility and denominator obstructions}
\label{p3:sec:gamma-arithmetic}

This section concerns the trajectory index $n$. Its conclusions are
independent of the differential factorization in the variable $z$.
We prove exact difference order and arithmetic obstructions for the
Rivoal recurrence recovered in \cite[\S8.4]{p3:KMLWSK26}. These are structural
deductions; they are not presented as resolutions of external open
conjectures or as a proof of the irrationality of $\gamma$.

\subsection{An elementary irreducibility certificate}

Write $Ef(n)=f(n+1)$, with the Ore rule $E c(n)=c(n+1)E$.
For a fixed complex parameter $x$, the recurrence in
\cite[(1.3)]{p3:Riv09} is $L_xu=0$, where
\begin{equation}\label{p3:eq:ga-L}
 L_x=\sum_{j=0}^3 a_j(n,x)E^j
\end{equation}
and
\begin{align*}
 a_0={}&(n+2)^2(3x-8n-22)(3x-8n-30),\\
 a_1={}&B_1(n,x)(3x-8n-30),\qquad
 a_2=-B_2(n,x)(3x-8n-14),\\
 a_3={}&-(n+3)^2(3x-8n-14)(3x-8n-22).
\end{align*}
Here the two cubic polynomials are
\begin{align*}
 B_1(n,x)={}&24n^3+(162-57x)n^2+(26x^2-258x+356)n\\
            &\quad-3x^3+59x^2-285x+254,\\
 B_2(n,x)={}&24n^3+(7x+210)n^2+(-6x^2+56x+600)n\\
            &\quad-15x^2+96x+564.
\end{align*}
Sequences in the general-parameter statements are understood as germs
at $n=+\infty$: changing finitely many terms is immaterial.

\begin{lemma}[A coefficient obstruction to a hypergeometric solution]
\label{p3:lem:ga-obstruction}
Let $L=\sum_{j=0}^3 b_j(n)E^j$ have polynomial coefficients
$b_j=\alpha_j n^m+\beta_j n^{m-1}+O(n^{m-2})$, where
\[
 \sum_{j=0}^3\alpha_j T^j=c(1-T)^3,\qquad
 c\ne0,\qquad \sum_{j=0}^3\beta_j\ne0.
\]
Then $L$ has no right factor of order one over $\mathbb C(n)$.
\end{lemma}
\begin{proof}
An order-one right factor can be made $E-r(n)$, where
$r\in\mathbb C(n)^*$, since the constant and leading coefficients
of $L$ are nonzero. Its nonzero solution has
$u_{n+j}/u_n=\prod_{i=0}^{j-1}r(n+i)$. Write
$r(n)=\lambda n^d(1+O(n^{-1}))$, with $d\in\mathbb Z$ and
$\lambda\ne0$. If $d>0$, the $j=3$ term of $Lu/u$ has uniquely
largest degree at infinity; if $d<0$, the $j=0$ term does. Both are
impossible. Hence $d=0$, and the leading coefficient equation is
$c(1-\lambda)^3=0$, so $\lambda=1$.
Now $r(n)=1+b/n+O(n^{-2})$, whence
$u_{n+j}/u_n=1+jb/n+O(n^{-2})$. The coefficient of $n^{m-1}$
in $Lu/u$ is
\[
 \sum_j\beta_j+b\sum_jj\alpha_j=\sum_j\beta_j\ne0,
\]
again a contradiction.
\end{proof}

\begin{theorem}[Exact difference order for every nonzero parameter]
\label{p3:thm:ga-irreducible}
For fixed $x\in\mathbb C$, $L_x$ is irreducible in
$\mathbb C(n)[E]$ if and only if $x\ne0$. For $x\ne0$,
every solution which is not eventually zero
has minimal recurrence order exactly three over $\mathbb C(n)$.
\end{theorem}
\begin{proof}
Expansion of the displayed coefficients gives
\[
 (\alpha_0,\alpha_1,\alpha_2,\alpha_3)=64(1,-3,3,-1),
\]
\[
 (\beta_0,\beta_1,\beta_2,\beta_3)
 =(672-48x,\;528x-2016,\;2016-16x,\;48x-672).
\]
In particular, $\sum\beta_j=512x\ne0$, so the lemma excludes
an order-one right factor. To exclude an order-one left factor,
use the formal adjoint, defined by $f(n)^*=f(n)$ and $E^*=E^{-1}$
in the Laurent Ore ring. Multiplying $L_x^*$ on the left by $E^3$
gives the polynomial operator
\[
 L_x^\vee=\sum_{j=0}^3 a_j(n+3-j,x)E^{3-j}.
\]
Its leading characteristic polynomial is
$64(T-1)^3$, and the sum of its next coefficients is
\[
 \sum_j\{\beta_j+4(3-j)\alpha_j\}=512x,
\]
since $\sum\alpha_j=\sum j\alpha_j=0$.
The lemma therefore also excludes an order-one right factor of
$L_x^\vee$. The adjoint reverses a factorization, so this excludes
an order-one left factor of $L_x$. Nonzero constant coefficients
ensure that shifting the adjoint factors to nonnegative powers of
$E$ introduces no loss of order. Every proper factorization of an
order-three operator has an order-one factor on one side.

Finally, choose a minimal monic annihilator of a nonzero solution
germ. Right division of $L_x$ by this annihilator has a remainder
of smaller order which still annihilates the germ. Minimality
forces the remainder to be zero. Irreducibility makes the minimal
annihilator have order three.

For completeness, at $x=0$ direct substitution gives $L_0(n+1)=0$.
Thus $E-(n+2)/(n+1)$ is a right factor and $L_0$ is reducible.
\end{proof}

The exceptional parameter $x=0$ is therefore genuine.
The theorem rules out an order-two recurrence for these particular
sequences. It does not rule out a different order-two construction
of rational approximations to $\gamma$.

\subsection{The integral solution lattice is zero}

At $x=1$, write the recurrence with positive leading coefficient as
\begin{align}
 &(n+3)^2(8n+11)(8n+19)u_{n+3}\notag\\
 &\quad=(n+3)(8n+11)(24n^2+145n+215)u_{n+2}\notag\\
 &\qquad\quad-(8n+27)(24n^3+105n^2+124n+25)u_{n+1}
 +(n+2)^2(8n+19)(8n+27)u_n.
 \label{p3:eq:ga-base}
\end{align}
Its coefficients of $u_n$ and $u_{n+3}$ do not vanish for $n\ge0$.
Thus its solution space is three-dimensional, and an eventually zero
solution is identically zero.

\begin{lemma}[Growth and the only possible rational generating functions]
\label{p3:lem:ga-growth}
For fixed $x\ne0$, every solution of $L_xu=0$ satisfies, for every
$\varepsilon>0$, $|u_n|\le C_\varepsilon e^{\varepsilon n}$.
The ordinary generating function of a solution which is not
eventually zero is D-finite and is not rational.
\end{lemma}
\begin{proof}
The companion matrices converge to a constant matrix all of whose
eigenvalues equal one, because the limiting characteristic polynomial
is $(T-1)^3$. For any $\delta>0$ there is a vector norm whose induced
matrix norm gives this limiting matrix norm at most $1+\delta$:
rescale the coordinates in each Jordan block to make its off-diagonal
entries arbitrarily small. In that norm the companion matrices
eventually have norm at most $1+2\delta$. Their products imply the
growth assertion, after decreasing $\delta$ and absorbing finitely
many terms into $C_\varepsilon$.

Let $F(t)=\sum_{n\ge0}u_nt^n$ and $\vartheta=t\,d/dt$.
Multiplication of the recurrence by $t^{n+3}$ and summation gives
\begin{equation}\label{p3:eq:ga-generating}
 \left(\sum_{j=0}^3t^{3-j}a_j(\vartheta-j,x)\right)F(t)=P(t),
\end{equation}
where $P$ is a polynomial accounting for initial terms. This proves
D-finiteness, by homogenizing the equation if necessary. The
coefficient of $F^{(4)}$ on the left is
\[
 t^4\sum_{j=0}^3\alpha_j t^{3-j}=64t^4(t-1)^3.
\]
If $F$ were rational, it could have no pole at an ordinary point of
this equation: its highest derivative would then have a pole of
strictly greater order than all the lower derivative terms.
It is analytic at zero, so its only possible finite pole is $t=1$.
Partial fractions imply that $u_n$ is eventually a polynomial in $n$.
A nonzero such polynomial supplies a rational quotient $u_{n+1}/u_n$,
contradicting Lemma~\ref{p3:lem:ga-obstruction}; a zero polynomial is
the excluded eventually zero case.
\end{proof}

\begin{theorem}[No nonzero bounded-denominator solution]
\label{p3:thm:ga-integral-lattice}
The only solution of \eqref{p3:eq:ga-base} taking values in $D^{-1}\mathbb Z$
for one fixed positive integer $D$ is the zero solution. In particular,
its integral solution lattice is $\{0\}$.
More quantitatively, for a nonzero rational solution put
\[
 d_N=\operatorname{lcm}\{\operatorname{den}(u_0),\ldots,
                         \operatorname{den}(u_N)\}.
\]
Then
\begin{equation}\label{p3:eq:ga-exponential-lcm}
 \limsup_{N\to\infty}\frac{\log d_N}{N}>0.
\end{equation}
\end{theorem}
\begin{proof}
First suppose the denominator is bounded. Then $DF$ is an
integer-coefficient D-finite series with radius of convergence at least
one, by the growth lemma. If it has infinitely many nonzero
coefficients, its radius cannot exceed one. The P\'olya--Carlson theorem
states that an integer series of radius one is rational or has the
unit circle as a natural boundary. D-finiteness excludes a natural
boundary because a polynomial differential equation has only finitely
many possible singular points. Thus $F$ is rational, contrary to the
lemma. The finite-series case is excluded by the same lemma.

For the stronger assertion, suppose instead that $\log d_N=o(N)$.
For a rational number $a$, its logarithmic height satisfies
$h(a)\le\log\operatorname{den}(a)+\log^+|a|$. Consequently
\[
 \max_{0\le n\le N}h(u_n)
 \le \log d_N+\max_{0\le n\le N}\log^+|u_n|=o(N),
\]
where the last estimate uses the growth lemma for every $\varepsilon$.
The univariate rationality theorem of
Bell--Chen--Nguyen--Zannier \cite[Proposition 2.1]{p3:BCNZ23} now says
that $F$ is rational: its hypotheses are precisely D-finiteness,
sublinear maximal height, and $\log d_N=o(N)$. This contradiction
proves \eqref{p3:eq:ga-exponential-lcm}. The limsup assertion is equivalent
to the failure of $\log d_N=o(N)$ because these quotients are
nonnegative.
\end{proof}

The same argument applies to every rational solution germ of $L_x$
with fixed $x\ne0$, whenever its coefficients are rational. The
resulting rational-coefficient series is D-finite over $\mathbb Q(t)$:
existence of a polynomial-coefficient differential relation descends
from $\mathbb C$ to $\mathbb Q$ by rational linear algebra on its
Taylor coefficients. The
formulation at $x=1$ avoids exceptional finite indices and is the
one relevant to Euler's constant. No effective universal constant
in \eqref{p3:eq:ga-exponential-lcm} is asserted.

\subsection{A complementary local obstruction from the Casoratian}

\begin{proposition}[Exact Casoratian and prime-by-prime rank bound]
\label{p3:prop:ga-casoratian}
For three solutions of \eqref{p3:eq:ga-base}, let
$W_N=\det(u^{(i)}_{N+j})_{1\le i\le3,\,0\le j\le2}$. Then
\begin{equation}\label{p3:eq:ga-casoratian}
 W_N=\frac{4(8N+11)(8N+19)}{209(N+2)^2}\,W_0.
\end{equation}
For every prime $p$, the rational solution space consisting of
sequences whose $p$-adic valuations are bounded below has dimension
at most two over $\mathbb Q$.
\end{proposition}
\begin{proof}
Taking the determinant of the companion update gives
\[
 \frac{W_{N+1}}{W_N}
   =\frac{(N+2)^2(8N+27)}{(N+3)^2(8N+11)}.
\]
This identity, interpreted without division when $W_N=0$, telescopes
to \eqref{p3:eq:ga-casoratian}. For an independent rational triple,
$W_0\ne0$. At $N=p^k-2$, for all sufficiently large $k$,
\[
 v_p(W_{p^k-2})
   =v_p(4W_0/209)+v_p(15)-2k.
\]
Indeed $(8N+11)(8N+19)=(8p^k-5)(8p^k+3)$ has valuation
$v_p(15)$ once $k$ is sufficiently large. The displayed valuation
tends to $-\infty$. If all three sequences had valuations bounded
below, every determinant term would have a uniform lower bound,
and so would their sum. This is a contradiction.
\end{proof}

This local statement does not say that every nonzero solution has an
unbounded denominator at every prime. Also $W_N\to256W_0/209$, not
zero. A direct ``integer determinant tends to zero'' argument is
therefore unavailable; the generating-function argument is what
strengthens a rank restriction to the zero integral lattice.

\subsection{Factorial denominators in the original CMF normalization}

The gauge transformation in \cite[\S8.4]{p3:KMLWSK26} permits a stronger
statement for the unrescaled dual matrix trajectory. Put
\[
 A(n)=\frac1{(n+1)^4}
 \begin{pmatrix}
 6n^2+13n+10&-8n-5&3\\
 3n^3+3n^2-3n-6&-3n^2+3n+3&n-2\\
 n^4-2n^3-6n^2-3n+3&6n^2+2n-1&1-2n
 \end{pmatrix},
\]
and define row-vector solutions by $u_{n+1}=u_n A(n)$.
The following explicit gauge is useful for making the normalization
unambiguous:
\[
 G(n)=\begin{pmatrix}
 0&\dfrac{8n+11}{(n+1)^2}
  &\dfrac{24n^2+97n+94}{(n+1)^2(n+2)}\\[5pt]
 -1&\dfrac{3n^2-n-7}{(n+1)^2}
  &\dfrac{15n^4+69n^3+66n^2-78n-111}{(n+1)^2(n+2)^2}\\[5pt]
 2&\dfrac{-6n^2-6n+3}{(n+1)^2}
  &\dfrac{8n^5+21n^4-38n^3-131n^2-48n+53}{(n+1)^2(n+2)^2}
 \end{pmatrix}.
\]
Let $C(n)$ be the row companion matrix of \eqref{p3:eq:ga-base}, so that
\[
 (w_n,w_{n+1},w_{n+2})C(n)=(w_{n+1},w_{n+2},w_{n+3}).
\]
Direct polynomial identities give
\begin{equation}\label{p3:eq:ga-gauge}
 \begin{gathered}
 C(n)=(n+1)^2G(n)^{-1}A(n)G(n+1),\\
 \det G(n)=\frac{(8n+11)(8n+19)}{(n+2)^2},\qquad
 \det A(n)=(n+1)^{-6}.
 \end{gathered}
\end{equation}
These identities may be checked by clearing denominators; no
asymptotic matching or numerical inference is required.

\begin{theorem}[Raw-vector factorial denominator barrier]
\label{p3:thm:ga-factorial}
Let $u_0\in\mathbb Q^3\setminus\{0\}$, and let $D_n$ be the least
positive integer for which $D_nu_n\in\mathbb Z^3$. For every
$\varepsilon>0$ there is $c_\varepsilon>0$ such that
\begin{equation}\label{p3:eq:ga-factorial-bound}
 D_n\ge c_\varepsilon(n!)^2e^{-\varepsilon n},
 \qquad D_n\mid D_0(n!)^4.
\end{equation}
In particular,
\[
 2\le\liminf_{n\to\infty}\frac{\log D_n}{n\log n}
 \le\limsup_{n\to\infty}\frac{\log D_n}{n\log n}\le4.
\]
\end{theorem}
\begin{proof}
By \eqref{p3:eq:ga-gauge},
$v_n=(n!)^2u_nG(n)$ satisfies $v_{n+1}=v_nC(n)$.
Its coordinates are consecutive terms of a scalar solution of
\eqref{p3:eq:ga-base}. Lemma~\ref{p3:lem:ga-growth}, together with the
polynomial growth of the rational matrix $G(n)^{-1}$, gives
\[
 \|u_n\|_\infty\le C_\varepsilon(n!)^{-2}e^{\varepsilon n}.
\]
All matrices $A(n)$ are invertible for $n\ge0$, so $u_n\ne0$.
The nonzero integer vector $D_nu_n$ has sup norm at least one.
Combining the two inequalities gives the lower bound in
\eqref{p3:eq:ga-factorial-bound}. The entries of $(n+1)^4A(n)$
are integers. Induction therefore gives the divisibility assertion.
Stirling's estimate $\log(n!)=n\log n+O(n)$ gives the last statement.
\end{proof}

\begin{remark}[What this denominator bound does not establish]
Theorem~\ref{p3:thm:ga-factorial} is tied to the raw vector $u_n$.
A common rescaling changes $D_n$, and taking a ratio of two linear
forms in $u_n$ can cancel large common factors. Hence the theorem
does not give the height of the reduced scalar approximants and
does not establish their reported irrationality measure parameter
$\delta=-1$. A gcd-controlled lower bound for the reduced fractions
would require an additional argument. Nor does the theorem alone
imply any irrationality statement for their limit.
\end{remark}

\paragraph{Prior-art boundary.}
Rivoal's original denominator conjectures for his specially normalized
approximants were proved by Hessami Pilehrood and Hessami Pilehrood
\cite{p3:HHP11}; they must not be counted as newly solved problems.
The exact-order statement above concerns the existing recurrence,
and the denominator results concern its full solution space or the
specified CMF normalization. No general CMF convergence conjecture
is settled by these arguments.


% Research fragment, 2026-10-18. Not a claim to a new general CF theorem.
\section{The all-order pairing problem of Xu and You}
\label{p3:sec:xu-you-pairing}

Xu and You~\cite{p3:XuYou2014}, \emph{Continued fraction inequalities for the
Euler--Mascheroni constant}, Journal of Inequalities and Applications
\textbf{2014}, 343, DOI \texttt{10.1186/1029-242X-2014-343}, formulate an
explicit ``Open problem'' immediately after Theorem~1: do the uniquely
successively optimized coefficients of their continued fraction satisfy
$a_{2j+1}=-a_{2j}$ for every $j\geq1$?  They compute through $a_{13}$;
Remark~4 also asks for a method to find the subsequent coefficients.
The following argument answers the coefficient question through the
classical correspondence between positive moments and Stieltjes fractions.
A first-priority claim requires a broader bibliographic check; the theorem
below is a mathematical resolution of the explicitly stated question.

Put $h(z)=\psi(z+1)-\log z$, so that
$H_n-\log n=\gamma+h(n)$.  Precisely, the fraction in question is
\[
 R_k(z)=\cfrac{a_1}{z+\cfrac{a_2z}{z+\cfrac{a_3z}{
 z+\ddots+\cfrac{a_{k-1}z}{z+a_k}}}},\qquad a_1=\frac12.
\]
For $k=1$ this means $R_1(z)=a_1/z$. Each finite fraction is interpreted
as a rational function after cancellation; a zero internal denominator
in an unsimplified nesting can therefore be removable. Successive optimization means
choosing the new coefficient to cancel the first remaining term of the
Poincare expansion as $z\to+\infty$.

\begin{theorem}[All-order coefficient pairing]\label{p3:thm:xu-pairing}
There is a unique infinite sequence of nonzero rational coefficients
$a_j$ such that
\[
 h(z)-R_k(z)=C_kz^{-k-1}+O(z^{-k-2})\quad(k\geq1),\qquad C_k\ne0.
\]
Here uniqueness refers to the successive optimization rule just stated.
Every even-indexed coefficient is positive and
\[
 a_{2j+1}=-a_{2j}\quad(j\geq1),\qquad
 C_k=(-1)^k\prod_{j=1}^{k+1}a_j.
\]
All the coefficients are given by the positive Hankel determinants below.
In particular,
\[
 a_{14}=\frac{3754087889491759}{2440840521848406},\qquad
 a_{15}=-a_{14}.
\]
\end{theorem}

\begin{proof}
Binet's digamma integral~\cite[5.9.15]{p3:DLMF}, together with
$\psi(z+1)=\psi(z)+1/z$, gives, for $z>0$,
\[
 h(z)=\frac1{2z}-2\int_0^\infty
       \frac{t\,dt}{(z^2+t^2)(e^{2\pi t}-1)}.
\]
Consequently, on writing $x=1/z$,
\begin{equation}\label{p3:eq:xy-stieltjes}
 F(x):=\frac{2h(1/x)}x=1-xU(x^2),\qquad
 U(y)=\int_0^\infty\frac{d\mu(s)}{1+ys},\quad
 d\mu(s)=\frac{2\,ds}{e^{2\pi\sqrt{s}}-1}.
\end{equation}
This positive measure has infinite support and all finite moments
\[
 \mu_n=\int_0^\infty s^n\,d\mu(s)
 =\frac{4\Gamma(2n+2)\zeta(2n+2)}{(2\pi)^{2n+2}}
 =\frac{2|B_{2n+2}|}{2n+2}\in\mathbb Q_{>0}.
\]
The geometric identity under the integral implies the all-orders
Poincare expansion $U(y)\sim\sum_{n\geq0}(-1)^n\mu_ny^n$ as
$y\downarrow0$; after any fixed number of terms its remainder is bounded
by the next moment times the corresponding power of $y$.

For $r\geq0$ and $\epsilon\in\{0,1\}$ define
\[
 D_r^{(\epsilon)}=\det(\mu_{i+j+\epsilon})_{0\leq i,j<r},
 \qquad D_0^{(\epsilon)}=1.
\]
Every such determinant is strictly positive: the associated matrix is
the Gram matrix of $1,s,\ldots,s^{r-1}$ in the positive measure
$s^\epsilon d\mu(s)$, and a nonzero polynomial cannot vanish on its
support.  The classical Stieltjes moment/continued-fraction theorem
\cite{p3:Stieltjes1894,p3:Wall1948} therefore gives the unique formal fraction
\begin{equation}\label{p3:eq:xy-positive-cf}
 U(y)\sim\cfrac{\mu_0}{1+\cfrac{\lambda_1y}{1+
                  \cfrac{\lambda_2y}{1+\ddots}}},\qquad \lambda_j>0,
\end{equation}
where the explicit determinant formulas are
\[
 \lambda_{2r-1}=\frac{D_r^{(1)}D_{r-1}^{(0)}}
                         {D_r^{(0)}D_{r-1}^{(1)}},\qquad
 \lambda_{2r}=\frac{D_{r+1}^{(0)}D_{r-1}^{(1)}}
                       {D_r^{(1)}D_r^{(0)}}\quad(r\geq1).
\]
The formal statement used here is the standard finite-moment version:
the fraction truncated after $\lambda_m$ agrees with the moment series
through degree $m$, and its next error coefficient is nonzero.  It follows
also from Gram orthogonalization and the two adjacent Pad\'e determinant
formulas.  No convergence of a divergent moment power series is assumed.
All $\lambda_j$ are positive rational numbers.

Set
\[
 c_1=\mu_0=\frac16,\qquad c_{j+1}=\frac{\lambda_j}{c_j}\quad(j\geq1).
\]
The finite pair-contraction identity, valid as a rational-function
identity for every $m\geq1$, is
\begin{equation}\label{p3:eq:xy-pair-contract}
 \frac{1}{1+\displaystyle\mathop{\mathrm K}_{j=1}^{2m}
                         \frac{b_jx}{1}}
 =1-\frac{c_1x}{1+\displaystyle\mathop{\mathrm K}_{j=1}^{m-1}
                         \frac{c_jc_{j+1}x^2}{1}},\qquad
 b_{2j-1}=c_j,\quad b_{2j}=-c_j.
\end{equation}
Here $\mathrm K$ is the usual finite generalized continued fraction,
with plus signs between successive levels, and an empty fraction is zero.
At $m=1$ the left side is $1/(1+c_1x/(1-c_1x))=1-c_1x$.
The general identity follows by clearing the denominators, or by
induction on continuants.
Applying it to \eqref{p3:eq:xy-positive-cf} and using $c_jc_{j+1}=\lambda_j$
shows, as a formal expansion to every finite order, that $F$ has the
normalized continued-fraction coefficients
\[
 c_1,-c_1,c_2,-c_2,c_3,-c_3,\ldots.
\]
Returning to $R_k$ therefore gives
$a_{2j}=c_j$ and $a_{2j+1}=-c_j$.

For completeness, a formal series $G(x)=1+O(x)$ has at most one
normalized continued-fraction expansion with nonzero coefficients:
if $G=1/(1+b x G_1)$, then the coefficient of $x$ determines $b$, and
$G_1=(G^{-1}-1)/(bx)$ determines the next tail.  Iteration proves
uniqueness.  The continuant determinant identity shows that truncation
after $b_1,\ldots,b_N$ has first omitted term
$(-1)^{N+1}b_1\cdots b_{N+1}x^{N+1}$.
Multiplication by $x/2$, with $N=k-1$ and $b_j=a_{j+1}$, proves the
stated error constant and its nonvanishing.  In particular these are
exactly the successively optimal coefficients of the source, rather
than merely another continued fraction for the same function.
\end{proof}

\begin{corollary}[Convergent bounds from the paired truncations]
For each fixed $z>0$, the odd truncations $R_{2m+1}(z)$ converge to
$h(z)$ as $m\to\infty$.  For every $m\geq1$,
\[
 (-1)^{m-1}\bigl(h(z)-R_{2m+1}(z)\bigr)>0.
\]
Thus $H_n-\log n-R_{2m+1}(n)$ gives alternating strict upper and lower
bounds for $\gamma$, converging to it with $m$ for every fixed positive
integer $n$.
\end{corollary}
\begin{proof}
The paired truncation is exactly
$R_{2m+1}(z)=x(1-xU_{m-1}(x^2))/2$, where $U_j$ is the
Stieltjes convergent through $\lambda_j$.  Positive Stieltjes fractions
have alternating strict upper and lower convergents at positive
arguments.  Moreover
$\mu_n^{-{1}/{(2n)}}\asymp n^{-1}$ by the displayed moment formula and
Stirling's formula.  Carleman's Stieltjes criterion~\cite{p3:Wall1948} therefore implies
moment determinacy, and the Stieltjes convergence theorem~\cite{p3:Stieltjes1894,p3:Wall1948} gives
$U_j(y)\to U(y)$ for every $y>0$.  These facts prove both assertions.
\end{proof}

\paragraph{Exact consistency check.}
Rational determinant arithmetic reproduces
\[
 (a_2,a_4,a_6,a_8,a_{10},a_{12})=
 \left(\frac16,\frac35,\frac{79}{126},\frac{7230}{6241},
 \frac{4146631}{3833346},\frac{306232774533}{179081182865}\right),
\]
the complete list in Xu--You's Theorem~1, and yields the stated $a_{14}$.
This calculation checks normalization; the proof for all orders uses
positivity and formal uniqueness, not a finite numerical test.


\paragraph{Prior-art distinction.}
Positive Stieltjes fractions for the digamma and polygamma remainders
predate Xu--You: see Shenton--Bowman~\cite{p3:ShentonBowman1971} (1971), DOI
\texttt{10.1137/0120057}, and Bonan-Hamada--Jones~\cite{p3:BonanJones2005} (2005), DOI
\texttt{10.1016/j.cam.2004.09.034}. In particular the latter paper's
Section~2 explicitly uses the positive moments and Carleman convergence
argument. The contribution here is the exact identification, via pair
contraction and formal uniqueness, with the particular optimized
coefficients in Xu--You's question. Neither the positive Stieltjes
representation nor its convergence is claimed as new.

% References to integrate:
% Xu, Hongmin; You, Xu. Continued fraction inequalities for the
% Euler--Mascheroni constant. J. Inequal. Appl. 2014,343.
% https://doi.org/10.1186/1029-242X-2014-343
% Binet digamma: NIST DLMF 5.9.15 plus psi(z+1)=psi(z)+1/z.
% T. J. Stieltjes, Recherches sur les fractions continues,
% Ann. Fac. Sci. Toulouse 8 (1894), J1--J122; 9 (1895), A5--A47.
% H. S. Wall, Analytic Theory of Continued Fractions, Van Nostrand,1948,
% classical positive Stieltjes fraction and moment theorem.


% Self-contained section fragment. Requires amsmath, amssymb, amsthm.
% Bibliography entries are supplied separately in moments_notes.md.
\section{A factorial Hankel conjecture and the boundary of moment positivity}
\label{p3:sec:factorial-hankel}

Ferguson~\cite[p.~6, final example of Section~2]{p3:Ferguson2020} defines
$a_0=0$, $a_k=(k-1)!$ for $k\ge1$, and conjectures the following identity
after checking it through index $50$:
\[
 -\frac{\det(a_{i+j})_{i,j=0}^{n+1}}
 {\det(a_{i+j+2})_{i,j=0}^{n}}=H_{n+1}.
\]
Here $H_N=\sum_{j=1}^N1/j$. We give an exact proof and a parameter extension.
The proof uses classical Laguerre orthogonality. Thus the result answers
the explicitly conjectural assertion in that source; no claim that this
determinant evaluation was previously inaccessible, or that no equivalent
evaluation occurs elsewhere, is needed.

\begin{theorem}[Bordered factorial Hankel determinants]
\label{p3:thm:factorial-border}
Let $\alpha>-1$ and $N\ge1$. Put
\[
 a_0=0,\qquad a_k=\Gamma(k+\alpha)\quad(k\ge1),
 \quad Q_N(\alpha)=\det(a_{i+j+2})_{i,j=0}^{N-1},
 \quad P_N(\alpha)=-\det(a_{i+j})_{i,j=0}^{N}.
\]
Then
\begin{align}
 Q_N(\alpha)&=\prod_{j=0}^{N-1}j!\,\Gamma(j+\alpha+2),
 \label{p3:eq:factorial-Q}\\
 R_N(\alpha):=\frac{P_N(\alpha)}{Q_N(\alpha)}
 &=\Gamma(\alpha+1)^2
       \sum_{j=0}^{N-1}\frac{j!}{\Gamma(j+\alpha+2)}.
 \label{p3:eq:factorial-R}
\end{align}
In particular $P_N(\alpha),Q_N(\alpha)>0$. If $\alpha\ne0$, then
\[
 R_N(\alpha)=\Gamma(\alpha)
 \left(1-\frac{\Gamma(\alpha+1)N!}{\Gamma(N+\alpha+1)}\right).
\]
At $\alpha=0$ the continuous value is $R_N(0)=H_N$, and hence
\[
 Q_N(0)=\prod_{j=0}^{N-1}j!(j+1)!,\qquad
 P_N(0)=H_N\prod_{j=0}^{N-1}j!(j+1)!.
\]
\end{theorem}

\begin{proof}
Let $M=(\Gamma(i+j+\alpha+2))_{i,j=0}^{N-1}$ and
$b=(\Gamma(j+\alpha+1))_{j=0}^{N-1}$. This is the Gram matrix of
$1,x,\ldots,x^{N-1}$ for the strictly positive weight
$x^{\alpha+1}e^{-x}$ on $(0,\infty)$, so it is positive definite.
The bordered determinant identity gives
\[
 P_N(\alpha)=\det M\,b^{\mathsf T}M^{-1}b.
\]
The monic Laguerre polynomials
$(-1)^j j!L_j^{(\alpha+1)}(x)$ have squared norms
$j!\Gamma(j+\alpha+2)$. Changing from monomials to this monic basis
has determinant one and gives~\eqref{p3:eq:factorial-Q}.

The vector $b$ represents the functional
$\ell(p)=\int_0^\infty p(x)x^\alpha e^{-x}\,dx$ on polynomials of
degree less than $N$. In the orthogonal basis $L_j^{(\alpha+1)}$,
\[
 b^{\mathsf T}M^{-1}b
 =\sum_{j=0}^{N-1}
 \frac{\ell(L_j^{(\alpha+1)})^2}
 {\Gamma(j+\alpha+2)/j!}.
\]
For all $j\ge0$,
\[
 \ell(L_j^{(\alpha+1)})=\Gamma(\alpha+1).
\]
For example, this follows by integrating the Laguerre generating
function: for $u$ sufficiently close to zero,
\[
 \int_0^\infty x^\alpha e^{-x}
   (1-u)^{-\alpha-2}
   \exp\left(-\frac{xu}{1-u}\right)dx
 =\frac{\Gamma(\alpha+1)}{1-u}.
\]
Coefficient comparison proves~\eqref{p3:eq:factorial-R}. These classical
Laguerre formulas may also be obtained directly from Rodrigues'
formula by integration by parts; see~\cite[Sections~18.3, 18.12]{p3:DLMF}.
Finally, for $\alpha\ne0$,
\[
 \frac{j!}{\Gamma(j+\alpha+2)}
 =\frac1\alpha\left(
 \frac{j!}{\Gamma(j+\alpha+1)}
 -\frac{(j+1)!}{\Gamma(j+\alpha+2)}\right).
\]
The sum telescopes. At $\alpha=0$ its summands are $1/(j+1)$.
\end{proof}

\begin{corollary}[Complete classification of the zeroth moment]
\label{p3:cor:factorial-completion}
Fix $\alpha>-1$ and prescribe the positive-index moments
$\mu_k=\Gamma(k+\alpha)$ for $k\ge1$. There is a finite positive
measure on $[0,\infty)$ with these moments and total mass $c$ if and
only if
\[
 \alpha>0\quad\hbox{and}\quad c\ge\Gamma(\alpha).
\]
When it exists it is unique and equals
\[
 d\mu(x)=x^{\alpha-1}e^{-x}\,dx+
                (c-\Gamma(\alpha))\,d\delta_0(x).
\]
Equivalently, all matrices $(\mu_{i+j})_{i,j=0}^{N}$, with
$\mu_0=c$, are positive semidefinite exactly under these conditions.
For each fixed $N\ge1$ this matrix has a negative eigenvalue exactly
when $c<R_N(\alpha)$; in that case it has precisely one.
\end{corollary}

\begin{proof}
Its lower-right block is $M$ from the proof above. Congruence by block
Gaussian elimination gives
\[
 (\mu_{i+j})_{i,j=0}^{N}\ \sim\
 \operatorname{diag}(c-R_N(\alpha),M).
\]
The claimed finite-dimensional positivity and inertia follow.
Formula~\eqref{p3:eq:factorial-R} is strictly increasing in $N$.
For $\alpha>0$ it tends to $\Gamma(\alpha)$. For $\alpha=0$ it is
$H_N\to\infty$. For $-1<\alpha<0$, its summands are asymptotic to
$\Gamma(\alpha+1)^2j^{-\alpha-1}$, so it again diverges to infinity.
This proves necessity, as well as the matrix formulation.

For $\alpha>0$ the displayed measure proves sufficiency. To verify
uniqueness, if $\mu$ is any such measure, the measure $x\,d\mu(x)$
has moments $\Gamma(k+\alpha+1)$, $k\ge0$. They determine that
measure uniquely. One direct verification uses an exponential
moment: for $0<t<1$, monotone convergence gives
\[
 \int e^{tx}x\,d\mu(x)
 =\sum_{k=0}^\infty\frac{t^k\Gamma(k+\alpha+1)}{k!}
 =\Gamma(\alpha+1)(1-t)^{-\alpha-1}.
\]
Uniqueness of the Laplace transform identifies $x\,d\mu$ with
$x^\alpha e^{-x}dx$. Division by $x$ on $(0,\infty)$ leaves only
an undetermined atom at zero, whose mass is fixed by $c$.
\end{proof}

In Ferguson's final example the functional has value $-\gamma$ at
$1$, and values $(k-1)!$ at $x^k$ for $k\ge1$. The exact value
$R_N(0)=H_N$ proves its asserted determinant identity and positivity,
and proves divergence of the proposed ratios. More strongly, no
finite replacement for the zeroth moment can turn those higher
moments into a positive moment functional. This explains the failure
structurally and does not imply anything about the arithmetic nature
of $\gamma$.

\subsection{The Gompertz determinant ratios are classical Laguerre approximants}

The same calculation identifies the Gompertz example of
\cite[Section~2]{p3:Ferguson2020}. This identification and its error
formula are classical orthogonal-polynomial consequences, included
to distinguish positive approximation from an irrationality proof.

\begin{proposition}[Exact denominator and exact positive error]
\label{p3:prop:gompertz-laguerre}
For $a>0$ put
\[
 F(a)=\int_0^\infty\frac{e^{-y}}{y+a}\,dy=e^aE_1(a),\qquad
 u_0=0,\quad
 u_k=\int_0^\infty(y+a)^{k-1}e^{-y}\,dy\quad(k\ge1).
\]
Let $Q_N(a)=\det(u_{i+j+2})_{i,j=0}^{N-1}$ and
$R_N(a)=-\det(u_{i+j})_{i,j=0}^{N}/Q_N(a)$. Then
\[
 Q_N(a)=N!L_N(-a)\prod_{j=0}^{N-1}(j!)^2
\]
and
\begin{align}
 F(a)-R_N(a)
 &=\frac1{L_N(-a)^2}
   \int_0^\infty\frac{e^{-y}L_N(y)^2}{y+a}\,dy\label{p3:eq:gompertz-square-error}\\
 &=\frac1{L_N(-a)}
   \int_0^\infty e^{-at}\frac{t^N}{(1+t)^{N+1}}\,dt>0.
 \label{p3:eq:gompertz-laplace-error}
\end{align}
Here $L_N=L_N^{(0)}$. In particular $a=1$ gives $F(1)=\delta$.
\end{proposition}

\begin{proof}
With $x=y+a$, use the positive measure
$d\nu(x)=e^{a-x}dx/x$ on $[a,\infty)$. The Schur complement
$F(a)-R_N(a)$ is the minimum of
\[
 \int_a^\infty(1-xq(x))^2\,d\nu(x),\qquad \deg q<N.
\]
The minimizing residual is orthogonal, in this measure, to
$x,x^2,\ldots,x^N$. Thus it is the degree-$N$ Laguerre polynomial
for the translated weight $e^{-y}$, normalized at $x=0$:
\[
 r(x)=\frac{L_N(x-a)}{L_N(-a)}.
\]
Since $L_N(-a)=\sum_{j=0}^N\binom Nj a^j/j!>0$, the normalization
is legitimate. This proves~\eqref{p3:eq:gompertz-square-error}.
The polynomial $(L_N(y)-L_N(-a))/(y+a)$ has degree $N-1$.
Its orthogonality to $L_N(y)$ for weight $e^{-y}$ gives
\[
 \int_0^\infty\frac{e^{-y}L_N(y)^2}{y+a}\,dy
 =L_N(-a)\int_0^\infty\frac{e^{-y}L_N(y)}{y+a}\,dy.
\]
Using $1/(y+a)=\int_0^\infty e^{-(y+a)t}dt$ and the finite
polynomial expansion of $L_N$ gives
\[
 \int_0^\infty e^{-(1+t)y}L_N(y)dy
 =\frac{t^N}{(1+t)^{N+1}},
\]
which proves~\eqref{p3:eq:gompertz-laplace-error}.

For the determinant formula, translation of the monomial basis has
determinant one, so $Q_N(a)$ is the size-$N$ Gram determinant for
$(y+a)e^{-y}$. The elementary Christoffel determinant identity is
\[
 \frac{\det\bigl(\int y^{i+j}(y+a)e^{-y}dy\bigr)_{i,j=0}^{N-1}}
 {\det\bigl(\int y^{i+j}e^{-y}dy\bigr)_{i,j=0}^{N-1}}
 =(-1)^Np_N(-a),
\]
where $p_N(y)=(-1)^NN!L_N(y)$ is monic. To see the identity,
expand the standard bordered Gram determinant for $p_N(-a)$
along its last row, or apply consecutive column operations to it.
The denominator is $\prod_{j=0}^{N-1}(j!)^2$, giving the result.
\end{proof}

\section{Why logarithmic Gompertz moments cannot inherit a finite recurrence}
\label{p3:sec:gompertz-nonholonomic}

Define
\[
 D_n(a)=\int_0^\infty e^{-t}\log^n(1+t/a)\,dt,
 \qquad a>0,\quad n\ge0.
\]
The logarithmic Gompertz moment sequence at $a=1$ is $D_n=D_n(1)$, with $D_0=1$ and
$D_1=\delta$. The substitution $x=\log(1+t/a)$ gives the positive
density
\[
 w_a(x)=a\exp(a+x-ae^x),\qquad x\ge0,
 \quad D_n(a)=\int_0^\infty x^nw_a(x)dx.
\]
Its moment positivity does not place it in the holonomic framework
of Franel or banana coefficient sequences.

\begin{theorem}[Nonexistence of polynomial-coefficient recurrences]
\label{p3:thm:gompertz-not-precursive}
For every fixed $a>0$, the sequence $(D_n(a))_{n\ge0}$ satisfies
no nonzero linear recurrence
\[
 \sum_{j=0}^r p_j(n)D_{n+j}(a)=0\qquad(n\ge n_0),
 \quad p_j\in\mathbb C[n].
\]
Equivalently its entire exponential generating function
\[
 M_a(z)=\sum_{n\ge0}D_n(a)\frac{z^n}{n!}
       =e^a a^{-z}\Gamma(z+1,a)
\]
is not D-finite over $\mathbb C(z)$. Nevertheless it obeys the
first-order inhomogeneous difference equation
\[
 aM_a(z+1)=(z+1)M_a(z)+a.
\]
\end{theorem}

\begin{proof}
The decay of $w_a(x)$ is faster than every exponential, giving
locally uniform convergence of its Laplace transform and all its
derivatives. This proves the entire generating-function formula.
Integration by parts in the incomplete gamma integral proves the
difference equation.

First, $w_a$ is not annihilated on $(0,\infty)$ by a nonzero
polynomial-coefficient differential operator. Indeed
\[
 \frac{w_a^{(j)}(x)}{w_a(x)}=B_j(e^x),
 \quad B_0(X)=1,\quad
 B_{j+1}(X)=XB_j'(X)+(1-aX)B_j(X).
\]
The polynomial $B_j$ has degree $j$ and leading coefficient
$(-a)^j$. If $\sum_{j=0}^rq_j(x)w_a^{(j)}(x)=0$ with
$q_r\ne0$, division by $w_a(x)e^{rx}$ gives
\[
 (-a)^rq_r(x)+\sum_{j<r}v_j(x)e^{-(r-j)x}=0
\]
for certain polynomials $v_j$. The sum tends to zero exponentially
times a polynomial as $x\to+\infty$, whereas a nonzero polynomial
$q_r$ cannot do so. This is a contradiction.

Now suppose $M_a$ satisfies
$\sum_{j=0}^s p_j(z)\partial_z^jM_a(z)=0$, not all $p_j$ zero.
Expanding the $p_j$ into powers of $z$ and integrating by parts
repeatedly gives
\[
 \int_0^\infty e^{zx}\,h(x)dx=P(z),\qquad
 h(x)=\sum_{j=0}^s p_j(-\partial_x)(x^jw_a(x)),
\]
where $P$ is a polynomial consisting of endpoint terms at zero.
All endpoint terms at infinity vanish, and $h$ has faster than
exponential decay. As real $z\to-\infty$ the integral tends to zero;
therefore $P=0$. Uniqueness of the Laplace transform gives $h=0$.
The operator
$\sum_j p_j(-\partial_x)x^j$ is nonzero: the substitution
$z\mapsto-\partial_x$, $\partial_z\mapsto x$ is an isomorphism of
Weyl algebras, since $[x,-\partial_x]=1$.
Thus $h=0$ contradicts the preceding paragraph.

Finally, a sequence is P-recursive if and only if its exponential
generating function is D-finite. In the direction needed here,
multiply an eventual recurrence by $z^n/n!$ and sum; the result
is a polynomial-coefficient inhomogeneous differential equation,
whose polynomial right-hand side is removed by sufficiently many
derivatives. The correspondence is injective at the operator level,
or follows directly by comparing coefficients. This contradiction
proves the recurrence assertion.
\end{proof}

\begin{remark}[Scope]
The last theorem is an obstruction to transporting a finite-order
recurrence argument to the raw logarithmic moments. It does not
exclude recurrences after introducing additional parameters, nonlinear
recurrences, or a different family of rational approximants. It gives
neither irrationality of $D_1=\delta$ nor algebraic independence of
the constants $D_n$. The differential transcendence theorem for the
two-variable incomplete gamma function in~\cite{p3:Arreche2013} concerns
a different statement; specializing a parameter in that theorem
would not by itself justify the result above.
\end{remark}

\subsection{A stronger obstruction from differential difference Galois theory}

The preceding elementary non-D-finiteness result can be strengthened.
The additional input is the theorem of Bostan--Di Vizio--Raschel
\cite[Theorem~3 and Corollary~40]{p3:BDR2024}: a nonrational formal
Laurent solution of a first-order inhomogeneous equation with rational
coefficients and a nontrivial parabolic M\"obius substitution is
differentially transcendental over the field of convergent meromorphic
germs at its fixed point. We apply that existing theorem to an
asymptotic expansion. This is a consequence of their theorem, not a
new solution of a conjecture posed in that paper.

\begin{theorem}[Differential transcendence of the fixed-cutoff moment function]
\label{p3:thm:gompertz-dtrans}
For every $a>0$, $M_a(z)=e^aa^{-z}\Gamma(z+1,a)$ is
differentially transcendental over $\mathbb C(z)$. More strongly, no
nonzero differential polynomial with coefficients meromorphic at
$z=\infty$ vanishes on this function. For each fixed $a>0$, both
parameter functions $s\mapsto\Gamma(s,a)$ and
$s\mapsto\gamma(s,a)$ are differentially transcendental over
$\mathbb C(s)$, where $\gamma(s,a)$ here denotes the lower incomplete
gamma function.
\end{theorem}

\begin{proof}
Let $w_a(x)=a\exp(a+x-ae^x)$ as above and put $g_a(t)=M_a(-1/t)$
for $t>0$. Repeated integration by parts in its Laplace integral gives
the full asymptotic expansion
\[
 g_a(t)\sim\widehat g_a(t):=
       \sum_{n=0}^\infty w_a^{(n)}(0)t^{n+1}\quad(t\downarrow0).
\]
This expansion is stable under differentiation. For completeness,
write $M_a^{(k)}(-s)=\int_0^\infty x^ke^{-sx}w_a(x)dx$.
Expand $x^kw_a(x)$ to an arbitrary finite order near zero and split
the integral at a fixed positive point. The Taylor remainder on the
first interval contributes the corresponding next inverse power of
$s$; the second interval is exponentially small. The chain rule
with $s=1/t$ then gives termwise differentiated expansions to every
fixed order. The same argument works uniformly in smaller right-half
sectors of $t$ if desired.

The exact difference equation for $M_a$ implies the formal identity
\begin{equation}\label{p3:eq:gompertz-parabolic}
 \widehat g_a\left(\frac{t}{1-t}\right)
   =\frac{t-1}{at}\widehat g_a(t)+1.
\end{equation}
The formal series $\widehat g_a$ is not rational. Otherwise
$R(z)=\widehat g_a(-1/z)$ would be a rational solution of
\[
 aR(z+1)=(z+1)R(z)+a.
\]
If $R$ has a finite pole $\beta$ of minimal real part, the left
side has a pole at $\beta-1$, while the right side is regular there,
a contradiction. Thus $R$ is a polynomial. A nonzero polynomial of
degree $d$ gives degrees $d$ and $d+1$ on the two sides, which is
impossible; the zero polynomial also fails since $a\ne0$.

The substitution $t\mapsto t/(1-t)$ is nontrivial and has its unique
fixed point at zero. Therefore the cited theorem applied to
\eqref{p3:eq:gompertz-parabolic} says that $\widehat g_a$ is
differentially transcendental over $\mathbb C(\{t\})$.
If $g_a$ satisfied an algebraic differential equation over that field,
its derivative-stable asymptotic expansion would satisfy the same
formal equation: asymptotic expansions respect finite sums, products,
and multiplication by convergent Laurent series. A zero function has
zero formal expansion. This contradicts the cited theorem.
The invertible change $z=-1/t$ proves the assertion for $M_a$,
including the stronger coefficient field. Since multiplication by
$e^aa^{-z}$ and translation of the variable preserve differential
algebraicity over rational functions, it also proves the claim for
the upper incomplete gamma function.

For the lower incomplete gamma function set
\[
 N_a(z)=e^aa^{-z}\gamma(z+1,a)
       =\int_0^\infty e^{-zx}w_a(-x)dx
       \qquad(\Re z>-1).
\]
The equality follows by the substitution $u=ae^{-x}$ in its defining
integral. Its derivative-stable expansion at $z=1/t\to+\infty$ is
\[
 N_a(1/t)\sim\sum_{n\ge0}(-1)^nw_a^{(n)}(0)t^{n+1}
            =-\widehat g_a(-t).
\]
Replacing $t$ by $-t$ and multiplying by $-1$ preserve differential
transcendence over $\mathbb C(\{t\})$. The same asymptotic argument
therefore proves differential transcendence of $N_a$ and hence of
$s\mapsto\gamma(s,a)$ over $\mathbb C(s)$.
\end{proof}

The derivatives $w_a^{(n)}(0)/a$ are precisely the coefficients
defined by the exponential generating function
$\exp(x-a(e^x-1))$, a specialization of the actuarial polynomials
already treated in~\cite[Table~2 and Corollary~4]{p3:BDR2024}.
This identifies the prior mechanism behind the stronger theorem.
It remains a statement about functions and does not establish
irrationality of any individual logarithmic moment.


\section{Rational intersections for exponential-integral values}
\label{p3:ef:section}

Let $\mathcal E$ be the ring of values of $E$-functions at algebraic
points, and write
\[
 \operatorname{Ein}(z)=\sum_{n\ge1}\frac{(-1)^{n-1}z^n}{n\,n!}.
\]
The function $\operatorname{Ein}$ is an $E$-function.  Rivoal's
published Problem~6 \cite[p.~286]{p3:RivoalXUPS2019}, restated as
the units conjecture of Fischler--Rivoal
\cite[Conjecture~1.2]{p3:FischlerRivoal2026Units},
asserts that $\mathcal E^\times$ consists exactly of the numbers
$\alpha e^\beta$, with $\alpha,\beta\in\overline{\mathbb Q}$ and
$\alpha\ne0$.  The following gives a complete intersection theorem
for a family of subfields of $\mathbb C$.  In particular, it settles
the unit question inside each indicated polynomial subalgebra.
It does not settle the full units conjecture.

\begin{theorem}[Polynomial intersection off the positive real axis]
\label{p3:ef:intersection}
For every $a\in\overline{\mathbb Q}\setminus[0,+\infty)$, put
$\eta_a=\operatorname{Ein}(a)$.  Then
\begin{equation}\label{p3:ef:intersection-equation}
 \overline{\mathbb Q}(\eta_a)\cap\mathcal E
 =\overline{\mathbb Q}[\eta_a].
\end{equation}
Consequently
\begin{equation}\label{p3:ef:units}
 \overline{\mathbb Q}[\eta_a]\cap\mathcal E^\times
 =\overline{\mathbb Q}^{\times}.
\end{equation}
Equivalently, for every nonconstant polynomial
$P\in\overline{\mathbb Q}[Y]$, the nonzero value $P(\eta_a)$ has
no reciprocal in $\mathcal E$.
\end{theorem}

The proof uses the injective ring morphisms
\[
 \Phi_\theta:\mathcal E\longrightarrow
 \left\{\sum_{j=1}^N b_j e^{\rho_jX}:
       b_j\in\mathbb C,\ \rho_j\in\overline{\mathbb Q}\right\},
 \qquad -\pi/2<\theta<\pi/2,
\]
constructed in \cite[Proposition~3.1]{p3:FischlerRivoal2026Units}.
They fix algebraic numbers and satisfy
$\Phi_\theta(\xi)(1)=\xi$.  If the canonical large-sector decomposition
of an $E$-function is $f(z)=\sum_\rho f_\rho(z)e^{\rho z}$,
then $\Phi_\theta(f(1))=\sum_\rho f_\rho(1)e^{\rho X}$.
The stated range of directions is essential.

\begin{lemma}[The two phases and their nonvanishing]
\label{p3:ef:phase}
For every $a\in\overline{\mathbb Q}\setminus[0,+\infty)$, an admissible
direction $\theta$ gives
\begin{equation}\label{p3:ef:phase-equation}
 \Phi_\theta(\eta_a)=C_a+B_a e^{-aX},
 \qquad C_a\notin\overline{\mathbb Q},\quad B_a\ne0.
\end{equation}
If $a$ is nonreal, one can take $\theta=0$ and
\begin{equation}\label{p3:ef:phase-coefficients}
 C_a=\gamma+\operatorname{Log}a,
 \qquad
 B_a=\int_0^\infty\frac{e^{-t}}{a+t}\,dt=e^aE_1(a),
\end{equation}
where $\operatorname{Log}$ denotes the principal branch.
If $a=-x$ with $x>0$, one can choose $\theta>0$ sufficiently small and
take
\begin{equation}\label{p3:ef:negative-phase}
 C_{-x}=\gamma+\log x-i\pi,
 \qquad
 B_{-x}=e^{-x}\bigl(\operatorname{Ein}(-x)-\gamma-\log x+i\pi\bigr).
\end{equation}
\end{lemma}
\begin{proof}
On the slit plane the classical identity \cite[\S6.2]{p3:DLMF}
\begin{equation}\label{p3:ef:ein-identity}
 \operatorname{Ein}(w)=\gamma+\operatorname{Log}w+E_1(w)
\end{equation}
follows either by differentiating both sides and comparing their
limits at zero, or directly from their convergent expansions.
For nonreal $a$, the zero-direction Borel sum of
\[
 \widehat B_a(z)=\sum_{n\ge0}
       \frac{(-1)^n n!}{a^{n+1}z^{n+1}}
\]
is
\[
 B_a(z)=\int_0^\infty\frac{e^{-z\xi}}{a+\xi}\,d\xi.
\]
Indeed its Borel transform is $1/(a+\xi)$, whose pole is off the
positive integration ray.  Small rotations of that ray provide the
canonical sum in a sector of opening greater than $\pi$ about the
positive $z$-axis.  The identity
\[
 e^{-az}B_a(z)=E_1(az)
\]
first holds where the principal integral representation applies and
then by analytic continuation with the branch fixed at $z=1$.
Thus \eqref{p3:ef:ein-identity} supplies the canonical phase-zero term
$\gamma+\operatorname{Log}a+\log z$ and phase-$(-a)$ term $B_a(z)$.
This proves \eqref{p3:ef:phase-coefficients} and the asserted phase map.
Moreover
\[
 \operatorname{Im}B_a
 =-\operatorname{Im}(a)
    \int_0^\infty\frac{e^{-t}}{|a+t|^2}\,dt\ne0.
\]

For $a=-x<0$ choose a small positive summation direction $\theta$;
the corresponding Laplace contour passes below the pole $\xi=x$.
The lower lateral value of $\operatorname{Log}(-x)$ is $\log x-i\pi$,
and \eqref{p3:ef:ein-identity} gives \eqref{p3:ef:negative-phase}.
These are again the canonical two phases, obtained by the same Borel
sum with its contour rotated below the pole.  Here
$\operatorname{Im}B_{-x}=\pi e^{-x}\ne0$.

It remains to prove transcendence of $C_a$.  In the nonreal case put
$\varphi=\arg(a)\in(-\pi,\pi)\setminus\{0\}$.  If $C_a$ were
algebraic, then
\[
 \varphi=\frac{C_a-\overline{C_a}}{2i}
\]
would be a nonzero algebraic real number.  But
$e^{i\varphi}=a/|a|$ is algebraic, contradicting the
Hermite--Lindemann theorem.  We used here that
$|a|=\sqrt{a\overline a}$ is algebraic and that $\gamma$ is real.
On the negative ray, algebraicity of $C_{-x}$ would similarly imply
algebraicity of $\pi$, again impossible.  No arithmetic assumption on
$\gamma$ was used.
\end{proof}

\begin{lemma}[One-phase-variable denominator test]
\label{p3:ef:denominator}
Suppose $\eta\in\mathcal E$ and an admissible phase morphism satisfies
\[
 \Phi_\theta(\eta)=C+B e^{\lambda X},
 \qquad B\ne0,\quad 0\ne\lambda\in\overline{\mathbb Q}.
\]
If coprime $P,Q\in\overline{\mathbb Q}[Y]$, with $P,Q\ne0$, satisfy
$P(\eta)/Q(\eta)\in\mathcal E$, then every complex root of $Q$
equals $C$.  In particular, if $C$ is transcendental, $Q$ is constant.
\end{lemma}
\begin{proof}
First $\eta$ is transcendental: if it were algebraic, its phase image
would be constant, contrary to $B\lambda\ne0$.
Write $H(X)=\Phi_\theta(P(\eta)/Q(\eta))$.  The ring-homomorphism
property gives, as an identity of meromorphic functions,
\[
 H(X)=\frac{P(C+B e^{\lambda X})}{Q(C+B e^{\lambda X})}.
\]
The function $H$ is entire.  Suppose $r$ is a root of $Q$ with
$r\ne C$.  Since the exponential map is onto $\mathbb C^\times$,
there is $X_0\in\mathbb C$ with
$e^{\lambda X_0}=(r-C)/B$.  Coprimality gives $P(r)\ne0$.
The right-hand side therefore has a pole at $X_0$, contradicting
entireness of $H$.  Thus every root of $Q$ equals $C$.
All roots of $Q\in\overline{\mathbb Q}[Y]$ are algebraic, proving the
last assertion.
\end{proof}

\begin{proof}[Proof of Theorem~\ref{p3:ef:intersection}]
Apply Lemma~\ref{p3:ef:denominator} to Lemma~\ref{p3:ef:phase}.
Every element of the field on the left of
\eqref{p3:ef:intersection-equation} has a reduced rational-function
representation in $\eta_a$, and its denominator is constant.
The reverse inclusion follows because $\mathcal E$ is a ring
containing $\overline{\mathbb Q}$ and $\eta_a$.
If $P(\eta_a)$ were a unit for a nonconstant $P$, its reciprocal
would lie in the field intersection and hence be a polynomial in
$\eta_a$.  The transcendence of $\eta_a$ would then give a polynomial
identity $P(Y)R(Y)=1$, which is impossible.  Nonzero algebraic numbers
are units, proving \eqref{p3:ef:units}.
\end{proof}

\begin{lemma}[Surjective phase images]
\label{p3:ef:surjective-phase}
Suppose $\xi\in\mathcal E$ and, for some admissible $\theta$, the
exponential polynomial $H=\Phi_\theta(\xi)$ has at least two distinct
nonzero phases with nonzero coefficients.  Then
\begin{equation}\label{p3:ef:surjective-intersection}
 \overline{\mathbb Q}(\xi)\cap\mathcal E
 =\overline{\mathbb Q}[\xi],\qquad
 \overline{\mathbb Q}[\xi]\cap\mathcal E^\times
 =\overline{\mathbb Q}^\times.
\end{equation}
\end{lemma}
\begin{proof}
We first show that $H$ takes every complex value.  If it omitted
$w\in\mathbb C$, the entire function $h=H-w$ would have no zeros.
It has a bound $|h(z)|\le A e^{B|z|}$ with $A,B>0$.
The zero-free case of Hadamard factorization then gives
$h(z)=e^{uz+v}$ for constants $u,v$.
Here this case can also be seen directly: write $h=e^g$ with $g$
entire, so $\operatorname{Re}g(z)\le\log A+B|z|$.
The Borel--Carath\'eodory inequality on disks of radii $R$ and $2R$
bounds $\max_{|z|\le R}|g(z)-g(0)|$ by $O(1+R)$.
Cauchy's inequalities, followed by $R\to\infty$, give
$g^{(n)}(0)=0$ for every $n\ge2$, proving that $g$ is affine.
But subtracting $w$ only changes the constant phase of $H$ and leaves
its two distinct nonzero phases.  Linear independence of distinct
exponential functions rules out $H-w=e^{uz+v}$.  Thus $H$ is surjective.

In particular $H$ is nonconstant, so $\xi$ is transcendental.
Let $P,Q\in\overline{\mathbb Q}[Y]$ be coprime and suppose
$R(\xi)=P(\xi)/Q(\xi)\in\mathcal E$.
Applying $\Phi_\theta$ gives the entire-function identity
\[
 Q(H(X))\Phi_\theta(R(\xi))(X)=P(H(X)).
\]
If $Q$ were nonconstant, choose one of its roots $r$ and then $X_0$
with $H(X_0)=r$.  The identity gives $P(r)=0$, contradicting
coprimality.  Hence $Q$ is constant, proving the field intersection.
The assertion about units follows exactly as in
Theorem~\ref{p3:ef:intersection}.
\end{proof}

\begin{theorem}[All finite combinations with at least two arguments]
\label{p3:ef:multiple-arguments}
Let $r\ge2$, let $a_1,\ldots,a_r$ be pairwise distinct nonzero
algebraic numbers, and let $c_1,\ldots,c_r$ be nonzero algebraic
numbers.  Put
\[
 \xi=\sum_{j=1}^r c_j\operatorname{Ein}(a_j).
\]
Then both conclusions in \eqref{p3:ef:surjective-intersection} hold.
In particular the result applies when all the arguments are positive.
\end{theorem}
\begin{proof}
Choose $\theta>0$ sufficiently small.  There are only finitely many
arguments $a_j$.  For every nonreal argument the interval from $0$ to
$\theta$ can therefore be chosen to contain no singular Borel-summation direction
of its tail with Borel transform $1/(a_j+\zeta)$; the phase coefficients
remain the principal ones in \eqref{p3:ef:phase-coefficients}.  For a
positive argument the same formula holds, and its integral $B_{a_j}$
is strictly positive.  For a negative argument the chosen positive
direction gives the lateral coefficients in \eqref{p3:ef:negative-phase}.
All tail coefficients are nonzero by Lemma~\ref{p3:ef:phase} and
positivity in the positive case.  Thus the one common morphism gives
\[
 \Phi_\theta(\xi)
 =\sum_{j=1}^r c_jC_{a_j}
   +\sum_{j=1}^r c_jB_{a_j}e^{-a_jX}.
\]
The $r$ nonzero phases $-a_j$ are distinct and their coefficients
$c_jB_{a_j}$ do not vanish.  Apply Lemma~\ref{p3:ef:surjective-phase}.
\end{proof}

For real examples, each of
\[
 \xi_+=\operatorname{Ein}(1)+\operatorname{Ein}(2),\qquad
 \xi_-=\operatorname{Ein}(1)-\operatorname{Ein}(2)
\]
has the full polynomial-intersection property.  The first equals
$2\gamma+\log2+E_1(1)+E_1(2)$ and retains Euler's constant explicitly;
the second cancels its coefficient.  These two cases are instances of
the same family theorem.

Theorem~\ref{p3:ef:intersection} also excludes
\[
 \frac1{\operatorname{Ein}(-1)-\alpha},\qquad
 \frac1{1+\operatorname{Ein}(i)^2}
 \quad(\alpha\in\overline{\mathbb Q})
\]
from $\mathcal E$, with no restriction on the differential order of a
possible representing $E$-function.  These examples are consequences
of one theorem, not separate resolutions of published conjectures.

\paragraph{What remains open at a positive argument.}
For positive algebraic $a$ the same phase computation gives
$C_a=\gamma+\log a$ and $B_a>0$.  Lemma~\ref{p3:ef:denominator} still
shows that a reduced denominator must be a constant multiple of
$(Y-C_a)^m$ if it is nonconstant.  The argument above cannot exclude
this possibility because $C_a$ is real.  At $a=1$ this is precisely
the exception controlled by the unknown arithmetic nature of
$\gamma$; the off-axis theorem does not resolve it.  Nor does it
decide the irrationality of the Euler--Gompertz constant.

% Bibliography input for integration:
% \bibitem{FischlerRivoal2026Units}
% S. Fischler and T. Rivoal, The ring of values of E-functions is not
% a field, arXiv:2610.16603v1, 12 October 2026 (manuscript dated
% 14 October 2026), especially Conjecture 1.2 and Proposition 3.1.
% Primary URL: https://arxiv.org/pdf/2610.16603
% \bibitem{RivoalXUPS2019} T. Rivoal, Les E-fonctions et G-fonctions
% de Siegel, Journees mathematiques X-UPS (2019), 197--298,
% doi:10.5802/xups.2019-03; Problem 6, p. 286.
% \bibitem{DLMF} NIST Digital Library of Mathematical Functions,
% https://dlmf.nist.gov/6.2 and https://dlmf.nist.gov/6.4.


\section{Unresolved targets and the arithmetic barrier}
\label{p3:sec:unresolved}

The results above give exact statements at the interfaces between
recurrences, analysis, and arithmetic. They do not supply the missing
arithmetic estimate needed to prove irrationality of $\gamma$.
This section records where the arguments stop, so that later work can
target a precise missing statement.

\begin{longtable}{@{}>{\raggedright\arraybackslash}p{0.35\textwidth}>{\raggedright\arraybackslash}p{0.57\textwidth}@{}}
\toprule
Target & Status and missing step \\
\midrule
\endhead
General scalar CMF convergence, all $r\ge3$, in~\cite[Conjecture 1]{p3:KMLWSK26}
& Proved here on every unequal rank-three repeated-minimum family.
Higher-rank scalar convergence is not established. The full-state assertion
is already false, so it cannot be used as the general intermediate claim.\\[6pt]
Equal lower increments $(\beta,\beta,\beta)$ with arbitrary nonpositive integral upper increments
& The source covers $\beta=1$. When both upper increments are divisible by
$\beta$, positive integral dilation follows by taking a subsequence of that
source trajectory. The present $d>0$ concentration proof does not address
the remaining equal-increment cases.\\[6pt]
Variable-argument CMF convergence and its full relation to Rivoal's parameter family~\cite[Section 9.2]{p3:KMLWSK26}
& The exact Meijer identity is proved. A matching identity for functions
does not prove a dual dynamical limit for every parameter or identify all
gauged recurrences.\\[6pt]
Order-two approximation construction for $\gamma$~\cite[Section 1]{p3:Riv09}
& The given order-three operator has no lower-order factor for nonzero
parameter. Existence of a different order-two construction remains open
here.\\[6pt]
Reduced-fraction height and the empirical approximation parameter in~\cite[Section 8.3]{p3:KMLWSK26}
& The factorial bound concerns a common denominator of a specified raw
vector. Cancellation in the scalar numerator and denominator has not been
controlled.\\[6pt]
Full $E$-value unit conjecture~\cite[Problem 6]{p3:RivoalXUPS2019}
& Polynomial intersections are proved in the stated single-generator
subfields. They do not classify all units of $\mathcal E$.\\[6pt]
Single positive argument $\operatorname{Ein}(a)$
& A possible denominator supported at $\gamma+\log a$ remains. This is an
explicit obstruction in the method, rather than a proved exceptional unit.\\[6pt]
Irrationality of $\gamma$
& Not proved. Neither convergence to $\gamma$ nor the functional Galois
calculations produce nonzero small integer linear forms with the requisite
denominator estimates.\\[6pt]
Transcendence of $\gamma$
& Not proved. Independence of functions over $\mathbb C(z)$ cannot be
specialized to arithmetic independence of their values without an
applicable specialization theorem.\\[6pt]
Irrationality of $\delta$
& Not proved. Positive Laguerre error formulas and differential
transcendence of the moment function do not decide this numerical question.\\
\bottomrule
\end{longtable}

This is a ledger of mathematical targets and obstructions. It is not a
list of ten solved published problems: several entries are components of
the same external conjecture, and others are questions exposed by the
present proofs. Counting individual parameters, signs, or corollaries as
independent resolutions would misstate the result.

\subsection{A concrete continuation of the recurrence method}

The scalar proof suggests an appropriate replacement for full-state
isolation. First identify the subspace killed by both extraction forms;
then compare the dual components in the quotient by that subspace. For
the family treated here this subspace contains $u=(1,-1,0)^T$, and the
relevant quotient coefficients are the two signed cofactors involving
the elementary row. Their dominant finite sums are positive before the
fixed seed is inserted. Concentration then controls the inserted
polynomial uniformly and proves eventual nonvanishing of the denominator.

In higher rank the analogous calculation would require control of several
cofactor minors, not merely a single positive weight. The exact contiguity
operators still make those minors accessible, but positivity and dominant
component cancellation must be proved rather than inferred from numerical
ratios. This is a concrete extension problem left by the manuscript.

For arithmetic progress one must add a second ingredient. If rational
approximants are written in lowest terms as $P_n/Q_n$, a route to
irrationality needs a nonzero error with sufficient decay after multiplying
by the actual reduced denominator. The normalization-dependent bounds in
Section~\ref{p3:sec:gamma-arithmetic} do not provide that ingredient. The
positive-moment and phase-morphism results likewise identify useful
structure without crossing this final arithmetic barrier.

\PartAppendix
\section{Exact checks and reproducibility}
\label{p3:sec:checks}

All limiting and nonexistence statements in the manuscript have written
proofs. Finite computations were used to check signs, indices, and source
normalizations. They are not substitutes for those proofs. The first
program below uses exact rational arithmetic to check factorial-Hankel
identities, translated Laguerre denominators, and the continued-fraction
coefficients. The second uses SymPy only to clear polynomial and rational
identities in the recurrence and gauge. Both are embedded here so that the
single TeX source contains the reported checks without auxiliary files.

Internal proof reviews independently rederived the full-state counterexample,
the scalar concentration argument, the recurrence irreducibility and height
obstructions, the Meijer contour shift, the $E$-value intersections, and the
use of the exact Bostan--Di Vizio--Raschel hypotheses. Such reviews are
fallible and do not amount to external refereeing or formal proof verification.



\subsection{Moment determinants and continued-fraction coefficients}
\begin{lstlisting}[language=Python]
"""Exact normalization checks; standard-library rational arithmetic only."""
from fractions import Fraction as F
from math import comb, factorial as fact


def determinant(rows):
    a = [[F(x) for x in row] for row in rows]
    result, size = F(1), len(a)
    for j in range(size):
        pivot = next((k for k in range(j, size) if a[k][j]), None)
        if pivot is None:
            return F(0)
        if pivot != j:
            a[j], a[pivot] = a[pivot], a[j]
            result = -result
        result *= a[j][j]
        for k in range(j + 1, size):
            ratio = a[k][j] / a[j][j]
            for ell in range(j + 1, size):
                a[k][ell] -= ratio * a[j][ell]
    return result


def product(values):
    result = F(1)
    for value in values:
        result *= value
    return result


def rising(a, k):
    return product(a + j for j in range(k))


for alpha in (F(-1, 2), F(0), F(1, 2), F(1), F(2), F(3)):
    for n in range(1, 7):
        # Divide each positive-index moment by Gamma(alpha+1).
        v = [F(0)] + [rising(alpha + 1, k - 1)
                      for k in range(1, 2 * n + 1)]
        q = determinant([[v[i+j+2] for j in range(n)] for i in range(n)])
        p = -determinant([[v[i+j] for j in range(n+1)] for i in range(n+1)])
        assert q == product(fact(j)*rising(alpha+1, j+1) for j in range(n))
        assert p / q == sum(F(fact(j))/rising(alpha+1, j+1) for j in range(n))
print("PASS: 36 normalized factorial-Hankel cases")

for a in (F(1, 2), F(1), F(2), F(3)):
    for n in range(1, 7):
        v = [F(0)] + [sum(comb(k-1, j)*a**(k-1-j)*fact(j)
                           for j in range(k)) for k in range(1, 2*n+1)]
        q = determinant([[v[i+j+2] for j in range(n)] for i in range(n)])
        laguerre = sum(F(comb(n, j), fact(j))*a**j for j in range(n+1))
        assert q == fact(n)*laguerre*product(fact(j)**2 for j in range(n))
print("PASS: 24 Gompertz denominator cases")

bernoulli = [F(1)]
for m in range(1, 15):
    bernoulli.append(-sum(comb(m+1, k)*bernoulli[k]
                          for k in range(m)) / (m+1))
mu = [abs(bernoulli[2*j+2]) / (j+1) for j in range(7)]


def hankel(n, shift):
    return determinant([[mu[i+j+shift] for j in range(n)] for i in range(n)])


c = [mu[0]]
for j in range(1, 7):
    r = (j+1)//2
    if j % 2:
        lam = hankel(r, 1)*hankel(r-1, 0)/(hankel(r, 0)*hankel(r-1, 1))
    else:
        lam = hankel(r+1, 0)*hankel(r-1, 1)/(hankel(r, 1)*hankel(r, 0))
    c.append(lam / c[-1])
expected = [F(1, 6), F(3, 5), F(79, 126), F(7230, 6241),
            F(4146631, 3833346), F(306232774533, 179081182865),
            F(3754087889491759, 2440840521848406)]
assert c == expected
# Independently extract the normalized continued fraction of 1-x U(x*x).
tail = [F(1)] + [F(0)] * 13
for j, moment in enumerate(mu):
    tail[2*j+1] = (-1)**(j+1)*moment
coefficients = []
while len(tail) > 1:
    inverse = [F(1)]
    for k in range(1, len(tail)):
        inverse.append(-sum(tail[j]*inverse[k-j] for j in range(1, k+1)))
    coefficient = inverse[1]
    coefficients.append(coefficient)
    tail = [v/coefficient for v in inverse[1:]]
assert coefficients == [value for q in c for value in (q, -q)][:13]
print("PASS: Xu--You a2 through a14 (13 coefficients, two constructions)")
print("a14 =", coefficients[-1])
\end{lstlisting}


\subsection{Recurrence layers, gauge, and Casoratian}
\begin{lstlisting}[language=Python]
import sympy as s

n, x, t = s.symbols('n x t')
a = [
    (n+2)**2*(3*x-8*n-22)*(3*x-8*n-30),
    (24*n**3+(162-57*x)*n**2+(26*x**2-258*x+356)*n
     -3*x**3+59*x**2-285*x+254)*(3*x-8*n-30),
    -(24*n**3+(7*x+210)*n**2+(-6*x**2+56*x+600)*n
      -15*x**2+96*x+564)*(3*x-8*n-14),
    -(n+3)**2*(3*x-8*n-14)*(3*x-8*n-22),
]
alpha = [s.expand(v).coeff(n, 4) for v in a]
beta = [s.expand(v).coeff(n, 3) for v in a]
assert alpha == [64, -192, 192, -64]
assert s.expand(sum(beta)) == 512*x
assert s.expand(sum(j*alpha[j] for j in range(4))) == 0
adjoint_beta = [s.expand(a[j].subs(n, n+3-j)).coeff(n, 3)
                for j in range(4)]
assert s.expand(sum(adjoint_beta)) == 512*x
assert s.expand(sum(alpha[j]*t**(3-j) for j in range(4))) == s.expand(64*(t-1)**3)
assert s.expand(sum(a[j].subs(x, 0)*(n+j+1) for j in range(4))) == 0

A = s.Matrix([
    [6*n**2+13*n+10, -8*n-5, 3],
    [3*n**3+3*n**2-3*n-6, -3*n**2+3*n+3, n-2],
    [n**4-2*n**3-6*n**2-3*n+3, 6*n**2+2*n-1, 1-2*n],
])/(n+1)**4
G = s.Matrix([
    [0, (8*n+11)/(n+1)**2, (24*n**2+97*n+94)/((n+1)**2*(n+2))],
    [-1, (3*n**2-n-7)/(n+1)**2,
     (15*n**4+69*n**3+66*n**2-78*n-111)/((n+1)**2*(n+2)**2)],
    [2, (-6*n**2-6*n+3)/(n+1)**2,
     (8*n**5+21*n**4-38*n**3-131*n**2-48*n+53)/((n+1)**2*(n+2)**2)],
])
C = s.zeros(3)
C[1, 0] = C[2, 1] = 1
for j in range(3):
    C[j, 2] = s.cancel(-a[j].subs(x, 1)/a[3].subs(x, 1))
assert all(s.cancel(v) == 0 for v in G*C-(n+1)**2*A*G.subs(n, n+1))
assert s.cancel(G.det()-(8*n+11)*(8*n+19)/(n+2)**2) == 0
assert s.cancel(A.det()-(n+1)**(-6)) == 0
cas = 4*(8*n+11)*(8*n+19)/(209*(n+2)**2)
assert s.cancel(cas.subs(n, n+1)/cas-C.det()) == 0
assert cas.subs(n, 0) == 1
print('PASS: leading layers, formal adjoint, OGF leading coefficient, gauge, determinants, Casoratian.')
\end{lstlisting}


\paragraph{Recorded exact outcomes.}
The moment program passed all 36 normalized factorial-Hankel cases, all 24 Gompertz denominator cases, and all 13 coefficients $a_2,\ldots,a_{14}$ by two constructions. The symbolic program passed the leading-layer, adjoint, generating-equation, gauge, determinant, and Casoratian identities.


\begin{thebibliography}{VIII.99}

\bibitem[III.1]{p3:KMLWSK26}
R.~Kalisch, T.~Monderer, E.~Leibtag, S.~Weinbaum, M.~Shalyt, and I.~Kaminer,
\emph{An Infinite Family of Ap\'ery-like Recurrences for Euler's Constant},
arXiv:2610.13919v1 (9 October 2026).
\url{https://arxiv.org/abs/2610.13919}.

\bibitem[III.2]{p3:Riv09}
T.~Rivoal, \emph{Rational approximations for values of derivatives of the
Gamma function}, Trans. Amer. Math. Soc. \textbf{361} (2009), 6115--6149.
Author manuscript: \url{https://rivoal.perso.math.cnrs.fr/articles/eulerconstant.pdf}.

\bibitem[III.3]{p3:XuYou2014}
H.~Xu and X.~You, \emph{Continued fraction inequalities for the
Euler--Mascheroni constant}, J. Inequal. Appl. \textbf{2014}, article 343.
\url{https://doi.org/10.1186/1029-242X-2014-343}.

\bibitem[III.4]{p3:Ferguson2020}
T.~Ferguson, \emph{Rational Approximations via Hankel Determinants},
arXiv:2003.10616v1 (2020).
\url{https://arxiv.org/abs/2003.10616}.

\bibitem[III.5]{p3:DLMF}
F.~W.~J. Olver et al., eds., \emph{NIST Digital Library of Mathematical
Functions}, Sections 5.9, 5.15, 6.2, 6.4, 8.6, 16.17, 18.3, and 18.12.
\url{https://dlmf.nist.gov/}.

\bibitem[III.6]{p3:Stieltjes1894}
T.~J. Stieltjes, \emph{Recherches sur les fractions continues},
Ann. Fac. Sci. Toulouse \textbf{8} (1894), J1--J122;
\textbf{9} (1895), A5--A47.

\bibitem[III.7]{p3:Wall1948}
H.~S. Wall, \emph{Analytic Theory of Continued Fractions},
D. Van Nostrand, New York, 1948.

\bibitem[III.8]{p3:ShentonBowman1971}
L.~R. Shenton and K.~O. Bowman, \emph{Continued Fractions for the PSI
Function and Its Derivatives}, SIAM J. Appl. Math. \textbf{20} (1971), 547--554.
\url{https://doi.org/10.1137/0120057}.

\bibitem[III.9]{p3:BonanJones2005}
C.~M. Bonan-Hamada and W.~B. Jones, \emph{Stieltjes continued fractions
for polygamma functions; speed of convergence},
J. Comput. Appl. Math. \textbf{179} (2005), 47--55.
\url{https://doi.org/10.1016/j.cam.2004.09.034}.

\bibitem[III.10]{p3:Srinivasan2008}
V.~R. Srinivasan, \emph{Extensions by Antiderivatives, Exponentials of
Integrals and by Iterated Logarithms}, arXiv:0811.3004 (2008),
Theorem 2.3.
\url{https://arxiv.org/abs/0811.3004}.

\bibitem[III.11]{p3:BCNZ23}
J.~P. Bell, S.~Chen, K.~D. Nguyen, and U.~Zannier,
\emph{D-finiteness, rationality, and height III: multivariate
P\'olya--Carlson dichotomy}, arXiv:2306.02590 (2023), Proposition 2.1.
\url{https://arxiv.org/abs/2306.02590}.

\bibitem[III.12]{p3:HHP11}
Kh.~Hessami Pilehrood and T.~Hessami Pilehrood,
\emph{Rational approximations to values of the digamma function and a
conjecture on denominators}, Math. Notes \textbf{90} (2011), 730--747;
arXiv:1004.0578.
\url{https://arxiv.org/abs/1004.0578}.

\bibitem[III.13]{p3:Arreche2013}
C.~E. Arreche, \emph{A Galois-theoretic proof of the differential
transcendence of the incomplete Gamma function},
J. Algebra \textbf{389} (2013), 119--127; arXiv:1304.1917.
\url{https://arxiv.org/abs/1304.1917}.

\bibitem[III.14]{p3:BDR2024}
A.~Bostan, L.~Di Vizio, and K.~Raschel,
\emph{Differential transcendence of Bell numbers and relatives:
a Galois theoretic approach}, Amer. J. Math. \textbf{146} (2024),
no.~6, 1577--1615; arXiv:2012.15292.
\url{https://arxiv.org/abs/2012.15292}.

\bibitem[III.15]{p3:FischlerRivoal2026Units}
S.~Fischler and T.~Rivoal,
\emph{The ring of values of E-functions is not a field},
arXiv:2610.16603v1 (2026), Conjecture 1.2 and Proposition 3.1.
\url{https://arxiv.org/abs/2610.16603}.

\bibitem[III.16]{p3:RivoalXUPS2019}
T.~Rivoal, \emph{Les E-fonctions et G-fonctions de Siegel},
Journ\'ees math\'ematiques X-UPS (2019), 197--298, Problem 6 on p.~286.
\url{https://doi.org/10.5802/xups.2019-03}.

\end{thebibliography}
% END PRESERVED SOURCE BODY 3
\phantomsection\label{part3:end}
\clearpage
\setcounter{section}{0}\setcounter{subsection}{0}
\renewcommand{\thesection}{\Roman{part}.\arabic{section}}
\renewcommand{\theHsection}{\arabic{part}.main.\arabic{section}}
\renewcommand{\refname}{References for Part IV}
\fancyhead[L]{\small Euler\textquotesingle s constant: Part IV}
\part{Twelve further irrationality attempts and their failure points}
\label{part4:start}
\noindent\textbf{Source:} {\small\path{gamma_attempts(2).tex}}.

Original PDF: 14 pages. Mathematical content and source-level qualifications are retained.

\noindent\textit{Original preparation note.} Prepared with the assistance of Claude (Anthropic), 18--19 October 2026. Computations used Python with \texttt{python-flint} (Arb) and SymPy.

\noindent\textbf{Editorial scope note.}
The broad holonomy verdict below is a record of the attempted method.
Its precise cancellation proposition must be read with its identifiability
hypotheses and with the later analysis in
Section~\ref{p6:sec:holonomy-scope} and Part VII. It is not a universal
theorem that every arithmetic-holonomy route to $\gamma$ is impossible.
The source author field does not identify a named author; the supplied
preparation credit is retained above.


% BEGIN PRESERVED SOURCE BODY 4
\begin{abstract}
We record twelve attempts to prove that Euler's constant $\gamma$ is irrational. They fall into four families, and each was chosen to avoid the routes already mapped in the companion manuscripts~\cite{p4:K1,p4:K2,p4:K3,p4:K4,p4:K5}. None succeeds. For each attempt we prove what can be proved and isolate the exact point of failure.

The most informative failure concerns the arithmetic holonomy method of Calegari, Dimitrov and Tang~\cite{p4:CDT}, currently the strongest irrationality technique for periods. Its analytic input is a cancellation between $G$-functions. By a theorem of Fischler and Rivoal~\cite{p4:FR14}, the nonzero constants that occur as cancellation coefficients in discs are exactly the elements of $\Frac(\bG)$, where $\bG$ is the ring of values of $G$-functions. We extend the necessity part to several constants and to meromorphic continuation along arbitrary paths (Proposition~\ref{p4:prop:cancel}). Consequently the method, in its present form, can reach $\gamma$ only if $\gamma\in\Frac(\bG)$. This is expected to be false, but it is open (Remarks~\ref{p4:rem:blind} and~\ref{p4:rem:expect}).

The other attempts fail at one of three walls:
\begin{itemize}[leftmargin=2em]
\item factorial denominators against sub-exponential analytic decay;
\item contamination by auxiliary transcendental constants;
\item blindness of the structure to $\gamma$.
\end{itemize}
These other attempts are denominator-free $\Gamma$-kernel forms, higher-rank Laplace kernels, $q$-deformations, quantum periods of Fano manifolds, o-minimal point counting, integer-valued functions, the Gauss map, Kronecker--Lerch--Chowla--Selberg, the Dickman--Buchstab functions, two-point Hermite--Pad\'e ``untying'', and the explicit formula with Baker's theory.

Two numerical studies are included. The first views Aptekarev's integer forms inside the denominator-free lattice: their gcd is the trivial factor $n!/d_n$ times a small integer (between $3$ and $39591$ for $n=20,30,\dots,300$), so no hidden arithmetic gain appears. The second treats the moments of the Dickman distribution. Their denominators are bounded by $e^{k\log k+O(k)}$, and numerically they are of that order.

\emph{Nothing in this note proves or disproves the irrationality of $\gamma$.}
\end{abstract}





%======================================================================
\section{Introduction}
%======================================================================

Let $\gamma=\lim_{n\to\infty}(H_n-\log n)=0.5772156649\ldots$, where $H_n=1+\frac12+\dots+\frac1n$. Whether $\gamma$ is irrational is open; see Lagarias' survey~\cite{p4:Lag} for the history.

\medskip\noindent\textbf{What is known.} The unconditional arithmetic results all concern $\gamma$ together with a companion constant.
\begin{itemize}[leftmargin=2em]
\item Mahler~\cite{p4:Mah} proved that $\frac{\pi}{2}\frac{Y_0(2)}{J_0(2)}-\gamma$ is transcendental.
\item From Shidlovskii's theorem and the identity
\begin{equation}\label{p4:eq:hardy}
\Ein(1)=\gamma+\delta/e,\qquad \Ein(z):=\int_0^z\frac{1-e^{-t}}{t}\,dt,\qquad \delta:=\int_0^\infty\frac{e^{-t}}{1+t}\,dt ,
\end{equation}
it follows that at least one of $\gamma$ and the Gompertz constant $\delta$ is transcendental. Rivoal~\cite{p4:Riv12} proved $\operatorname{trdeg}_\Q\Q(e,\gamma,\delta)\ge2$.
\item Murty--Saradha proved that for each $q>1$ at most one of $\gamma$ and the $\psi(a/q)$ with $(a,q)=1$ is algebraic~\cite{p4:MS07}, with an analogue for Euler--Lehmer constants~\cite{p4:MS10}. Chatterjee--Gun extended this to the whole set $\{\gamma\}\cup\{\psi(r/n):n>1,\ 0<r<n,\ (r,n)=1\}$~\cite{p4:CG}.
\item Fischler and Rivoal showed that the asymptotic constants of $E$-functions at infinity lie in the $\bG$-module generated by the derivatives $\Gamma^{(s)}(a)$ at rational $a$~\cite{p4:FR16}; this is where $\gamma=-\Gamma'(1)$ enters. They also showed that $\gamma\in\bE+e^{-1}\bD$~\cite{p4:FR24}. Their conjecture $\bE\cap\bD=\Qbar$ would make $\gamma$ transcendental.
\end{itemize}

\medskip\noindent\textbf{Rational approximations.} All known rational approximation schemes have errors $e^{-cn^\alpha}$ with $\alpha<1$ and denominators at least of size $n!$:
\begin{itemize}[leftmargin=2em]
\item Hessami Pilehrood: $\alpha=\tfrac12$~\cite{p4:HP};
\item Aptekarev: $\alpha=\tfrac12$~\cite{p4:Apt};
\item Rivoal: $\alpha=\tfrac23$~\cite{p4:Riv09};
\item Van Assche--Wolfs: $\alpha=\tfrac34$~\cite{p4:VAW}.
\end{itemize}
Sondow's double integrals~\cite{p4:Son} decay geometrically but contain logarithms of primes. If $\gamma=p/q$, then $q>10^{244663}$ (Papanikolaou; see~\cite{p4:Lag}). Several unrefereed proof claims have appeared since 2023; none has been accepted.

\medskip\noindent\textbf{The companion manuscripts.} The manuscripts~\cite{p4:K1,p4:K2,p4:K3,p4:K4,p4:K5} audit many dozens of routes. Their conclusions are summarized in~\S\ref{p4:sec:constraints}; we use them only as constraints on what a new attempt must avoid. In particular:
\begin{itemize}[leftmargin=2em]
\item the \emph{shear/twist} barrier of~\cite{p4:K1};
\item the \emph{tie theorem} of~\cite{p4:K1}: the small forms of Rivoal type see $\gamma$ only through the $E$-value $e\Ein(1)$;
\item the two failure mechanisms isolated throughout~\cite{p4:K2,p4:K3,p4:K4,p4:K5}: eliminating auxiliary constants costs more height than the analytic decay provides, and fully rational approximants have factorial denominators against sub-exponential errors.
\end{itemize}

\medskip\noindent\textbf{This note.} Twelve attempts are organized in four families:
\begin{itemize}[leftmargin=2em]
\item new kernels (\S\ref{p4:sec:kernels});
\item the strongest current tool (\S\ref{p4:sec:holonomy});
\item structures that turn out to be $\gamma$-blind (\S\ref{p4:sec:blind});
\item new bridges (\S\ref{p4:sec:bridges}).
\end{itemize}
For each we state what is proved, give computations where useful, and pin down the failure. Section~\ref{p4:sec:synthesis} collects the obstructions in one table and states what seems to be missing.

Statements labelled Theorem, Proposition, Lemma or Corollary are proved here or quoted with a reference. \emph{Observations} are numerical, and \emph{Remarks} and the ``Where it breaks'' paragraphs contain heuristic discussion. We were unable to consult the October 2026 preprint arXiv:2610.13919 (``An infinite family of Ap\'ery-like recurrences for Euler's constant''); nothing below depends on it.

\medskip\noindent\textbf{Notation.}
\begin{itemize}[leftmargin=2em]
\item $d_n=\lcm(1,\dots,n)$.
\item $\psi=\Gamma'/\Gamma$.
\item $E_1(x)=\int_x^\infty e^{-t}\,dt/t$, so that $\Ein(x)=\gamma+\log x+E_1(x)$ for $x>0$.
\item $\bE$ and $\bD$ are the rings of values at algebraic points of $E$-functions and of Borel sums of arithmetic Gevrey series~\cite{p4:FR16,p4:FR18,p4:FR24}.
\item $\bG$ is the ring of values at algebraic points of $G$-functions with algebraic coefficients, including values of their analytic continuations along paths avoiding singularities. Let $\bG_\Q$ denote the values of $G$-functions with rational coefficients at rational points of their discs of convergence. By~\cite[Thm.~1]{p4:FR14}, $\bG=\bG_\Q+i\bG_\Q$; it contains $\Qbar$, $\pi$, and the logarithms of nonzero algebraic numbers.
\end{itemize}

%======================================================================
\section{Two constraints every attempt must respect}\label{p4:sec:constraints}
%======================================================================

\subsection{The ledger}

\begin{lemma}\label{p4:lem:ledger}
Let $(A_n,B_n)\in\Q^2$ and $D_n\in\N$ with $D_nA_n,D_nB_n\in\Z$. Put $\Lambda_n=A_n-\gamma B_n$. If $0<D_n|\Lambda_n|\to0$ along a subsequence, then $\gamma\notin\Q$.
\end{lemma}

\begin{proof}
If $\gamma=a/b$, then $bD_n\Lambda_n=bD_nA_n-aD_nB_n$ is an integer. It is nonzero, so its absolute value is at least $1$, whereas it tends to $0$.
\end{proof}

We call $\log D_n+\log|\Lambda_n|$ the \emph{ledger} of the family. A proof by linear forms needs a family with certified nonvanishing whose ledger tends to $-\infty$. (Proofs of other types, such as the holonomy method of~\S\ref{p4:sec:holonomy}, are not of this form.)

\subsection{Bridges and the twist}

The following observation is the shear-invariance principle of~\cite{p4:K1}, stated in the language of $\Gamma$.

\begin{lemma}[Twist]\label{p4:lem:twist}
For $\tau\in\R$ put $\Gamma_\tau(s):=e^{\tau(1-s)}\Gamma(s)$ and $\psi_\tau=\Gamma_\tau'/\Gamma_\tau=\psi-\tau$.
\begin{enumerate}[label=(\roman*),leftmargin=2em]
\item $\Gamma_\tau(1)=1$, and $\Gamma_\tau$ is log-convex on $(0,\infty)$.
\item $\psi_\tau$ satisfies the same identities as $\psi$:
\begin{itemize}[leftmargin=1.5em]
\item $\psi_\tau(s+1)=\psi_\tau(s)+1/s$;
\item $\psi_\tau(1-s)-\psi_\tau(s)=\pi\cot\pi s$;
\item $\sum_{k=0}^{m-1}\psi_\tau(s+k/m)=m\psi_\tau(ms)-m\log m$;
\item Gauss's digamma theorem, with $\gamma$ replaced by $\gamma_\tau:=-\psi_\tau(1)=\gamma+\tau$.
\end{itemize}
\item The $\Gamma$-level normalizations change:
\begin{itemize}[leftmargin=1.5em]
\item $\Gamma_\tau(\tfrac12)=e^{\tau/2}\sqrt\pi$;
\item $\Gamma_\tau(s)\Gamma_\tau(1-s)=e^{\tau}\pi/\sin\pi s$;
\item $\int_0^1\log\Gamma_\tau=\tfrac12\log2\pi+\tfrac\tau2$;
\item $\Gamma_\tau(s+1)=e^{-\tau}s\,\Gamma_\tau(s)$.
\end{itemize}
\end{enumerate}
\end{lemma}

\begin{proof}
Direct computation. Each identity in (ii) is unchanged when every $\psi$ is replaced by $\psi-\tau$. In Gauss's theorem, $\psi(r/q)-\psi(1)$ is $\tau$-invariant.
\end{proof}

Consequently, an argument that uses only identities of type (ii), together with a characterization of $\psi$ up to an additive constant, proves the same statement for $\gamma+\tau$ for every $\tau$. It therefore cannot prove that $\gamma\notin\Q$. A proof must use a \emph{bridge}: an identity of type (iii), or an asymptotic normalization such as $\psi(x)-\log x\to0$.

The bridges available, and the arithmetic cost of each, are:
\begin{enumerate}[label=(B\arabic*),leftmargin=3em]
\item $\gamma=\lim(H_n-\log n)$ and its Euler--Maclaurin refinements: logarithms of integers.
\item $\gamma=\Ein(x)-\log x-E_1(x)$: $E$-values at $x$, the logarithm $\log x$, and the $\bD$-value $e^xE_1(x)$. For $x=1$ this is~\eqref{p4:eq:hardy}.
\item Stirling, Raabe, Wallis: $\log\pi$, $\log 2\pi$, and logarithms of integers.
\item Lerch and Chowla--Selberg: $\log\Gamma(a/d)$ and Euler--Kronecker constants (Attempt~9).
\item de Bruijn's identity $\int_0^\infty\rho=e^\gamma$: the number $e^\gamma$ (Attempt~10).
\item The $\Gamma$-kernel itself, $\gamma=-\int_0^\infty e^{-t}\log t\,dt$ (Attempt~1).
\end{enumerate}

%======================================================================
\section{Family I: new kernels}\label{p4:sec:kernels}
%======================================================================

\subsection{\texorpdfstring{Attempt 1: denominator-free forms from the $\Gamma$-kernel}{Attempt 1: denominator-free forms from the Gamma-kernel}}

The usual ledger is ruined by the denominators $d_n$ of harmonic numbers. They disappear completely if one integrates integer polynomials against the $\Gamma$-kernel.

\begin{lemma}\label{p4:lem:poch}
For $P(t)=\sum_{k=0}^np_kt^k\in\Z[t]$,
\[
\int_0^\infty P(t)\,e^{-t}\log t\,dt = A(P)-\gamma B(P),\qquad A(P)=\sum_k p_k\,c_k,\quad B(P)=\sum_k p_k\,k!,
\]
where $c_k:=k!H_k=|s(k+1,2)|\in\Z$ are unsigned Stirling numbers of the first kind. The map $P\mapsto(A(P),B(P))$ is already surjective from polynomials of degree $\le1$ onto $\Z^2$.
\end{lemma}

\begin{proof}
We have $\int_0^\infty t^ke^{-t}\log t\,dt=\Gamma'(k+1)=k!\,(H_k-\gamma)$. For surjectivity, $P=(B-A)+At$ gives $(A(P),B(P))=(A,B)$.
\end{proof}

\begin{corollary}\label{p4:cor:poch}
$\gamma\notin\Q$ if and only if there are $P_n\in\Z[t]$ with $0<\bigl|\int_0^\infty P_ne^{-t}\log t\,dt\bigr|\to0$.
\end{corollary}

\begin{proof}
For the ``if'' direction, apply Lemma~\ref{p4:lem:ledger} with $D_n=1$. For the ``only if'' direction, use Dirichlet's theorem and the surjectivity in Lemma~\ref{p4:lem:poch}.
\end{proof}

\noindent\textbf{Where it breaks.} The lattice has no denominators, but the integral carries no analytic information beyond the pair $(A,B)$: every pair is represented by a polynomial of degree one. Smallness must be \emph{forced} by a construction in an auxiliary variable, such as vanishing conditions (Pad\'e) or Rodrigues formulas. All such constructions known for $\gamma$ pass through the confluent saddle of $e^{-t}$ and decay like $e^{-cn^\alpha}$ with $\alpha<1$.

In the lattice of Lemma~\ref{p4:lem:poch}, the normalizations appear as gcd's, and one can test for hidden arithmetic gains. Aptekarev's forms, in the integer normalization, are
\[
\tilde q_n=\sum_{k=0}^n\binom nk^2(n+k)!,\qquad
\tilde p_n=\sum_{k=0}^n\binom nk^2(n+k)!\,\bigl(H_{n+k}+2H_{n-k}-2H_k\bigr).
\]
Both are integers, and both are divisible by the trivial factor $n!/d_n$.

\begin{observation}\label{p4:obs:apt}
For $n=20,30,\dots,300$, exact computation (partly shown in Table~\ref{p4:tab:apt}) gives three things:
\begin{itemize}[leftmargin=2em]
\item $\gcd(\tilde q_n,\tilde p_n)=(n!/d_n)\cdot\epsilon_n$ with $\epsilon_n$ a small integer, $3\le\epsilon_n\le39591$ (so $\log\epsilon_n\le10.6$), against the $\log|\Lambda_n|\approx3217$ that would have to be absorbed at $n=300$; there is no arithmetic gain beyond the trivial factor;
\item $\log|\gamma-\tilde p_n/\tilde q_n|+2\sqrt{2n}$ increases slowly towards $\log 2\pi=1.838$, consistent with the error $\approx2\pi e^{-2\sqrt{2n}}$;
\item the reduced form $|\tilde q_n\gamma-\tilde p_n|/\gcd$ grows like $e^{2096}$ at $n=300$.
\end{itemize}
\end{observation}

\begin{table}[ht]
\centering
\small
\begin{tabular}{rrrrrr}
\toprule
$n$ & $\log|\Lambda_n|$ & $\log g_n$ & $\log(n!/d_n)$ & $\log(|\Lambda_n|/g_n)$ & $\log|\gamma-\tilde p_n/\tilde q_n|+2\sqrt{2n}$\\
\midrule
20 & 103.4 & 26.9 & 23.1 & 76.6 & 1.519\\
40 & 264.1 & 82.0 & 74.1 & 182.1 & 1.612\\
80 & 642.0 & 198.8 & 194.2 & 443.2 & 1.678\\
120 & 1062.6 & 347.1 & 340.4 & 715.5 & 1.707\\
200 & 1979.5 & 661.5 & 657.1 & 1318.0 & 1.737\\
300 & 3216.7 & 1120.7 & 1115.7 & 2096.0 & 1.755\\
\bottomrule
\end{tabular}
\caption{Aptekarev's forms in the denominator-free lattice, with $\Lambda_n=\tilde p_n-\tilde q_n\gamma$ and $g_n=\gcd(\tilde q_n,\tilde p_n)$ (natural logarithms). Exact integer arithmetic; $\gamma$ computed in Arb.}\label{p4:tab:apt}
\end{table}

\noindent\textbf{Verdict.} The lattice relocates the $d_n$-wall into the gcd and changes nothing. Discarded.

\subsection{\texorpdfstring{Attempt 2: higher-rank Laplace kernels $e^{-t^r}$}{Attempt 2: higher-rank Laplace kernels}}

Replacing $e^{-t}$ by $e^{-t^r}$ changes the moments from $k!$ to $\Gamma(\frac{k+1}r)/r$, roughly $(k!)^{1/r}$. Pad\'e-type decay at an irregular singularity of Poincar\'e rank $r$ improves heuristically to $e^{-cn^{r/(r+1)}}$. Letting $r$ grow with $n$ would make both sides of the ledger exponential.

\begin{proposition}\label{p4:prop:rank}
Let $r\ge2$ and $P\in\Q[t]$. There are $\Q$-linear functionals $R_a,S_a$ ($1\le a\le r$) with
\[
\int_0^\infty P(t)e^{-t^r}\log t\,dt=\frac1{r^2}\sum_{a=1}^r\Gamma\!\left(\tfrac ar\right)\Bigl(R_a(P)\,\psi\!\left(\tfrac ar\right)+S_a(P)\Bigr).
\]
For $1\le a<r$,
\[
\psi(a/r)=-\gamma-\log2r-\tfrac\pi2\cot\tfrac{\pi a}r+2\sum_{1\le j\le(r-1)/2}\cos\tfrac{2\pi ja}r\log\sin\tfrac{\pi j}r .
\]
The component $a=r$ is the only one that does not involve $\Gamma(a/r)$ with $a<r$. On that component the substitution $u=t^r$ reduces the form to the rank-one kernel of Attempt~1.
\end{proposition}

\begin{proof}
First,
\[
\int_0^\infty t^ke^{-t^r}\log t\,dt=r^{-2}\Gamma'\bigl(\tfrac{k+1}r\bigr).
\]
Write $\frac{k+1}r=m+\frac ar$ with $1\le a\le r$, $m\ge0$, and use
\[
\Gamma(m+\tfrac ar)=\Gamma(\tfrac ar)(\tfrac ar)_m,\qquad
\psi(m+\tfrac ar)=\psi(\tfrac ar)+\sum_{j<m}(\tfrac ar+j)^{-1}.
\]
The formula for $\psi(a/r)$ is Gauss's digamma theorem.
\end{proof}

\noindent\textbf{Where it breaks.} Any gain from $r>1$ comes with new constants:
\begin{itemize}[leftmargin=2em]
\item the values $\Gamma(a/r)$, which are expected to be algebraically independent with $\pi$ modulo the standard relations (Lang--Rohrlich);
\item logarithms of cyclotomic units;
\item $\pi\cot(\pi a/r)$.
\end{itemize}
For $r=2$ these are $\sqrt\pi$ and $\log2$; for growing $r$ there are about $\varphi(r)/2$ independent ones. This is exactly the contamination wall of~\cite{p4:K2,p4:K3}. Discarded.

\subsection{\texorpdfstring{Attempt 3: $q$-deformation and the limit $q\to1$}{Attempt 3: q-deformation and the limit q -> 1}}

Put $\zeta_q(1):=\sum_{n\ge1}q^n/(1-q^n)=\sum_{n\ge1}d(n)q^n$ for $0<q<1$.

\begin{proposition}\label{p4:prop:q}
As $q\to1^-$,
\[
(1-q)\,\zeta_q(1)+\log(1-q)=\gamma+O\bigl((1-q)\log\tfrac1{1-q}\bigr).
\]
\end{proposition}

\begin{proof}
With $q=e^{-t}$, the Mellin transform $\Gamma(s)\zeta(s)^2$ gives
\[
\sum d(n)e^{-nt}=\frac{\gamma-\log t}t+\frac14+O(t).
\]
Then use $1-q=t-t^2/2+O(t^3)$.
\end{proof}

The irrationality of $\zeta_q(1)$ is known for $q=1/p$ with $p\in\Z$, $|p|\ge2$: Erd\H{o}s for $p=2$, B\'ezivin~\cite{p4:Bez} and Borwein~\cite{p4:Bor} in general, with irrationality exponent at most $2.46498$ uniformly in $p$ (Zudilin~\cite{p4:Zud}). Kurokawa and Wakayama's $q$-analogue of $\gamma$~\cite{p4:KW} is of this type.

\noindent\textbf{Where it breaks.} Two independent obstructions:
\begin{itemize}[leftmargin=2em]
\item These results describe the limit $q\to0$ along $1/p$, not $q\to1$. The proofs use that the $q$-analogues of $d_n$ (products of cyclotomic values) are integers when $1/q\in\Z$. We know of no irrationality result for $\zeta_q(1)$ at $q=1-1/m$.
\item Even with such results, Proposition~\ref{p4:prop:q} at $q=1-1/m$ approximates $\gamma$ by $\frac1m\zeta_{1-1/m}(1)-\log m$, so logarithms of integers return.
\end{itemize}
Discarded.

%======================================================================
\section{Family II: the strongest current tool and its blind spot}\label{p4:sec:holonomy}
%======================================================================

\subsection{Attempt 4: arithmetic holonomy bounds}

Calegari, Dimitrov and Tang~\cite{p4:CDT} proved that $1,\zeta(2),L(2,\chi_{-3})$ are linearly independent over $\Q$. Their method bounds the dimension of spaces of holonomic power series with controlled denominators, of integrality type $x^n/[1,\dots,n]^2$, which continue meromorphically to a domain of large conformal size. The constants enter as connection constants of a holonomic local system of $G$-type (Zagier's functions on $\C\smallsetminus\{0,\tfrac19,1\}$). If such a constant $\xi$ were rational, a combination $A-\xi B$ of $G$-functions would continue further than its parts, producing too many functions of that kind.

The analytic input of the method is therefore a \emph{cancellation}: constants $\xi_1,\dots,\xi_r$ such that $A-\sum_i\xi_iB_i$ continues further than the $B_i$ do. The next two statements identify which constants can occur in such cancellations.

\begin{theorem}[Fischler--Rivoal {\cite[Thm.~2]{p4:FR14}}, case $K=\Qbar$]\label{p4:thm:FR}
For $\xi\in\C^\times$ the following are equivalent:
\begin{enumerate}[label=(\roman*),leftmargin=2em]
\item $\xi\in\Frac(\bG)$;
\item there are coefficient sequences $(a_n),(b_n)$ of $G$-functions with $b_n\ne0$ for large $n$ and $a_n/b_n\to\xi$;
\item for every $R\ge1$ there are $G$-functions $A,B\in\Qbar[[z]]$ of radius $1$ such that $A-\xi B$ has radius $>R$.
\end{enumerate}
\end{theorem}

\noindent In~\cite{p4:FR14} the coefficient field is an arbitrary algebraic extension $K$ of $\Q$ and (i) reads $\xi\in\Frac(\bG_K)$. For $K=\Qbar$ one has $\bG_{\Qbar}=\bG_\Q+i\bG_\Q=\bG$ by~\cite[Thm.~1]{p4:FR14}. The assumption $\xi\neq0$ is needed: for $\xi=0$, (i) holds but (iii) fails.

The method of~\cite{p4:CDT} uses meromorphic continuation to general domains rather than discs, and several constants at once. The following proposition covers that case.

\begin{proposition}[Cancellation along a path]\label{p4:prop:cancel}
Let $A,B_1,\dots,B_r\in\Qbar[[z]]$ be $G$-functions and $\xi_1,\dots,\xi_r\in\C$. Let $\sigma$ be a path from $0$ to a point $\zeta\in\C$ along which $A$ and the $B_i$ continue analytically, up to (but not including) $\zeta$. Suppose that:
\begin{itemize}[leftmargin=2em]
\item the continuation of $A-\sum_i\xi_iB_i$ is meromorphic at $\zeta$;
\item no nontrivial $\C$-linear combination of the continuations of the $B_i$ is meromorphic at $\zeta$.
\end{itemize}
Then $\zeta\in\Qbar$ and $\xi_1,\dots,\xi_r\in\Frac(\bG)$. The same holds with ``meromorphic'' replaced by ``holomorphic'' in both hypotheses.
\end{proposition}

\begin{proof}[Proof (from standard facts)]
\emph{Step 1: the operator.} Let $L$ be the least common left multiple of the minimal operators of $A,B_1,\dots,B_r$. By the theorem of Chudnovsky and Andr\'e~\cite{p4:Chu,p4:And}, $L$ is a $G$-operator. By the theorems of Katz and Andr\'e~\cite{p4:And,p4:And00} it is Fuchsian with rational exponents, and near each of its singular points $\zeta$ it has a basis of solutions
\[
u_j=\sum_k(z-\zeta)^{e_j}\log^k(z-\zeta)\,f_{j,k}(z-\zeta),\qquad e_j\in\Q,
\]
with $f_{j,k}$ $G$-functions. The point $\zeta$ must be a singular point of $L$: at a regular point all local solutions are holomorphic, which contradicts the second hypothesis (note $r\ge1$ there). Singular points of $L$ are algebraic, so $\zeta\in\Qbar$.

\emph{Step 2: the connection coefficients.} Modifying the final segment of $\sigma$ inside a small disc around $\zeta$ if necessary, choose an algebraic point $z_0$ on $\sigma$ with three properties: $z_0$ is a regular point of $L$; $|z_0-\zeta|$ is smaller than the radii of convergence of all the $f_{j,k}$; and the disc around $z_0$ used below avoids $\zeta$. Fix the branches of $(z-\zeta)^{e_j}$ and $\log(z-\zeta)$ near $z_0$ by continuation along $\sigma$. On a small disc around $z_0$ write
\[
A=\sum_j\alpha_ju_j,\qquad B_i=\sum_j\beta_{ij}u_j .
\]
By Cramer's rule, $w(z_0)\alpha_j$ and $w(z_0)\beta_{ij}$ are polynomials in the entries of the Wronskian matrix $W_u(z_0)$ and in the values $A^{(l)}(z_0)$, $B_i^{(l)}(z_0)$. Here $w(z_0)=\det W_u(z_0)\ne0$.
\begin{itemize}[leftmargin=1.5em]
\item The values $A^{(l)}(z_0)$ and $B_i^{(l)}(z_0)$ lie in $\bG$ by~\cite[Thm.~1]{p4:FR14}.
\item The entries of $W_u(z_0)$ are polynomials in algebraic numbers $(z_0-\zeta)^{e_j-m}$, in $\log(z_0-\zeta)$, and in values $f_{j,k}^{(m)}(z_0-\zeta)\in\bG$. The logarithm of a nonzero algebraic number lies in $\bG$ for every branch, because $\log(1-z)$ is a $G$-function and $2\pi i\in\bG$.
\end{itemize}
Hence $w(z_0)\alpha_j$ and $w(z_0)\beta_{ij}$ lie in $\bG$.

\emph{Step 3: the cancellation.} A local solution $\sum_jc_ju_j$ is meromorphic at $\zeta$ if and only if, in its expansion, the coefficients of all terms $(z-\zeta)^e\log^k(z-\zeta)$ with $e\notin\Z$ or $k\ge1$ vanish. It is holomorphic if, in addition, the coefficients with $e\in\Z_{<0}$ vanish. The Taylor coefficients of the $f_{j,k}$ are algebraic, so these conditions form a system $\ell(c)=0$, $\ell\in\mathcal L$, of $\Qbar$-linear equations; finitely many suffice. Put $\alpha=(\alpha_j)$ and $\beta_i=(\beta_{ij})_j$.
\begin{itemize}[leftmargin=1.5em]
\item The first hypothesis gives $\ell(\alpha)=\sum_i\xi_i\,\ell(\beta_i)$ for all $\ell\in\mathcal L$.
\item The second hypothesis says that the matrix $\bigl(\ell(\beta_i)\bigr)_{\ell\in\mathcal L,\,1\le i\le r}$ has rank $r$.
\end{itemize}
Choose $\ell_1,\dots,\ell_r\in\mathcal L$ with $\det\bigl(\ell_k(\beta_i)\bigr)\ne0$ and solve for the $\xi_i$ by Cramer's rule. After multiplying numerators and denominators by $w(z_0)^r$, all entries lie in $\bG$, so $\xi_i\in\Frac(\bG)$.
\end{proof}

\begin{remark}[The $G$-blindness principle]\label{p4:rem:blind}
Proposition~\ref{p4:prop:cancel} has the following consequence. Consider any argument whose analytic input is a cancellation of this kind between $G$-functions; the arithmetic holonomy method of~\cite{p4:CDT}, in its present form, is of this type. Such an argument can only produce statements about constants in $\Frac(\bG)$. Conversely, by Theorem~\ref{p4:thm:FR}, every nonzero element of $\Frac(\bG)$ admits such cancellations. In particular, applying the method to $\gamma$ would require a cancellation as input, and Proposition~\ref{p4:prop:cancel} would then force $\gamma\in\Frac(\bG)$.

This is a principle about a class of methods, not a theorem about all conceivable proofs. It says nothing about arguments whose input is of a different kind.
\end{remark}

\begin{remark}\label{p4:rem:expect}
It is expected that $\gamma\notin\Frac(\bG)$, but this is not known to follow from any standard functional conjecture.
\begin{itemize}[leftmargin=2em]
\item Fischler and Rivoal conjecture that $\bG$ consists of periods with $1/\pi$ adjoined~\cite{p4:FR14}.
\item Kontsevich and Zagier~\cite{p4:KZ} expect that $\gamma$ is not a period.
\item Fres\'an and Jossen~\cite{p4:FJ} realize $\gamma$ as a period of a non-classical exponential motive.
\end{itemize}
So the method reaches $\gamma$ only if $\gamma\in\Frac(\bG)$, which would itself be a major surprise. A proof of the holonomy type would need a theorem for objects whose connection data contain $\gamma$. By~\cite[Thm.~3]{p4:FR16} these include the asymptotic constants of $E$-functions at their irregular singularity $\infty$, which involve the derivatives $\Gamma^{(s)}(a)$ at rational $a$ (for instance $\gamma=-\Gamma'(1)$). They also include Borel sums of arithmetic Gevrey series~\cite{p4:FR18}. No such theorem is known.
\end{remark}

\noindent\textbf{Verdict.} The method can reach $\gamma$ only if $\gamma\in\Frac(\bG)$, which is not expected. Discarded.

\subsection{Attempt 5: quantum periods of Fano manifolds and the Gamma conjecture}

Ap\'ery-type limits occur naturally in the quantum differential equations of Fano manifolds (Golyshev). The Gamma conjecture~\cite{p4:GGI} expresses their asymptotics through
\[
\widehat\Gamma_X=\exp\Bigl(-\gamma\,c_1(X)+\sum_{k\ge2}(-1)^k(k-1)!\,\zeta(k)\,\mathrm{ch}_k(T_X)\Bigr),
\]
in which $\gamma$ does appear.

\noindent\textbf{Where it breaks.} Take $X$ with a Laurent-polynomial mirror $f$.
\begin{itemize}[leftmargin=2em]
\item The regularized quantum period $\widehat G_X(t)=\sum_d[f^d]_0\,t^d$, with $[\,\cdot\,]_0$ the constant term, has integer coefficients and is a period of the mirror family. Hence it is a $G$-function.
\item Its Ap\'ery limits are limits of ratios of coefficient sequences of $G$-functions, so they lie in $\Frac(\bG)$ by Theorem~\ref{p4:thm:FR}(ii).
\item The constant $\gamma$ enters $\widehat\Gamma_X$ only through $e^{-\gamma c_1}$. This factor is the Laplace normalization relating $\widehat G_X$ to the $E$-function $G_X(t)=\sum_d[f^d]_0\,t^d/d!$, in accordance with~\cite[Thm.~3]{p4:FR16}.
\end{itemize}

\noindent\textbf{Verdict.} The Ap\'ery mechanism of quantum cohomology lives on the $G$-side. By Remark~\ref{p4:rem:blind} it can reach $\gamma$ only if $\gamma\in\Frac(\bG)$. Discarded.

%======================================================================
\section{\texorpdfstring{Family III: structures that do not see $\gamma$}{Family III: structures that do not see gamma}}\label{p4:sec:blind}
%======================================================================

\subsection{Attempt 6: o-minimal point counting}

The function $\Gamma|_{(0,\infty)}$, and hence $\psi|_{[1,\infty)}$, is definable in an o-minimal expansion of $\R$ (van den Dries--Speissegger~\cite{p4:vdDS}).

\begin{proposition}\label{p4:prop:omin}
\begin{enumerate}[label=(\alph*),leftmargin=2em]
\item Suppose $\gamma=a/b$. Then the rational points of $\operatorname{graph}(\psi|_{[1,\infty)})$ are exactly the points $(n+1,H_n-a/b)$, $n\ge0$. Their number $N(T)$ of height $\le T$ satisfies $N(T)\asymp\log T$.
\item Suppose $\gamma\notin\Q$. Then the graph has no rational point over an integer, and all its rational points lie over a single class $r/q+\Z$ with $0<r<q$, $(r,q)=1$. In particular $N(T)=O(\log T)$ in this case too.
\end{enumerate}
\end{proposition}

\begin{proof}
\emph{Rational abscissae.} Write a rational abscissa as $x=m+r/q$ with $m\in\Z$ and $0\le r<q$. Then $\psi(x)-\psi(r/q)\in\Q$ (for $r>0$) and $\psi(m+1)=H_m-\gamma$. For $0<r<q$, the number $\psi(r/q)+\gamma$ is, by Gauss's digamma theorem, a nonzero linear form in logarithms of algebraic numbers and $\pi$ with algebraic coefficients, hence transcendental by Baker's theorem~\cite{p4:MS07}. If $\gamma=a/b$, this shows that $\psi(x)$ is transcendental at every non-integral rational $x$, which proves the first claim of (a). If $\gamma\notin\Q$, then $\psi(m+1)\notin\Q$, and by Chatterjee--Gun~\cite{p4:CG} at most one value $\psi(r/q)$ with $0<r<q$, $(r,q)=1$ is algebraic; rational points can therefore only lie over the corresponding class $r/q+\Z$.

\emph{Heights.} The primes in $(n/2,n]$ divide $\den H_n$, and $\den H_n\mid d_n$. Hence $e^{(1/2+o(1))n}\le\den H_n\le e^{(1+o(1))n}$, which gives $N(T)\asymp\log T$ in (a). The same estimate applies to $\psi(r/q)+\sum_{0\le j<m}(r/q+j)^{-1}$ in (b).
\end{proof}

\noindent\textbf{Where it breaks.} The Pila--Wilkie theorem~\cite{p4:PW} bounds rational points on the transcendental part by $O_\varepsilon(T^\varepsilon)$. Even Wilkie-type bounds $(\log T)^\kappa$ must allow $\kappa\ge1$, since the graph of $2^x$ already has $\asymp\log T$ rational points. The ``rational world'' produces too few points to be refuted by counting. Discarded.

\subsection{Attempt 7: integer-valued functions}

\begin{proposition}\label{p4:prop:intval}
Let $\N=\{0,1,2,\dots\}$. For $c\in\C$ put $F_c(z)=\Gamma'(1+z)+c\,\Gamma(1+z)$. Then $F_c(n)=n!(H_n-\gamma+c)$ for $n\in\N$. Moreover, $F_c(\N)\subset\Z$ if and only if $c\in\gamma+\Z$.
\end{proposition}

\begin{proof}
We have $\psi(n+1)=H_n-\gamma$ and $n!H_n\in\Z$; take $n=0$ for the converse.
\end{proof}

\noindent\textbf{Where it breaks.} Integrality on $\N$ detects the coset $\gamma+\Z$, never whether $\gamma\in\Q$. The canonical function $F_\gamma=\Gamma(1+z)H_z$, where $H_z:=\psi(1+z)+\gamma$, is integer-valued on $\N$ in every world.
\begin{itemize}[leftmargin=2em]
\item Rationality $\gamma=a/b$ only adds the statement $bF_\gamma\in\Z\Gamma'(1+\cdot)+\Z\Gamma(1+\cdot)$.
\item No P\'olya-type rigidity holds at the $\Gamma$-scale, since $\Gamma(1+z)$ is itself integer-valued of the same growth.
\item Carlson/P\'olya density versions require bounded integer outputs $bF_\gamma(n)-(\cdots)$, which is the ledger again.
\end{itemize}
Discarded.

\subsection{Attempt 8: the Gauss map}

Let $T(x)=\{1/x\}$ and let $\mathcal L f(y)=\sum_{n\ge1}f\bigl(\frac1{n+y}\bigr)(n+y)^{-2}$ be the Gauss--Kuzmin--Wirsing operator.

\begin{proposition}\label{p4:prop:gauss}
\begin{enumerate}[label=(\alph*),leftmargin=2em]
\item $1-\gamma=\int_0^1\{1/x\}\,dx=\int_0^1y\,\psi'(1+y)\,dy=\langle y,\mathcal L1\rangle$.
\item $\int_0^1\{1/x\}^2dx=\log2\pi-\gamma-1$.
\item $\int_0^1T^k(x)\,dx=\langle y,\mathcal L^k1\rangle$, and this tends to $\frac1{\log2}-1$ at the geometric rate $|\lambda_2|^k$, where $\lambda_2=-0.3036630\ldots$ is Wirsing's constant~\cite{p4:Wir}.
\end{enumerate}
\end{proposition}

\begin{proof}
For (a), substitute $x=1/(n+t)$ and use $\mathcal L1=\psi'(1+\cdot)$. Integrating by parts gives $\psi(2)-\log\Gamma(2)+\log\Gamma(1)=1-\gamma$. Part (b) is similar and was checked numerically to $25$ digits. Part (c) is the spectral theory of~\cite{p4:Wir}.
\end{proof}

\noindent\textbf{Where it breaks.} $\gamma$ is the first term of a geometrically convergent sequence whose limit $\frac1{\log2}-1$ lies in $\Q(\log2)$. The intermediate terms are governed by the spectrum of $\mathcal L$, and no arithmetic is known for that spectrum. Discarded.

\subsection{Attempt 9: Kronecker, Lerch and Chowla--Selberg}

For an imaginary quadratic field $K$ of discriminant $-d$, with character $\chi$, class number $h$ and $w$ roots of unity, let $\gamma_K$ be the Euler--Kronecker constant (the constant term of $\zeta_K'/\zeta_K$ at $s=1$, see~\cite{p4:Iha}).

\begin{proposition}\label{p4:prop:EK}
\[
\gamma_K=2\gamma+\log2\pi-\frac w{2h}\sum_{a=1}^{d-1}\chi(a)\log\Gamma\!\left(\tfrac ad\right).
\]
\end{proposition}

\begin{proof}
We have $\gamma_K=\gamma+L'(1,\chi)/L(1,\chi)$. Lerch's formula $\partial_s\zeta(0,x)=\log\Gamma(x)-\frac12\log2\pi$ gives $L'(0,\chi)$. The functional equation, with $\psi(1)=-\gamma$ and $\psi(\frac12)=-\gamma-2\log2$, transfers this to $s=1$, using $L(0,\chi)=2h/w$. The formula was checked numerically to $30$ digits for $d=3,4,7,8,11,15,23$.
\end{proof}

\noindent\textbf{Where it breaks.} Lerch's formula is a genuine bridge, of type (B4). By Chowla--Selberg~\cite{p4:CS}, however, the $\Gamma$-product is a period up to algebraic factors and powers of $\pi$. Hence $\gamma\equiv\gamma_K/2$ modulo logarithms of periods: one trades $\gamma$ for a Laurent constant of the same nature. Discarded.

%======================================================================
\section{Family IV: new bridges}\label{p4:sec:bridges}
%======================================================================

\subsection{\texorpdfstring{Attempt 10: the Dickman--Buchstab functions and $e^\gamma$}{Attempt 10: the Dickman-Buchstab functions and exp(gamma)}}

Let $\rho$ be Dickman's function: $\rho=1$ on $[0,1]$ and $u\rho'(u)=-\rho(u-1)$ for $u>1$.

\begin{proposition}[de Bruijn; see~\cite{p4:dB,p4:Ten}]\label{p4:prop:dick}
\begin{enumerate}[label=(\alph*),leftmargin=2em]
\item The Laplace transform is
\[
\hat\rho(s)=\int_0^\infty e^{-su}\rho(u)\,du=\exp\bigl(\gamma-\Ein(s)\bigr).
\]
Hence $\hat\rho(s)=e^{-E_1(s)}/s$ for $s>0$, which is free of $\gamma$.
\item $\int_0^\infty u^k\rho(u)\,du=e^\gamma m_k$, where
\[
m_0=1,\qquad m_k=\frac1k\sum_{i<k}\binom ki m_i,\qquad\text{equivalently}\qquad m_k=\sum_{\pi}\prod_{B\in\pi}\frac1{|B|},
\]
the sum running over all set partitions $\pi$ of $\{1,\dots,k\}$. Thus $m_1=1$, $m_2=\frac32$, $m_3=\frac{17}6$, $m_4=\frac{19}3$.
\end{enumerate}
\end{proposition}

\begin{proof}
Laplace-transforming the delay equation gives $s\hat\rho'=(e^{-s}-1)\hat\rho$, with normalization $\hat\rho(s)\sim1/s$ as $s\to\infty$. Comparing coefficients in $s\,u'=(e^{-s}-1)u$ for $u=e^{-\Ein}$ gives the recurrence. The partition formula is the exponential formula for $\exp\bigl(\sum_n(-1)^ns^n/(n\cdot n!)\bigr)$. The moments were confirmed numerically by integrating the delay equation.
\end{proof}

\begin{proposition}\label{p4:prop:dickden}
For every prime $p$ and $k\ge1$, $v_p(\den m_k)\le\lfloor k/p\rfloor$.
\end{proposition}

\begin{proof}
Each term of the partition formula has denominator $\prod_B|B|$. If $v=v_p(m)\ge1$ then $m\ge p^v\ge vp$, so $v_p(m)\le\lfloor m/p\rfloor$. Summing over blocks gives $\sum_B v_p(|B|)\le\sum_B\lfloor|B|/p\rfloor\le\lfloor k/p\rfloor$.
\end{proof}

\begin{observation}\label{p4:obs:dick}
For $1\le k\le300$, exact computation gives the following.
\begin{enumerate}[label=(\alph*),leftmargin=2em]
\item Equality $v_p(\den m_k)=\lfloor k/p\rfloor$ holds for $9822$ of the $10088$ pairs $(k,p)$ with $3\le p\le k$, $p$ prime.
\item $\log\den m_k$ agrees closely with $S_k:=\sum_{3\le p\le k}\lfloor k/p\rfloor\log p=k\log k+O(k)$. For example, $(\log\den m_k,S_k)=(347.3,347.3)$ at $k=120$ and $(1110.7,1112.6)$ at $k=300$.
\item Curiously, $v_2(\den m_k)=v_2(\lfloor k/2\rfloor+1)$ for all $k\le300$. This plays no role below.
\end{enumerate}
\end{observation}

\begin{proposition}\label{p4:prop:dickpos}
There is $\varepsilon(n)\to0$ such that every nonzero $Q\in\Z[u]$ of degree $n$ with $Q\ge0$ on $[1,\infty)$ satisfies
\[
\int_1^\infty\rho\,Q\ \ge\ \exp\bigl((1-\varepsilon(n))\,n\log n\bigr).
\]
\end{proposition}

\begin{proof}
For $L\ge1$,
\[
\int_1^\infty\rho Q\ \ge\ \rho(2L)\int_L^{2L}|Q|\ \ge\ \rho(2L)\,L(L/4)^n .
\]
The second inequality is the $L^1$ Chebyshev bound: monic polynomials of degree $n$ on an interval of length $L$ have $L^1$-norm at least $L(L/4)^n$, and $Q$ has leading coefficient of absolute value at least $1$. Take $L=n/(\log n)^2$ and use $\log\rho(u)=-u\log u\,(1+o(1))$~\cite{p4:dB}.
\end{proof}

\noindent\textbf{Where it breaks.} The forms are
\[
\int_1^\infty\rho\,Q=e^\gamma\,Y(Q)-\int_0^1Q,\qquad Y(Q)=\sum q_km_k .
\]
They are clean linear forms in $1$ and $e^\gamma$, but they meet three obstructions:
\begin{itemize}[leftmargin=2em]
\item Their denominators are bounded by $e^{k\log k+O(k)}$ (Proposition~\ref{p4:prop:dickden}), and numerically they are of that order (Observation~\ref{p4:obs:dick}). This is the same factorial wall as for the $\Gamma$-kernel.
\item By Proposition~\ref{p4:prop:dickpos}, positivity-certified forms are huge.
\item Most decisively, the direction is wrong. If $\gamma\in\Qbar$, then $e^\gamma$ is transcendental by Hermite--Lindemann, so an irrationality proof for $e^\gamma$ would say nothing about $\gamma$.
\end{itemize}
Buchstab's function $\omega(u)\to e^{-\gamma}$ (super-exponentially fast) has the same structure: its values $\omega(n)$ are periods of weight $\le n-2$. Discarded.

\subsection{Attempt 11: untying by two-point Hermite--Pad\'e}

By the tie theorem of~\cite{p4:K1}, small forms in $(1,\gamma,\delta)$ built from $E$-function Pad\'e data at $z=1$ approximate $e\Ein(1)=e\gamma+\delta$. The natural fix is to use data at $0$ and at $\infty$ simultaneously.

\begin{proposition}\label{p4:prop:untie}
Let $x>0$ be algebraic and $\lambda,\mu\in\Qbar$ with $\mu x>0$. Then
\[
e^{\lambda x}\Ein(\mu x)=e^{\lambda x}\bigl(\gamma+\log\mu x\bigr)+e^{(\lambda-\mu)x}\,G_0(\mu x),\qquad G_0(y):=e^{y}E_1(y)=\int_0^\infty\frac{e^{-t}\,dt}{y+t}.
\]
\end{proposition}

\begin{proof}
Use $\Ein(y)=\gamma+\log y+E_1(y)$ for $y>0$.
\end{proof}

\begin{remark}
Proposition~\ref{p4:prop:untie} suggests the following picture, which we have not made into a theorem. Values at a single point see $\gamma$ only through $\gamma+\log(\mu x)$, accompanied by the $\bD$-value $G_0(\mu x)$. To isolate $\gamma$ one needs points with $\prod(\mu_ix_i)^{c_i}=1$ and $\sum c_i\ne0$. The values $G_0(\mu_ix_i)$ then remain, and they become negligible only when $\mu_ix_i\to\infty$.
\end{remark}

\noindent\textbf{Where it breaks.} Untying always costs either a logarithm or a $\bD$-value. Removing both forces the moving-point schemes of Sweeney and Brent--McMillan~\cite{p4:BM}, or the dyadic levels of~\cite{p4:K5}. Their heights grow like $e^{cx^2\log x}$ (or doubly exponentially) against errors $e^{-x}$. Discarded.

\subsection{Attempt 12: the explicit formula and Baker's theory}

Since $-\zeta'/\zeta(s)=\frac1{s-1}-\gamma+O(s-1)$, the prime number theorem gives
\[
\sum_{n\le x}\frac{\Lambda(n)}n=\log x-\gamma+O\bigl(e^{-c\sqrt{\log x}}\bigr).
\]
Take $x\ge2$ an integer, so that $\log x$ is an integral combination of the $\log p$. If $\gamma=a/b$, the formula above is then a linear form in the $\pi(x)$ logarithms $\log p$ ($p\le x$) with rational coefficients of denominator $d_x=e^{(1+o(1))x}$.

\noindent\textbf{Where it breaks.} After clearing denominators the form has size about $d_x\,e^{-c\sqrt{\log x}}$, which tends to infinity, so it is not small. Even under the Riemann hypothesis the error $x^{-1/2+\varepsilon}$ is too large. Baker--W\"ustholz lower bounds~\cite{p4:BW} for $\pi(x)$ logarithms are astronomically smaller and therefore irrelevant. Discarded.

%======================================================================
\section{Synthesis}\label{p4:sec:synthesis}
%======================================================================

\begin{table}[ht]
\centering
\small
\begin{tabular}{rl>{\raggedright\arraybackslash}p{3.6cm}>{\raggedright\arraybackslash}p{5.4cm}}
\toprule
\# & route & where $\gamma$ enters & failure\\
\midrule
1 & $\Gamma$-kernel lattice & $\Gamma'(1)$ & no intrinsic analytic content; Pad\'e reintroduces the walls\\
2 & kernels $e^{-t^r}$ & $\psi(a/r)$ & $\Gamma(a/r)$, cyclotomic logs, $\pi\cot$\\
3 & $q\to1$ & $(1-q)\zeta_q(1)+\log(1-q)$ & $\log(1-q)$; $q$-integrality needs $1/q\in\Z$\\
4 & arithmetic holonomy & cancellation constants & reaches $\gamma$ only if $\gamma\in\Frac(\bG)$ (Prop.~\ref{p4:prop:cancel})\\
5 & Fano / Gamma conjecture & $\widehat\Gamma_X$ & Ap\'ery limits lie in $\Frac(\bG)$\\
6 & o-minimal counting & graph of $\psi$ & only $\asymp\log T$ points\\
7 & integer-valued functions & $\Gamma'+\gamma\Gamma$ & sees only $\gamma+\Z$\\
8 & Gauss map & $\int\{1/x\}$ & non-arithmetic spectrum\\
9 & Kronecker--Lerch & $\gamma_K$ & trades $\gamma$ for $\gamma_K$\\
10 & Dickman--Buchstab & $e^{\gamma}$ & wrong direction; factorial-type denominators\\
11 & two-point untying & $\gamma+\log x$ & log or $\bD$-value; moving-point heights\\
12 & explicit formula & Mertens & forms not small; $\pi(x)$ logarithms\\
\bottomrule
\end{tabular}
\caption{The twelve attempts and their obstructions.}\label{p4:tab:synth}
\end{table}

\medskip\noindent\textbf{What the attempts show.} The twelve attempts confirm, from new directions, the two walls isolated in~\cite{p4:K1,p4:K2,p4:K3,p4:K4,p4:K5}. They add two structural points.

First, in its present form the strongest modern irrationality method can only produce statements about constants in $\Frac(\bG)$ (Proposition~\ref{p4:prop:cancel}, Remark~\ref{p4:rem:blind}). It reaches $\gamma$ only if $\gamma\in\Frac(\bG)$, which is expected to be false but is open.

Second, every new bridge we found has one of four properties:
\begin{itemize}[leftmargin=2em]
\item it pays with logarithms, $\Gamma$-values or $\bD$-values (Attempts~2, 3, 9, 11, 12);
\item it points the wrong way, as with $e^\gamma$ (Attempt~10);
\item it is invisible to the arithmetic of $\gamma$ (Attempts~6, 7);
\item it sees $\gamma$ but gives no arithmetic handle on it (Attempt~8).
\end{itemize}

\medskip\noindent\textbf{What seems to be missing.} A proof that $\gamma\notin\Q$ appears to require a theorem of a new type. It would be an arithmetic statement about connection data at an \emph{irregular} singular point, equivalently about the rings $\bE$ and $\bD$. One example is Fischler and Rivoal's conjecture $\bE\cap\bD=\Qbar$, which would imply that $\gamma$ is transcendental~\cite{p4:FR24}. The work of Calegari--Dimitrov--Tang shows that holonomy theorems can be proved in the regular-singular ($G$) world, and Proposition~\ref{p4:prop:cancel} shows that this world reaches $\gamma$ only if $\gamma\in\Frac(\bG)$, which is not expected. Whether an irregular-singular analogue exists is, in our view, the real question behind the irrationality of $\gamma$.

\medskip\noindent\textbf{Practical consequence.} The concrete open targets closest to $\gamma$ that we can identify are those of~\cite{p4:K1}:
\begin{itemize}[leftmargin=2em]
\item the $\Qbar$-linear independence of $1,e,e\gamma$;
\item the implication $\delta\in\bE\Rightarrow\delta\in\Qbar$.
\end{itemize}
Both remain at the level of the main problem. For a genuinely new theorem in the near term, the holonomy methods of~\cite{p4:CDT} are better aimed at constants that are periods.

\begin{thebibliography}{VIII.99}

\bibitem[IV.1]{p4:And} Y.~Andr\'e, $G$-functions and Geometry, Aspects of Mathematics E13, Vieweg, Braunschweig, 1989.

\bibitem[IV.2]{p4:And00} Y.~Andr\'e, S\'eries Gevrey de type arithm\'etique I. Th\'eor\`emes de puret\'e et de dualit\'e, Ann. of Math. (2) 151 (2000), 705--740.

\bibitem[IV.3]{p4:Apt} A.~I. Aptekarev, On linear forms containing the Euler constant, arXiv:0902.1768 (2009); see also A.~I. Aptekarev (ed.), Rational approximants for Euler constant and recurrence relations, Sovrem. Probl. Mat.~9, Steklov Inst., Moscow, 2007.

\bibitem[IV.4]{p4:BW} A.~Baker and G.~W\"ustholz, Logarithmic forms and group varieties, J. reine angew. Math. 442 (1993), 19--62.

\bibitem[IV.5]{p4:Bez} J.-P.~B\'ezivin, Ind\'ependance lin\'eaire des valeurs des solutions transcendantes de certaines \'equations fonctionnelles, Manuscripta Math. 61 (1988), 103--129.

\bibitem[IV.6]{p4:Bor} P.~Borwein, On the irrationality of $\sum 1/(q^n+r)$, J. Number Theory 37 (1991), 253--259.

\bibitem[IV.7]{p4:BM} R.~P. Brent and E.~M. McMillan, Some new algorithms for high-precision computation of Euler's constant, Math. Comp. 34 (1980), 305--312.

\bibitem[IV.8]{p4:CDT} F.~Calegari, V.~Dimitrov and Y.~Tang, The linear independence of $1$, $\zeta(2)$, and $L(2,\chi_{-3})$, arXiv:2408.15403 (2024).

\bibitem[IV.9]{p4:CG} T.~Chatterjee and S.~Gun, The digamma function, Euler--Lehmer constants and their $p$-adic counterparts, Acta Arith. 162 (2014), 197--208.

\bibitem[IV.10]{p4:Chu} D.~V. Chudnovsky and G.~V. Chudnovsky, Applications of Pad\'e approximations to Diophantine inequalities in values of $G$-functions, in: Number Theory (New York, 1983--84), Lecture Notes in Math. 1135, Springer, 1985, 9--51.

\bibitem[IV.11]{p4:CS} S.~Chowla and A.~Selberg, On Epstein's zeta-function, J. reine angew. Math. 227 (1967), 86--110.

\bibitem[IV.12]{p4:dB} N.~G. de Bruijn, On the number of positive integers $\le x$ and free of prime factors $>y$, Indag. Math. 13 (1951), 50--60.

\bibitem[IV.13]{p4:vdDS} L.~van den Dries and P.~Speissegger, The field of reals with multisummable series and the exponential function, Proc. London Math. Soc. (3) 81 (2000), 513--565.

\bibitem[IV.14]{p4:FR14} S.~Fischler and T.~Rivoal, On the values of $G$-functions, Comment. Math. Helv. 89 (2014), 313--341.

\bibitem[IV.15]{p4:FR16} S.~Fischler and T.~Rivoal, Arithmetic theory of $E$-operators, J. \'Ec. polytech. Math. 3 (2016), 31--65.

\bibitem[IV.16]{p4:FR18} S.~Fischler and T.~Rivoal, Microsolutions of differential operators and values of arithmetic Gevrey series, Amer. J. Math. 140 (2018), 317--348.

\bibitem[IV.17]{p4:FR24} S.~Fischler and T.~Rivoal, Relations between values of arithmetic Gevrey series, and applications to values of the Gamma function, J. Number Theory 261 (2024), 36--54.

\bibitem[IV.18]{p4:FJ} J.~Fres\'an and P.~Jossen, Exponential motives, book manuscript.

\bibitem[IV.19]{p4:GGI} S.~Galkin, V.~Golyshev and H.~Iritani, Gamma classes and quantum cohomology of Fano manifolds: Gamma conjectures, Duke Math. J. 165 (2016), 2005--2077.

\bibitem[IV.20]{p4:HP} Kh.~Hessami Pilehrood and T.~Hessami Pilehrood, On a continued fraction expansion for Euler's constant, J. Number Theory 133 (2013), 769--786; arXiv:1010.1420.

\bibitem[IV.21]{p4:Iha} Y.~Ihara, On the Euler--Kronecker constants of global fields and primes with small norms, in: Algebraic Geometry and Number Theory, Progr. Math. 253, Birkh\"auser, 2006, 407--451.

\bibitem[IV.22]{p4:K1} KEIO, The two-world search for Euler's constant: shear invariance, collision worlds, forgeries, and the map of possible worlds, manuscript, 11 September 2026.

\bibitem[IV.23]{p4:K2} KEIO, Euler's constant: complete integrated research manuscript, cumulative edition, 10 September 2026.

\bibitem[IV.24]{p4:K3} KEIO, Arithmetic alternatives for Euler's constant: integrated research manuscript, 23 September 2026.

\bibitem[IV.25]{p4:K4} KEIO, Exponential-integral periods across arithmetic, Pad\'e, Hankel, reciprocal geometry, and Frobenius framing, version 6.3, 31 August 2026.

\bibitem[IV.26]{p4:K5} KEIO, The polyexponential trace-tower program: moving levels, Faber and inverse monodromy, positive spectra, and arithmetic barriers, version 3.1, 1 September 2026.

\bibitem[IV.27]{p4:KZ} M.~Kontsevich and D.~Zagier, Periods, in: Mathematics Unlimited---2001 and Beyond, Springer, 2001, 771--808.

\bibitem[IV.28]{p4:KW} N.~Kurokawa and M.~Wakayama, On $q$-analogues of the Euler constant and Lerch's limit formula, Proc. Amer. Math. Soc. 132 (2004), 935--943.

\bibitem[IV.29]{p4:Lag} J.~C. Lagarias, Euler's constant: Euler's work and modern developments, Bull. Amer. Math. Soc. 50 (2013), 527--628.

\bibitem[IV.30]{p4:Mah} K.~Mahler, Applications of a theorem by A.~B. Shidlovski, Proc. Roy. Soc. London Ser. A 305 (1968), 149--173.

\bibitem[IV.31]{p4:MS07} M.~R. Murty and N.~Saradha, Transcendental values of the digamma function, J. Number Theory 125 (2007), 298--318.

\bibitem[IV.32]{p4:MS10} M.~R. Murty and N.~Saradha, Euler--Lehmer constants and a conjecture of Erd\H{o}s, J. Number Theory 130 (2010), 2671--2682.

\bibitem[IV.33]{p4:PW} J.~Pila and A.~J. Wilkie, The rational points of a definable set, Duke Math. J. 133 (2006), 591--616.

\bibitem[IV.34]{p4:Riv09} T.~Rivoal, Rational approximations for values of derivatives of the Gamma function, Trans. Amer. Math. Soc. 361 (2009), 6115--6149.

\bibitem[IV.35]{p4:Riv12} T.~Rivoal, On the arithmetic nature of the values of the gamma function, Euler's constant, and Gompertz's constant, Michigan Math. J. 61 (2012), 239--254.

\bibitem[IV.36]{p4:Son} J.~Sondow, Criteria for irrationality of Euler's constant, Proc. Amer. Math. Soc. 131 (2003), 3335--3344.

\bibitem[IV.37]{p4:Ten} G.~Tenenbaum, Introduction to Analytic and Probabilistic Number Theory, 3rd ed., Grad. Stud. Math. 163, Amer. Math. Soc., 2015.

\bibitem[IV.38]{p4:VAW} W.~Van Assche and T.~Wolfs, Rational approximation of Euler's constant using multiple orthogonal polynomials, Combin. Number Theory 14 (2025), 141--162.

\bibitem[IV.39]{p4:Wir} E.~Wirsing, On the theorem of Gauss--Kusmin--L\'evy and a Frobenius-type theorem for function spaces, Acta Arith. 24 (1974), 507--528.

\bibitem[IV.40]{p4:Zud} W.~Zudilin, Heine's basic transform and a permutation group for $q$-harmonic series, Acta Arith. 111 (2004), 153--164.

\end{thebibliography}
% END PRESERVED SOURCE BODY 4
\phantomsection\label{part4:end}
\clearpage
\setcounter{section}{0}\setcounter{subsection}{0}
\renewcommand{\thesection}{\Roman{part}.\arabic{section}}
\renewcommand{\theHsection}{\arabic{part}.main.\arabic{section}}
\renewcommand{\refname}{References for Part V}
\fancyhead[L]{\small Euler\textquotesingle s constant: Part V}
\part{Monodromy, primitive lattices, and positive rational forms}
\label{part5:start}
\noindent\textbf{Source:} {\small\path{Keio_Gamma_Monodromy_Arithmetic_20261019.tex}}.

Original PDF: 30 pages. Mathematical content and source-level qualifications are retained.


% BEGIN PRESERVED SOURCE BODY 5
\begin{abstract}
We investigate several new routes relative to the supplied research
corpus toward the arithmetic of Euler's constant. The strongest
structural result extends phase embeddings of the ring of $E$-function
values by analytic continuation around the origin. Two different
windings eliminate all nonconstant rational denominators in a single
exponential-integral value, including positive algebraic arguments
previously left exceptional. Monodromy also supplies an additional
necessary condition for a value to be a unit and gives polynomial
intersection theorems for irreducible rational-parameter Kummer
values and integer-order Bessel values at nonzero algebraic points.

A finite-characteristic argument proves primitivity of integer Euler
jets and determines the exact image lattice of rank-three numerator
and denominator extraction. For equal upper decrements, its residual
index is proved to be at most exponential in the trajectory index,
despite the factorial size of the original determinant.
Independently, an integer affine Hankel
pencil selects the actual Euler constant by positivity, with a proved
finite-depth resolution cost. A rational differential reduction
produces positive pure gamma forms and allows precise obstructions
for completely monotone, Bernstein-positive, and fixed-order moving-pole
constructions. We also record the uniform arithmetic estimate needed
for a $q$-limit and a small-prime correction in the Diamond digamma
normalization. All unsuccessful transfers are identified explicitly.
No irrationality or transcendence proof for Euler's constant is obtained.
\end{abstract}
\noindent\textbf{Status.} AI-assisted research manuscript with full
proofs of its stated results and internal independent checks.
External peer review and an exhaustive priority search remain outstanding.
The results are not presented as ten resolutions of previously open
problems; that original objective has not been met.

\section{The target, the new results, and the proof standard}

Throughout,
\[
 \gamma=\lim_{n\to\infty}(H_n-\log n),\qquad
 H_n=\sum_{j=1}^n\frac1j,\qquad
 \delta=\int_0^\infty\frac{e^{-t}}{1+t}\,dt.
\]
These are the canonically normalized Euler and Euler--Gompertz
constants, not freely variable symbols in a formal comparison matrix.
Write $k=\overline{\mathbb Q}$ with its fixed embedding into
$\mathbb C$. An $E$-function is an entire series
$f(z)=\sum_{n\ge0}a_nz^n/n!$ with algebraic coefficients, satisfying
a linear differential equation over $k(z)$, for which all conjugates
of $a_n$ and a common denominator of $a_0,\ldots,a_n$ have exponential
growth bounds. The ring of its values at algebraic points is denoted
by $\mathcal E$. Rescaling the argument permits every such value to
be represented as $f(1)$.

The central arithmetic task is to produce integers $A_n,B_n$ with
\begin{equation}\label{p5:intro:integer-target}
 0<|A_n-B_n\gamma|\longrightarrow0.
\end{equation}
This condition is sufficient for irrationality: if $\gamma=a/b$,
multiplication by $b$ would give nonzero integers tending to zero.
An equivalent formulation for rational approximants $p_n/q_n$ is
\[
 0<q_n|\gamma-p_n/q_n|\longrightarrow0,
\]
where $q_n$ is the denominator after reduction. Convergence without
this arithmetic scale, or a result about functional transcendence,
does not supply \eqref{p5:intro:integer-target}. Transcendence requires
stronger information; no claim of that kind for $\gamma$ is made here.

\subsection{Results that survive the attempted proof}
\begin{center}
\small
\begin{tabular}{>{\raggedright\arraybackslash}p{.23\textwidth}>{\raggedright\arraybackslash}p{.66\textwidth}}
\toprule
Direction & Proved output and exact boundary\\
\midrule
Monodromy of phase maps &
$k(\operatorname{Ein}(a))\cap\mathcal E=k[\operatorname{Ein}(a)]$
for every nonzero algebraic $a$. The positive-argument denominator
exception is removed. The corresponding arithmetic classification
of $\gamma$ does not follow. The method also treats the Kummer and
Bessel families stated in Section~\ref{p5:km:section}.\\[4pt]
Finite-characteristic jets &
Primitivity and the exact cyclic image quotient of the integral
extraction matrix; for equal upper decrements its index divides
$\lcm(1,\ldots,S+1)^8$. Seed-dependent heights remain uncontrolled.\\[4pt]
Canonical positivity &
An integer affine moment pencil is positive at all depths exactly at
$c=\gamma$. Every finite depth still accepts an interval of nearby
values, with an explicit and an asymptotically stronger lower bound.\\[4pt]
Rational positive kernels &
Exact removal of auxiliary constants, a sharp $1/4$ barrier for one
family, and further certified obstructions after that family is
enlarged. None provides \eqref{p5:intro:integer-target}.\\[4pt]
Boundary and local transfers &
A quantified sufficient $q$-arithmetic bridge, whose hypothesis
remains unproved, and the precise small-prime exceptions in a
published $p$-adic statement.\\
\bottomrule
\end{tabular}
\end{center}

The first direction builds on Fischler--Rivoal's phase-map method
\cite{p5:FischlerRivoal2026Units}. The extension here continues its
component functions around the origin, while keeping the original
summation direction fixed. This distinction is essential. The
rank-three construction is motivated by the conservative matrix
fields of Kalisch et al.\ \cite{p5:Kalisch2026}, but the jet and lattice
proofs below are given explicitly.

\subsection{What is and is not inherited from the earlier manuscripts}
The supplied archive contains 31 TeX manuscripts. We use it as a
record of prior attempts, not as a reason to declare an entire area
of mathematics unavailable. In particular, the two-model argument
in \emph{The Two-World Search for Euler's Constant} explicitly removes
the analytic definitions of the comparison symbol. It is not a proof
that the actual number $\gamma$ can be changed while preserving all
of its defining properties. The later \emph{Arithmetic Alternatives
for Euler's Constant} already distinguishes these scopes.

For example, the exact real Gamma recurrence rules out nonzero real
twists $\Gamma_\tau(s)=\Gamma(s)e^{\tau(1-s)}$, since
$\Gamma_\tau(s+1)=se^{-\tau}\Gamma_\tau(s)$. Similarly, the recessive
boundary condition at infinity excludes a nonzero additive constant
in $E_1$. Fixing the normalization is therefore possible; extracting
small primitive integer forms from it is the unsolved part.

The previous scalar CMF convergence, reciprocal-Gamma sampling,
polynomial Laplace kernels, isolated-atom moment constructions, and
formal real/$p$-adic interpolation arguments are not counted again
as new discoveries. The present note is a complement to that corpus,
not a complete reproduction or merger of it. ``New'' in descriptions
of the research process means new in this round relative to the
inspected corpus. Publication priority has not been established.


% Standalone fragment; no preamble.  Original extension of cited inputs.
\section{Monodromy continuation of phase embeddings}
\label{p5:pm:section}

Put $k=\overline{\mathbb Q}$, let $\mathcal E$ denote the ring of
$E$-function values, and let
\[
 \mathcal R=\left\{\sum_{\rho\in k}c_\rho e^{\rho X}:
        c_\rho\in\mathbb C,\quad\text{finite support}\right\}.
\]
Fix $-\pi/2<\theta<\pi/2$.  We use the right-lateral convention
of \cite[Proposition~3.1]{p5:FischlerRivoal2026Units} at an exceptional
direction.  For an $E$-function $f$, write its canonical decomposition
as $f(z)=\sum_\rho f_\rho(z)e^{\rho z}$ near $z=1$.
The analytic input needed below is stronger than sectorial existence:
each $e^{\rho z}f_\rho(z)$ solves an $E$-operator whose only possible
finite singularity is zero
\cite[\S2, immediately before Theorem~2.3]{p5:FischlerRivoal2026Units}.
Thus every $f_\rho$ continues holomorphically along every path in
$\mathbb C^\times$.  Let $M^j$ denote continuation around a based
loop at $1$ of counterclockwise winding number $j\in\mathbb Z$.

\begin{theorem}[Continued phase embeddings]
\label{p5:pm:embedding}
For every integer $j$, the prescription
\begin{equation}\label{p5:pm:definition}
 \Phi_\theta^{[j]}(f(1))(X)
   =\sum_\rho(M^j f_\rho)(1)e^{\rho X}
\end{equation}
defines an injective $k$-algebra homomorphism
$\mathcal E\longrightarrow\mathcal R$.  It satisfies
\[
 \Phi_\theta^{[j]}(\xi)(1)=\xi,
 \qquad \Phi_\theta^{[j]}(e^a)=e^{aX}\quad(a\in k).
\]
The case $j=0$ is the original phase embedding.
\end{theorem}
\begin{proof}
Suppose $f(1)=g(1)$ for two $E$-functions.  Beukers' zero-removal
theorem supplies an $E$-function $h$ with $f-g=(z-1)h$.
Choose one sufficiently small positive perturbation of $\theta$
for the decompositions of $f,g,h$.  Uniqueness of these
decompositions gives
\[
 f_\rho-g_\rho=(z-1)h_\rho
\]
for each phase, inserting zero components when necessary.  Continue
this identity along the same based loop.  All continued components,
including $M^j h_\rho$, are holomorphic at the endpoint $1$.
Consequently $(M^j f_\rho)(1)=(M^j g_\rho)(1)$, proving that
\eqref{p5:pm:definition} is independent of the representative.

The component decomposition of a sum is the sum of the component
decompositions.  For a product $q=fg$ it is
\[
 q_\rho=\sum_{\alpha+\beta=\rho}f_\alpha g_\beta.
\]
These identities follow initially from uniqueness of the canonical
expansion, and remain true after analytic continuation.  They prove
additivity and multiplicativity.  Constant functions and pure
exponentials have constant components, so the stated normalizations
hold.  Finally, the continuation of the total sum is the continuation
of the entire function $f$.  Therefore
\[
 \sum_\rho(M^j f_\rho)(1)e^\rho=(M^j f)(1)=f(1),
\]
which proves evaluation compatibility and injectivity.
\end{proof}

This construction does not replace $\theta$ by $\theta+2\pi j$.
The direction remains admissible.  The continued components need
not form the canonical asymptotic decomposition in any prescribed
sector; only their analytic continuations are used.

\begin{lemma}[Two omitted values exclude denominators]
\label{p5:pm:two-centers}
Let $\xi\in\mathcal E$.  Suppose two injective $k$-algebra maps
$\Psi_i:\mathcal E\to\mathcal R$, $i=1,2$, have images
\[
 \Psi_i(\xi)=C_i+B_i e^{a_iX},\qquad
 B_i a_i\ne0,\quad C_1\ne C_2.
\]
Then $\xi$ is transcendental and
\[
 k(\xi)\cap\mathcal E=k[\xi],
 \qquad k[\xi]\cap\mathcal E^\times=k^\times.
\]
\end{lemma}
\begin{proof}
A nonconstant image excludes $\xi\in k$, because both maps fix $k$.
Take relatively prime $P,Q\in k[Y]$ with
$R=P(\xi)/Q(\xi)\in\mathcal E$.  Applying either map to the
identity $Q(\xi)R=P(\xi)$ gives
\[
 Q(C_i+B_i e^{a_iX})\Psi_i(R)(X)
       =P(C_i+B_i e^{a_iX}).
\]
If $\alpha$ is a root of $Q$ different from $C_i$, the exponential
attains $(\alpha-C_i)/B_i\ne0$.  Evaluating at such a point gives
$P(\alpha)=0$, contrary to coprimality.  Thus every root of $Q$
would have to equal both $C_1$ and $C_2$.  Hence $Q$ is constant.
The reverse inclusion is immediate.  Applying this conclusion to
$1/P(\xi)$ proves the assertion about units.
\end{proof}

\begin{theorem}[All nonzero algebraic exponential-integral arguments]
\label{p5:pm:ein}
For every $a\in k^\times$, let $\eta_a=\operatorname{Ein}(a)$, where
\[
 \operatorname{Ein}(z)=\sum_{n\ge1}\frac{(-1)^{n-1}z^n}{n\,n!}.
\]
Then $\eta_a$ is transcendental and
\begin{equation}\label{p5:pm:ein-intersection}
 k(\eta_a)\cap\mathcal E=k[\eta_a],
 \qquad k[\eta_a]\cap\mathcal E^\times=k^\times.
\end{equation}
In particular this includes every positive algebraic $a$, including
$a=1$, without any arithmetic hypothesis on $\gamma$.
\end{theorem}
\begin{proof}
The classical identity
$\operatorname{Ein}(w)=\gamma+\log w+E_1(w)$
\cite[\S6.2]{p5:DLMF} gives the two components of
$f(z)=\operatorname{Ein}(az)$:
\[
 f_0(z)=C+\log z,\qquad
 f_{-a}(z)=e^{az}\bigl(f(z)-C-\log z\bigr).
\]
Here $C=\gamma+\log a$ with the branch determined by a chosen
admissible canonical sum.  To see that these are the canonical
components, the formal tail for $f_{-a}$ has Borel transform
$1/(a+t)$; a nonsingular integration ray gives its sum in a sector
of opening greater than $\pi$.  The displayed identity then fixes
the compatible logarithm.  If the positive ray meets the pole,
take either sufficiently small lateral direction.

Put $B=f_{-a}(1)=e^a(\eta_a-C)$ and $\tau=2\pi i$.
Since $f$ and $e^{az}$ are single-valued whereas
$M^j\log z=\log z+j\tau$, Theorem~\ref{p5:pm:embedding} gives
\begin{equation}\label{p5:pm:ein-windings}
 \Phi_\theta^{[j]}(\eta_a)
   =C+j\tau+(B-j\tau e^a)e^{-aX}.
\end{equation}
The coefficient $B-j\tau e^a$ vanishes for at most one integer
$j$.  Choose two distinct integers for which it does not vanish.
Their constant terms differ.  Lemma~\ref{p5:pm:two-centers} proves
all the assertions.
\end{proof}

\begin{corollary}[Conditional restrictions on $\gamma$ and $\delta$]
\label{p5:pm:gamma-delta}
Let $\delta=e\operatorname{Ein}(1)-e\gamma$ be the Euler--Gompertz
constant.  Then
\begin{align*}
 \gamma\in k&\ \Longrightarrow\ 
 \delta\in\mathcal E\setminus k,\quad
 k(\delta)\cap\mathcal E=k[\delta],\quad
 k[\delta]\cap\mathcal E^\times=k^\times;\\
 \delta\in k&\ \Longrightarrow\ 
 \gamma\in\mathcal E\setminus k,\quad
 k(\gamma)\cap\mathcal E=k[\gamma],\quad
 k[\gamma]\cap\mathcal E^\times=k^\times.
\end{align*}
\end{corollary}
\begin{proof}
For $a=1$, choose direction zero and the principal branch.  Then
\[
 \Phi_0^{[j]}(\operatorname{Ein}(1))
    =\gamma+j\tau+(\delta-j\tau e)e^{-X}.
\]
If $\gamma\in k$, then $\delta\in\mathcal E$ and
\[
 \Phi_0^{[j]}(\delta)=\delta-j\tau e+j\tau e^X.
\]
The cases $j=1,2$ satisfy Lemma~\ref{p5:pm:two-centers}.
If instead $\delta\in k$, use $e^{-1}\in\mathcal E$ to obtain
$\gamma\in\mathcal E$ and
\[
 \Phi_0^{[j]}(\gamma)=\gamma+j\tau-j\tau e e^{-X}.
\]
Again $j=1,2$ prove the claim.
\end{proof}

\begin{theorem}[A monodromy obstruction to units]
\label{p5:pm:units}
Let $0\ne\xi=f(1)\in\mathcal E$.  For the finitely many component
solution spaces of $f$, choose an integer $L\ge1$ such that their
local monodromies raised to the $L$th power are unipotent.  Such an
$L$ exists because $E$-operators are regular singular at zero with
rational exponents
\cite[\S1, equation~(1.1)]{p5:FischlerRivoal2016Eoperators}.
Multiplication by the single-valued factor $e^{-\rho z}$ does not
change local monodromy, so this applies also to the component $f_\rho$.
If $\xi\in\mathcal E^\times$, there is one $\rho\in k$ for which
\begin{equation}\label{p5:pm:unit-period}
 \Phi_\theta^{[Lm]}(\xi)(X)=\xi e^{\rho(X-1)}
 \qquad(m\in\mathbb Z).
\end{equation}
More generally the images are periodic in the winding number,
with period $L$.
\end{theorem}
\begin{proof}
The units of $\mathcal R$ are exactly $c e^{\rho X}$ with
$c\ne0$.  Indeed, the finitely many phases in an element and a
putative inverse generate a finitely generated free abelian group.
The associated group algebra is a Laurent polynomial ring, since
distinct exponentials are linearly independent.  Its units are
scalar monomials.

For each phase, the coefficient
$A_\rho(m)=(M^{Lm}f_\rho)(1)$ is a polynomial in $m$:
if $M^L=I+N$ on a finite-dimensional solution space with $N$
nilpotent, then
$(I+N)^m=\sum_{s\ge0}\binom ms N^s$, a finite sum valid also for
negative integers $m$.  If $\xi$ is a unit, every continued image
is a unit of $\mathcal R$ and therefore has exactly one nonzero
phase.  Thus, for distinct phases $\rho,\sigma$, the polynomial
$A_\rho(m)A_\sigma(m)$ vanishes at every integer and is identically
zero.  At most one $A_\rho$ is a nonzero polynomial.  There is at
least one, because evaluation at $X=1$ is the nonzero number $\xi$.
For the surviving phase, evaluation gives
$A_\rho(m)e^\rho=\xi$ at every integer; hence $A_\rho$ is constant.
This proves \eqref{p5:pm:unit-period}.  Apply the same argument to
$M^j f_\rho$ for each residue class $j\bmod L$ to obtain periodicity.
\end{proof}

\begin{lemma}[A nonzero difference cannot be pulled back]
\label{p5:pm:no-common-image}
If $\xi\in\mathcal E$ and
$H=\Phi_\theta^{[j]}(\xi)-\Phi_\theta^{[\ell]}(\xi)\ne0$,
then $H$ belongs to none of the images
$\Phi_\vartheta^{[r]}(\mathcal E)$, for any admissible
$\vartheta$ and any integer $r$.
\end{lemma}
\begin{proof}
One has $H(1)=\xi-\xi=0$.  If
$H=\Phi_\vartheta^{[r]}(\zeta)$ for some $\zeta\in\mathcal E$,
evaluation gives $\zeta=0$, whence $H=0$, a contradiction.
\end{proof}

For example, with $\eta=\operatorname{Ein}(1)$ and $j\ne0$,
\[
 \Phi_0^{[j]}(\eta)-\Phi_0^{[0]}(\eta)
       =2\pi i j(1-e e^{-X})\ne0.
\]
Its coefficients do not become $E$-values merely because it is a
difference of two phase images.  There is no inverse phase map on
this difference.  Nor may one invoke an established separation
$\mathcal E\cap\mathcal G=k$: this is stated as a conjecture in
\cite[p.~38]{p5:FischlerRivoal2024Mixed}, where $\mathcal G$ is the
ring of analytically continued $G$-function values.

\paragraph{Scope of the result.}
Equation~\eqref{p5:pm:ein-intersection} is a restricted case of the
unit problem for $\mathcal E$, with a complete rational-intersection
statement.  The unit obstruction also tests the variation of
component coefficients under monodromy, rather than only their
initial phase support.  None of these statements proves that
$\gamma$ is irrational or transcendental.  In particular,
Corollary~\ref{p5:pm:gamma-delta} is consistent with algebraicity of
$\gamma$: it would make $\delta$ a transcendental nonunit of
$\mathcal E$, and no independent theorem places $1/\delta$ in
$\mathcal E$.
These extensions are new to this compilation; their priority in the
literature has not been established, and no literature-novelty claim
is made.

% Bibliography addition:
% \bibitem{FischlerRivoal2016Eoperators}
% S. Fischler and T. Rivoal, Arithmetic theory of E-operators,
% J. Ec. polytech. Math. 3 (2016), 31--65, DOI 10.5802/jep.28.
% \bibitem{FischlerRivoal2024Mixed}
% S. Fischler and T. Rivoal, Relations between values of arithmetic
% Gevrey series, and applications to values of the Gamma function,
% J. Number Theory 261 (2024), 36--54, DOI 10.1016/j.jnt.2024.02.016.
% https://www.imo.universite-paris-saclay.fr/~stephane.fischler/ssmixte.pdf

% Requires phase_monodromy.tex; no preamble.
\section{Rational intersections for irreducible Kummer values}
\label{p5:km:section}

The following is one application of the continued phase embeddings.
It extends the restricted parameter range of the Kummer nonunit
examples in \cite[Theorem~1.1 and \S4.3]{p5:FischlerRivoal2026Units}.
It is a theorem about a single parameterized family, not a collection
of independently solved conjectures.  Literature priority for this
extension has not been established.

\begin{theorem}[Kummer polynomial intersection]
\label{p5:km:intersection}
Let $a,b\in\mathbb Q$ satisfy
\[
 b\notin\mathbb Z_{\le0},\qquad a\notin\mathbb Z,
 \qquad b-a\notin\mathbb Z,
\]
Let $\alpha\in\overline{\mathbb Q}^{\times}$ and put
$\xi={}_1F_1(a;b;\alpha)=M(a,b,\alpha)$.  Then $\xi$ is
transcendental and, with $k=\overline{\mathbb Q}$,
\[
 k(\xi)\cap\mathcal E=k[\xi],\qquad
 k[\xi]\cap\mathcal E^\times=k^\times.
\]
In particular $M(a,b,\alpha)$ is not a unit of $\mathcal E$.
\end{theorem}
\begin{proof}
Set $f(z)=M(a,b,\alpha z)$, an $E$-function solving
\begin{equation}\label{p5:km:equation}
 zf''+(b-\alpha z)f'-\alpha af=0.
\end{equation}
Choose a sufficiently small admissible direction $\theta$ such that
$\alpha z$ lies in an upper or lower canonical sector, and denote
the corresponding sign by $\varepsilon\in\{1,-1\}$.  Its two canonical
phase components, including their exponentials, will be denoted by
$g_0$ and $g_\alpha$, so that $f=g_0+g_\alpha$ and
\[
 \Phi_\theta^{[m]}(\xi)
   =(M^m g_0)(1)+e^{-\alpha}(M^m g_\alpha)(1)e^{\alpha X}.
\]
Here $M$ denotes monodromy, not the Kummer function when used as an
operator.  The connection formula
\cite[13.2.41]{p5:DLMF}, with the compatible sign, gives
\begin{equation}\label{p5:km:components}
 g_0(z)=\frac{\Gamma(b)e^{\varepsilon\pi ia}}{\Gamma(b-a)}
             U(a,b,\alpha z),
 \qquad g_\alpha(z)=f(z)-g_0(z).
\end{equation}
All occurrences of $U$ and of $(\alpha z)^{1-b}$ below use the
same germ determined in that sector and continued to $z=1$.
The standard asymptotics of $U$ identify these as the canonical
phase-zero and phase-$\alpha$ components
\cite[13.7.2--13.7.3]{p5:DLMF}.  Both solve the same equation
\eqref{p5:km:equation}; this also follows by applying its differential
operator to the canonical decomposition and using uniqueness, since
summation commutes with differentiation and polynomial multiplication.
They are independent: both formal phase coefficients
$\Gamma(b)e^{\varepsilon\pi ia}/\Gamma(b-a)$ and $\Gamma(b)/\Gamma(a)$
are nonzero, so proportionality would contradict uniqueness of the
two distinct phases.  Equivalently, their Wronskian is a nonzero
multiple of $e^{\alpha z}z^{-b}$ \cite[13.2.38]{p5:DLMF}.

First, $\xi\ne0$.  Otherwise $\Phi_\theta(\xi)=0$ forces
$g_0(1)=g_\alpha(1)=0$ by independence of $1$ and $e^{\alpha X}$.
Two independent
solutions of a second-order equation cannot both vanish at its
ordinary point $1$, since their Wronskian there would vanish.

\emph{Case 1: $b$ is a positive integer.}
The exponents at zero are $0$ and $1-b$.  The holomorphic solution
space there is the line $\mathbb C f$, and local monodromy is
nontrivial unipotent.  For completeness, if $b\ge2$, a putative
logarithm-free second solution with leading term $w^{1-b}$ in the
coordinate $w=\alpha z$ would
have coefficients $c_0\ne0$ satisfying
\[
 n(n+1-b)c_n=(n+a-b)c_{n-1}.
\]
Up to $n=b-2$ all these coefficients are nonzero, whereas at
$n=b-1$ the equation would require
$0=(a-1)c_{b-2}\ne0$.  When $b=1$, a meromorphic solution cannot
have a pole, since the only exponent is zero, and the holomorphic
solution is again unique up to scalar.  Trivial monodromy would
give two meromorphic solutions at this regular singularity with
integer exponents, contradicting these facts.

Thus $\operatorname{im}(M-I)=\ker(M-I)=\mathbb C f$.
Since $g_0$ is not proportional to $f$, there is a scalar
$\kappa\ne0$ such that
\[
 (M-I)g_0=\kappa f,\qquad (M-I)g_\alpha=-\kappa f.
\]
Writing $C=g_0(1)$ yields
\begin{equation}\label{p5:km:integer-orbit}
 \Phi_\theta^{[m]}(\xi)
   =C+m\kappa\xi+e^{-\alpha}(\xi-C-m\kappa\xi)e^{\alpha X}.
\end{equation}
The constant terms are distinct because $\kappa\xi\ne0$, and at
most one exponential coefficient vanishes.  Lemma~\ref{p5:pm:two-centers}
proves all the conclusions in this case.

\emph{Case 2: $b$ is not an integer.}
Put
\[
 u=e^{2\pi ia},\quad v=e^{2\pi ib},\quad\lambda=v^{-1},
 \quad s_+=\frac{u-v}{1-v},\qquad s_-=u^{-1}s_+,
 \qquad s=s_\varepsilon.
\]
These numbers are algebraic, $\lambda\ne1$, and
$s\notin\{0,1\}$ for either sign: the excluded values would give respectively
$a-b\in\mathbb Z$ or $a\in\mathbb Z$.
The local connection formula \cite[13.2.42]{p5:DLMF}, combined with
\eqref{p5:km:components} and the gamma reflection formula, gives
\[
 g_0=s f+t h,\qquad
 h(z)=(\alpha z)^{1-b}M(a-b+1,2-b,\alpha z)
\]
for a nonzero constant $t$.  Indeed the coefficient of $f$ is
\[
 \frac{\Gamma(b)\Gamma(1-b)e^{\varepsilon\pi ia}}
      {\Gamma(b-a)\Gamma(1+a-b)}
   =\frac{e^{\varepsilon\pi ia}\sin\pi(b-a)}{\sin\pi b}
   =s_\varepsilon.
\]
Since $Mf=f$ and $Mh=\lambda h$, the phase images are
\begin{equation}\label{p5:km:finite-orbit}
 \Phi_\theta^{[m]}(\xi)=C_m+e^{-\alpha}(\xi-C_m)e^{\alpha X},
 \qquad C_m=s\xi+\lambda^m(C_0-s\xi).
\end{equation}
This formula first proves $\xi\notin k$.  If $\xi$ were algebraic,
every image would be the constant $\xi$, so $C_m=\xi$ for all
$m$.  Comparing $m=0,1$ gives $(1-\lambda)(1-s)\xi=0$,
contrary to the established nonvanishing factors.

Now take relatively prime $P,Q\in k[Y]$ with
$P(\xi)/Q(\xi)\in\mathcal E$.  The denominator argument of
Lemma~\ref{p5:pm:two-centers} says that any root of $Q$ must equal
every $C_m$ for which $C_m\ne\xi$.
If $C_0=s\xi$, all the $C_m$ equal the transcendental number
$s\xi\ne\xi$; hence $Q$ has no algebraic root and is constant.
Otherwise the $C_m$ form an orbit of $q$ distinct values, where
$q\ge2$ is the order of the root of unity $\lambda$.
For $q\ge3$, at least two distinct values differ from $\xi$,
so $Q$ is constant by the two-center argument.
For $q=2$, the same holds unless one orbit value equals $\xi$.
In that last case the other equals $(2s-1)\xi$.  But $v=-1$ gives
$2s-1=u^\varepsilon\in k^\times$, so this other value is transcendental and
cannot be a root of $Q\in k[Y]$.  Again $Q$ is constant.
This proves the field intersection.  The statement about units
follows by applying it to reciprocals of nonconstant polynomials.
\end{proof}

\begin{corollary}[Integer-order Bessel values]
\label{p5:km:bessel}
For $n\in\mathbb Z_{\ge0}$ and
$t\in\overline{\mathbb Q}^{\times}$, let $\beta=J_n(t)$.
Then $\beta$ is transcendental and
\[
 k(\beta)\cap\mathcal E=k[\beta],\qquad
 k[\beta]\cap\mathcal E^\times=k^\times.
\]
\end{corollary}
\begin{proof}
The classical Kummer--Bessel identity
\cite[13.6.9]{p5:DLMF}, together with $I_n(it)=i^nJ_n(t)$, gives
\[
 \xi=M\left(n+\tfrac12,2n+1,2it\right)
       =n!e^{it}(t/2)^{-n}J_n(t).
\]
Put $d=(t/2)^n/n!\in k^\times$.  The integer-$b$ case of the
proof above, with $\alpha=2it$, gives affine coefficients
$C_m=C_0+m\kappa\xi$ with $\kappa\xi\ne0$.  Multiplicativity
and the exponential normalization of the continued embeddings give
\[
 \Phi_\theta^{[m]}(\beta)
   =d C_m e^{-itX}
       +d e^{-2it}(\xi-C_m)e^{itX}.
\]
Choose $m$ avoiding the at most two zeros of its two coefficients.
Then this image is $A e^{itX}+B e^{-itX}$ with $AB\ne0$.
It is surjective onto $\mathbb C$: for any $w\in\mathbb C$,
the polynomial $AZ^2-wZ+B$ has a nonzero complex root, and
$Z=e^{itX}$ attains every nonzero complex number.
The image is nonconstant, so $\beta$ is transcendental.  If
$P(\beta)/Q(\beta)\in\mathcal E$ in reduced form, apply this
embedding to the product identity and evaluate at a preimage of
any root of $Q$.  Coprimality is contradicted unless $Q$ is constant.
The assertion about units follows as before.
\end{proof}

The transcendence of these irreducible Kummer values is classical;
the content here is the rational-intersection conclusion and its
unit consequence.  The parameter exclusions are material.  For
example, $M(b,b,\alpha)=e^\alpha$ is a unit, and terminating cases can give
algebraic values.  No conclusion about the irrationality or
transcendence of $\gamma$ follows from this application.

\section{Primitive elementary jets and the exact CMF image lattice}
\label{p5:sec:primitive-cmf-jet}

The denominator of a raw vector is not the denominator of a reduced
approximant. We examine the integer matrix before selecting a seed.
The following finite-characteristic argument determines a universal
content exactly. It leaves seed-specific gcds as a separate problem.

\begin{lemma}[Primitive jets of integer Euler products]
\label{p5:lem:primitive-euler-jet}
Let $\theta=z\,d/dz$, $D=\theta-z$, and
$P(X)=\pm\prod_{\nu=1}^m(X-a_\nu)$ with $a_\nu\in\mathbb Z$;
the empty product is allowed. Put $R(z)=P(D)1\in\mathbb Z[z]$.
Suppose $e^{-z}R(z)$ satisfies a differential equation
\[
 \left(\theta^r+\sum_{j=0}^{r-1}c_j(z)\theta^j\right)
       (e^{-z}R(z))=0,\qquad c_j\in\mathbb Z[z].
\]
Then the integer vector
\[
 \bigl(R(1),(DR)(1),\ldots,(D^{r-1}R)(1)\bigr)
\]
is primitive: the gcd of its coordinates is one.
\end{lemma}
\begin{proof}
Fix a prime $\ell$. Work in the $\ell$-dimensional vector space
$V=\mathbb F_\ell[z]/(z^\ell-1)$, which is also
$\mathbb F_\ell[z]/((z-1)^\ell)$. Differentiation is well defined on
this quotient and its $\ell$th power is zero. With $T=\partial_z-1$,
the relation $Tz-zT=1$ gives
\[
 D(D-1)\cdots(D-\ell+1)=z^\ell T^\ell.
\]
In characteristic $\ell$, the left side is $D^\ell-D$, whereas
$T^\ell=\partial_z^\ell-1=-1$ on $V$. Consequently
\[
 D^\ell-D=-1\qquad\text{on }V.
\]
Each $D-a$, $a\in\mathbb F_\ell$, is invertible: the negative product
of the other $\ell-1$ factors is its inverse. Therefore $P(D)$ is
invertible, and $R=P(D)1$ is nonzero in $V$.

Suppose that the displayed $r$ jet coordinates were all divisible by
$\ell$. The operators $D^j$ have leading ordinary derivative term
$z^j\partial_z^j$; the lower terms have integer polynomial coefficients.
At $z=1$ this is a triangular change of derivative coordinates with
diagonal entries one. Thus $R^{(j)}(1)=0$ modulo $\ell$ for $j<r$.
If $\ell\le r$, the first $\ell$ derivatives already imply
$R=0$ in $V$, since $0!,\ldots,(\ell-1)!$ are invertible.

If $\ell>r$, conjugating the given differential equation by $e^{-z}$
gives
\[
 \left(D^r+\sum_{j<r}c_j(z)D^j\right)R=0.
\]
Its coefficient of $\partial_z^r$ is $z^r$, equal to one at $z=1$.
Differentiate this identity successively zero through $\ell-r-1$
times, and evaluate at one. Induction forces
$R^{(r)}(1),\ldots,R^{(\ell-1)}(1)$ to vanish modulo $\ell$:
the highest derivative has coefficient one, and all other derivatives
in that identity have already vanished. Again $R=0$ in $V$, a
contradiction. This holds for every prime, proving primitivity.
\end{proof}

\begin{corollary}[All integral rank-three CMF trajectories]
\label{p5:cor:primitive-cmf-jet}
Let $A_1,A_2,b_1,b_2,b_3$ be nonnegative integers. Starting at the zero
parameter point, let $\Phi_n$ be the row-state transition matrix for
the trajectory $(-A_1,-A_2;b_1,b_2,b_3)$, with the scalar contiguous
operator normalization. Then $\Phi_n$ is a nonsingular integer matrix,
and
\[
 E_n=(1,-1,0)\Phi_n\in\mathbb Z^3
 \quad\text{satisfies}\quad \gcd(E_{n,0},E_{n,1},E_{n,2})=1.
\]
Moreover
\begin{equation}\label{p5:eq:primitive-cmf-det}
 \Delta_n:=\det\Phi_n
 =\prod_{i=1}^2\prod_{j=1}^3((A_i+b_j)n)!>0.
\end{equation}
\end{corollary}
\begin{proof}
The product of contiguous scalar operators is
\[
 P_n(\theta)=(\theta+1)_{A_1n}(\theta+1)_{A_2n}
                 \prod_{j=1}^3(-\theta)_{b_jn}.
\]
These are commuting polynomials in $\theta$, each a product of
integer-linear factors with leading coefficient $1$ or $-1$.
The transformed elementary solution is
$f_n(z)=e P_n(\theta)e^{-z}=e^{1-z}R_n(z)$.
Its state at one is $E_n$, since the initial state of $e^{1-z}$
is $(1,-1,0)$. Its final scalar equation is
\[
 \left[\prod_{j=1}^3(\theta-b_jn)
       +z\prod_{i=1}^2(\theta+A_in+1)\right]f_n=0.
\]
This equation is monic of degree three in $\theta$ with integer
polynomial coefficients. The preceding lemma proves the gcd assertion.

For clarity, the determinant formula can be obtained without the gcd
argument. A single upper decrement at parameters $a_i,b_j$ has
determinant $\prod_j(b_j-a_i+1)$, and a single lower increment has
determinant $\prod_i(b_j-a_i+1)$. These follow from the rank-three
companion polynomial at $z=1$, evaluated at $a_i-1$ or $b_j$.
Take all upper decrements first and then all lower increments.
For each pair $(i,j)$ the factors multiply to $((A_i+b_j)n)!$.
All factors are positive. Contiguous matrices have integer entries
because the old monic equation reduces powers of $\theta$ using
integer polynomial coefficients at one. Thus every step is defined,
integral, and nonsingular, and multiplication proves the formula.
\end{proof}

\begin{lemma}[Content and image of a primitive complementary pair]
\label{p5:lem:cmf-content}
Let $\Phi\in M_3(\mathbb Z)$ have determinant $\Delta\ne0$.
Let $p,q,u\in\mathbb Z^3$ satisfy $p\times q=u$, with $u$ primitive,
and suppose $E=u^T\Phi$ is primitive. Put
\[
 H=\operatorname{adj}(\Phi)^T,\qquad
 h=\gcd\{H_{ij}\},\qquad
 B=\begin{pmatrix}p^TH\\q^TH\end{pmatrix}.
\]
Then the gcd of the six entries of $B$ is exactly $h$.
The matrix $\widehat B=B/h$ has image lattice
$\Omega=\widehat B\mathbb Z^3\subset\mathbb Z^2$ with
\begin{equation}\label{p5:eq:cmf-image-index}
 [\mathbb Z^2:\Omega]=\frac{|\Delta|}{h^2},\qquad
 \mathbb Z^2/\Omega\simeq\mathbb Z/(|\Delta|/h^2)\mathbb Z,
 \qquad \ker_{\mathbb Z}\widehat B=\mathbb Z E^T.
\end{equation}
\end{lemma}
\begin{proof}
Choose $S,T\in\mathrm{SL}_3(\mathbb Z)$ so that the first row of $S$
is $u^T$ and $ET=(1,0,0)$. Then
\[
 \Psi=S\Phi T=\begin{pmatrix}1&0&0\\a&b&c\\d&e&f\end{pmatrix},
 \qquad
 \operatorname{adj}(\Psi)^T=
 \begin{pmatrix}bf-ce&cd-af&ae-bd\\0&f&-e\\0&-c&b\end{pmatrix}.
\]
The gcd of all nine entries of this adjugate is
$\gcd(b,c,e,f)$, which is also the gcd of its bottom two rows.
Unimodular left and right transformations preserve the gcd of entries.
Since $p,q$ form a basis of the saturated lattice $u^\perp\cap\mathbb Z^3$
(their cross product is primitive), their images under $S$ form a
unimodular basis of $\{0\}\times\mathbb Z^2$.
The identity
$H=S^T\operatorname{adj}(\Psi)^T T^T$ therefore proves the asserted
equality of contents.

The cross-product transformation law gives
\[
 (p^TH)\times(q^TH)=\Delta\,u^T\Phi=\Delta E.
\]
Because $E$ is primitive, the gcd of the three maximal minors of
$B$ is $|\Delta|$; that of $\widehat B$ is $|\Delta|/h^2$.
The elementary Smith normal form formulas identify this gcd with the
image index. The first invariant factor of $\widehat B$ is one,
since its entries are relatively prime, proving cyclicity. The
nonzero cross product identifies its rational kernel as the line
$\mathbb QE^T$; primitivity of $E$ gives the integer kernel.
\end{proof}

\begin{theorem}[No further universal cancellation in gamma extraction]
\label{p5:thm:cmf-universal-content}
In Corollary~\ref{p5:cor:primitive-cmf-jet}, take
$p=(1,1,0)^T$, $q=(1,1,1)^T$, and let $h_n$ be the content of
$\operatorname{adj}(\Phi_n)$. After removing this matrix content,
the two rows extracting the numerator and denominator have joint
content exactly one. Their image is a rank-two integer lattice with
cyclic quotient of order $\Delta_n/h_n^2$.
If $A_1=A_2=A$, set
\[
 N_n=\prod_{j=1}^3((A+b_j)n)!.
\]
Then $h_n\mid N_n$, and the image index is the perfect square
$(N_n/h_n)^2$.
\end{theorem}
\begin{proof}
Here $p\times q=(1,-1,0)^T$, and its transported row is primitive by
Corollary~\ref{p5:cor:primitive-cmf-jet}. Apply the preceding lemma.
When $A_1=A_2$, formula~\eqref{p5:eq:primitive-cmf-det} is
$\Delta_n=N_n^2$. Since $h_n^2\mid\Delta_n$, prime valuations imply
$h_n\mid N_n$.
\end{proof}

This theorem concerns cancellation common to every seed. A particular
$v\in\mathbb Z^3$ can still produce a pair $\widehat B_nv$ with a
large gcd. Thus the statement does not control reduced scalar heights,
does not exclude a specially designed varying-seed construction, and
does not prove an irrationality measure for $\gamma$.

\subsection{An exponential bound for the quotient index}

Write $d_M=\operatorname{lcm}(1,\ldots,M)$, with $d_0=1$.
The square index above admits a substantially smaller arithmetic bound
than its defining factorial determinant suggests.

\begin{theorem}[Only exponential index remains]
\label{p5:thm:cmf-lcm-index}
In Theorem~\ref{p5:thm:cmf-universal-content}, assume $A_1=A_2=A$ and put
\[
 m_j=(A+b_j)n,\qquad S=m_1+m_2+m_3,\qquad
 N=m_1!m_2!m_3!,\qquad J=N/h_n.
\]
Then
\begin{equation}\label{p5:eq:cmf-lcm-index}
 J\mid d_{S+1}^{\,4},\qquad
 [\mathbb Z^2:\Omega_n]=J^2\mid d_{S+1}^{\,8}.
\end{equation}
For a fixed trajectory, the logarithm of the image index is therefore
$O(n)$.
\end{theorem}
\begin{proof}
We first treat $A=0$. Let
$P(X)=\prod_{j=1}^3(-X)_{m_j}$, of degree $S$.
In the triangular source system
\[
 \theta f=-zf+g,\qquad \theta g=h,\qquad \theta h=0,
\]
there are polynomials $R,U,V\in\mathbb Z[z]$ such that
\begin{equation}\label{p5:eq:cmf-triangular-action}
 P(\theta)f=R(z)f+U(z)g+V(z)h.
\end{equation}
Indeed one application of $\theta$ sends the coefficient triple
$(R,U,V)$ to
\[
 (\theta R-zR,\ R+\theta U,\ U+\theta V),
\]
so this assertion follows by applying the integer-linear factors of $P$.
The polynomial $R=e^zP(\theta)e^{-z}$ has degree $S$. Write
$R(z)=\sum_{k=0}^S r_kz^k$. Expanding the exponential gives
\[
 k!r_k=\sum_{j=0}^k(-1)^j\binom kjP(j)\in N\mathbb Z.
\]
The divisibility holds because
$P(j)/N=\prod_\nu(-1)^{m_\nu}\binom{j}{m_\nu}$ is an integer.

We next identify $U,V$ and their denominators after division by $N$.
All calculations in this paragraph are formal at infinity; no
convergence assertion is used. The Laurent series
\[
 F(z)=\sum_{k\ge0}k!z^{-k-1},\qquad
 G(z)=\sum_{k\ge0}k!z^{-k-1}(\log z-H_k)
\]
satisfy $(\theta+z)F=1$ and $(\theta+z)G=\log z$.
Apply \eqref{p5:eq:cmf-triangular-action} first with
$(f,g,h)=(F,1,0)$ and then with $(G,\log z,1)$.
Both $P(\theta)F$ and $P(\theta)G$ have only strictly negative
powers of $z$, with logarithmic coefficients in the second case.
Taking polynomial parts therefore gives, for $0\le i<S$,
\begin{equation}\label{p5:eq:cmf-U-V}
 U_i=-\sum_{k=i+1}^S r_k(k-i-1)!,\qquad
 V_i=\sum_{k=i+1}^S r_k(k-i-1)!H_{k-i-1}.
\end{equation}
For $0\le i<k$, the integer
$k\binom{k-1}{i}$ divides $d_k$. To see this, its valuation at a
prime $\ell$ is
\[
 \sum_{a\ge1}\left(
  \left\lfloor\frac{k}{\ell^a}\right\rfloor
 -\left\lfloor\frac{i}{\ell^a}\right\rfloor
 -\left\lfloor\frac{k-i-1}{\ell^a}\right\rfloor\right).
\]
Every summand is zero or one, and it is zero for $\ell^a>k$.
Since
$i!(k-i-1)!/k!=1/[k\binom{k-1}{i}]$, and the denominator of
$H_j$ divides $d_S$ for $j\le S$, formula~\eqref{p5:eq:cmf-U-V}
implies
\begin{equation}\label{p5:eq:cmf-weighted-integrality}
 k!r_k\in N\mathbb Z,\qquad
 i!U_i\in \frac{N}{d_S}\mathbb Z,\qquad
 i!V_i\in \frac{N}{d_S^2}\mathbb Z.
\end{equation}

We must bound the degrees of the actual extraction polynomials, not
of their uncancelled products. Set $D=\theta-z$ and define three
polynomial rows
\begin{align*}
 \mathcal R&=(R,DR,D^2R),\\
 \mathcal U&=(U,R+\theta U,2\theta R-zR+\theta^2U),\\
 \mathcal V&=(V,U+\theta V,R+2\theta U+\theta^2V).
\end{align*}
Their evaluations at one are the transition rows in the source
coordinates $(f,g,h)$. The two numerator--denominator extraction
rows of the adjugate are, up to a sign,
\[
 \mathcal W_1(1),\ \mathcal W_2(1),\qquad
 \mathcal W_1=\mathcal R\times\mathcal U,
 \mathcal W_2=\mathcal R\times\mathcal V.
\]
For example the original state at one is
$(f,\theta f,\theta^2f)=(f,g,h)
 \left(\begin{smallmatrix}1&-1&0\\0&1&-1\\0&0&1\end{smallmatrix}\right)$;
the extraction vectors select $g,h$, respectively. This confirms
the row conventions and the absence of an extra scalar factor.

Here is a formal proof that every component of $\mathcal W_1$ and
$\mathcal W_2$ has degree at most $S+1$.
The jets of $P(\theta)F=RF+U$ are
$F\mathcal R+\mathcal U$. The jets of
$P(\theta)G=RG+U\log z+V$ are
$G\mathcal R+(\log z)\mathcal U+\mathcal V$.
Taking their cross products with $\mathcal R$ gives the exact identities
\[
 \mathcal R\times\operatorname{Jet}(P(\theta)F)=\mathcal W_1,\qquad
 \mathcal R\times\operatorname{Jet}(P(\theta)G)
   =(\log z)\mathcal W_1+\mathcal W_2.
\]
But each of the first two formal jets has strictly negative powers
of $z$ (and possibly one power of $\log z$), whereas
$\mathcal R$ has degree at most $S+2$. Consequently
\[
 \deg\mathcal W_1\le S+1,\qquad
 \deg\mathcal W_2\le S+1;
\]
for the second assertion compare the constant and logarithmic
coefficients separately. This is a statement about formal Laurent
series equal to polynomials, so it establishes exact degree bounds.

Multiplication of polynomials whose factorial-weighted coefficients
are integral preserves that property, since the coefficient formula
uses binomial coefficients. The same is true of $\theta$ and of
multiplication by $z$. Hence \eqref{p5:eq:cmf-weighted-integrality} gives
\[
 k![z^k]\mathcal W_i\in\frac{N^2}{d_S^2}\mathbb Z^3
 \qquad(i=1,2).
\]
Using their degree bounds, evaluation at one yields
\[
 v_\ell(h_n)\ge
 2v_\ell(N)-v_\ell((S+1)!)-2v_\ell(d_S)
\]
for every prime $\ell$; here the equality of extraction content and
adjugate content is essential. Thus
\begin{align*}
 v_\ell(J)
 &\le v_\ell((S+1)!/N)+2v_\ell(d_S)\\
 &\le 4\left\lfloor\log_\ell(S+1)\right\rfloor.
\end{align*}
For the last inequality, write the valuation directly as
\[
 v_\ell((S+1)!/N)=\sum_{a\ge1}\left(
 \left\lfloor\frac{S+1}{\ell^a}\right\rfloor
 -\sum_{j=1}^3\left\lfloor\frac{m_j}{\ell^a}\right\rfloor\right).
\]
Every summand is zero, one, or two: the sum of three residues
modulo $\ell^a$, plus one, is at most $3\ell^a-2<3\ell^a$.
There are at most $\lfloor\log_\ell(S+1)\rfloor$ nonzero
summands, while $v_\ell(d_S)\le\lfloor\log_\ell(S+1)\rfloor$.
The case $S=0$ is immediate. This proves the divisibility for $A=0$.

Finally equal upper decrements reduce to this case by an integral
unimodular gauge. Put $a=An$. On the initial solution space,
\[
 (-\theta)_a^3=z^a(\theta+1)_a^2,
\]
as follows inductively from the initial equation
$(-\theta)^3f=z(\theta+1)^2f$ and $\theta z=z(\theta+1)$.
Commuting the remaining Euler factors gives
\[
 \prod_j(-\theta)_{a+b_jn}f
   =z^a(\theta+1)_a^2\prod_j(-\theta)_{b_jn}f.
\]
At $z=1$, multiplication by $z^a$ acts on jets through
\[
 U_a=\begin{pmatrix}1&a&a^2\\0&1&2a\\0&0&1\end{pmatrix}
 \in\mathrm{SL}_3(\mathbb Z).
\]
It leaves adjugate content and the image lattice invariant, proving
the claimed reduction. The classical bound $\log d_M=O(M)$ then
gives the last assertion.
\end{proof}

The factor $4$ in \eqref{p5:eq:cmf-lcm-index} is a proved upper bound,
not an optimality claim. Exact experiments suggest the stronger
$J\mid d_{\max m_j}$, but that sharper statement is not established
here. Even an exponential image index does not supply small-error
lifts of a primitive target; controlling such lifts remains necessary
for an irrationality argument.

\section{A canonical integer moment pencil and its finite resolution}
\label{p5:sec:canonical-pencil}

The exact Gamma functional equation and its real positive log-convex
normalization fix the analytic function.  An artificial shear of a
comparison constant need not respect these data.  We therefore keep
the canonical normalization and ask whether positivity supplies an
integer certificate.  The construction below does supply such a
certificate system.  It also permits a quantitative analysis of what
a finite part of the system cannot yet distinguish.  No irrationality
conclusion is asserted.

Put $H_0=0$ and, for $k\geq0$, define
\begin{equation}\label{p5:eq:canonical-moments}
 \mu_k(c)=k!\{k(H_k-c)+1\}=a_k-b_kc,
 \qquad a_k,b_k\in\mathbb Z.
\end{equation}
Here $a_k=k!(kH_k+1)$ and $b_k=k\,k!$.  For a real polynomial
$Q(X)=\sum_iq_iX^i$ set
\[
 \mathcal F_c(Q)=\sum_{i,j}q_iq_j\mu_{i+j}(c),\qquad
 M_N(c)=(\mu_{i+j}(c))_{0\leq i,j<N}.
\]
Thus $\mathcal F_c(Q)=A_Q-B_Qc$ with $A_Q,B_Q\in\mathbb Z$
when $Q\in\mathbb Z[X]$.  Differentiation of the canonical Gamma
integral gives
\[
 \int_0^\infty x^k\log x\,e^{-x}\,dx=k!(H_k-\gamma).
\]
Subtracting consecutive moments proves the identity
\begin{equation}\label{p5:eq:canonical-integral}
 \mathcal F_c(Q)=\int_0^\infty Q(x)^2
          (x-1)(\log x+\gamma-c)e^{-x}\,dx.
\end{equation}
The coefficient construction in \eqref{p5:eq:canonical-moments} itself
uses no value of $\gamma$.  Its identification with
\eqref{p5:eq:canonical-integral} is where the canonical analytic
normalization enters.

\begin{theorem}[Exact positivity parameter]\label{p5:thm:canonical-positivity}
For a real number $c$, the following conditions are equivalent:
\begin{enumerate}
 \item $c=\gamma$;
 \item $\mathcal F_c(Q)\geq0$ for every $Q\in\mathbb R[X]$;
 \item $\mathcal F_c(Q)\geq0$ for every $Q\in\mathbb Z[X]$;
 \item $M_N(c)$ is positive semidefinite for every $N\geq1$.
\end{enumerate}
At $c=\gamma$, every nonzero polynomial gives strict positivity.
For $c=a/b$ with integers $a,b$ and $b>0$, a negative certificate is
the entirely integer inequality $bA_Q-aB_Q<0$.
\end{theorem}
\begin{proof}
The weight $(x-1)\log x\,e^{-x}$ is positive for all $x>0$
except $x=1$.  This proves strict positivity at $c=\gamma$.
If $c\ne\gamma$, the weight
\[
 w_c(x)=(x-1)(\log x+\gamma-c)e^{-x}
\]
is negative on the interval with endpoints $1$ and $e^{c-\gamma}$.
We show that polynomials are dense in $L^2(|w_c(x)|\,dx)$.
Suppose $g$ is orthogonal to every polynomial.  By
Cauchy--Schwarz, the integral
\[
 T(z)=\int_0^\infty g(x)|w_c(x)|e^{zx}\,dx
\]
is analytic in $\operatorname{Re}z<1/2$: on compact subsets of
that half-plane the necessary polynomially weighted integrals
converge absolutely.  All derivatives of $T$ at zero vanish by
orthogonality.  The identity theorem makes $T$ zero throughout
the half-plane.  The Fourier transform of the finite signed or
complex measure $g(x)|w_c(x)|\,dx$ is consequently zero.
Uniqueness of Fourier transforms of finite measures gives $g=0$
almost everywhere for $|w_c|\,dx$, proving density.

Choose a nonzero continuous function supported in a compact
subinterval where $w_c<0$, and approximate it in this $L^2$ space
by real polynomials $Q_j$.  The inequality
\[
 \int |Q_j^2-f^2|\,|w_c|
 \leq \|Q_j-f\|_{L^2(|w_c|)}
       (\|Q_j\|_{L^2(|w_c|)}+\|f\|_{L^2(|w_c|)})
\]
shows that $\mathcal F_c(Q_j)<0$ for some $j$.
Continuity in the finitely many polynomial coefficients allows
rational coefficients instead.  Clearing denominators gives a
negative integer-polynomial certificate.  This establishes the
equivalence, since the matrix condition is precisely positivity
on each finite-dimensional polynomial space.
\end{proof}

\begin{remark}[The remaining arithmetic step]\label{p5:rem:canonical-gap}
To deduce irrationality one would have to construct a negative
certificate for every rational $c$.  The proof above obtains one
only after assuming $c\ne\gamma$; it therefore does not supply that
step.  The generating function also entails no arithmetic
contradiction at rational $c$:
\[
 \sum_{k\geq0}\frac{\mu_k(c)}{k!}z^k
    =\frac{1-cz-z\log(1-z)}{(1-z)^2},\qquad |z|<1.
\]
The moments form an integer affine pencil, unlike the prior
isolated-atom construction whose entries involve special-function
values.  The two-sided positivity selection is related to, but
different from, the prior one-sided rational spectral-edge method.
\end{remark}

\subsection{An explicit estimate valid at every depth}

We use ordinary Laguerre polynomials
\[
 L_j(x)=\sum_{k=0}^j(-1)^k\binom jk\frac{x^k}{k!},
 \qquad
 \int_0^\infty L_i(x)L_j(x)e^{-x}\,dx=\delta_{ij}.
\]
The classical inequality $e^{-x/2}|L_j(x)|\leq1$ for $x\geq0$
is the case $\alpha=0$ of DLMF (18.14.8),
\url{https://dlmf.nist.gov/18.14.E8}.

\begin{lemma}[Normalized positive-form estimate]\label{p5:lem:canonical-energy}
For a polynomial $Q$ of degree less than $N$ normalized by
$\int_0^\infty Q^2e^{-x}\,dx=1$, write
\[
 I=\mathcal F_\gamma(Q),\qquad
 b=\int_0^\infty(x-1)Q(x)^2e^{-x}\,dx.
\]
Then
\begin{equation}\label{p5:eq:canonical-energy-bound}
 I\geq\nu_N:=\frac1{32N^2(1+1/(4N))},\qquad
 |b|\leq\sqrt{\frac2{\nu_N}}\,I.
\end{equation}
\end{lemma}
\begin{proof}
Expand $Q=\sum_{j<N}u_jL_j$ with $\sum u_j^2=1$.
Cauchy--Schwarz and the Laguerre inequality imply
\[
 e^{-x}Q(x)^2\leq e^{-x}\sum_{j<N}L_j(x)^2\leq N.
\]
For $\varepsilon=1/(4N)$, the interval
$[1-\varepsilon,1+\varepsilon]$ thus has probability mass at most
$2N\varepsilon=1/2$ under $Q^2e^{-x}\,dx$.
Outside this interval,
\[
 f(x):=(x-1)\log x\geq\frac{\varepsilon^2}{1+\varepsilon}.
\]
Indeed, $f$ decreases on $(0,1)$ and increases on $(1,\infty)$;
at $1+\varepsilon$ use
$\log(1+\varepsilon)\geq\varepsilon/(1+\varepsilon)$,
and at $1-\varepsilon$ use
$-\log(1-\varepsilon)\geq\varepsilon$.
At least half of the probability mass lies outside, proving the
lower bound for $I$.

Since $f''(x)=(x+1)/x^2>0$, Jensen's inequality gives
\begin{equation}\label{p5:eq:canonical-jensen}
 I\geq f(1+b)=b\log(1+b),\qquad b>-1.
\end{equation}
For $-1<b<0$, $-\log(1+b)\geq-b$ gives a lower bound $b^2$.
For $0\leq b\leq1$, $\log(1+b)\geq b/(1+b)\geq b/2$.
Thus $I\geq b^2/2$ for $|b|\leq1$.  Combining this with
$I\geq\nu_N$ yields
\[
 \frac{|b|}{I}\leq
 \frac{|b|}{\max(\nu_N,b^2/2)}\leq\sqrt{2/\nu_N}.
\]
For $b\geq1$, \eqref{p5:eq:canonical-jensen} instead gives
$b/I\leq1/\log2<\sqrt{2/\nu_N}$.  These cases prove the
second inequality.
\end{proof}

\begin{theorem}[Finite-depth ambiguity]\label{p5:thm:canonical-radius}
Set
\[
 r_N=\frac1{8N\sqrt{1+1/(4N)}}.
\]
The matrix $M_N(c)$ is positive definite for $|c-\gamma|<r_N$
and positive semidefinite for $|c-\gamma|\leq r_N$.
In particular, a negative certificate of degree less than $N$
requires $|c-\gamma|>r_N$.

For $N\geq2$, the feasible set
$\mathcal I_N=\{c\in\mathbb R:M_N(c)\succeq0\}$ is a compact
interval $[\ell_N,u_N]$.  These intervals are nested, their
intersection is $\{\gamma\}$, and
\begin{equation}\label{p5:eq:canonical-endpoints}
 \ell_N\leq\gamma-r_N,\quad
 u_N\geq\gamma+r_N,\quad
 \ell_N\longrightarrow\gamma,\quad u_N\longrightarrow\gamma.
\end{equation}
Both endpoints are real algebraic numbers.
\end{theorem}
\begin{proof}
Under the normalization of Lemma~\ref{p5:lem:canonical-energy},
$\mathcal F_c(Q)=I-(c-\gamma)b$.  The second estimate in
\eqref{p5:eq:canonical-energy-bound} yields the stated strict and
nonstrict positivity.  Homogeneity removes the normalization.

Each feasible set is closed and convex, since it is an intersection
of affine closed half-lines.  Moreover
\[
 \det M_2(c)=
 \det\begin{pmatrix}1&2-c\\2-c&8-4c\end{pmatrix}=4-c^2,
\]
so $\mathcal I_N\subseteq[-2,2]$ for $N\geq2$.
The intervals have nonempty interior by the first part, and they
are nested by taking principal submatrices.
Theorem~\ref{p5:thm:canonical-positivity} identifies their intersection.
Compactness and nesting then give convergence of the endpoints.
At either endpoint the matrix must be singular, since positive
definiteness is open.  The determinant is a nonzero polynomial
in $\mathbb Z[c]$, as it is positive at $c=\gamma$.
This proves algebraicity.
\end{proof}

The estimate concerns the raw finite pencil.  It does not prohibit
accelerated combinations, rational-function tests, or additional
arithmetic information.

\subsection{A separate refinement from compact-uniform Laguerre asymptotics}

\begin{lemma}[Local density]\label{p5:lem:canonical-local-density}
For every compact interval $[a,b]\subset(0,\infty)$,
\begin{equation}\label{p5:eq:canonical-local-density}
 e^{-x}\sum_{j=0}^{N-1}L_j(x)^2
   =\frac{\sqrt N}{\pi\sqrt x}+O_{a,b}(\log N)
\end{equation}
uniformly for $x\in[a,b]$ as $N\to\infty$.
\end{lemma}
\begin{proof}
The compact-uniform expansion in DLMF (18.15.14)--(18.15.15),
at $\alpha=0$, gives
\[
 e^{-x/2}L_j(x)=\frac1{\sqrt\pi\,(jx)^{1/4}}
 \left\{\cos(2\sqrt{jx}-\pi/4)+O_{a,b}(j^{-1/2})\right\}.
\]
The expansion and the compact-uniformity statement are recorded at
\url{https://dlmf.nist.gov/18.15#iv}.  Squaring gives
\[
 e^{-x}L_j(x)^2=
 \frac{\cos^2(2\sqrt{jx}-\pi/4)}{\pi\sqrt{jx}}
       +O_{a,b}(j^{-1}).
\]
The constant part of $\cos^2 t=(1+\cos2t)/2$ sums to
$\sqrt N/(\pi\sqrt x)+O_{a,b}(1)$.
For the oscillatory part, put
$g_x(t)=t^{-1/2}e^{4i\sqrt{tx}}$.  Uniformly for $a\leq x\leq b$,
\[
 \left|\int_1^N g_x(t)\,dt\right|=O_{a,b}(1),\qquad
 |g_x'(t)|\leq\tfrac12t^{-3/2}+2\sqrt b\,t^{-1}.
\]
The integral bound follows on substituting $t=u^2$.
Comparing the sum of $g_x(j)$ to its integral interval by interval
therefore gives $O_{a,b}(1+\int_1^N|g_x'(t)|dt)=O_{a,b}(\log N)$.
The squared-expansion error also sums to $O_{a,b}(\log N)$.
The finitely many initial indices, including $j=0$, contribute
only a bounded term.  This proves \eqref{p5:eq:canonical-local-density}.
\end{proof}

\begin{theorem}[Asymptotic finite-depth cost]\label{p5:thm:canonical-root-radius}
The endpoints in Theorem~\ref{p5:thm:canonical-radius} satisfy
\begin{equation}\label{p5:eq:canonical-root-radius}
 \liminf_{N\to\infty}\sqrt N
    \min(\gamma-\ell_N,u_N-\gamma)\geq\frac\pi8.
\end{equation}
Consequently, as $\delta=|c-\gamma|\to0$, a negative certificate
of degree less than $N$ requires
$N\geq(\pi^2/64-o(1))\delta^{-2}$.
The constant is a lower bound and is not claimed to be optimal.
\end{theorem}
\begin{proof}
Fix $0<\eta<1$ and put
$\varepsilon_N=(1-\eta)\pi/(4\sqrt N)$.
Normalize $Q$ as in Lemma~\ref{p5:lem:canonical-energy}.
Cauchy--Schwarz bounds its probability density by the kernel in
\eqref{p5:eq:canonical-local-density}.  On the interval
$[1-\varepsilon_N,1+\varepsilon_N]$, that estimate gives mass at most
\[
 2\varepsilon_N\{\sqrt N/\pi+O(\log N+1)\}
       =(1-\eta)/2+o(1)\leq1/2
\]
for sufficiently large $N$, uniformly over the polynomial $Q$.
The preceding energy argument consequently yields
\[
 I\geq\frac{\varepsilon_N^2}{2(1+\varepsilon_N)}
   =\frac{(1-\eta)^2\pi^2}{32N}(1+o(1)).
\]
Using exactly the Jensen argument from
Lemma~\ref{p5:lem:canonical-energy}, a positive lower bound $I\geq\nu$
of this size gives the feasible radius $\sqrt{\nu/2}$.
Thus
\[
 \liminf_{N\to\infty}\sqrt N
       \min(\gamma-\ell_N,u_N-\gamma)
       \geq(1-\eta)\pi/8.
\]
Letting $\eta$ decrease to zero proves
\eqref{p5:eq:canonical-root-radius}.  The certificate-degree statement
follows by inversion; Theorem~\ref{p5:thm:canonical-radius} already
ensures $N\to\infty$ when $\delta\to0$.
\end{proof}

\subsection{The integer-square budget and its limitation}

\begin{corollary}[Raw integer forms]\label{p5:cor:canonical-integer-budget}
If $0\ne Q\in\mathbb Z[X]$ has degree $D$, then
\begin{equation}\label{p5:eq:canonical-integer-budget}
 \mathcal F_\gamma(Q)\geq
 \frac{(D!)^2}{32(D+1)^2(1+1/(4(D+1)))}\geq\frac1{144}.
\end{equation}
Uniformly over integer polynomials of degree $D$, the asymptotic
refinement is
\begin{equation}\label{p5:eq:canonical-integer-budget-asymptotic}
 \mathcal F_\gamma(Q)\geq
       \left(\frac{\pi^2}{32}-o(1)\right)
                \frac{(D!)^2}{D+1}.
\end{equation}
These statements concern the raw coefficients $A_Q,B_Q$, before
division by their common content.
\end{corollary}
\begin{proof}
The monic polynomials $P_j=(-1)^jj!L_j$ belong to $\mathbb Z[X]$.
Their unitriangular coefficient matrix implies
$Q=\sum_{j=0}^Dc_jP_j$ with $c_j\in\mathbb Z$ and $c_D\ne0$.
Orthogonality gives
\[
 \int_0^\infty Q^2e^{-x}\,dx
      =\sum_{j=0}^Dc_j^2(j!)^2\geq(D!)^2.
\]
Apply the homogeneous energy bound with $N=D+1$.
Its lower bound is $(D!)^2/[8(D+1)(4D+5)]$.
At $D=0$ it is $1/40$ and at $D=1$ it is $1/144$;
the ratio of the bound at $D+1$ to that at $D$ is
\[
 \frac{(D+1)^3(4D+5)}{(D+2)(4D+9)}>1\quad(D\geq1).
\]
This proves the universal lower bound.
For \eqref{p5:eq:canonical-integer-budget-asymptotic}, use the energy
estimate in the proof of Theorem~\ref{p5:thm:canonical-root-radius}
and then let $\eta\downarrow0$.  All estimates there are uniform
over the normalized polynomial.
\end{proof}

If $g_Q=\gcd(A_Q,B_Q)$, the primitive integer form is
$(A_Q-B_Q\gamma)/g_Q$.  The bounds above do not control $g_Q$ and
do not rule out small primitive forms.  They identify a specific
arithmetic task: demonstrate sufficient common divisibility while
preserving a small, strictly nonzero normalized integral.
Canonical uniqueness, or the failure of a formal shear, cannot
replace that divisibility analysis.

\section{Rational logarithmic Laplace kernels and exact positivity obstructions}
\label{p5:sec:rational-log-laplace}

We seek pure rational linear forms in $\gamma$ by cancelling auxiliary
exponential integrals identically. The following reduction produces such
forms with rational, nonpolynomial kernels. Two explicit constructions are
tested below. Neither gives small primitive integer forms. Rational
reduction and creative telescoping are established methods; the discussion
makes no literature-wide priority claim for that methodology.

\subsection{A rational differential reduction with two local obstructions}
Define
\[
 \mathcal D U=tU''+(1-2t)U'+(t-1)U
       =(\partial_t-1)t(\partial_t-1)U.
\]
For rational $U$ without poles on $[0,\infty)$, all of its derivatives have
at most polynomial growth there, and
\begin{align}
 \int_0^\infty e^{-t}\mathcal D U(t)\,dt&=0,\label{p5:pk:mass}\\
 \int_0^\infty e^{-t}\mathcal D U(t)\log t\,dt&=U(0),\label{p5:pk:log}\\
 \int_0^\infty te^{-t}\mathcal D U(t)\,dt
       &=\int_0^\infty e^{-t}U(t)\,dt.\label{p5:pk:first}
\end{align}
To prove these identities, put $g=e^{-t}U$, so that
$e^{-t}\mathcal D U=(tg')'$. For $\Re s>0$, integration by parts twice gives
\[
 \int_0^\infty t^s e^{-t}\mathcal D U(t)\,dt
       =s^2\int_0^\infty t^{s-1}e^{-t}U(t)\,dt.
\]
The integral on the right has residue $U(0)$ at $s=0$, whereas the left
side is analytic near zero. Its value and derivative at zero yield
\eqref{p5:pk:mass}--\eqref{p5:pk:log}; integrating $t(tg')'$ proves
\eqref{p5:pk:first}. Consequently
\begin{equation}
 \boxed{\int_0^\infty e^{-t}[A+\mathcal D U(t)]\log t\,dt
                 =U(0)-A\gamma.}\label{p5:pk:pure}
\end{equation}
No logarithms of algebraic numbers or auxiliary $E_1$ values remain.

\begin{proposition}[Exact rational normal form]\label{p5:pk:normal}
Every rational function $R$ regular at zero has a unique decomposition
\[
 R=A+\mathcal D U+
 \sum_{\beta}\left(\frac{c_\beta}{t-\beta}
                       +\frac{d_\beta}{(t-\beta)^2}\right),
\]
where $U$ is rational and regular at zero, and its finite poles are among
the nonzero poles $\beta$ of $R$. The residual coefficients at $\beta$
vanish if and only if
\begin{equation}
 \operatorname{Res}_{t=\beta}(e^{-t}R)=
 \operatorname{Res}_{t=\beta}(e^{-t}R\log t)=0.\label{p5:pk:residues}
\end{equation}
Here any local analytic branch of $\log t$ near $\beta$ may be used.
The normal form respects conjugation and Galois action; in particular its
coefficients are rational when the poles and coefficients of $R$ are
rational.
\end{proposition}
\begin{proof}
At a nonzero pole $\beta$, applying $\mathcal D$ to $(t-\beta)^{-k}$
produces a pole of order $k+2$ with leading coefficient
$\beta k(k+1)\ne0$. Remove all orders at least three successively.
On polynomials $\mathcal D$ raises degree by one, preserving the leading
coefficient. Remove all positive polynomial degrees and leave $A$.
The same leading-order comparisons prove uniqueness.

For $g=e^{-t}U$, both residues of $e^{-t}\mathcal D U=(tg')'$ vanish,
since
\[
 \operatorname{Res}((tg')'\log t)
       =-\operatorname{Res}(g')=0.
\]
The two residues of the residual simple and double poles are
\[
 e^{-\beta}(c_\beta-d_\beta),\qquad
 e^{-\beta}\left((c_\beta-d_\beta)\log\beta+\frac{d_\beta}{\beta}\right).
\]
Their coefficient matrix is nonsingular. A change of logarithm branch
adds a constant multiple of the first residue to the second, so the
simultaneous vanishing condition is branch independent.
\end{proof}
This classifies cancellation by rational differential exactness. It does
not assert that residual special-function values are arithmetically
independent. No such assertion is used below.

\subsection{Positive densities and a sharp completely monotone obstruction}
Let $V$ be rational without poles on $[0,\infty)$, and set
\[
 U=V-V',\qquad R=1+\mathcal D(V-V'),\qquad h(t)=t-1-\log t.
\]
If $R\ge0$ there, then \eqref{p5:pk:mass} makes $e^{-t}R(t)\,dt$ a
probability measure. Equations \eqref{p5:pk:log}--\eqref{p5:pk:first} give
\begin{equation}
 \boxed{\gamma+V'(0)=\int_0^\infty e^{-t}R(t)h(t)\,dt>0.}
 \label{p5:pk:positive}
\end{equation}
Indeed its first moment is $1+V(0)$ and its logarithmic moment is
$V(0)-V'(0)-\gamma$. Strict positivity follows because $h(t)>0$ except
at $t=1$, and a continuous probability density cannot be supported at
that point. For $V\in\mathbb Q(t)$ this gives a nonzero pure rational
linear form, certified without assuming irrationality of $\gamma$.

\begin{theorem}[Exact admissibility and the sharp $1/4$ ceiling]
\label{p5:pk:cm}
Suppose the nonzero rational function $V$ has a positive Laplace
representation
\[
 V(t)=\int_{[0,\infty)}e^{-st}\,d\mu(s)
\]
with finite moments through order three. Put
$A=V(0)$, $B=-V'(0)$, $C=V''(0)$ and $K=A+2B+C$.
For $a\ge0$, the density factor
\[
 R_a=1+a\mathcal D(V-V')
\]
is nonnegative on $[0,\infty)$ exactly when $a\le1/K$.
The largest target supplied by this scaling is
\[
 \frac{B}{A+2B+C}\le\frac14.
\]
The bound is sharp as a supremum over rational functions of increasing
pole order, but equality is impossible for a nonzero rational $V$.
\end{theorem}
\begin{proof}
Necessity follows from $R_a(0)=1-aK$. With $W=V-V'$, write
\[
 \mathcal D W=t(W-2W'+W'')-(W-W').
\]
The Laplace representation gives $W-2W'+W''\ge0$ and
$0\le W-W'\le K$. This proves sufficiency. Furthermore,
\[
 \frac{B}{K}
 =\frac{\int s\,d\mu(s)}{\int(1+s)^2\,d\mu(s)}\le\frac14,
\]
since $(1+s)^2-4s=(s-1)^2$. Equality would force $V$ to be a
nonzero multiple of $e^{-t}$, which is not rational.
For sharpness take $V_n(t)=(1+t/n)^{-n}$. Then $A=B=1$,
$C=1+1/n$, and $B/K=n/(4n+1)\to1/4$.
\end{proof}
Finite positive mixtures of $(t+b)^{-m}$ with positive rational
parameters are explicit instances of the hypotheses. The functions
$V_n$ in the proof have positive Gamma-distribution Laplace measures.
Applying \eqref{p5:pk:positive} to $a_nV_n$, where $a_n=n/(4n+1)$, gives
\[
 \gamma-\frac{n}{4n+1}>0.
\]
This fraction is reduced. Thus its primitive integer form is
\begin{equation}
 (4n+1)\gamma-n=(4\gamma-1)n+\gamma\longrightarrow+\infty.
 \label{p5:pk:cm-failure}
\end{equation}
For an elementary strict bound, the standard integral for $\gamma$ and
$\coth u>1/u$ give
\[
 \gamma=\int_0^\infty e^{-t}
       \left(\frac1{1-e^{-t}}-\frac1t\right)dt>\frac12.
\]
Hence \eqref{p5:pk:cm-failure} exceeds $n+1/2$. Arbitrary pole order and
exact logarithmic cancellation do not overcome the sharp analytic
obstruction in this family.

\subsection{Signed partial fractions can cross the ceiling}
The preceding obstruction does not apply to all rational kernels. Put
$x=(1+t)^{-1}$ and
\[
 V=-x+\frac{129}{94}x^2-\frac{45}{94}x^3.
\]
Direct differentiation yields
\[
 1+\mathcal D(V-V')
 =\frac{t^2(35t^3+552t^2+2763t+4946)}{94(1+t)^6}>0\quad(t>0).
\]
Since $-V'(0)=29/94>1/4$, \eqref{p5:pk:positive} becomes the exact identity
\begin{equation}
 \gamma-\frac{29}{94}
 =\frac1{94}\int_0^\infty
 \frac{e^{-t}t^2(35t^3+552t^2+2763t+4946)}{(1+t)^6}
       (t-1-\log t)\,dt>0.\label{p5:pk:explicit}
\end{equation}
Nevertheless $94\gamma-29>18$, so this positive integral is far from
an irrationality certificate.

For systematic exploration, let $V=\sum_{k=0}^n v_kx^k$ and require
$P(x)=xR(t)$ to have nonnegative Bernstein coefficients of degree
$d=n+4$. With $U=V+x^2V_x$, the exact polynomial identity is
\begin{equation}
 P=x+x^4(1-x)U_{xx}+x^2(2-x-2x^2)U_x+(1-2x)U.\label{p5:pk:search}
\end{equation}
A floating linear program proposes a support, after which rational
linear algebra reconstructs the coefficients and checks both
\eqref{p5:pk:search} and every coefficient inequality exactly. The retained
script \texttt{stein\_density\_probe.py} and certificate file
\texttt{positive\_kernel\_certificates.json} record this computation.
The asserted certificates do not depend on numerical optimality.
The separate script \texttt{positive\_kernel\_verify.py} checks all retained
certificates and both explicit symbolic identities without evaluating
$\gamma$ or rerunning numerical optimization.
Selected reduced targets and exact lower bounds on their integer forms
are
\[
\begin{array}{c|c|c}
 n&P/Q&Q\gamma-P>Q/2-P\\\hline
 1&1/5&3/2\\
 2&11/46&12\\
 3&29/94&18\\
 6&9983629/24987068&2519905\\
 10&95700655335119039/212285756553929320&10442222941845621
\end{array}
\]
Every exactly reconstructed certificate at
$n=1,2,3,4,5,6,8,10,12,14$ has $Q/2-P>1$.
These are failed finite attempts, not evidence for a universal
reduced-denominator theorem.

\begin{theorem}[A uniform error floor in the Bernstein-positive cone]
\label{p5:pk:bernstein-floor}
Suppose $R\ge0$, $\int_0^\infty e^{-t}R(t)dt=1$, and
$xR(t)$, with $x=(1+t)^{-1}$, has nonnegative Bernstein coefficients
of degree $d\ge1$. Then
\begin{equation}
 \int_0^\infty e^{-t}R(t)(t-1-\log t)dt
       \ge\frac1{4224d^2}.\label{p5:pk:floor}
\end{equation}
\end{theorem}
\begin{proof}
The density is a convex mixture of normalized positive functions
\[
 h_k(t)=e^{-t}t^k(1+t)^{1-d},\qquad 0\le k\le d.
\]
On $[1/2,3/2]$ their logarithmic derivatives have absolute value at
most $3d$. For $u\in[3/4,5/4]$, integrate over
$|t-u|\le1/(12d)$ to obtain
\[
 \frac{h_k(u)}{\int_0^\infty h_k(t)dt}
       \le6e^{1/4}d<8d.
\]
The same estimate holds for the mixture. With
$\epsilon=1/(32d)$, at most half its mass lies in
$[1-\epsilon,1+\epsilon]$. On the complement,
\[
 t-1-\log t\ge\epsilon-\log(1+\epsilon)
          \ge\frac{\epsilon^2}{2(1+\epsilon)}.
\]
The integral is therefore at least
$[4096d^2(1+1/(32d))]^{-1}\ge(4224d^2)^{-1}$.
\end{proof}
For the gamma targets in \eqref{p5:pk:positive}, this excludes exponential
accuracy in the certificate degree. It does not prove arithmetic
impossibility: a sequence satisfying $Q(\gamma-P/Q)\to0$ would necessarily
have $Q=o(d^2)$, and no reduced-denominator lower bound excluding that
possibility for the whole cone is proved here. The exact tested targets
have very different denominator growth.

A further sparse attempt imposed the two residue conditions of
Proposition~\ref{p5:pk:normal}, together with
$\operatorname{Res}_{t=-1}(te^{-t}R)=0$, on
\[
 R_n(t)=\frac{t^n(a_0+a_1t+a_2t^2+a_3t^3)}{(1+t)^{2n+2}}.
\]
The one-dimensional rational coefficient space has positive generators
at $n=1,2,3$, but at $n=4$ a primitive generator is
\[
 (-187805251068,-59288222773,-5243087734,12508971).
\]
Its opposite constant and leading signs forbid global positivity under
any nonzero scaling. The exact script
\texttt{stein\_sparse\_probe.py} verifies these residue systems.
Thus the first positive instances cannot be extrapolated into a positive
all-index family.

The construction has accomplished exact elimination and certified
nonvanishing. The missing ingredient is primitive arithmetic smallness.
The completely monotone subfamily is sharply obstructed; the larger
rational cone crosses that obstruction but does not currently supply a
sequence proving irrationality of $\gamma$.

\subsection{Moving conjugate poles: a fixed-order obstruction}
\label{p5:pk:moving-poles}
The preceding Bernstein cone keeps its poles at $-1$. A different attempt
places conjugate poles close to the desired concentration point $t=1$.
It too has a precise obstruction when the rational complexity is fixed.
Write
\[
 \mathcal T=\mathcal D(1-\partial_t)
 =-t\partial_t^3+(3t-1)\partial_t^2+(2-3t)\partial_t+(t-1),
 \qquad q_\epsilon(t)=(t-1)^2+\epsilon^2.
\]

For a first explicit test, take
\[
 V_\epsilon(t)=b\frac{t-1}{q_\epsilon(t)},\qquad 0<\epsilon<1.
\]
Its rational target is
\[
 -V_\epsilon'(0)
       =b\frac{1-\epsilon^2}{(1+\epsilon^2)^2}.
\]
At $t=1+\epsilon$, direct differentiation gives
\[
 \mathcal T V_\epsilon(1+\epsilon)
 =-\frac{b(3+5\epsilon+3\epsilon^2-\epsilon^4)}{2\epsilon^4}.
\]
Thus nonnegative target and nonnegative density $1+\mathcal T V_\epsilon$
require
\[
 0\le b\le\frac{2\epsilon^4}{3+5\epsilon+3\epsilon^2-\epsilon^4},
 \qquad 0\le -V_\epsilon'(0)\le\frac23\epsilon^4.
\]
The target tends to zero as the poles approach the real axis. The
higher derivatives force a negative lobe unless the amplitude is reduced
by at least this fourth-power scale.

\begin{theorem}[Fixed-order conjugate poles cannot concentrate]
\label{p5:pk:fixed-moving}
Fix nonnegative integers $m,d$. There is a constant $\eta_{m,d}>0$
such that for every $0<\epsilon\le1$ and every real polynomial $P$
of degree at most $d$, if
\[
 V(t)=\frac{P(t)}{q_\epsilon(t)^m},\qquad
 R(t)=1+\mathcal T V(t)\ge0\quad(t\ge0),
\]
then
\[
 \gamma+V'(0)
 =\int_0^\infty e^{-t}R(t)(t-1-\log t)dt\ge\eta_{m,d}.
\]
In particular, rational coefficients and rational $\epsilon$ in this
fixed-complexity family cannot give primitive integer forms tending to
zero, regardless of coefficient growth.
\end{theorem}
\begin{proof}
First, $\mathcal T V=C$ has no nonzero rational solution for any constant
$C$. At a nonzero pole $\beta$ of order $k$, the leading pole of
$\mathcal T V$ has order $k+3$ and coefficient
$\beta k(k+1)(k+2)$ times the leading coefficient of $V$. At a pole at
zero, its leading order is $k+2$ with nonzero coefficient
$k(k+1)^2$ times that coefficient. Thus a rational solution must be a
polynomial. For a nonzero polynomial of degree $j$, $\mathcal T V$ has
degree $j+1$ with the same leading coefficient, so it cannot be constant.

On $J=[2,3]$, the functions
\[
 1,\quad \mathcal T\left(\frac{t^j}{q_\epsilon(t)^m}\right),
       \qquad 0\le j\le d,
\]
are therefore linearly independent for every $\epsilon\in[0,1]$,
including $\epsilon=0$. They depend continuously on $(\epsilon,t)$
on the compact set $[0,1]\times J$, since $q_\epsilon(t)\ge1$ there.
Their linear combinations with coefficient vector on the Euclidean unit
sphere consequently have a strictly positive uniform minimum $c_{m,d}$
of their $L^1(J)$ norms: otherwise compactness would produce a nontrivial
zero combination at a limiting parameter.

The coefficient vector of $R$ has first coordinate $1$, hence Euclidean
norm at least one. Thus $\int_J|R(t)|dt\ge c_{m,d}$. Positivity turns this
into an ordinary mass bound. Since
\[
 b_J=\min_{t\in J}e^{-t}(t-1-\log t)>0,
\]
its error integral is at least $b_Jc_{m,d}$. Set
$\eta_{m,d}=b_Jc_{m,d}$. For positive $\epsilon$, all integrations in
\eqref{p5:pk:positive} remain legitimate. The limiting singular function
at $\epsilon=0$ is used only on $J$, never inside the full integral.
\end{proof}
The integers $m,d$ are fixed in this theorem. It does not exclude growing
pole order, growing numerator degree, or a different family of moving
poles. The constant is an existence bound from a finite-dimensional
compactness argument, not a claimed sharp numerical value.

\section{Quantitative confluence and a local exceptional-case correction}
\label{p5:conf:sec}

This section tests two transfers aimed at the real Euler constant.  The
$q$-limit below has the correct target, but requires arithmetic estimates
uniform in a moving parameter.  The local calculation corrects an
unrestricted published statement; it does not prove irrationality of the
real Euler constant or of the $p$-adic Euler constant individually.

\subsection{The quantitative statement required at the \texorpdfstring{$q$}{q}-boundary}
Put
\[
 L(q)=\sum_{k\ge1}\frac{q^k}{1-q^k},\qquad
 G(q)=\log(1-q)-\log q\,L(q),\qquad h=-\log q>0.
\]
Thus $G(q)=-\psi_q(1)$.  Its asymptotic expansion is known
\cite{p5:conf:BW}; we record a remainder formulation suited to an
arithmetic application.

\begin{lemma}[A controlled finite-part expansion]\label{p5:conf:expansion}
For every fixed integer $m\ge0$,
\[
 G(e^{-h})=\gamma-\frac h4+
 \sum_{j=1}^m c_jh^{2j}+O_m(h^{2m+2}),\qquad
 c_j=\frac{B_{2j}(1-B_{2j})}{2j(2j)!}.
\]
In particular $c_1=5/144$ and $c_2=-31/86400$.
If $C_m(q)=G(q)+h/4-\sum_{j=1}^m c_jh^{2j}$, then
$|C_m(q)-\gamma|\le K_mh^{2m+2}$, where one may take
\[
 K_m=|c_{m+1}|+
 \frac{2\zeta(2m+2)}{(2\pi)^{2m+2}}
 \int_0^\infty|f^{(2m+2)}(x)|\,dx,
 \qquad f(x)=\frac1{e^x-1}-\frac{e^{-x}}x.
\]
\end{lemma}
\begin{proof}
The singularity of $f$ at zero is removable, $f(0)=1/2$, and every
derivative is integrable at infinity.  Direct summation gives
$G(e^{-h})=h\sum_{k\ge1}f(kh)$.  Also
\[
 H_N-\log N-\int_0^\infty f(x)\,dx
 =\int_0^\infty e^{-Nx}
 \left(\frac1x-\frac1{e^x-1}\right)dx.
\]
The parenthesis is between zero and one, so the integral of $f$ equals
$\gamma$.  Euler--Maclaurin gives the boundary term $-hf(0)/2$ and
$-B_{2j}h^{2j}f^{(2j-1)}(0)/(2j)!$.  The Taylor expansion gives
$f^{(2j-1)}(0)=(B_{2j}-1)/(2j)$.
The periodic-Bernoulli remainder through order $2m+2$ is bounded by
the integral in $K_m$ times $h^{2m+2}$; including the next boundary
term proves the stated inequality.
\end{proof}

\begin{proposition}[A sufficient uniform arithmetic bridge]
\label{p5:conf:bridge}
Set $q_N=1-1/N$.  Fix $m\ge0$ and $\kappa,\mu,\rho\ge0$.
Suppose that, on an infinite unbounded set of integers $N$, there
are integers $A_N\ge1,B_N$ and positive $u_N\to0$ such that
\[
 A_N\le N^\rho,\qquad
 N^{-\kappa}A_N^{-\mu}
 \le |A_NC_m(q_N)-B_N|\le u_N.
\]
If $2m+2>\kappa+\rho(\mu+1)$, then $\gamma$ is irrational.
\end{proposition}
\begin{proof}
Since $N^{-1}\le-\log q_N\le2N^{-1}$, the ratio of
$|A_N(C_m(q_N)-\gamma)|$ to the assumed lower bound is at most a
constant times $N^{\kappa+\rho(\mu+1)-(2m+2)}$, hence tends to zero.
Therefore $A_N\gamma-B_N$ is eventually nonzero and tends to zero.
For rational $\gamma=a/b$, multiplying this sequence by $b$ gives
nonzero integers tending to zero, which is impossible.
\end{proof}

The arithmetic hypothesis is \emph{not established here}.
For example, the Lambert-series irrationality theorems of
Coussement--Smet \cite[Theorems 1.2--1.3]{p5:conf:CS} allow a denominator
parameter $v=1/b$, $b\in\mathbb Z_{\ge2}$; these parameters cannot
tend to one.  Varying their numerator parameter $u$ instead gives
\[
 \sum_{k\ge1}\frac{u^k}{1-v^k}-\frac{u}{1-u}
 =\sum_{k\ge1}\frac{u^kv^k}{1-v^k}\longrightarrow L(v),
\]
by dominated convergence, not $\gamma$.
Moreover, irrationality of $L(q)$ does not give the required lower
bound for the mixture of Lambert values and logarithms in $C_m(q)$.

There is an exact warning about direct truncation.  Its omitted tail
satisfies, by monotone integral comparison,
\[
 -\log(1-e^{-(K+1)h})\le
 T_K(h):=h\sum_{k>K}\frac1{e^{kh}-1}
 \le-\log(1-e^{-Kh}).
\]
Thus $T_K(h)\le h^r$ implies $K\ge r\log(1/h)/h-1$.
At $q_N=1-1/N$ this requires order $N\log N$ terms.  The naive
clearing factor
\[
 D_{N,K}=\prod_{k=1}^K\bigl(N^k-(N-1)^k\bigr)
 \ge N^{K(K-1)/2}
\]
overwhelms the tail: $D_{N,K}T_K\to\infty$ for fixed $N$ as
$K\to\infty$, and along the preceding high-accuracy diagonals.
This concerns this particular clearing factor, not the reduced
denominator of an optimized Pad\'e construction.

\subsection{Diamond normalization and the small-prime exception}
Fix the Iwasawa logarithm on $\mathbb C_p^\times$, so that
$\log_p p=0$ and it vanishes on roots of unity.  Use exactly the
normalization of \cite[p.~355]{p5:conf:MS}:
\begin{align*}
 \gamma_p&=-\frac p{p-1}\lim_{k\to\infty}
 \frac1{p^k}\sum_{\substack{1\le m<p^k\\p\nmid m}}\log_p m,\\
 \psi_p(x)&=\lim_{k\to\infty}\frac1{p^k}
 \sum_{n=0}^{p^k-1}\log_p(x+n),\qquad x\notin\mathbb Z_p.
\end{align*}
The latter function is Diamond's logarithmic derivative outside
$\mathbb Z_p$.  It is not the differently normalized Morita
logarithmic derivative on $\mathbb Z_p$.

\begin{lemma}[Distribution and reflection from the defining limits]
\label{p5:conf:distribution}
For every prime $p$,
\[
 \sum_{r=1}^{p-1}(\psi_p(r/p)+\gamma_p)=0,
 \qquad \psi_p(r/p)=\psi_p(1-r/p)\quad(1\le r<p).
\]
Consequently $\psi_2(1/2)+\gamma_2=0$ and
$\psi_3(1/3)+\gamma_3=\psi_3(2/3)+\gamma_3=0$.
\end{lemma}
\begin{proof}
Summing the finite approximants over $r$ replaces their numerators
by the logarithms of all $m<p^{k+1}$ coprime to $p$, since
$\log_p(m/p)=\log_p m$.  Comparison with the definition of $\gamma_p$
proves the first identity.

For the second, write $a=r/p$ and
$F_k(a)=p^{-k}\sum_{j=0}^{p^k-1}\log_p(a+j)$.
Reversing the summation gives
\[
 F_k(1-a)-F_k(a)=p^{-k}\sum_{j=0}^{p^k-1}
 \log_p\left(1-\frac{p^k}{a+j}\right).
\]
The linear logarithm term is $-\sum_j(a+j)^{-1}$.
Expanding $(a+j)^{-1}=\sum_{\ell\ge0}(-j)^\ell/a^{\ell+1}$,
the $\ell=0$ term has valuation $k+1$, while those for $\ell\ge1$
have valuation at least $k+\ell$.
Here $\sum_{j=0}^{p^k-1}j^\ell$ is divisible by $p^{k-1}$;
this follows by induction, splitting $j=j_0+bp^{k-1}$.
The logarithm terms of degree $d\ge2$ have valuation at least
$k(d-1)+d-v_p(d)$ even without cancellation in the sum.
These bounds tend to infinity, uniformly over the respective
indices, proving reflection.  Distribution and reflection imply
the asserted small-prime identities.
\end{proof}

\begin{theorem}[Corrected prime-denominator classification]
\label{p5:conf:classification}
With the preceding normalization and $1\le r<p$,
\[
 \psi_p(r/p)+\gamma_p=
 \begin{cases}
 0,&p=2\text{ or }p=3,\\
 \text{a transcendental number},&p\ge5.
 \end{cases}
\]
\end{theorem}
\begin{proof}
The small-prime cases follow from Lemma~\ref{p5:conf:distribution}.
For $p\ge5$, put $m=(p-1)/2$, fix a primitive $p$th root $\zeta$,
and set
\[
 x=1-\zeta,\qquad u_a=\frac{1-\zeta^a}{1-\zeta}
 \quad(2\le a\le m).
\]
The $u_a$ are multiplicatively independent algebraic units
\cite[Lemma 14]{p5:conf:MS}.  They are independent also modulo roots
of unity: a torsion product would become an ordinary relation on
raising to the order of that root.  Equivalently one can replace
$u_a$ by the real unit $\zeta^{k_a}u_a$, where
$2k_a\equiv1-a\pmod p$; this does not change its $p$-adic logarithm.

The identity $\prod_{a=1}^{p-1}(1-\zeta^a)=p$ and conjugate pairing
give
\[
 (p-1)\log_p x+2\sum_{a=2}^m\log_p u_a=0.
\]
Substitution into the $p$-adic Gauss formula
\cite[Lemma 11]{p5:conf:MS} gives
\begin{equation}\label{p5:conf:unit-form}
 \psi_p(r/p)+\gamma_p=
 \sum_{a=2}^m\left(\frac2{p-1}+\zeta^{ar}+\zeta^{-ar}\right)
 \log_p u_a.
\end{equation}
Every coefficient is nonzero: $\zeta^{ar}+\zeta^{-ar}$ is an
algebraic integer, whereas $-2/(p-1)$ is a rational noninteger.
Since all $u_a$ are $p$-adic units, a common positive power puts
them arbitrarily close to one.  These powers remain
multiplicatively independent.  The $p$-adic linear-form theorem
\cite[Corollary 13]{p5:conf:MS}, applied to these powers, makes
\eqref{p5:conf:unit-form} transcendental.
\end{proof}

\begin{remark}[Exact scope of the source correction]
In the journal-hosted 2008 PDF, Theorem~2(i) on p.~351 asserts
transcendence of all the displayed values for $q=p$ without a
small-prime exclusion.  Lemma~\ref{p5:conf:distribution}, derived
directly from the same definitions, supplies the exceptions.
One cannot raise the nonunit $1-\zeta_p$ to powers approaching one;
the corrected proof first eliminates its logarithm using the norm
relation.  The independent-unit sum is empty at the exceptional
primes.  No first-priority claim is made for this correction, and
it is not counted as solving a previously open conjecture.
The later prime-power result \cite[Theorem 1.2]{p5:conf:CG} allows
an exception.  None of these conclusions establishes the
arithmetic nature of $\gamma_p$ individually.
\end{remark}

\subsection{A failed adelic identification of the real limit}
For each prime $p$, the term $1/p^n$ is the unique summand of
$H_{p^n}$ with valuation $-n$.  Hence
\[
 v_p(H_{p^n})=-n.
\]
The defining real approximants $H_{p^n}-\log(p^n)\to\gamma$
therefore become $H_{p^n}$ under the Iwasawa logarithm and diverge
in $\mathbb Q_p$.  They do not converge to $\gamma_p$.
Moreover a unital field embedding cannot send the nonzero real
$\log p$ to $\log_p p=0$.  A comparison on a smaller ring may
have a kernel, but compatibility with $\gamma$ would still require
an actual theorem.  The local classification above does not
supply such a comparison or prove irrationality of real $\gamma$.

\section{What remains between these results and irrationality}
\label{p5:sec:remaining}

The research routes above improve several structural statements, but
none has produced \eqref{p5:intro:integer-target}. The exact gaps matter
more than the number of auxiliary theorems.

\paragraph{The monodromy route.}
If $\gamma$ were algebraic, the identity
$\delta=e\operatorname{Ein}(1)-e\gamma$ would put $\delta$ in
$\mathcal E$. The proved consequences make it a transcendental
nonunit, with no nonconstant rational denominator in
$k(\delta)\cap\mathcal E$. Those conclusions are consistent.
There is no independently proved assertion that $1/\delta$ must
belong to $\mathcal E$. Nor can one extract a new $E$-value from
the difference of two monodromy images: that nonzero difference
vanishes at $X=1$, while evaluation there is injective on each image.
Thus the new denominator result is a completed result in its own
setting, not a hidden proof of algebraicity or transcendence of
$\gamma$.

\paragraph{The integer moment route.}
For $Q\in\mathbb Z[X]$, positivity gives
$A_Q-B_Q\gamma>0$. The raw forms have a uniform positive lower
bound and in high degree grow at least factorially. Only reduction
by $g_Q=\gcd(A_Q,B_Q)$ could turn these into small primitive forms.
A sufficient genuinely arithmetic advance would be an explicit
sequence $Q_n$ for which
\[
 \frac{A_{Q_n}-B_{Q_n}\gamma}{g_{Q_n}}\longrightarrow0.
\]
The positivity proof would then supply nonvanishing without assuming
irrationality. No such common-divisor estimate is established here.
The characterization of $\gamma$ by infinitely many integer
inequalities cannot be used to assume that every rational parameter
eventually fails one of them.

\paragraph{The rational-density route.}
Here strict positivity is also available before any arithmetic
assumption. The completely monotone construction is definitively
excluded by its $1/4$ ceiling. In the larger Bernstein-positive
family, a candidate sequence with reduced denominator $Q_n$ and
certificate degree $d_n$ would require $Q_n=o(d_n^2)$ just to be
compatible with a vanishing primitive form. That necessary condition
is not sufficient, and no theorem here rules it out for every
possible sequence. Fixed-order poles approaching $1$ are separately
excluded by compactness. Allowing the degree or pole order to grow
escapes that last theorem and requires new estimates.

\paragraph{The CMF and confluence routes.}
The exact image lattice concerns all integer seeds together. A
particular seed can have further cancellation. The proved bound
$J\mid\lcm(1,\ldots,S+1)^4$ shows that the index obstruction is
only exponential for equal upper decrements; it does not produce
small-height lifts with a sufficiently small, nonzero gamma error.
The lattice theorem does not replace that further control. Likewise, a
quantitatively controlled analytic $q$-limit does not supply uniform
integer linear forms in the moving parameter. Both missing
statements are explicitly arithmetic, and neither is implied by
functional identities or numerical accuracy.

These are research targets, not hypotheses silently adopted in any
of the proofs. The manuscript leaves the irrationality and
transcendence of the real Euler constant unresolved.

\PartAppendix
\section{Exact checks and their limited role}

The proofs do not rely on floating-point optimization or on checking
a finite number of indices. The following diagnostics were used to
reject tempting generalizations and to check formulas.

For the canonical pencil, exact rational Gram recurrences found
negative integer-polynomial certificates for $c=0$ and $c=1$ at
degrees 38 and 41. Their largest coefficients have 544 and 622
decimal digits, respectively. These exclusions are elementary by
other means; their role here is to expose the cost of this specific
certificate system. They are not counted as new irrationality
results.

For positive rational kernels, a numerical linear program suggested
supports; rational linear algebra then reconstructed the coefficients
and checked every positivity certificate exactly. The displayed
fraction $29/94$ is accompanied by a fully factored positive kernel,
so that no numerical optimizer is needed to verify it. The selected
higher-degree fractions are exploratory data, with exact simple
lower bounds on their primitive errors.

For the jet theorem, integer matrix calculations and a separate
finite-characteristic derivation were compared. A conjectural rule
excluding all fresh large primes from reduced recurrence contents
was discarded: in the tested normalized Rivoal pair at index 288,
the pair after clearing its least common denominator has common
divisor 709, which exceeds
$2\cdot288$. This finite counterexample is not an asymptotic
denominator theorem. Here $P_n,Q_n$ are the solutions of the
recurrence in \cite[Section 1]{p5:Rivoal2009}, specialized at $x=1$,
with initials
\[
 (P_0,P_1,P_2)=(-1,4,77/4),\qquad
 (Q_0,Q_1,Q_2)=(1,7,65/2).
\]
The exact recurrence and the gcd check are included in the script
below, so that the normalization of this finite observation is explicit.

\subsection{A compact exact verification script}
The following Python code uses SymPy and verifies the identities
for the moment pencil, the rational telescoper, and the explicit
positive example. It is supplementary to their proofs. It contains
no test that presumes the arithmetic nature of $\gamma$.
\begin{lstlisting}[language=Python]
import sympy as s
t, x, c = s.symbols('t x c')

def D(u):
    return t*s.diff(u,t,2)+(1-2*t)*s.diff(u,t)+(t-1)*u

def mu(k):
    return s.factorial(k)*(k*(s.harmonic(k)-c)+1)

# Integral of a monomial against the signed canonical weight.
for k in range(12):
    from_integral = (s.factorial(k+1)*(s.harmonic(k+1)-c)
                     -s.factorial(k)*(s.harmonic(k)-c))
    assert s.expand(from_integral-mu(k)) == 0
assert s.det(s.Matrix([[mu(0),mu(1)],[mu(1),mu(2)]])) == 4-c*c

# e^(-t) D(U) = (t (e^(-t) U)')'.
U = (t+2)/(t*t+3*t+7)
assert s.simplify(s.exp(-t)*D(U)
        -s.diff(t*s.diff(s.exp(-t)*U,t),t)) == 0

# A positive nonpolynomial kernel beyond the 1/4 ceiling.
y = 1/(1+t)
V = -y+s.Rational(129,94)*y*y-s.Rational(45,94)*y**3
R = 1+D(V-s.diff(V,t))
target = -s.diff(V,t).subs(t,0)
explicit = t*t*(35*t**3+552*t*t+2763*t+4946)/(94*(1+t)**6)
assert target == s.Rational(29,94)
assert s.cancel(R-explicit) == 0
assert target > s.Rational(1,4)

# Finite characteristic identity in the jet proof.
for p in (2,3,5,7,11):
    B = s.zeros(p)
    for j in range(p):
        B[j,j] = j
        B[(j+1)%p,j] -= 1
    C = B**p-B+s.eye(p)
    assert all(int(v)%p == 0 for v in C)

# A finite counterexample to the proposed small-prime gcd shortcut.
from fractions import Fraction as F
from math import gcd, lcm
P, Q = [F(-1),F(4),F(77,4)], [F(1),F(7),F(65,2)]
for n in range(286):
    den = (n+3)**2*(8*n+11)*(8*n+19)
    c2 = (n+3)*(8*n+11)*(24*n*n+145*n+215)
    c1 = -(8*n+27)*(24*n**3+105*n*n+124*n+25)
    c0 = (n+2)**2*(8*n+19)*(8*n+27)
    for seq in (P,Q):
        seq.append((c2*seq[-1]+c1*seq[-2]+c0*seq[-3])/den)
lam = lcm(P[288].denominator,Q[288].denominator)
a288, b288 = int(lam*P[288]), int(lam*Q[288])
assert gcd(a288,b288) == 709
assert 709 > 2*288
print('Exact supplementary identities passed.')
\end{lstlisting}

\paragraph{Verification status.}
The monodromy, primitive-jet, canonical-pencil, and rational-kernel
arguments received independent internal mathematical checks by
separate AI research threads, followed by integration review.
This is not external refereeing or machine-checked formal proof.
Unproved experimental patterns are not promoted to theorems.


\begin{thebibliography}{VIII.99}
\bibitem[V.1]{p5:FischlerRivoal2026Units}
S.~Fischler and T.~Rivoal,
\emph{The ring of values of $E$-functions is not a field},
arXiv:2610.16603 (2026), especially Sections 2--3 and Conjecture 1.2.
\url{https://arxiv.org/abs/2610.16603}.

\bibitem[V.2]{p5:FischlerRivoal2016Eoperators}
S.~Fischler and T.~Rivoal,
\emph{Arithmetic theory of $E$-operators},
J. Ec. polytech. Math. \textbf{3} (2016), 31--65.
\url{https://doi.org/10.5802/jep.28}.

\bibitem[V.3]{p5:FischlerRivoal2024Mixed}
S.~Fischler and T.~Rivoal,
\emph{Relations between values of arithmetic Gevrey series,
and applications to values of the Gamma function},
J. Number Theory \textbf{261} (2024), 36--54.
\url{https://doi.org/10.1016/j.jnt.2024.02.016}.

\bibitem[V.4]{p5:Beukers2006}
F.~Beukers, \emph{A refined version of the Siegel--Shidlovskii theorem},
Ann. of Math. \textbf{163} (2006), 369--379.
\url{https://doi.org/10.4007/annals.2006.163.369}.

\bibitem[V.5]{p5:Kalisch2026}
R.~Kalisch, T.~Monderer, E.~Leibtag, S.~Weinbaum, M.~Shalyt,
and I.~Kaminer, \emph{An Infinite Family of Ap\'ery-like Recurrences
for Euler's Constant}, arXiv:2610.13919 (2026).
\url{https://arxiv.org/abs/2610.13919}.

\bibitem[V.6]{p5:DLMF}
NIST Digital Library of Mathematical Functions,
Chapters 5, 6, 13, and 18.
\url{https://dlmf.nist.gov/}. Consulted 19 October 2026.

\bibitem[V.7]{p5:Sondow2003}
J.~Sondow, \emph{Criteria for irrationality of Euler's constant},
Proc. Amer. Math. Soc. \textbf{131} (2003), 3335--3344.
\url{https://arxiv.org/abs/math/0209070}.

\bibitem[V.8]{p5:BostanReduction}
A.~Bostan, S.~Chen, F.~Chyzak, Z.~Li, and G.~Xin,
\emph{Hermite reduction and creative telescoping for hyperexponential
functions}, arXiv:1301.5038 (2013).
\url{https://arxiv.org/abs/1301.5038}.

\bibitem[V.9]{p5:VanAsscheWolfs2024}
W.~Van Assche and T.~Wolfs, \emph{Rational approximation of Euler's
constant using multiple orthogonal polynomials}, arXiv:2404.09799.
\url{https://arxiv.org/abs/2404.09799}.

\bibitem[V.10]{p5:Rivoal2009}
T.~Rivoal, \emph{Rational approximations for values of derivatives
of the Gamma function}, Trans. Amer. Math. Soc. \textbf{361}
(2009), 6115--6149.
\url{https://rivoal.perso.math.cnrs.fr/articles/eulerconstant.pdf}.

\bibitem[V.11]{p5:conf:BW} S.~Banerjee and B.~Wilkerson,
\emph{Lambert series and $q$-functions near $q=1$},
arXiv:1602.01085.
\url{https://arxiv.org/abs/1602.01085}.

\bibitem[V.12]{p5:conf:CS} J.~Coussement and C.~Smet,
\emph{Irrationality proof of certain Lambert series using little
$q$-Jacobi polynomials}, arXiv:math/0701345.
\url{https://arxiv.org/abs/math/0701345}.

\bibitem[V.13]{p5:conf:MS} M.~Ram Murty and N.~Saradha,
\emph{Transcendental values of the $p$-adic digamma function},
Acta Arith. \textbf{133} (2008), 349--362.
\url{https://doi.org/10.4064/aa133-4-4}.
Journal PDF: \url{https://www.impan.pl/shop/publication/transaction/download/product/82486}.

\bibitem[V.14]{p5:conf:CG} T.~Chatterjee and S.~Garg,
\emph{Transcendental nature of $p$-adic digamma values},
arXiv:2404.07690.
\url{https://arxiv.org/abs/2404.07690}.

\end{thebibliography}
% END PRESERVED SOURCE BODY 5
\phantomsection\label{part5:end}
\clearpage
\setcounter{section}{0}\setcounter{subsection}{0}
\renewcommand{\thesection}{\Roman{part}.\arabic{section}}
\renewcommand{\theHsection}{\arabic{part}.main.\arabic{section}}
\renewcommand{\refname}{References for Part VI}
\fancyhead[L]{\small Euler\textquotesingle s constant: Part VI}
\part{Rational certificates, ramification, and arithmetic obstructions}
\label{part6:start}
\noindent\textbf{Source:} {\small\path{Keio_Gamma_Rational_Certificates_20261019.tex}}.

Original PDF: 27 pages. Mathematical content and source-level qualifications are retained.

\noindent\textit{Original preparation note.} Research draft prepared with AI assistance. New arguments were checked independently within this research session; this is not external peer review or formal proof verification. No claim of literature priority is made.


% BEGIN PRESERVED SOURCE BODY 6
\begin{abstract}
We revisit the attached manuscript \emph{Twelve further attempts at the irrationality of Euler's constant, and where each one breaks}, and test six further routes. The main constructive result is completeness of a positive rational-certificate system: with poles confined to the single pair $1\pm i$ and unbounded order, its exactly attainable rational targets are all $r<\gamma$. The remaining primitive integer-error infimum is zero exactly when $\gamma$ is irrational. Thus analytic existence of these certificates is no longer the missing step; their denominator-versus-error control is.

Ramified exponential-integral differences remove logarithmic monodromy while retaining a nonzero multiple of $\gamma$, correcting the scope of an equal-winding obstruction. We describe affine and algebraic loci of regular-singular cancellation and distinguish them from constructions made only under a rationality hypothesis. We also prove an exact decomposition of individual lattice gcds, a signed prime-periodicity law, divergence of fully primitive Hankel polynomials at $\gamma$, and a torsion-character obstruction to transporting Dickman moments to compact $p$-adic distributions. Exact finite certificates and reproducible computations accompany the proofs.

\emph{The irrationality and transcendence of $\gamma$ remain unproved. No previously open conjecture is counted as solved, and no collection of subsidiary statements is presented as ten solutions of published open problems.}
\end{abstract}



\section{Purpose, inputs, and the arithmetic endpoint}
The starting point is the user-supplied \texttt{gamma\_attempts(1).tex}
and its matching PDF \cite{p6:Attached}. We read all twelve attempts.
The companion manuscript on monodromy, arithmetic and canonical
positivity \cite{p6:Companion} supplies explicitly identified earlier
inputs. This document is a new follow-up; it does not reproduce the
entire preceding corpus or silently promote its open targets to theorems.

For reference, the October 2026 recurrence paper \cite{p6:Kalisch2026}
still describes the irrationality of $\gamma$ as unproved. Its
construction and the recent $E$-value-ring work \cite{p6:FR2026} motivate
some of the present routes. The distinction between a new proof in
this document and a result known to be new in the literature is kept
throughout: the latter has not been established.

Write
\[
 \gamma=\lim_{n\to\infty}(H_n-\log n),\qquad
 \Ein(z)=\sum_{n\ge1}\frac{(-1)^{n-1}z^n}{n\,n!}.
\]
On the positive real axis,
$E_1(x)=\int_x^\infty e^{-t}dt/t
=\Ein(x)-\gamma-\log x$.
A nonzero integer linear form $q\gamma-p$ tending to zero would
prove irrationality. A finite list of very small such forms does not.
Likewise an equivalent criterion is a research reduction, not a
solution unless its remaining assertion is proved.

Our changes of direction are concrete. We prescribe the target
rational number before constructing a positive kernel; replace
linear forms by primitive determinant polynomials; replace linear
pullbacks by ramified ones; examine nonidentifiable and nonlinear
cancellation parameters; isolate the exact seed divisor; and test
a real-to-$p$-adic moment transfer. Each route below states what
survives and what the unsuccessful final step would still require.


\section{Growing-order rational kernels: completeness and primitive certificates}
\label{p6:sec:next-kernel}
The preceding fixed-order compactness obstruction leaves open the case in
which the pole order increases. We keep the two poles fixed at $1\pm i$,
and prove completeness of this growing-order family. This result does not
prove that $\gamma$ is irrational. It identifies exactly what remains of
the arithmetic problem after positivity and logarithmic cancellation have
been imposed. Two finite certificates also show that the height of the
kernel coefficients need not survive in the primitive linear form.

Put $q(t)=(t-1)^2+1$, $h(t)=t-1-\log t$, and
\[
 T=-t\partial_t^3+(3t-1)\partial_t^2+(2-3t)\partial_t+(t-1).
\]
For a rational $V$ with no pole on $[0,\infty)$, let $R_V=1+TV$.
Integration by parts, or the identity
$T(e^t v)=-e^t(tv'')'$, gives
\begin{align}
 \int_0^\infty e^{-t}R_V(t)\,dt&=1,\label{p6:nk:mass}\\
 \gamma+V'(0)&=\int_0^\infty e^{-t}R_V(t)h(t)\,dt.
 \label{p6:nk:error}
\end{align}
For completeness, write $U=V-V'$ and $g=e^{-t}U$; then
$e^{-t}TV=(tg')'$. The mass correction integrates to zero.
The logarithmic moment correction is $U(0)$, obtained by differentiating
$\int t^s(tg')'dt=s^2\int t^{s-1}g(t)dt$ at $s=0$.
The first moment correction is $\int e^{-t}U(t)dt=V(0)$.
Subtracting these moments proves \eqref{p6:nk:error}. These integrations also
hold for a smooth $V$ and its derivatives with polynomial growth. When
$R_V\geq0$, the right side of \eqref{p6:nk:error} is strictly positive:
$h$ vanishes only at $1$, whereas a continuous probability density cannot
be supported at one point.

\begin{theorem}[Completeness with one fixed conjugate pole pair]
\label{p6:nk:complete}
Let $\mathcal C$ consist of the rational functions $V\in\mathbb Q(t)$
whose only possible finite poles are $1\pm i$, for which $V(\infty)=0$
and $1+TV>0$ on $[0,\infty)$. Then
\[
 \boxed{\{-V'(0):V\in\mathcal C\}=\mathbb Q\cap(-\infty,\gamma).}
\]
Indeed, every rational $r<\gamma$ has a representative with
$1+TV\geq\delta$ on $[0,\infty)$ for some $\delta>0$.
\end{theorem}
\begin{proof}
The inclusion from left to right follows from \eqref{p6:nk:error}. We prove
the converse by constructing smooth models, approximating them by rational
functions while retaining a positive margin, and interpolating their
rational targets.

Choose a nonnegative smooth function $b$ with compact support in
$(0,\infty)$, normalized by $\int e^{-t}b(t)dt=1$, and choose $0<\rho<1$.
Set $R_*(t)=\rho+(1-\rho)b(t)$. Thus $R_*\geq\rho$ and its weighted mass
is one. Define
\[
 F(t)=e^{-t}(1-R_*(t)),\quad
 H(t)=\int_0^tF(s)\,ds=-\int_t^\infty F(s)\,ds,
\]
\[
 v(t)=\int_t^\infty(s-t)\frac{H(s)}s\,ds,
 \qquad V_*(t)=e^t v(t).
\]
Since $H(0)=0$, the quotient $H(t)/t$ is smooth at zero; it decays
exponentially at infinity. Differentiation gives $v''=H/t$, hence
$TV_*=-e^tH'=R_*-1$. Beyond the support of $b$,
\begin{equation}
 V_*(t)=-(1-\rho)[1-t e^tE_1(t)],\qquad
 E_1(t)=\int_t^\infty e^{-s}\frac{ds}{s}.
 \label{p6:nk:tail}
\end{equation}
In particular $V_*(\infty)=0$. Its extension in the coordinate $w=1/t$
is smooth at $w=0$: the identity
\[
 t e^tE_1(t)=\int_0^\infty\frac{e^{-u}}{1+uw}\,du
\]
allows differentiation to every order for $w\geq0$, by domination with
$k!u^ke^{-u}$. No convergence of its formal asymptotic series is assumed.
All integrations used in \eqref{p6:nk:error} are therefore justified for
$V_*$. Its target is
\begin{equation}
 r_*=-V_*'(0)=\gamma-\rho\gamma
       -(1-\rho)\int_0^\infty e^{-t}b(t)h(t)\,dt.
 \label{p6:nk:modeltarget}
\end{equation}
By taking $b$ supported in a shrinking neighborhood of $1$ and then
$\rho\downarrow0$, these targets approach $\gamma$ from below. By moving
the support of $b$ to arbitrarily large $t$, with any fixed $\rho<1$,
they tend to $-\infty$.

We next approximate a given $V_*$ without losing positivity. Set
\[
 z(t)=\frac{t-1+i}{t-1-i}=e^{i\theta},\qquad
 0\leq\theta\leq\frac{3\pi}{2},\qquad
 t=1+\cot(\theta/2).
\]
Here $t=\infty$ corresponds to $\theta=0$, and $t=0$ corresponds to
$\theta=3\pi/2$. The function $f(\theta)=V_*(t(\theta))$ is smooth on
this closed arc and $f(0)=0$. In this coordinate
\[
 \partial_t=a(\theta)\partial_\theta,
 \qquad a(\theta)=\cos\theta-1.
\]
For any $C^3$ function $u$ on the arc with $u(0)=0$, put
$V_u(t)=u(\theta(t))$. Direct substitution gives the uniform bound
\begin{equation}
 \|TV_u\|_{L^\infty(0,\infty)}
 \leq90\max_{1\leq j\leq3}\|u^{(j)}\|_\infty.
 \label{p6:nk:Tbound}
\end{equation}
Here are explicit bounds verifying that the apparent singularity at
$\theta=0$ causes no loss. We have $|a|\leq2$,
$|a'|,|a''|\leq1$, $t|a|\leq3$, $ta^2\leq6$, and
$t|a|^3\leq12$. The coefficients of $u'$, $u''$, and $u'''$ in the
three derivative terms of $T$ are bounded by $33$, $40$, and $12$,
respectively. Finally
$|(t-1)u|\leq5\|u'\|_\infty$, since
$\theta|\cot(\theta/2)|<5$ on the arc. Their sum is $90$.

Extend $f$ to a $C^3$ periodic function on the circle: on the remaining
arc a Hermite polynomial can match its endpoint derivatives of orders
zero through three. The Fej\'er means of this extension converge with
all derivatives through order three, because differentiation commutes
with convolution and Fej\'er means converge uniformly for continuous
periodic functions. Subtracting the value at zero, and then approximating
the finitely many trigonometric coefficients by rationals, proves that
$f$ can be approximated in $C^3$ by functions
\begin{equation}
 g(\theta)=\sum_{k=1}^m
       \bigl[A_k(\cos k\theta-1)+B_k\sin k\theta\bigr],
 \qquad A_k,B_k\in\mathbb Q.
 \label{p6:nk:trig}
\end{equation}
The associated $V_g$ is a real rational function with rational
coefficients, poles only at $1\pm i$, and value zero at infinity.
Taking the approximation sufficiently close in \eqref{p6:nk:Tbound} gives
$1+TV_g\geq\rho/2$. Moreover
$-V_g'(0)=g'(3\pi/2)\in\mathbb Q$ converges to $r_*$.

Now fix any rational $r<\gamma$. Choose two smooth models whose targets
lie strictly below and strictly above $r$, respectively. Approximate them
closely enough to obtain $V_a,V_b$ with rational targets $a<r<b$ and
strictly positive uniform margins. Then
\[
 V=\frac{b-r}{b-a}V_a+\frac{r-a}{b-a}V_b
\]
has rational coefficients and target exactly $r$. Its density is the
same convex combination of the two positive densities, so it retains
a positive uniform margin. This proves the theorem.
\end{proof}

\begin{corollary}[Finite certificates and the remaining arithmetic condition]
\label{p6:nk:arithmetic}
For each rational $r<\gamma$, enumeration of $m$ and rational polynomials
$P$ of degree less than $2m$, followed by exact polynomial sign tests,
eventually finds $V=P/q^m$ with $-V'(0)=r$ and $1+TV>0$.
If $r=p/s$ is reduced with $s>0$, the resulting primitive linear form is
$s\gamma-p$. Furthermore
\[
 \gamma\notin\mathbb Q
 \quad\Longleftrightarrow\quad
 \inf_{V\in\mathcal C}
       \operatorname{den}(-V'(0))\,[\gamma+V'(0)]=0.
\]
\end{corollary}
\begin{proof}
Every member of $\mathcal C$ has the stated denominator shape. Its
positivity reduces to that of a polynomial, since $q>0$; Sturm's theorem
decides strict positivity on the positive half-line. The theorem ensures
that an enumeration finds a certificate whenever $r<\gamma$.
If $\gamma=a/b$ is rational in lowest terms, every positive primitive
form $s\gamma-p$ is at least $1/b$. This lower bound is attained: choose
integers $s>0,p$ with $sa-pb=1$. They are coprime, and completeness
represents $p/s<\gamma$. Thus the infimum is exactly $1/b$ in this case. If $\gamma$ is irrational, the
continued-fraction convergents below it include infinitely many
$p/s$ with $0<s\gamma-p<1/s$; completeness represents all of them.
\end{proof}
The arithmetic equivalence in the last sentence uses the standard
continued-fraction criterion. The substantive input here is that all
lower rational approximants admit positive kernels with the same two
fixed poles and exact logarithmic cancellation. Completeness gives no
useful upper bound for the required pole order or coefficient height.
It supplies neither a uniform constructive sequence of small primitive
forms nor a proof of irrationality. A finite certificate for a lower
bound is not a solution of that open problem.

\subsection{Two exact certificates with large removable common factors}
The accompanying exact data give two rational functions
$V_m=P_m/q^m$, with $m=8,12$, for which $R_m=1+TV_m>0$ on
$[0,\infty)$ and
\[
 -V_8'(0)=\frac9{16},\qquad -V_{12}'(0)=\frac47.
\]
The $m=12$ witness has value zero at infinity and belongs to the exact
class $\mathcal C$ of Theorem~\ref{p6:nk:complete}. It was obtained from a
bounded witness by subtracting its constant limit $597/5000000$; this
does not change its target, and the resulting density remains strictly
positive by the exact sign test below. The $m=8$ witness is bounded with
a nonzero limit at infinity, which still suffices for \eqref{p6:nk:error}.
The construction first fixes the small primitive target, then chooses
kernel coefficients, so it does not rely on a fortuitous common divisor.

Here is a precise verification rule. For a given $P=P_m$, set $N_0=P$ and
\[
 N_{j+1}=qN_j'-(m+j)q'N_j\quad(0\leq j<3).
\]
Then $V^{(j)}=N_j/q^{m+j}$ and
\begin{equation}
 R_m=\frac{q^{m+3}-tN_3+(3t-1)qN_2
                +(2-3t)q^2N_1+(t-1)q^3P}{q^{m+3}}.
 \label{p6:nk:numerator}
\end{equation}
After multiplication of its numerator by a positive common denominator,
the resulting integer polynomial has positive value at zero and exactly
zero real roots on $(0,\infty)$, by exact Sturm computation. Its degrees
are $23$ and $30$, respectively. Thus this is a global certificate,
not a sampled positivity assertion.

\begin{center}
\begin{tabular}{c|r|r|r|c}
$m$ & raw coefficient of $\gamma$ & raw integer coefficient
 & common divisor & primitive form\\\hline
8&409600000000&230400000000&25600000000&$16\gamma-9$\\
12&183500800000000&104857600000000&26214400000000&$7\gamma-4$
\end{tabular}
\end{center}
In both rows the common divisor is the full greatest common divisor.
The exact positive identities imply
\[
 0<16\gamma-9<\frac14,\qquad
 0<7\gamma-4<\frac1{20}.
\]
For an elementary verification of the upper bounds, use
$\gamma<289/500$. Indeed, if $a_n=H_n-\log n$, then
\[
 a_n-\gamma=
 \sum_{k=n}^\infty\left[\log(1+1/k)-\frac1{k+1}\right]
 >\frac1{2(n+1)}.
\]
The estimate follows by integrating $1/x-1/(k+1)$ on $[k,k+1]$
and then comparing the sum of $1/[2(k+1)^2]$ with its integral.
Also $\log2>2\sum_{j=0}^3[(2j+1)3^{2j+1}]^{-1}=:L$.
Exact rational arithmetic gives
$H_{32}-5L-1/66<289/500$, which proves the claim.

The files \texttt{moving\_certificates.json} and
\texttt{verify\_kernel\_certificates.py} contain the complete rational
coefficients and an independent exact verification. The polynomial
coefficients are also printed below, in increasing powers of $t$ after
multiplication by their displayed denominators. These bounds on
$\gamma$ are not asserted to be numerically new; the result is an exact
positive-kernel realization with a deliberately small primitive form.


\subsection{Complete coefficient appendix and executable sign certificate}
Write $P_m(t)=L_m^{-1}\sum_{j=0}^{2m}c_{m,j}t^j$. The denominators are
\[L_8=800000000,\qquad L_{12}=22400000000.\]
Blank entries mean zero. All entries below are exact integers.
\begin{center}\small
\begin{tabular}{r|r|r}
$j$&$c_{8,j}$&$c_{12,j}$\\\hline
0&12739606016&720764073984\\
1&-217116848128&-61077968887808\\
2&1309352582208&673998911908864\\
3&-4573301234560&-3884834464481280\\
4&10912087626016&14961446716067328\\
5&-19234037607136&-42751574553974784\\
6&26110432892400&95931077322792704\\
7&-27904440731776&-174960053696659456\\
8&23692954202800&265073025752304640\\
9&-15976408455696&-338200122492639872\\
10&8483955133932&366245497685868544\\
11&-3490637535048&-337813701541174272\\
12&1084266045698&265446174967026432\\
13&-244100083718&-177290672020359456\\
14&37180970397&100228396732354432\\
15&-3353835512&-47680867158749696\\
16&127145584&18940508192803272\\
17&&-6218385700629880\\
18&&1663548645882260\\
19&&-355169721881216\\
20&&58579277695970\\
21&&-7060837833734\\
22&&558246990041\\
23&&-21941765888\\
24&&\\
\end{tabular}
\end{center}
The following Python/SymPy program reconstructs both certificates directly
from these coefficients. Its sign decisions use exact rational arithmetic
and Sturm root counts; it does not invoke the floating-point search.
\begin{lstlisting}
import sympy as S
from math import gcd
t = S.symbols("t")
q = S.Poly((t-1)**2+1,t)
certificates = [
    (8, 800000000, S.Rational("9/16"), [
        12739606016, -217116848128, 1309352582208,
        -4573301234560, 10912087626016, -19234037607136,
        26110432892400, -27904440731776, 23692954202800,
        -15976408455696, 8483955133932, -3490637535048,
        1084266045698, -244100083718, 37180970397,
        -3353835512, 127145584,
    ]),
    (12, 22400000000, S.Rational("4/7"), [
        720764073984, -61077968887808, 673998911908864,
        -3884834464481280, 14961446716067328, -42751574553974784,
        95931077322792704, -174960053696659456, 265073025752304640,
        -338200122492639872, 366245497685868544, -337813701541174272,
        265446174967026432, -177290672020359456, 100228396732354432,
        -47680867158749696, 18940508192803272, -6218385700629880,
        1663548645882260, -355169721881216, 58579277695970,
        -7060837833734, 558246990041, -21941765888,
    ]),
]
for m,L,r,c in certificates:
    P = S.Poly(sum(S.Rational(a,L)*t**j
                   for j,a in enumerate(c)),t)
    assert -(2*P.diff().eval(0)+2*m*P.eval(0))/2**(m+1)==r
    N=[P]
    for j in range(3):
        N.append(q*N[-1].diff()-(m+j)*q.diff()*N[-1])
    B=(q**(m+3)-S.Poly(t,t)*N[3]
       +S.Poly(3*t-1,t)*q*N[2]
       +S.Poly(2-3*t,t)*q**2*N[1]
       +S.Poly(t-1,t)*q**3*P)
    _,B=B.clear_denoms(convert=True)
    assert B.eval(0)>0 and B.count_roots(0,S.oo)==0
    D=L*2**(m+1); p=int(D*r); g=gcd(D,p)
    assert (D//g,p//g)==(int(S.denom(r)),int(S.numer(r)))
    print(m, B.degree(), D//g, p//g, "verified")
\end{lstlisting}


\subsection{A further exact primitive certificate: \texorpdfstring{$26\gamma-15$}{26 gamma - 15}}
\label{p6:nk:15-26}
The next finite certificate improves the primitive error, and uses the
fixed pole pair $1\pm i/2$. It is independent of the completeness theorem's
choice $1\pm i$. Put
\[
 Q(t)=4(t-1)^2+1,\quad
 z(t)=\frac{2(t-1)+i}{2(t-1)-i},\qquad
 V(t)=\sum_{k=1}^{40}\{a_k(\Re(z(t)^k)-1)+b_k\Im(z(t)^k)\}.
\]
The exact coefficients are $c_{2k-2}=a_k$, $c_{2k-1}=b_k$ in the
zero-indexed list \texttt{c} in the executable certificate below.
In particular $V=P/Q^{40}$ with $\deg P<80$, so $V(\infty)=0$.
Exact differentiation gives $-V'(0)=15/26$.
The numerator of $R=1+TV$, over $Q^{43}$ and after positive denominator
clearing, is an integer polynomial $B$ of degree $86$.
The verification below proves its strict positivity on $[0,\infty)$.
Consequently
\[
 \boxed{0<26\gamma-15<\frac1{100}.}
\]
For the upper bound, the same harmonic-sum estimate used above gives
\[
 \gamma<H_{128}-7L-\frac1{258}<\frac{1501}{2600},\qquad
 L=2\sum_{j=0}^7\frac1{(2j+1)3^{2j+1}}<\log2.
\]
Both final comparisons are exact rational calculations in the code.
This finite inequality is not asserted to be numerically new.

Here is a short algebraic explanation of the global sign certificate.
Write $B(t)=\sum_{k=0}^{n}B_k t^k$, $n=86$. Under $x=t/(1+t)$,
\[
 (1-x)^nB\!\left(\frac{x}{1-x}\right)
 =\sum_{k=0}^n\frac{B_k}{\binom nk}
               \binom nk x^k(1-x)^{n-k}.
\]
Thus the numbers $B_k/\binom nk$ are its Bernstein coefficients on
$[0,1]$. Midpoint subdivision is the exact de Casteljau algorithm:
repeated averages produce the Bernstein coefficients on the two halves.
The routine \texttt{split} implements these averages after multiplying
all output coefficients by the common positive factor $2^n$.
Multiplication by positive factors does not change their signs.
After $34$ dyadic subintervals, of maximum depth $11$, every coefficient
on every interval is strictly positive. The Bernstein basis is
nonnegative and sums to one, including at both endpoints. The polynomial
is therefore strictly positive on $[0,1]$, and $B(t)>0$ for every finite
$t\geq0$. This is an exact global proof, with no floating-point sign test.

The following listing is self-contained; the coefficient list is the
complete certificate, and no search output or external data file is
needed to verify it.
\begingroup\small
\begin{lstlisting}
import sympy as S
from math import comb,lcm
t=S.symbols("t"); m=40; Q=S.Poly(4*(t-1)**2+1,t)
a1=S.Rational(
    253680339833956228611786925344465790449083,
    945874489843845367431640625000000000000000)
integers=[
    -9212384918587,10384847593798,-14278998630324,3085415136785,
    -11360345279171,-738559184268,-8832187772878,-2676395526916,
    -6192423671432,-3567107918961,-4069735992231,-3768019736809,
    -2212080899728,-3511982101329,-775955612605,-2951380108521,
    324729349222,-2227199615677,1030747902366,-1459122845145,
    1409827388935,-750583296281,1484372769611,-173000124064,
    1349531831870,234694291741,1073178464829,468574893037,
    747558209639,550746939535,429434328495,520771849028,
    169775817002,423236565949,-15619353075,299267927856,
    -121509496856,179685644252,-163685913715,83217985364,
    -159109492605,16936071851,-129852242501,-20500974617,
    -91077249192,-35467040493,-55224081475,-35854616717,
    -27149265537,-28894235243,-9019495515,-19818304346,
    974084002,-11717080502,4819873639,-5830415500,
    5322959971,-2236679339,4162878054,-412148723,
    2741366730,287544954,1504064412,412862235,
    721871281,316274428,269031677,185957829,
    78762343,88693716,2706880,34901675,
    -7954641,10700954,-8052289,2440388,
    -2569058,268188,-1056300,
]
c=[a1]+[S.Rational(v,10**14) for v in integers]
assert len(c)==80
Y=S.Poly(2*(t-1),t)
re,im=S.Poly(1,t),S.Poly(0,t); P=S.Poly(0,t)
for j in range(1,81):
    re,im=Y*re-im,re+Y*im
    if j%2==0:
        k=j//2
        P+=c[2*k-2]*(re*Q**(m-k)-Q**m)
        P+=c[2*k-1]*im*Q**(m-k)
assert P.degree()<2*m
r=-(5*P.diff().eval(0)+8*m*P.eval(0))/5**(m+1)
assert r==S.Rational(15,26)
N=[P]
for j in range(3):
    N.append(Q*N[-1].diff()-(m+j)*Q.diff()*N[-1])
B=(Q**(m+3)-S.Poly(t,t)*N[3]
   +S.Poly(3*t-1,t)*Q*N[2]
   +S.Poly(2-3*t,t)*Q**2*N[1]
   +S.Poly(t-1,t)*Q**3*P)
_,B=B.clear_denoms(convert=True)
n=B.degree(); assert n==86
M=lcm(*[comb(n,k) for k in range(n+1)])
b=[int(B.nth(k))*(M//comb(n,k)) for k in range(n+1)]
def split(b):
    left=[b[0]*(1<<n)]; right=[b[-1]*(1<<n)]
    for j in range(1,n+1):
        b=[x+y for x,y in zip(b[:-1],b[1:])]
        left.append(b[0]*(1<<(n-j)))
        right.append(b[-1]*(1<<(n-j)))
    return left,right[::-1]
stack=[(b,0)]; leaves=0; maxdepth=0
while stack:
    b,d=stack.pop()
    if min(b)>0:
        leaves+=1; maxdepth=max(maxdepth,d); continue
    assert b[0]>0 and b[-1]>0 and d<12
    left,right=split(b)
    stack.extend([(left,d+1),(right,d+1)])
assert (leaves,maxdepth)==(34,11)
H=sum(S.Rational(1,k) for k in range(1,129))
L=2*sum(S.Rational(1,(2*j+1)*3**(2*j+1))
        for j in range(8))
assert H-7*L-S.Rational(1,258)<S.Rational(1501,2600)
print("Verified: 0 < 26*gamma-15 < 1/100")
\end{lstlisting}
\endgroup


% Standalone fragment: no preamble.  Analytic conclusions, not a proof about
% the irrationality of gamma.  Labels have prefix phn:.
\section{Ramified cancellation and the limits of extracting the Euler constant}
\label{p6:phn:section}

This section examines two possible ways of removing the logarithmic
monodromy of an exponential-integral tail.  Ramification really does
preserve a nonzero coefficient of $\gamma$ after the monodromy has
cancelled; an unconditional obstruction remains in the coefficient
arithmetic.  Exponentiating the tail removes its monodromy as well,
but every nonconstant rational function of the resulting function
fails to be D-finite.  These are precise tests of proposed extraction
steps, rather than proofs of irrationality or transcendence of $\gamma$.
They do not assert that the elementary identities below are new in the
literature.

We use the entire function
\[
 \operatorname{Ein}(z)=\sum_{n\geq1}\frac{(-1)^{n-1}z^n}{n\,n!}
\]
and the classical identity
\begin{equation}\label{p6:phn:identity}
 E_1(z)=\operatorname{Ein}(z)-\log z-\gamma.
\end{equation}
Initially the logarithm and $E_1$ have their principal branches; all
subsequent statements about continuation use compatible branches.
Identity~\eqref{p6:phn:identity} and its integral normalization are
\cite[\S6.2, equations (6.2.1), (6.2.3), (6.2.4)]{p6:DLMF}.
D-finiteness below means satisfaction of a nonzero homogeneous linear
differential equation over $\mathbb C(z)$; where asserted, the equation
can be chosen over $\mathbb Q(z)$.

\subsection{A ramified escape from cancellation of the common constant}

\begin{theorem}[Entire recessive connections with a gamma initial value]
\label{p6:phn:ramification}
For an integer $m\geq2$, set, initially near the positive real axis,
\[
 F_m(z)=E_1(z^m)-mE_1(z).
\]
Then $F_m$ extends to an entire function and
\begin{equation}\label{p6:phn:ramification-identity}
 F_m(z)=(m-1)\gamma+\operatorname{Ein}(z^m)
                    -m\operatorname{Ein}(z).
\end{equation}
It has the following properties.
\begin{enumerate}
\item $F_m(0)=(m-1)\gamma$ and
      $F_m(x)\sim-m e^{-x}/x$ as $x\to+\infty$.
\item Its minimal differential order over $\mathbb C(z)$ is exactly
      three, and it admits an equation of order three over
      $\mathbb Q(z)$.
\item It is neither an $E$-function nor a $G$-function, independently
      of any assumption on $\gamma$.
\end{enumerate}
\end{theorem}
\begin{proof}
On a neighborhood of the positive axis, $\log(z^m)=m\log z$.
Substitution into~\eqref{p6:phn:identity} proves
\eqref{p6:phn:ramification-identity} and the entire extension.
For $x>0$, integration by parts gives
\[
 E_1(x)=\frac{e^{-x}}x-
              \int_x^\infty\frac{e^{-t}}{t^2}\,dt
       =\frac{e^{-x}}x\bigl(1+O(x^{-1})\bigr).
\]
The term $E_1(x^m)$ is negligible relative to $e^{-x}/x$, proving
assertion~(1).

Put $A=e^{-z}$ and $B=e^{-z^m}$.  Direct differentiation yields
\begin{equation}\label{p6:phn:system}
 F_m'=\frac m z(A-B),\qquad A'=-A,\qquad
 B'=-mz^{m-1}B.
\end{equation}
This rational differential system gives an equation over
$\mathbb Q(z)$ of order at most three.  The functions $A,B$ are
linearly independent over $\mathbb C(z)$: otherwise
$A/B=\exp(z^m-z)$ would be rational, contrary to its essential
singularity at infinity.  Moreover, the determinant expressing
$F_m',F_m''$ in the basis $A,B$ is
\[
 \frac{m^2}{z^2}(mz^{m-1}-1),
\]
which is not identically zero.  Hence $F_m',F_m''$ are independent.
If $F_m,F_m',F_m''$ were dependent, we would have
$F_m=rA+sB$ with $r,s\in\mathbb C(z)$.  Differentiating and using the
independence of $A,B$ would give
\[
 r'-r=\frac m z.
\]
There is no such rational $r$: a pole of $r$ at zero makes $r'-r$
have a pole of order at least two; absence of that pole makes
$r'-r$ holomorphic at zero.  Neither alternative produces $m/z$.
This proves assertion~(2).

Write $F_m(z)=\sum_{n\geq0}c_nz^n$.  If $m\nmid n$ and $n\geq1$,
\begin{equation}\label{p6:phn:off-block}
 c_n=-\frac{m(-1)^{n-1}}{n\,n!}.
\end{equation}
Its reduced denominator is at least $n\,n!/m$.  The common
denominators of the first $n$ coefficients therefore cannot be
bounded exponentially along these indices.  This contradicts the
necessary denominator condition for a $G$-function.  If $c_0$ is
not algebraic, that condition already fails at the constant term;
if it is algebraic, rational coefficients such as
\eqref{p6:phn:off-block} still give the stated obstruction in any
number field, since a rational algebraic integer is an integer.

At the indices $n=mk$ one has
\begin{equation}\label{p6:phn:on-block}
 (mk)!c_{mk}=\frac{(-1)^{k-1}(mk)!}{k\,k!}
                  -\frac{(-1)^{mk-1}}k.
\end{equation}
The absolute values in~\eqref{p6:phn:on-block} grow faster than any
exponential in $mk$, by the elementary factorial estimate or
Stirling's formula.  This contradicts the necessary exponential
size bound on the coefficients of an $E$-function written as
$\sum a_nz^n/n!$.  Assertion~(3) follows.
\end{proof}

Thus the statement ``removing the tail monodromy also removes
$\gamma$'' is false without a restriction on ramification.  A loop
around zero changes $E_1(z^m)$ by $-2\pi i m$ and changes $mE_1(z)$
by the same amount, whereas their coefficients of $-\gamma$ are
$1$ and $m$.  At $z=1$, however,
\[
 F_m(1)=-(m-1)E_1(1)
       =(m-1)\bigl(\gamma-\operatorname{Ein}(1)\bigr).
\]
The value identity therefore returns the same unproved arithmetic
bridge between the Euler constant and the resummed tail.  Under the
hypothesis that $\gamma$ is algebraic, the Taylor germ of $F_m$ has
algebraic coefficients and is D-finite; neither fact is contradictory.

A simple family of factorial transforms makes the remaining
arithmetic mismatch explicit.  This next assertion concerns precisely
the transforms defined here, not all possible transforms of a
D-finite function.

\begin{proposition}[Failure of block factorial normalization]
\label{p6:phn:block-borel}
Let $H_m=F_m-F_m(0)$, and for an integer $d\geq1$ define the formal
coefficient transform
\[
 \mathcal B_d\left(\sum_{n\geq0}c_nz^n\right)
       =\sum_{n\geq0}\lfloor n/d\rfloor!\,c_nz^n.
\]
For every $d$, $\mathcal B_dH_m$ is not a $G$-function.  If $d<m$
its radius of convergence is zero.  If $d\geq m$, its rational
coefficients already violate the exponential common-denominator
condition.
\end{proposition}
\begin{proof}
For $n=mk$, formula~\eqref{p6:phn:on-block} gives
\[
 \lfloor mk/d\rfloor!\,c_{mk}
  =\frac{(-1)^{k-1}\lfloor mk/d\rfloor!}{k\,k!}
     -\frac{(-1)^{mk-1}\lfloor mk/d\rfloor!}{k\,(mk)!}.
\]
If $d<m$, the first term has superexponential size in $mk$ and the
second is negligible relative to it.  The radius is zero.
If $d\geq m$, use instead the infinitely many $n$ with $m\nmid n$.
Since $\lfloor n/d\rfloor!$ divides $n!$, the reduced denominator
of the corresponding coefficient is at least
\[
 \frac{n\,n!}{m\lfloor n/d\rfloor!}.
\]
Its logarithm is $(1-1/d)n\log n+O(n)$, which is superlinear.
This proves the denominator obstruction.
\end{proof}

\subsection{Exactly what rational monodromy cancellation removes}

\begin{theorem}[Rational invariants for tails with prescribed winding]
\label{p6:phn:fixed-field}
Let $T_1,\ldots,T_r$ be multivalued meromorphic germs continuing
around a common based loop with
\[
 M T_j=T_j-w_j\tau,\qquad 0\ne w_j\in\mathbb Z,
 \qquad \tau=2\pi i.
\]
Let $K$ be a field of single-valued meromorphic germs fixed by this
continuation, containing $\tau$.  Set
\[
 U_j=T_j/w_j,\qquad
 L=K(U_2-U_1,\ldots,U_r-U_1),\qquad
 F=K(T_1,\ldots,T_r).
\]
Then $U_1$ is transcendental over $L$, and the fixed field is exactly
\begin{equation}\label{p6:phn:fixed-field-eq}
 F^{M}=L.
\end{equation}
If $K$ is a differential field containing all $T_j'$, the same
statement classifies rational differential expressions in the $T_j$
that are fixed by the loop.
\end{theorem}
\begin{proof}
The field $L$ is fixed, and $F=L(U_1)$.  If $U_1$ were algebraic over
$L$, continue a nonzero polynomial relation around the loop any
integer number $n$ of times.  The same polynomial would vanish at
all the distinct elements $U_1-n\tau$, an impossibility.  Thus $F$
is a rational function field in $U_1$ over $L$.

It remains to determine the rational functions fixed by the
translation $U_1\mapsto U_1-\tau$.  Over an algebraic closure of
$L$, the finite set of finite poles of an invariant rational function
must be translation invariant.  Such a nonempty finite set cannot
be invariant under translation by a nonzero element in characteristic
zero.  The rational function is therefore a polynomial.  Comparing
the leading nonconstant coefficient of $p(X-\tau)-p(X)$ shows that
an invariant polynomial is constant.  This proves
\eqref{p6:phn:fixed-field-eq}.  Finally all derivatives of positive order
of the $T_j$ lie in $K$ under the differential hypothesis, so adjoining
those derivatives does not enlarge $F$.
\end{proof}

For the equally ramified tails $T_j(z)=E_1(a_jz)$ with $a_j\ne0$,
the windings are all $w_j=1$, and
\begin{equation}\label{p6:phn:same-winding}
 T_j-T_1=\operatorname{Ein}(a_jz)-\operatorname{Ein}(a_1z)
                  -\log a_j+\log a_1.
\end{equation}
This eliminates the common formal additive parameter $-\gamma$ as
well as $-\log z$.  The assertion is a classification of analytic
invariants.  In particular, if $K$ contains all complex constants,
it does not imply that $\gamma$ is absent from the arithmetic field
of coefficients: it already belongs to that $K$.

For $T_j(z)=E_1(a_jz^{m_j})$ with $a_j\ne0$ and $m_j\geq1$, use instead
$w_j=m_j$.  Formula~\eqref{p6:phn:fixed-field-eq} then permits
\[
 \frac{T_j}{m_j}-\frac{T_1}{m_1}
 =\left(\frac1{m_1}-\frac1{m_j}\right)\gamma
 +\frac{\operatorname{Ein}(a_jz^{m_j})-\log a_j}{m_j}
 -\frac{\operatorname{Ein}(a_1z^{m_1})-\log a_1}{m_1}.
\]
The ramified example in Theorem~\ref{p6:phn:ramification} is thus an
invariant allowed by the exact fixed-field theorem.

There is a useful direct version for differential cancellation.
If $h_j$ and the coefficients of differential operators $P_j$ are
single-valued along the loop, continuation commutes with
multiplication and differentiation and gives
\begin{equation}\label{p6:phn:differential-cancellation}
 (M-I)\sum_jP_j(h_jT_j)=-\tau\sum_j w_jP_j(h_j).
\end{equation}
For equal windings, a single-valued output forces the coefficient
of the common formal additive constant to vanish.  For unequal
windings only the weighted sum vanishes.  This proves the precise
scope of that linear cancellation principle.

Even entire D-finiteness after cancellation does not supply the
missing arithmetic hypothesis.  Unconditionally,
\[
 E_1(z)-E_1(2z)=\log2+\operatorname{Ein}(z)
                            -\operatorname{Ein}(2z)
\]
is entire and D-finite, but is not an $E$-function because its
constant coefficient $\log2$ is transcendental, by
Lindemann--Weierstrass.  This is a statement about the function;
it is not a proof that its value at $z=1$ lies outside the ring of
$E$-function values.

\subsection{Periodicization removes monodromy but destroys holonomy}

\begin{lemma}[Two growth facts]
\label{p6:phn:growth}
An entire D-finite function has finite order.  If $g$ is entire
and $\exp(g)$ has finite order, then $g$ is a polynomial.
\end{lemma}
\begin{proof}
For the first assertion, turn a polynomial-coefficient scalar
equation for $f$ into a first-order system for
$V=(f,f',\ldots,f^{(r-1)})$.  Outside a sufficiently large circle,
the rational coefficient matrix is bounded by $C|z|^d$ for some
integer $d\geq0$.  The initial norms of $V$ on that circle are
uniformly bounded.  Gronwall's inequality along every radial segment
then gives $|f(z)|\leq C_1\exp(C_2|z|^{d+1})$.

For the second assertion, finite order supplies an integer $N$ and
$C>0$ such that $\Re(g(z)-g(0))\leq C R^N$ on $|z|\leq R$ for
all large $R$.  Put $h=g-g(0)$ and $M=C R^N$, enlarging $C$ if
necessary.  The holomorphic function $h/(2M-h)$ maps this disk into
the closed unit disk and vanishes at zero: indeed
$|2M-h|^2-|h|^2=4M(M-\Re h)\geq0$.
Schwarz's lemma bounds its absolute value by $1/2$ on
$|z|\leq R/2$, and hence $|h(z)|\leq2M$ there.
Cauchy's estimate gives $h^{(j)}(0)=0$ for every integer $j>N$
by letting $R\to\infty$.  Therefore $g$ is a polynomial.
\end{proof}

\begin{theorem}[No rational periodicization of the tail is D-finite]
\label{p6:phn:periodicization}
Define the single-valued meromorphic function
\[
 Y(z)=\exp(E_1(z))=\frac{e^{-\gamma}}z
                               \exp(\operatorname{Ein}(z)).
\]
For every nonconstant $R\in\mathbb C(U)$, the meromorphic function
$R(Y(z))$ is not D-finite over $\mathbb C(z)$.
In particular, $\exp(\operatorname{Ein}(z))$ and
$\exp(-\operatorname{Ein}(z))$ are not D-finite.
\end{theorem}
\begin{proof}
The displayed expression defines $Y$ globally: it has a simple
pole at zero, no other finite pole, and no zero.  Its singularity
at infinity is essential.  Otherwise $Y$ would be rational and
$zY=e^{-\gamma}\exp(\operatorname{Ein})$ would be a zero-free
polynomial, hence constant, which is impossible.

Suppose first that $R$ has a finite pole $\beta\ne0$.  In every
punctured neighborhood of infinity the function $Y$ omits both
zero and infinity.  Great Picard's theorem therefore says that
it assumes each nonzero finite $\beta$ infinitely often there.
Every such preimage is a pole of $R(Y)$, since a reduced rational
function really has a pole at $\beta$.  A meromorphic D-finite
function has no pole at an ordinary point of its differential
equation, and thus has only finitely many finite poles.  This is a
contradiction.

In the remaining case $R$ has poles only at zero or infinity,
so it is a nonconstant Laurent polynomial.  Consequently $R(Y)$
has at most a pole of finite order at zero and no other finite
pole.  If it were D-finite, multiplication by a sufficiently large
power of $z$ would give an entire D-finite function, hence a
finite-order function by Lemma~\ref{p6:phn:growth}.  Thus $R(Y)$
has an exponential-of-a-power upper bound outside a large disk.

Write $R(U)=\sum_{j=\ell}^q c_jU^j$ with $c_\ell c_q\ne0$.
If $q>0$, dominance of the leading term for large $|U|$ gives
\[
 |U|^q\leq C(1+|R(U)|) \qquad(U\ne0).
\]
(The bounded region requires only enlarging $C$.)  This implies
that the entire function $zY=e^{-\gamma}\exp(\operatorname{Ein})$
has finite order.  If $q\leq0$, nonconstancy implies $\ell<0$;
apply the same argument to the polynomial $R(1/V)$ to deduce
finite order for $1/Y=e^\gamma z\exp(-\operatorname{Ein})$,
and therefore for $\exp(-\operatorname{Ein})$ as well.
Lemma~\ref{p6:phn:growth} would then make $\operatorname{Ein}$ or
$-\operatorname{Ein}$ a polynomial.  Their power series have
infinitely many nonzero coefficients, giving the final contradiction.
The last assertions also follow directly from that lemma (or from
$R(U)=U$ and $R(U)=U^{-1}$ after multiplication by rational
functions).
\end{proof}

For example, the entire Dickman-transform expression
$\exp(\gamma-\operatorname{Ein}(z))=1/(zY(z))$ is not D-finite.
This settles only the proposed passage through D-finite arithmetic
Gevrey functions.  It places no restriction on genuinely
nonholonomic methods for studying its values.

\subsection{Arithmetic status of the extraction step}

The phase embeddings of Fischler--Rivoal take $E$-values into
exponential polynomials with complex coefficients; the cited paper
places the coefficients in a larger module involving resummed
arithmetic Gevrey values, and leaves the image problem open
\cite[Proposition~3.1 and \S6]{p6:FR2026}.
Theorems~\ref{p6:phn:ramification} and \ref{p6:phn:periodicization} do not
supply a map sending an individual phase coefficient back to an
$E$-value.  Nor do they prove a special case of the conjectural
intersection statements for arithmetic Gevrey value rings in
\cite{p6:FischlerRivoal2024Gevrey}.  What is now exact is the distinction:
equal-winding rational cancellation removes a common additive
parameter; ramification can retain the Euler constant but violates
the relevant coefficient bounds; rational periodicization leaves
D-finiteness altogether.  None of these statements resolves the
irrationality or transcendence of $\gamma$.


\section{What regular-singular cancellation does and does not determine}
\label{p6:sec:holonomy-scope}

Write $\mathcal K=\operatorname{Frac}(\mathbf G)$, where $\mathbf G$
is the ring of values of $G$-functions at algebraic points, allowing
analytic continuation.  The path-cancellation proposition in the
attached manuscript is valid with its full-rank hypothesis.  The
hypothesis is essential: it asks that the singular germs of the
functions being cancelled be linearly independent, not merely that
the functions themselves be independent.  We give the precise
replacement when this fails, and then treat polynomial and quotient
dependence on several parameters.  These are structural consequences
of established $G$-function theory, with no claim of publication
priority or of a proof of the irrationality of $\gamma$.

The input from Fischler--Rivoal is the regular-singular expansion and
the arithmetic of connection coefficients; see Theorems~3--4 of
\emph{On the values of $G$-functions},
\url{https://arxiv.org/pdf/1103.6022}.
Their Theorem~2 identifies nonzero fixed scalar cancellation
coefficients with $\mathcal K$.  The local arguments below also
cover several branches and several points.

\subsection{The full affine cancellation locus}

Fix finitely many $G$-functions $F_1,\ldots,F_m$, an algebraic point
$\zeta$, and specified continuations to a punctured neighborhood of
$\zeta$.  Set $x=z-\zeta$.  Their local expansions have the form
\begin{equation}\label{p6:hol:expansion}
 F_j(z)=\sum_{e\in E}\sum_{k=0}^{q}
       x^e(\log x)^k f_{j,e,k}(x),
\end{equation}
where $E\subset\mathbb Q$ is finite, the $f_{j,e,k}$ are convergent
power series, and all their coefficients belong to $\mathcal K$.
One can combine exponents differing by integers so that each
coefficient of each actual power and logarithm is unambiguous.
This follows either directly from the connection theorem, or by
Cramer's rule at a nearby algebraic ordinary point: local Frobenius
bases have algebraic Taylor coefficients, and their connection
coefficients lie in $\mathcal K$.

\begin{lemma}[Finite singular-coefficient matrix]\label{p6:hol:matrix}
There is a matrix $C\in M_{s,m}(\mathcal K)$ such that, for every
$c\in\mathbb C^m$, the germ $\sum_jc_jF_j$ is meromorphic at
$\zeta$ if and only if $Cc=0$.  The same statement holds with
holomorphic in place of meromorphic.  Finitely many prescribed
points and branches can be treated simultaneously by stacking
such matrices.
\end{lemma}
\begin{proof}
Meromorphicity of \eqref{p6:hol:expansion} is equivalent to vanishing
of the coefficients of nonintegral powers and of all positive
logarithmic powers.  Holomorphicity also requires vanishing of
negative integral powers.  These are linear equations in $c$
with coefficients in $\mathcal K$.  Their rows span a subspace
of $\mathcal K^m$, so a finite basis of rows suffices.  This gives
$C$.  Adding finitely many other tests amounts to adding rows.
\end{proof}

\begin{theorem}[Affine classification and identifiable constants]
\label{p6:hol:affine}
Let $A,B_1,\ldots,B_r$ be fixed $G$-functions with the continuations
just specified.  The set of vectors $\xi\in\mathbb C^r$ for which
$A-\sum_i\xi_iB_i$ is meromorphic at the prescribed points is
either empty or an affine space
\[
 \xi_0+\ker(M:\mathbb C^r\longrightarrow\mathbb C^s),
 \qquad \xi_0\in\mathcal K^r,\quad M\in M_{s,r}(\mathcal K).
\]
If a coordinate is constant on this affine space, that coordinate
belongs to $\mathcal K$.  In particular, a unique cancellation
vector belongs to $\mathcal K^r$.

These conclusions remain valid for an arbitrary collection of
exact linear cancellation tests in the same fixed finite vector
of parameters, even if different $G$-functions and paths are used
in different tests.
\end{theorem}
\begin{proof}
Lemma~\ref{p6:hol:matrix} makes the equations $M\xi=b$ with
$M,b$ over $\mathcal K$.  Consistency over $\mathbb C$ implies
consistency over $\mathcal K$ by row reduction.  Choosing such a
solution gives the displayed affine space.  If coordinate $i$
is constant, its value is $(\xi_0)_i\in\mathcal K$.

For an arbitrary collection, take a basis of all the augmented
rows $(M_{j,*},-b_j)\in\mathcal K^{r+1}$.  At most $r+1$ rows
are needed.  Their equations are equivalent to the entire
collection, proving the last assertion.  Holomorphic cancellation
has the identical proof.
\end{proof}

\begin{example}[Function independence is insufficient]
\label{p6:hol:free-example}
Take
\[
 A(z)=B_1(z)=\log(1-z),\qquad B_2(z)=\log(1-z)+z.
\]
The functions $B_1,B_2$ are linearly independent over $\mathbb C$.
Nevertheless
\[
 A-\xi B_1-(1-\xi)B_2=-(1-\xi)z
\]
is holomorphic at $z=1$ for every $\xi\in\mathbb C$, including
$\xi=\gamma$.  The entire cancellation locus is $\xi_1+\xi_2=1$.
Thus the presence of $\gamma$ as one coefficient in a cancellation
does not imply $\gamma\in\mathcal K$.  This is not a counterexample
to the attached proposition: $B_2-B_1=z$ violates its second
hypothesis.  It is a counterexample to omitting that hypothesis
from the subsequent interpretation.  Indeed $\mathcal K$ is
countable (there are countably many algebraic differential equations,
algebraic initial data, algebraic evaluation points and continuation
classes), whereas the example allows every complex $\xi$.
\end{example}

\subsection{Polynomial dependence and infinitely many exact tests}

\begin{theorem}[Algebraic cancellation loci]\label{p6:hol:polynomial}
Let $t=(t_1,\ldots,t_r)$ and consider any collection of tests of
meromorphic or holomorphic continuation of functions
\[
 H_\nu(t,z)=\sum_{j=1}^{m_\nu}P_{\nu,j}(t)F_{\nu,j}(z),
 \qquad P_{\nu,j}\in\mathcal K[t_1,\ldots,t_r],
\]
where, in each test, the $F_{\nu,j}$ are fixed $G$-functions and
the endpoints and continuations are fixed.  The common parameter
set is an affine algebraic set defined over $\mathcal K$.
Even an infinite collection of tests is equivalent to finitely
many polynomial equations.

Every isolated complex solution has coordinates algebraic over
$\mathcal K$.  If the entire common solution set is one complex
point, its coordinates belong to $\mathcal K$ itself.
\end{theorem}
\begin{proof}
For each test, Lemma~\ref{p6:hol:matrix} gives finitely many polynomial
equations with coefficients in $\mathcal K$ after substitution
of the $P_{\nu,j}$.  All these polynomials generate an ideal in
the Noetherian ring $\mathcal K[t_1,\ldots,t_r]$.
A finite set generates that ideal; each generator is a finite
combination of the original equations.  This proves the assertion
also for an infinite collection.

An isolated complex point is a zero-dimensional irreducible
component after base extension to an algebraic closure of
$\mathcal K$.  Zero-dimensional components of an affine
algebraic set over a field have finite residue extensions of
that field, by Zariski's lemma; their coordinates are therefore
algebraic over $\mathcal K$.  Equivalently, elimination on the
zero-dimensional components supplies a nonzero polynomial over
$\mathcal K$ for each coordinate.

If there is only one complex point, every embedding into
$\mathbb C$ of its finite coordinate extension over $\mathcal K$
must preserve that point.  Characteristic zero makes the
extension separable, so its degree is one.  The coordinates
then lie in $\mathcal K$.
\end{proof}

Products of finitely many $G$-functions cause no difficulty: they
are again $G$-functions, so they can be included among the $F_{\nu,j}$.
The theorem concerns exact equations in a fixed finite parameter
vector.  It does not cover changing parameter vectors of unbounded
dimension, approximate cancellation inequalities, or limits of
different cancellation parameters.

\begin{example}[An isolated nonlinear parameter]
\label{p6:hol:square-example}
Define rational-coefficient $G$-functions
\[
 B(z)=-\log(1-z),\qquad
 A(z)=\int_0^z\frac{-\log(1-u/2)}{1-u}\,du.
\]
The $G$-function assertion follows from closure under products
and integration; directly, the coefficient denominators have
only exponential growth.  Differentiation gives
\[
 (A-T^2B)'(z)=\frac{-\log(1-z/2)-T^2}{1-z}.
\]
The numerator is holomorphic at $1$, and its value there is
$\log2-T^2$.  Thus $A-T^2B$ is meromorphic at $1$ if and only
if $T^2=\log2$, and in that case it is holomorphic there.
The two isolated parameters are $T=\pm\sqrt{\log2}$.
This demonstrates the algebraic equation arising from nonlinear
cancellation.  We do not assert that either root lies outside
$\mathcal K$: that would need an additional arithmetic theorem.
Selecting the positive root by an inequality is not the same as
a globally unique complex algebraic solution.
\end{example}

\subsection{Quotients and one-turn monodromy}

\begin{theorem}[Quotient cancellation]\label{p6:hol:quotient}
Fix an algebraic point $\zeta$ and branches of finitely many
$G$-functions there.  Let $N_t,D_t$ be polynomial combinations
in $t\in\mathbb C^r$ of these functions and of their finite
products, with coefficients in $\mathcal K$.  On the parameter
set where $D_t$ is not the zero germ, the set for which $N_t/D_t$
extends meromorphically to $\zeta$ is constructible over
$\mathcal K$.  In particular, any isolated complex parameter
in that set is algebraic over $\mathcal K$.
\end{theorem}
\begin{proof}
Superscript $+$ will denote one positive turn around $\zeta$.
Local expansions \eqref{p6:hol:expansion} and their products have
rational exponents and finitely many logarithmic powers.
Continuation replaces $x^e$ by $e^{2\pi ie}x^e$ and $\log x$
by $\log x+2\pi i$.  All new coefficient constants remain in
$\mathcal K$, because $e^{2\pi ie}$ is algebraic and
$2\pi i\in\mathcal K$.

The identity
\begin{equation}\label{p6:hol:monodromy-identity}
 N_t^+D_t-N_tD_t^+=0
\end{equation}
is therefore equivalent to polynomial equations over
$\mathcal K$ in $t$: equate all local coefficients and use
Noetherianity to retain finitely many.  The condition $D_t=0$
is likewise algebraic.  Its complement is a Zariski open set.

It remains to show that \eqref{p6:hol:monodromy-identity} is
sufficient, not merely necessary, for meromorphic extension.
For fixed admissible $t$, the nonzero denominator has an
expansion with a leading term
\[
 D_t(x)=c x^a(\log x)^b(1+o(1)),\qquad
 c\ne0,\quad a\in\mathbb Q,\quad b\in\mathbb Z_{\geq0}.
\]
The estimate is uniform as $x\to0$ with its argument in any
fixed bounded interval: first choose the smallest actual power
of $x$ with nonzero coefficient and then the largest logarithmic
power at that exponent.  The remaining exponents have a
positive gap, and the convergent series remainders are uniform.
Thus $D_t$ has no zeros in a sufficiently small fundamental
angular sector and the quotient has at most power growth there.
Identity \eqref{p6:hol:monodromy-identity} makes the quotient
single-valued on the punctured disk.  The growth estimate over
an angular interval of length $2\pi$ implies
$|N_t/D_t|\leq C|x|^{-M}$ there for some integer $M$.
Multiplication by $x^M$ and the removable singularity theorem
give meromorphic extension.

The set is consequently an algebraic zero set minus the
algebraic set where $D_t=0$, hence constructible.  At an
isolated point the denominator nonvanishing condition is
locally open, so the point is an isolated algebraic solution
as in Theorem~\ref{p6:hol:polynomial}.
\end{proof}

No holonomicity of the quotient itself is assumed here.
Consequently this theorem is a description of local cancellation
geometry, not a license to insert a quotient into an arithmetic
holonomy theorem whose holonomicity hypothesis must still be checked.

\subsection{The logical boundary of the holonomy inference}

The arithmetic holonomy theorem is a dimension bound for holonomic
series with specified arithmetic and analytic properties; it is
not formulated as a theorem that every target constant must first
be a uniquely determined scalar connection coefficient.
The 2025 sequel of Calegari--Dimitrov--Tang, Theorem~2.1 and
Theorem~2.6, explicitly includes power systems in a target parameter:
\url{https://arxiv.org/html/2510.04156v1}.
These are primary sources for the actual theorem hypotheses.

The justified restriction is this: an \emph{unconditional}, exact,
identifiable linear cancellation of a fixed constant between
$G$-functions forces that constant into $\mathcal K$.
Finite nonlinear cancellation has the algebraic-locus consequences
proved above.  These facts do not establish that every conceivable
application of an arithmetic holonomy bound to $\gamma$ requires
an unconditional representation of that kind.

In particular, if the needed $G$-structure or extra continuation
is constructed only under the assumption $\gamma\in\mathbb Q$,
then a conditional deduction $\gamma\in\mathcal K$ is tautological.
It does not rule out a subsequent contradiction from the dimension
bound.  No such conditional construction is supplied here; the point
is to preserve this logical research possibility rather than mark
it as excluded by the path-cancellation proposition.

Nor does membership in $\mathcal K$ pass to limits.  The field
contains $\mathbb Q$ and is dense, so a convergent sequence of its
elements need not determine the arithmetic nature of its limit.
An explicit regular-to-irregular family illustrates this issue.
For an integer $M\geq1$ put
\begin{equation}\label{p6:hol:confluence-family}
 F_M(z)=\int_0^z\frac{1-(1+t/M)^{-M}}{t}\,dt.
\end{equation}
It is a rational-coefficient $G$-function for each fixed $M$, and
on every compact subset of the complex plane it converges to
$\operatorname{Ein}(z)$ as $M\to\infty$.
Indeed $(1+t/M)^{-M}\to e^{-t}$ uniformly on compact sets,
and the quotients by $t$ have removable singularities at zero
and converge uniformly there as well.

Writing $u=(1+z/M)^{-1}$ and differentiating verifies exactly
\[
 F_M(z)=\log(1+z/M)+\sum_{j=1}^{M-1}\frac{1-u^j}{j}.
\]
Consequently
\[
 \lim_{z\to+\infty}\{F_M(z)-\log z\}
       =H_{M-1}-\log M\longrightarrow\gamma.
\]
For finite $M$ this is a $G$-value constant involving a logarithm;
the limit is a canonical irregular connection constant.
This example does not solve an arithmetic problem.  It shows
exactly why a moving degeneration requires uniform arithmetic
and analytic estimates, and cannot be dismissed by a theorem
about a single fixed cancellation system.  Obtaining a useful
uniform holonomy estimate through such a degeneration remains
an unproved step of a possible new route.

There is an explicit arithmetic cost in this particular normalization.
Writing $F_M(z)=\sum_{n\geq1}a_{M,n}z^n$, binomial expansion gives
\[
 a_{M,n}=\frac{(-1)^{n+1}}{nM^n}\binom{M+n-1}{n}.
\]
If $M=p$ is prime and $1\leq n<p$, the numerator of the binomial
coefficient contains precisely one factor of $p$, whereas $n!$
and $n$ contain none.  Hence
\[
 v_p(a_{p,n})=1-n.
\]
The common denominator of the first $N$ coefficients is therefore
at least $p^{N-1}$ for $1\leq N<p$, at the same time as the
distance to the nearest singularity grows to $p$.
Scaling $z=pu$ removes these powers of $p$ from the coefficients:
the first $N$ denominators then divide $\operatorname{lcm}(1,\ldots,N)$.
But the nearest singularity is now at $u=-1$.
This exact calculation specifies an input cost to be handled by
any uniform degeneration argument.  It is not a comparison with
an optimal conformal holonomy bound and proves no impossibility
for other normalizations or constructions.


\section{Individual cancellation as congruence to the exact kernel}
\label{p6:sec:seed-content}
This section continues the earlier primitive-jet and image-index results;
those results are inputs, not claimed again as new. The elementary Smith
normal form argument below identifies the remaining seed-specific divisor.
It does not establish irrationality of Euler's constant.

\subsection{An exact decomposition of the individual gcd}
For an integer vector $x$, write $\operatorname{cont}(x)$ for the gcd of its
coordinates, taken nonnegative. Let $B\in M_{2,3}(\mathbb Z)$ have rank two,
content one, primitive kernel generator $E\in\mathbb Z^3$, and image index
$I=[\mathbb Z^2:B\mathbb Z^3]$. For $v\notin\mathbb ZE$ set
\[
 c_E(v)=\operatorname{cont}(E\times v),\qquad
 g_B(v)=\operatorname{cont}(Bv).
\]

\begin{theorem}[Individual gcd and quotient content]
\label{p6:thm:seed-content}
One has
\begin{equation}\label{p6:eq:seed-content}
 c_E(v)\mid g_B(v),\qquad g_B(v)/c_E(v)\mid I.
\end{equation}
More precisely, there are unimodular integer matrices $U,V$ with
$UBV=\left(\begin{smallmatrix}1&0&0\\0&I&0\end{smallmatrix}\right)$
and $E=\pm Ve_3$. If $V^{-1}v=(x,y,z)^T$, $c=\gcd(x,y)$,
$x=cx_0$, $y=cy_0$, then
\[
 c_E(v)=c,\qquad g_B(v)=c\gcd(x_0,I).
\]
Moreover $c_E(v)$ is the largest positive integer $d$ such that
$v\in d\mathbb Z^3+\mathbb ZE$.
\end{theorem}
\begin{proof}
The Smith invariant factors are $1,I$, because the matrix content is one.
The integer kernel is generated by $Ve_3$. Under a unimodular coordinate
change the exterior product is transformed by another unimodular matrix,
so its content is unchanged. The product $e_3\times(x,y,z)$ has content
$\gcd(x,y)$. Unimodular left multiplication preserves the content of an
output, whence
\[
 g_B(v)=\gcd(x,Iy)=c\gcd(x_0,Iy_0)=c\gcd(x_0,I),
\]
since $\gcd(x_0,y_0)=1$. Finally $v\in d\mathbb Z^3+\mathbb ZE$
if and only if $d\mid x,y$, which proves the last assertion.
\end{proof}

Consequently every seed admits $v=cw+tE$ with $w\in\mathbb Z^3$ and
$c_E(w)=1$, and $Bv=cBw$. Its rational output is unchanged by this
quotient primitivization. The price is that the new representative $w$
need not have a small height. In the Euclidean metric one can choose it so
that
\begin{equation}\label{p6:eq:lift-height}
 \|w\|_2^2\le
 \frac{\operatorname{dist}(v,\mathbb RE)^2}{c^2}
       +\frac{\|E\|_2^2}{4}.
\end{equation}
Indeed all representatives differ by integer multiples of $E$; choose the
parallel component in the interval $[-\|E\|_2/2,\|E\|_2/2]$. Their
perpendicular component is exactly $v_\perp/c$. This upper bound retains
possibly factorial kernel height and therefore does not solve the lift
problem.

For a prime $p\nmid I$ and $a\ge1$, Theorem~\ref{p6:thm:seed-content} gives
\begin{equation}\label{p6:eq:seed-local}
 p^a\mid g_B(v)\quad\Longleftrightarrow\quad
 v\equiv tE\pmod {p^a}\quad\text{for some }t\in\mathbb Z.
\end{equation}
At primes dividing $I$, the two valuations differ by at most $v_p(I)$.
Thus arbitrarily large individual cancellation outside the image index
is exactly a congruence to the kernel, not an unlocated determinant effect.

\subsection{The CMF arithmetic target}
Use the scalar-contiguous normalization of the rank-three CMF, with trajectory
\[(-A_1,-A_2;b_1,b_2,b_3),\qquad A_i,b_j\in\mathbb Z_{\ge0}.\] Let $\Phi_n$ be its integral transition matrix,
$H_n=\operatorname{adj}(\Phi_n)^T$, $h_n=\operatorname{cont}(H_n)$, and
\[
 B_n=\frac1{h_n}
 \begin{pmatrix}(1,1,0)H_n\\(1,1,1)H_n\end{pmatrix},
 \qquad E_n=(1,-1,0)\Phi_n.
\]
The preceding manuscript established that $E_n$ is primitive, $B_n$ has
content one and kernel $\mathbb ZE_n^T$, and its image quotient is cyclic.
For $A_1=A_2=A$, with
\[
 S=n\sum_{j=1}^3(A+b_j),\quad
 N=\prod_{j=1}^3((A+b_j)n)!,\quad J=N/h_n,
\]
it also established $I_n=J^2\mid d_{S+1}^{\,8}$, where
$d_m=\operatorname{lcm}(1,\ldots,m)$.
The new specialization of \eqref{p6:eq:seed-content} is
\begin{equation}\label{p6:eq:cmf-individual}
 \gcd(B_nv)=\gcd(E_n\times v)\,\rho_n(v),
 \qquad \rho_n(v)\mid I_n\mid d_{S+1}^{\,8}.
\end{equation}
Here vectors are identified with their transposes when needed.
Every quotient-primitive seed, meaning $\gcd(E_n\times v)=1$, thus has
only an exponentially bounded individual gcd. Any factorial-scale gain
must come from this explicit exterior-product gcd.

If $\lambda_n(v)=(-1,\gamma)B_nv$, the primitive form satisfies exactly
\[
 \frac{|\lambda_n(v)|}{I_nc_{E_n}(v)}
 \le \frac{|\lambda_n(v)|}{g_{B_n}(v)}
 \le \frac{|\lambda_n(v)|}{c_{E_n}(v)}.
\]
These inequalities require no assumption about $\gamma$. In particular,
for any seed family with $c_{E_n}(v_n)\le e^{O(n)}$ and a proved lower bound
$|\lambda_n(v_n)|\ge\exp(\kappa n\log n-O(n))$, $\kappa>0$, its primitive
forms diverge. Neither of these two additional estimates is asserted for
all varying CMF seeds.

\subsection{A finite-characteristic reduction for every fixed prime}
Put $D=z\,d/dz-z$ and
\[
 P_n(X)=(X+1)_{A_1n}(X+1)_{A_2n}
             \prod_{j=1}^3(-X)_{b_jn},\qquad R_n=P_n(D)1.
\]
Then $E_n=(R_n(1),DR_n(1),D^2R_n(1))$.

\begin{theorem}[Signed prime-periodicity]
\label{p6:thm:seed-period}
For every prime $p$ and every $n\ge0$,
\[
 E_{n+p}\equiv(-1)^{A_1+A_2}E_n\pmod p.
\]
If $v(n)\in\mathbb Z[n]^3$, the condition
$p\mid\gcd(E_n\times v(n))$ depends only on $n\bmod p$.
It is exactly the condition that the two nonzero or zero vectors
$E_r,v(r)$ are collinear in $\mathbb F_p^3$, $r\equiv n\pmod p$;
$E_r$ itself is always nonzero.
\end{theorem}
\begin{proof}
In $V_p=\mathbb F_p[z]/(z^p-1)$, the operator $D$ is well defined and
satisfies $D^p-D=-1$. For completeness, put $T=d/dz-1$. The Weyl relation
$[T,z]=1$ gives $\prod_{j=0}^{p-1}(zT-j)=z^pT^p$.
On this quotient $(d/dz)^p=0$, hence $T^p=-1$ and $z^p=1$.
For every integer $a$, it follows that
\[
 \prod_{j=0}^{p-1}(D+a+j)=-1,\qquad
 \prod_{j=0}^{p-1}(-D+a+j)=1.
\]
Increasing $n$ by $p$ adds $A_1+A_2$ upper blocks and $\sum b_j$ lower
blocks. All factors commute, so
$P_{n+p}(D)1=(-1)^{A_1+A_2}P_n(D)1$ in $V_p$.
The ideal $(z^p-1)$ is stable under $D$, and evaluation at $z=1$
annihilates it. Applying $D^j$, $j=0,1,2$, therefore proves the jet
congruence, including $p=2$. Polynomial seeds are periodic modulo $p$,
and $E_r\ne0$ follows from the earlier primitive-jet theorem.
\end{proof}

This is an exact finite computation for each fixed prime, not a bound
uniformly over all primes. In particular a frequently tempting shortcut,
that every fixed coordinate seed is quotient-primitive, is false.
For the balanced trajectory $(0,0;1,1,1)$ direct integer products give
\[
 E_7=(-154606591,-255280789,8251597802),
 \quad
 \bigl(c_{E_7}(e_1),c_{E_7}(e_2),c_{E_7}(e_3)\bigr)=(17,1,19),
\]
where the coordinate vectors here are indexed $1,2,3$.
Likewise $c_{E_{14}}(e_1)=139$.
The theorem proves that $17\mid c_{E_n}(e_1)$ for every $n\equiv7\pmod{17}$,
and $139\mid c_{E_n}(e_1)$ for every $n\equiv14\pmod{139}$.
In particular $2363\mid c_{E_n}(e_1)$ whenever $n\equiv1265\pmod{2363}$.
These are infinite congruence families, but they provide only fixed
factors, not the growing factorial cancellation required for irrationality.

\subsection{Aptekarev comparison: what the finite experiment does establish}
The attached note's Aptekarev normalization and explicit formulas agree
with Corollary~4 of the Hessami Pilehroods' paper, arXiv:1010.0429.
The common factor $n!/d_n$ is prior art, stated explicitly on page~2
of that source; it is not a gain created by a new kernel representation.
Recomputing the published integer recurrence exactly verifies the attached
range $3\le\epsilon_n\le39591$ for $n=20,30,\ldots,300$, where
\[
 \epsilon_n=\frac{\gcd(p_n,q_n)}{n!/d_n}.
\]
Extending to every $1\le n\le3000$ gives, for example,
\[
 \epsilon_{300}=147,\quad \epsilon_{1000}=9,\quad
 \epsilon_{2717}=2060427942214275.
\]
The latter is the largest value in this finite range; its logarithm is
about $35.26$, versus about $25295.68$ for the reduced denominator at
that index. Thus the small sampled gcd bound cannot be promoted to a
uniform theorem. These exact computations prove neither an asymptotic
upper bound on $\epsilon_n$ nor the existence of the exceptionally large
cancellation that a successful irrationality proof would need.


\subsection{Reproducible exact certificate for the Aptekarev comparison}
The primary recurrence is equation~(8) of
\url{https://arxiv.org/pdf/0902.1768}; the binomial formulas and the
$n!/d_n$ divisor are in
\url{https://arxiv.org/pdf/1010.0429}.
For $u=p,q$ it reads
\begin{align*}
 (16n-15)u_{n+1}={}&(128n^3+40n^2-82n-45)u_n\\
 &-n^2(256n^3-240n^2+64n-7)u_{n-1}\\
 &+n^2(n-1)^2(16n+1)u_{n-2},\qquad n\ge2,
\end{align*}
with $(p_0,p_1,p_2)=(0,2,31)$ and $(q_0,q_1,q_2)=(1,3,50)$.
The following Python code uses only exact integers and reproduces the
stated exceptional gcd. No value of $\gamma$ enters this certificate.
\begin{lstlisting}
from math import gcd, lcm, factorial
p, q = [0,2,31], [1,3,50]
for n in range(2,2717):
    a = 128*n**3 + 40*n*n - 82*n - 45
    b = -n*n*(256*n**3 - 240*n*n + 64*n - 7)
    c = n*n*(n-1)**2*(16*n+1)
    for u in (p,q):
        v = a*u[-1] + b*u[-2] + c*u[-3]
        assert v % (16*n-15) == 0
        u.append(v // (16*n-15))
d = lcm(*range(1,2718))
base = factorial(2717) // d
g = gcd(p[2717],q[2717])
assert g % base == 0
assert g // base == 2060427942214275
\end{lstlisting}
The complete calculation through $3000$, together with $2000$ exact
checks of \eqref{p6:eq:seed-content} for five CMF trajectories, is in
\texttt{work/gamma\_next\_20261019/probe\_seed\_arithmetic.py}.
The script also independently recomputes $p_{2717},q_{2717}$ from the
binomial--harmonic formulas of Corollary~4 and checks exact agreement
with the recurrence. These finite checks support the reported examples;
the infinite assertions above are established by the written proofs.


\section{A nonlinear attempt: primitive Hankel polynomials}
\label{p6:sec:determinant}

The preceding manuscript constructed the integer affine moments
\[
 \mu_k(c)=k!\{k(H_k-c)+1\},\qquad H_0=0.
\]
Their defining coefficients require no numerical value of $\gamma$.
The canonical Gamma integral identifies them as
\begin{equation}\label{p6:eq:det-moment}
 \mu_k(c)=\int_0^\infty x^k(x-1)(\log x+\gamma-c)e^{-x}\,dx.
\end{equation}
Instead of a linear form, we now try the integer polynomial
\[
 D_N(c)=\det(\mu_{i+j}(c))_{0\leq i,j<N},\qquad
 P_N(c)=D_N(c)/g_N,
\]
where $g_N>0$ is the greatest common divisor of all coefficients of
$D_N$.  Positivity in \eqref{p6:eq:det-moment} gives $P_N(\gamma)>0$.
If these values had $N$th root tending to zero, the family would prove
irrationality: for $\gamma=a/b$, the nonzero integer
$b^{\deg P_N}P_N(a/b)$ has absolute value at least one.
This is a genuinely nonlinear test, with nonvanishing already supplied
by analysis.  It fails even after the \emph{entire} polynomial content
has been removed.

\begin{theorem}[Primitive determinant growth]\label{p6:thm:det-growth}
For the polynomials above,
\[
 P_N(\gamma)>0,\qquad
 \lim_{N\to\infty}P_N(\gamma)^{1/N}=+\infty.
\]
This assertion includes arbitrary rational changes of basis in the
space of polynomials of degree less than $N$, followed by clearing
denominators and removing the full content: their primitive
determinant is $P_N$ up to sign.
\end{theorem}

\begin{proof}
Write $F_N=\prod_{j=0}^{N-1}(j!)^2$.  Let $L_j$ be the ordinary
Laguerre polynomials, orthonormal for $e^{-x}dx$, with leading
coefficient $(-1)^j/j!$.  In this basis the moment matrix is
$A_N-cB_N$, and
\begin{equation}\label{p6:eq:det-basis}
 D_N(c)=F_N\det(A_N-cB_N).
\end{equation}
The matrix $B_N$ is the compression of multiplication by $x-1$:
its diagonal is $2j$, and its adjacent off-diagonal entries are
$-(j+1)$.  In particular $\|B_N\|\leq4N$ and
$\operatorname{rank}B_N\geq N-1$.  The latter follows immediately
by solving the tridiagonal homogeneous system from its first row.

The elementary logarithmic Laguerre integrals are
\begin{equation}\label{p6:eq:det-loglag}
 \int_0^\infty L_i(x)L_j(x)(\log x+\gamma)e^{-x}\,dx
 =\begin{cases}H_i,&i=j,\\-1/|i-j|,&i\ne j.\end{cases}
\end{equation}
For completeness, differentiate the Laguerre orthogonality relation
for the weight $x^\alpha e^{-x}$ at $\alpha=0$, using
$\partial_\alpha L_n^{(\alpha)}|_{\alpha=0}
=\sum_{k<n}L_k/(n-k)$.  The diagonal comes from
$\partial_\alpha[\Gamma(n+\alpha+1)/n!]|_{\alpha=0}=H_n-\gamma$;
the off-diagonal follows from the same differentiated relation.
These identities also follow by integrating the finite polynomial
expansions term by term.
Multiplication by $x-1$ and the three-term Laguerre recurrence now
give $\|A_N\|\leq4N^2(H_N+1)$.

If $B_N$ is nonsingular, the leading coefficient of $D_N$ is nonzero
and has absolute value at most $F_N(4N)^N$.  If $B_N$ is singular,
it has one-dimensional kernel.  On this kernel $A_N$ has positive
quadratic form, since $A_N-\gamma B_N$ is positive definite by
\eqref{p6:eq:det-moment}.  Therefore the coefficient of $c^{N-1}$ is
nonzero.  Its absolute value is bounded by
$F_N N\|A_N\|\|B_N\|^{N-1}$.  Consequently, in both cases,
\begin{equation}\label{p6:eq:det-content}
 g_N\leq F_N N^2(H_N+1)(4N)^N.
\end{equation}
This estimate does not require determining the exact content, or
assuming that $L_N(1)$ is nonzero.

For a nonnegative weight $w$, let
$\mathcal H_N[w]=\det(\int_0^\infty t^{i+j}w(t)dt)_{i,j<N}$.
Its Vandermonde integral representation has nonnegative integrand.
Restrict that representation for $D_N(\gamma)$ to $x_i\geq1$ and
put $x_i=t_i+1$.  For every fixed $a>1$,
\[
 \log(1+t)=\int_0^\infty
       \frac{t\,ds}{(1+s)(1+s+t)}
 \geq \frac{t\log a}{t+a},\qquad t\geq0.
\]
It follows that
\begin{equation}\label{p6:eq:det-lower-weight}
 D_N(\gamma)\geq(e^{-1}\log a)^N
     \mathcal H_N[t^2e^{-t}/(t+a)].
\end{equation}

Normalize the Vandermonde probability density with weight
$t^\alpha e^{-t}$, and denote expectation by $\mathbb E_\alpha$.
Laguerre norms and the characteristic-polynomial identity give
\[
 \mathcal H_N[te^{-t}]=N!F_N,\qquad
 \frac{\mathcal H_N[t^2e^{-t}]}{\mathcal H_N[te^{-t}]}=(N+1)!,
 \qquad
 \mathbb E_2\prod_{i=1}^N(t_i+a)=N!L_N^{(2)}(-a).
\]
The last identity follows directly from the Vandermonde integral:
the average monic characteristic polynomial is orthogonal to every
polynomial of degree less than $N$.  Cauchy--Schwarz, in the form
$\mathbb E X\,\mathbb E X^{-1}\geq1$, therefore yields
\begin{align*}
 \frac{\mathcal H_N[t^2e^{-t}/(t+a)]}{N!F_N}
 &=(N+1)!\,\mathbb E_2\prod_{i=1}^N(t_i+a)^{-1}\\
 &\geq\frac{N+1}{L_N^{(2)}(-a)}
 \geq\frac{2}{N+2}e^{-2\sqrt{aN}}.
\end{align*}
Indeed, the finite Laguerre expansion implies
\[
 L_N^{(2)}(-a)
 =\sum_{k=0}^N\binom{N+2}{k+2}\frac{a^k}{k!}
 \leq\binom{N+2}{2}\sum_{k=0}^N\frac{(aN)^k}{(k!)^2}
 \leq\binom{N+2}{2}e^{2\sqrt{aN}}.
\]
For the final inequality, compare the series with the product
$e^{\sqrt{aN}}e^{\sqrt{aN}}$.

Combining the bounds gives the fully explicit estimate
\begin{equation}\label{p6:eq:det-primitive-lower}
 P_N(\gamma)\geq
 \frac{2e^{-2\sqrt{aN}}}{(N+2)N^2(H_N+1)}
 \left(\frac{\log a}{e}\right)^N\frac{N!}{(4N)^N}.
\end{equation}
Taking $N$th roots and applying Stirling's formula gives
$\liminf P_N(\gamma)^{1/N}\geq\log a/(4e^2)$.
The number $a>1$ was arbitrary and fixed before the limit.  Letting
$a$ grow proves the assertion.  A rational basis change multiplies
the determinant by a nonzero rational square.  Primitive integer
polynomials spanning the same rational line agree up to sign,
which proves the basis-invariance statement.
\end{proof}

\begin{remark}[What this excludes]
This result excludes the full Hankel determinant of this specific
moment pencil as a source of vanishing primitive values, even after
optimal scalar normalization.  It does not exclude proper minors,
different moment functionals, differences of determinants, or
simultaneous constructions with varying auxiliary data.  In
particular it is not an impossibility theorem for nonlinear
irrationality methods in general.
\end{remark}


% Fragment: generalized Dickman moments and an obstruction to p-adic transfer.
\section{Dickman moments, prime congruences, and a failed adelic transfer}
\label{p6:dn:section}

The uploaded manuscript's Dickman attempt concerns rational moments of
a probability density whose normalization contains $e^{-\gamma}$.
Here we test a different bridge: transporting those rational moments
to a compactly supported $p$-adic distribution, then seeking arithmetic
information from its normalization.  We prove that the proposed transport
does not exist, even after rational affine changes of scale.  We also
derive exact prime congruences that sharpen the manuscript's numerical
denominator observations.  These results do not prove the irrationality
or transcendence of $\gamma$; literature priority is not asserted.

For $\vartheta>0$, let $D_\vartheta$ have the generalized Dickman law.
Its moment generating function and density on $(0,1]$ are
\begin{equation}\label{p6:dn:mgf}
 M_\vartheta(z)=\mathbb E e^{zD_\vartheta}
 =\exp\left(\vartheta\sum_{j\ge1}\frac{z^j}{j\,j!}\right),
 \qquad
 p_\vartheta(x)=\frac{e^{-\vartheta\gamma}}{\Gamma(\vartheta)}
 x^{\vartheta-1}\quad(0<x\le1).
\end{equation}
These are standard identities; see \cite[\S1, (1.1)]{p6:PinskyDickman2016}.
In particular $\Pr(D_1\le1)=e^{-\gamma}$.  Write
$m_n(\vartheta)=\mathbb E D_\vartheta^n$.  Formally the same definition
is meaningful for every rational $\vartheta$.  Differentiating
\eqref{p6:dn:mgf} or using the exponential formula gives
\begin{equation}\label{p6:dn:moments}
 m_0(\vartheta)=1,\qquad
 m_n(\vartheta)=\frac{\vartheta}{n}
       \sum_{i<n}\binom ni m_i(\vartheta),\qquad
 m_n(\vartheta)=\sum_{\pi\in\Pi_n}
       \prod_{B\in\pi}\frac{\vartheta}{|B|}.
\end{equation}
Here $\Pi_n$ is the set of partitions of an $n$-element set.

\subsection{Exact congruences behind the moment denominators}

\begin{theorem}[Prime congruences for generalized Dickman moments]
\label{p6:dn:congruence}
Let $p$ be an odd prime and let $0\ne\vartheta\in\mathbb Q\cap\mathbb Z_p$.
Put $c=v_p(\vartheta)$, $q=p^{c+1}$, and
$u=\vartheta/p^c\in\mathbb Z_p^\times$.  If $n=aq+r$ with
$0\le r<q$, then
\begin{equation}\label{p6:dn:congruence-eq}
 p^a m_n(\vartheta)\in\mathbb Z_p,
 \qquad p^a m_n(\vartheta)\equiv u^a m_r(\vartheta)\pmod p.
\end{equation}
In particular $v_p(m_{aq}(\vartheta))=-a$ for every $a\ge0$.
For $\vartheta=1$ this becomes
\[
 p^{\lfloor n/p\rfloor}m_n(1)
       \equiv m_{n\bmod p}(1)\pmod p.
\]
Consequently equality in the elementary denominator bound
$v_p(\operatorname{den}m_n(1))\le\lfloor n/p\rfloor$ holds
whenever $m_{n\bmod p}(1)$ is a $p$-adic unit.
\end{theorem}
\begin{proof}
A block of size $j$ contributes valuation $c-v_p(j)$.
If $w=v_p(j)-c>0$, then
\[
 j\ge p^{c+w}\ge q w.
\]
Since $p\ge3$, equality is possible only for $w=1$ and $j=q$;
every other positive-$w$ block has $j-qw\ge q$.
It follows that the valuation of every term in the partition sum is
at least $-a$.  Terms with valuation $-a$ can occur only among
partitions having exactly $a$ blocks of size $q$, with their remaining
$r$ elements partitioned arbitrarily.  Indeed a different block with
positive $w$, or a total negative valuation exceeding $a$, would use
at least $(a+1)q$ elements.  Any remaining block of positive valuation
vanishes modulo $p$ after the common scaling.

The sum over the indicated partitions is
\[
 N\left(\frac{\vartheta}{q}\right)^a m_r(\vartheta),
 \qquad N=\frac{n!}{(q!)^a a!r!}.
\]
We have $N\in\mathbb Z_p^\times$ and $N\equiv1\pmod p$.
For clarity, Legendre's formula gives $v_p(N)=0$, and the elementary
unit-factorial congruence
\[
 p^{-v_p(k!)}k!\equiv
 (-1)^{v_p(k!)}\prod_i k_i!\pmod p
 \quad\left(k=\sum_i k_i p^i,\ 0\le k_i<p\right)
\]
shows that its unit is $1$: the base-$p$ digits of $n=aq+r$ are
the concatenation of those of $a$ and $r$, while $q$ has a single
digit $1$.  The sign exponents cancel because $v_p(N)=0$.
This unit-factorial formula also follows immediately by repeatedly
removing multiples of $p$ and applying Wilson's theorem.
Multiplying the displayed partition sum by $p^a$ proves
\eqref{p6:dn:congruence-eq}.  When $r=0$ its right side is a unit.
\end{proof}

The residue condition is substantive.  For example
$m_2(1)=3/2$, so the leading congruence vanishes when
$p=3$ and $n\equiv2\pmod3$.  No assertion of equality in all
residue classes, or of the uploaded manuscript's observed
$2$-adic formula, is being made.

\begin{corollary}[Exact radius of the moment generating series]
\label{p6:dn:radius}
Under the hypotheses of Theorem~\ref{p6:dn:congruence}, the $p$-adic
radius of convergence of
$M_\vartheta(z)=\sum_{n\ge0}m_n(\vartheta)z^n/n!$ is exactly
\[
 R_p(M_\vartheta)=p^{-1/q-1/(p-1)}.
\]
\end{corollary}
\begin{proof}
The valuation bound in \eqref{p6:dn:congruence-eq} and
$v_p(n!)/n\to1/(p-1)$ give one inequality in the
Cauchy--Hadamard formula.  The exact valuations at $n=aq$ give
the other inequality.
\end{proof}

\subsection{Torsion characters obstruct every rational affine scaling}

The preceding denominators already exclude a bounded measure with
unscaled moments.  Merely multiplying the random variable by a large
power of $p$ can remove that immediate obstruction.  The following
argument survives such changes.

For rational $\vartheta>0$, $c\ne0$, and $d$, define the rational
polynomial functional
\[
 \mathcal L_{\vartheta,c,d}(P)=\mathbb E P(d+cD_\vartheta),
 \qquad
 a_n=\mathcal L_{\vartheta,c,d}\left(\binom Xn\right),
 \qquad A(T)=\sum_{n\ge0}a_nT^n.
\]
This uses only rational moments, not a real-to-$p$-adic embedding.

\begin{theorem}[No compact $p$-adic distribution with Dickman moments]
\label{p6:dn:no-transfer}
Let $p$ be any prime, $\vartheta\in\mathbb Q_{>0}$,
and $c,d\in\mathbb Q\cap\mathbb Z_p$ with $c\ne0$.
Put $j=v_p(c)$.  The formal Amice series above has radius
\begin{equation}\label{p6:dn:amice-radius}
 R_p(A)\le p^{-1/(p^j(p-1))}<1.
\end{equation}
Hence $\mathcal L_{\vartheta,c,d}$ is not the restriction of a
locally analytic $p$-adic distribution on $\mathbb Z_p$.
In particular it is not a bounded $p$-adic measure.
There is no compactly supported locally analytic distribution on
$\mathbb Q_p$ whose moments are $m_n(\vartheta)$.
\end{theorem}
\begin{proof}
All of the following identities first hold as formal power series
over $\mathbb Q$.  The binomial generating function gives
\[
 A(T)=(1+T)^dM_\vartheta\bigl(c\log(1+T)\bigr).
\]
Since $zM_\vartheta'(z)=\vartheta(e^z-1)M_\vartheta(z)$,
formal differentiation gives
\begin{equation}\label{p6:dn:amice-ode}
 (1+T)\log(1+T)A'(T)=
 \left[d\log(1+T)+\vartheta\bigl((1+T)^c-1\bigr)\right]A(T).
\end{equation}
Both coefficients in this equation are analytic for $|T|_p<1$;
here $c,d\in\mathbb Z_p$ ensures convergence of the binomial powers.

Let $\zeta$ be a primitive $p^{j+1}$-st root of unity and set
$T_0=\zeta-1$.  Then
\[
 |T_0|_p=p^{-1/(p^j(p-1))},\qquad
 \log\zeta=0,
\]
and $\zeta^c$ is a primitive $p$-th root of unity.  The latter
assertion follows by approximating $c$ by integers in $\mathbb Z_p$;
the value depends on $c$ modulo $p^{j+1}$.
Suppose that the radius of $A$ were greater than $|T_0|_p$.
The formal identity \eqref{p6:dn:amice-ode} would then hold in a disk
containing $T_0$.  Since $A(0)=1$, the analytic function $A$ is
not identically zero and has a finite order $k\in\mathbb Z_{\ge0}$
at $T_0$.  The factor $(1+T)\log(1+T)$ has a simple zero there,
with derivative $1$.  Comparing the lowest terms at $T_0$ gives
\begin{equation}\label{p6:dn:resonance}
 k=\vartheta(\zeta^c-1).
\end{equation}
For odd $p$ the right side is not rational, because a primitive
$p$-th root of unity is not rational and $\vartheta\ne0$ is rational.
For $p=2$ it equals $-2\vartheta<0$.
Both alternatives contradict $k\in\mathbb Z_{\ge0}$.
This proves \eqref{p6:dn:amice-radius}.

The Amice transform of any locally analytic distribution on
$\mathbb Z_p$ is analytic on the whole open unit disk
\cite[Theorem~2.2]{p6:SchneiderTeitelbaum2001}; its Taylor coefficients
are its values on $\binom Xn$.  Thus such a distribution cannot
restrict to $\mathcal L_{\vartheta,c,d}$.  The bounded-measure
case also follows elementarily: $\binom xn\in\mathbb Z_p$ for
$x\in\mathbb Z_p$, so a bounded measure gives bounded $a_n$ and
therefore radius at least $1$.
Finally, a compact subset of $\mathbb Q_p$ lies in $p^{-h}\mathbb Z_p$
for some $h\ge0$.  Pushing a putative distribution forward by
$x\mapsto p^h x$ would give the forbidden moments with
$c=p^h,d=0$ on $\mathbb Z_p$.
\end{proof}

This proof identifies the obstruction as a local exponent at a torsion
character, rather than a denominator estimate.  Quantitatively,
\eqref{p6:dn:amice-radius} implies
\[
 \limsup_{n\to\infty}|a_n|_p^{1/n}
   \ge p^{1/(p^j(p-1))}>1.
\]
Multiplication of the functional by any nonzero scalar cannot change
this conclusion.  Thus a scalar normalization involving $e^{-\gamma}$
would not repair the failure even if that scalar had independently
been assigned a $p$-adic value.

\subsection{What remains necessary for an arithmetic bridge}

The real identity $\Pr(D_1\le1)=e^{-\gamma}$ remains exact, but the
rational moment functional cannot be transported to compact $p$-adic
measure theory with those moments intact.  A proposed adelic argument
must therefore change its local objects or supply genuinely different
arithmetic data.  No conclusion here excludes noncompact distributions,
formal functionals, or other nonarchimedean constructions.

The direction of the original bridge also matters: transcendence of
$e^\gamma$ would not imply irrationality of $\gamma$.  One precise
alternative would be to prove
\[
 \Pr(D_1\le1)\ne e^{-r}\quad\text{for every }r\in\mathbb Q,
\]
or the corresponding exclusion for algebraic $r$ to prove
transcendence of $\gamma$.  These are exactly the unresolved
arithmetic assertions, not consequences of the distributional
identities.  The new congruences and the no-transfer theorem identify
what the attempted $p$-adic route can and cannot supply.

% Bibliography entries to add to the combined manuscript:
% \bibitem{PinskyDickman2016} R. G. Pinsky, On the strange domain of attraction
% to generalized Dickman distributions for sums of independent random variables,
% arXiv:1611.07207 (2016), especially Section 1.
% \bibitem{SchneiderTeitelbaum2001} P. Schneider and J. Teitelbaum,
% p-adic Fourier theory, Documenta Math. 6 (2001), 447--481;
% arXiv:math/0102012, Theorem 2.2.


\section{What has changed, and what remains unproved}

The constructive result is Theorem~\ref{p6:nk:complete}: the positive
rational-kernel family is analytically complete when pole order grows.
Thus failure at every tested finite order was not evidence that this
family can never approach the correct target. The theorem also makes
the remaining arithmetic condition exact. It does not estimate the
required order or the relation between a reduced target denominator
and its error.

The ramified connection in Theorem~\ref{p6:phn:ramification} keeps
$\gamma$ after logarithmic monodromy cancellation. Its coefficient
arithmetic fails both the $E$ and $G$ conditions, and the tested block
factorial transforms do not cure this. The holonomy analysis preserves
a different opening: a construction made only under a rationality
hypothesis is not excluded by an unconditional classification of
fixed cancellation constants. No such contradiction-producing
conditional construction has been obtained.

\begin{center}
\small
\begin{tabular}{@{}p{0.28\textwidth}p{0.32\textwidth}p{0.32\textwidth}@{}}
\toprule
Route & Established here & Missing arithmetic endpoint\\
\midrule
Fixed conjugate poles, growing order & All rational targets below
$\gamma$ have strict positive certificates. & A sequence with
$\operatorname{den}(r)(\gamma-r)\to0$.\\[4pt]
Ramified exponential integrals & Entire recessive order-three
connections retain $(m-1)\gamma$ at zero. & An arithmetic theorem
compatible with the mixed coefficient growth.\\[4pt]
Regular-singular holonomy & Exact affine/algebraic cancellation loci;
the conditional construction is not ruled out. & An actual conditional
system meeting all arithmetic and analytic bounds.\\[4pt]
CMF integer lattices & Individual gcd equals quotient content times
a divisor of the image index. & Uniform seed heights, kernel
congruence content and nonzero error estimates.\\[4pt]
Canonical Hankel polynomials & Fully primitive values at $\gamma$
have $N$th root tending to infinity. & This particular determinant
family is unsuitable for shrinking primitive values.\\[4pt]
Dickman moment transfer & Exact odd-prime congruences; compact
$p$-adic realization is impossible after any rational rescaling. &
A different local object or genuinely different arithmetic data.\\
\bottomrule
\end{tabular}
\end{center}

The first row is an equivalence to irrationality, not an easier
assertion proved by the preceding analytic argument. The last two
rows exclude specific proposed constructions, not all determinant
or adelic approaches. No statement in this manuscript proves
$\gamma\notin\mathbb Q$ or $\gamma\notin\overline{\mathbb Q}$.

\subsection{Verification and reproducibility}
Each of the six main proof sections was read independently in
addition to its author's derivation. This internal verification is
not peer review and does not establish priority. Exact arithmetic
checks include the kernel numerator identities, Sturm root counts, exact
Bernstein subdivision for the $15/26$ certificate,
the Aptekarev exceptional gcd at $n=2717$ from both the recurrence and
the independent binomial--harmonic sums, small Hankel/Laguerre
identities, and 6335 instances of the Dickman prime congruences.
Numerical discovery was used to select candidates; it is not used
to establish any infinite-family conclusion.

The standalone source contains the explicit kernel coefficients and
the main exact recurrence verification. The following short routine
independently checks the generalized Dickman congruence without any
decimal evaluation of $\gamma$.
\begin{lstlisting}[language=Python]
from fractions import Fraction as Q
from math import comb

def moments(theta, N):
    out = [Q(1)]
    for n in range(1, N + 1):
        out.append(theta / n * sum(comb(n, k) * out[k]
                                   for k in range(n)))
    return out

def vp(x, p):
    if x == 0:
        return float('inf')
    a, b = abs(x.numerator), x.denominator
    v = 0
    while a % p == 0:
        a //= p; v += 1
    while b % p == 0:
        b //= p; v -= 1
    return v

for theta in map(Q, ['1', '2', '3', '9', '5/2', '-3']):
    m = moments(theta, 100)
    for p in [3, 5, 7, 11, 13]:
        c = vp(theta, p)
        if c < 0:
            continue
        q, u = p**(c + 1), theta / p**c
        for n in range(101):
            a, r = divmod(n, q)
            assert vp(p**a * m[n], p) >= 0
            assert vp(p**a * m[n] - u**a * m[r], p) >= 1
print('Dickman congruence checks passed')
\end{lstlisting}


\begin{thebibliography}{VIII.99}
\bibitem[VI.1]{p6:Attached} User-supplied research manuscript,
\emph{Twelve further attempts at the irrationality of Euler's constant,
and where each one breaks}, 19 October 2026;
\texttt{gamma\_attempts(1).tex} and matching PDF.
\bibitem[VI.2]{p6:Companion} KEIO Research Notes,
\emph{Monodromy, primitive lattices, and positive rational forms
for Euler's constant},
research manuscript, 19 October 2026;
\texttt{Keio\_Gamma\_Monodromy\_Arithmetic\_20261019.tex}.
\bibitem[VI.3]{p6:FR2014} S. Fischler and T. Rivoal,
\emph{On the values of $G$-functions}, Comment. Math. Helv. 89 (2014),
313--341; \url{https://arxiv.org/abs/1103.6022}.
\bibitem[VI.4]{p6:FR2016} S. Fischler and T. Rivoal,
\emph{Arithmetic theory of $E$-operators}, J. Ec. polytech. Math.
3 (2016), 31--65; \url{https://doi.org/10.5802/jep.28}.
\bibitem[VI.5]{p6:FR2026} S. Fischler and T. Rivoal,
\emph{The ring of values of $E$-functions is not a field},
12 October 2026; \url{https://arxiv.org/abs/2610.16603}.
\bibitem[VI.6]{p6:FischlerRivoal2024Gevrey} S. Fischler and T. Rivoal,
\emph{Relations between values of arithmetic Gevrey series, and
applications to values of the Gamma function}, J. Number Theory
261 (2024), 36--54; \url{https://doi.org/10.1016/j.jnt.2024.02.016}.
\bibitem[VI.7]{p6:CDT2025} F. Calegari, V. Dimitrov and Y. Tang,
\emph{Arithmetic holonomy bounds and effective Diophantine approximation},
2025; \url{https://arxiv.org/abs/2510.04156}.
\bibitem[VI.8]{p6:Kalisch2026} R. Kalisch, T. Monderer, E. Leibtag, S. Weinbaum,
M. Shalyt and I. Kaminer,
\emph{An infinite family of Ap\'ery-like recurrences for Euler's constant},
9 October 2026; \url{https://arxiv.org/abs/2610.13919}.
\bibitem[VI.9]{p6:Aptekarev2009} A. I. Aptekarev,
\emph{On linear forms containing the Euler constant}, 2009;
\url{https://arxiv.org/abs/0902.1768}.
\bibitem[VI.10]{p6:HHP2010} Kh. Hessami Pilehrood and T. Hessami Pilehrood,
\emph{Rational approximations for the quotient of gamma values},
Indag. Math. (N.S.) 20 (2009), 583--601;
\url{https://arxiv.org/abs/1010.0429}.
\bibitem[VI.11]{p6:PinskyDickman2016} R. G. Pinsky,
\emph{On the strange domain of attraction to generalized Dickman
distributions for sums of independent random variables},
2016, revised 2017; \url{https://arxiv.org/abs/1611.07207}.
\bibitem[VI.12]{p6:SchneiderTeitelbaum2001} P. Schneider and J. Teitelbaum,
\emph{$p$-adic Fourier theory}, Documenta Math. 6 (2001), 447--481;
\url{https://arxiv.org/abs/math/0102012}.
\bibitem[VI.13]{p6:DLMF} NIST Digital Library of Mathematical Functions,
Chapters 6 and 18, especially Laguerre orthogonality and recurrence;
\url{https://dlmf.nist.gov/18.3}, \url{https://dlmf.nist.gov/18.9}.
\end{thebibliography}
% END PRESERVED SOURCE BODY 6
\phantomsection\label{part6:end}
\clearpage
\setcounter{section}{0}\setcounter{subsection}{0}
\renewcommand{\thesection}{\Roman{part}.\arabic{section}}
\renewcommand{\theHsection}{\arabic{part}.main.\arabic{section}}
\renewcommand{\refname}{References for Part VII}
\fancyhead[L]{\small Euler\textquotesingle s constant: Part VII}
\part{A logarithmic Meijer identity and exact arithmetic structure}
\label{part7:start}
\noindent\textbf{Source:} {\small\path{gamma_exact_structures(1).tex}}.

Original PDF: 28 pages. Mathematical content and source-level qualifications are retained.


% BEGIN PRESERVED SOURCE BODY 7
\begin{abstract}
We continue the supplied attempts on Euler's constant by replacing several
numerical or structural obstacles with exact statements. A Mellin--Barnes
contour shift gives an arbitrary-parameter residue tower whose first
nonconstant output is $\gamma+\Log z$. For a rank-three trajectory in the
October 2026 Meijer construction, we determine the entire projected integer
lattice: it is the image of
$n!\left(\begin{smallmatrix}1&3H_n\\0&1\end{smallmatrix}\right)$.
The proof combines a characteristic-$p$ differential identity with
logarithmic Frobenius solutions and an adjoint pairing. We also determine
the integer kernel of the elementary Gamma integral, prove sharp
pointwise-positivity and uniform absolute-integral barriers, and obtain an
odd-prime congruence yielding exact common-denominator growth for Dickman
moments and their fixed polynomial transforms.

These are proved identities and arithmetic results, not a proof that
$\gamma$ is irrational or transcendental. In particular, unrestricted
selection of an integer initial vector in the exact lattice encodes
arbitrary rational approximation and does not preserve a fixed
trajectory's error estimate. A separate height-controlled,
nonvanishing small-form construction remains missing. The results were
derived and checked for this continuation; historical priority is not
claimed.
\end{abstract}

\medskip
\noindent\textbf{Status of the main problem.}
No irrationality or transcendence proof is claimed anywhere in this
supplement. Statements marked as computations are finite exact checks or
ordinary numerical quadratures; the latter are not interval certificates.

\medskip
\noindent\textbf{Relation to the supplied files.}
This is a self-contained supplement to \texttt{gamma\_attempts.tex} and the
companion manuscript bundle supplied with the request. It does not silently
replace their statements or treat their unpublished claims as established
theorems.



\section{What a successful construction must still achieve}
\label{p7:sec:introduction}

Write
\[
 H_0=0,\qquad H_n=\sum_{j=1}^n\frac1j,\qquad
 \gamma=\lim_{n\to\infty}(H_n-\log n).
\]
The current preprint used here does not resolve the irrationality of
$\gamma$ \cite{p7:Kalisch2026}. We therefore separate two tasks: constructing
new exact objects, and proving that their arithmetic and analytic
properties contradict a rationality assumption. Only the former is
completed in this note.

\begin{lemma}[The required arithmetic conclusion]
\label{p7:lem:smallform}
Suppose $a_j,b_j\in\Z$ and
\[
 0<|a_j-b_j\gamma|\longrightarrow0
\]
along an infinite sequence. Then $\gamma\notin\Q$.
More generally, for $A_j,B_j\in\Q$ and positive integers $D_j$ with
$D_jA_j,D_jB_j\in\Z$, it suffices to prove
$0<D_j|A_j-B_j\gamma|\to0$.
\end{lemma}
\begin{proof}
If $\gamma=u/v$ with $u,v\in\Z$ and $v>0$, then
$v(a_j-b_j\gamma)$ is a nonzero integer tending to zero, a contradiction.
Apply this to $(D_jA_j,D_jB_j)$ for the second assertion.
\end{proof}

The strict inequality matters. Geometry of numbers can supply a small
linear form that is allowed to vanish; a determinant can certify that two
independent forms cannot both vanish. Neither assertion alone supplies
nonvanishing and smallness for the same selected form.

Our strongest arithmetic calculation is
\begin{equation}
\label{p7:eq:intro-lattice}
 \{(p_nv,q_nv):v\in\Z^3\}
 =
 \{(n!a+3n!H_nb,\ n!b):a,b\in\Z\},
\end{equation}
where the rows $p_n,q_n$ come from the explicit matrix product in
Section~\ref{p7:sec:lattice}. This resolves the common-factor question for
the full family of projected outputs, rather than for a few selected
columns. Yet, for $b\ne0$, its ratios are $3H_n+a/b$, which cover all of
$\Q$ at every fixed $n$. Thus a useful adaptive choice of $v$ must carry
additional, proved control of its size and its direction. Without such
control, choosing a good output ratio merely restates the original
Diophantine problem.

The analytic construction has a different purpose. In
Section~\ref{p7:sec:residue} we construct a tower with outputs
$\gamma+\Log z$, its powers, and zeta values. This gives the
identity extension suggested in \cite[Section 9.2(1)]{p7:Kalisch2026}.
It does not establish convergence for all parameter-dependent
trajectories. We then try a coefficient Borel transform, including under
the hypothesis that $\gamma$ is algebraic, and identify precisely why
its logarithmic singularity does not cancel.

The two remaining constructions test potential arithmetic escapes.
The integer Gamma kernel permits large signed changes of integrands
without changing their linear forms; its exact quotient and positivity
bounds appear in Section~\ref{p7:sec:kernel}. Dickman moments offer rational
moments of a distribution whose normalization contains $e^\gamma$;
Section~\ref{p7:sec:dickman} determines the denominator cost of fixed
polynomial transforms and integer convolution shapes. Each conclusion is
restricted to its stated family.



\section{A logarithmic Meijer residue tower}
\label{p7:sec:residue}

Here $G_p$ denotes a \emph{Meijer} function, defined by its
Mellin--Barnes integral. It must not be confused with a \emph{Siegel
$G$-function}, which has additional arithmetic coefficient conditions.
The integral convention agrees with the standard Meijer convention
\cite{p7:DLMF16}. Put $\theta=z\,d/dz$.

\begin{definition}
For an integer $p\ge1$, let
\[
 \Phi_p(s)=\Gamma(-s)^{p+1}\Gamma(1+s)^p,\qquad
 G_p(z)=\frac1{2\pi i}\int_{c-i\infty}^{c+i\infty}
                  \Phi_p(s)z^s\dd s,\quad -1<c<0.
\]
Initially take the principal logarithm on $\C\setminus(-\infty,0]$.
More generally, the same integral is valid on the logarithmic covering
of $\C^\times$ whenever
$|\operatorname{Im}\Log z|<(p+\tfrac12)\pi$.
\end{definition}

\begin{theorem}[Residue tower]
\label{p7:thm:residue}
For $0\le k\le p-1$, define
\[
 \mathcal L_{p,k}
 =z(1+\theta)^{p-1-k}+(-1)^p\theta^{p-k},
\]
with multiplication by $z$ on the left. Then
\begin{equation}
\label{p7:eq:residue}
 \boxed{\quad
 \mathcal L_{p,k}G_p(z)
 =[t^k]\,
 \Gamma(1-t)\left(\frac{\pi t}{\sin\pi t}\right)^p
 e^{t\Log z}.
 \quad}
\end{equation}
In particular,
\begin{align}
 \bigl[z(1+\theta)^{p-1}+(-1)^p\theta^p\bigr]G_p(z)
 &=1 &&(p\ge1),\label{p7:eq:residue-one}\\
 \bigl[z(1+\theta)^{p-2}+(-1)^p\theta^{p-1}\bigr]G_p(z)
 &=\gamma+\Log z &&(p\ge2).\label{p7:eq:residue-gamma}
\end{align}
The identities hold throughout the direct integral domain and continue
analytically along compatible branches.
\end{theorem}

\begin{proof}
For $s=c+iy$, Stirling's estimate gives
\[
 |\Phi_p(c+iy)|
 =O\left((1+|y|)^{-c-1/2}
             e^{-(p+1/2)\pi|y|}\right).
\]
Also $|z^{c+iy}|=|z|^c e^{-y\operatorname{Im}\Log z}$. Hence the
integral and the finitely many differentiated integrals below converge
absolutely and locally uniformly in the stated domain. The same estimate,
uniform on $c-1\le\operatorname{Re}s\le c$, makes the horizontal sides
of the contour shift tend to zero.

The Gamma recurrence gives the exact relation
\[
 \Phi_p(s-1)=(-1)^{p+1}s\Phi_p(s).
\]
For a polynomial $Q$, set
\[
 I_d[Q]=\frac1{2\pi i}\int_{d-i\infty}^{d+i\infty}
                           Q(s)\Phi_p(s)z^s\dd s.
\]
Applying the preceding recurrence and then substituting $s=u+1$ yields
\[
 (-1)^p\theta Q(\theta-1)G_p=-zI_{c-1}[Q].
\]
It follows that
\begin{equation}
\label{p7:eq:shift-difference}
 \{zQ(\theta)+(-1)^p\theta Q(\theta-1)\}G_p
 =z\{I_c[Q]-I_{c-1}[Q]\}.
\end{equation}
The rectangle with right side traversed upwards has positive
orientation. Its strip contains just the pole $s=-1$, so the
right-minus-left difference is the residue there, with positive sign.

Choose $Q(s)=(s+1)^{p-1-k}$ and put $t=s+1$. Using the reflection
formula, the integrand in the residue, including the outside factor $z$,
is
\[
 z(s+1)^{p-1-k}\Phi_p(s)z^s
 =t^{-k-1}\Gamma(1-t)
      \left(\frac{\pi t}{\sin\pi t}\right)^p e^{t\Log z}.
\]
Its residue is the coefficient in \eqref{p7:eq:residue}, and
$\theta Q(\theta-1)=\theta^{p-k}$, proving the theorem.
\end{proof}

\subsection{Explicit output polynomials and the operator factorization}

Set $X=\gamma+\Log z$ and $a_p=2p+1$. Define polynomials $C_{p,k}$ by
\begin{equation}
\label{p7:eq:appell-generating}
 \sum_{k\ge0}C_{p,k}(X)t^k
 =
 \exp\left(
 Xt+\sum_{m\ge2}
       \frac{1+p(1+(-1)^m)}{m}\zeta(m)t^m
 \right).
\end{equation}
The Taylor expansion of $\log\Gamma(1-t)$ \cite{p7:DLMF5} and the sine
product identify the right side with the analytic factor in
\eqref{p7:eq:residue}. Thus $\mathcal L_{p,k}G_p=C_{p,k}(X)$ for
$k\le p-1$. In low degrees,
\begin{align*}
 C_{p,0}(X)&=1,\\
 C_{p,1}(X)&=X,\\
 C_{p,2}(X)&=\frac{X^2+a_p\zeta(2)}2,\\
 C_{p,3}(X)&=\frac{X^3+3a_pX\zeta(2)+2\zeta(3)}6,\\
 C_{p,4}(X)&=\frac{
 X^4+6a_pX^2\zeta(2)+8X\zeta(3)
 +3a_p^2\zeta(2)^2+6a_p\zeta(4)}{24}.
\end{align*}
The polynomial in degree $k$ is defined for every $k$, but its displayed
differential realization requires $p\ge k+1$. No negative power of
$1+\theta$ is being introduced.

The commutation rule $\theta z=z(\theta+1)$ gives
\[
 \theta\mathcal L_{p,k}=\mathcal L_{p,k-1}\quad(1\le k\le p-1),
 \qquad C'_{p,k}=C_{p,k-1},
\]
and, with $\sigma=(-1)^p$,
\begin{equation}
\label{p7:eq:module-factor}
 \mathcal H_p
 :=z(1+\theta)^p+\sigma\theta^{p+1}
 =\theta^p\circ(z+\sigma\theta),\qquad
 \mathcal H_pG_p=0.
\end{equation}
In particular, the top member of the residue tower is the first-order
inhomogeneous equation
\begin{equation}
\label{p7:eq:first-order}
 (z+\sigma\theta)G_p
 =C_{p,p-1}(\gamma+\Log z).
\end{equation}
The order of the factors in \eqref{p7:eq:module-factor} is essential.

At $z=1$, the constant and linear outputs recover the identity pattern
used by the Meijer construction \cite{p7:Kalisch2026}. Formula
\eqref{p7:eq:residue-gamma} provides its requested arbitrary-$z$ identity
extension. Its coefficients contain $z$, so the fixed coefficient
vectors at $z=1$ cannot simply be reused unchanged. At algebraic $z$ the
coefficients are algebraic; at rational $z$ they can be cleared to
integers. This is an exact evaluation statement. It supplies no
automatic convergence theorem or height bound for the associated
parameter-dependent matrix trajectories.

\subsection{A positive Laplace representation}

\begin{proposition}
\label{p7:prop:positive-meijer}
For $\operatorname{Re}z>0$,
\[
 G_p(z)=\int_0^\infty e^{-zt}h_p(t)\dd t,
\]
where
\begin{align}
 h_{2m+1}(t)
 &=\frac{\displaystyle\prod_{j=1}^{m}
       ((\log t)^2+(2j-1)^2\pi^2)}
       {(2m)!(1+t)} &&(m\ge0),\label{p7:eq:h-odd}\\
 h_{2m}(t)
 &=\frac{\displaystyle\log t\prod_{j=1}^{m-1}
       ((\log t)^2+(2j)^2\pi^2)}
       {(2m-1)!(t-1)} &&(m\ge1).\label{p7:eq:h-even}
\end{align}
The value at $t=1$ in the even case is removable. Every $h_p$ is strictly
positive on $(0,\infty)$, and hence
$(-1)^jG_p^{(j)}(z)>0$ for $z>0$ and $j\ge0$.
\end{proposition}
\begin{proof}
For $0<\operatorname{Re}u<1$, the Mellin transform of
$h_1(t)=1/(1+t)$ is $F_1(u)=\pi/\sin\pi u$. Mellin convolution gives
\[
 h_2(t)=\int_0^\infty\frac{\dd v}{(1+v)(v+t)}
       =\frac{\log t}{t-1},
\]
with transform $F_2(u)=(\pi/\sin\pi u)^2$. For
$F_p(u)=(\pi/\sin\pi u)^p$, differentiation gives
\[
 F_p''(u)+p^2\pi^2F_p(u)=p(p+1)F_{p+2}(u).
\]
Thus
\[
 h_{p+2}(t)
 =\frac{(\log t)^2+p^2\pi^2}{p(p+1)}h_p(t)
\]
has Mellin transform $F_{p+2}$. Iteration proves
\eqref{p7:eq:h-odd}--\eqref{p7:eq:h-even}. The Mellin transform in $z$ of the
claimed Laplace integral equals
\[
 \Gamma(u)F_p(1-u)
 =\Gamma(u)^{p+1}\Gamma(1-u)^p
 \qquad(0<\operatorname{Re}u<1),
\]
which is the defining Mellin--Barnes transform of $G_p$ under $u=-s$.
Mellin inversion proves the identity; absolute convergence extends it
throughout $\operatorname{Re}z>0$. Positivity and differentiation under
the integral give the final assertion.
\end{proof}

For $p=2$ there is also an elementary check of the main identity:
\[
 (z+\theta)G_2(z)
 =z\int_0^\infty e^{-zt}(1-t)\frac{\log t}{t-1}\dd t
 =-z\int_0^\infty e^{-zt}\log t\dd t
 =\gamma+\log z \quad(z>0).
\]
Thus the appearance of $\gamma+\log z$ has a direct integral explanation.
The positive kernels do not by themselves make the differentiated
integer value forms both nonzero and small.

\section{Trying to remove the irregularity}
\label{p7:sec:borel}

A natural next attempt is to place the new identity in the arithmetic
theory of Siegel $G$-functions. The exact factorization above lets us
test this possibility without assuming anything about the arithmetic
nature of $\gamma$.

\begin{theorem}[No Siegel $G$-solution of this particular module]
\label{p7:thm:no-siegel}
For every $p\ge1$, the equation
\[
 \mathcal H_pf=\theta^p(z+(-1)^p\theta)f=0
\]
has no nonzero solution which is a Siegel $G$-function germ. The statement
includes germs at an algebraic ordinary point, after translating that
point to the origin.
\end{theorem}
\begin{proof}
The minimal differential equation of a Siegel $G$-function is Fuchsian
with rational exponents; in particular its solutions have convergent
Frobenius-logarithmic expansions at infinity
\cite[Theorem 6 and the paragraph following it]{p7:FR2014}. This means a
finite sum of rational powers of $1/z$, powers of $\Log z$, and power
series in $1/z$ with positive convergence radius.

Suppose $f\ne0$ were such a germ and put $\sigma=(-1)^p$.
On a simply connected logarithmic domain,
\[
 (z+\sigma\theta)f=P(\Log z)
\]
for some polynomial $P$ of degree at most $p-1$.
If $P=0$, then $f=Ce^{-\sigma z}$ with $C\ne0$, which has an irregular
singularity at infinity and cannot have the required convergent
Frobenius expansion.

If $P\ne0$, collect the powers in a convergent Frobenius expansion of
$f$ into exponent classes modulo $\Z$. A nonintegral exponent class
would have to solve the homogeneous equation: its largest power of
$z$ produces a still larger power under multiplication by $z$ that
$\theta$ cannot cancel. Such a class must therefore vanish.
In the integral class, the same leading-power comparison forces
\[
 f(z)=\sum_{n\ge1}z^{-n}A_n(\Log z),
\]
where the $A_n$ have uniformly bounded degree and each series of
coefficients of a fixed log power converges. Comparison in the
first-order equation gives, with $D=d/dX$,
\begin{equation}
\label{p7:eq:factorial-rec}
 A_1(X)=P(X),\qquad
 A_{n+1}(X)=\sigma(n-D)A_n(X).
\end{equation}
If $P$ has degree $d$ and leading coefficient $a\ne0$, the coefficient
of $X^d$ in $A_{n+1}$ is $\sigma^n n!a$. Its series in $1/z$ has radius
zero, contradicting convergence. Translation of an algebraic ordinary
point does not change this required behavior at infinity.
\end{proof}

\begin{remark}
This theorem excludes nonzero Siegel $G$-functions in one explicit
solution space. It does not determine the arithmetic nature of any
individual value $G_p(\alpha)$, nor does it exclude other differential
modules, transforms, or arithmetic-holonomy constructions. The
factorization even has a regular-singular quotient represented by
polynomials in $\Log z$.
\end{remark}

\subsection{A coefficient Borel transform under an algebraicity hypothesis}

Equation~\eqref{p7:eq:factorial-rec} suggests dividing the coefficients by
$n!$. This genuinely changes the function class, so it is not excluded
by Theorem~\ref{p7:thm:no-siegel}. Keep $X$ as an independent polynomial
symbol and define
\[
 \widehat f(z)=
 \sum_{n\ge0}z^{-n-1}\sigma^n(1-D)_nP(X),
 \qquad
 (1-D)_n=\prod_{j=1}^n(j-D).
\]

\begin{proposition}[Exact coefficient Borel symbol]
\label{p7:prop:borel-symbol}
As a formal power series in $t$ with polynomial coefficients in $X$,
\begin{equation}
\label{p7:eq:borel-symbol}
 \sum_{n\ge0}
 \frac{\sigma^n(1-D)_nP(X)}{n!}t^n
 =
 \frac{P(X+\log(1-\sigma t))}{1-\sigma t}.
\end{equation}
\end{proposition}
\begin{proof}
The binomial identity gives
$\sum_{n\ge0}(1-D)_n(\sigma t)^n/n!=(1-\sigma t)^{D-1}$.
Since $\exp(aD)$ translates a polynomial by $a$, applying the operator
to $P$ gives \eqref{p7:eq:borel-symbol}.
\end{proof}

This is a \emph{coefficient} Borel symbol. A Borel transformation of a
full logarithmic transseries also acts on $\Log z$; no such
transformation or summability theorem is implicitly asserted here.

For $p=2$ one has $\sigma=1$ and $P(X)=\gamma+X$, so
\[
 \widehat G_2(z)
 =\sum_{n\ge0}\frac{n!}{z^{n+1}}
        (\gamma+\Log z-H_n).
\]
Its coefficient Borel symbol is
$(\gamma+X+\log(1-t))/(1-t)$. Now make the actual reductio assumption
$\gamma=a\in\Qbar$ and specialize the independent symbol $X$ to zero.
The candidate becomes
\begin{equation}
\label{p7:eq:conditional-G}
 B_a(t)=\frac{a+\log(1-t)}{1-t}
       =\sum_{n\ge0}(a-H_n)t^n.
\end{equation}
For every algebraic $a$, this is a Siegel $G$-function: it is
holonomic, its coefficients have at most logarithmic growth apart from
the fixed algebraic constant, and their denominators are bounded by a
fixed denominator times $\lcm(1,\ldots,n)$, hence exponentially.

However, its logarithmic singularity at $t=1$ is nonzero for
\emph{every} $a$. A positive loop about $1$ changes $B_a$ by
$2\pi i/(1-t)$, independently of $a$. The parameter only changes the
simple-pole term. Thus the algebraicity assumption does produce a
$G$-function, but it creates no extra holomorphic cancellation at the
singularity. This particular attempted escape is exhausted by the
explicit formula; it does not prove that every hypothesis-dependent
$G$-function construction must fail.



\section{Exact arithmetic of a rank-three matrix trajectory}
\label{p7:sec:lattice}

We determine the complete image lattice of the rank-three trajectory
$(0,0;1,1,1)$ in \cite[Section~8.4]{p7:Kalisch2026}.  The calculation concerns
all integer input vectors at a fixed depth.  It therefore distinguishes
the arithmetic freedom of changing the input from the convergence of a
particular sequence of inputs.

For $n\geq0$, let
\begin{equation}\label{p7:eq:lat-M}
M_n=
\begin{pmatrix}
1&n^3&3(n+1)(n-1)(n^2+n+1)\\
2&1-3n^2&-(n+1)(8n^2+n+5)\\
1&3n+2&6n(n+1)
\end{pmatrix}
\end{equation}
and
\begin{equation}\label{p7:eq:lat-N}
N_n=
\begin{pmatrix}
6n^2+13n+10&-8n-5&3\\
3n^3+3n^2-3n-6&-3n^2+3n+3&n-2\\
n^4-2n^3-6n^2-3n+3&6n^2+2n-1&1-2n
\end{pmatrix}.
\end{equation}
Direct polynomial multiplication gives
\begin{equation}\label{p7:eq:lat-matrix-identities}
\det M_n=\det N_n=(n+1)^6,\qquad
N_n=(n+1)^4M_n^{-T},\qquad
\operatorname{adj}N_n=(n+1)^2M_n^T.
\end{equation}
Products are ordered from left to right with increasing index:
\[
P_n=N_0\cdots N_{n-1},\qquad
U_n=M_0\cdots M_{n-1},\qquad P_0=U_0=I_3.
\]
Define
\[
R=\begin{pmatrix}1&1&0\\1&1&1\end{pmatrix},\qquad
\begin{pmatrix}\mathbf p_n\\\mathbf q_n\end{pmatrix}=RP_n,
\qquad \Lambda_n=RP_n\Z^3\subset\Z^2.
\]
We put $H_0=0$ and $H_n=\sum_{j=1}^n1/j$ for $n\geq1$.

\begin{theorem}[Exact projected lattice]\label{p7:thm:lat-main}
For every $n\geq0$,
\begin{equation}\label{p7:eq:lat-exact}
\Lambda_n
=\bigl\{(n!a+3n!H_nb,\ n!b)^T:a,b\in\Z\bigr\}.
\end{equation}
In particular, $[\Z^2:\Lambda_n]=(n!)^2$.  Its column Hermite normal
form, with the upper-right entry in $[0,n!)$, is
\begin{equation}\label{p7:eq:lat-HNF}
\begin{pmatrix}
n!&n!\bigl(3H_n-\lfloor3H_n\rfloor\bigr)\\
0&n!
\end{pmatrix}.
\end{equation}
If $D_n=\den(3H_n)$ is the positive denominator in lowest terms
(with $\den(0)=1$) and $C_n=3D_nH_n$, then
\begin{equation}\label{p7:eq:lat-content}
\gcd(q_{n,1},q_{n,2},q_{n,3})=n!,\qquad
\gcd(p_{n,1},p_{n,2},p_{n,3})=
\gcd(\text{all entries of }RP_n)=\frac{n!}{D_n}.
\end{equation}
After division by this common gcd, the image lattice is exactly
\begin{equation}\label{p7:eq:lat-normalized}
\bigl\{(D_na+C_nb,\ D_nb)^T:a,b\in\Z\bigr\},
\end{equation}
and has index $D_n^2$ in $\Z^2$.
\end{theorem}

The proof has two independent arithmetic components: an exact covolume
calculation and integral formulas for the two rows after a harmonic shear.

\subsection{A primitive polynomial companion and the covolume}

\begin{lemma}\label{p7:lem:lat-covolume}
The gcd of the three maximal minors of $RP_n$ is $(n!)^2$.
Moreover, $\det P_n=(n!)^6$, and the gcd of all $2\times2$ minors of
$P_n$ is also $(n!)^2$.
\end{lemma}

\begin{proof}
Equation~\eqref{p7:eq:lat-matrix-identities} implies
\[
P_n=(n!)^4U_n^{-T},\qquad \det P_n=(n!)^6.
\]
The cross product of the rows of $R$ is $\mathbf v_0=(1,-1,0)$.
The transformation rule for cross products of row vectors therefore gives
\begin{equation}\label{p7:eq:lat-cross}
\operatorname{cross}(\mathbf p_n,\mathbf q_n)
=(n!)^2\mathbf b_n,\qquad \mathbf b_n=\mathbf v_0U_n\in\Z^3.
\end{equation}
It remains to prove that $\mathbf b_n$ is primitive.

Set $\theta=z\dd/\dd z$, $\delta=\theta-z$, and define integer
polynomials
\begin{equation}\label{p7:eq:lat-K}
K_n(z)=\left[\prod_{j=0}^{n-1}(j-\delta)\right]^3 1,
\qquad K_0=1.
\end{equation}
All factors in the product are polynomials in the same operator and hence
commute.  In characteristic zero, the gauge identity
$e^z\theta e^{-z}=\delta$ gives
\[
F_n(z):=e^{-z}K_n(z)
=\left[\prod_{j=0}^{n-1}(j-\theta)\right]^3e^{-z}
=z^n\sum_{k\geq0}\frac{(-1)^k(n+k)!^2}{k!^3}z^k.
\]
The series directly verifies
\begin{equation}\label{p7:eq:lat-forward}
L_nF_n=0,\qquad F_{n+1}=z(\theta+1)^2F_n,
\qquad L_n=(\theta-n)^3+z(\theta+1)^2.
\end{equation}
At $n=0$ the noncommutative factorization is
$L_0=\theta^2(\theta+z)$, since $\theta z=z(\theta+1)$.
Differentiate the second identity twice and eliminate $\theta^3F_n$ and
$\theta^4F_n$ using the first.  At $z=1$ the resulting row of three
jets transforms by right multiplication by $M_n$ in
\eqref{p7:eq:lat-M}.  Its initial value, multiplied by $e$, is
$(1,-1,0)$.  Consequently,
\begin{equation}\label{p7:eq:lat-integer-jets}
\mathbf b_n=\bigl(K_n(1),\delta K_n(1),\delta^2K_n(1)\bigr).
\end{equation}
This is an identity of integer vectors.  The exponential gauge is used
only to establish it in characteristic zero.

Fix a prime $\ell$ and now regard $\delta$ as an operator on
$\mathbb F_\ell[z]$.  We claim that
\begin{equation}\label{p7:eq:lat-frobenius}
\delta^\ell-\delta=-z^\ell.
\end{equation}
Indeed, $\delta z=z(\delta+1)$ gives
$[\delta^\ell-\delta,z]=0$ in characteristic $\ell$.  An operator
commuting with multiplication by $z$ is multiplication by its value on
$1$.  The elementary Stirling-number identity
\[
\delta^m1=\sum_{j=1}^m S(m,j)(-z)^j
\]
shows that this value is $-z^\ell$: for $1<j<\ell$,
\[
j!S(\ell,j)=\sum_{i=0}^j(-1)^{j-i}\binom ji i^\ell\equiv0\pmod\ell,
\]
because Fermat's identity replaces $i^\ell$ by $i$, whose $j$th finite
difference vanishes.  The endpoint Stirling numbers both equal $1$.
This proves \eqref{p7:eq:lat-frobenius}, including $\ell=2$.

It follows that
\[
\prod_{j=0}^{\ell-1}(j-\delta)=z^\ell.
\]
Multiplication by $z^\ell$ commutes with $\delta$.  Write
$n=a\ell+r$ with $0\leq r<\ell$.  Splitting the operator product into
blocks of length $\ell$ in \eqref{p7:eq:lat-K} yields
\[
K_n(z)\equiv z^{3a\ell}K_r(z)\pmod\ell,
\qquad \mathbf b_n\equiv\mathbf b_r\pmod\ell.
\]
Every $M_j$ with $0\leq j<r$ is invertible over $\mathbb F_\ell$,
since $\det M_j=(j+1)^6$.  Hence
$\mathbf b_r=\mathbf v_0M_0\cdots M_{r-1}$ is nonzero modulo $\ell$.
No prime divides all entries of $\mathbf b_n$, proving primitivity.
Equation~\eqref{p7:eq:lat-cross} proves the assertion about $RP_n$.

Finally, $\operatorname{adj}N_j=(j+1)^2M_j^T$ shows that every
$2\times2$ minor of $N_j$ is divisible by $(j+1)^2$.  Taking exterior
squares of the product proves divisibility of every $2\times2$ minor of
$P_n$ by $(n!)^2$.  By Cauchy--Binet, every maximal minor of $RP_n$ is an
integer combination of these minors.  The equality already proved for
$RP_n$ forces their gcd to be exactly $(n!)^2$.
\end{proof}

The index of the column image of a rank-two integer matrix is the gcd of
its maximal minors.  Thus Lemma~\ref{p7:lem:lat-covolume} already determines
$[\Z^2:\Lambda_n]$; the following calculation determines its precise
position inside $\Z^2$.

\subsection{Integral polynomial rows and their adjoint interpretation}

Put $D=u\dd/\dd u$ and define
\begin{align}
c_{n,k}&=\binom nk^3k!,&
Q_n(u)&=\sum_{k=0}^n c_{n,k}u^k,\label{p7:eq:lat-Q}\\
A_n(u)&=\sum_{k=0}^n c_{n,k}
 \bigl[3(H_{n-k}-H_n)-2H_k\bigr]u^k.\label{p7:eq:lat-A}
\end{align}
Both polynomials belong to $\Z[u]$.  For $A_n$, integrality holds
termwise: $c_{n,k}H_k$ is integral because $k!H_k$ is integral, while
\[
c_{n,k}=\binom nk^2\frac{n!}{(n-k)!}
\]
is divisible by each integer $n-k+1,\ldots,n$.  Therefore
$c_{n,k}(H_n-H_{n-k})$ is integral as well.

For brevity set $a_n=1-3n$ and $d_n=1+3n^2$, and introduce the rows
\begin{align}
\mathcal R_{Q,n}
&=\left.\bigl((D^2+a_nD+d_n)Q_n,\ (D+a_n)Q_n,\ Q_n\bigr)
\right|_{u=1},\label{p7:eq:lat-RQ}\\
\mathcal R_{A,n}
&=\left.\bigl((D^2+a_nD+d_n)A_n+2DQ_n+a_nQ_n,\right.\notag\\[-2mm]
&\hspace{39mm}\left.(D+a_n)A_n+Q_n,\ A_n\bigr)
\right|_{u=1}.\label{p7:eq:lat-RA}
\end{align}
These are integer rows.

\begin{proposition}[Exact row formulas]\label{p7:prop:lat-row-formulas}
For every $n\geq0$,
\begin{equation}\label{p7:eq:lat-row-result}
\mathcal R_{Q,n}=\frac{\mathbf q_n}{n!},\qquad
\mathcal R_{A,n}=\frac{\mathbf p_n-3H_n\mathbf q_n}{n!}.
\end{equation}
\end{proposition}

\begin{proof}
We establish the covariance of the rows through exact adjoint solutions.
The adjoint of $L_n$ with respect to $\dd z/z$ is
\[
L_n^*=z\theta^2-(\theta+n)^3.
\]
For any function $w$ define
\begin{equation}\label{p7:eq:lat-concomitant-row}
r_n(w)=\bigl(\theta^2w-(z-3n)\theta w+(z+3n^2)w,
             -\theta w+(z-3n)w,\ w\bigr).
\end{equation}
Direct differentiation shows that
\[
\mathcal C_n(f,w):=r_n(w)\,(f,\theta f,\theta^2f)^T
\]
is constant in $z$ whenever $L_nf=L_n^*w=0$.

The two needed adjoint solutions are
\begin{equation}\label{p7:eq:lat-adjoint-solutions}
w_{Q,n}(z)=Q_n(1/z),\qquad
w_{A,n}(z)=A_n(1/z)-\Log z\,Q_n(1/z).
\end{equation}
Here and below a single branch of $\Log$ is fixed.  To verify these
solutions, let $c_0(\varepsilon)=1$ and
\[
c_{k+1}(\varepsilon)
=\frac{(n-\varepsilon-k)^3}{(\varepsilon+k+1)^2}
 c_k(\varepsilon)\quad(0\leq k<n),\qquad
W(\varepsilon,z)=z^{-\varepsilon}
 \sum_{k=0}^n c_k(\varepsilon)z^{-k}.
\]
Coefficient comparison, including both endpoints, gives
\begin{equation}\label{p7:eq:lat-epsilon-residual}
\bigl((\theta+n)^3-z\theta^2\bigr)W
=-\varepsilon^2z^{1-\varepsilon}
 -\varepsilon^3c_n(\varepsilon)z^{-n-\varepsilon}.
\end{equation}
Thus $W(0,z)$ and $\partial_\varepsilon W(0,z)$ solve the homogeneous
adjoint equation.  Moreover,
\[
c_k(0)=c_{n,k},\qquad
c_k'(0)=c_{n,k}\bigl[3(H_{n-k}-H_n)-2H_k\bigr],
\]
so they are precisely \eqref{p7:eq:lat-adjoint-solutions}.  Evaluating
\eqref{p7:eq:lat-concomitant-row} at $z=1$ gives
\begin{equation}\label{p7:eq:lat-adjoint-rows}
r_n(w_{Q,n})(1)=\mathcal R_{Q,n},\qquad
r_n(w_{A,n})(1)=\mathcal R_{A,n}.
\end{equation}
For example, at $z=1$ the first two derivatives of $w_{A,n}$ are
$-DA_n-Q_n$ and $D^2A_n+2DQ_n$, accounting for every additional term
in \eqref{p7:eq:lat-RA}.

We next compute the concomitants at the regular singular point $0$.
Every solution of $L_nf=0$ has the form
\begin{equation}\label{p7:eq:lat-Frobenius-leading}
f(z)=z^n\left[P(\Log z)+O\bigl(z(1+|\Log z|^2)\bigr)\right],
\qquad P(X)=\frac{\alpha}{2}X^2+\beta X+\chi,
\end{equation}
with arbitrary $\alpha,\beta,\chi$.  To see this directly, write
$f=z^n\sum_{j\geq0}z^jP_j(\Log z)$.  Its coefficient equations are
\[
\partial_X^3P_0=0,\qquad
(j+\partial_X)^3P_j=-(n+j+\partial_X)^2P_{j-1}\quad(j\geq1).
\]
On polynomials of degree at most two the operator on the left is
invertible.  The successive operator norms are $O(1/j)$ for fixed $n$,
which ensures convergence of the coefficient series.  This constructs
three independent solutions and proves \eqref{p7:eq:lat-Frobenius-leading},
also after applying $\theta$ twice.

Since $c_{n,n}=n!$, the adjoint solutions have leading terms
\[
w_{Q,n}=n!z^{-n}+O(z^{1-n}),\qquad
w_{A,n}=n!z^{-n}(-\Log z-5H_n)
 +O\bigl(z^{1-n}(1+|\Log z|)\bigr).
\]
For leading forms $f=z^nP(X)$ and $w=n!z^{-n}V(X)$, the constant
term of the concomitant is
$n!(VP''-V'P'+V''P)$.  All terms with an additional factor $z$ tend
to zero.  With $V=1$ and $V=-X-5H_n$ we obtain the exact constants
\begin{equation}\label{p7:eq:lat-concomitant-values}
\mathcal C_n(f,w_{Q,n})=n!\alpha,\qquad
\mathcal C_n(f,w_{A,n})=n!(\beta-5H_n\alpha).
\end{equation}

Let $T=z(\theta+1)^2$ and put $a=n+1$.  The identity
$L_{n+1}T=TL_n$ follows from $(\theta-n-1)T=T(\theta-n)$.
The leading polynomial of $Tf$ is $(a+\partial_X)^2P$, so
\[
\alpha_{\rm new}=a^2\alpha,\qquad
\beta_{\rm new}=a^2\beta+2a\alpha.
\]
Using $H_{n+1}=H_n+1/a$ in
\eqref{p7:eq:lat-concomitant-values} therefore gives
\begin{align*}
\mathcal C_{n+1}(Tf,w_{Q,n+1})
 &=a^3\mathcal C_n(f,w_{Q,n}),\\
\mathcal C_{n+1}(Tf,w_{A,n+1})
 &=a^3\mathcal C_n(f,w_{A,n})
   -3a^2\mathcal C_n(f,w_{Q,n}).
\end{align*}
At $z=1$ the state of $Tf$ is $M_n^T(f,\theta f,\theta^2f)^T$.
The state of $f$ there is arbitrary.  Hence these equalities, together
with \eqref{p7:eq:lat-adjoint-rows}, are row identities:
\[
\mathcal R_{Q,n+1}M_n^T=a^3\mathcal R_{Q,n},\qquad
\mathcal R_{A,n+1}M_n^T
 =a^3\mathcal R_{A,n}-3a^2\mathcal R_{Q,n}.
\]
Multiplication by $M_n^{-T}=N_n/a^4$ yields
\begin{equation}\label{p7:eq:lat-row-recurrences}
\mathcal R_{Q,n+1}=\frac{\mathcal R_{Q,n}N_n}{n+1},\qquad
\mathcal R_{A,n+1}
 =\frac{\mathcal R_{A,n}N_n}{n+1}
  -\frac{3\mathcal R_{Q,n+1}}{n+1}.
\end{equation}
Finally, $Q_0=1$ and $A_0=0$ give
$\mathcal R_{Q,0}=(1,1,1)$ and
$\mathcal R_{A,0}=(1,1,0)$.  Induction in
\eqref{p7:eq:lat-row-recurrences} proves \eqref{p7:eq:lat-row-result}.
\end{proof}

\subsection{Completion of the lattice calculation and its scope}

\begin{proof}[Proof of Theorem~\ref{p7:thm:lat-main}]
Let $B_n$ be the integer matrix with rows
$\mathcal R_{A,n},\mathcal R_{Q,n}$.  Proposition~\ref{p7:prop:lat-row-formulas}
gives the exact factorization
\begin{equation}\label{p7:eq:lat-shear-factorization}
RP_n=
\begin{pmatrix}n!&3n!H_n\\0&n!\end{pmatrix}B_n.
\end{equation}
Thus $\Lambda_n$ is contained in the lattice on the right of
\eqref{p7:eq:lat-exact}.  The latter has covolume $(n!)^2$.
Lemma~\ref{p7:lem:lat-covolume} gives the same covolume for $\Lambda_n$,
so the inclusion is equality.

The triangular matrix in \eqref{p7:eq:lat-shear-factorization} is an exact
lattice basis, but its upper-right entry need not be reduced.  Subtracting
$\lfloor3H_n\rfloor$ times its first column from its second gives
\eqref{p7:eq:lat-HNF}.  Its entries are integers because $n!H_n\in\Z$.
The set of second coordinates is $n!\Z$, whereas the set of first
coordinates is $n!\Z+3n!H_n\Z$.  This proves
\eqref{p7:eq:lat-content}, since
\[
\gcd(n!,3n!H_n)=\frac{n!}{\den(3H_n)}.
\]
In particular $D_n\mid n!$.  Dividing the basis in
\eqref{p7:eq:lat-shear-factorization} by $n!/D_n$ proves
\eqref{p7:eq:lat-normalized} and its index assertion.
\end{proof}

\begin{corollary}[Unrestricted input ratios]\label{p7:cor:lat-all-rationals}
For each fixed $n\geq0$,
\[
\left\{\frac{\mathbf p_n\mathbf v}{\mathbf q_n\mathbf v}:
 \mathbf v\in\Z^3,\ \mathbf q_n\mathbf v\ne0\right\}=\Q.
\]
\end{corollary}

\begin{proof}
By \eqref{p7:eq:lat-exact}, the ratios are precisely
$3H_n+a/b$ with $a,b\in\Z$ and $b\ne0$.  Translation by the rational
number $3H_n$ maps $\Q$ onto itself.
\end{proof}

\begin{remark}[What adaptive inputs do not inherit]\label{p7:rem:lat-scope}
For an arbitrary input, the exactly normalized linear form has the shape
\[
(D_na+C_nb)-\gamma D_nb
=D_n\bigl[a-b(\gamma-3H_n)\bigr].
\]
The lattice description alone gives neither a bound on the norm of an
input vector realizing $(a,b)$ nor a bound on this linear form.  In
particular, replacing a fixed input by a depth-dependent vector does not
inherit an error estimate proved for the original trajectory solution.
Since all rational ratios are already available at each depth, any claimed
improvement from adaptive inputs must include such additional height and
error estimates, together with a nonvanishing argument.

This limitation is compatible with both arithmetic possibilities for
$\gamma$.  If $\gamma=r/s\in\Q$, every nonzero normalized form lies
in $s^{-1}\Z$ and therefore has absolute value at least $1/s$, while
the lattice also contains nonzero pairs whose form is exactly zero.
Consequently, a geometry-of-numbers argument that allows the small form
to vanish cannot by itself prove irrationality.  The exact lattice
theorem establishes its arithmetic structure without asserting a new
irrationality bound.
\end{remark}


\begin{corollary}[Growth of the normalized lattice]
\label{p7:cor:lat-denominator-growth}
Let $D_n=\den(3H_n)$ and $h_n=n!/D_n$. Then, as $n\to\infty$,
\[
 \log D_n=n+o(n),\qquad
 \log h_n=\log(n!)-n+o(n).
\]
Consequently, the lattice obtained from $\Lambda_n$ by dividing by its
common coordinate gcd has covolume
$D_n^2=\exp(2n+o(n))$.
\end{corollary}
\begin{proof}
Set $K=\lfloor\sqrt{\log n}\rfloor$ and take a prime
$n/K<p\le n$. For large $n$, one has $p^2>n$ and $p>K$.
Put $k=\lfloor n/p\rfloor$, so $1\le k<K$. Separating the multiples of
$p$ gives, in the localization of $\Z$ at $p$,
\[
 pH_n
 =H_k+p\sum_{\substack{1\le j\le n\\p\nmid j}}\frac1j
 \equiv H_k\pmod p.
\]
The positive integer $k!H_k$ satisfies $0<k!H_k\le K!H_K<p$
for all large $n$, because
\[
 \log(K!H_K)\le K\log K+\log(1+\log K)=o(\log n),
 \qquad \log p\ge\log n-\log K.
\]
Since $p\nmid k!$, $H_k$ is a unit at $p$. Thus $v_p(H_n)=-1$.
These primes eventually exceed $3$, so every one of them divides $D_n$.
Writing $\vartheta(x)=\sum_{p\le x}\log p$, and using
$D_n\mid\lcm(1,\ldots,n)$, we obtain
\[
 \vartheta(n)-\vartheta(n/K)
 \le\log D_n\le\log\lcm(1,\ldots,n).
\]
The prime number theorem gives $n+o(n)$ for both bounds, since
$K\to\infty$. The other assertions follow from $h_n=n!/D_n$ and
Theorem~\ref{p7:thm:lat-main}.
\end{proof}

This is a covolume asymptotic for the full normalized image. It does not
bound the size of an individual fixed-trajectory output or its error.



\section{The integer Gamma kernel: its quotient and two positivity barriers}
\label{p7:sec:kernel}

Let $H_0=0$ and $H_k=\sum_{j=1}^k j^{-1}$ for $k\geq1$.
For a real polynomial $P(t)=\sum_{k=0}^d p_k t^k$, put
\begin{align}
 B(P)&=\int_0^\infty P(t)e^{-t}\dd t
       =\sum_{k=0}^d p_k k!,\label{p7:eq:kernel-B}\\
 A(P)&=\sum_{k=0}^d p_k k!H_k,\qquad
 L(P)=\int_0^\infty P(t)e^{-t}\log t\dd t
       =A(P)-\gamma B(P).\label{p7:eq:kernel-AL}
\end{align}
The last identity follows from
$\Gamma'(k+1)=k!(H_k-\gamma)$.
In particular, $A(P),B(P)\in\Z$ when $P\in\Z[t]$.
This integrality concerns the two coefficients separately and makes no
assumption about the arithmetic nature of $\gamma$.

\subsection{An exact quotient over the integers}

Define the $\Z$-linear differential operator
\[
 T=t\frac{\dd}{\dd t}+1-t.
\]
Its image is understood as an additive subgroup of $\Z[t]$.

\begin{lemma}\label{p7:lem:kernel-T-identities}
For every $P\in\Z[t]$,
\[
 B(TP)=0,\qquad A(TP)=-B(P),\qquad
 A(T^2P)=B(T^2P)=0.
\]
\end{lemma}

\begin{proof}
The identity
\[
 e^{-t}TP(t)=\frac{\dd}{\dd t}\bigl(te^{-t}P(t)\bigr)
\]
gives $B(TP)=0$.  Integration by parts also gives
\[
 L(TP)=-\int_0^\infty P(t)e^{-t}\dd t=-B(P).
\]
The logarithmic boundary term vanishes at zero because
$tP(t)\log t\to0$, and at infinity because of the exponential factor.
Equation~\eqref{p7:eq:kernel-AL} now gives $A(TP)=-B(P)$.
Applying these two identities once more proves the assertion for $T^2P$.
\end{proof}

\begin{proposition}[Integer normal form]\label{p7:prop:kernel-normal-form}
Every $P\in\Z[t]$ has a unique representation
\begin{equation}\label{p7:eq:kernel-normal-form}
 P=(B(P)-A(P))+A(P)t+T^2R,\qquad R\in\Z[t].
\end{equation}
Consequently,
\[
 \ker\bigl((A,B):\Z[t]\longrightarrow\Z^2\bigr)=T^2\Z[t],
 \qquad
 \Z[t]/T^2\Z[t]\simeq\Z^2
\]
as abelian groups.
\end{proposition}

\begin{proof}
For $k\geq0$ a direct calculation gives
\begin{align*}
 T(t^k)&=(k+1)t^k-t^{k+1},\\
 T^2(t^k)&=t^{k+2}-(2k+3)t^{k+1}+(k+1)^2t^k.
\end{align*}
The second polynomial is monic.  Successively subtracting integer
multiples of these polynomials reduces any $P\in\Z[t]$ to $c+dt$ without
introducing denominators.  The reduction is unique: if $R\ne0$, then
$\deg(T^2R)=\deg R+2$, so $T^2R$ cannot have degree at most one.
Lemma~\ref{p7:lem:kernel-T-identities} and
$A(c+dt)=d$, $B(c+dt)=c+d$ identify $d=A(P)$ and $c=B(P)-A(P)$.
These observations also prove the kernel assertion.  Every pair
$(a,b)\in\Z^2$ is represented by $(b-a)+at$, proving the stated
isomorphism.
\end{proof}

\begin{remark}\label{p7:rem:kernel-not-scalar}
The quotient above is an additive quotient, not a quotient ring.
For example, $T^2(1)=t^2-3t+1$ belongs to the pair kernel, whereas
$tT^2(1)$ has $(A,B)=(3,1)$ and does not.
It is equally essential to distinguish this pair kernel from
$\ker L$.  If $\gamma=a/b\in\Q$ with $b>0$, then
$(b-a)+at$ lies in $\ker L$ but has nonzero pair $(A,B)=(a,b)$.
Thus identifying the two kernels would already assume the irrationality
that the construction is intended to investigate.
\end{remark}

\subsection{A sharp lower bound under pointwise positivity}

We first record a real-polynomial estimate.  For $\alpha\in\{0,1\}$,
write
\[
 \|R\|_\alpha^2=
 \int_0^\infty R(t)^2t^\alpha e^{-t}\dd t.
\]

\begin{lemma}\label{p7:lem:kernel-laguerre-estimate}
If $R\in\mathbb R[t]$ has degree $m$ and leading coefficient $a$, then
\begin{align}
 \|R\|_\alpha^2&\geq a^2m!\Gamma(m+\alpha+1),\label{p7:eq:kernel-monic-norm}\\
 \int_0^\infty tR(t)^2t^\alpha e^{-t}\dd t
 &\leq(4m+2\alpha+2)\|R\|_\alpha^2,\label{p7:eq:kernel-multiplication-bound}\\
 \int_0^\infty\frac{R(t)^2}{1+t}t^\alpha e^{-t}\dd t
 &\geq\frac{a^2m!\Gamma(m+\alpha+1)}{4m+2\alpha+3}.
 \label{p7:eq:kernel-rational-norm}
\end{align}
\end{lemma}

\begin{proof}
The Rodrigues formula
\[
 L_j^{(\alpha)}(t)=
 \frac{t^{-\alpha}e^t}{j!}
 \frac{\dd^j}{\dd t^j}\bigl(e^{-t}t^{j+\alpha}\bigr)
\]
gives leading coefficient $(-1)^j/j!$ and, by integration by parts,
\[
 \int_0^\infty L_j^{(\alpha)}(t)L_k^{(\alpha)}(t)
                t^\alpha e^{-t}\dd t
 =\frac{\Gamma(j+\alpha+1)}{j!}\,\delta_{jk}.
\]
The monic polynomial $\ell_m=(-1)^m m!L_m^{(\alpha)}$ therefore has
norm squared $m!\Gamma(m+\alpha+1)$.  Writing
$R=a\ell_m+R_0$ with $\deg R_0<m$ and using orthogonality proves
\eqref{p7:eq:kernel-monic-norm}.

Choose the orthonormal Laguerre polynomials with positive leading
coefficients and denote them by $p_j$.  Their recurrence is
\[
 tp_j=\sqrt{(j+1)(j+\alpha+1)}\,p_{j+1}
      +(2j+\alpha+1)p_j+\sqrt{j(j+\alpha)}\,p_{j-1},
\]
where $p_{-1}=0$.
Expand $R=\sum_{j=0}^m c_jp_j$ and apply
$2|c_jc_{j+1}|\leq c_j^2+c_{j+1}^2$ to the off-diagonal terms.
The bounds
\[
 \sqrt{j(j+\alpha)}\leq j+\frac\alpha2,
 \qquad
 \sqrt{(j+1)(j+\alpha+1)}\leq j+1+\frac\alpha2
\]
show that the coefficient of each $c_j^2$ is at most
$4j+2\alpha+2\leq4m+2\alpha+2$.
This proves~\eqref{p7:eq:kernel-multiplication-bound}.
Finally, Cauchy--Schwarz gives
\[
 \int_0^\infty\frac{R^2}{1+t}t^\alpha e^{-t}\dd t
 \geq
 \frac{\|R\|_\alpha^4}
 {\displaystyle\int_0^\infty(1+t)R^2t^\alpha e^{-t}\dd t}
 \geq\frac{\|R\|_\alpha^2}{4m+2\alpha+3}.
\]
Combine this with~\eqref{p7:eq:kernel-monic-norm}.
\end{proof}

\begin{theorem}[Pointwise positivity barrier]\label{p7:thm:kernel-positive-gap}
Suppose that $0\ne P\in\Z[t]$ and
\[
 P(t)\log t\geq0\qquad(t>0).
\]
Then $L(P)\geq1$, and equality holds if and only if $P=t-1$.
More generally, if $P(t)\log t$ has one fixed sign on $(0,\infty)$,
then $|L(P)|\geq1$, with equality only for $P=\pm(t-1)$.
\end{theorem}

\begin{proof}
Continuity and the change of sign of $\log t$ force $P(1)=0$.
Division by the monic polynomial $t-1$ gives
$P=(t-1)Q$, with $Q\in\Z[t]$ nonzero and nonnegative on
$[0,\infty)$.  Set
\[
 w(t)=(t-1)\log t\,e^{-t},\qquad L(P)=\int_0^\infty w(t)Q(t)\dd t.
\]

Every real polynomial nonnegative on this half-line has a representation
\begin{equation}\label{p7:eq:kernel-half-line-sos}
 Q=A^2+tB^2,\qquad A,B\in\mathbb R[t].
\end{equation}
For completeness, factor $Q$ over $\mathbb R$.  Positive real roots have
even multiplicity; their factors are squares.  A negative-root factor
$t+c$, $c>0$, is $(\sqrt c)^2+t$, and $t$ itself has the required form.
An irreducible monic quadratic $t^2+bt+c$ has $c>0$ and
$b+2\sqrt c>0$, and can be written as
\[
 (t-\sqrt c)^2+t(b+2\sqrt c).
\]
The identity
\[
 (A^2+tB^2)(C^2+tD^2)=(AC-tBD)^2+t(AD+BC)^2
\]
completes the factorization argument, including the positive leading
constant.  There is no cancellation between the leading terms of the
two summands in~\eqref{p7:eq:kernel-half-line-sos}.  Thus if
$\deg Q=2d$ and $a=\lc(Q)$, then $\deg A=d$ and $\lc(A)^2=a$;
if $\deg Q=2d+1$, then $\deg B=d$ and $\lc(B)^2=a$.
In both cases $a$ is an integer at least one.

The elementary inequality
\begin{equation}\label{p7:eq:kernel-log-bound}
 (t-1)\log t\geq\frac{2(t-1)^2}{1+t}\qquad(t>0)
\end{equation}
follows because $f(t)=\log t-2(t-1)/(t+1)$ satisfies
$f(1)=0$ and $f'(t)=(t-1)^2/[t(t+1)^2]\geq0$.
Apply Lemma~\ref{p7:lem:kernel-laguerre-estimate} to $(t-1)A$ with
$\alpha=0$, and to $(t-1)B$ with $\alpha=1$.
When the relevant polynomial has degree $d$, the resulting bounds are
\begin{align}
 \int_0^\infty wA^2\dd t
 &\geq\frac{2\lc(A)^2((d+1)!)^2}{4d+7},\label{p7:eq:kernel-even-bound}\\
 \int_0^\infty wtB^2\dd t
 &\geq\frac{2\lc(B)^2(d+1)!(d+2)!}{4d+9}.
 \label{p7:eq:kernel-odd-bound}
\end{align}
The right side of~\eqref{p7:eq:kernel-even-bound}, with leading coefficient
squared at least one, exceeds one for $d\geq2$: its value at $d=2$ is
$72/15$, and it increases thereafter.
The corresponding right side of~\eqref{p7:eq:kernel-odd-bound} exceeds one
for $d\geq1$, starting at $24/13$ and increasing.
These estimates settle every case with $\deg Q\geq3$ and leave only
degrees zero, one and two.

The moments needed for these remaining cases are
\begin{equation}\label{p7:eq:kernel-three-moments}
 \int_0^\infty w\dd t=1,\qquad
 \int_0^\infty tw\dd t=2-\gamma,\qquad
 \int_0^\infty t^2w\dd t=8-4\gamma.
\end{equation}
They follow directly from~\eqref{p7:eq:kernel-AL}.
The elementary harmonic-sum bounds give $0<\gamma<1$.
If $Q=c$ has degree zero, then $c\in\N$ and $L(P)=c$.
If $Q=at+b$ has degree one, its nonnegativity gives integers
$a\geq1$, $b\geq0$, whence $L(P)=a(2-\gamma)+b>1$.
Finally, in degree two write $A=ut+v$ in
\eqref{p7:eq:kernel-half-line-sos}; here $u^2=a\geq1$.
Equation~\eqref{p7:eq:kernel-three-moments} gives
\[
 \int_0^\infty w(ut+v)^2\dd t
 =\bigl(v+u(2-\gamma)\bigr)^2+u^2(4-\gamma^2)>3a.
\]
The other summand $\int wtB^2$ is nonnegative.
Every nonconstant $Q$ therefore gives a strict inequality, and the
constant case has equality precisely when $Q=1$.
Applying the result to $-P$ proves the assertion for the opposite sign.
\end{proof}

\begin{remark}\label{p7:rem:kernel-positive-scope}
The hypothesis concerns the sign of the original integrand, not merely
the sign of its integral.  For example, $P=t$ gives
$L(P)=1-\gamma\in(0,1)$ while $P(t)\log t$ changes sign.
The theorem excludes a sequence of small integer forms certified by
pointwise positivity in this particular polynomial Gamma kernel.
It does not exclude positivity obtained after a transformation that
uses cancellation, or positivity in a different integral representation.
\end{remark}

\subsection{A uniform obstruction to absolute-integral bounds}

\begin{lemma}[A monic $L^1$ estimate]\label{p7:lem:kernel-chebyshev-L1}
If a real polynomial $q$ has degree $d$ and leading coefficient $a$,
then
\[
 \int_{-1}^1|q(x)|\dd x\geq |a|2^{1-d}.
\]
Consequently, every real polynomial $P$ satisfies
$\int_2^6|P(t)|\dd t\geq4|\lc(P)|$.
\end{lemma}

\begin{proof}
Let $U_d$ be the Chebyshev polynomial determined by
$U_d(\cos\theta)=\sin((d+1)\theta)/\sin\theta$; its leading
coefficient is $2^d$.  Put $s(x)=\operatorname{sign}U_d(x)$,
with arbitrary values at its zeros.
This bounded function is orthogonal, for Lebesgue measure on $[-1,1]$,
to every polynomial of degree less than $d$.
Here is a finite verification.  Set $N=d+1$.
After $x=\cos\theta$, each such polynomial multiplied by
$\sin\theta$ is a linear combination of $\sin(k\theta)$ with
$1\leq k<N$.  Integrating over the $N$ intervals on which
$\operatorname{sign}\sin(N\theta)$ is constant gives
\[
 \int_0^\pi\sin(k\theta)\operatorname{sign}\sin(N\theta)\dd\theta
 =\frac1k\operatorname{Re}
   \frac{(1+r)(1-r^N)}{1-r}=0,
 \qquad r=-e^{ik\pi/N}.
\]
Indeed, $r\ne1$, $|r|=1$, and $r^N=(-1)^{N+k}$ is real, whereas
$(1+r)/(1-r)$ is purely imaginary.

Write $q=a2^{-d}U_d+q_0$, where $\deg q_0<d$.
The orthogonality just proved and
\[
 \int_{-1}^1|U_d(x)|\dd x
 =\int_0^\pi|\sin(N\theta)|\dd\theta=2
\]
give $\int_{-1}^1 q(x)s(x)\dd x=a2^{1-d}$.
Since $|s|\leq1$, taking absolute values proves the estimate.
Apply it to $q(x)=P(4+2x)$ and use $\dd t=2\dd x$ to obtain
the assertion on $[2,6]$.
\end{proof}

\begin{theorem}[Absolute-integral barrier]\label{p7:thm:kernel-absolute-gap}
For every nonzero $P\in\Z[t]$,
\begin{equation}\label{p7:eq:kernel-absolute-gap}
 \int_0^\infty|P(t)\log t|e^{-t}\dd t
 \geq c_0,\qquad
 c_0=4\log6\,e^{-6}=0.0177653107376464\ldots.
\end{equation}
More generally, if $0\ne P\in\Q[t]$ and $D\in\N$ satisfies
$DP\in\Z[t]$, then
\[
 D\int_0^\infty|P(t)\log t|e^{-t}\dd t\geq c_0.
\]
\end{theorem}

\begin{proof}
On $[2,6]$ the function $\log t\,e^{-t}$ is decreasing and is at
least $\log6\,e^{-6}$.  Restricting the integral to this interval and
applying Lemma~\ref{p7:lem:kernel-chebyshev-L1} gives
\[
 \int_0^\infty|P(t)\log t|e^{-t}\dd t
 \geq \log6\,e^{-6}\int_2^6|P(t)|\dd t
 \geq4\log6\,e^{-6}|\lc(P)|.
\]
For a nonzero integer polynomial $|\lc(P)|\geq1$.
Applying this conclusion to $DP$ proves the rational-coefficient case.
\end{proof}

\begin{remark}[Exact scope of the obstruction]\label{p7:rem:kernel-absolute-scope}
Suppose that a proposed construction clears the coefficients of $P_n$
with $D_nP_n\in\Z[t]$ and then estimates its form solely by
\[
 D_n|L(P_n)|\leq
 D_n\int_0^\infty|P_n(t)\log t|e^{-t}\dd t.
\]
The right-hand side cannot tend to zero, regardless of the degrees or
coefficients chosen.  This remains true after adding any integer
null-kernel term $T^2R_n$ to an integer polynomial.
The conclusion concerns this absolute-integral estimate; it is not a
lower bound for the signed form $|L(P_n)|$.
Nor does it automatically apply when the actual denominators of
$A(P_n),B(P_n)$ are smaller than the denominators needed to clear the
coefficients of $P_n$.  A procedure using that arithmetic discrepancy,
or using cancellation after integration by parts, requires a separate
analysis.  Together with Theorem~\ref{p7:thm:kernel-positive-gap}, the result
identifies two precise classes of kernel estimates that cannot supply
the required small nonzero integer forms.
\end{remark}


\section{Exact arithmetic of Dickman moments}
\label{p7:sec:dickman-arithmetic}
\label{p7:sec:dickman}

The Dickman construction leads naturally to rational moments and to the
normalization $e^{-\gamma}$.  Here we derive exact local congruences and
common-denominator estimates for those moments.  The results also quantify
the effects of fixed polynomial transformations and of an integer convolution
parameter.  These are arithmetic consequences of the classical moment
formula, without a claim of historical priority.

Let $D$ have the normalized Dickman--de Bruijn distribution, with density
$e^{-\gamma}\rho(x)$ on $[0,\infty)$, where $\rho(x)=1$ for $0\le x\le1$.
Write $m_n=\mathbb E[D^n]$.  The classical generating function and
set-partition formula are
\begin{align}
 M(z)&=\sum_{n\ge0}m_n\frac{z^n}{n!}
       =\exp\left(\sum_{k\ge1}\frac{z^k}{k\,k!}\right),
       \label{p7:dick:eq:mgf}\\
 m_n&=\sum_{\pi\in\Pi_n}\prod_{B\in\pi}\frac1{|B|},
       \label{p7:dick:eq:partitions}
\end{align}
where $\Pi_n$ is the set of partitions of $\{1,\ldots,n\}$ and the empty
partition gives $m_0=1$; see \cite[pp.~257--258]{p7:CellarosiSinai2011}.
In particular,
\[
 (m_0,\ldots,m_6)
 =\left(1,1,\frac32,\frac{17}6,\frac{19}3,\frac{81}5,
                \frac{8351}{180}\right).
\]
Differentiating \eqref{p7:dick:eq:mgf} gives the rational recurrence
\begin{equation}\label{p7:dick:eq:recurrence}
 n m_n=\sum_{j=0}^{n-1}\binom nj m_j\qquad(n\ge1).
\end{equation}
For a prime $p$, let $v_p(p)=1$ and
$\Z_{(p)}=\{a/b\in\Q:p\nmid b\}$.  We use $v_p(0)=+\infty$ and
$\den(0)=1$.

\subsection{An exact congruence at every odd prime}

\begin{lemma}\label{p7:dick:lem:den-bound}
For every prime $p$ and every $n\ge0$,
\[
 v_p(m_n)\ge-\lfloor n/p\rfloor.
\]
\end{lemma}
\begin{proof}
For every positive integer $s$, $v_p(s)\le s/p$: if $v_p(s)=a\ge1$,
then $s\ge p^a\ge pa$.  Thus each term of
\eqref{p7:dick:eq:partitions}, with block sizes $s_1,\ldots,s_t$, has
denominator valuation at most
$\sum_i v_p(s_i)\le\lfloor n/p\rfloor$.  Taking a sum preserves this bound.
\end{proof}

\begin{theorem}[Odd-prime moment congruence]\label{p7:dick:thm:congruence}
Let $p$ be an odd prime, and write $n=qp+r$, where $q\ge0$ and
$0\le r<p$.  Then
\begin{equation}\label{p7:dick:eq:congruence}
 p^q m_n\equiv m_r\pmod{p\Z_{(p)}}.
\end{equation}
Consequently,
\begin{equation}\label{p7:dick:eq:exact-valuation-test}
 v_p(m_n)=-q
 \quad\Longleftrightarrow\quad
 p\nmid\operatorname{num}(m_r),
\end{equation}
where the numerator is taken in lowest terms.
\end{theorem}
\begin{proof}
For a partition with block sizes $s_1,\ldots,s_t$, put
$v=\sum_i v_p(s_i)\le q$.  If $v<q$, its contribution to $p^q m_n$
vanishes modulo $p\Z_{(p)}$.  If $v=q$, then
\begin{equation}\label{p7:dick:eq:deficit}
 \sum_i\bigl(s_i-pv_p(s_i)\bigr)=r<p.
\end{equation}
Each summand is nonnegative.  Moreover, if $p\mid s$ and $s\ne p$, then
$s-pv_p(s)\ge p$.  Indeed, for $v_p(s)=1$ one has $s\ge2p$; for
$v_p(s)=a\ge2$ one has
$s-pa\ge p(p^{a-1}-a)\ge p$, since $p\ge3$.
Equation~\eqref{p7:dick:eq:deficit} therefore forces all $p$-divisible block
sizes to equal $p$.  There are exactly $q$ such blocks, and the remaining
blocks partition a set of size $r$.

The total contribution of these partitions is
\[
 p^{-q}C_{n,p}m_r,
 \qquad C_{n,p}=\frac{n!}{r!\,q!\,(p!)^q}.
\]
The coefficient counts the choice of the remaining $r$ elements and the
partition of the other elements into $q$ unordered blocks of size $p$.
To compute its residue one must account for the possible divisibility of
$q!$ by $p$.  Separating the multiples of $p$ in $n!$ gives the exact identity
\begin{equation}\label{p7:dick:eq:unit-factor}
 C_{n,p}=
 \frac{\displaystyle\prod_{\substack{1\le j\le n\\p\nmid j}}j}
      {r!\,((p-1)!)^q}.
\end{equation}
The numerator and denominator on the right are units in $\Z_{(p)}$.
The $q$ complete blocks of nonmultiples of $p$ each have residue product
$(p-1)!$, and the final incomplete block has residue product $r!$.
Hence $C_{n,p}\equiv1\pmod p$.  All other partitions vanish after
multiplication by $p^q$ and reduction, proving
\eqref{p7:dick:eq:congruence}.  Since $m_r\in\Z_{(p)}$ for $r<p$, the
valuation criterion follows.
\end{proof}

\begin{remark}\label{p7:dick:rem:exceptions}
The exceptional residue classes have an exact finite description for each
$p$: they are the $r\in\{0,\ldots,p-1\}$ with
$p\mid\operatorname{num}(m_r)$.  For example, the sole exceptional classes
for $p=3$ and $p=7$ are respectively $r=2$ and $r=6$.
Theorem~\ref{p7:dick:thm:congruence} does not determine the higher valuation
within an exceptional class.  Its restriction to odd primes is essential
to the proof: at $p=2$, both block sizes $2$ and $4$ satisfy
$s-2v_2(s)=0$, allowing additional cancellations.
\end{remark}

\subsection{Common denominators and fixed polynomial transformations}

\begin{corollary}[Exact odd common denominator]\label{p7:dick:cor:common-den}
Let $L_N=\lcm(\den(m_0),\ldots,\den(m_N))$, and let $L_N^{\mathrm{odd}}$
denote its odd part.  Then
\begin{equation}\label{p7:dick:eq:odd-den}
 L_N^{\mathrm{odd}}=
 \prod_{\substack{3\le p\le N\\p\ \mathrm{prime}}}p^{\lfloor N/p\rfloor},
\end{equation}
and
\begin{equation}\label{p7:dick:eq:den-growth}
 \log L_N^{\mathrm{odd}}=N\log N+O(N),
 \qquad \log L_N=N\log N+O(N).
\end{equation}
\end{corollary}
\begin{proof}
Lemma~\ref{p7:dick:lem:den-bound} bounds the exponent of $p$ in $L_N$ by
$\lfloor N/p\rfloor$.  For odd $p\le N$, take
$n=p\lfloor N/p\rfloor$.  The remainder is zero, so
Theorem~\ref{p7:dick:thm:congruence}, with $m_0=1$, attains that exponent.
This proves \eqref{p7:dick:eq:odd-den}.

The standard Mertens and Chebyshev estimates
\[
 \sum_{p\le x}\frac{\log p}{p}=\log x+O(1),
 \qquad \sum_{p\le x}\log p=O(x)
\]
give
\[
 \sum_{3\le p\le N}\lfloor N/p\rfloor\log p
 =N\sum_{3\le p\le N}\frac{\log p}{p}+O(N)
 =N\log N+O(N).
\]
Finally $v_2(L_N)\le\lfloor N/2\rfloor$, so the power of $2$ contributes
only $O(N)$ to the logarithm.
\end{proof}

\begin{theorem}[Polynomial transformation law]\label{p7:dick:thm:polynomial}
Fix a nonconstant $g\in\Q[X]$ of degree $d\ge1$, and put
\[
 u_n=\mathbb E[g(D)^n],\qquad
 L_N(g)=\lcm(\den(u_0),\ldots,\den(u_N)).
\]
Then
\begin{equation}\label{p7:dick:eq:polynomial-growth}
 \log L_N(g)=dN\log N+O_g(N).
\end{equation}
More precisely, write $g=h/A$, where $A\ge1$ is an integer and
$h\in\Z[X]$ has leading coefficient $c\ne0$.  For every odd prime
$p\nmid Ac$ and every $q\ge0$,
\begin{equation}\label{p7:dick:eq:polynomial-local}
 v_p(u_{pq})=-dq.
\end{equation}
\end{theorem}
\begin{proof}
Expand
$h(X)^n=\sum_{k=0}^{dn}b_{n,k}X^k$, with $b_{n,k}\in\Z$ and
$b_{n,dn}=c^n$.  For every prime $p$,
\begin{equation}\label{p7:dick:eq:polynomial-upper}
 v_p(u_n)\ge-nv_p(A)-\lfloor dn/p\rfloor.
\end{equation}
It follows simultaneously for $n\le N$ that
\[
 L_N(g)\mid A^N\prod_{p\le dN}p^{\lfloor dN/p\rfloor}.
\]
The estimates used in Corollary~\ref{p7:dick:cor:common-den} therefore give
$\log L_N(g)\le dN\log N+O_g(N)$.

For the reverse bound, let $p\nmid2Ac$ and put $n=pq$.  The leading term
of $p^{dq}u_n$ has residue
\[
 A^{-n}c^n p^{dq}m_{dpq}\equiv(c/A)^n\not\equiv0\pmod p.
\]
For every lower degree $k\le dpq-1$, one has
$\lfloor k/p\rfloor\le dq-1$.  Thus
$p^{dq}b_{n,k}m_k\in p\Z_{(p)}$ by
Lemma~\ref{p7:dick:lem:den-bound}; multiplication by $A^{-n}$ preserves
this divisibility.  All lower terms vanish modulo $p$, proving
\eqref{p7:dick:eq:polynomial-local}.
Taking $q=\lfloor N/p\rfloor$ now yields
\[
 \log L_N(g)\ge
 d\sum_{\substack{p\le N\\p\nmid2Ac}}\lfloor N/p\rfloor\log p
 =dN\log N+O_g(N).
\]
Omitting the finitely many primes dividing $2Ac$ costs $O_g(N)$.
The lower and upper bounds match.
\end{proof}

\begin{corollary}[Exact local radii]\label{p7:dick:cor:radii}
With the notation of Theorem~\ref{p7:dick:thm:polynomial}, for every odd
prime $p\nmid Ac$ the formal series
\[
 U_g(z)=\sum_{n\ge0}u_nz^n,
 \qquad V_g(z)=\sum_{n\ge0}u_n\frac{z^n}{n!}
\]
have exact $p$-adic convergence radii
\begin{equation}\label{p7:dick:eq:radii}
 R_p(U_g)=p^{-d/p},\qquad
 R_p(V_g)=p^{-d/p-1/(p-1)}.
\end{equation}
In particular, for the classical Dickman moment generating function,
$R_p(M)=p^{-1/p-1/(p-1)}$ at every odd prime.
\end{corollary}
\begin{proof}
Equations~\eqref{p7:dick:eq:polynomial-upper} and
\eqref{p7:dick:eq:polynomial-local} imply
$\liminf_{n\to\infty}v_p(u_n)/n=-d/p$.
The $p$-adic Cauchy--Hadamard formula gives the first radius.
Legendre's formula gives $v_p(n!)/n\to1/(p-1)$; subtracting this limit
gives the second.
\end{proof}

\begin{remark}\label{p7:dick:rem:scope}
The common-denominator law is stronger than a finite observation about
individual moments, but its quantifiers matter.  It does not give a lower
bound for every individual denominator or every selected linear combination
of moments.  Fixed rational translations and rescalings retain the leading
cost $N\log N$; a fixed degree-$d$ polynomial transformation changes it to
$dN\log N$.  Transformations depending on $N$ are outside
Theorem~\ref{p7:dick:thm:polynomial}.  For $d\ge2$, the series $V_g$ need not
converge for positive real $z$; Corollary~\ref{p7:dick:cor:radii} treats it as
a formal rational series and asserts only its $p$-adic radius.
\end{remark}

\subsection{Moving integer shape and the remaining normalization}

For a positive integer $\theta$, let $D_\theta$ be the sum of $\theta$
independent copies of $D$, and write
$m_n(\theta)=\mathbb E[D_\theta^n]$.  Taking the $\theta$th power in
\eqref{p7:dick:eq:mgf} shows that its cumulants are
\begin{equation}\label{p7:dick:eq:shape-cumulants}
 \kappa_n(\theta)=\frac{\theta}{n}\qquad(n\ge1).
\end{equation}

\begin{theorem}[Exact shape cost of integer moments]\label{p7:dick:thm:shape}
For $\theta\in\Z_{>0}$ and $N\ge1$,
\begin{equation}\label{p7:dick:eq:shape-integrality}
 m_0(\theta),\ldots,m_N(\theta)\in\Z
 \quad\Longleftrightarrow\quad
 \lcm(1,\ldots,N)\mid\theta.
\end{equation}
The least positive such shape is therefore $\lcm(1,\ldots,N)$, whose
logarithm is $N+o(N)$.
\end{theorem}
\begin{proof}
The moment--cumulant formula
\[
 m_n=\sum_{\pi\in\Pi_n}\prod_{B\in\pi}\kappa_{|B|}
\]
is a polynomial with integer coefficients.  Thus divisibility by
$\lcm(1,\ldots,N)$ makes all cumulants in
\eqref{p7:dick:eq:shape-cumulants}, and hence all moments through $N$,
integers.  Conversely, the inverse formula
\[
 \kappa_n=\sum_{\pi\in\Pi_n}
   (-1)^{|\pi|-1}(|\pi|-1)!\prod_{B\in\pi}m_{|B|}
\]
has integer coefficients as well.  Integer moments through $N$ force
$\theta/n\in\Z$ for every $1\le n\le N$, which is precisely the required
divisibility.  The logarithmic asymptotic for the least common multiple
is the usual prime number theorem.
\end{proof}

This construction does remove moment denominators, but its connection to
Euler's constant retains an exponential normalization.  Let $f_\theta$
be the density of $D_\theta$.  For $0<x\le1$, every coordinate in its
convolution simplex lies in $[0,1]$, where the original density is
$e^{-\gamma}$.  The simplex volume therefore gives the exact formula
\begin{equation}\label{p7:dick:eq:shape-density}
 f_\theta(x)=\frac{e^{-\gamma\theta}x^{\theta-1}}{(\theta-1)!}.
\end{equation}
For $P(X)=\sum_k a_kX^k\in\Q[X]$, it follows that
\begin{equation}\label{p7:dick:eq:shape-form}
 \int_1^\infty P(x)f_\theta(x)\dd x
 =\mathbb E[P(D_\theta)]
  -\frac{e^{-\gamma\theta}}{(\theta-1)!}
     \sum_k\frac{a_k}{\theta+k}.
\end{equation}
If $P\in\Z[X]$, $\deg P\le N$, and the divisibility condition in
Theorem~\ref{p7:dick:thm:shape} holds, the unrestricted expectation in
\eqref{p7:dick:eq:shape-form} is an integer.  This produces explicit rational
linear forms in $1$ and $e^{-\gamma\theta}$.

\begin{remark}[Arithmetic scope of the construction]\label{p7:dick:rem:gamma-gap}
No estimate making the forms \eqref{p7:dick:eq:shape-form} sufficiently small
has been established here.  Moreover, their direct target is
$e^{-\gamma\theta}$: if $\gamma$ were a nonzero algebraic number, this
target would already be transcendental by Hermite--Lindemann.  Proving only
its irrationality or transcendence would consequently not exclude
algebraic $\gamma$.  Taking logarithms does not preserve the rational
linear-form structure or its denominator bounds.  Thus the exact
congruence, transformation laws, and shape construction resolve specific
arithmetic questions about the moment method while leaving the required
transfer to $\gamma$ unproved.
\end{remark}

For a finite consistency check, recurrence~\eqref{p7:dick:eq:recurrence} was
evaluated with exact rational arithmetic.  Among the $10{,}088$ pairs
$(n,p)$ with $1\le n\le300$ and $3\le p\le n$ prime, exactly $9{,}822$
attained the denominator exponent $\lfloor n/p\rfloor$.  Each of the
remaining $266$ pairs satisfies
$p\mid\operatorname{num}(m_{n\bmod p})$, as predicted by
Theorem~\ref{p7:dick:thm:congruence}.  A further $47{,}094$ congruence checks,
covering all odd primes $p\le500$ and $0\le n\le500$, had no failures.
These computations corroborate the formulas; the proofs above supply
their unrestricted quantifiers.


\section{What the attempted escapes achieved, and where they stop}
\label{p7:sec:assessment}

\subsection{Adaptive lattice compression: a concrete failed combination}

The exact content of the projected trajectory is much larger than one.
This suggested removing all common factors and then using the freedom
of the initial vector to obtain a short output pair. Theorems above
carry out the arithmetic part completely. What fails is an automatic
transfer of the approximation estimate.

There is a quantitative version of this distinction. With
$D=D_n$ and $C=C_n$, the normalized forms are
\[
 \Lambda=D\left[a-b\left(\gamma-\frac CD\right)\right],
 \qquad (a,b)\in\Z^2,\quad \gcd(C,D)=1.
\]
Dirichlet's elementary approximation theorem, with $1\le b\le B$,
gives an output satisfying $|\Lambda|\le D/B$, but allows $\Lambda=0$.
If $\gamma=r/s$ with $s>0$, an exact zero implies
\[
 D\mid sCb,\qquad\text{hence}\qquad D\mid sb.
\]
Thus $0<|b|<D/s$ would certify nonvanishing under that hypothesis.
The two proposed bounds pull in opposite directions: making the
Dirichlet upper bound tend to zero requires $B/D\to\infty$, whereas
this modular nonvanishing certificate requires $B<D/s$.

This calculation discards the combination of \emph{these} two estimates.
It does not rule out better approximation estimates or a different
nonvanishing argument. Nor does it rule out carefully restricted initial
vectors for which an analytic estimate survives. Such a construction
would have to prove the restriction's arithmetic and analytic effects
together.

\subsection{Finite evaluation of a fixed trajectory}

For a reproducible numerical comparison, let $v=(0,0,1)^T$ and reduce
\[
 \frac{\mathbf p_nv}{\mathbf q_nv}=\frac{a_n}{b_n},
 \qquad b_n>0,\quad \gcd(a_n,b_n)=1.
\]
The matrix product and the reduction are exact integer operations.
The following real logarithms were then evaluated with 200-decimal-digit
arithmetic; the displayed values are rounded.
\begin{center}
\small
\begin{tabular}{@{}rrrrr@{}}
\toprule
$n$ & $\log_{10}b_n$ &
$\log_{10}|\gamma-a_n/b_n|$ &
$\log_{10}|b_n\gamma-a_n|$ & $\tau_n$\\
\midrule
10  & 13.2581  & $-7.6011$  & 5.6571   & $-0.4267$\\
20  & 30.9021  & $-11.9023$ & 18.9998  & $-0.6148$\\
40  & 74.0625  & $-20.5465$ & 53.5159  & $-0.7226$\\
80  & 171.6475 & $-33.4311$ & 138.2165 & $-0.8052$\\
120 & 274.2720 & $-43.8412$ & 230.4308 & $-0.8402$\\
200 & 499.1524 & $-62.1680$ & 436.9844 & $-0.8755$\\
300 & 794.2474 & $-82.3402$ & 711.9072 & $-0.8963$\\
\bottomrule
\end{tabular}
\end{center}
Here
\[
 \tau_n=-1-
 \frac{\log|\gamma-a_n/b_n|}{\log b_n}.
\]
The table distinguishes a small rational-approximation error from the
required small integer linear form. For this sample, the latter is
large despite increasingly accurate rational values. These finite data
do not prove an asymptotic lower bound, do not exclude favorable
subsequences, and do not describe every trajectory or every input vector.

The analytic exponent relevant to this rank-three/Rivoal trajectory is
$n^{2/3}$, not the $n^{3/4}$ associated with a different approximation
family. No conclusion here multiplies a mere upper error bound by a
denominator and treats the result as a lower bound.

\subsection{Corrections when reusing the supplied arguments}

Several statements in the supplied \texttt{gamma\_attempts.tex} need a
narrower reading or a small correction before they can be reused in a
proof.

Here $\bG$ is the ring of finite values of Siegel $G$-functions at
algebraic points, allowing analytic continuation; $\bE$ is the analogous
ring of $E$-function values, and $\bD$ denotes the ring of Borel-summed
arithmetic Gevrey values used in \cite{p7:FR2024}.

\begin{enumerate}
\item \textbf{A cancellation template is not a universal obstruction.}
The expectation $\gamma\notin\Frac(\bG)$ is unproved. Moreover,
$\gamma\in\Frac(\bG)$ follows automatically under the reductio assumption
$\gamma\in\Qbar$. A membership prerequisite therefore cannot, by itself,
exclude a construction made only under that assumption.
Theorem~\ref{p7:thm:no-siegel} is unconditional but concerns only its
specified differential module.

\item \textbf{The cited $G$-function theorem needs its hypotheses.}
The ratio theorem in \cite{p7:FR2014} assumes a nonzero constant $\xi$.
At $\xi=0$, the function $A-\xi B=A$ cannot simultaneously have radius
$1$ and radius greater than $R\ge1$. In continuation arguments, choose
algebraic points in open neighborhoods, not necessarily on an arbitrary
given path. Algebraicity of a Wronskian normalization also needs the
normalized Frobenius basis; a rational logarithmic derivative alone
leaves an arbitrary multiplicative constant.

\item \textbf{Twist invariance has a defined scope.}
For $\Gamma_\tau(s)=e^{\tau(1-s)}\Gamma(s)$, the digamma difference
identities can remain invariant. The original Gamma recurrence does not:
\[
 \Gamma_\tau(s+1)=e^{-\tau}s\Gamma_\tau(s).
\]
For real $\tau$, requiring the original normalized recurrence forces
$\tau=0$. The invariance argument therefore restricts arguments using
only the invariant identities, not every normalized Gamma construction.

\item \textbf{Denominator quantifiers cannot be interchanged.}
An upper bound on a common denominator is not a lower bound on an
individual reduced denominator. A finite table does not settle every
degree. The Dickman results here prove an exact odd common denominator
and a local congruence; they deliberately do not assert the same leading
denominator exponent for every individual moment or arbitrary linear
combination. Similarly, the exact content formula concerns $RP_n$.
Equality with the content of the full $3\times3$ matrix $P_n$ is not
proved or used.

\item \textbf{The Gevrey statements are conditional at different levels.}
The conjecture $\bE\cap\bD=\Qbar$ implies the transcendence of $\gamma$.
The stronger exclusion $\gamma\notin\bE$ in the cited work uses an
additional divergent-series zero-division conjecture
\cite{p7:FR2024}. Neither conjecture is proved in this supplement.
\end{enumerate}

\subsection{What remains open after these calculations}

The Meijer identity now provides exact coefficient vectors for
$\gamma+\Log z$; the matrix calculation provides the full projected
lattice and its exact content; the kernel calculation provides all
integer null directions; and the moment calculation provides exact
local denominator costs. None supplies both strict nonvanishing and
decay after the necessary arithmetic normalization.

A further construction would need an ingredient beyond the discarded
estimates: for example, an explicitly restricted adaptive input whose
height and approximation direction can both be proved, or a transformed
integral whose true coefficient pair has much smaller denominators than
its integrand coefficients and whose signed cancellation can be
controlled. These are descriptions of missing mechanisms, not new
sufficient conditions advertised as a proof. The present work makes the
tested routes more explicit and their unresolved arithmetic step
checkable; the irrationality and transcendence of $\gamma$ remain
unproved here.

\PartAppendix
\section{An independent algebraic certificate for the lattice rows}
\label{p7:app:telescoping}

For completeness, the row formulas in
Proposition~\ref{p7:prop:lat-row-formulas} have a second verification using
only rational functions and polynomial algebra. This also permits
checking the harmonic coefficient without relying on the adjoint
calculation.

Put $a=n+1$ and define polynomial rows
\[
 h_n(x)=\bigl(x^2+(1-3n)x+3n^2+1,\ x+1-3n,\ 1\bigr)
\]
and $V_n=(V_{n,1},V_{n,2},V_{n,3})$, where
\begin{align*}
 V_{n,1}(x)
 &=-a^2(10a^2+3a+6)+a(15a^2+3a+8)x
                     -(6a^2+a+3)x^2,\\
 V_{n,2}(x)
 &=3a^2(5a-2)-a(21a-8)x+(8a-3)x^2,\\
 V_{n,3}(x)&=-6a^2+8ax-3x^2.
\end{align*}
Direct polynomial multiplication proves
\begin{equation}
\label{p7:eq:tel-certificate}
 (a-x)^3V_n(x+1)-x^2V_n(x)
 =a^4h_{n+1}(x)-(a-x)^3h_n(x)N_n.
\end{equation}

For a formal parameter $\varepsilon$ set
\[
 W_n(j;\varepsilon)
 =\frac{n!^3}{(1+\varepsilon)_j^2(1-\varepsilon)_{n-j}^3},
 \qquad
 S_n(\varepsilon)
 =\sum_{j=0}^n W_n(j;\varepsilon)h_n(j+\varepsilon),
\]
where $(x)_j=x(x+1)\cdots(x+j-1)$ and $(x)_0=1$.
When $j=n+1,n+2$, use
$(x)_{-r}=1/[(x-r)\cdots(x-1)]$; the resulting $W_n$ are rational
functions vanishing to order three at $\varepsilon=0$.

The exact adjacent-index ratios are
\[
 \frac{W_{n+1}(j;\varepsilon)}{W_n(j;\varepsilon)}
 =\frac{a^3}{(a-j-\varepsilon)^3},\qquad
 \frac{W_n(j+1;\varepsilon)}{W_n(j;\varepsilon)}
 =\frac{(n-j-\varepsilon)^3}{(j+1+\varepsilon)^2}.
\]
Multiply \eqref{p7:eq:tel-certificate} by
$W_n(j;\varepsilon)/(a-j-\varepsilon)^3$ and take
$x=j+\varepsilon$. Its two sides become the exact telescoping identity
\begin{align*}
 &aW_{n+1}(j;\varepsilon)h_{n+1}(j+\varepsilon)
  -W_n(j;\varepsilon)h_n(j+\varepsilon)N_n\\
 &\hspace{30mm}=T_n(j+1;\varepsilon)-T_n(j;\varepsilon),
\end{align*}
with
\[
 T_n(j;\varepsilon)=
 \frac{W_n(j;\varepsilon)(j+\varepsilon)^2
                  V_n(j+\varepsilon)}
      {(a-j-\varepsilon)^3}.
\]
Sum from $j=0$ to $n+1$. The lower boundary is $O(\varepsilon^2)$,
the upper boundary at $n+2$ is $O(\varepsilon^3)$, and the extra
$W_n(n+1;\varepsilon)$ in the second sum is $O(\varepsilon^3)$.
These statements include the apparent interior pole at $j=n+1$,
which is handled as an identity of rational functions before summing.
Therefore
\[
 (n+1)S_{n+1}(\varepsilon)-S_n(\varepsilon)N_n
 =O(\varepsilon^2).
\]
Both constant and first-order coefficients obey the divided matrix
recurrence. At $n=0$ their initial rows are respectively
$(1,1,1)$ and $(1,1,0)$. Since
\[
 W_n(j;0)=\binom nj^3j!,\qquad
 W_n'(j;0)=\binom nj^3j!\,(3H_{n-j}-2H_j),
\]
we obtain
\begin{align*}
 \frac{\mathbf q_n}{n!}
 &=\sum_{j=0}^n\binom nj^3j!\,h_n(j),\\
 \frac{\mathbf p_n}{n!}
 &=\sum_{j=0}^n\binom nj^3j!\,
   \bigl[(3H_{n-j}-2H_j)h_n(j)+h_n'(j)\bigr].
\end{align*}
These are precisely the row formulas already proved by the adjoint
method. Their integral harmonic shear follows term by term as in
Section~\ref{p7:sec:lattice}.



\section{Reproducible finite checks}
\label{p7:app:computation}

The code below checks the exact row identities, integral shear, projected
contents and maximal-minor gcds for all $0\le n\le80$. It also checks
$47{,}094$ Dickman congruences and five exact odd common denominators,
and reproduces the fixed-column table through $n=300$.
The first two routines use only Python's exact integers and rational
fractions. The last two use \texttt{mpmath}; the reported run used
version 1.3.0. The table uses 200 decimal digits.

For the residue tower, 24 independent vertical-line quadratures at
$p=2,3,4$, $z\in\{1/2,1,2\}$ and $0\le k\le\min(2,p-1)$ gave
maximum absolute residual $6.0185311\times10^{-36}$ at 35-digit precision.
The numerical integral is truncated at imaginary height 20. This is
ordinary high-precision corroboration, not a rigorous interval bound
on roundoff or tails. The contour proof establishes the identity.

The assertions compare independent representations rather than infer
universal statements from finite agreement. Neither the quadratures nor
the finite arithmetic checks are used to prove irrationality.

Copy the following listing into a Python file and run it. It is embedded
here so the accompanying \texttt{.tex} file remains self-contained.

\begin{lstlisting}
"""Exact checks and optional numerical checks for the supplement.
Run with Python 3.11+. Only the last two functions require mpmath.
These finite checks corroborate, and do not replace, the proofs.
"""
from fractions import Fraction as F
from math import comb, factorial, gcd, lcm, prod
from itertools import combinations

def N(n):
    return [
        [6*n*n + 13*n + 10, -8*n - 5, 3],
        [3*n**3 + 3*n*n - 3*n - 6, -3*n*n + 3*n + 3, n-2],
        [n**4 - 2*n**3 - 6*n*n - 3*n + 3,
         6*n*n + 2*n - 1, 1-2*n],
    ]

def multiply(A, B):
    return [[sum(x*y for x, y in zip(row, col))
             for col in zip(*B)] for row in A]

def harmonic(n):
    return sum((F(1, j) for j in range(1, n+1)), F(0))

def row_checks(limit=80):
    P = [[1,0,0], [0,1,0], [0,0,1]]
    R = [[1,1,0], [1,1,1]]
    H = [harmonic(j) for j in range(limit+1)]
    for n in range(limit+1):
        p, q = multiply(R, P)
        f = factorial(n)
        qsum, psum = [F(0)]*3, [F(0)]*3
        for k in range(n+1):
            w = comb(n,k)**3 * factorial(k)
            h = [k*k+(1-3*n)*k+3*n*n+1, k+1-3*n, 1]
            dh = [2*k+1-3*n, 1, 0]
            for j in range(3):
                qsum[j] += w*h[j]
                psum[j] += w*((3*H[n-k]-2*H[k])*h[j]+dh[j])
        assert [f*x for x in qsum] == q
        assert [f*x for x in psum] == p
        shifted = [F(x,f)-3*H[n]*F(y,f) for x,y in zip(p,q)]
        assert all(x.denominator == 1 for x in shifted)
        minors = [p[i]*q[j]-p[j]*q[i]
                  for i,j in combinations(range(3),2)]
        assert gcd(*minors) == f*f
        D = (3*H[n]).denominator
        assert gcd(*p, *q) == f//D
        assert gcd(*q) == f and gcd(*p) == f//D
        P = multiply(P, N(n))
    print("Exact lattice checks:", limit+1, "depths")

def moment_checks(limit=500):
    m = [F(1)]
    for n in range(1, limit+1):
        m.append(sum((comb(n,j)*m[j] for j in range(n)), F(0))/n)
    primes = [p for p in range(3, limit+1)
              if all(p%d for d in range(2, int(p**0.5)+1))]
    total = maximal = pairs300 = exceptions = 0
    for p in primes:
        for n in range(limit+1):
            q, r = divmod(n, p)
            delta = p**q*m[n]-m[r]
            assert delta.denominator % p
            assert delta.numerator % p == 0
            total += 1
            if 1 <= n <= 300 and p <= n:
                pairs300 += 1
                is_maximal = (p**q*m[n]).numerator % p != 0
                assert is_maximal == (m[r].numerator % p != 0)
                maximal += is_maximal
                exceptions += not is_maximal
    for depth in (10, 40, 100, 300, 500):
        if depth > limit:
            continue
        L = lcm(*(x.denominator for x in m[:depth+1]))
        while L % 2 == 0:
            L //= 2
        assert L == prod(p**(depth//p) for p in primes if p <= depth)
    print("Dickman congruences:", total)
    print("Pairs through 300, maximal, exceptional:",
          pairs300, maximal, exceptions)

def fixed_column_table(limit=300):
    import mpmath as mp
    mp.mp.dps = 200
    P = [[1,0,0], [0,1,0], [0,0,1]]
    selected = {10,20,40,80,120,200,300}
    for n in range(1, limit+1):
        P = multiply(P, N(n-1))
        if n not in selected:
            continue
        p = P[0][2]+P[1][2]
        q = p+P[2][2]
        a, b = F(p,q).as_integer_ratio()
        err = abs(mp.mpf(a)/b-mp.euler)
        height, logerr = mp.log10(b), mp.log10(err)
        print(n, *(mp.nstr(x,12) for x in
                   (height,logerr,height+logerr,-1-logerr/height)))

def mellin_checks():
    import mpmath as mp
    with mp.workdps(35):
        residuals = []
        for p in (2,3,4):
            for z in (mp.mpf("0.5"),mp.mpf(1),mp.mpf(2)):
                X = mp.euler+mp.log(z)
                rhs = [mp.mpf(1), X,
                       (X*X+(2*p+1)*mp.zeta(2))/2]
                for k in range(min(p,3)):
                    def integrand(y):
                        s = mp.mpc("-0.5",y)
                        phi = mp.gamma(-s)**(p+1)*mp.gamma(1+s)**p
                        op = z*(1+s)**(p-1-k)+(-1)**p*s**(p-k)
                        return mp.re(phi*mp.power(z,s)*op)/mp.pi
                    value = mp.quad(integrand,[0,1,3,8,20])
                    residuals.append(abs(value-rhs[k]))
        print("Mellin quadratures:",len(residuals),
              "maximum residual",mp.nstr(max(residuals),8))
        assert max(residuals) < mp.mpf("1e-28")

if __name__ == "__main__":
    row_checks()
    moment_checks()
    fixed_column_table()
    mellin_checks()


\end{lstlisting}

\clearpage
\begin{thebibliography}{VIII.99}

\bibitem[VII.1]{p7:Kalisch2026}
R. Kalisch, T. Monderer, E. Leibtag, S. Weinbaum, M. Shalyt, and
I. Kaminer,
\emph{An Infinite Family of Ap\'ery-like Recurrences for Euler's Constant},
arXiv:2610.13919v1, 9 October 2026.
\url{https://arxiv.org/abs/2610.13919}.
The trajectory is taken from Section 8.4; Section 9.2(1) motivates the
arbitrary-$z$ identity extension.

\bibitem[VII.2]{p7:FR2014}
S. Fischler and T. Rivoal,
\emph{On the values of $G$-functions},
Commentarii Mathematici Helvetici \textbf{89} (2014), 313--341.
\url{https://doi.org/10.4171/CMH/321}.
Theorem numbering in this supplement follows the published version:
the Fuchsian/Frobenius theorem is Theorem 6, and the Wronskian statement
is Lemma 9.

\bibitem[VII.3]{p7:FR2024}
S. Fischler and T. Rivoal,
\emph{Relations between values of arithmetic Gevrey series, and
applications to values of the Gamma function},
Journal of Number Theory \textbf{261} (2024), 36--54.
\url{https://arxiv.org/abs/2301.13518}.
Author text:
\url{https://rivoal.perso.math.cnrs.fr/articles/ssmixte.pdf}.

\bibitem[VII.4]{p7:CellarosiSinai2011}
F. Cellarosi and Ya. G. Sinai,
\emph{The M\"obius function and statistical mechanics},
Bulletin of Mathematical Sciences \textbf{1} (2011), 245--275.
\url{https://doi.org/10.1007/s13373-011-0011-6}.
The normalized Dickman density and moment formulas are on
pp.~257--258.

\bibitem[VII.5]{p7:DLMF16}
NIST Digital Library of Mathematical Functions,
Section 16.17, \emph{Meijer $G$-Function: Definition}.
\url{https://dlmf.nist.gov/16.17}.

\bibitem[VII.6]{p7:DLMF5}
NIST Digital Library of Mathematical Functions,
Section 5.7, \emph{Gamma Function: Series Expansions}.
\url{https://dlmf.nist.gov/5.7}.

\end{thebibliography}
% END PRESERVED SOURCE BODY 7
\phantomsection\label{part7:end}
\clearpage
\setcounter{section}{0}\setcounter{subsection}{0}
\renewcommand{\thesection}{\Roman{part}.\arabic{section}}
\renewcommand{\theHsection}{\arabic{part}.main.\arabic{section}}
\renewcommand{\refname}{References for Part VIII}
\fancyhead[L]{\small Euler\textquotesingle s constant: Part VIII}
\part{Fractional degrees and Dickman-related asymptotics}
\label{part8:start}
\noindent\textbf{Source:} {\small\path{Keio_Gamma_Fractional_Degrees_and_Dickman_20261019.tex}}.

Original PDF: 29 pages. Mathematical content and source-level qualifications are retained.


% BEGIN PRESERVED SOURCE BODY 8
\begin{abstract}
For the signed remainders $R_N$ in Stirling's expansion of $\log\Gamma$,
we prove
\[
 2N+m-2<\deg_{\rm cm}[(-1)^mR_N^{(m)}]<2N+m-1
 \quad(N\ge2,\ m=0,1,2),
\]
and show that the positive fractional excess tends to zero as
$N\to\infty$. A spectral representation and a stability theorem for
small fractional derivatives supply the lower bounds; zero integral
moments force the strict upper bounds. Two boundary families also fail
at every derivative order $m\ge3$. These results refute five real-degree
formula families explicitly conjectured by Mahmoud and Qi (2022).
For Dickman-related questions, we prove divergence of a formal delay
equation solution and construct counterexamples to a conjectured
tilted-density limit, including one with a purely absolutely continuous
L\'evy measure and a proper density. A polynomial intertwiner proves
Atia's (2024) conjectured inverse Bessel recurrence. We also prove a
half-integral Laguerre coefficient statement printed as a conjecture
in a 2007 preprint, explain a factorial-error divergence observation,
and establish restricted Borel-division and lattice-obstruction results.
The source ledger distinguishes these outcomes and their priority
limitations. No irrationality or transcendence proof for Euler's
constant is obtained, and no claim of ten independent, previously open
problems newly solved is made.
\end{abstract}


\section{Results, provenance, and the arithmetic objective}
\label{p8:sec:ledger}
The starting objective is to obtain new information about the arithmetic
nature of
\[
 \gamma=\lim_{n\to\infty}\left(\sum_{k=1}^n\frac1k-\log n\right).
\]
The investigation instead yields several complete analytic and algebraic
arguments in adjacent areas, and identifies precise obstructions to
turning them into an arithmetic proof. The direct connections are
$\psi(1)=-\gamma$, the Dickman normalization $e^{-\gamma}$, and the
Laguerre--Gompertz approximation problem. The Bessel inverse recurrence
is an auxiliary special-function problem reached through these
asymptotic methods; it is not itself an irrationality criterion.

\subsection{Ten source-linked statements, with different statuses}
The ledger below records statements, not a claim of ten independent
major breakthroughs. Rows G1--G5 are five formula families in a single
published conjectural programme. Rows P1, P2, and A2 concern three other
explicit questions or conjectures. Row A1 is historical: its conjecture
label was removed from the final journal article. Row O1 is an explicit
numerical observation, not the source's irrationality conjecture.
Thus there are eight explicit conjectural/question statements in the
main source versions, one historical statement, and one observation;
they do not certify the requested total of ten independent currently
open problems newly resolved. The complete ledger is on the following page.

\newpage
\small
\begin{longtable}{@{}>{\raggedright\arraybackslash}p{.07\textwidth}>{\raggedright\arraybackslash}p{.31\textwidth}>{\raggedright\arraybackslash}p{.36\textwidth}>{\raggedright\arraybackslash}p{.17\textwidth}@{}}
\toprule
ID & Exact source & Conclusion established here & Qualification\\
\midrule
\endfirsthead
\toprule ID & Exact source & Conclusion established here & Qualification\\
\midrule
\endhead
G1 & Mahmoud--Qi \cite{p8:MahmoudQi2022}, Eq.~(2.5), $N\ge2$
 & $\deg_{\rm cm}R_N\in(2N-2,2N-1)$; conjectured equality is false.
 & One family, Theorem~\ref{p8:frontier:strip}.\\[5pt]
G2 & Same source, Eq.~(2.7), $N\ge2$
 & $\deg_{\rm cm}(-R_N')\in(2N-1,2N)$; conjectured equality is false.
 & Same structural method.\\[5pt]
G3 & Same source, first equality of Eq.~(2.8)
 & $\deg_{\rm cm}[(-1)^mR_0^{(m)}]<m$ for all $m\ge3$.
 & Theorem~\ref{p8:frontier:boundary}; $m=2$ is not refuted.\\[5pt]
G4 & Same source, second equality of Eq.~(2.8)
 & $\deg_{\rm cm}[(-1)^mR_1^{(m)}]<m+1$ for all $m\ge3$.
 & Same boundary argument.\\[5pt]
G5 & Same source, Eq.~(2.9), $m,N\ge2$
 & Entire $m=2$ subfamily has degree in $(2N,2N+1)$, refuting the universal equality.
 & Cases $m\ge3,N\ge2$ not classified.\\[5pt]
P1 & Dibl\'ik--Medina \cite{p8:DiblikMedina2018}, Remark~1, p.~104
 & Every nonzero formal solution in the stated inverse-power form diverges at each finite nonzero argument.
 & Journal question; Theorem~\ref{p8:pf:divergence}.\\[5pt]
P2 & Arratia--Kochman--Zabell \cite{p8:ArratiaKochmanZabell2016}, \S7, Conjecture~1
 & Density/local-limit assertion is false under its printed hypotheses, even for pure-AC infinite-activity L\'evy measure.
 & 2016 preprint; two counterexamples count once.\\[5pt]
A1 & Borwein--Borwein--Crandall \cite{p8:BBC2007}, Conjecture~1, p.~38
 & Half-integral coefficient factorization, with a stronger lower bound on surviving powers.
 & Historical status uncertain; Theorem~\ref{p8:thm:bbc-half-integral}.\\[5pt]
A2 & Atia \cite{p8:Atia2024}, \S5.4, Eq.~(24)
 & Inverse-linearization recurrence proved for every $a\ne1/2$.
 & Journal conjecture; Theorem~\ref{p8:thm:atia-recurrence}.\\[5pt]
O1 & Murabayashi--Yoshida \cite{p8:MurabayashiYoshida}, p.~15
 & Factorially cleared even and odd approximation errors tend to infinity.
 & Numerical observation; not Conjecture~7.5.\\
\bottomrule
\end{longtable}
\normalsize

All proofs are given in full below. The Laguerre proof uses an established
finite coefficient formula from its source, and the arithmetic division
argument explicitly invokes Lindemann--Weierstrass and the established
$E$-function division theorem. Independent internal readers checked the
arguments and source matching. This is not external peer review or a
machine-checked formal proof. Targeted follow-up searches did not locate
the same resolutions, but those searches cannot certify priority.

\subsection{Additional results and limits of the investigation}
The positive fractional degree corrections vanish at large index
(Theorem~\ref{p8:thm:degree-asymptotics}). The divergent Dickman series has
an entire Borel transform and a Cauchy-transform solution whose jump
is $2\pi i e^{-\gamma}\rho(x)$ (Proposition~\ref{p8:pf:moments}).
The lattice construction shows that smooth size ratios and primitivity
cannot alone force the balance condition in a proposed route to
$\gamma$ (Theorem~\ref{p8:ntf:counterfamily}). Rational Borel transforms
with positive algebraic poles satisfy a restricted arithmetic division
theorem (Theorem~\ref{p8:sf:real-sum}). These are substantive additional
statements, but none is relabeled as a complete solution of an
unrestricted published conjecture.

This manuscript is a self-contained account of the present investigation.
It does not reproduce every theorem from earlier drafts. In particular,
earlier positivity certificates, local parameter-identification results,
and failed irrationality routes are background rather than additional
entries in the count.


\section{Noninteger complete-monotonicity degrees and false boundary conjectures}
\label{p8:sec:frontier-inequalities}

\subsection{The exact published conjecture and the scope of the result}
Let $B_j$ be the Bernoulli numbers, with $B_2=1/6$, and define
\[
 R_N(x)=(-1)^N\left[\log\Gamma(x)-\left(x-\frac12\right)\log x+x
 -\frac12\log(2\pi)
 -\sum_{k=1}^N\frac{B_{2k}}{(2k)(2k-1)x^{2k-1}}\right].
\]
Qi's 2021 paper \emph{Bounds for completely monotonic degree of a remainder
for an asymptotic expansion of the trigamma function},
\url{https://doi.org/10.1080/25765299.2021.1962060}, conjectures that
$\deg_{\rm cm}R_2''=4$ and proves only $4\leq\deg_{\rm cm}R_2''\leq5$.
The author's manuscript is \url{https://arxiv.org/abs/2003.05300v2},
Eq.~(2.5). The larger conjecture, for all $m\geq2$ and $N\geq2$, is
\[
 \deg_{\rm cm}\bigl[(-1)^mR_N^{(m)}\bigr]=m+2(N-1).
\]
Mahmoud and Qi explicitly restate this as Eq.~(2.9), page~294, in their
2022 paper \emph{Bounds for completely monotonic degrees of remainders in
asymptotic expansions of the digamma function},
\url{https://doi.org/10.7153/mia-2022-25-17}.
The same paper also conjectures $\deg_{\rm cm}R_N=2N-2$ in
Eq.~(2.5) and $\deg_{\rm cm}(-R_N')=2N-1$ in Eq.~(2.7).
The theorem below disproves both of those families and the entire
$m=2$, $N\geq2$ subfamily of Eq.~(2.9).

For a nonzero completely monotone $F$ tending to zero at infinity, write
\[
 \deg_{\rm cm}F=\sup\{\alpha\in\mathbb R:x^\alpha F(x)
                 \text{ is completely monotone on }(0,\infty)\}.
\]
Here complete monotonicity means $(-1)^jF^{(j)}(x)\geq0$ for every
integer $j\geq0$ and every $x>0$.

\Needspace{13\baselineskip}
\begin{theorem}\label{p8:frontier:strip}
For every integer $N\geq2$, put
\[
 Q_N(x)=R_N''(x)
 =(-1)^N\left[\psi'(x)-\frac1x-\frac1{2x^2}
                  -\sum_{k=1}^N\frac{B_{2k}}{x^{2k+1}}\right].
\]
Then, for $m=0,1,2$,
\[
 \boxed{2N+m-2<\deg_{\rm cm}\bigl[(-1)^mR_N^{(m)}\bigr]
                      <2N+m-1.}
\]
In particular, all three families have noninteger real degrees. The
theorem does not specify their exact fractional parts.
\end{theorem}

\subsection{A positivity lemma for small fractional derivatives}
\begin{lemma}\label{p8:frontier:fractional}
Suppose $f\in C^1([0,\infty))$, $f(0)=0$, $f(t)>0$ for $t>0$, and
$f'(t)>0$ on some $(0,\delta)$. Suppose also that $f(t)\to L>0$ and
$|f'(t)|\leq C(1+t)^{-p}$ for all sufficiently large $t$, where
$C>0$ and $p>2$. Then there is $\varepsilon>0$ such that
$x^\varepsilon\int_0^\infty e^{-xt}f(t)\,dt$ is completely monotone.
\end{lemma}
\begin{proof}
For $0<\varepsilon<1$, define
\[
 D^\varepsilon f(t)=\frac1{\Gamma(1-\varepsilon)}
     \int_0^t(t-u)^{-\varepsilon}f'(u)\,du.
\]
We show that it is positive for all $t>0$ when $\varepsilon$ is
sufficiently small. The sign is that of
\[
 J_\varepsilon(t)=\int_0^t f'(u)(1-u/t)^{-\varepsilon}\,du.
\]
It is positive for $0<t\leq\delta$ because $f'>0$ there. Fix any
$0<\varepsilon_0<1$. For $0\leq\varepsilon\leq\varepsilon_0$, splitting
the integral at $t/2$ gives
\[
 J_\varepsilon(t)=L+O(t^{-1})
                  +O(t^{1-p})
                  \qquad(t\to\infty),
\]
uniformly in $\varepsilon$. Indeed, on $0\leq u\leq t/2$ the difference
$(1-u/t)^{-\varepsilon}-1$ is bounded by a constant times $u/t$;
$\int_0^\infty u|f'(u)|du<\infty$. On the other half the polynomial decay
bound and $\int_{1/2}^1(1-v)^{-\varepsilon_0}dv<\infty$ apply.
Consequently there is $T$ such that $J_\varepsilon(t)>L/2$ for $t\geq T$
and all these $\varepsilon$.

On $[\delta,T]$, if $M=\max_{0\leq u\leq T}|f'(u)|$, then
\[
 |J_\varepsilon(t)-f(t)|
 \leq Mt\left[\frac1{1-\varepsilon}-1\right]
 \leq MT\frac{\varepsilon}{1-\varepsilon}.
\]
Since $\min_{[\delta,T]}f>0$, this proves positivity on the remaining
compact interval for sufficiently small $\varepsilon>0$.
Finally, absolute Fubini and $f(0)=0$ give
\[
 \int_0^\infty e^{-xt}D^\varepsilon f(t)dt
 =x^{\varepsilon-1}\int_0^\infty e^{-xt}f'(t)dt
 =x^\varepsilon\int_0^\infty e^{-xt}f(t)dt.
\]
The positive Laplace density proves complete monotonicity.
\end{proof}

\subsection{The spectral kernel and its unavoidable overshoot}
\begin{proof}[Proof of Theorem~\ref{p8:frontier:strip}]
Put $g(t)=t/(e^t-1)$, with its smooth value at zero, and
\[
 h_N(t)=(-1)^N\left[g(t)-1+\frac t2
                   -\sum_{k=1}^N\frac{B_{2k}t^{2k}}{(2k)!}\right].
\]
The standard Laplace representation of $\psi'$ gives $Q_N=\mathcal Lh_N$.
Every derivative of $h_N$ of order at most $2N+1$ vanishes at zero.
Set $f_N=h_N^{(2N)}$. Repeated integration by parts gives
\begin{equation}
 x^{2N}Q_N(x)=\int_0^\infty e^{-xt}f_N(t)dt,
 \qquad
 f_N(t)=|B_{2N}|+(-1)^N g^{(2N)}(t).
 \label{p8:frontier:kernel}
\end{equation}

We require strict positivity, including the absence of interior zeros.
The Mittag--Leffler expansion of $\coth$ gives the exact identity
\begin{equation}
 g(t)+\frac t2
 =1+\int_0^\infty(1-\cos(tu))\nu(u)du,
 \qquad \nu(u)=\frac{\pi}{\sinh^2(\pi u)}.
 \label{p8:frontier:spectral}
\end{equation}
To check the normalization, expand
$\nu(u)=4\pi\sum_{j\geq1}j e^{-2\pi ju}$ and use Tonelli:
\[
 \int_0^\infty(1-\cos(tu))\nu(u)du
 =2t^2\sum_{j\geq1}\frac1{t^2+4\pi^2j^2}
 =\frac t2\coth\frac t2-1.
\]
The last equality is the classical partial-fraction expansion of $\coth$
(NIST DLMF~4.36.3, \url{https://dlmf.nist.gov/4.36.E3}).
Because $u^{2N}\nu(u)$ is integrable for $N\geq1$, differentiation
of \eqref{p8:frontier:spectral} through the even order $2N$, including at
zero, is legitimate. It yields
\begin{equation}
 f_N(t)=\int_0^\infty u^{2N}(1-\cos(tu))\nu(u)du>0\quad(t>0).
 \label{p8:frontier:strictpositive}
\end{equation}
Strictness follows because $\nu(u)>0$ and the zeros of $1-\cos(tu)$
form a discrete set. Moreover,
\[
 f_N(t)=c_Nt^2+O(t^4),\qquad
 c_N=\frac12\int_0^\infty u^{2N+2}\nu(u)du>0,
\]
so $f_N'>0$ near zero. Equation~\eqref{p8:frontier:kernel}, or the series
$g(t)=t\sum_{j\geq1}e^{-jt}$, shows that
\[
 f_N(t)\longrightarrow L_N:=|B_{2N}|>0,
 \qquad f_N'(t)=O_N((1+t)e^{-t}).
\]
Lemma~\ref{p8:frontier:fractional} applies. Together with
\eqref{p8:frontier:kernel}, it proves that $x^{2N+\varepsilon_N}Q_N(x)$ is
completely monotone for some $\varepsilon_N>0$. This proves the strict
lower bound.

For the upper bound the condition $N\geq2$ is essential. Since
$B_{2N-1}=0$, absolute integration gives
\begin{equation}
 \int_0^\infty[f_N(t)-L_N]dt
 =(-1)^N\int_0^\infty g^{(2N)}(t)dt
 =-(-1)^NB_{2N-1}=0.
 \label{p8:frontier:zerointegral}
\end{equation}
The integrand is negative near zero, because $f_N(0)=0$. Hence it must
be positive somewhere. Thus $f_N$ exceeds its limiting value $L_N$ at
some point, and $f_N'$ is negative somewhere. Since $f_N(0)=0$,
\[
 x^{2N+1}Q_N(x)=\int_0^\infty e^{-xt}f_N'(t)dt.
\]
Uniqueness of Laplace transforms excludes complete monotonicity: a
continuous signed density which is negative on an interval cannot equal
a nonnegative representing measure. Uniqueness applies after multiplying
both measures by $e^{-at}$ for any $a>0$, making them finite.

If some exponent larger than $2N+1$ were admissible, multiplication by
a completely monotone negative power of $x$ would make $2N+1$
admissible as well. Finally, the admissible set of exponents is closed:
for fixed $x>0$ and integer $j\geq0$, the quantity
$(-1)^j\partial_x^j[x^\alpha Q_N(x)]$ is continuous in $\alpha$.
The supremum therefore cannot equal the excluded endpoint $2N+1$.
This proves the strict upper bound for $m=2$.

To transfer the argument to $m=1,0$, put $r=2-m\in\{1,2\}$ and
$d=2N-r$. The standard Binet formula and differentiation give
\[
 (-1)^mR_N^{(m)}(x)=\int_0^\infty e^{-xt}\frac{h_N(t)}{t^r}dt.
\]
Set $k_r=h_N/t^r$ and $f_r=k_r^{(d)}$. Taylor's integral formula
and the vanishing derivatives of $h_N$ at zero give
\begin{equation}
 f_r(t)=\frac1{(r-1)!}\int_0^1(1-s)^{r-1}s^d f_N(st)ds.
 \label{p8:frontier:transfer}
\end{equation}
Consequently $f_r(t)>0$ for $t>0$, $f_r(t)\sim c_rt^2$ at zero
with $c_r>0$, and
\[
 L_r:=\lim_{t\to\infty}f_r(t)=L_N\frac{d!}{(2N)!}>0.
\]
Differentiating \eqref{p8:frontier:transfer} and substituting $u=st$
shows
\[
 |f_r'(t)|=O_N(t^{-d-2}),\qquad
 \int_0^\infty u^{d+1}|f_N'(u)|du<\infty.
\]
Here $d\geq2$, so
the fractional positivity lemma applies once again. Repeated integration
by parts gives $x^d(-1)^mR_N^{(m)}=\mathcal Lf_r$, proving strict
lower bound $d$.

The same transfer preserves the zero integral responsible for the
strict upper bound. Absolute Fubini is legitimate because
$f_N-L_N$ is integrable and
$\int_0^1(1-s)^{r-1}s^{d-1}ds<\infty$. Thus
\[
 \int_0^\infty[f_r(t)-L_r]dt
 =\frac1{(r-1)!}\int_0^1(1-s)^{r-1}s^{d-1}ds
           \int_0^\infty[f_N(u)-L_N]du=0.
\]
Since $f_r(0)=0<L_r$, this forces an overshoot of $L_r$, and hence
a negative value of $f_r'$. The preceding Laplace-uniqueness and
closedness argument excludes $d+1$ with a strict gap. This proves
the remaining two cases of the theorem.
\end{proof}

\subsection{Both boundary conjectures fail at every derivative order beyond two}
The same 2022 paper, Eq.~(2.8), proposes
\[
 \deg_{\rm cm}\bigl[(-1)^mR_0^{(m)}\bigr]=m,
 \qquad
 \deg_{\rm cm}\bigl[(-1)^mR_1^{(m)}\bigr]=m+1
 \quad(m\geq2).
\]
Remark~4.5 explicitly leaves these equalities open. Although they are
correct at $m=2$, every subsequent case fails.

\begin{theorem}\label{p8:frontier:boundary}
For every integer $m\geq3$,
\[
 \deg_{\rm cm}\bigl[(-1)^mR_0^{(m)}\bigr]<m,
 \qquad
 \deg_{\rm cm}\bigl[(-1)^mR_1^{(m)}\bigr]<m+1.
\]
\end{theorem}
\begin{proof}
Let
\[
 a_m(t)=\frac{t^{m-1}}{e^t-1},\quad
 k_0(t)=a_m(t)-t^{m-2}+\frac12t^{m-1},\quad
 k_1(t)=-a_m(t)+t^{m-2}-\frac12t^{m-1}+\frac1{12}t^m.
\]
The Binet representation gives
$(-1)^mR_j^{(m)}=\mathcal Lk_j$, $j=0,1$. These functions are nonzero
and completely monotone: the partial-fraction expansion above gives
$0<h_0(t)<t^2/12$ for $t>0$, while $h_1(t)=t^2/12-h_0(t)$ and
$k_j=t^{m-2}h_j$. The expansion at zero is
\[
 a_m(t)=t^{m-2}-\frac12t^{m-1}+\frac1{12}t^m
                 -\frac1{720}t^{m+2}+O(t^{m+4}).
\]
Thus $k_0=O(t^m)$ and $k_1=O(t^{m+2})$. All boundary terms needed
for the following integrations by parts vanish:
\[
 x^m(-1)^mR_0^{(m)}=\mathcal L A_m,\qquad
 x^{m+1}(-1)^mR_1^{(m)}=\mathcal L C_m,
 \quad A_m=a_m^{(m)},\quad C_m=-a_m^{(m+1)}.
\]
Both densities have exponentially decaying tails, and at zero
\[
 A_m(t)=\frac{m!}{12}+O(t^2)>0,\qquad
 C_m(t)=\frac{(m+2)!}{720}t+O(t^3)>0.
\]
However, $a_m^{(m-3)}(0)=0$ for every $m\geq3$. Repeated integration
by parts, first against $t^2$ and then against $t^3$, gives the exact
vanishing moments
\[
 \int_0^\infty t^2 A_m(t)dt=-2a_m^{(m-3)}(0)=0,
 \qquad
 \int_0^\infty t^3 C_m(t)dt=-6a_m^{(m-3)}(0)=0.
\]
Since each density is positive on an initial interval, it must be
negative somewhere else. Laplace uniqueness therefore excludes complete
monotonicity at each conjectured endpoint. Downward closure and
closedness of the admissible exponents, proved above, make both upper
bounds strict.

As an independent tail check, for odd $m\geq3$ the two densities are
already negative eventually:
\[
 A_m(t)=(-1)^m t^{m-1}e^{-t}\bigl(1+O_m(t^{-1})\bigr),\qquad
 C_m(t)=(-1)^m t^{m-1}e^{-t}\bigl(1+O_m(t^{-1})\bigr).
\]
The zero-moment argument also covers every even $m\geq4$.
\end{proof}

\subsection{Priority distinctions and rejected candidates}
Ling Zhu's 2020 paper \emph{Completely monotonic integer degrees for a
class of special functions}, \url{https://doi.org/10.3934/math.2020224},
conjectures the integer degree $2N$ in the $m=2$ case of Eq.~(6.3).
Theorem~\ref{p8:frontier:strip} is consistent with, and proves, this integer
statement, together with the corresponding $m=0,1$ integer statements.
It contradicts the separately stated real-degree equalities.
These consequences concern one structural family and must not be
counted as many unrelated solved problems.

There is an earlier relevant observation. Qi and Guo's
\emph{Integral representations and complete monotonicity of remainders
of the Binet and Stirling formulas for the gamma function},
RACSAM~111 (2017), 425--434,
\url{https://doi.org/10.1007/s13398-016-0302-6}, Remark~3.1,
reports the absence of a concrete nontrivial example having noninteger
degree. Theorem~\ref{p8:frontier:strip} supplies three infinite Gamma
remainder families with this property. This contextual observation is
not counted as another independently resolved conjecture.

Two other explicit questions were screened. Section~3 of Bulboac\u{a}
and Zayed's \emph{Monotonic nature of the Gamma function},
\url{https://doi.org/10.1186/s13660-025-03425-0}, asks for an analytic
proof that the continuously extended quotient
\[
 \frac{\log\Gamma(x+1)}{\log(x^2+6)-\log(x+6)}
\]
has exactly one minimum on $(-1,\infty)$. Open Problem~1.5 of Szab\'o's
\emph{Completely monotone functions in general and some applications},
\url{https://arxiv.org/abs/2411.17670}, asks for the sharp exponent in
\[
 (x+a)^\alpha\left[\psi(x+b)-\psi(x+a)-\frac{b-a}{x+a}\right],
 \qquad 0<b-a<1.
\]
The answer to the latter is $\alpha_0=1$: sufficiency is proved in the
source, and necessity follows because the bracket is
$(1-b+a)/(x+a)+O(1)$ at $x=-a$. For $\alpha>1$ the function would tend
to zero at its left endpoint, impossible for a nonzero completely
monotone function. However, a public repository already lists solution
claims for both of these exact questions (entries APP-0029/0034 and
APP-0021 at
\url{https://github.com/pudim-project/pudim-ai-demo-zetalaw}).
They are excluded from any claim of a new solution here.

The targeted search also inspected the 2022 Mahmoud--Qi paper,
Zhu's integer-degree paper, the author's 2024 degree bibliography,
and relevant 2024--2025 followups. No prior resolution of the
$Q_N$ real-degree equality was found. This is a bounded literature audit,
not a claim that priority has been exhaustively established. The proof
above is independent of the repository's mathematical claims.

\subsection{The fractional parts tend to zero}
\label{p8:subsec:degree-asymptotics}
The strict intervals in Theorem~\ref{p8:frontier:strip} do not merely give
counterexamples at a few small indices. Their noninteger corrections
vanish in the large-index limit.

\begin{theorem}\label{p8:thm:degree-asymptotics}
For each fixed $m\in\{0,1,2\}$,
\[
 \deg_{\rm cm}\bigl[(-1)^m R_N^{(m)}\bigr]
       =2N+m-2+o(1)\qquad(N\longrightarrow\infty),
\]
where the $o(1)$ is strictly positive for every $N\ge2$.
\end{theorem}
\begin{proof}
Put $r=2-m$, $d=2N-r$, $b_N=(2N+1)/(2\pi)$, and
$L_N=|B_{2N}|$. Let $U_N>0$ have probability density
$u^{2N}\nu(u)/L_N$, where $\nu$ is defined in
\eqref{p8:frontier:spectral}. Expanding $\nu(u)=4\pi\sum_{j\ge1}je^{-2\pi ju}$
and integrating its moments yields, with $W_N=U_N/b_N$,
\[
 \mathbb E W_N=\frac{\zeta(2N+1)}{\zeta(2N)},\qquad
 \mathbb E W_N^2=\frac{2N+2}{2N+1}
                    \frac{\zeta(2N+2)}{\zeta(2N)}.
\]
Thus $W_N\to1$ in $L^2$, and hence in $L^1$.
For $r=0$ set $S_N=1$. For $r=1,2$ let $S_N$ be independent of $U_N$
and have beta density proportional to
$s^{2N-r}(1-s)^{r-1}$ on $(0,1)$. Its mean is
$(2N-r+1)/(2N+1)$, so $S_N\to1$ in $L^1$ as well. Consequently
$A_N=S_NW_N\to1$ in $L^1$.

For $r=0$ write $f_{N,0}=f_N$, and for $r=1,2$ use $f_{N,r}=f_r$
from \eqref{p8:frontier:transfer}. Let $L_{N,r}$ be the corresponding
positive limit at infinity. The spectral and beta representations give
\[
 G_N(v):=\frac{f_{N,r}(v/b_N)}{L_{N,r}}
       =\mathbb E[1-\cos(v A_N)],\qquad
 G_N'(v)=\mathbb E[A_N\sin(v A_N)].
\]
Differentiation is justified by $\mathbb E A_N<\infty$. For every $T>0$,
\[
 \sup_{0\le v\le T}|G_N'(v)-\sin v|
 \le(1+T)\mathbb E|A_N-1|\longrightarrow0.
\]
Fix any $\varepsilon\in(0,1)$. At $T=2\pi$, the integrable fractional
kernel and this uniform convergence give
\[
 D^\varepsilon G_N(2\pi)\longrightarrow
 \frac1{\Gamma(1-\varepsilon)}
 \int_0^{2\pi}\sin u\,(2\pi-u)^{-\varepsilon}\,du<0.
\]
The strict sign follows on pairing $u$ with $2\pi-u$:
\[
 \int_0^{2\pi}\sin u\,(2\pi-u)^{-\varepsilon}\,du
 =\int_0^\pi\sin u\bigl[(2\pi-u)^{-\varepsilon}
                            -u^{-\varepsilon}\bigr]du<0.
\]
Scaling preserves the sign of a fractional derivative, so
$D^\varepsilon f_{N,r}$ is negative somewhere for all sufficiently
large $N$. By the same Laplace uniqueness argument used above,
$x^{d+\varepsilon}(-1)^mR_N^{(m)}(x)$ is then not completely monotone.
Downward closure of the admissible exponents gives
$\deg_{\rm cm}[(-1)^mR_N^{(m)}]\le d+\varepsilon$ eventually.
Theorem~\ref{p8:frontier:strip} supplies the strict lower bound $d$.
Since $\varepsilon>0$ was arbitrary, the asserted limit follows.
\end{proof}

\subsection{What this does, and does not, say about Euler's constant}
At $N=2$, the three completely monotone functions take the values
\[
 R_2(1)=\frac{331}{360}-\frac12\log(2\pi),\qquad
 -R_2'(1)=\gamma-\frac{23}{40},\qquad
 R_2''(1)=\frac{\pi^2}{6}-\frac{49}{30}.
\]
In particular, the middle function contains $\gamma$ directly, rather
than merely sharing its name with the gamma function. The method transfers
between its Stirling, digamma, and trigamma forms through positive
integral kernels. The new information is an exact exclusion of integer
values for their real complete-monotonicity degrees, together with an
asymptotic determination up to $o(1)$.

Complete monotonicity is an analytic sign condition. It does not provide
integer coefficients with controlled denominators. None of these degree
results implies that $\gamma$ is irrational or transcendental, and no
such implication is asserted.

% Fragment: requires amsmath, amssymb, amsthm, hyperref and theorem environments.
\section{Two explicit Dickman-related questions from the literature}
\label{p8:sec:probability-frontier}

This section addresses two precisely identified statements. The first is a
question in a refereed 2018 article; the second is a conjecture in a publicly
available 2016 research preprint. These publication statuses are different.
The proofs below answer the statements as printed. A search of the accessible
subsequent literature did not locate these answers, but does not establish
priority or certify that no unpublished or unindexed resolution exists.
Neither result decides the arithmetic nature of Euler's constant.

\subsection{Divergence of the Dibl\'ik--Medina formal solution}

Dibl\'ik and Medina \cite[Remark~1, p.~104]{p8:DiblikMedina2018} ask whether the
formal solution
\begin{equation}\label{p8:pf:formal}
 S(t)=\sum_{n\geq1}C_nt^{-n},\qquad
 (n-1)C_n=\sum_{s=1}^{n-1}\binom{n-1}{s-1}C_s\quad(n\geq2),
\end{equation}
of $t x'(t)=-x(t-1)$ converges or diverges. The initial coefficient $C_1$ is
arbitrary. Their question has the following direct answer.

\begin{theorem}\label{p8:pf:divergence}
If $C_1\ne0$, the power series $\sum_{n\geq1}C_nz^n$ has radius of
convergence zero. In particular, the series \eqref{p8:pf:formal} diverges at
every finite nonzero real or complex $t$. More explicitly,
\begin{equation}\label{p8:pf:factorial}
 |C_{2k+2}|\geq |C_1|k!\qquad(k\geq0).
\end{equation}
If $C_1=0$, all coefficients vanish.
\end{theorem}
\begin{proof}
Set $a_n=C_n/C_1$ when $C_1\ne0$. The recurrence gives $a_1=a_2=1$ and
$a_n>0$ for every $n$. Keeping just the summand $s=n-2$ gives, for $n\geq3$,
\[
 a_n\geq \frac{1}{n-1}\binom{n-1}{n-3}a_{n-2}
       =\frac{n-2}{2}a_{n-2}.
\]
In particular $a_{2k+2}\geq k a_{2k}$, which proves \eqref{p8:pf:factorial}
by induction. For fixed $t\ne0$, the lower bound
$|C_1|k!|t|^{-2k-2}$ tends to infinity, so the terms of the proposed series
do not tend to zero. The case $C_1=0$ follows by induction directly from
the recurrence.
\end{proof}

The next identification explains why a moment problem appears in this
delay equation. It is additional structure, not a second independent
resolution of a published question.

\begin{proposition}\label{p8:pf:moments}
Let $D$ have the standard Dickman distribution, with density
$p(x)=e^{-\gamma}\rho(x)$ and moment generating function
\[
 M(z)=\exp\left(\int_0^1\frac{e^{zx}-1}{x}\,dx\right).
\]
Then $C_n=C_1\mathbb E[D^{n-1}]$, and the formal Borel transform is entire:
\begin{equation}\label{p8:pf:borel}
 \sum_{n\geq1}\frac{C_n}{(n-1)!}z^{n-1}=C_1M(z).
\end{equation}
For $w\in\mathbb C\setminus[0,\infty)$, the Cauchy transform
\[
 F(w)=\int_0^\infty\frac{p(x)}{w-x}\,dx
\]
is an actual holomorphic solution of
\begin{equation}\label{p8:pf:cauchy-dde}
 w F'(w)=-F(w-1).
\end{equation}
On every closed sector bounded away from the positive real axis it has
the asymptotic expansion $F(w)\sim\sum_{n\geq1}(C_n/C_1)w^{-n}$.
Its boundary values satisfy
\begin{equation}\label{p8:pf:jump}
 F(x-i0)-F(x+i0)=2\pi i e^{-\gamma}\rho(x),\qquad x>0.
\end{equation}
\end{proposition}
\begin{proof}
The standard distribution and normalization are recalled in
\cite[Section~1]{p8:PinskyDickman2016}. If $m_j=\mathbb E[D^j]$, then
$zM'(z)=(e^z-1)M(z)$ gives
$j m_j=\sum_{i=0}^{j-1}\binom ji m_i$. This is precisely
\eqref{p8:pf:formal} after $j=n-1$, with the same initial value. The
integrand defining $\log M$ has a removable singularity at $x=0$ and
defines an entire function of $z$, proving \eqref{p8:pf:borel}.

For the Cauchy transform, differentiation and integration by parts are
valid off $[0,\infty)$ because $p$ has all moments, is locally absolutely
continuous on $(0,\infty)$, is constant on $(0,1)$, and decreases faster
than any power. The relation $xp'(x)=-p(x-1)$, with $p=0$ on the negative
axis, gives
\[
 wF'(w)=-\int_0^\infty\frac{wp(x)}{(w-x)^2}\,dx
       =\int_0^\infty\frac{xp'(x)}{w-x}\,dx
       =-\int_0^\infty\frac{p(y)}{w-1-y}\,dy.
\]
For every $N\geq1$, geometric division gives the exact remainder
\[
 F(w)-\sum_{j=0}^{N-1}\frac{m_j}{w^{j+1}}
   =\frac1{w^N}\int_0^\infty\frac{x^Np(x)}{w-x}\,dx.
\]
In the stated sectors $|w-x|\geq c|w|$ for $x\geq0$, so the remainder
is at most $c^{-1}m_N|w|^{-N-1}$. Finally, the usual Cauchy boundary-value
formula gives \eqref{p8:pf:jump}; it also follows directly by subtracting
the two kernels and applying the Poisson approximate identity to the
continuous density $p$ at each $x>0$.
\end{proof}

Thus the divergent dominant expansion and the rapidly decaying Dickman
solution are connected by an exact Cauchy-transform jump. The factor
$e^{-\gamma}$ is a known normalization, and this identity alone gives
no rationality or transcendence conclusion about $\gamma$.

\subsection{A counterexample to a tilted-density conjecture}

Conjecture~1 of \cite[Section~7, p.~29]{p8:ArratiaKochmanZabell2016} permits a
L\'evy measure $\mu$ on $(0,1]$ with atoms, requires
$\int x\,\mu(dx)<\infty$, and assumes that its absolutely continuous part
$g(x)\,dx$ is bounded below by a positive constant on some nonempty
interval. With
\[
 C(\beta)=\int_{(0,1]}(e^{\beta x}-1)\,\mu(dx),\quad
 u=C'(\beta),\quad \sigma_\beta^2=C''(\beta),
\]
the conjecture asserts a pointwise Gaussian density limit for the
standardized exponentially tilted law, and consequently the density
asymptotic
\begin{equation}\label{p8:pf:conjectured}
 f(u)\sim G(u):=\frac{\exp(C(\beta(u))-u\beta(u))}
                    {\sqrt{2\pi}\,\sigma_{\beta(u)}}
 \qquad(u\longrightarrow\infty).
\end{equation}
The version inspected is arXiv:1606.03524v1 (2016); no journal publication
was identified. The permitted atomic part creates a counterexample even
when the resulting law has a proper absolutely continuous density.

\begin{theorem}\label{p8:pf:counterexample}
Let
\begin{equation}\label{p8:pf:levy}
 \mu(dx)=\delta_1(dx)+\frac12\mathbf1_{(0,1/2]}(x)\frac{dx}{x}.
\end{equation}
This measure satisfies all the hypotheses of the cited conjecture.
Its infinitely divisible law has a proper density, but
\eqref{p8:pf:conjectured} fails. In fact there are intervals $I_n$ of positive
length, tending to infinity, on which $f(u)/G(u)$ tends uniformly to
infinity. Thus altering a density on a null set cannot restore the
conjectured asymptotic. For the canonical tilted densities, even the
pointwise local limit at the fixed standardized point $y=0$ fails.
\end{theorem}
\begin{proof}
Let $N\sim\operatorname{Poisson}(1)$ and let $D_{1/2}$ have the generalized
Dickman law, independently. Then $T=N+X$, where $X=D_{1/2}/2$, has the
L\'evy measure \eqref{p8:pf:levy}. Its mean is $1+1/4$; its absolutely
continuous L\'evy density is at least $1$ on $(1/4,1/2)$. In particular
the hypotheses hold, and the L\'evy measure even has infinite total mass.

The generalized Dickman density formula
\cite[Section~1]{p8:PinskyDickman2016} gives
\[
 f_X(x)=A x^{-1/2}\quad(0<x<1/2),\qquad
 A=\frac{\sqrt2\,e^{-\gamma/2}}{\sqrt\pi}>0.
\]
Consequently $T$ is absolutely continuous, with density
\[
 f(t)=e^{-1}\sum_{k\geq0}\frac{f_X(t-k)}{k!},
 \qquad f_X(s)=0\quad(s\leq0).
\]
On each bounded interval this sum has only finitely many nonzero terms.
In particular, for every integer $n\geq0$ and $0<\varepsilon<1/2$,
\begin{equation}\label{p8:pf:spike}
 f(n+\varepsilon)\geq\frac{e^{-1}A}{n!\sqrt\varepsilon}.
\end{equation}

Here
\[
 C(\beta)=e^\beta-1+
       \frac12\int_0^{1/2}\frac{e^{\beta x}-1}{x}\,dx,
 \quad
 C'(\beta)=e^\beta+\frac12\int_0^{1/2}e^{\beta x}\,dx.
\]
Since $C''(\beta)>0$, $C'$ maps $[0,\infty)$ bijectively onto
$[5/4,\infty)$. For $u\geq5/4$ its inverse $\beta(u)$ is nonnegative,
and $C''(\beta)\geq e^\beta\geq1$. Convexity, together with $C(0)=0$,
gives $C(\beta)-\beta C'(\beta)\leq0$. It follows that
\begin{equation}\label{p8:pf:Gbound}
 0<G(u)\leq\frac1{\sqrt{2\pi}}\qquad(u\geq5/4).
\end{equation}
Set $\varepsilon_n=e^{-n^3}$ and
$I_n=(n+\varepsilon_n,\,n+2\varepsilon_n)$ for $n\geq2$.
Combining \eqref{p8:pf:spike} and \eqref{p8:pf:Gbound} yields, throughout $I_n$,
\[
 \frac{f(u)}{G(u)}\geq
       \frac{\sqrt\pi\,e^{-1}A}{n!}\exp(n^3/2)
       \longrightarrow\infty.
\]
Every $I_n$ has positive length, so this remains incompatible with a
pointwise asymptotic after any modification on a fixed null set.

For completeness, let $u_n=n+\varepsilon_n$ and
$\beta_n=\beta(u_n)$. The standardized tilted variable
$Y_\beta=(T_\beta-C'(\beta))/\sigma_\beta$ has its natural density
\[
 f_{Y_\beta}(y)=\sigma_\beta
 e^{\beta(u+\sigma_\beta y)-C(\beta)}
 f(u+\sigma_\beta y).
\]
At $y=0$ this is
$f(u)/(\sqrt{2\pi}\,G(u))$. Hence
$f_{Y_{\beta_n}}(0)\to\infty$, instead of $(2\pi)^{-1/2}$.
Here $\beta_n\to\infty$, because $u_n\to\infty$ and $C'$ is increasing.
\end{proof}

\begin{remark}
This refutation concerns the density assertion in the full generality
of the stated conjecture. It does not contradict the paper's theorems
under stronger regularity assumptions, nor its weak central limit
theorem. It also does not assert failure of the corresponding tail
probability asymptotic. The obstruction is an arbitrarily narrow cusp
near every integer, created by convolving a Poisson lattice law with
a generalized Dickman density singular at zero.
\end{remark}

% Bibliography entries to merge into the main manuscript:
% \bibitem{DiblikMedina2018}
% J. Dibl\'ik and R. Medina, Exact asymptotics of positive solutions to
% Dickman equation, Discrete Contin. Dyn. Syst. B 23 (2018), 101--121.
% doi:10.3934/dcdsb.2018007.
% \bibitem{ArratiaKochmanZabell2016}
% R. Arratia, F. Kochman and S. Zabell, Large deviation asymptotics for a
% random variable with L\'evy measure supported by $[0,1]$,
% arXiv:1606.03524v1 (2016), Section 7, Conjecture 1.
% \bibitem{PinskyDickman2016}
% R. G. Pinsky, On the strange domain of attraction to generalized Dickman
% distributions for sums of independent random variables,
% arXiv:1611.07207v3 (2017), Section 1.

% Stronger alternative to probability_nonatomic.tex; same conjecture, no new count.
\subsection{A purely continuous L\'evy measure and a proper density}

The preceding counterexample to the Arratia--Kochman--Zabell conjecture
can be strengthened: neither atoms in the L\'evy measure nor atoms in
the resulting law are necessary. This is a second construction refuting
the same conjecture and is not counted as a separate resolution.

\begin{proposition}\label{p8:pf:nonatomic-proper}
The conclusion \eqref{p8:pf:conjectured} can fail when the L\'evy measure is
purely absolutely continuous, has infinite total mass, and the resulting
law has a proper density. The failure occurs on positive-length
intervals tending to infinity.
\end{proposition}
\begin{proof}
Set $a=1/2$, $b=1/4$, $\delta=e^{-2}$ and
\[
 h(v)=\frac{\mathbf1_{(0,\delta)}(v)}{v(\log v)^2},
 \qquad
 \mu(dx)=\left\{h(x-a)+\frac{\mathbf1_{(0,b]}(x)}{2x}\right\}dx.
\]
This measure is supported in $(0,1)$, has finite first moment, and has
density at least $2$ on $(b/2,b)$. It is purely absolutely continuous
and has infinite total mass. The translated component has mass
$\lambda=\int_0^\delta h(v)\,dv=1/2$.

Let $J$ be the compound Poisson variable with L\'evy density $h(x-a)$,
and independently let $X=bD_{1/2}$. Then $T=J+X$ has L\'evy measure
$\mu$ and a proper density. The generalized Dickman density gives
\[
 f_X(x)=A x^{-1/2}\quad(0<x<b),\qquad
 A=\frac{2e^{-\gamma/2}}{\sqrt\pi}>0.
\]
Writing $h^{*0}$ for unit mass at zero, a density of $T$ is
\[
 f(t)=e^{-\lambda}\sum_{k\geq0}
          \frac{(f_X*h^{*k})(t-ka)}{k!}.
\]
For each $t$ the sum has only finitely many nonzero terms, because every
jump in $J$ exceeds $a$.

For $n\geq2$ and $0<\varepsilon<\min(b,\delta)$, restrict all $n$
integration variables in $(f_X*h^{*n})(\varepsilon)$ to
$(0,\varepsilon/(2n))$. Their sum is at most $\varepsilon/2$, so the
remaining argument of $f_X$ lies in $[\varepsilon/2,\varepsilon]$.
Using $\int_0^r h(v)\,dv=1/\log(1/r)$, we obtain
\[
 f(na+\varepsilon)\geq
 \frac{e^{-1/2}A\varepsilon^{-1/2}}
 {n!\{\log(1/\varepsilon)+\log(2n)\}^{n}}.
\]
Consequently, uniformly for
$\varepsilon\in(e^{-n^3},2e^{-n^3})$,
\[
 f(na+\varepsilon)\geq
 \frac{e^{-1/2}A e^{n^3/2}}
 {\sqrt2\,n!\{n^3+\log(2n)\}^{n}}
 \longrightarrow\infty.
\]

For the cumulant $C$ of this law, $C'$ is strictly increasing with range
$(0,\infty)$ and $\beta(u)\geq0$ for all sufficiently large $u$.
The translated component alone gives
$C''(\beta)\geq C''(0)\geq a^2\lambda=1/8$ for $\beta\geq0$.
Convexity and $C(0)=0$ give $C(\beta)-\beta C'(\beta)\leq0$, hence
$G(u)\leq2/\sqrt\pi$. Therefore $f/G\to\infty$ uniformly on
$(na+e^{-n^3},na+2e^{-n^3})$. These intervals have positive length,
so no change of the density on a null set can restore
\eqref{p8:pf:conjectured}.
\end{proof}

% Fragment for the parent manuscript; theorem environments are supplied there.
\section{Half-integral vanishing in the Laguerre asymptotic coefficients}
\label{p8:sec:bbc-laguerre}

\paragraph{Source and scope.}
The publicly posted version of D.~Borwein, J.~M.~Borwein and
R.~E.~Crandall, \emph{Effective Laguerre asymptotics}, dated
18 January 2007, states Conjecture~1 on printed page~38:
for $j\geq1$ and $a\in\{-1/2,-3/2,\ldots,-j+1/2\}$,
the asymptotic coefficient $C_j(a,z)$ is a positive power of $z$
times a polynomial over $\mathbb Q$, with total $z$-degree $3j/2$.
The source is
\url{https://carmamaths.org/resources/jon/laguerre-talk.pdf}.
The final article is \emph{SIAM J. Numer. Anal.} \textbf{46}
(2008), 3285--3312, DOI~\url{https://doi.org/10.1137/07068031X};
the authors' 28-page version is
\url{https://carmamaths.org/jon/Laguerre.pdf}.
The final version omits the numbered conjecture and describes the
coefficient structure as a topic for further research.
Accordingly, the result below proves the precise conjecture in the
2007 preprint; it is not presented as evidence that the final published
article still labels this statement open.  No separate subsequent proof
was located in the searches accompanying this note.

Set $m=n+1$.  We use exactly the normalization in the source:
\begin{equation}\label{p8:eq:bbc-normalization}
 L_n^{(-a)}(-z)\sim
 \frac{e^{-z/2}e^{2\sqrt{mz}}}
 {2\sqrt\pi\,z^{1/4-a/2}m^{1/4+a/2}}
 \sum_{j\geq0}C_j(a,z)m^{-j/2},\qquad C_0=1.
\end{equation}
The branch is the principal one on $z\notin(-\infty,0]$.
This normalization matters: replacing $m$ by $n$ changes the coefficients.

\begin{theorem}[Stronger half-integral vanishing]
\label{p8:thm:bbc-half-integral}
Let $j\geq1$ and $0\leq r<j$ be integers.  Let $d$ be the least
integer strictly larger than $r$ with $d\equiv j\pmod2$.
Then
\begin{equation}\label{p8:eq:bbc-half-factor}
 C_j(-r-1/2,z)=z^{d/2}P_{j,r}(z),\qquad
 P_{j,r}\in\mathbb Q[z],\qquad
 \deg P_{j,r}=\frac{3j-d}{2}.
\end{equation}
The coefficient of $z^{3j/2}$ in $C_j$ is $1/(12^j j!)$.
In particular every term has $z$-exponent strictly larger than $r/2$.
The assertion does not require $P_{j,r}(0)\ne0$.
\end{theorem}

\begin{proof}
We first express the finite coefficient formula in a form suitable for
tracking powers.  Define $A_{h,u}(a)$ by the formal generating function
\begin{equation}\label{p8:eq:bbc-A-generating}
 \sum_{h\geq0}\sum_{u=0}^h A_{h,u}(a)t^h y^u
 =(1+t/2)^{-a-1}
 \exp\!\left(
 y\sum_{\substack{\ell\geq1\\\ell\text{ odd}}}
 \frac{t^\ell}{2^{\ell+1}(\ell+2)}\right).
\end{equation}
All its coefficients belong to $\mathbb Q[a]$.
Each occurrence of $y$ is accompanied by a positive power of $t$,
so $A_{h,u}=0$ for $u>h$.
In the notation of equations~(39)--(44) of the final source,
\[
 \alpha_{h,u}(a,z)=A_{h,u}(a)z^{h/2+u}.
\]
This identity follows directly by expanding the exponential defining
their $A_h$: its first part is $-(a+1)\log(1+t/2)$,
and its second part has precisely the coefficients in
\eqref{p8:eq:bbc-A-generating}.  Thus the source's finite formula~(44) reads
\begin{equation}\label{p8:eq:bbc-finite}
 C_j(a,z)=\sum_{k=0}^j\sum_{u=0}^{j-k}
 \frac{A_{j-k,u}(a)}{16^k k!}
 z^{(j-2k)/2+u}
 \prod_{\ell=1}^k
 \bigl((2\ell-1)^2-4(a+j-k+2u)^2\bigr).
\end{equation}
Empty products have value one.  Formula~\eqref{p8:eq:bbc-finite}
is an exact finite identity for the uniquely normalized asymptotic
coefficients, not a numerical approximation.

Put $h=j-k$, $a=-r-1/2$, and $e=h+u-j/2$.
Consider a summand for which $e\leq r/2$.
Writing $v=h+2u$, this inequality gives
\[
 v\leq r+k.
\]
Also $r\leq j-1$ and $u\geq0$ give
\[
 v\geq r+1-k,
\]
since the latter is equivalent to $2u\geq r+1-j$.
Consequently
\[
 -k+\tfrac12\leq a+v\leq k-\tfrac12.
\]
The number $a+v$ is a half integer, so
$|a+v|=\ell-1/2$ for some $1\leq\ell\leq k$.
The corresponding factor in \eqref{p8:eq:bbc-finite} is zero.
(The case $k=0$ cannot satisfy $e\leq r/2$.)
Every surviving summand therefore has $e>r/2$.

All $2e$ are integers congruent to $j$ modulo two.  Hence all
surviving exponents belong to $d/2+\mathbb Z_{\geq0}$, proving
the factorization and rationality in \eqref{p8:eq:bbc-half-factor}.
Finally $h+u\leq2h\leq2j$, with equality only for
$h=u=j$ and $k=0$.  In \eqref{p8:eq:bbc-A-generating},
the coefficient of $t^j y^j$ is obtained by taking $j$ copies of
the term $ty/12$ in the exponential.  It is $1/(12^j j!)$.
Thus the highest exponent is exactly $3j/2$ and its coefficient
is as asserted.  This also proves the degree statement.
\end{proof}

\begin{corollary}[A coefficient divisibility pattern]
\label{p8:cor:bbc-divisibility}
Write
\[
 C_j(a,z)=\sum_{s=0}^{2j}c_{j,s}(a)z^{s-j/2},
 \qquad c_{j,s}\in\mathbb Q[a].
\]
For $0\leq s\leq j-1$, put $b=\max(0,2s-j)$.  Then
\[
 \prod_{r=b}^{j-1}(a+r+1/2)\quad\text{divides}\quad c_{j,s}(a)
 \quad\text{in }\mathbb Q[a].
\]
\end{corollary}
\begin{proof}
For every $b\leq r\leq j-1$ one has $s-j/2\leq r/2$.
The theorem gives $c_{j,s}(-r-1/2)=0$.
These roots are distinct, so their product divides $c_{j,s}$.
\end{proof}

\paragraph{Interpretation.}
The cancellation is termwise in the Bessel expansion: at half-integral
order, its elementary finite product vanishes.  No cancellation between
large symbolic expressions, and no asymptotic error estimate beyond the
source's established coefficient formula, is needed.
The divisibility corollary is a consequence of the same theorem, not a
second independent resolution of a previously open problem.
The connection to Euler's constant is through Laguerre expansions of
$\Gamma(0,z)=E_1(z)$ and the asymptotic methods underlying rational
approximants; this result by itself gives no irrationality statement.

\section{An inverse-linearization conjecture for Bessel polynomials}
\label{p8:sec:atia-recurrence}

\paragraph{Published question.}
M.~J.~Atia, \emph{On the Inverse of the Linearization Coefficients of
Bessel Polynomials}, \emph{Symmetry} \textbf{16} (2024), 737,
\url{https://doi.org/10.3390/sym16060737}, states an unnumbered
conjecture in Section~5.4, printed page~16, displayed equation~(24).
It asserts the recurrence proved below.  A distinct conjecture on the
same page asserts positivity of the inverse matrix outside $[0,1]$;
we do not claim that positivity statement here.

Define the normalized reverse Bessel polynomials by
\[
 q_n(u)=\sum_{k=0}^n
 \frac{n!(2n-k)!\,2^k}{(n-k)!(2n)!k!}u^k,
 \qquad q_0(u)=1.
\]
For $N\geq0$ let $M_N(a)$ be the matrix defined by
\begin{equation}\label{p8:eq:atia-M-definition}
 q_{N-j}(au)q_j((1-a)u)
 =\sum_{i=0}^N (M_N(a))_{j,i}q_i(u),\quad 0\leq j\leq N.
\end{equation}
Atia denotes this matrix by $M_N^\beta$.
Put $G_N=M_N^{-1}$ when it exists, and follow the source's convention
\[
 \gamma_k^{(N-m,m)}(a)=(G_N(a))_{m,k}.
\]
In particular the superscript of $\gamma$ indexes a row of the inverse,
not a row of the original matrix.

\begin{lemma}[A polynomial intertwiner]
\label{p8:lem:atia-intertwiner}
Put $T_N=2N+u-u\frac{d}{du}$ and define the $N\times(N+1)$ matrix
$A_N$ by
\[
 (A_N)_{j,j}=2N-2j-1,\qquad
 (A_N)_{j,j+1}=2j+1\quad(0\leq j<N),
\]
with all other entries zero.  Then
\begin{equation}\label{p8:eq:atia-intertwining}
 A_NM_N=M_{N-1}A_N.
\end{equation}
\end{lemma}
\begin{proof}
Comparison of coefficients in the finite formula for $q_j$ gives
\begin{equation}\label{p8:eq:atia-raising}
 (2j+1)q_{j+1}(u)=(2j+1+u)q_j(u)-u q_j'(u),
 \qquad j\geq0.
\end{equation}
Consequently
\[
 T_Nq_j=(2j+1)q_{j+1}+(2N-2j-1)q_j.
\]
Set $B_{N,j}(u)=q_{N-j}(au)q_j((1-a)u)$.
Apply \eqref{p8:eq:atia-raising} to each factor and use the product rule.
Since $a+(1-a)=1$, this gives
\[
 T_NB_{N-1,j}=(2j+1)B_{N,j+1}+(2N-2j-1)B_{N,j}.
\]
Computing $T_N(B_{N-1,j})$ first in the $q$ basis and then in the
$B$ basis proves \eqref{p8:eq:atia-intertwining}.
\end{proof}

\begin{proposition}[Spectrum and invertibility]
\label{p8:prop:atia-spectrum}
For every $N\geq0$,
\[
 \det(\lambda I-M_N(a))
 =\prod_{j=0}^N\bigl(\lambda-(1-2a)^j\bigr).
\]
In particular, for $N\geq1$, the matrix is invertible exactly when
$a\ne1/2$, and
$\det M_N(a)=(1-2a)^{N(N+1)/2}$.
\end{proposition}
\begin{proof}
The matrix $A_N$ has rank $N$.  Define a linear functional on polynomials
of degree at most $N$ by
\[
 \ell_N(u^k)=(-1)^k\frac{(2N)!}{(2N-k)!}\quad(0\leq k\leq N).
\]
For $k<N$,
\[
 \ell_N(T_Nu^k)=(2N-k)\ell_N(u^k)+\ell_N(u^{k+1})=0.
\]
Thus the nonzero column $w=(\ell_N(q_0),\ldots,\ell_N(q_N))^T$
spans $\ker A_N$; its first entry is one.
By \eqref{p8:eq:atia-intertwining}, $\ker A_N$ is $M_N$ invariant.
Its eigenvalue is the first entry of $M_Nw$, namely
\[
 \ell_N(q_N(au))
 =\sum_{k=0}^N(-2a)^k\binom Nk=(1-2a)^N.
\]
The map $A_N$ identifies the induced action on the quotient
$\mathbb C^{N+1}/\ker A_N$ with $M_{N-1}$.
The characteristic polynomial therefore acquires precisely the displayed
additional linear factor at each step.  Starting with $M_0=(1)$ proves
the assertion, including cases of coincident eigenvalues.
\end{proof}

\begin{theorem}[Atia's equation~(24)]
\label{p8:thm:atia-recurrence}
Let $n,m\geq1$, $1\leq k\leq n+m$, and $a\ne1/2$.
Then
\begin{align*}
 &(2m-1)\gamma_k^{(n,m)}(a)
 +(2n+1)\gamma_k^{(n+1,m-1)}(a)\\
 &\qquad=(2k-1)\gamma_{k-1}^{(n,m-1)}(a)
 +(2n+2m-2k-1)\gamma_k^{(n,m-1)}(a),
\end{align*}
where inverse-matrix entries outside their column range are zero.
\end{theorem}
\begin{proof}
Multiplying \eqref{p8:eq:atia-intertwining} on the left by $G_{N-1}$
and on the right by $G_N$ gives
$G_{N-1}A_N=A_NG_N$.
Take $N=n+m$ and its entry with row $m-1$ and column $k$.
The two nonzero entries in each relevant row or column of $A_N$
give exactly the stated identity.
\end{proof}

\paragraph{Status and counting.}
This proves one explicitly conjectured inverse recurrence.
The spectrum proposition independently justifies a result already stated
as Proposition~2 in Atia's article; it is not counted as a second open
problem.  The displayed proof there attempts to infer nesting of spectra
from a trace formula, which alone does not imply that nesting.
The intertwining identity supplies the missing invariant-subspace
argument.  We use the defining inverse matrix throughout; the signs of
some printed small inverse matrices in the source are inconsistent with
that definition.

% Standalone fragment, no preamble. Labels ntf:.
\section[Congruence obstructions for balanced lattices]{A congruence obstruction to inferring balanced lattices from smooth coefficients}
\label{p8:ntf:section}

Elsner's \emph{Euler's Constant and Series Transformations with Bases of
Balanced Lattices}, available from the author's website, formulates a
specific conjecture as Lemma~13.  It concerns coefficients coming from
truncated Ser series, not arbitrary positive integers.  With $N>e^{81000}$,
$32N/51<n<16N/25$, and the associated graph lattice $\Lambda$, the
conjecture asks for $N+1$ independent vectors, ordered by their maximum
norms $T_1\leq\cdots\leq T_{N+1}$, such that
\begin{equation}\label{p8:ntf:elsner-goal}
 T_{N+1}^{N/2}<\prod_{j=1}^{N+1}T_j
              <N^{2\varepsilon N/5}\det\Lambda.
\end{equation}
The author applies this with $\varepsilon=1/50000$.  Lemma~14 of the
same source proves three elementary size and ratio estimates for the
special coefficients, presented as possible preliminary input toward
Lemma~13.  We show that those estimates, even together with primitivity,
are insufficient: there are explicit integer coefficient families
satisfying every one of them for which \eqref{p8:ntf:elsner-goal} is impossible.
This does not disprove Lemma~13 for its stipulated Ser coefficients.

\subsection{A general obstruction from near divisibility}

For $a=(a_0,\ldots,a_N)\in\mathbb Z^{N+1}$, put
\[
 \Lambda(a)=\{(x_0,\ldots,x_N,y)\in\mathbb Z^{N+2}:
                  a\cdot x=y\}.
\]
This graph lattice has rank $N+1$ and Euclidean covolume
\begin{equation}\label{p8:ntf:determinant}
 \det\Lambda(a)=\sqrt{1+\sum_{j=0}^N a_j^2}.
\end{equation}
Indeed its standard basis has Gram matrix $I+aa^{\mathsf T}$,
whose determinant is $1+a^{\mathsf T}a$.  All successive minima below
refer to the maximum norm restricted to its real span.

\begin{lemma}[Near-divisibility forces a large last minimum]
\label{p8:ntf:congruence}
Suppose $a_j=Bb_j+r_j$, where $B$ is a positive integer,
$b,r\in\mathbb Z^{N+1}$ and $b\ne0$.  Set
$R=1+\sum_{j=0}^N|r_j|$.  Then
\[
 \lambda_{N+1}(\Lambda(a))\geq B/R.
\]
Consequently a necessary condition for \eqref{p8:ntf:elsner-goal} is
\begin{equation}\label{p8:ntf:necessary}
 (B/R)^{N/2}<N^{2\varepsilon N/5}\det\Lambda(a).
\end{equation}
\end{lemma}
\begin{proof}
If $(x,y)\in\Lambda(a)$ has maximum norm at most $H<B/R$, then
\[
 B(b\cdot x)=y-r\cdot x,\qquad |y-r\cdot x|\leq RH<B.
\]
The left side is an integer multiple of $B$, so $b\cdot x=0$.
Thus all such short vectors belong to the proper codimension-one
subspace $b\cdot x=0$ of the graph space.  They cannot span rank
$N+1$.  This proves the minimum bound.  Any $N+1$ independent
lattice vectors have largest norm at least this last minimum,
which together with \eqref{p8:ntf:elsner-goal} gives
\eqref{p8:ntf:necessary}.  Allowing a full-rank sublattice does not change
this argument: its independent vectors are still in $\Lambda(a)$.
\end{proof}

\subsection{Explicit counterexamples satisfying all three preparatory estimates}

\begin{theorem}[Smooth coefficient ratios do not imply balance]
\label{p8:ntf:counterfamily}
For every integer $N>e^{81000}$ define
\[
 B=2^N,\qquad K=2N,\qquad
 A=\left\lceil2N\left(\log N+\frac35\right)\right\rceil,
 \qquad n=\left\lfloor\frac{63N}{100}\right\rfloor,
\]
and take
\begin{equation}\label{p8:ntf:coefficients}
 a_0=BK,\qquad a_\nu=B(A+\nu)+1\quad(1\leq\nu\leq N).
\end{equation}
These positive integers have greatest common divisor one and satisfy
$32N/51<n<16N/25$, as well as all the numerical estimates in
Elsner's Lemma~14:
\begin{align}
 a_0<a_1<\cdots<a_N&<n!\,(2^{8409})^n,
                                      \label{p8:ntf:size}\\
 0.45+\log N<\frac{a_\nu}{a_0}
                &<1.23+\log N\qquad(1\leq\nu\leq N),
                                      \label{p8:ntf:ratio}\\
 1+\frac{0.499}{N\log N}<\frac{a_{\nu+1}}{a_\nu}
                &<1+\frac{3.390}{N\log N}
                         \qquad(1\leq\nu<N).
                                      \label{p8:ntf:successive-ratio}
\end{align}
Nevertheless
\begin{equation}\label{p8:ntf:minima}
 \lambda_1(\Lambda(a))=1,\qquad
 \lambda_{N+1}(\Lambda(a))\geq\frac{2^N}{N+1}.
\end{equation}
For any fixed $\varepsilon\geq0$, all sufficiently large $N$ in
this family admit no vectors satisfying \eqref{p8:ntf:elsner-goal}.
For the value $\varepsilon=1/50000$ used by Elsner, this failure
holds for every $N>e^{81000}$.
\end{theorem}
\begin{proof}
First, $n\leq0.63N<0.64N=16N/25$.  Also
\[
 0.63-32/51=13/5100,
\]
so $n\geq0.63N-1>32N/51$ for the stated range of $N$.
Successive coefficients differ by $B$.  Any common divisor of all
$a_j$ therefore divides $B$, and then divides
$a_1-B(A+1)=1$.  Thus their greatest common divisor is one.
The ordering follows immediately from the formula and $A>K$.

The elementary ceiling bounds give
\[
 \log N+\frac35+\frac{\nu}{2N}
       <\frac{a_\nu}{a_0}
       <\log N+\frac35+\frac{\nu+1}{2N}+\frac1{2NB}.
\]
The first expression is greater than $\log N+0.45$ and the last
is at most $\log N+1.1+1/N<\log N+1.23$.  This proves
\eqref{p8:ntf:ratio}.  Next,
\[
 \frac{a_{\nu+1}}{a_\nu}-1=\frac1{A+\nu+1/B}.
\]
The denominator is greater than $2N\log N$, giving an upper
bound $1/(2N\log N)<3.390/(N\log N)$.  It is also less than
$2N\log N+2.2N+2$.  Since $\log N>81000$, we have
\[
 0.499(2N\log N+2.2N+2)<N\log N,
\]
which proves the lower bound in \eqref{p8:ntf:successive-ratio}.

For completeness, $a_N<4BN\log N$ in this range.  Also
$n>3N/5$, and hence
\[
 (2^{8409})^n>2^{5000N}>4\,2^N N\log N>a_N.
\]
This proves the stronger version of \eqref{p8:ntf:size} with the
factorial omitted.

Apply Lemma~\ref{p8:ntf:congruence} with
$b=(K,A+1,\ldots,A+N)$ and $r=(0,1,\ldots,1)$.
Here $R=N+1$, yielding the second bound in \eqref{p8:ntf:minima}.
For the first, the vector with coordinates $x_1=1,x_2=-1,x_3=-1,
x_4=1$, all other $x_j=0$, and $y=0$, belongs to $\Lambda(a)$
by \eqref{p8:ntf:coefficients}.  It has maximum norm one, which is
the smallest possible norm of a nonzero integral vector.

Using \eqref{p8:ntf:determinant} and $a_N<4BN\log N$, we have
\begin{equation}\label{p8:ntf:det-upper}
 \det\Lambda(a)<5B N^{3/2}\log N.
\end{equation}
For fixed $\varepsilon$, the logarithm of the ratio between the
left side of \eqref{p8:ntf:necessary} and the upper bound for its
right side is at least
\[
 \left(\frac{N^2}2-N\right)\log2
 -\frac N2\log(N+1)-\frac{2\varepsilon N}5\log N
 -\frac32\log N-\log\log N-\log5.
\]
It tends to positive infinity, proving eventual impossibility.
For $\varepsilon=1/50000$ and $N>e^{81000}$ one can bound
$\log(N+1)<N/10$ and $\log N<N/10$.  The positive term is
at least $N^2/8$ since $N\geq4$ and $\log2>1/2$.
The two negative terms proportional to $N$ are at most
$N^2/20+N^2/1250000$, respectively.  The remaining three
negative terms are less than $N^2/100$ in this range.
The displayed logarithm is therefore strictly positive.
Condition \eqref{p8:ntf:necessary} fails, completing the proof.
\end{proof}

The obstruction is arithmetic: the coefficients have nearly equal
residues modulo $2^N$.  Their archimedean ratios conceal this common
modular structure.  The construction does not give the Ser-series
coefficients from Elsner's equation (2.5.4), and therefore is neither
a counterexample to his actual conjecture nor an irrationality proof
for $\gamma$.  It shows rigorously that a proof using only
Lemma~14's bounds and primitivity cannot establish that conjecture.
Additional information on the coefficients modulo large integers
is indispensable for such an argument.

\subsection{An elementary answer to a stated numerical observation}

On page 15 of the September 2024 version of
\url{https://arxiv.org/pdf/2210.06768}, Murabayashi and Yoshida
state that computation suggests divergence of
$|k!q_{2k}\delta-k!p_{2k}|$.  The assertion can be proved directly.
It is distinct from their Conjecture~7.5 and is not presented here
as a new result about classical Laguerre asymptotics.

\begin{proposition}[The factorially cleared errors diverge]
\label{p8:ntf:laguerre-error}
Let $\delta=\int_0^\infty e^{-t}(1+t)^{-1}\,dt$, and define the
numerators and denominators in the normalization of
Murabayashi--Yoshida by
\[
 p_{-1}=1,\quad p_0=0,\quad q_{-1}=0,\quad q_0=1,
 \qquad v_j=c_jv_{j-1}+v_{j-2},
\]
for $v=p,q$, where $c_{2k-1}=1$ and $c_{2k}=1/k$.
Then for $k\geq0$,
\begin{equation}\label{p8:ntf:even-error}
 \delta q_{2k}-p_{2k}=J_k
      :=\int_0^\infty\frac{e^{-t}t^k}{(1+t)^{k+1}}\,dt>0,
\end{equation}
and for $k\geq1$,
\begin{equation}\label{p8:ntf:odd-error}
 \delta q_{2k-1}-p_{2k-1}
     =-k\int_0^\infty\frac{e^{-t}t^{k-1}}{(1+t)^{k+1}}\,dt<0.
\end{equation}
In particular both
\[
 k!\,|\delta q_{2k}-p_{2k}|\longrightarrow\infty,
 \qquad
 (k-1)!\,|\delta q_{2k-1}-p_{2k-1}|\longrightarrow\infty.
\]
\end{proposition}
\begin{proof}
The recurrence is the specialization $x=1$ of Section 3 in the
cited paper.  Eliminating the odd indices shows that both even
subsequences satisfy
\[
 (k+1)v_{2k+2}=(2k+2)v_{2k}-kv_{2k-2}\qquad(k\geq1).
\]
For the proposed integral sequence, differentiation gives
\[
 \frac{d}{dt}\left(\frac{e^{-t}t^k}{(1+t)^{k+1}}\right)
   =\frac{e^{-t}t^{k-1}}{(1+t)^{k+2}}(k-2t-t^2).
\]
Its integral is zero for $k\geq1$, which is exactly
$(k+1)J_{k+1}-(2k+2)J_k+kJ_{k-1}=0$.
The initial values are $J_0=\delta$ and $J_1=2\delta-1$.
For the second, integrate the derivative of $e^{-t}/(1+t)$
to obtain $\int_0^\infty e^{-t}(1+t)^{-2}\,dt=1-\delta$.
These values equal $\delta q_0-p_0$ and $\delta q_2-p_2$.
Uniqueness of the recurrence proves \eqref{p8:ntf:even-error}.
The relation $v_{2k-1}=k(v_{2k}-v_{2k-2})$ now proves
\eqref{p8:ntf:odd-error} by subtracting the two integrands.

Finally restrict the even integral to
$\sqrt{k}\leq t\leq\sqrt{k}+1$.  The inequality
$\log(1+u)\leq u$ for $u>0$ gives
\[
 e^{-t}\left(\frac{t}{1+t}\right)^k
       \geq e^{-2\sqrt{k}-1},\qquad
 J_k\geq\frac{e^{-2\sqrt{k}-1}}{\sqrt{k}+2}.
\]
For the odd integral, the additional factor is $k/t$, giving
\[
 |\delta q_{2k-1}-p_{2k-1}|
  \geq\frac{k e^{-2\sqrt{k}-1}}
               {(\sqrt{k}+1)(\sqrt{k}+2)}.
\]
Multiplication by $k!$ and $(k-1)!$, respectively, proves the two
limits, for example using $k!\geq(k/2)^{k/2}$ for $k\geq2$.
\end{proof}

This verifies the paper's stated numerical divergence and explains
why its direct factorial denominator clearing cannot prove
irrationality.  It gives no monotone subsequence of the fractional-part
errors required in Conjecture~7.5.  The error integrals are classical
Laguerre--Stieltjes objects, so this proof should not be counted as
a newly resolved longstanding conjecture.

\subsection{External targets and current proof status}

The following source-defined questions were inspected without
counting proved results or routine reformulations as new closures.
\begin{enumerate}
\item Sondow, \emph{Criteria for irrationality of Euler's constant},
Proc. Amer. Math. Soc.\ 131 (2003), 3335--3344;
\url{https://arxiv.org/pdf/math/0209070}, unnumbered observation
following Table~1 (preprint pp.\ 4--5): the fractional parts
$\{\log S_n\}$ appear dense in $(0,1)$, with Ces\`aro average
$1/2$.  Neither assertion is proved here.  The published rationality
criteria are already theorems and cannot be counted as conjectures
resolved by this investigation.
\item Elsner, \emph{Euler's Constant and Series Transformations
with Bases of Balanced Lattices}, author manuscript,
\url{https://www.continuedfractions.net/Euler_s_Constant.pdf},
Lemma~13 and equation (2.5.6).  The special-coefficient conjecture
remains unresolved here.  Theorem~\ref{p8:ntf:counterfamily} refutes
only its proposed extension to all coefficient families satisfying
Lemma~14's preparatory estimates.
\item Murabayashi--Yoshida,
\url{https://arxiv.org/pdf/2210.06768}, Conjecture~7.5:
the even sequence $\delta\{q_{2k}\}-\{p_{2k}\}$ has an infinite
nondecreasing subsequence, and the odd sequence
$\delta\{q_{2k-1}\}-\{p_{2k-1}\}$ has an infinite nonincreasing
subsequence, with the approximants normalized as in that paper.
Here $\delta$ is the Euler--Gompertz constant.  The paper's
Theorem~7.4 makes either stated type of subsequence sufficient for
irrationality of $\delta$; this conjecture is not proved here.
\item Dixit--Murty, \emph{On Ihara's conjectures for Euler--Kronecker
constants}, Acta Arith.\ 210 (2023), 95--123;
\url{https://www.imsc.res.in/~anupdixit/Ihara-final-revised.pdf},
equation (15) and the following sentence ask for an unconditional
proof of
\[
 \sum_{\substack{q\leq Q\\q\ \mathrm{prime}}}
       \frac{|\gamma_q|}{q-1}=\log Q+o(\log Q).
\]
The source proves it assuming Elliott--Halberstam.  No
unconditional proof is supplied here.
\end{enumerate}

The sources establish the precise historical targets.  The status
statements describe the retrieved versions and this investigation;
a comprehensive proof of absence of all subsequent unpublished work
is not asserted.  The denominator conjectures of Rivoal proved by
Hessami Pilehrood--Hessami Pilehrood in
\url{https://arxiv.org/abs/1004.0578} are explicitly excluded from
the list of unresolved targets.

% Research fragment: no preamble.  Partial results, not claimed new in the literature.
\section{A restricted division theorem for rational Borel transforms}
\label{p8:sf:section}

\paragraph{Published targets and accounting.}
The following four targets were checked against their primary sources.
The first three occur in Fischler--Rivoal,
\emph{Relations between values of arithmetic Gevrey series, and
applications to values of the Gamma function}, J. Number Theory
261 (2024), 36--54,
\url{https://arxiv.org/abs/2301.13518}.
Conjecture~2 asks whether a vanishing non-Stokes sum at $1$ of an
arithmetic Gevrey series of order $1$ permits division by $z-1$
within the same class. Conjecture~3 asks the corresponding question
for a mixed function $F(z)+\mathfrak f(1/z)$, with $F$ an $E$-function.
Section~2.2 asks about membership of $\Gamma^{(s)}(a)$ in
$\mathbf M_G=\mathbf E+\mathbf D$; the Euler-constant case is
$-\Gamma'(1)=\gamma$.
The fourth target is Conjecture~1.2 of Fischler--Rivoal,
\emph{The ring of values of $E$-functions is not a field},
12 October 2026, \url{https://arxiv.org/abs/2610.16603}:
are the units of $\mathbf E$ precisely
$\alpha e^\beta$, with $\alpha,\beta\in\overline{\mathbb Q}$
and $\alpha\ne0$?
No complete resolution of these four targets was located or proved
in this investigation. The theorem below proves a restricted case
of the two division conjectures; it contributes \emph{zero} to the
count of completely solved published open problems. Its priority
in the literature has not been established.

\subsection{The precise restricted class}
Put $k=\overline{\mathbb Q}\cap\mathbb R$. For a rational function
$B(t)=\sum_{n\ge0}b_nt^n\in k(t)$ regular at zero, write
\[
 \mathcal L_B(u)=\sum_{n\ge0}n!b_nu^n.
\]
This is an arithmetic Gevrey series of order $1$: its Borel
transform is the rational $G$-function $B$. Fix $0<\varphi<\pi/2$.
If every finite pole of $B$ is positive real, the two lateral
Borel--Laplace values at $u=1$ are
\[
 I_+(B)=\int_{0}^{e^{i\varphi}\infty}e^{-t}B(t)\,dt,
 \qquad
 I_-(B)=\int_{0}^{e^{-i\varphi}\infty}e^{-t}B(t)\,dt.
\]
They are independent of $\varphi$ on the indicated side.
The polynomial growth of $B$ and exponential decay on those rays
give absolute convergence. These integrals are the two lateral
$1$-sums, up to the harmless exchange of the side labels resulting
from the reciprocal-variable convention.

\begin{theorem}[Reality forces a rational telescoper]
\label{p8:sf:real-sum}
Suppose $B\in k(t)$ is regular at $0$ and all its finite poles
belong to $(0,\infty)$. The following conditions are equivalent:
\begin{enumerate}
\item One of $I_+(B),I_-(B)$ is real.
\item Both lateral values are equal.
\item There is an $S\in k(t)$, regular at $0$, such that
\begin{equation}\label{p8:sf:telescoper}
 B=S-S'.
\end{equation}
\end{enumerate}
When these conditions hold, $S$ is unique, its finite poles are
among those of $B$, and
\[
 I_+(B)=I_-(B)=S(0)\in k.
\]
In particular the Borel transform of
\begin{equation}\label{p8:sf:quotient}
 \frac{\mathcal L_B(u)-S(0)}{u-1}
\end{equation}
is the rational function $S'(t)$. Hence the quotient is again an
arithmetic Gevrey series of order $1$.
\end{theorem}

\begin{proof}
Real coefficients give $I_-(B)=\overline{I_+(B)}$, so conditions
1 and 2 are equivalent. Close the sector between the two rays.
The connecting arcs at infinity contribute zero, and the residue
theorem gives, with an immaterial orientation sign,
\[
 I_+(B)-I_-(B)
   =\pm2\pi i\sum_{a}\operatorname{Res}_{t=a}(e^{-t}B(t)).
\]
Here the sum runs through the distinct positive algebraic poles.
For each pole there is a $c_a\in k$ such that
\[
 \operatorname{Res}_{t=a}(e^{-t}B(t))=e^{-a}c_a.
\]
Indeed, if the principal part is
$\sum_{j=1}^{m}b_{a,j}(t-a)^{-j}$, then
\[
 c_a=\sum_{j=1}^{m}\frac{(-1)^{j-1}}{(j-1)!}b_{a,j}.
\]
Under condition 2 the Lindemann--Weierstrass theorem, applied to
the distinct algebraic exponents $-a$, implies $c_a=0$ for every
pole separately.

We give the rational integration step explicitly. On principal
parts at a fixed pole of order at most $m$, the map
$S\mapsto S-S'$ sends principal parts of order at most $m-1$
injectively into the $m$-dimensional space of principal parts of
order at most $m$. Its image is contained in the kernel of the
nonzero linear functional
$B\mapsto\operatorname{Res}_{t=a}(e^{-t}B)$, because
$e^{-t}(S-S')=-(e^{-t}S)'$. Both the image and this kernel have
dimension $m-1$, and therefore they coincide. This argument also
applies to $m=1$, where both spaces are zero. All matrices are
defined over $k$, after removing the common nonzero factor
$e^{-a}$ from the residue functional.
Thus the vanishing of every $c_a$ produces the required principal
part of $S$ at every pole. The remaining polynomial part is also
in the image: on polynomials, $(1-d/dt)^{-1}$ is the finite sum
$\sum_{j\ge0}(d/dt)^j$. This proves existence of $S\in k(t)$
with no additional finite poles. In particular $S$ is regular at
zero. Uniqueness follows since a rational solution of $S'=S$ is
zero: a nonzero finite pole would increase its order after
differentiation, and a nonzero polynomial cannot equal its
derivative.

Conversely, \eqref{p8:sf:telescoper} allows integration along either
ray without crossing a pole. Since $e^{-t}S(t)\to0$ at infinity,
\[
 I_\pm(B)=-[e^{-t}S(t)]_0^{e^{\pm i\varphi}\infty}=S(0).
\]
This proves all equivalences and the value formula.

Finally write $S(t)=\sum_{n\ge0}s_nt^n$ and
$\mathcal L_S(u)=\sum_{n\ge0}n!s_nu^n$.
Taking coefficients in $B=S-S'$ gives the exact formal identity
\[
 \mathcal L_B(u)-s_0
 =\frac{u-1}{u}\bigl(\mathcal L_S(u)-s_0\bigr).
\]
Consequently \eqref{p8:sf:quotient} equals
$\sum_{n\ge0}(n+1)!s_{n+1}u^n$, whose Borel transform is
$\sum_{n\ge0}(n+1)s_{n+1}t^n=S'(t)$. Rational Borel transforms
give the claimed arithmetic Gevrey property.
\end{proof}

\begin{corollary}[A restricted case of Conjecture 2]
\label{p8:sf:conj2}
Under the hypotheses of Theorem~\ref{p8:sf:real-sum}, if a lateral
$1$-sum of $\mathcal L_B$ vanishes at $1$, then
$\mathcal L_B(u)/(u-1)$ is an arithmetic Gevrey series of order
$1$ with rational Borel transform.
\end{corollary}

The hypothesis is not vacuous. Take
\[
 S(t)=\frac{t}{t-1},\qquad
 B(t)=\frac{t}{t-1}+\frac1{(t-1)^2}.
\]
Then $S(0)=0$ and
\[
 \mathcal L_B(u)=1+\sum_{n\ge1}n\,n!u^n,
 \qquad
 \frac{\mathcal L_B(u)}{u-1}
      =-\sum_{n\ge0}(n+1)!u^n.
\]
Both lateral values of the first, genuinely divergent series are
zero. This example follows from the telescoper and is not an
independent solved open problem.

\begin{corollary}[A restricted case of Conjecture 3]
\label{p8:sf:conj3}
Let $F$ be an $E$-function with real algebraic Taylor coefficients,
and let $B$ satisfy the hypotheses of Theorem~\ref{p8:sf:real-sum}.
Set
\[
 \Psi(z)=F(z)+\mathcal L_B(1/z).
\]
If a lateral sum of $\Psi$ vanishes at $z=1$, then
$\Psi(z)/(z-1)$ is a mixed function.
\end{corollary}
\begin{proof}
Since $F(1)$ is real, the assumed relation makes the lateral value
of $\mathcal L_B$ real. Theorem~\ref{p8:sf:real-sum} gives a rational
telescoper $S$ and $a=S(0)=-F(1)\in k$. The established
Andr\'e--Beukers division theorem makes
$(F(z)+a)/(z-1)$ an $E$-function.
Put $h(u)=(\mathcal L_B(u)-a)/(u-1)$, which is an arithmetic
Gevrey series of order $1$ by Theorem~\ref{p8:sf:real-sum}. Then
\[
 \frac{\Psi(z)}{z-1}
   =\frac{F(z)+a}{z-1}-\frac1z h(1/z),
\]
which is the required mixed decomposition.
\end{proof}

\subsection{Exact limitation at the Gompertz example}
The positive-pole assumption is used at the decisive step: the
upper and lower contours enclose every finite pole. For the
Gompertz series
\[
 \sum_{n\ge0}(-1)^nn!u^n,
 \qquad B(t)=\frac1{1+t},
\]
the only pole is $-1$. Its positive-direction sum is already real,
but the contour difference sees no pole. The residue argument
therefore supplies no rational telescoper and no arithmetic
conclusion about its value $\delta$. Allowing negative poles would
require a genuinely new input, rather than a change of wording in
Theorem~\ref{p8:sf:real-sum}.

The prior positivity certificates establish only explicit positive
linear forms in $\gamma$. The prior cancellation-locus results
control which parameters can be identified from specified local
functional data. Neither supplies the missing arithmetic division
theorem for the unrestricted Borel class, proves membership or
nonmembership of $\gamma$ in $\mathbf E+\mathbf D$, or classifies
all units of $\mathbf E$.

% Suggested bibliography entries if not already present:
% FischlerRivoal2024Gevrey: DOI 10.1016/j.jnt.2024.02.016,
% https://arxiv.org/abs/2301.13518
% FR2026: https://arxiv.org/abs/2610.16603
% Ferguson2021: T. Ferguson, Algebraic properties of arithmetic
% Gevrey series of order 1, JNT229 (2021),168--178,
% DOI10.1016/j.jnt.2021.04.010.  Published title uses the Cyrillic
% character reverse-E rather than the descriptive phrase above.

\section{What remains on the route to irrationality}
\label{p8:sec:remaining}
The results reveal several distinct kinds of information. Positive
spectral kernels control signs of derivatives. Their fractional
derivatives detect noninteger complete-monotonicity degrees.
Moment recurrences control formal asymptotic expansions, and local
singularities can defeat a density approximation even when global
probabilistic hypotheses hold. Polynomial intertwiners turn a
special-function recurrence conjecture into an invariant-subspace
identity. None of these mechanisms, by itself, controls the arithmetic
denominators of approximations to $\gamma$.

A direct irrationality argument would be obtained from rational
approximants with a genuinely small nonzero cleared error: for example,
integers $A_n,B_n$, with $B_n\ne0$, satisfying
\[
 0<|B_n\gamma-A_n|\longrightarrow0.
\]
Indeed, if $\gamma=p/q$, every nonzero such error is at least $1/q$.
The analytic positivity results here do not construct this sequence.
An integral with a positive integrand is not enough when the integer
multiplier needed to clear denominators grows faster than its decay.
Proposition~\ref{p8:ntf:laguerre-error} supplies an exact example of this
failure for the neighboring Gompertz problem.

The division result has a different obstruction. Its decisive step is
that the two lateral contours together detect every pole, after which
Lindemann--Weierstrass separates the residues. The negative pole of the
Gompertz Borel transform is invisible to that contour difference.
Extending the theorem to such poles requires another arithmetic input;
the present proof cannot be continued merely by weakening a hypothesis.

Several concrete questions now remain within the methods developed here:
\begin{enumerate}
\item Determine the exact fractional degree of $R_2$, $-R_2'$, or
$R_2''$, or give explicit quantitative upper and lower bounds for its
fractional part. The existence and large-index limit proved here do not
supply an exact value.
\item Classify the degrees of $(-1)^mR_N^{(m)}$ for $N\ge2$, $m\ge3$.
The derivative and beta-transfer proof settles $m=0,1,2$ only.
\item Identify regularity conditions on a general L\'evy density that
are sufficient for the local tilted-density limit while excluding the
integrable cusp construction. Such a corrected theorem must be compared
with the regularity theorems already proved in the AKZ source.
\item Obtain arithmetic information on the actual Ser coefficients
strong enough to control the last successive minimum. The generic
counterfamily shows exactly why the preparatory ratio estimates do not
provide this information.
\item Extend the rational-Borel division argument to an arithmetic
class containing negative real poles without assuming the desired
irrationality or a stronger transcendence conjecture.
\end{enumerate}
These are explicitly new research questions or remaining source targets,
not proved conclusions. The manuscript therefore records real progress
on the neighboring conjectures together with a precise boundary on its
arithmetic implications. It does not prove irrationality or transcendence
of $\gamma$ or of the Euler--Gompertz constant $\delta$.

\PartAppendix
\section{Reproducible exact checks}
\label{p8:app:checks}
The proofs are symbolic and do not rely on numerical evidence. The
following exact-arithmetic checks were run as safeguards against a
normalization or transpose error. They are included in this standalone
source so that no external code file is required. With Python 3 and
SymPy, the first script verifies 24 size/parameter cases through $N=6$,
including the singular parameter $a=1/2$. The second verifies 36
half-integral cases through $j=8$ and the four $C_4$ polynomials printed
in the Laguerre source. All checks passed.

\subsection{Bessel intertwiner and inverse orientation}


\begin{lstlisting}
"""Exact normalization/orientation checks for the Bessel intertwiner proof."""
import sympy as s

u, lam = s.symbols('u lam')

def q(n):
    return sum(s.Rational(s.factorial(n)*s.factorial(2*n-k)*2**k,
                          s.factorial(n-k)*s.factorial(2*n)*s.factorial(k))*u**k
               for k in range(n+1))

def matrix(N, a):
    Q=s.Matrix([[q(i).coeff(u,k) for i in range(N+1)] for k in range(N+1)])
    P=s.Matrix([[s.Poly(q(N-j).subs(u,a*u)*q(j).subs(u,(1-a)*u),u).nth(k)
                 for j in range(N+1)] for k in range(N+1)])
    return (Q.inv()*P).T

checks=0
for a in [s.Integer(2),s.Integer(-1),s.Rational(2,3),s.Rational(1,2)]:
    previous=s.eye(1)
    for N in range(1,7):
        M=matrix(N,a)
        A=s.zeros(N,N+1)
        for j in range(N):
            A[j,j]=2*N-2*j-1
            A[j,j+1]=2*j+1
        assert A*M==previous*A
        assert s.expand(M.charpoly(lam).as_expr()-s.prod(lam-(1-2*a)**j
                         for j in range(N+1)))==0
        if a!=s.Rational(1,2):
            assert previous.inv()*A==A*M.inv()
        previous=M
        checks+=1
print(f'PASS: {checks} exact parameter/size cases, including the singular a=1/2 case.')
\end{lstlisting}

\subsection{Laguerre coefficients and half-integral powers}
\begin{lstlisting}
"""Exact checks of the finite BBC coefficient normalization; not used as proof."""
from math import factorial
import sympy as s

a, z, v = s.symbols('a z v')
# Work with x=sqrt(z) to avoid branch simplification issues.
x = s.symbols('x')
N = 8
A = [s.Integer(1)]
for h in range(1, N + 1):
    value = 0
    for ell in range(1, h + 1):
        b = (a + 1) * (-v*x/2)**ell
        b += s.Rational((1 - (-1)**ell)*ell, ell + 2) * (v*x/2)**(ell+2)
        value += A[h-ell] * b
    A.append(s.expand(value/h))

def C(j, avalue):
    value = 0
    for k in range(j+1):
        h = j-k
        for u in range(h+1):
            alpha = A[h].coeff(v, h+2*u).subs(a, avalue)
            g = s.prod((2*ell-1)**2 - 4*(avalue+h+2*u)**2
                       for ell in range(1,k+1))
            value += alpha*g/(16**k*factorial(k)*x**k)
    return s.expand(value)

for j in range(1, N+1):
    for r in range(j):
        value = C(j, -s.Rational(1,2)-r)
        poly = s.Poly(value,x)
        assert poly.degree() == 3*j
        assert poly.LC() == s.Rational(1,12**j*factorial(j))
        assert all(power[0] > r and (power[0]-j)%2 == 0
                   for power, coefficient in poly.terms())

expected = [
    z*(5*z**5-60*z**4-918*z**3+6804*z**2+16605*z-48600),
    z*(5*z**5+60*z**4-1038*z**3-6804*z**2+17253*z+29160),
    z**2*(5*z**4+180*z**3+882*z**2-10692*z-39771),
    z**2*(5*z**4+300*z**3+4842*z**2+14580*z-63747),
]
for r, numerator in enumerate(expected):
    assert s.expand(C(4,-s.Rational(1,2)-r)-numerator.subs(z,x*x)/2488320)==0
print('PASS: 36 exact half-integral cases through j=8; all four source C4 examples agree.')
\end{lstlisting}
\begin{thebibliography}{VIII.99}
\bibitem[VIII.1]{p8:MahmoudQi2022}
M. Mahmoud and F. Qi, Bounds for completely monotonic degrees of
remainders in asymptotic expansions of the digamma function,
\emph{Math. Inequal. Appl.} \textbf{25} (2022), 291--306.
\url{https://doi.org/10.7153/mia-2022-25-17}.

\bibitem[VIII.2]{p8:Qi2021}
F. Qi, Bounds for completely monotonic degree of a remainder for an
asymptotic expansion of the trigamma function (2021).
\url{https://doi.org/10.1080/25765299.2021.1962060};
author manuscript \url{https://arxiv.org/abs/2003.05300v2}.

\bibitem[VIII.3]{p8:Zhu2020}
L. Zhu, Completely monotonic integer degrees for a class of special
functions, \emph{AIMS Mathematics} (2020).
\url{https://doi.org/10.3934/math.2020224}.

\bibitem[VIII.4]{p8:QiGuo2017}
F. Qi and B.-N. Guo, Integral representations and complete monotonicity
of remainders of the Binet and Stirling formulas for the gamma function,
\emph{RACSAM} \textbf{111} (2017), 425--434.
\url{https://doi.org/10.1007/s13398-016-0302-6}.

\bibitem[VIII.5]{p8:DLMF}
NIST Digital Library of Mathematical Functions, Eq.~4.36.3,
partial-fraction expansion of $\coth$.
\url{https://dlmf.nist.gov/4.36.E3}.

\bibitem[VIII.6]{p8:DiblikMedina2018}
J. Dibl\'ik and R. Medina, Exact asymptotics of positive solutions to
Dickman equation, \emph{Discrete Contin. Dyn. Syst. B}
\textbf{23} (2018), 101--121.
\url{https://doi.org/10.3934/dcdsb.2018007}.

\bibitem[VIII.7]{p8:ArratiaKochmanZabell2016}
R. Arratia, F. Kochman and S. Zabell, Large deviation asymptotics for a
random variable with L\'evy measure supported by $[0,1]$,
arXiv:1606.03524v1 (2016).
\url{https://arxiv.org/abs/1606.03524v1}.

\bibitem[VIII.8]{p8:PinskyDickman2016}
R. G. Pinsky, On the strange domain of attraction to generalized Dickman
distributions for sums of independent random variables,
arXiv:1611.07207v3 (2017).
\url{https://arxiv.org/abs/1611.07207v3}.

\bibitem[VIII.9]{p8:BBC2007}
D. Borwein, J. M. Borwein and R. E. Crandall, Effective Laguerre
asymptotics, 18 January 2007 preprint, Conjecture~1, p.~38.
\url{https://carmamaths.org/resources/jon/laguerre-talk.pdf}.

\bibitem[VIII.10]{p8:BBC2008}
D. Borwein, J. M. Borwein and R. E. Crandall, Effective Laguerre
asymptotics, \emph{SIAM J. Numer. Anal.} \textbf{46} (2008), 3285--3312.
\url{https://doi.org/10.1137/07068031X}; author version
\url{https://carmamaths.org/jon/Laguerre.pdf}.

\bibitem[VIII.11]{p8:Atia2024}
M. J. Atia, On the Inverse of the Linearization Coefficients of Bessel
Polynomials, \emph{Symmetry} \textbf{16} (2024), 737.
\url{https://doi.org/10.3390/sym16060737}.

\bibitem[VIII.12]{p8:Elsner}
C. Elsner, Euler's Constant and Series Transformations with Bases of
Balanced Lattices, author manuscript.
\url{https://www.continuedfractions.net/Euler_s_Constant.pdf}.

\bibitem[VIII.13]{p8:MurabayashiYoshida}
N. Murabayashi and H. Yoshida, Rational function approximations of the
special function $e^xE_1(x)$ and applications to irrationality of
Euler--Gompertz constant $\delta$, arXiv:2210.06768v2,
22 September 2024.
\url{https://arxiv.org/abs/2210.06768v2}.

\bibitem[VIII.14]{p8:FischlerRivoal2024}
S. Fischler and T. Rivoal, Relations between values of arithmetic
Gevrey series, and applications to values of the Gamma function,
\emph{J. Number Theory} \textbf{261} (2024), 36--54.
\url{https://doi.org/10.1016/j.jnt.2024.02.016};
\url{https://arxiv.org/abs/2301.13518}.

\bibitem[VIII.15]{p8:FischlerRivoal2026}
S. Fischler and T. Rivoal, The ring of values of $E$-functions is not
a field, arXiv:2610.16603 (2026).
\url{https://arxiv.org/abs/2610.16603}.

\bibitem[VIII.16]{p8:Sondow2003}
J. Sondow, Criteria for irrationality of Euler's constant,
\emph{Proc. Amer. Math. Soc.} \textbf{131} (2003), 3335--3344.
\url{https://arxiv.org/abs/math/0209070}.

\bibitem[VIII.17]{p8:DixitMurty2023}
A. B. Dixit and M. Ram Murty, On Ihara's conjectures for Euler--Kronecker
constants, \emph{Acta Arith.} \textbf{210} (2023), 95--123.
\url{https://doi.org/10.4064/aa220729-19-3}.
\end{thebibliography}
% END PRESERVED SOURCE BODY 8
\phantomsection\label{part8:end}
\end{document}
